% Colección didáctica ampliable. Fuente editable en LaTeX.
\documentclass[11pt,a4paper,landscape]{article}
\usepackage[utf8]{inputenc}
\usepackage[T1]{fontenc}
\usepackage{lmodern,amsmath,amssymb,array,booktabs,multicol,graphicx,xcolor,colortbl,fancyhdr,hyperref}
\usepackage[margin=1.45cm]{geometry}
\hypersetup{colorlinks=true,linkcolor=black,urlcolor=blue!50!black,pdftitle={Colección 1.1 y 1.2 - Álgebra lineal}}
\pagestyle{fancy}\fancyhf{}
\fancyfoot[L]{\small Álgebra lineal | Colección didáctica | Secciones 1.1 y 1.2}
\fancyfoot[R]{\small\thepage}
\renewcommand{\headrulewidth}{0pt}
\setlength{\headheight}{0pt}\setlength{\parindent}{0pt}\setlength{\parskip}{4pt}
\setlength{\abovedisplayskip}{5pt}\setlength{\belowdisplayskip}{5pt}\setlength{\abovedisplayshortskip}{3pt}\setlength{\belowdisplayshortskip}{3pt}
\setlength{\columnsep}{1cm}\setlength{\multicolsep}{6pt}\raggedcolumns
\newcommand{\titulo}[1]{{\LARGE\bfseries #1}\par\smallskip\hrule\medskip}
\newcommand{\respuesta}[1]{\par\smallskip\fcolorbox{black!30}{black!3}{\parbox{\dimexpr\linewidth-2\fboxsep-2\fboxrule\relax}{\textbf{Resultado.} #1}}\par\smallskip}
\begin{document}
\pdfbookmark[0]{Inicio}{inicio}
\titulo{Álgebra lineal | Colección de ejercicios}
{\Large Secciones 1.1 y 1.2 | Sistemas, matrices y MATLAB}\par\bigskip
\begin{minipage}[t]{.48\textwidth}
\textbf{Una guía para aprender, no solamente para copiar respuestas}\par
Cada ejercicio conserva su número y su enunciado. Después aparecen la idea, las operaciones, la solución y su comprobación. Los problemas de aplicación comienzan con las variables y las unidades.\par
\textbf{Cómo estudiar}\par
1. Lee el enunciado y trata de identificar las incógnitas.\par
2. Sigue una operación por vez. Antes de continuar, intenta hacerla tú.\par
3. Comprueba la respuesta en todas las ecuaciones originales.\par
4. En los problemas, explica qué significan los números y sus unidades.\par
\textbf{Material recibido}\par
Se recibieron las figuras 1.2 y 1.3. Ambas se incluyen con sus explicaciones en los ejercicios 21 y 22, respectivamente. Todos los enunciados recibidos de este bloque están cubiertos.
\end{minipage}\hfill
\begin{minipage}[t]{.48\textwidth}
\textbf{Sección 1.1: sistemas lineales}\par
1--14. Sistemas por eliminación: solución única, ninguna o infinitas.\par
15--16. Parámetros y propiedades de las soluciones.\par
17--18. Sustitución directa sin eliminación.\par
19--20. Verificación de una solución propuesta.\par
21--22. Interpretación de intersecciones de planos.\par
23--26. Modelos de producción y alimentación.\par
27. Una parábola a partir de tres puntos.\par
28. Inversiones e intereses.\par
T.1--T.4. Demostraciones de equivalencia y sistemas homogéneos.\par\medskip\textbf{Sección 1.2: matrices}\par 1--15. Entradas, igualdad, operaciones, transposición y aritmética binaria.\par T.1--T.18. Demostraciones, conteos y conmutadores.\par ML.1--ML.5. Códigos MATLAB y laboratorios visuales.\par\medskip
\textbf{Colección ampliable}\par
Las secciones 1.1 y 1.2 conservan su numeración. La portada web permite abrir cada sección. Los siguientes bloques podrán añadirse conservando los anteriores; el código fuente separa los datos, los laboratorios y la presentación.
\end{minipage}
\newpage
\titulo{Antes de empezar | Cómo leer una solución}
\begin{multicols}{2}
\textbf{1. Una letra representa un número por descubrir}\par
En \(x+2y=8\), todavía no sabemos cuánto valen \(x\) e \(y\). La escritura \(2y\) significa \(2\cdot y\): dos veces el valor de \(y\).\par
\textbf{2. El signo igual funciona como una balanza}\par
Los dos lados deben valer lo mismo. Para conservar esa igualdad, sumamos, restamos, multiplicamos o dividimos por la misma cantidad en ambos lados. No dividimos entre cero.\par
\textbf{3. La llave significa ``cumplir todo a la vez''}\par
\[\left\{\begin{aligned}x+2y&=8\\3x-4y&=4\end{aligned}\right.\]
Una solución es la misma pareja de números que cumple las dos condiciones. No elegimos valores distintos para cada renglón.\par
\textbf{4. Sustituir significa reemplazar una letra}\par
Si ya sabemos que \(y=2\), donde aparece \(y\) escribimos 2:\[x+2y=8\quad\longrightarrow\quad x+2(2)=8.\]
Luego calculamos \(2\cdot2=4\) y queda \(x+4=8\).\par
\columnbreak
\textbf{5. Eliminar una letra es cancelar sus términos}\par
Si aparecen \(3x\) y \(3x\), al restarlos obtenemos cero:\[3x-3x=0.\]
Esto no dice que \(x\) valga cero. Solo dice que esos dos términos se cancelaron en la cuenta. Así podremos resolver otra letra.\par
\textbf{6. Multiplicar una ecuación afecta a todo}\par
Al multiplicar \(x+2y=8\) por 3, calculamos \(3x\), \(3(2y)=6y\) y \(3(8)=24\). La nueva forma es \(3x+6y=24\).\par
\textbf{7. Dos reglas de signos que necesitaremos}\par
Restar 6 a menos 4 nos lleva a menos 10: \(-4-6=-10\). Si debes 4 y tu deuda aumenta en 6, debes 10.\par
Restar un negativo equivale a sumar su opuesto: \(-16-(-36)=-16+36=20\).\par
\textbf{8. Comprobar es volver al enunciado}\par
Ponemos nuestros valores en cada ecuación y calculamos. Si una no se cumple, no tenemos una solución del sistema.\par
\textbf{Cómo usar cada explicación}\par
Lee qué buscamos y por qué elegimos la operación. Sigue la cuenta de la tabla; después mira qué conseguimos. La pregunta de repaso sirve para comprobar si entendiste la razón, además del resultado.
\end{multicols}
\newpage\titulo{Índice de ejercicios}
\begin{multicols}{2}

\hyperlink{1.1-1}{1.1 / 1 | Sistema lineal por eliminación}\par
\hyperlink{1.1-2}{1.1 / 2 | Sistema lineal por eliminación}\par
\hyperlink{1.1-3}{1.1 / 3 | Sistema lineal por eliminación}\par
\hyperlink{1.1-4}{1.1 / 4 | Sistema lineal por eliminación}\par
\hyperlink{1.1-5}{1.1 / 5 | Sistema lineal por eliminación}\par
\hyperlink{1.1-6}{1.1 / 6 | Sistema lineal por eliminación}\par
\hyperlink{1.1-7}{1.1 / 7 | Sistema lineal por eliminación}\par
\hyperlink{1.1-8}{1.1 / 8 | Sistema lineal por eliminación}\par
\hyperlink{1.1-9}{1.1 / 9 | Sistema lineal por eliminación}\par
\hyperlink{1.1-10}{1.1 / 10 | Sistema lineal por eliminación}\par
\hyperlink{1.1-11}{1.1 / 11 | Sistema lineal por eliminación}\par
\hyperlink{1.1-12}{1.1 / 12 | Sistema lineal por eliminación}\par
\hyperlink{1.1-13}{1.1 / 13 | Sistema lineal por eliminación}\par
\hyperlink{1.1-14}{1.1 / 14 | Sistema lineal por eliminación}\par
\hyperlink{1.1-15}{1.1 / 15 | Un parámetro que decide la compatibilidad}\par
\hyperlink{1.1-16}{1.1 / 16 | Sumar y multiplicar soluciones}\par
\hyperlink{1.1-17}{1.1 / 17 | Sustitución hacia atrás}\par
\hyperlink{1.1-18}{1.1 / 18 | Sustitución hacia adelante}\par
\hyperlink{1.1-19}{1.1 / 19 | Un mismo valor debe cumplir todo el sistema}\par
\hyperlink{1.1-20}{1.1 / 20 | Comprobamos antes de concluir}\par
\hyperlink{1.1-21}{1.1 / 21 | Tres planos: punto, vacío o recta}\par
\hyperlink{1.1-22}{1.1 / 22 | ¿Qué comparten los tres planos?}\par
\hyperlink{1.1-23}{1.1 / 23 | Producción de gasolina}\par
\hyperlink{1.1-24}{1.1 / 24 | Dos tipos de plástico}\par
\hyperlink{1.1-25}{1.1 / 25 | Una dieta con tres alimentos}\par
\hyperlink{1.1-26}{1.1 / 26 | Reveladores: una familia de planes de producción}\par
\hyperlink{1.1-27}{1.1 / 27 | Una parábola determinada por tres puntos}\par
\hyperlink{1.1-28}{1.1 / 28 | Reparto de una herencia e intereses}\par
\hyperlink{1.1-T.1}{1.1 / T.1 | Intercambiar dos ecuaciones conserva las soluciones}\par
\hyperlink{1.1-T.2}{1.1 / T.2 | Multiplicar una ecuación por una constante no nula}\par
\hyperlink{1.1-T.3}{1.1 / T.3 | Sumar un múltiplo de otra ecuación}\par
\hyperlink{1.1-T.4}{1.1 / T.4 | Todo sistema homogéneo tiene solución}\par
\hyperlink{1.2-1}{1.2 / 1 | Leer una entrada de una matriz}\par
\hyperlink{1.2-2}{1.2 / 2 | Igualdad de matrices: hallar a, b, c y d}\par
\hyperlink{1.2-3}{1.2 / 3 | Igualdad de matrices: hallar a, b, c y d}\par
\hyperlink{1.2-4}{1.2 / 4 | Operaciones con matrices}\par
\hyperlink{1.2-5}{1.2 / 5 | Operaciones con matrices}\par
\hyperlink{1.2-6}{1.2 / 6 | Transposición y dimensiones}\par
\hyperlink{1.2-7}{1.2 / 7 | Transposición y dimensiones}\par
\hyperlink{1.2-8}{1.2 / 8 | ¿Es una combinación lineal?}\par
\hyperlink{1.2-9}{1.2 / 9 | ¿Es una combinación lineal?}\par
\hyperlink{1.2-10}{1.2 / 10 | Una matriz que depende de $\lambda$}\par
\hyperlink{1.2-11}{1.2 / 11 | Operaciones con matrices binarias}\par
\hyperlink{1.2-12}{1.2 / 12 | Operaciones con matrices binarias}\par
\hyperlink{1.2-13}{1.2 / 13 | Encontrar matrices binarias B y C}\par
\hyperlink{1.2-14}{1.2 / 14 | Completar un vector binario}\par
\hyperlink{1.2-15}{1.2 / 15 | Completar un vector binario}\par
\hyperlink{1.2-T.1}{1.2 / T.1 | Suma y diferencia de matrices diagonales}\par
\hyperlink{1.2-T.2}{1.2 / T.2 | Suma y diferencia de matrices escalares}\par
\hyperlink{1.2-T.3}{1.2 / T.3 | Una matriz simbólica y su transpuesta}\par
\hyperlink{1.2-T.4}{1.2 / T.4 | Si kA es cero, k es cero o A es cero}\par
\hyperlink{1.2-T.5}{1.2 / T.5 | Matrices triangulares y diagonales}\par
\hyperlink{1.2-T.6}{1.2 / T.6 | La transpuesta intercambia los tipos triangulares}\par
\hyperlink{1.2-T.7}{1.2 / T.7 | La diagonal de A-A\textsuperscript{T}}\par
\hyperlink{1.2-T.8}{1.2 / T.8 | El vector cero es neutro para la suma}\par
\hyperlink{1.2-T.9}{1.2 / T.9 | Todos los 2-vectores binarios}\par
\hyperlink{1.2-T.10}{1.2 / T.10 | Todos los 3-vectores binarios}\par
\hyperlink{1.2-T.11}{1.2 / T.11 | Todos los 4-vectores binarios}\par
\hyperlink{1.2-T.12}{1.2 / T.12 | Contar vectores binarios de longitud n}\par
\hyperlink{1.2-T.13}{1.2 / T.13 | Todas las matrices binarias de orden 2}\par
\hyperlink{1.2-T.14}{1.2 / T.14 | Contar matrices binarias de orden 3}\par
\hyperlink{1.2-T.15}{1.2 / T.15 | Contar matrices binarias de orden n}\par
\hyperlink{1.2-T.16}{1.2 / T.16 | Interruptores: obtener todos OFF}\par
\hyperlink{1.2-T.17}{1.2 / T.17 | Interruptores: obtener todos ON}\par
\hyperlink{1.2-T.18}{1.2 / T.18 | Invertir todos los interruptores con la misma máscara}\par
\hyperlink{1.2-ML.1}{1.2 / ML.1 | MATLAB: entradas y formatos}\par
\hyperlink{1.2-ML.2}{1.2 / ML.2 | MATLAB: matriz de Hilbert}\par
\hyperlink{1.2-ML.3}{1.2 / ML.3 | MATLAB: bingen para vectores binarios}\par
\hyperlink{1.2-ML.4}{1.2 / ML.4 | MATLAB: enumerar matrices binarias}\par
\hyperlink{1.2-ML.5}{1.2 / ML.5 | MATLAB: binadd y aritmética módulo 2}\par
\end{multicols}
\newpage\fancyfoot[L]{\small Sección 1.1 | Ejercicio 1 | Paso a paso}\hypertarget{1.1-1}{}
\pdfbookmark[1]{1.1 / 1}{bm1.1-1}
\titulo{1 | Sistema lineal por eliminación}
\textbf{Enunciado}\par Resuelva el sistema lineal dado por medio del método de eliminación.
\[\left\{\begin{aligned}x + 2y &= 8\\3x - 4y &= 4\end{aligned}\right.\]
\begin{multicols*}{2}
\noindent\begin{minipage}{\linewidth}\textbf{1. Primero: ¿qué significa resolver este sistema?}\\[2pt]Las letras x, y representan números que todavía no conocemos. Una solución es una elección de esos números que hace verdaderas TODAS las ecuaciones al mismo tiempo. La llave de la izquierda significa que las condiciones se deben cumplir juntas.
\par\smallskip\textbf{Qué buscamos.} Encontrar una misma pareja de valores.
\par\smallskip El número 3 delante de x indica tres veces ese mismo valor; es decir, sumar tres copias:
\[3x=x+x+x\]
\par\smallskip Ejemplo numérico de la notación: si x valiera 2, tres veces ese valor sería 6. Este 2 es solo un ejemplo; todavía no es la respuesta del ejercicio. El punto centrado indica multiplicación:
\[3\cdot 2=2+2+2=6\]
\par\smallskip\textbf{Comprueba que lo entiendes.} ¿Basta con que los valores funcionen en una sola ecuación?\par\textit{Respuesta:} No. Deben funcionar en todas las ecuaciones del sistema, usando los mismos valores.
\end{minipage}\par\medskip
\noindent\begin{minipage}{\linewidth}\textbf{2. Nuestro plan: ir quitando letras}\\[2pt]Una ecuación con una sola letra es más fácil de resolver. Vamos a combinar ecuaciones para que una letra desaparezca. Cuando sepamos un valor, lo pondremos en las ecuaciones que faltan.
\par\smallskip\textbf{Qué buscamos.} Primero quitar x para calcular y. Después calcular x.
\par\smallskip\textbf{Por qué.} Si una cantidad se resta de sí misma, queda cero. Usaremos esta idea con términos que tengan la misma letra.
\par\smallskip Ejemplo de la idea que utilizaremos:
\[3x-3x=0\]
\par\smallskip\textbf{Qué aprendimos.} Eliminar una letra de una ecuación NO significa que su valor sea cero; significa que sus términos se cancelaron en esa cuenta.
\end{minipage}\par\medskip
\noindent\begin{minipage}{\linewidth}\textbf{3. ¿Cómo haremos desaparecer x?}\\[2pt]Trabajaremos con las ecuaciones 1 y 2. Queremos que los términos con x tengan el mismo tamaño. Así podremos restarlos.
\par\smallskip\textbf{Qué buscamos.} Conseguir términos con 3 veces x en las dos ecuaciones.
\par\smallskip\textbf{Por qué.} Tener 3 veces la misma letra en ambos lados de la resta nos permitirá cancelarla.
\par\smallskip {\small\renewcommand{\arraystretch}{1.3}\begin{tabular}{>{\raggedright\arraybackslash}p{2.58cm}>{\raggedright\arraybackslash}p{2.47cm}>{\raggedright\arraybackslash}p{6.17cm}}\hline\rowcolor{black!7} \textbf{Dónde miramos} & \textbf{Qué término hay} & \textbf{Qué necesitamos hacer}\\\hline Ecuación 1 & $x$ & Multiplicar TODA esta ecuación por 3\\ Ecuación 2 & $3x$ & La dejamos igual.\\ \hline\end{tabular}}\par\smallskip
\par\smallskip\textbf{Qué aprendimos.} Estos multiplicadores se eligen para igualar los términos con x. Los demás términos también se multiplican para conservar la igualdad.
\end{minipage}\par\medskip
\noindent\begin{minipage}{\linewidth}\textbf{4. Multiplicamos la ecuación 1 por 3}\\[2pt]Multiplicar una ecuación por 3 significa multiplicar CADA término de la izquierda y también el número de la derecha. La igualdad se conserva porque hacemos lo mismo en ambos lados.
\par\smallskip Esta es la ecuación antes de multiplicar:
\[x + 2y = 8\]
\par\smallskip {\small\renewcommand{\arraystretch}{1.3}\begin{tabular}{>{\raggedright\arraybackslash}p{2.80cm}>{\raggedright\arraybackslash}p{1.91cm}>{\raggedright\arraybackslash}p{6.50cm}}\hline\rowcolor{black!7} \textbf{Parte de la ecuación} & \textbf{Antes} & \textbf{La cuenta y su resultado}\\\hline Término con x & $x$ & $3\cdot\left(x\right)=3x$\\ Término con y & $2y$ & $3\cdot\left(2y\right)=6y$\\ Número a la derecha del igual & $8$ & $3\cdot\left(8\right)=24$\\ \hline\end{tabular}}\par\smallskip
\par\smallskip Después de hacer todas las multiplicaciones queda:
\[3x + 6y = 24\]
\par\smallskip\textbf{Qué aprendimos.} La misma condición ahora está escrita con los términos multiplicados. Ya tenemos el tamaño que buscábamos para x.
\par\smallskip\textbf{Comprueba que lo entiendes.} ¿Podríamos multiplicar solamente el término con x?\par\textit{Respuesta:} No. Eso cambiaría la condición. Debemos multiplicar todos los términos, también el número que está después del igual.
\end{minipage}\par\medskip
\noindent\begin{minipage}{\linewidth}\textbf{5. Restamos las dos ecuaciones preparadas}\\[2pt]Restamos la ecuación 1 preparada a la ecuación 2 preparada. Juntamos cada término solo con otro que tenga la misma letra. Los números sin letra se calculan por separado.
\par\smallskip {\small\renewcommand{\arraystretch}{1.3}\begin{tabular}{>{\raggedright\arraybackslash}p{2.80cm}>{\raggedright\arraybackslash}p{5.38cm}>{\raggedright\arraybackslash}p{3.03cm}}\hline\rowcolor{black!7} \textbf{Qué agrupamos} & \textbf{Cuenta} & \textbf{Qué significa}\\\hline Los términos con x & $3x - \left(3x\right)=0$ & Se cancela esta letra.\\ Los términos con y & $-4y - \left(6y\right)=-10y$ & Quedan -10y.\\ Los números a la derecha del igual & $4 - \left(24\right)=-20$ & Este será el número de la nueva igualdad.\\ \hline\end{tabular}}\par\smallskip
\par\smallskip Primera igualdad que usaremos:
\[3x - 4y = 4\]
\par\smallskip Le restamos esta igualdad:
\[3x + 6y = 24\]
\par\smallskip\textbf{Qué aprendimos.} Los términos con x dan cero. Los términos con letras distintas no se mezclan.
\end{minipage}\par\medskip
\noindent\begin{minipage}{\linewidth}\textbf{6. Mira lo que hemos conseguido}\\[2pt]La nueva ecuación es el resultado de la cuenta anterior. Solo contiene y: después podremos despejarla.
\par\smallskip Escribimos la nueva condición:
\[-10y = -20\]
\[\left\{\begin{aligned}x + 2y &= 8\\-10y &= -20\end{aligned}\right.\]
\par\smallskip\textbf{Qué aprendimos.} Conservamos la ecuación 1 y reemplazamos la 2 por esta nueva condición. De esa manera conservamos toda la información necesaria para resolver el sistema.
\end{minipage}\par\medskip
\noindent\begin{minipage}{\linewidth}\textbf{7. Ya tenemos una ecuación con una sola letra}\\[2pt]En esta condición solamente falta conocer y. El número -10 está multiplicando a esa letra. Ahora podemos deshacer esa multiplicación.
\par\smallskip Esta es la condición que vamos a resolver:
\[-10y = -20\]
\end{minipage}\par\medskip
\noindent\begin{minipage}{\linewidth}\textbf{8. Calculamos cuánto vale y}\\[2pt]Para deshacer una multiplicación por -10, dividimos AMBOS lados entre -10. En el lado izquierdo queda una sola y. En el derecho hacemos la división.
\par\smallskip Dividimos los dos lados entre el mismo número:
\[\frac{-10y}{-10}=\frac{-20}{-10}\]
\par\smallskip Ya podemos leer el valor de la letra:
\[y=\frac{-20}{-10}=2\]
\par\smallskip\textbf{Qué aprendimos.} Un número negativo dividido entre otro negativo da positivo.
\end{minipage}\par\medskip
\noindent\begin{minipage}{\linewidth}\textbf{9. Usamos lo que ya sabemos para hallar x}\\[2pt]Ya sabemos y = 2. Sustituir significa escribir ese valor cada vez que aparezca su letra. Usamos paréntesis para que se vea qué número estamos multiplicando.
\par\smallskip\textbf{Qué buscamos.} Dejar solamente x como letra desconocida en esta ecuación.
\par\smallskip Tomamos esta condición que obtuvimos antes:
\[x + 2y = 8\]
\par\smallskip Reemplazamos las letras conocidas por sus valores:
\[x + 2\left(2\right)=8\]
\end{minipage}\par\medskip
\noindent\begin{minipage}{\linewidth}\textbf{10. Calculamos las multiplicaciones que ya se pueden hacer}\\[2pt]La letra x sigue sin conocerse. Las otras partes ya tienen un valor o una expresión conocida y podemos efectuar sus productos.
\par\smallskip {\small\renewcommand{\arraystretch}{1.3}\begin{tabular}{>{\raggedright\arraybackslash}p{3.36cm}>{\raggedright\arraybackslash}p{4.48cm}>{\raggedright\arraybackslash}p{3.36cm}}\hline\rowcolor{black!7} \textbf{Parte conocida} & \textbf{La multiplicación} & \textbf{Resultado}\\\hline El término 2y & $2\cdot\left(2\right)$ & $4$\\ \hline\end{tabular}}\par\smallskip
\par\smallskip Al reunir esas partes conocidas, la ecuación queda:
\[x + 4=8\]
\end{minipage}\par\medskip
\noindent\begin{minipage}{\linewidth}\textbf{11. Dejamos el término con x solo a un lado}\\[2pt]Para quitar la parte conocida, restamos 4 en AMBOS lados. Piensa en una balanza: si hacemos lo mismo a ambos lados, la igualdad se mantiene.
\par\smallskip Hacemos la misma operación a la izquierda y a la derecha:
\[x + 4 - 4=8 - 4\]
\par\smallskip Las partes que se compensan dan cero; queda:
\[x=4\]
\par\smallskip\textbf{Qué aprendimos.} «Pasar al otro lado» es una forma abreviada de decir que hicimos la misma suma o resta en los dos lados.
\end{minipage}\par\medskip
\noindent\begin{minipage}{\linewidth}\textbf{12. Leemos el valor de x}\\[2pt]La letra ya está sola: su número multiplicador es 1, que no se escribe. La igualdad nos dice directamente cuánto vale.
\par\smallskip Anotamos el valor encontrado:
\[x=4\]
\end{minipage}\par\medskip
\noindent\begin{minipage}{\linewidth}\textbf{13. Reunimos la respuesta; todavía falta comprobarla}\\[2pt]Escribimos todos los valores en el orden x, y. Estos valores deben funcionar juntos, no cada uno en una ecuación diferente.
\[\left\{\begin{aligned}x&=4\\y&=2\end{aligned}\right.\]
\par\smallskip\textbf{Qué aprendimos.} El siguiente paso es volver a las ecuaciones del enunciado y comprobar cada condición.
\end{minipage}\par\medskip
\noindent\begin{minipage}{\linewidth}\textbf{14. ¿La respuesta cumple la ecuación 1?}\\[2pt]Comprobar significa volver a la condición original y calcularla con nuestra respuesta. Si los dos lados dan el mismo número, esta condición sí se cumple. Deben cumplirse todas.
\par\smallskip La condición del enunciado era:
\[x + 2y = 8\]
\par\smallskip Ponemos nuestros valores donde estaban las letras:
\[\left(4\right) + 2\left(2\right)=8\]
\par\smallskip Hacemos las multiplicaciones y luego la suma:
\[4 + 4=8\]
\par\smallskip Comparamos los dos lados:
\[8=8\]
\par\smallskip\textbf{Qué aprendimos.} La igualdad final es verdadera. Nuestra respuesta sí cumple la ecuación 1.
\end{minipage}\par\medskip
\noindent\begin{minipage}{\linewidth}\textbf{15. ¿La respuesta cumple la ecuación 2?}\\[2pt]Comprobar significa volver a la condición original y calcularla con nuestra respuesta. Si los dos lados dan el mismo número, esta condición sí se cumple. Deben cumplirse todas.
\par\smallskip La condición del enunciado era:
\[3x - 4y = 4\]
\par\smallskip Ponemos nuestros valores donde estaban las letras:
\[3\left(4\right) - 4\left(2\right)=4\]
\par\smallskip Hacemos las multiplicaciones y luego la suma:
\[12 - 8=4\]
\par\smallskip Comparamos los dos lados:
\[4=4\]
\par\smallskip\textbf{Qué aprendimos.} La igualdad final es verdadera. Nuestra respuesta sí cumple la ecuación 2.
\end{minipage}\par\medskip
\respuesta{x = 4; y = 2.}
\noindent\begin{minipage}{\linewidth}\centering\includegraphics[width=.96\linewidth,height=6.8cm,keepaspectratio]{../recursos/grafico_1_1-1.png}\\{\small Lectura visual complementaria; la demostración está en las operaciones.}\end{minipage}
\end{multicols*}
\newpage\fancyfoot[L]{\small Sección 1.1 | Ejercicio 2 | Paso a paso}\hypertarget{1.1-2}{}
\pdfbookmark[1]{1.1 / 2}{bm1.1-2}
\titulo{2 | Sistema lineal por eliminación}
\textbf{Enunciado}\par Resuelva el sistema lineal dado por medio del método de eliminación.
\[\left\{\begin{aligned}2x - 3y + 4z &= -12\\x - 2y + z &= -5\\3x + y + 2z &= 1\end{aligned}\right.\]
\begin{multicols*}{2}
\noindent\begin{minipage}{\linewidth}\textbf{1. Primero: ¿qué significa resolver este sistema?}\\[2pt]Las letras x, y, z representan números que todavía no conocemos. Una solución es una elección de esos números que hace verdaderas TODAS las ecuaciones al mismo tiempo. La llave de la izquierda significa que las condiciones se deben cumplir juntas.
\par\smallskip\textbf{Qué buscamos.} Encontrar un mismo grupo de valores para las 3 letras.
\par\smallskip El número 3 delante de x indica tres veces ese mismo valor; es decir, sumar tres copias:
\[3x=x+x+x\]
\par\smallskip Ejemplo numérico de la notación: si x valiera 2, tres veces ese valor sería 6. Este 2 es solo un ejemplo; todavía no es la respuesta del ejercicio. El punto centrado indica multiplicación:
\[3\cdot 2=2+2+2=6\]
\par\smallskip\textbf{Comprueba que lo entiendes.} ¿Basta con que los valores funcionen en una sola ecuación?\par\textit{Respuesta:} No. Deben funcionar en todas las ecuaciones del sistema, usando los mismos valores.
\end{minipage}\par\medskip
\noindent\begin{minipage}{\linewidth}\textbf{2. Nuestro plan: ir quitando letras}\\[2pt]Una ecuación con una sola letra es más fácil de resolver. Vamos a combinar ecuaciones para que una letra desaparezca. Cuando sepamos un valor, lo pondremos en las ecuaciones que faltan.
\par\smallskip\textbf{Qué buscamos.} Primero quitar x de las otras ecuaciones. Después seguir con las letras restantes.
\par\smallskip\textbf{Por qué.} Si una cantidad se resta de sí misma, queda cero. Usaremos esta idea con términos que tengan la misma letra.
\par\smallskip Ejemplo de la idea que utilizaremos:
\[3x-3x=0\]
\par\smallskip\textbf{Qué aprendimos.} Eliminar una letra de una ecuación NO significa que su valor sea cero; significa que sus términos se cancelaron en esa cuenta.
\end{minipage}\par\medskip
\noindent\begin{minipage}{\linewidth}\textbf{3. Ponemos una ecuación más cómoda arriba}\\[2pt]Cambiamos de lugar las ecuaciones 1 y 2. Las condiciones siguen siendo las mismas; solamente cambia el orden en que las leemos.
\par\smallskip\textbf{Por qué.} La ecuación que subimos tiene un número más pequeño delante de x. Eso facilita las multiplicaciones que haremos a continuación.
\[\left\{\begin{aligned}x - 2y + z &= -5\\2x - 3y + 4z &= -12\\3x + y + 2z &= 1\end{aligned}\right.\]
\end{minipage}\par\medskip
\noindent\begin{minipage}{\linewidth}\textbf{4. ¿Cómo haremos desaparecer x?}\\[2pt]Trabajaremos con las ecuaciones 1 y 2. Queremos que los términos con x tengan el mismo tamaño. Así podremos restarlos.
\par\smallskip\textbf{Qué buscamos.} Conseguir términos con 2 veces x en las dos ecuaciones.
\par\smallskip\textbf{Por qué.} Tener 2 veces la misma letra en ambos lados de la resta nos permitirá cancelarla.
\par\smallskip {\small\renewcommand{\arraystretch}{1.3}\begin{tabular}{>{\raggedright\arraybackslash}p{2.58cm}>{\raggedright\arraybackslash}p{2.47cm}>{\raggedright\arraybackslash}p{6.17cm}}\hline\rowcolor{black!7} \textbf{Dónde miramos} & \textbf{Qué término hay} & \textbf{Qué necesitamos hacer}\\\hline Ecuación 1 & $x$ & Multiplicar TODA esta ecuación por 2\\ Ecuación 2 & $2x$ & La dejamos igual.\\ \hline\end{tabular}}\par\smallskip
\par\smallskip\textbf{Qué aprendimos.} Estos multiplicadores se eligen para igualar los términos con x. Los demás términos también se multiplican para conservar la igualdad.
\end{minipage}\par\medskip
\noindent\begin{minipage}{\linewidth}\textbf{5. Multiplicamos la ecuación 1 por 2}\\[2pt]Multiplicar una ecuación por 2 significa multiplicar CADA término de la izquierda y también el número de la derecha. La igualdad se conserva porque hacemos lo mismo en ambos lados.
\par\smallskip Esta es la ecuación antes de multiplicar:
\[x - 2y + z = -5\]
\par\smallskip {\small\renewcommand{\arraystretch}{1.3}\begin{tabular}{>{\raggedright\arraybackslash}p{2.80cm}>{\raggedright\arraybackslash}p{1.91cm}>{\raggedright\arraybackslash}p{6.50cm}}\hline\rowcolor{black!7} \textbf{Parte de la ecuación} & \textbf{Antes} & \textbf{La cuenta y su resultado}\\\hline Término con x & $x$ & $2\cdot\left(x\right)=2x$\\ Término con y & $-2y$ & $2\cdot\left(-2y\right)=-4y$\\ Término con z & $z$ & $2\cdot\left(z\right)=2z$\\ Número a la derecha del igual & $-5$ & $2\cdot\left(-5\right)=-10$\\ \hline\end{tabular}}\par\smallskip
\par\smallskip Después de hacer todas las multiplicaciones queda:
\[2x - 4y + 2z = -10\]
\par\smallskip\textbf{Qué aprendimos.} La misma condición ahora está escrita con los términos multiplicados. Ya tenemos el tamaño que buscábamos para x.
\par\smallskip\textbf{Comprueba que lo entiendes.} ¿Podríamos multiplicar solamente el término con x?\par\textit{Respuesta:} No. Eso cambiaría la condición. Debemos multiplicar todos los términos, también el número que está después del igual.
\end{minipage}\par\medskip
\noindent\begin{minipage}{\linewidth}\textbf{6. Restamos las dos ecuaciones preparadas}\\[2pt]Restamos la ecuación 1 preparada a la ecuación 2 preparada. Juntamos cada término solo con otro que tenga la misma letra. Los números sin letra se calculan por separado.
\par\smallskip {\small\renewcommand{\arraystretch}{1.3}\begin{tabular}{>{\raggedright\arraybackslash}p{2.80cm}>{\raggedright\arraybackslash}p{5.38cm}>{\raggedright\arraybackslash}p{3.03cm}}\hline\rowcolor{black!7} \textbf{Qué agrupamos} & \textbf{Cuenta} & \textbf{Qué significa}\\\hline Los términos con x & $2x - \left(2x\right)=0$ & Se cancela esta letra.\\ Los términos con y & $-3y - \left(-4y\right)=y$ & Quedan y.\\ Los términos con z & $4z - \left(2z\right)=2z$ & Quedan 2z.\\ Los números a la derecha del igual & $-12 - \left(-10\right)=-12+10=-2$ & Este será el número de la nueva igualdad.\\ \hline\end{tabular}}\par\smallskip
\par\smallskip Primera igualdad que usaremos:
\[2x - 3y + 4z = -12\]
\par\smallskip Le restamos esta igualdad:
\[2x - 4y + 2z = -10\]
\par\smallskip\textbf{Qué aprendimos.} Los términos con x dan cero. Recuerda: restar un número negativo equivale a sumar su opuesto.
\end{minipage}\par\medskip
\noindent\begin{minipage}{\linewidth}\textbf{7. Mira lo que hemos conseguido}\\[2pt]La nueva ecuación es el resultado de la cuenta anterior. Todavía contiene y, z. Seguiremos reduciendo las condiciones.
\par\smallskip Escribimos la nueva condición:
\[y + 2z = -2\]
\[\left\{\begin{aligned}x - 2y + z &= -5\\y + 2z &= -2\\3x + y + 2z &= 1\end{aligned}\right.\]
\par\smallskip\textbf{Qué aprendimos.} Conservamos la ecuación 1 y reemplazamos la 2 por esta nueva condición. De esa manera conservamos toda la información necesaria para resolver el sistema.
\end{minipage}\par\medskip
\noindent\begin{minipage}{\linewidth}\textbf{8. ¿Cómo haremos desaparecer x?}\\[2pt]Trabajaremos con las ecuaciones 1 y 3. Queremos que los términos con x tengan el mismo tamaño. Así podremos restarlos.
\par\smallskip\textbf{Qué buscamos.} Conseguir términos con 3 veces x en las dos ecuaciones.
\par\smallskip\textbf{Por qué.} Tener 3 veces la misma letra en ambos lados de la resta nos permitirá cancelarla.
\par\smallskip {\small\renewcommand{\arraystretch}{1.3}\begin{tabular}{>{\raggedright\arraybackslash}p{2.58cm}>{\raggedright\arraybackslash}p{2.47cm}>{\raggedright\arraybackslash}p{6.17cm}}\hline\rowcolor{black!7} \textbf{Dónde miramos} & \textbf{Qué término hay} & \textbf{Qué necesitamos hacer}\\\hline Ecuación 1 & $x$ & Multiplicar TODA esta ecuación por 3\\ Ecuación 3 & $3x$ & La dejamos igual.\\ \hline\end{tabular}}\par\smallskip
\par\smallskip\textbf{Qué aprendimos.} Estos multiplicadores se eligen para igualar los términos con x. Los demás términos también se multiplican para conservar la igualdad.
\end{minipage}\par\medskip
\noindent\begin{minipage}{\linewidth}\textbf{9. Multiplicamos la ecuación 1 por 3}\\[2pt]Multiplicar una ecuación por 3 significa multiplicar CADA término de la izquierda y también el número de la derecha. La igualdad se conserva porque hacemos lo mismo en ambos lados.
\par\smallskip Esta es la ecuación antes de multiplicar:
\[x - 2y + z = -5\]
\par\smallskip {\small\renewcommand{\arraystretch}{1.3}\begin{tabular}{>{\raggedright\arraybackslash}p{2.80cm}>{\raggedright\arraybackslash}p{1.91cm}>{\raggedright\arraybackslash}p{6.50cm}}\hline\rowcolor{black!7} \textbf{Parte de la ecuación} & \textbf{Antes} & \textbf{La cuenta y su resultado}\\\hline Término con x & $x$ & $3\cdot\left(x\right)=3x$\\ Término con y & $-2y$ & $3\cdot\left(-2y\right)=-6y$\\ Término con z & $z$ & $3\cdot\left(z\right)=3z$\\ Número a la derecha del igual & $-5$ & $3\cdot\left(-5\right)=-15$\\ \hline\end{tabular}}\par\smallskip
\par\smallskip Después de hacer todas las multiplicaciones queda:
\[3x - 6y + 3z = -15\]
\par\smallskip\textbf{Qué aprendimos.} La misma condición ahora está escrita con los términos multiplicados. Ya tenemos el tamaño que buscábamos para x.
\par\smallskip\textbf{Comprueba que lo entiendes.} ¿Podríamos multiplicar solamente el término con x?\par\textit{Respuesta:} No. Eso cambiaría la condición. Debemos multiplicar todos los términos, también el número que está después del igual.
\end{minipage}\par\medskip
\noindent\begin{minipage}{\linewidth}\textbf{10. Restamos las dos ecuaciones preparadas}\\[2pt]Restamos la ecuación 1 preparada a la ecuación 3 preparada. Juntamos cada término solo con otro que tenga la misma letra. Los números sin letra se calculan por separado.
\par\smallskip {\small\renewcommand{\arraystretch}{1.3}\begin{tabular}{>{\raggedright\arraybackslash}p{2.80cm}>{\raggedright\arraybackslash}p{5.38cm}>{\raggedright\arraybackslash}p{3.03cm}}\hline\rowcolor{black!7} \textbf{Qué agrupamos} & \textbf{Cuenta} & \textbf{Qué significa}\\\hline Los términos con x & $3x - \left(3x\right)=0$ & Se cancela esta letra.\\ Los términos con y & $y - \left(-6y\right)=7y$ & Quedan 7y.\\ Los términos con z & $2z - \left(3z\right)=-z$ & Quedan -z.\\ Los números a la derecha del igual & $1 - \left(-15\right)=1+15=16$ & Este será el número de la nueva igualdad.\\ \hline\end{tabular}}\par\smallskip
\par\smallskip Primera igualdad que usaremos:
\[3x + y + 2z = 1\]
\par\smallskip Le restamos esta igualdad:
\[3x - 6y + 3z = -15\]
\par\smallskip\textbf{Qué aprendimos.} Los términos con x dan cero. Recuerda: restar un número negativo equivale a sumar su opuesto.
\end{minipage}\par\medskip
\noindent\begin{minipage}{\linewidth}\textbf{11. Mira lo que hemos conseguido}\\[2pt]La nueva ecuación es el resultado de la cuenta anterior. Todavía contiene y, z. Seguiremos reduciendo las condiciones.
\par\smallskip Escribimos la nueva condición:
\[7y - z = 16\]
\[\left\{\begin{aligned}x - 2y + z &= -5\\y + 2z &= -2\\7y - z &= 16\end{aligned}\right.\]
\par\smallskip\textbf{Qué aprendimos.} Conservamos la ecuación 1 y reemplazamos la 3 por esta nueva condición. De esa manera conservamos toda la información necesaria para resolver el sistema.
\end{minipage}\par\medskip
\noindent\begin{minipage}{\linewidth}\textbf{12. ¿Cómo haremos desaparecer y?}\\[2pt]Trabajaremos con las ecuaciones 2 y 3. Queremos que los términos con y tengan el mismo tamaño. Así podremos restarlos.
\par\smallskip\textbf{Qué buscamos.} Conseguir términos con 7 veces y en las dos ecuaciones.
\par\smallskip\textbf{Por qué.} Tener 7 veces la misma letra en ambos lados de la resta nos permitirá cancelarla.
\par\smallskip {\small\renewcommand{\arraystretch}{1.3}\begin{tabular}{>{\raggedright\arraybackslash}p{2.58cm}>{\raggedright\arraybackslash}p{2.47cm}>{\raggedright\arraybackslash}p{6.17cm}}\hline\rowcolor{black!7} \textbf{Dónde miramos} & \textbf{Qué término hay} & \textbf{Qué necesitamos hacer}\\\hline Ecuación 2 & $y$ & Multiplicar TODA esta ecuación por 7\\ Ecuación 3 & $7y$ & La dejamos igual.\\ \hline\end{tabular}}\par\smallskip
\par\smallskip\textbf{Qué aprendimos.} Estos multiplicadores se eligen para igualar los términos con y. Los demás términos también se multiplican para conservar la igualdad.
\end{minipage}\par\medskip
\noindent\begin{minipage}{\linewidth}\textbf{13. Multiplicamos la ecuación 2 por 7}\\[2pt]Multiplicar una ecuación por 7 significa multiplicar CADA término de la izquierda y también el número de la derecha. La igualdad se conserva porque hacemos lo mismo en ambos lados.
\par\smallskip Esta es la ecuación antes de multiplicar:
\[y + 2z = -2\]
\par\smallskip {\small\renewcommand{\arraystretch}{1.3}\begin{tabular}{>{\raggedright\arraybackslash}p{2.80cm}>{\raggedright\arraybackslash}p{1.91cm}>{\raggedright\arraybackslash}p{6.50cm}}\hline\rowcolor{black!7} \textbf{Parte de la ecuación} & \textbf{Antes} & \textbf{La cuenta y su resultado}\\\hline Término con y & $y$ & $7\cdot\left(y\right)=7y$\\ Término con z & $2z$ & $7\cdot\left(2z\right)=14z$\\ Número a la derecha del igual & $-2$ & $7\cdot\left(-2\right)=-14$\\ \hline\end{tabular}}\par\smallskip
\par\smallskip Después de hacer todas las multiplicaciones queda:
\[7y + 14z = -14\]
\par\smallskip\textbf{Qué aprendimos.} La misma condición ahora está escrita con los términos multiplicados. Ya tenemos el tamaño que buscábamos para y.
\par\smallskip\textbf{Comprueba que lo entiendes.} ¿Podríamos multiplicar solamente el término con y?\par\textit{Respuesta:} No. Eso cambiaría la condición. Debemos multiplicar todos los términos, también el número que está después del igual.
\end{minipage}\par\medskip
\noindent\begin{minipage}{\linewidth}\textbf{14. Restamos las dos ecuaciones preparadas}\\[2pt]Restamos la ecuación 2 preparada a la ecuación 3 preparada. Juntamos cada término solo con otro que tenga la misma letra. Los números sin letra se calculan por separado.
\par\smallskip {\small\renewcommand{\arraystretch}{1.3}\begin{tabular}{>{\raggedright\arraybackslash}p{2.80cm}>{\raggedright\arraybackslash}p{5.38cm}>{\raggedright\arraybackslash}p{3.03cm}}\hline\rowcolor{black!7} \textbf{Qué agrupamos} & \textbf{Cuenta} & \textbf{Qué significa}\\\hline Los términos con y & $7y - \left(7y\right)=0$ & Se cancela esta letra.\\ Los términos con z & $-z - \left(14z\right)=-15z$ & Quedan -15z.\\ Los números a la derecha del igual & $16 - \left(-14\right)=16+14=30$ & Este será el número de la nueva igualdad.\\ \hline\end{tabular}}\par\smallskip
\par\smallskip Primera igualdad que usaremos:
\[7y - z = 16\]
\par\smallskip Le restamos esta igualdad:
\[7y + 14z = -14\]
\par\smallskip\textbf{Qué aprendimos.} Los términos con y dan cero. Recuerda: restar un número negativo equivale a sumar su opuesto.
\end{minipage}\par\medskip
\noindent\begin{minipage}{\linewidth}\textbf{15. Mira lo que hemos conseguido}\\[2pt]La nueva ecuación es el resultado de la cuenta anterior. Solo contiene z: después podremos despejarla.
\par\smallskip Escribimos la nueva condición:
\[-15z = 30\]
\[\left\{\begin{aligned}x - 2y + z &= -5\\y + 2z &= -2\\-15z &= 30\end{aligned}\right.\]
\par\smallskip\textbf{Qué aprendimos.} Conservamos la ecuación 2 y reemplazamos la 3 por esta nueva condición. De esa manera conservamos toda la información necesaria para resolver el sistema.
\end{minipage}\par\medskip
\noindent\begin{minipage}{\linewidth}\textbf{16. Ya tenemos una ecuación con una sola letra}\\[2pt]En esta condición solamente falta conocer z. El número -15 está multiplicando a esa letra. Ahora podemos deshacer esa multiplicación.
\par\smallskip Esta es la condición que vamos a resolver:
\[-15z = 30\]
\end{minipage}\par\medskip
\noindent\begin{minipage}{\linewidth}\textbf{17. Calculamos cuánto vale z}\\[2pt]Para deshacer una multiplicación por -15, dividimos AMBOS lados entre -15. En el lado izquierdo queda una sola z. En el derecho hacemos la división.
\par\smallskip Dividimos los dos lados entre el mismo número:
\[\frac{-15z}{-15}=\frac{30}{-15}\]
\par\smallskip Ya podemos leer el valor de la letra:
\[z=\frac{30}{-15}=-2\]
\par\smallskip\textbf{Qué aprendimos.} Este valor se usará en las ecuaciones donde todavía falten letras.
\end{minipage}\par\medskip
\noindent\begin{minipage}{\linewidth}\textbf{18. Usamos lo que ya sabemos para hallar y}\\[2pt]Ya sabemos z = -2. Sustituir significa escribir ese valor cada vez que aparezca su letra. Usamos paréntesis para que se vea qué número estamos multiplicando.
\par\smallskip\textbf{Qué buscamos.} Dejar solamente y como letra desconocida en esta ecuación.
\par\smallskip Tomamos esta condición que obtuvimos antes:
\[y + 2z = -2\]
\par\smallskip Reemplazamos las letras conocidas por sus valores:
\[y + 2\left(-2\right)=-2\]
\end{minipage}\par\medskip
\noindent\begin{minipage}{\linewidth}\textbf{19. Calculamos las multiplicaciones que ya se pueden hacer}\\[2pt]La letra y sigue sin conocerse. Las otras partes ya tienen un valor o una expresión conocida y podemos efectuar sus productos.
\par\smallskip {\small\renewcommand{\arraystretch}{1.3}\begin{tabular}{>{\raggedright\arraybackslash}p{3.36cm}>{\raggedright\arraybackslash}p{4.48cm}>{\raggedright\arraybackslash}p{3.36cm}}\hline\rowcolor{black!7} \textbf{Parte conocida} & \textbf{La multiplicación} & \textbf{Resultado}\\\hline El término 2z & $2\cdot\left(-2\right)$ & $-4$\\ \hline\end{tabular}}\par\smallskip
\par\smallskip Al reunir esas partes conocidas, la ecuación queda:
\[y - 4=-2\]
\end{minipage}\par\medskip
\noindent\begin{minipage}{\linewidth}\textbf{20. Dejamos el término con y solo a un lado}\\[2pt]Para quitar la parte conocida, sumamos 4 en AMBOS lados. Piensa en una balanza: si hacemos lo mismo a ambos lados, la igualdad se mantiene.
\par\smallskip Hacemos la misma operación a la izquierda y a la derecha:
\[y - 4 + 4=-2 + 4\]
\par\smallskip Las partes que se compensan dan cero; queda:
\[y=2\]
\par\smallskip\textbf{Qué aprendimos.} «Pasar al otro lado» es una forma abreviada de decir que hicimos la misma suma o resta en los dos lados.
\end{minipage}\par\medskip
\noindent\begin{minipage}{\linewidth}\textbf{21. Leemos el valor de y}\\[2pt]La letra ya está sola: su número multiplicador es 1, que no se escribe. La igualdad nos dice directamente cuánto vale.
\par\smallskip Anotamos el valor encontrado:
\[y=2\]
\end{minipage}\par\medskip
\noindent\begin{minipage}{\linewidth}\textbf{22. Usamos lo que ya sabemos para hallar x}\\[2pt]Ya sabemos y = 2; z = -2. Sustituir significa escribir ese valor cada vez que aparezca su letra. Usamos paréntesis para que se vea qué número estamos multiplicando.
\par\smallskip\textbf{Qué buscamos.} Dejar solamente x como letra desconocida en esta ecuación.
\par\smallskip Tomamos esta condición que obtuvimos antes:
\[x - 2y + z = -5\]
\par\smallskip Reemplazamos las letras conocidas por sus valores:
\[x - 2\left(2\right) + \left(-2\right)=-5\]
\end{minipage}\par\medskip
\noindent\begin{minipage}{\linewidth}\textbf{23. Calculamos las multiplicaciones que ya se pueden hacer}\\[2pt]La letra x sigue sin conocerse. Las otras partes ya tienen un valor o una expresión conocida y podemos efectuar sus productos.
\par\smallskip {\small\renewcommand{\arraystretch}{1.3}\begin{tabular}{>{\raggedright\arraybackslash}p{3.36cm}>{\raggedright\arraybackslash}p{4.48cm}>{\raggedright\arraybackslash}p{3.36cm}}\hline\rowcolor{black!7} \textbf{Parte conocida} & \textbf{La multiplicación} & \textbf{Resultado}\\\hline El término -2y & $-2\cdot\left(2\right)$ & $-4$\\ El término z & $1\cdot\left(-2\right)$ & $-2$\\ \hline\end{tabular}}\par\smallskip
\par\smallskip Al reunir esas partes conocidas, la ecuación queda:
\[x - 6=-5\]
\end{minipage}\par\medskip
\noindent\begin{minipage}{\linewidth}\textbf{24. Dejamos el término con x solo a un lado}\\[2pt]Para quitar la parte conocida, sumamos 6 en AMBOS lados. Piensa en una balanza: si hacemos lo mismo a ambos lados, la igualdad se mantiene.
\par\smallskip Hacemos la misma operación a la izquierda y a la derecha:
\[x - 6 + 6=-5 + 6\]
\par\smallskip Las partes que se compensan dan cero; queda:
\[x=1\]
\par\smallskip\textbf{Qué aprendimos.} «Pasar al otro lado» es una forma abreviada de decir que hicimos la misma suma o resta en los dos lados.
\end{minipage}\par\medskip
\noindent\begin{minipage}{\linewidth}\textbf{25. Leemos el valor de x}\\[2pt]La letra ya está sola: su número multiplicador es 1, que no se escribe. La igualdad nos dice directamente cuánto vale.
\par\smallskip Anotamos el valor encontrado:
\[x=1\]
\end{minipage}\par\medskip
\noindent\begin{minipage}{\linewidth}\textbf{26. Reunimos la respuesta; todavía falta comprobarla}\\[2pt]Escribimos todos los valores en el orden x, y, z. Estos valores deben funcionar juntos, no cada uno en una ecuación diferente.
\[\left\{\begin{aligned}x&=1\\y&=2\\z&=-2\end{aligned}\right.\]
\par\smallskip\textbf{Qué aprendimos.} El siguiente paso es volver a las ecuaciones del enunciado y comprobar cada condición.
\end{minipage}\par\medskip
\noindent\begin{minipage}{\linewidth}\textbf{27. ¿La respuesta cumple la ecuación 1?}\\[2pt]Comprobar significa volver a la condición original y calcularla con nuestra respuesta. Si los dos lados dan el mismo número, esta condición sí se cumple. Deben cumplirse todas.
\par\smallskip La condición del enunciado era:
\[2x - 3y + 4z = -12\]
\par\smallskip Ponemos nuestros valores donde estaban las letras:
\[2\left(1\right) - 3\left(2\right) + 4\left(-2\right)=-12\]
\par\smallskip Hacemos las multiplicaciones y luego la suma:
\[2 - 6 - 8=-12\]
\par\smallskip Comparamos los dos lados:
\[-12=-12\]
\par\smallskip\textbf{Qué aprendimos.} La igualdad final es verdadera. Nuestra respuesta sí cumple la ecuación 1.
\end{minipage}\par\medskip
\noindent\begin{minipage}{\linewidth}\textbf{28. ¿La respuesta cumple la ecuación 2?}\\[2pt]Comprobar significa volver a la condición original y calcularla con nuestra respuesta. Si los dos lados dan el mismo número, esta condición sí se cumple. Deben cumplirse todas.
\par\smallskip La condición del enunciado era:
\[x - 2y + z = -5\]
\par\smallskip Ponemos nuestros valores donde estaban las letras:
\[\left(1\right) - 2\left(2\right) + \left(-2\right)=-5\]
\par\smallskip Hacemos las multiplicaciones y luego la suma:
\[1 - 4 - 2=-5\]
\par\smallskip Comparamos los dos lados:
\[-5=-5\]
\par\smallskip\textbf{Qué aprendimos.} La igualdad final es verdadera. Nuestra respuesta sí cumple la ecuación 2.
\end{minipage}\par\medskip
\noindent\begin{minipage}{\linewidth}\textbf{29. ¿La respuesta cumple la ecuación 3?}\\[2pt]Comprobar significa volver a la condición original y calcularla con nuestra respuesta. Si los dos lados dan el mismo número, esta condición sí se cumple. Deben cumplirse todas.
\par\smallskip La condición del enunciado era:
\[3x + y + 2z = 1\]
\par\smallskip Ponemos nuestros valores donde estaban las letras:
\[3\left(1\right) + \left(2\right) + 2\left(-2\right)=1\]
\par\smallskip Hacemos las multiplicaciones y luego la suma:
\[3 + 2 - 4=1\]
\par\smallskip Comparamos los dos lados:
\[1=1\]
\par\smallskip\textbf{Qué aprendimos.} La igualdad final es verdadera. Nuestra respuesta sí cumple la ecuación 3.
\end{minipage}\par\medskip
\respuesta{x = 1; y = 2; z = -2.}
\noindent\begin{minipage}{\linewidth}\centering\includegraphics[width=.96\linewidth,height=6.8cm,keepaspectratio]{../recursos/grafico_1_1-2.png}\\{\small Lectura visual complementaria; la demostración está en las operaciones.}\end{minipage}
\end{multicols*}
\newpage\fancyfoot[L]{\small Sección 1.1 | Ejercicio 3 | Paso a paso}\hypertarget{1.1-3}{}
\pdfbookmark[1]{1.1 / 3}{bm1.1-3}
\titulo{3 | Sistema lineal por eliminación}
\textbf{Enunciado}\par Resuelva el sistema lineal dado por medio del método de eliminación.
\[\left\{\begin{aligned}3x + 2y + z &= 2\\4x + 2y + 2z &= 8\\x - y + z &= 4\end{aligned}\right.\]
\begin{multicols*}{2}
\noindent\begin{minipage}{\linewidth}\textbf{1. Primero: ¿qué significa resolver este sistema?}\\[2pt]Las letras x, y, z representan números que todavía no conocemos. Una solución es una elección de esos números que hace verdaderas TODAS las ecuaciones al mismo tiempo. La llave de la izquierda significa que las condiciones se deben cumplir juntas.
\par\smallskip\textbf{Qué buscamos.} Encontrar un mismo grupo de valores para las 3 letras.
\par\smallskip El número 3 delante de x indica tres veces ese mismo valor; es decir, sumar tres copias:
\[3x=x+x+x\]
\par\smallskip Ejemplo numérico de la notación: si x valiera 2, tres veces ese valor sería 6. Este 2 es solo un ejemplo; todavía no es la respuesta del ejercicio. El punto centrado indica multiplicación:
\[3\cdot 2=2+2+2=6\]
\par\smallskip\textbf{Comprueba que lo entiendes.} ¿Basta con que los valores funcionen en una sola ecuación?\par\textit{Respuesta:} No. Deben funcionar en todas las ecuaciones del sistema, usando los mismos valores.
\end{minipage}\par\medskip
\noindent\begin{minipage}{\linewidth}\textbf{2. Nuestro plan: ir quitando letras}\\[2pt]Una ecuación con una sola letra es más fácil de resolver. Vamos a combinar ecuaciones para que una letra desaparezca. Cuando sepamos un valor, lo pondremos en las ecuaciones que faltan.
\par\smallskip\textbf{Qué buscamos.} Primero quitar x de las otras ecuaciones. Después seguir con las letras restantes.
\par\smallskip\textbf{Por qué.} Si una cantidad se resta de sí misma, queda cero. Usaremos esta idea con términos que tengan la misma letra.
\par\smallskip Ejemplo de la idea que utilizaremos:
\[3x-3x=0\]
\par\smallskip\textbf{Qué aprendimos.} Eliminar una letra de una ecuación NO significa que su valor sea cero; significa que sus términos se cancelaron en esa cuenta.
\end{minipage}\par\medskip
\noindent\begin{minipage}{\linewidth}\textbf{3. Ponemos una ecuación más cómoda arriba}\\[2pt]Cambiamos de lugar las ecuaciones 1 y 3. Las condiciones siguen siendo las mismas; solamente cambia el orden en que las leemos.
\par\smallskip\textbf{Por qué.} La ecuación que subimos tiene un número más pequeño delante de x. Eso facilita las multiplicaciones que haremos a continuación.
\[\left\{\begin{aligned}x - y + z &= 4\\4x + 2y + 2z &= 8\\3x + 2y + z &= 2\end{aligned}\right.\]
\end{minipage}\par\medskip
\noindent\begin{minipage}{\linewidth}\textbf{4. ¿Cómo haremos desaparecer x?}\\[2pt]Trabajaremos con las ecuaciones 1 y 2. Queremos que los términos con x tengan el mismo tamaño. Así podremos restarlos.
\par\smallskip\textbf{Qué buscamos.} Conseguir términos con 4 veces x en las dos ecuaciones.
\par\smallskip\textbf{Por qué.} Tener 4 veces la misma letra en ambos lados de la resta nos permitirá cancelarla.
\par\smallskip {\small\renewcommand{\arraystretch}{1.3}\begin{tabular}{>{\raggedright\arraybackslash}p{2.58cm}>{\raggedright\arraybackslash}p{2.47cm}>{\raggedright\arraybackslash}p{6.17cm}}\hline\rowcolor{black!7} \textbf{Dónde miramos} & \textbf{Qué término hay} & \textbf{Qué necesitamos hacer}\\\hline Ecuación 1 & $x$ & Multiplicar TODA esta ecuación por 4\\ Ecuación 2 & $4x$ & La dejamos igual.\\ \hline\end{tabular}}\par\smallskip
\par\smallskip\textbf{Qué aprendimos.} Estos multiplicadores se eligen para igualar los términos con x. Los demás términos también se multiplican para conservar la igualdad.
\end{minipage}\par\medskip
\noindent\begin{minipage}{\linewidth}\textbf{5. Multiplicamos la ecuación 1 por 4}\\[2pt]Multiplicar una ecuación por 4 significa multiplicar CADA término de la izquierda y también el número de la derecha. La igualdad se conserva porque hacemos lo mismo en ambos lados.
\par\smallskip Esta es la ecuación antes de multiplicar:
\[x - y + z = 4\]
\par\smallskip {\small\renewcommand{\arraystretch}{1.3}\begin{tabular}{>{\raggedright\arraybackslash}p{2.80cm}>{\raggedright\arraybackslash}p{1.91cm}>{\raggedright\arraybackslash}p{6.50cm}}\hline\rowcolor{black!7} \textbf{Parte de la ecuación} & \textbf{Antes} & \textbf{La cuenta y su resultado}\\\hline Término con x & $x$ & $4\cdot\left(x\right)=4x$\\ Término con y & $-y$ & $4\cdot\left(-y\right)=-4y$\\ Término con z & $z$ & $4\cdot\left(z\right)=4z$\\ Número a la derecha del igual & $4$ & $4\cdot\left(4\right)=16$\\ \hline\end{tabular}}\par\smallskip
\par\smallskip Después de hacer todas las multiplicaciones queda:
\[4x - 4y + 4z = 16\]
\par\smallskip\textbf{Qué aprendimos.} La misma condición ahora está escrita con los términos multiplicados. Ya tenemos el tamaño que buscábamos para x.
\par\smallskip\textbf{Comprueba que lo entiendes.} ¿Podríamos multiplicar solamente el término con x?\par\textit{Respuesta:} No. Eso cambiaría la condición. Debemos multiplicar todos los términos, también el número que está después del igual.
\end{minipage}\par\medskip
\noindent\begin{minipage}{\linewidth}\textbf{6. Restamos las dos ecuaciones preparadas}\\[2pt]Restamos la ecuación 1 preparada a la ecuación 2 preparada. Juntamos cada término solo con otro que tenga la misma letra. Los números sin letra se calculan por separado.
\par\smallskip {\small\renewcommand{\arraystretch}{1.3}\begin{tabular}{>{\raggedright\arraybackslash}p{2.80cm}>{\raggedright\arraybackslash}p{5.38cm}>{\raggedright\arraybackslash}p{3.03cm}}\hline\rowcolor{black!7} \textbf{Qué agrupamos} & \textbf{Cuenta} & \textbf{Qué significa}\\\hline Los términos con x & $4x - \left(4x\right)=0$ & Se cancela esta letra.\\ Los términos con y & $2y - \left(-4y\right)=6y$ & Quedan 6y.\\ Los términos con z & $2z - \left(4z\right)=-2z$ & Quedan -2z.\\ Los números a la derecha del igual & $8 - \left(16\right)=-8$ & Este será el número de la nueva igualdad.\\ \hline\end{tabular}}\par\smallskip
\par\smallskip Primera igualdad que usaremos:
\[4x + 2y + 2z = 8\]
\par\smallskip Le restamos esta igualdad:
\[4x - 4y + 4z = 16\]
\par\smallskip\textbf{Qué aprendimos.} Los términos con x dan cero. Los términos con letras distintas no se mezclan.
\end{minipage}\par\medskip
\noindent\begin{minipage}{\linewidth}\textbf{7. Mira lo que hemos conseguido}\\[2pt]La nueva ecuación es el resultado de la cuenta anterior. Todavía contiene y, z. Seguiremos reduciendo las condiciones.
\par\smallskip Escribimos la nueva condición:
\[6y - 2z = -8\]
\[\left\{\begin{aligned}x - y + z &= 4\\6y - 2z &= -8\\3x + 2y + z &= 2\end{aligned}\right.\]
\par\smallskip\textbf{Qué aprendimos.} Conservamos la ecuación 1 y reemplazamos la 2 por esta nueva condición. De esa manera conservamos toda la información necesaria para resolver el sistema.
\end{minipage}\par\medskip
\noindent\begin{minipage}{\linewidth}\textbf{8. Escribimos la misma condición con números más pequeños}\\[2pt]Todos los números de esta ecuación se pueden dividir entre 2 sin dejar restos. Dividimos ambos lados entre ese mismo número. Así la condición no cambia y las siguientes cuentas serán más sencillas.
\par\smallskip Antes de simplificar:
\[6y - 2z = -8\]
\par\smallskip Después de dividir todos los términos entre 2:
\[3y - z = -4\]
\end{minipage}\par\medskip
\noindent\begin{minipage}{\linewidth}\textbf{9. ¿Cómo haremos desaparecer x?}\\[2pt]Trabajaremos con las ecuaciones 1 y 3. Queremos que los términos con x tengan el mismo tamaño. Así podremos restarlos.
\par\smallskip\textbf{Qué buscamos.} Conseguir términos con 3 veces x en las dos ecuaciones.
\par\smallskip\textbf{Por qué.} Tener 3 veces la misma letra en ambos lados de la resta nos permitirá cancelarla.
\par\smallskip {\small\renewcommand{\arraystretch}{1.3}\begin{tabular}{>{\raggedright\arraybackslash}p{2.58cm}>{\raggedright\arraybackslash}p{2.47cm}>{\raggedright\arraybackslash}p{6.17cm}}\hline\rowcolor{black!7} \textbf{Dónde miramos} & \textbf{Qué término hay} & \textbf{Qué necesitamos hacer}\\\hline Ecuación 1 & $x$ & Multiplicar TODA esta ecuación por 3\\ Ecuación 3 & $3x$ & La dejamos igual.\\ \hline\end{tabular}}\par\smallskip
\par\smallskip\textbf{Qué aprendimos.} Estos multiplicadores se eligen para igualar los términos con x. Los demás términos también se multiplican para conservar la igualdad.
\end{minipage}\par\medskip
\noindent\begin{minipage}{\linewidth}\textbf{10. Multiplicamos la ecuación 1 por 3}\\[2pt]Multiplicar una ecuación por 3 significa multiplicar CADA término de la izquierda y también el número de la derecha. La igualdad se conserva porque hacemos lo mismo en ambos lados.
\par\smallskip Esta es la ecuación antes de multiplicar:
\[x - y + z = 4\]
\par\smallskip {\small\renewcommand{\arraystretch}{1.3}\begin{tabular}{>{\raggedright\arraybackslash}p{2.80cm}>{\raggedright\arraybackslash}p{1.91cm}>{\raggedright\arraybackslash}p{6.50cm}}\hline\rowcolor{black!7} \textbf{Parte de la ecuación} & \textbf{Antes} & \textbf{La cuenta y su resultado}\\\hline Término con x & $x$ & $3\cdot\left(x\right)=3x$\\ Término con y & $-y$ & $3\cdot\left(-y\right)=-3y$\\ Término con z & $z$ & $3\cdot\left(z\right)=3z$\\ Número a la derecha del igual & $4$ & $3\cdot\left(4\right)=12$\\ \hline\end{tabular}}\par\smallskip
\par\smallskip Después de hacer todas las multiplicaciones queda:
\[3x - 3y + 3z = 12\]
\par\smallskip\textbf{Qué aprendimos.} La misma condición ahora está escrita con los términos multiplicados. Ya tenemos el tamaño que buscábamos para x.
\par\smallskip\textbf{Comprueba que lo entiendes.} ¿Podríamos multiplicar solamente el término con x?\par\textit{Respuesta:} No. Eso cambiaría la condición. Debemos multiplicar todos los términos, también el número que está después del igual.
\end{minipage}\par\medskip
\noindent\begin{minipage}{\linewidth}\textbf{11. Restamos las dos ecuaciones preparadas}\\[2pt]Restamos la ecuación 1 preparada a la ecuación 3 preparada. Juntamos cada término solo con otro que tenga la misma letra. Los números sin letra se calculan por separado.
\par\smallskip {\small\renewcommand{\arraystretch}{1.3}\begin{tabular}{>{\raggedright\arraybackslash}p{2.80cm}>{\raggedright\arraybackslash}p{5.38cm}>{\raggedright\arraybackslash}p{3.03cm}}\hline\rowcolor{black!7} \textbf{Qué agrupamos} & \textbf{Cuenta} & \textbf{Qué significa}\\\hline Los términos con x & $3x - \left(3x\right)=0$ & Se cancela esta letra.\\ Los términos con y & $2y - \left(-3y\right)=5y$ & Quedan 5y.\\ Los términos con z & $z - \left(3z\right)=-2z$ & Quedan -2z.\\ Los números a la derecha del igual & $2 - \left(12\right)=-10$ & Este será el número de la nueva igualdad.\\ \hline\end{tabular}}\par\smallskip
\par\smallskip Primera igualdad que usaremos:
\[3x + 2y + z = 2\]
\par\smallskip Le restamos esta igualdad:
\[3x - 3y + 3z = 12\]
\par\smallskip\textbf{Qué aprendimos.} Los términos con x dan cero. Los términos con letras distintas no se mezclan.
\end{minipage}\par\medskip
\noindent\begin{minipage}{\linewidth}\textbf{12. Mira lo que hemos conseguido}\\[2pt]La nueva ecuación es el resultado de la cuenta anterior. Todavía contiene y, z. Seguiremos reduciendo las condiciones.
\par\smallskip Escribimos la nueva condición:
\[5y - 2z = -10\]
\[\left\{\begin{aligned}x - y + z &= 4\\3y - z &= -4\\5y - 2z &= -10\end{aligned}\right.\]
\par\smallskip\textbf{Qué aprendimos.} Conservamos la ecuación 1 y reemplazamos la 3 por esta nueva condición. De esa manera conservamos toda la información necesaria para resolver el sistema.
\end{minipage}\par\medskip
\noindent\begin{minipage}{\linewidth}\textbf{13. ¿Cómo haremos desaparecer y?}\\[2pt]Trabajaremos con las ecuaciones 2 y 3. Queremos que los términos con y tengan el mismo tamaño. Así podremos restarlos.
\par\smallskip\textbf{Qué buscamos.} Conseguir términos con 15 veces y en las dos ecuaciones.
\par\smallskip\textbf{Por qué.} Tener 15 veces la misma letra en ambos lados de la resta nos permitirá cancelarla.
\par\smallskip {\small\renewcommand{\arraystretch}{1.3}\begin{tabular}{>{\raggedright\arraybackslash}p{2.58cm}>{\raggedright\arraybackslash}p{2.47cm}>{\raggedright\arraybackslash}p{6.17cm}}\hline\rowcolor{black!7} \textbf{Dónde miramos} & \textbf{Qué término hay} & \textbf{Qué necesitamos hacer}\\\hline Ecuación 2 & $3y$ & Multiplicar TODA esta ecuación por 5\\ Ecuación 3 & $5y$ & Multiplicar TODA esta ecuación por 3\\ \hline\end{tabular}}\par\smallskip
\par\smallskip\textbf{Qué aprendimos.} Estos multiplicadores se eligen para igualar los términos con y. Los demás términos también se multiplican para conservar la igualdad.
\end{minipage}\par\medskip
\noindent\begin{minipage}{\linewidth}\textbf{14. Multiplicamos la ecuación 2 por 5}\\[2pt]Multiplicar una ecuación por 5 significa multiplicar CADA término de la izquierda y también el número de la derecha. La igualdad se conserva porque hacemos lo mismo en ambos lados.
\par\smallskip Esta es la ecuación antes de multiplicar:
\[3y - z = -4\]
\par\smallskip {\small\renewcommand{\arraystretch}{1.3}\begin{tabular}{>{\raggedright\arraybackslash}p{2.80cm}>{\raggedright\arraybackslash}p{1.91cm}>{\raggedright\arraybackslash}p{6.50cm}}\hline\rowcolor{black!7} \textbf{Parte de la ecuación} & \textbf{Antes} & \textbf{La cuenta y su resultado}\\\hline Término con y & $3y$ & $5\cdot\left(3y\right)=15y$\\ Término con z & $-z$ & $5\cdot\left(-z\right)=-5z$\\ Número a la derecha del igual & $-4$ & $5\cdot\left(-4\right)=-20$\\ \hline\end{tabular}}\par\smallskip
\par\smallskip Después de hacer todas las multiplicaciones queda:
\[15y - 5z = -20\]
\par\smallskip\textbf{Qué aprendimos.} La misma condición ahora está escrita con los términos multiplicados. Ya tenemos el tamaño que buscábamos para y.
\par\smallskip\textbf{Comprueba que lo entiendes.} ¿Podríamos multiplicar solamente el término con y?\par\textit{Respuesta:} No. Eso cambiaría la condición. Debemos multiplicar todos los términos, también el número que está después del igual.
\end{minipage}\par\medskip
\noindent\begin{minipage}{\linewidth}\textbf{15. Multiplicamos la ecuación 3 por 3}\\[2pt]Multiplicar una ecuación por 3 significa multiplicar CADA término de la izquierda y también el número de la derecha. La igualdad se conserva porque hacemos lo mismo en ambos lados.
\par\smallskip Esta es la ecuación antes de multiplicar:
\[5y - 2z = -10\]
\par\smallskip {\small\renewcommand{\arraystretch}{1.3}\begin{tabular}{>{\raggedright\arraybackslash}p{2.80cm}>{\raggedright\arraybackslash}p{1.91cm}>{\raggedright\arraybackslash}p{6.50cm}}\hline\rowcolor{black!7} \textbf{Parte de la ecuación} & \textbf{Antes} & \textbf{La cuenta y su resultado}\\\hline Término con y & $5y$ & $3\cdot\left(5y\right)=15y$\\ Término con z & $-2z$ & $3\cdot\left(-2z\right)=-6z$\\ Número a la derecha del igual & $-10$ & $3\cdot\left(-10\right)=-30$\\ \hline\end{tabular}}\par\smallskip
\par\smallskip Después de hacer todas las multiplicaciones queda:
\[15y - 6z = -30\]
\par\smallskip\textbf{Qué aprendimos.} La misma condición ahora está escrita con los términos multiplicados. Ya tenemos el tamaño que buscábamos para y.
\par\smallskip\textbf{Comprueba que lo entiendes.} ¿Podríamos multiplicar solamente el término con y?\par\textit{Respuesta:} No. Eso cambiaría la condición. Debemos multiplicar todos los términos, también el número que está después del igual.
\end{minipage}\par\medskip
\noindent\begin{minipage}{\linewidth}\textbf{16. Restamos las dos ecuaciones preparadas}\\[2pt]Restamos la ecuación 2 preparada a la ecuación 3 preparada. Juntamos cada término solo con otro que tenga la misma letra. Los números sin letra se calculan por separado.
\par\smallskip {\small\renewcommand{\arraystretch}{1.3}\begin{tabular}{>{\raggedright\arraybackslash}p{2.80cm}>{\raggedright\arraybackslash}p{5.38cm}>{\raggedright\arraybackslash}p{3.03cm}}\hline\rowcolor{black!7} \textbf{Qué agrupamos} & \textbf{Cuenta} & \textbf{Qué significa}\\\hline Los términos con y & $15y - \left(15y\right)=0$ & Se cancela esta letra.\\ Los términos con z & $-6z - \left(-5z\right)=-z$ & Quedan -z.\\ Los números a la derecha del igual & $-30 - \left(-20\right)=-30+20=-10$ & Este será el número de la nueva igualdad.\\ \hline\end{tabular}}\par\smallskip
\par\smallskip Primera igualdad que usaremos:
\[15y - 6z = -30\]
\par\smallskip Le restamos esta igualdad:
\[15y - 5z = -20\]
\par\smallskip\textbf{Qué aprendimos.} Los términos con y dan cero. Recuerda: restar un número negativo equivale a sumar su opuesto.
\end{minipage}\par\medskip
\noindent\begin{minipage}{\linewidth}\textbf{17. Mira lo que hemos conseguido}\\[2pt]La nueva ecuación es el resultado de la cuenta anterior. Solo contiene z: después podremos despejarla.
\par\smallskip Escribimos la nueva condición:
\[-z = -10\]
\[\left\{\begin{aligned}x - y + z &= 4\\3y - z &= -4\\-z &= -10\end{aligned}\right.\]
\par\smallskip\textbf{Qué aprendimos.} Conservamos la ecuación 2 y reemplazamos la 3 por esta nueva condición. De esa manera conservamos toda la información necesaria para resolver el sistema.
\end{minipage}\par\medskip
\noindent\begin{minipage}{\linewidth}\textbf{18. Ya tenemos una ecuación con una sola letra}\\[2pt]En esta condición solamente falta conocer z. El número -1 está multiplicando a esa letra. Ahora podemos deshacer esa multiplicación.
\par\smallskip Esta es la condición que vamos a resolver:
\[-z = -10\]
\end{minipage}\par\medskip
\noindent\begin{minipage}{\linewidth}\textbf{19. Calculamos cuánto vale z}\\[2pt]Para deshacer una multiplicación por -1, dividimos AMBOS lados entre -1. En el lado izquierdo queda una sola z. En el derecho hacemos la división.
\par\smallskip Dividimos los dos lados entre el mismo número:
\[\frac{-z}{-1}=\frac{-10}{-1}\]
\par\smallskip Ya podemos leer el valor de la letra:
\[z=\frac{-10}{-1}=10\]
\par\smallskip\textbf{Qué aprendimos.} Un número negativo dividido entre otro negativo da positivo.
\end{minipage}\par\medskip
\noindent\begin{minipage}{\linewidth}\textbf{20. Usamos lo que ya sabemos para hallar y}\\[2pt]Ya sabemos z = 10. Sustituir significa escribir ese valor cada vez que aparezca su letra. Usamos paréntesis para que se vea qué número estamos multiplicando.
\par\smallskip\textbf{Qué buscamos.} Dejar solamente y como letra desconocida en esta ecuación.
\par\smallskip Tomamos esta condición que obtuvimos antes:
\[3y - z = -4\]
\par\smallskip Reemplazamos las letras conocidas por sus valores:
\[3y - \left(10\right)=-4\]
\end{minipage}\par\medskip
\noindent\begin{minipage}{\linewidth}\textbf{21. Calculamos las multiplicaciones que ya se pueden hacer}\\[2pt]La letra y sigue sin conocerse. Las otras partes ya tienen un valor o una expresión conocida y podemos efectuar sus productos.
\par\smallskip {\small\renewcommand{\arraystretch}{1.3}\begin{tabular}{>{\raggedright\arraybackslash}p{3.36cm}>{\raggedright\arraybackslash}p{4.48cm}>{\raggedright\arraybackslash}p{3.36cm}}\hline\rowcolor{black!7} \textbf{Parte conocida} & \textbf{La multiplicación} & \textbf{Resultado}\\\hline El término -z & $-1\cdot\left(10\right)$ & $-10$\\ \hline\end{tabular}}\par\smallskip
\par\smallskip Al reunir esas partes conocidas, la ecuación queda:
\[3y - 10=-4\]
\end{minipage}\par\medskip
\noindent\begin{minipage}{\linewidth}\textbf{22. Dejamos el término con y solo a un lado}\\[2pt]Para quitar la parte conocida, sumamos 10 en AMBOS lados. Piensa en una balanza: si hacemos lo mismo a ambos lados, la igualdad se mantiene.
\par\smallskip Hacemos la misma operación a la izquierda y a la derecha:
\[3y - 10 + 10=-4 + 10\]
\par\smallskip Las partes que se compensan dan cero; queda:
\[3y=6\]
\par\smallskip\textbf{Qué aprendimos.} «Pasar al otro lado» es una forma abreviada de decir que hicimos la misma suma o resta en los dos lados.
\end{minipage}\par\medskip
\noindent\begin{minipage}{\linewidth}\textbf{23. Calculamos cuánto vale y}\\[2pt]Para deshacer una multiplicación por 3, dividimos AMBOS lados entre 3. En el lado izquierdo queda una sola y. En el derecho hacemos la división.
\par\smallskip Dividimos los dos lados entre el mismo número:
\[\frac{3y}{3}=\frac{6}{3}\]
\par\smallskip Ya podemos leer el valor de la letra:
\[y=\frac{6}{3}=2\]
\par\smallskip\textbf{Qué aprendimos.} Este valor se usará en las ecuaciones donde todavía falten letras.
\end{minipage}\par\medskip
\noindent\begin{minipage}{\linewidth}\textbf{24. Usamos lo que ya sabemos para hallar x}\\[2pt]Ya sabemos y = 2; z = 10. Sustituir significa escribir ese valor cada vez que aparezca su letra. Usamos paréntesis para que se vea qué número estamos multiplicando.
\par\smallskip\textbf{Qué buscamos.} Dejar solamente x como letra desconocida en esta ecuación.
\par\smallskip Tomamos esta condición que obtuvimos antes:
\[x - y + z = 4\]
\par\smallskip Reemplazamos las letras conocidas por sus valores:
\[x - \left(2\right) + \left(10\right)=4\]
\end{minipage}\par\medskip
\noindent\begin{minipage}{\linewidth}\textbf{25. Calculamos las multiplicaciones que ya se pueden hacer}\\[2pt]La letra x sigue sin conocerse. Las otras partes ya tienen un valor o una expresión conocida y podemos efectuar sus productos.
\par\smallskip {\small\renewcommand{\arraystretch}{1.3}\begin{tabular}{>{\raggedright\arraybackslash}p{3.36cm}>{\raggedright\arraybackslash}p{4.48cm}>{\raggedright\arraybackslash}p{3.36cm}}\hline\rowcolor{black!7} \textbf{Parte conocida} & \textbf{La multiplicación} & \textbf{Resultado}\\\hline El término -y & $-1\cdot\left(2\right)$ & $-2$\\ El término z & $1\cdot\left(10\right)$ & $10$\\ \hline\end{tabular}}\par\smallskip
\par\smallskip Al reunir esas partes conocidas, la ecuación queda:
\[x + 8=4\]
\end{minipage}\par\medskip
\noindent\begin{minipage}{\linewidth}\textbf{26. Dejamos el término con x solo a un lado}\\[2pt]Para quitar la parte conocida, restamos 8 en AMBOS lados. Piensa en una balanza: si hacemos lo mismo a ambos lados, la igualdad se mantiene.
\par\smallskip Hacemos la misma operación a la izquierda y a la derecha:
\[x + 8 - 8=4 - 8\]
\par\smallskip Las partes que se compensan dan cero; queda:
\[x=-4\]
\par\smallskip\textbf{Qué aprendimos.} «Pasar al otro lado» es una forma abreviada de decir que hicimos la misma suma o resta en los dos lados.
\end{minipage}\par\medskip
\noindent\begin{minipage}{\linewidth}\textbf{27. Leemos el valor de x}\\[2pt]La letra ya está sola: su número multiplicador es 1, que no se escribe. La igualdad nos dice directamente cuánto vale.
\par\smallskip Anotamos el valor encontrado:
\[x=-4\]
\end{minipage}\par\medskip
\noindent\begin{minipage}{\linewidth}\textbf{28. Reunimos la respuesta; todavía falta comprobarla}\\[2pt]Escribimos todos los valores en el orden x, y, z. Estos valores deben funcionar juntos, no cada uno en una ecuación diferente.
\[\left\{\begin{aligned}x&=-4\\y&=2\\z&=10\end{aligned}\right.\]
\par\smallskip\textbf{Qué aprendimos.} El siguiente paso es volver a las ecuaciones del enunciado y comprobar cada condición.
\end{minipage}\par\medskip
\noindent\begin{minipage}{\linewidth}\textbf{29. ¿La respuesta cumple la ecuación 1?}\\[2pt]Comprobar significa volver a la condición original y calcularla con nuestra respuesta. Si los dos lados dan el mismo número, esta condición sí se cumple. Deben cumplirse todas.
\par\smallskip La condición del enunciado era:
\[3x + 2y + z = 2\]
\par\smallskip Ponemos nuestros valores donde estaban las letras:
\[3\left(-4\right) + 2\left(2\right) + \left(10\right)=2\]
\par\smallskip Hacemos las multiplicaciones y luego la suma:
\[-12 + 4 + 10=2\]
\par\smallskip Comparamos los dos lados:
\[2=2\]
\par\smallskip\textbf{Qué aprendimos.} La igualdad final es verdadera. Nuestra respuesta sí cumple la ecuación 1.
\end{minipage}\par\medskip
\noindent\begin{minipage}{\linewidth}\textbf{30. ¿La respuesta cumple la ecuación 2?}\\[2pt]Comprobar significa volver a la condición original y calcularla con nuestra respuesta. Si los dos lados dan el mismo número, esta condición sí se cumple. Deben cumplirse todas.
\par\smallskip La condición del enunciado era:
\[4x + 2y + 2z = 8\]
\par\smallskip Ponemos nuestros valores donde estaban las letras:
\[4\left(-4\right) + 2\left(2\right) + 2\left(10\right)=8\]
\par\smallskip Hacemos las multiplicaciones y luego la suma:
\[-16 + 4 + 20=8\]
\par\smallskip Comparamos los dos lados:
\[8=8\]
\par\smallskip\textbf{Qué aprendimos.} La igualdad final es verdadera. Nuestra respuesta sí cumple la ecuación 2.
\end{minipage}\par\medskip
\noindent\begin{minipage}{\linewidth}\textbf{31. ¿La respuesta cumple la ecuación 3?}\\[2pt]Comprobar significa volver a la condición original y calcularla con nuestra respuesta. Si los dos lados dan el mismo número, esta condición sí se cumple. Deben cumplirse todas.
\par\smallskip La condición del enunciado era:
\[x - y + z = 4\]
\par\smallskip Ponemos nuestros valores donde estaban las letras:
\[\left(-4\right) - \left(2\right) + \left(10\right)=4\]
\par\smallskip Hacemos las multiplicaciones y luego la suma:
\[-4 - 2 + 10=4\]
\par\smallskip Comparamos los dos lados:
\[4=4\]
\par\smallskip\textbf{Qué aprendimos.} La igualdad final es verdadera. Nuestra respuesta sí cumple la ecuación 3.
\end{minipage}\par\medskip
\respuesta{x = -4; y = 2; z = 10.}
\noindent\begin{minipage}{\linewidth}\centering\includegraphics[width=.96\linewidth,height=6.8cm,keepaspectratio]{../recursos/grafico_1_1-3.png}\\{\small Lectura visual complementaria; la demostración está en las operaciones.}\end{minipage}
\end{multicols*}
\newpage\fancyfoot[L]{\small Sección 1.1 | Ejercicio 4 | Paso a paso}\hypertarget{1.1-4}{}
\pdfbookmark[1]{1.1 / 4}{bm1.1-4}
\titulo{4 | Sistema lineal por eliminación}
\textbf{Enunciado}\par Resuelva el sistema lineal dado por medio del método de eliminación.
\[\left\{\begin{aligned}x + y &= 5\\3x + 3y &= 10\end{aligned}\right.\]
\begin{multicols*}{2}
\noindent\begin{minipage}{\linewidth}\textbf{1. Primero: ¿qué significa resolver este sistema?}\\[2pt]Las letras x, y representan números que todavía no conocemos. Una solución es una elección de esos números que hace verdaderas TODAS las ecuaciones al mismo tiempo. La llave de la izquierda significa que las condiciones se deben cumplir juntas.
\par\smallskip\textbf{Qué buscamos.} Encontrar una misma pareja de valores.
\par\smallskip El número 3 delante de x indica tres veces ese mismo valor; es decir, sumar tres copias:
\[3x=x+x+x\]
\par\smallskip Ejemplo numérico de la notación: si x valiera 2, tres veces ese valor sería 6. Este 2 es solo un ejemplo; todavía no es la respuesta del ejercicio. El punto centrado indica multiplicación:
\[3\cdot 2=2+2+2=6\]
\par\smallskip\textbf{Comprueba que lo entiendes.} ¿Basta con que los valores funcionen en una sola ecuación?\par\textit{Respuesta:} No. Deben funcionar en todas las ecuaciones del sistema, usando los mismos valores.
\end{minipage}\par\medskip
\noindent\begin{minipage}{\linewidth}\textbf{2. Nuestro plan: ir quitando letras}\\[2pt]Una ecuación con una sola letra es más fácil de resolver. Vamos a combinar ecuaciones para que una letra desaparezca. Cuando sepamos un valor, lo pondremos en las ecuaciones que faltan.
\par\smallskip\textbf{Qué buscamos.} Primero quitar x para calcular y. Después calcular x.
\par\smallskip\textbf{Por qué.} Si una cantidad se resta de sí misma, queda cero. Usaremos esta idea con términos que tengan la misma letra.
\par\smallskip Ejemplo de la idea que utilizaremos:
\[3x-3x=0\]
\par\smallskip\textbf{Qué aprendimos.} Eliminar una letra de una ecuación NO significa que su valor sea cero; significa que sus términos se cancelaron en esa cuenta.
\end{minipage}\par\medskip
\noindent\begin{minipage}{\linewidth}\textbf{3. ¿Cómo haremos desaparecer x?}\\[2pt]Trabajaremos con las ecuaciones 1 y 2. Queremos que los términos con x tengan el mismo tamaño. Así podremos restarlos.
\par\smallskip\textbf{Qué buscamos.} Conseguir términos con 3 veces x en las dos ecuaciones.
\par\smallskip\textbf{Por qué.} Tener 3 veces la misma letra en ambos lados de la resta nos permitirá cancelarla.
\par\smallskip {\small\renewcommand{\arraystretch}{1.3}\begin{tabular}{>{\raggedright\arraybackslash}p{2.58cm}>{\raggedright\arraybackslash}p{2.47cm}>{\raggedright\arraybackslash}p{6.17cm}}\hline\rowcolor{black!7} \textbf{Dónde miramos} & \textbf{Qué término hay} & \textbf{Qué necesitamos hacer}\\\hline Ecuación 1 & $x$ & Multiplicar TODA esta ecuación por 3\\ Ecuación 2 & $3x$ & La dejamos igual.\\ \hline\end{tabular}}\par\smallskip
\par\smallskip\textbf{Qué aprendimos.} Estos multiplicadores se eligen para igualar los términos con x. Los demás términos también se multiplican para conservar la igualdad.
\end{minipage}\par\medskip
\noindent\begin{minipage}{\linewidth}\textbf{4. Multiplicamos la ecuación 1 por 3}\\[2pt]Multiplicar una ecuación por 3 significa multiplicar CADA término de la izquierda y también el número de la derecha. La igualdad se conserva porque hacemos lo mismo en ambos lados.
\par\smallskip Esta es la ecuación antes de multiplicar:
\[x + y = 5\]
\par\smallskip {\small\renewcommand{\arraystretch}{1.3}\begin{tabular}{>{\raggedright\arraybackslash}p{2.80cm}>{\raggedright\arraybackslash}p{1.91cm}>{\raggedright\arraybackslash}p{6.50cm}}\hline\rowcolor{black!7} \textbf{Parte de la ecuación} & \textbf{Antes} & \textbf{La cuenta y su resultado}\\\hline Término con x & $x$ & $3\cdot\left(x\right)=3x$\\ Término con y & $y$ & $3\cdot\left(y\right)=3y$\\ Número a la derecha del igual & $5$ & $3\cdot\left(5\right)=15$\\ \hline\end{tabular}}\par\smallskip
\par\smallskip Después de hacer todas las multiplicaciones queda:
\[3x + 3y = 15\]
\par\smallskip\textbf{Qué aprendimos.} La misma condición ahora está escrita con los términos multiplicados. Ya tenemos el tamaño que buscábamos para x.
\par\smallskip\textbf{Comprueba que lo entiendes.} ¿Podríamos multiplicar solamente el término con x?\par\textit{Respuesta:} No. Eso cambiaría la condición. Debemos multiplicar todos los términos, también el número que está después del igual.
\end{minipage}\par\medskip
\noindent\begin{minipage}{\linewidth}\textbf{5. Restamos las dos ecuaciones preparadas}\\[2pt]Restamos la ecuación 1 preparada a la ecuación 2 preparada. Juntamos cada término solo con otro que tenga la misma letra. Los números sin letra se calculan por separado.
\par\smallskip {\small\renewcommand{\arraystretch}{1.3}\begin{tabular}{>{\raggedright\arraybackslash}p{2.80cm}>{\raggedright\arraybackslash}p{5.38cm}>{\raggedright\arraybackslash}p{3.03cm}}\hline\rowcolor{black!7} \textbf{Qué agrupamos} & \textbf{Cuenta} & \textbf{Qué significa}\\\hline Los términos con x & $3x - \left(3x\right)=0$ & Se cancela esta letra.\\ Los términos con y & $3y - \left(3y\right)=0$ & Se cancela esta letra.\\ Los números a la derecha del igual & $10 - \left(15\right)=-5$ & Este será el número de la nueva igualdad.\\ \hline\end{tabular}}\par\smallskip
\par\smallskip Primera igualdad que usaremos:
\[3x + 3y = 10\]
\par\smallskip Le restamos esta igualdad:
\[3x + 3y = 15\]
\par\smallskip\textbf{Qué aprendimos.} Los términos con x dan cero. Los términos con letras distintas no se mezclan.
\end{minipage}\par\medskip
\noindent\begin{minipage}{\linewidth}\textbf{6. Mira lo que hemos conseguido}\\[2pt]La nueva ecuación es el resultado de la cuenta anterior. Ya no contiene letras. Tenemos que decidir si la igualdad numérica es verdadera o imposible.
\par\smallskip Escribimos la nueva condición:
\[0 = -5\]
\[\left\{\begin{aligned}x + y &= 5\\0 &= -5\end{aligned}\right.\]
\par\smallskip\textbf{Qué aprendimos.} Conservamos la ecuación 1 y reemplazamos la 2 por esta nueva condición. De esa manera conservamos toda la información necesaria para resolver el sistema.
\end{minipage}\par\medskip
\noindent\begin{minipage}{\linewidth}\textbf{7. Esta igualdad nos dice que no hay solución}\\[2pt]Al hacer las cuentas, todas las letras desaparecieron y quedó la igualdad que ves abajo. Cero no puede ser igual a -5. Como esa condición tendría que cumplirse junto con las demás, ningún grupo de valores puede resolver el sistema.
\par\smallskip Esta es la contradicción encontrada:
\[0 = -5\]
\par\smallskip\textbf{Qué aprendimos.} El resultado es «sin solución». No significa que debamos responder con ceros: significa que no existe una respuesta que cumpla todas las condiciones.
\par\smallskip\textbf{Comprueba que lo entiendes.} ¿Responder x = 0, y = 0 resolvería esta contradicción?\par\textit{Respuesta:} No. La contradicción ya no depende de las letras. Ningún valor puede convertir cero en otro número.
\end{minipage}\par\medskip
\respuesta{Sin solución.}
\noindent\begin{minipage}{\linewidth}\centering\includegraphics[width=.96\linewidth,height=6.8cm,keepaspectratio]{../recursos/grafico_1_1-4.png}\\{\small Lectura visual complementaria; la demostración está en las operaciones.}\end{minipage}
\end{multicols*}
\newpage\fancyfoot[L]{\small Sección 1.1 | Ejercicio 5 | Paso a paso}\hypertarget{1.1-5}{}
\pdfbookmark[1]{1.1 / 5}{bm1.1-5}
\titulo{5 | Sistema lineal por eliminación}
\textbf{Enunciado}\par Resuelva el sistema lineal dado por medio del método de eliminación.
\[\left\{\begin{aligned}2x + 4y + 6z &= -12\\2x - 3y - 4z &= 15\\3x + 4y + 5z &= -8\end{aligned}\right.\]
\begin{multicols*}{2}
\noindent\begin{minipage}{\linewidth}\textbf{1. Primero: ¿qué significa resolver este sistema?}\\[2pt]Las letras x, y, z representan números que todavía no conocemos. Una solución es una elección de esos números que hace verdaderas TODAS las ecuaciones al mismo tiempo. La llave de la izquierda significa que las condiciones se deben cumplir juntas.
\par\smallskip\textbf{Qué buscamos.} Encontrar un mismo grupo de valores para las 3 letras.
\par\smallskip El número 3 delante de x indica tres veces ese mismo valor; es decir, sumar tres copias:
\[3x=x+x+x\]
\par\smallskip Ejemplo numérico de la notación: si x valiera 2, tres veces ese valor sería 6. Este 2 es solo un ejemplo; todavía no es la respuesta del ejercicio. El punto centrado indica multiplicación:
\[3\cdot 2=2+2+2=6\]
\par\smallskip\textbf{Comprueba que lo entiendes.} ¿Basta con que los valores funcionen en una sola ecuación?\par\textit{Respuesta:} No. Deben funcionar en todas las ecuaciones del sistema, usando los mismos valores.
\end{minipage}\par\medskip
\noindent\begin{minipage}{\linewidth}\textbf{2. Nuestro plan: ir quitando letras}\\[2pt]Una ecuación con una sola letra es más fácil de resolver. Vamos a combinar ecuaciones para que una letra desaparezca. Cuando sepamos un valor, lo pondremos en las ecuaciones que faltan.
\par\smallskip\textbf{Qué buscamos.} Primero quitar x de las otras ecuaciones. Después seguir con las letras restantes.
\par\smallskip\textbf{Por qué.} Si una cantidad se resta de sí misma, queda cero. Usaremos esta idea con términos que tengan la misma letra.
\par\smallskip Ejemplo de la idea que utilizaremos:
\[3x-3x=0\]
\par\smallskip\textbf{Qué aprendimos.} Eliminar una letra de una ecuación NO significa que su valor sea cero; significa que sus términos se cancelaron en esa cuenta.
\end{minipage}\par\medskip
\noindent\begin{minipage}{\linewidth}\textbf{3. ¿Cómo haremos desaparecer x?}\\[2pt]Trabajaremos con las ecuaciones 1 y 2. Queremos que los términos con x tengan el mismo tamaño. Así podremos restarlos.
\par\smallskip\textbf{Qué buscamos.} Conseguir términos con 2 veces x en las dos ecuaciones.
\par\smallskip\textbf{Por qué.} Tener 2 veces la misma letra en ambos lados de la resta nos permitirá cancelarla.
\par\smallskip {\small\renewcommand{\arraystretch}{1.3}\begin{tabular}{>{\raggedright\arraybackslash}p{2.58cm}>{\raggedright\arraybackslash}p{2.47cm}>{\raggedright\arraybackslash}p{6.17cm}}\hline\rowcolor{black!7} \textbf{Dónde miramos} & \textbf{Qué término hay} & \textbf{Qué necesitamos hacer}\\\hline Ecuación 1 & $2x$ & La dejamos igual.\\ Ecuación 2 & $2x$ & La dejamos igual.\\ \hline\end{tabular}}\par\smallskip
\par\smallskip\textbf{Qué aprendimos.} Estos multiplicadores se eligen para igualar los términos con x. Los demás términos también se multiplican para conservar la igualdad.
\end{minipage}\par\medskip
\noindent\begin{minipage}{\linewidth}\textbf{4. Restamos las dos ecuaciones preparadas}\\[2pt]Restamos la ecuación 1 preparada a la ecuación 2 preparada. Juntamos cada término solo con otro que tenga la misma letra. Los números sin letra se calculan por separado.
\par\smallskip {\small\renewcommand{\arraystretch}{1.3}\begin{tabular}{>{\raggedright\arraybackslash}p{2.80cm}>{\raggedright\arraybackslash}p{5.38cm}>{\raggedright\arraybackslash}p{3.03cm}}\hline\rowcolor{black!7} \textbf{Qué agrupamos} & \textbf{Cuenta} & \textbf{Qué significa}\\\hline Los términos con x & $2x - \left(2x\right)=0$ & Se cancela esta letra.\\ Los términos con y & $-3y - \left(4y\right)=-7y$ & Quedan -7y.\\ Los términos con z & $-4z - \left(6z\right)=-10z$ & Quedan -10z.\\ Los números a la derecha del igual & $15 - \left(-12\right)=15+12=27$ & Este será el número de la nueva igualdad.\\ \hline\end{tabular}}\par\smallskip
\par\smallskip Primera igualdad que usaremos:
\[2x - 3y - 4z = 15\]
\par\smallskip Le restamos esta igualdad:
\[2x + 4y + 6z = -12\]
\par\smallskip\textbf{Qué aprendimos.} Los términos con x dan cero. Recuerda: restar un número negativo equivale a sumar su opuesto.
\end{minipage}\par\medskip
\noindent\begin{minipage}{\linewidth}\textbf{5. Mira lo que hemos conseguido}\\[2pt]La nueva ecuación es el resultado de la cuenta anterior. Todavía contiene y, z. Seguiremos reduciendo las condiciones.
\par\smallskip Escribimos la nueva condición:
\[-7y - 10z = 27\]
\[\left\{\begin{aligned}2x + 4y + 6z &= -12\\-7y - 10z &= 27\\3x + 4y + 5z &= -8\end{aligned}\right.\]
\par\smallskip\textbf{Qué aprendimos.} Conservamos la ecuación 1 y reemplazamos la 2 por esta nueva condición. De esa manera conservamos toda la información necesaria para resolver el sistema.
\end{minipage}\par\medskip
\noindent\begin{minipage}{\linewidth}\textbf{6. ¿Cómo haremos desaparecer x?}\\[2pt]Trabajaremos con las ecuaciones 1 y 3. Queremos que los términos con x tengan el mismo tamaño. Así podremos restarlos.
\par\smallskip\textbf{Qué buscamos.} Conseguir términos con 6 veces x en las dos ecuaciones.
\par\smallskip\textbf{Por qué.} Tener 6 veces la misma letra en ambos lados de la resta nos permitirá cancelarla.
\par\smallskip {\small\renewcommand{\arraystretch}{1.3}\begin{tabular}{>{\raggedright\arraybackslash}p{2.58cm}>{\raggedright\arraybackslash}p{2.47cm}>{\raggedright\arraybackslash}p{6.17cm}}\hline\rowcolor{black!7} \textbf{Dónde miramos} & \textbf{Qué término hay} & \textbf{Qué necesitamos hacer}\\\hline Ecuación 1 & $2x$ & Multiplicar TODA esta ecuación por 3\\ Ecuación 3 & $3x$ & Multiplicar TODA esta ecuación por 2\\ \hline\end{tabular}}\par\smallskip
\par\smallskip\textbf{Qué aprendimos.} Estos multiplicadores se eligen para igualar los términos con x. Los demás términos también se multiplican para conservar la igualdad.
\end{minipage}\par\medskip
\noindent\begin{minipage}{\linewidth}\textbf{7. Multiplicamos la ecuación 1 por 3}\\[2pt]Multiplicar una ecuación por 3 significa multiplicar CADA término de la izquierda y también el número de la derecha. La igualdad se conserva porque hacemos lo mismo en ambos lados.
\par\smallskip Esta es la ecuación antes de multiplicar:
\[2x + 4y + 6z = -12\]
\par\smallskip {\small\renewcommand{\arraystretch}{1.3}\begin{tabular}{>{\raggedright\arraybackslash}p{2.80cm}>{\raggedright\arraybackslash}p{1.91cm}>{\raggedright\arraybackslash}p{6.50cm}}\hline\rowcolor{black!7} \textbf{Parte de la ecuación} & \textbf{Antes} & \textbf{La cuenta y su resultado}\\\hline Término con x & $2x$ & $3\cdot\left(2x\right)=6x$\\ Término con y & $4y$ & $3\cdot\left(4y\right)=12y$\\ Término con z & $6z$ & $3\cdot\left(6z\right)=18z$\\ Número a la derecha del igual & $-12$ & $3\cdot\left(-12\right)=-36$\\ \hline\end{tabular}}\par\smallskip
\par\smallskip Después de hacer todas las multiplicaciones queda:
\[6x + 12y + 18z = -36\]
\par\smallskip\textbf{Qué aprendimos.} La misma condición ahora está escrita con los términos multiplicados. Ya tenemos el tamaño que buscábamos para x.
\par\smallskip\textbf{Comprueba que lo entiendes.} ¿Podríamos multiplicar solamente el término con x?\par\textit{Respuesta:} No. Eso cambiaría la condición. Debemos multiplicar todos los términos, también el número que está después del igual.
\end{minipage}\par\medskip
\noindent\begin{minipage}{\linewidth}\textbf{8. Multiplicamos la ecuación 3 por 2}\\[2pt]Multiplicar una ecuación por 2 significa multiplicar CADA término de la izquierda y también el número de la derecha. La igualdad se conserva porque hacemos lo mismo en ambos lados.
\par\smallskip Esta es la ecuación antes de multiplicar:
\[3x + 4y + 5z = -8\]
\par\smallskip {\small\renewcommand{\arraystretch}{1.3}\begin{tabular}{>{\raggedright\arraybackslash}p{2.80cm}>{\raggedright\arraybackslash}p{1.91cm}>{\raggedright\arraybackslash}p{6.50cm}}\hline\rowcolor{black!7} \textbf{Parte de la ecuación} & \textbf{Antes} & \textbf{La cuenta y su resultado}\\\hline Término con x & $3x$ & $2\cdot\left(3x\right)=6x$\\ Término con y & $4y$ & $2\cdot\left(4y\right)=8y$\\ Término con z & $5z$ & $2\cdot\left(5z\right)=10z$\\ Número a la derecha del igual & $-8$ & $2\cdot\left(-8\right)=-16$\\ \hline\end{tabular}}\par\smallskip
\par\smallskip Después de hacer todas las multiplicaciones queda:
\[6x + 8y + 10z = -16\]
\par\smallskip\textbf{Qué aprendimos.} La misma condición ahora está escrita con los términos multiplicados. Ya tenemos el tamaño que buscábamos para x.
\par\smallskip\textbf{Comprueba que lo entiendes.} ¿Podríamos multiplicar solamente el término con x?\par\textit{Respuesta:} No. Eso cambiaría la condición. Debemos multiplicar todos los términos, también el número que está después del igual.
\end{minipage}\par\medskip
\noindent\begin{minipage}{\linewidth}\textbf{9. Restamos las dos ecuaciones preparadas}\\[2pt]Restamos la ecuación 1 preparada a la ecuación 3 preparada. Juntamos cada término solo con otro que tenga la misma letra. Los números sin letra se calculan por separado.
\par\smallskip {\small\renewcommand{\arraystretch}{1.3}\begin{tabular}{>{\raggedright\arraybackslash}p{2.80cm}>{\raggedright\arraybackslash}p{5.38cm}>{\raggedright\arraybackslash}p{3.03cm}}\hline\rowcolor{black!7} \textbf{Qué agrupamos} & \textbf{Cuenta} & \textbf{Qué significa}\\\hline Los términos con x & $6x - \left(6x\right)=0$ & Se cancela esta letra.\\ Los términos con y & $8y - \left(12y\right)=-4y$ & Quedan -4y.\\ Los términos con z & $10z - \left(18z\right)=-8z$ & Quedan -8z.\\ Los números a la derecha del igual & $-16 - \left(-36\right)=-16+36=20$ & Este será el número de la nueva igualdad.\\ \hline\end{tabular}}\par\smallskip
\par\smallskip Primera igualdad que usaremos:
\[6x + 8y + 10z = -16\]
\par\smallskip Le restamos esta igualdad:
\[6x + 12y + 18z = -36\]
\par\smallskip\textbf{Qué aprendimos.} Los términos con x dan cero. Recuerda: restar un número negativo equivale a sumar su opuesto.
\end{minipage}\par\medskip
\noindent\begin{minipage}{\linewidth}\textbf{10. Mira lo que hemos conseguido}\\[2pt]La nueva ecuación es el resultado de la cuenta anterior. Todavía contiene y, z. Seguiremos reduciendo las condiciones.
\par\smallskip Escribimos la nueva condición:
\[-4y - 8z = 20\]
\[\left\{\begin{aligned}2x + 4y + 6z &= -12\\-7y - 10z &= 27\\-4y - 8z &= 20\end{aligned}\right.\]
\par\smallskip\textbf{Qué aprendimos.} Conservamos la ecuación 1 y reemplazamos la 3 por esta nueva condición. De esa manera conservamos toda la información necesaria para resolver el sistema.
\end{minipage}\par\medskip
\noindent\begin{minipage}{\linewidth}\textbf{11. Escribimos la misma condición con números más pequeños}\\[2pt]Todos los números de esta ecuación se pueden dividir entre 4 sin dejar restos. Dividimos ambos lados entre ese mismo número. Así la condición no cambia y las siguientes cuentas serán más sencillas.
\par\smallskip Antes de simplificar:
\[-4y - 8z = 20\]
\par\smallskip Después de dividir todos los términos entre 4:
\[-y - 2z = 5\]
\end{minipage}\par\medskip
\noindent\begin{minipage}{\linewidth}\textbf{12. Ponemos una ecuación más cómoda arriba}\\[2pt]Cambiamos de lugar las ecuaciones 2 y 3. Las condiciones siguen siendo las mismas; solamente cambia el orden en que las leemos.
\par\smallskip\textbf{Por qué.} La ecuación que subimos tiene un número más pequeño delante de y. Eso facilita las multiplicaciones que haremos a continuación.
\[\left\{\begin{aligned}2x + 4y + 6z &= -12\\-y - 2z &= 5\\-7y - 10z &= 27\end{aligned}\right.\]
\end{minipage}\par\medskip
\noindent\begin{minipage}{\linewidth}\textbf{13. ¿Cómo haremos desaparecer y?}\\[2pt]Trabajaremos con las ecuaciones 2 y 3. Queremos que los términos con y tengan el mismo tamaño. Así podremos restarlos.
\par\smallskip\textbf{Qué buscamos.} Conseguir términos con 7 veces y en las dos ecuaciones.
\par\smallskip\textbf{Por qué.} Tener 7 veces la misma letra en ambos lados de la resta nos permitirá cancelarla.
\par\smallskip {\small\renewcommand{\arraystretch}{1.3}\begin{tabular}{>{\raggedright\arraybackslash}p{2.58cm}>{\raggedright\arraybackslash}p{2.47cm}>{\raggedright\arraybackslash}p{6.17cm}}\hline\rowcolor{black!7} \textbf{Dónde miramos} & \textbf{Qué término hay} & \textbf{Qué necesitamos hacer}\\\hline Ecuación 2 & $-y$ & Multiplicar TODA esta ecuación por 7\\ Ecuación 3 & $-7y$ & La dejamos igual.\\ \hline\end{tabular}}\par\smallskip
\par\smallskip\textbf{Qué aprendimos.} Estos multiplicadores se eligen para igualar los términos con y. Los demás términos también se multiplican para conservar la igualdad.
\end{minipage}\par\medskip
\noindent\begin{minipage}{\linewidth}\textbf{14. Multiplicamos la ecuación 2 por 7}\\[2pt]Multiplicar una ecuación por 7 significa multiplicar CADA término de la izquierda y también el número de la derecha. La igualdad se conserva porque hacemos lo mismo en ambos lados.
\par\smallskip Esta es la ecuación antes de multiplicar:
\[-y - 2z = 5\]
\par\smallskip {\small\renewcommand{\arraystretch}{1.3}\begin{tabular}{>{\raggedright\arraybackslash}p{2.80cm}>{\raggedright\arraybackslash}p{1.91cm}>{\raggedright\arraybackslash}p{6.50cm}}\hline\rowcolor{black!7} \textbf{Parte de la ecuación} & \textbf{Antes} & \textbf{La cuenta y su resultado}\\\hline Término con y & $-y$ & $7\cdot\left(-y\right)=-7y$\\ Término con z & $-2z$ & $7\cdot\left(-2z\right)=-14z$\\ Número a la derecha del igual & $5$ & $7\cdot\left(5\right)=35$\\ \hline\end{tabular}}\par\smallskip
\par\smallskip Después de hacer todas las multiplicaciones queda:
\[-7y - 14z = 35\]
\par\smallskip\textbf{Qué aprendimos.} La misma condición ahora está escrita con los términos multiplicados. Ya tenemos el tamaño que buscábamos para y.
\par\smallskip\textbf{Comprueba que lo entiendes.} ¿Podríamos multiplicar solamente el término con y?\par\textit{Respuesta:} No. Eso cambiaría la condición. Debemos multiplicar todos los términos, también el número que está después del igual.
\end{minipage}\par\medskip
\noindent\begin{minipage}{\linewidth}\textbf{15. Restamos las dos ecuaciones preparadas}\\[2pt]Restamos la ecuación 2 preparada a la ecuación 3 preparada. Juntamos cada término solo con otro que tenga la misma letra. Los números sin letra se calculan por separado.
\par\smallskip {\small\renewcommand{\arraystretch}{1.3}\begin{tabular}{>{\raggedright\arraybackslash}p{2.80cm}>{\raggedright\arraybackslash}p{5.38cm}>{\raggedright\arraybackslash}p{3.03cm}}\hline\rowcolor{black!7} \textbf{Qué agrupamos} & \textbf{Cuenta} & \textbf{Qué significa}\\\hline Los términos con y & $-7y - \left(-7y\right)=0$ & Se cancela esta letra.\\ Los términos con z & $-10z - \left(-14z\right)=4z$ & Quedan 4z.\\ Los números a la derecha del igual & $27 - \left(35\right)=-8$ & Este será el número de la nueva igualdad.\\ \hline\end{tabular}}\par\smallskip
\par\smallskip Primera igualdad que usaremos:
\[-7y - 10z = 27\]
\par\smallskip Le restamos esta igualdad:
\[-7y - 14z = 35\]
\par\smallskip\textbf{Qué aprendimos.} Los términos con y dan cero. Los términos con letras distintas no se mezclan.
\end{minipage}\par\medskip
\noindent\begin{minipage}{\linewidth}\textbf{16. Mira lo que hemos conseguido}\\[2pt]La nueva ecuación es el resultado de la cuenta anterior. Solo contiene z: después podremos despejarla.
\par\smallskip Escribimos la nueva condición:
\[4z = -8\]
\[\left\{\begin{aligned}2x + 4y + 6z &= -12\\-y - 2z &= 5\\4z &= -8\end{aligned}\right.\]
\par\smallskip\textbf{Qué aprendimos.} Conservamos la ecuación 2 y reemplazamos la 3 por esta nueva condición. De esa manera conservamos toda la información necesaria para resolver el sistema.
\end{minipage}\par\medskip
\noindent\begin{minipage}{\linewidth}\textbf{17. Ya tenemos una ecuación con una sola letra}\\[2pt]En esta condición solamente falta conocer z. El número 4 está multiplicando a esa letra. Ahora podemos deshacer esa multiplicación.
\par\smallskip Esta es la condición que vamos a resolver:
\[4z = -8\]
\end{minipage}\par\medskip
\noindent\begin{minipage}{\linewidth}\textbf{18. Calculamos cuánto vale z}\\[2pt]Para deshacer una multiplicación por 4, dividimos AMBOS lados entre 4. En el lado izquierdo queda una sola z. En el derecho hacemos la división.
\par\smallskip Dividimos los dos lados entre el mismo número:
\[\frac{4z}{4}=\frac{-8}{4}\]
\par\smallskip Ya podemos leer el valor de la letra:
\[z=\frac{-8}{4}=-2\]
\par\smallskip\textbf{Qué aprendimos.} Este valor se usará en las ecuaciones donde todavía falten letras.
\end{minipage}\par\medskip
\noindent\begin{minipage}{\linewidth}\textbf{19. Usamos lo que ya sabemos para hallar y}\\[2pt]Ya sabemos z = -2. Sustituir significa escribir ese valor cada vez que aparezca su letra. Usamos paréntesis para que se vea qué número estamos multiplicando.
\par\smallskip\textbf{Qué buscamos.} Dejar solamente y como letra desconocida en esta ecuación.
\par\smallskip Tomamos esta condición que obtuvimos antes:
\[-y - 2z = 5\]
\par\smallskip Reemplazamos las letras conocidas por sus valores:
\[-y - 2\left(-2\right)=5\]
\end{minipage}\par\medskip
\noindent\begin{minipage}{\linewidth}\textbf{20. Calculamos las multiplicaciones que ya se pueden hacer}\\[2pt]La letra y sigue sin conocerse. Las otras partes ya tienen un valor o una expresión conocida y podemos efectuar sus productos.
\par\smallskip {\small\renewcommand{\arraystretch}{1.3}\begin{tabular}{>{\raggedright\arraybackslash}p{3.36cm}>{\raggedright\arraybackslash}p{4.48cm}>{\raggedright\arraybackslash}p{3.36cm}}\hline\rowcolor{black!7} \textbf{Parte conocida} & \textbf{La multiplicación} & \textbf{Resultado}\\\hline El término -2z & $-2\cdot\left(-2\right)$ & $4$\\ \hline\end{tabular}}\par\smallskip
\par\smallskip Al reunir esas partes conocidas, la ecuación queda:
\[-y + 4=5\]
\end{minipage}\par\medskip
\noindent\begin{minipage}{\linewidth}\textbf{21. Dejamos el término con y solo a un lado}\\[2pt]Para quitar la parte conocida, restamos 4 en AMBOS lados. Piensa en una balanza: si hacemos lo mismo a ambos lados, la igualdad se mantiene.
\par\smallskip Hacemos la misma operación a la izquierda y a la derecha:
\[-y + 4 - 4=5 - 4\]
\par\smallskip Las partes que se compensan dan cero; queda:
\[-y=1\]
\par\smallskip\textbf{Qué aprendimos.} «Pasar al otro lado» es una forma abreviada de decir que hicimos la misma suma o resta en los dos lados.
\end{minipage}\par\medskip
\noindent\begin{minipage}{\linewidth}\textbf{22. Calculamos cuánto vale y}\\[2pt]Para deshacer una multiplicación por -1, dividimos AMBOS lados entre -1. En el lado izquierdo queda una sola y. En el derecho hacemos la división.
\par\smallskip Dividimos los dos lados entre el mismo número:
\[\frac{-y}{-1}=\frac{1}{-1}\]
\par\smallskip Ya podemos leer el valor de la letra:
\[y=\frac{1}{-1}=-1\]
\par\smallskip\textbf{Qué aprendimos.} Este valor se usará en las ecuaciones donde todavía falten letras.
\end{minipage}\par\medskip
\noindent\begin{minipage}{\linewidth}\textbf{23. Usamos lo que ya sabemos para hallar x}\\[2pt]Ya sabemos y = -1; z = -2. Sustituir significa escribir ese valor cada vez que aparezca su letra. Usamos paréntesis para que se vea qué número estamos multiplicando.
\par\smallskip\textbf{Qué buscamos.} Dejar solamente x como letra desconocida en esta ecuación.
\par\smallskip Tomamos esta condición que obtuvimos antes:
\[2x + 4y + 6z = -12\]
\par\smallskip Reemplazamos las letras conocidas por sus valores:
\[2x + 4\left(-1\right) + 6\left(-2\right)=-12\]
\end{minipage}\par\medskip
\noindent\begin{minipage}{\linewidth}\textbf{24. Calculamos las multiplicaciones que ya se pueden hacer}\\[2pt]La letra x sigue sin conocerse. Las otras partes ya tienen un valor o una expresión conocida y podemos efectuar sus productos.
\par\smallskip {\small\renewcommand{\arraystretch}{1.3}\begin{tabular}{>{\raggedright\arraybackslash}p{3.36cm}>{\raggedright\arraybackslash}p{4.48cm}>{\raggedright\arraybackslash}p{3.36cm}}\hline\rowcolor{black!7} \textbf{Parte conocida} & \textbf{La multiplicación} & \textbf{Resultado}\\\hline El término 4y & $4\cdot\left(-1\right)$ & $-4$\\ El término 6z & $6\cdot\left(-2\right)$ & $-12$\\ \hline\end{tabular}}\par\smallskip
\par\smallskip Al reunir esas partes conocidas, la ecuación queda:
\[2x - 16=-12\]
\end{minipage}\par\medskip
\noindent\begin{minipage}{\linewidth}\textbf{25. Dejamos el término con x solo a un lado}\\[2pt]Para quitar la parte conocida, sumamos 16 en AMBOS lados. Piensa en una balanza: si hacemos lo mismo a ambos lados, la igualdad se mantiene.
\par\smallskip Hacemos la misma operación a la izquierda y a la derecha:
\[2x - 16 + 16=-12 + 16\]
\par\smallskip Las partes que se compensan dan cero; queda:
\[2x=4\]
\par\smallskip\textbf{Qué aprendimos.} «Pasar al otro lado» es una forma abreviada de decir que hicimos la misma suma o resta en los dos lados.
\end{minipage}\par\medskip
\noindent\begin{minipage}{\linewidth}\textbf{26. Calculamos cuánto vale x}\\[2pt]Para deshacer una multiplicación por 2, dividimos AMBOS lados entre 2. En el lado izquierdo queda una sola x. En el derecho hacemos la división.
\par\smallskip Dividimos los dos lados entre el mismo número:
\[\frac{2x}{2}=\frac{4}{2}\]
\par\smallskip Ya podemos leer el valor de la letra:
\[x=\frac{4}{2}=2\]
\par\smallskip\textbf{Qué aprendimos.} Este valor se usará en las ecuaciones donde todavía falten letras.
\end{minipage}\par\medskip
\noindent\begin{minipage}{\linewidth}\textbf{27. Reunimos la respuesta; todavía falta comprobarla}\\[2pt]Escribimos todos los valores en el orden x, y, z. Estos valores deben funcionar juntos, no cada uno en una ecuación diferente.
\[\left\{\begin{aligned}x&=2\\y&=-1\\z&=-2\end{aligned}\right.\]
\par\smallskip\textbf{Qué aprendimos.} El siguiente paso es volver a las ecuaciones del enunciado y comprobar cada condición.
\end{minipage}\par\medskip
\noindent\begin{minipage}{\linewidth}\textbf{28. ¿La respuesta cumple la ecuación 1?}\\[2pt]Comprobar significa volver a la condición original y calcularla con nuestra respuesta. Si los dos lados dan el mismo número, esta condición sí se cumple. Deben cumplirse todas.
\par\smallskip La condición del enunciado era:
\[2x + 4y + 6z = -12\]
\par\smallskip Ponemos nuestros valores donde estaban las letras:
\[2\left(2\right) + 4\left(-1\right) + 6\left(-2\right)=-12\]
\par\smallskip Hacemos las multiplicaciones y luego la suma:
\[4 - 4 - 12=-12\]
\par\smallskip Comparamos los dos lados:
\[-12=-12\]
\par\smallskip\textbf{Qué aprendimos.} La igualdad final es verdadera. Nuestra respuesta sí cumple la ecuación 1.
\end{minipage}\par\medskip
\noindent\begin{minipage}{\linewidth}\textbf{29. ¿La respuesta cumple la ecuación 2?}\\[2pt]Comprobar significa volver a la condición original y calcularla con nuestra respuesta. Si los dos lados dan el mismo número, esta condición sí se cumple. Deben cumplirse todas.
\par\smallskip La condición del enunciado era:
\[2x - 3y - 4z = 15\]
\par\smallskip Ponemos nuestros valores donde estaban las letras:
\[2\left(2\right) - 3\left(-1\right) - 4\left(-2\right)=15\]
\par\smallskip Hacemos las multiplicaciones y luego la suma:
\[4 + 3 + 8=15\]
\par\smallskip Comparamos los dos lados:
\[15=15\]
\par\smallskip\textbf{Qué aprendimos.} La igualdad final es verdadera. Nuestra respuesta sí cumple la ecuación 2.
\end{minipage}\par\medskip
\noindent\begin{minipage}{\linewidth}\textbf{30. ¿La respuesta cumple la ecuación 3?}\\[2pt]Comprobar significa volver a la condición original y calcularla con nuestra respuesta. Si los dos lados dan el mismo número, esta condición sí se cumple. Deben cumplirse todas.
\par\smallskip La condición del enunciado era:
\[3x + 4y + 5z = -8\]
\par\smallskip Ponemos nuestros valores donde estaban las letras:
\[3\left(2\right) + 4\left(-1\right) + 5\left(-2\right)=-8\]
\par\smallskip Hacemos las multiplicaciones y luego la suma:
\[6 - 4 - 10=-8\]
\par\smallskip Comparamos los dos lados:
\[-8=-8\]
\par\smallskip\textbf{Qué aprendimos.} La igualdad final es verdadera. Nuestra respuesta sí cumple la ecuación 3.
\end{minipage}\par\medskip
\respuesta{x = 2; y = -1; z = -2.}
\noindent\begin{minipage}{\linewidth}\centering\includegraphics[width=.96\linewidth,height=6.8cm,keepaspectratio]{../recursos/grafico_1_1-5.png}\\{\small Lectura visual complementaria; la demostración está en las operaciones.}\end{minipage}
\end{multicols*}
\newpage\fancyfoot[L]{\small Sección 1.1 | Ejercicio 6 | Paso a paso}\hypertarget{1.1-6}{}
\pdfbookmark[1]{1.1 / 6}{bm1.1-6}
\titulo{6 | Sistema lineal por eliminación}
\textbf{Enunciado}\par Resuelva el sistema lineal dado por medio del método de eliminación.
\[\left\{\begin{aligned}x + y - 2z &= 5\\2x + 3y + 4z &= 2\end{aligned}\right.\]
\begin{multicols*}{2}
\noindent\begin{minipage}{\linewidth}\textbf{1. Primero: ¿qué significa resolver este sistema?}\\[2pt]Las letras x, y, z representan números que todavía no conocemos. Una solución es una elección de esos números que hace verdaderas TODAS las ecuaciones al mismo tiempo. La llave de la izquierda significa que las condiciones se deben cumplir juntas.
\par\smallskip\textbf{Qué buscamos.} Encontrar un mismo grupo de valores para las 3 letras.
\par\smallskip El número 3 delante de x indica tres veces ese mismo valor; es decir, sumar tres copias:
\[3x=x+x+x\]
\par\smallskip Ejemplo numérico de la notación: si x valiera 2, tres veces ese valor sería 6. Este 2 es solo un ejemplo; todavía no es la respuesta del ejercicio. El punto centrado indica multiplicación:
\[3\cdot 2=2+2+2=6\]
\par\smallskip\textbf{Comprueba que lo entiendes.} ¿Basta con que los valores funcionen en una sola ecuación?\par\textit{Respuesta:} No. Deben funcionar en todas las ecuaciones del sistema, usando los mismos valores.
\end{minipage}\par\medskip
\noindent\begin{minipage}{\linewidth}\textbf{2. Nuestro plan: ir quitando letras}\\[2pt]Una ecuación con una sola letra es más fácil de resolver. Vamos a combinar ecuaciones para que una letra desaparezca. Cuando sepamos un valor, lo pondremos en las ecuaciones que faltan.
\par\smallskip\textbf{Qué buscamos.} Primero quitar x de las otras ecuaciones. Después seguir con las letras restantes.
\par\smallskip\textbf{Por qué.} Si una cantidad se resta de sí misma, queda cero. Usaremos esta idea con términos que tengan la misma letra.
\par\smallskip Ejemplo de la idea que utilizaremos:
\[3x-3x=0\]
\par\smallskip\textbf{Qué aprendimos.} Eliminar una letra de una ecuación NO significa que su valor sea cero; significa que sus términos se cancelaron en esa cuenta.
\end{minipage}\par\medskip
\noindent\begin{minipage}{\linewidth}\textbf{3. ¿Cómo haremos desaparecer x?}\\[2pt]Trabajaremos con las ecuaciones 1 y 2. Queremos que los términos con x tengan el mismo tamaño. Así podremos restarlos.
\par\smallskip\textbf{Qué buscamos.} Conseguir términos con 2 veces x en las dos ecuaciones.
\par\smallskip\textbf{Por qué.} Tener 2 veces la misma letra en ambos lados de la resta nos permitirá cancelarla.
\par\smallskip {\small\renewcommand{\arraystretch}{1.3}\begin{tabular}{>{\raggedright\arraybackslash}p{2.58cm}>{\raggedright\arraybackslash}p{2.47cm}>{\raggedright\arraybackslash}p{6.17cm}}\hline\rowcolor{black!7} \textbf{Dónde miramos} & \textbf{Qué término hay} & \textbf{Qué necesitamos hacer}\\\hline Ecuación 1 & $x$ & Multiplicar TODA esta ecuación por 2\\ Ecuación 2 & $2x$ & La dejamos igual.\\ \hline\end{tabular}}\par\smallskip
\par\smallskip\textbf{Qué aprendimos.} Estos multiplicadores se eligen para igualar los términos con x. Los demás términos también se multiplican para conservar la igualdad.
\end{minipage}\par\medskip
\noindent\begin{minipage}{\linewidth}\textbf{4. Multiplicamos la ecuación 1 por 2}\\[2pt]Multiplicar una ecuación por 2 significa multiplicar CADA término de la izquierda y también el número de la derecha. La igualdad se conserva porque hacemos lo mismo en ambos lados.
\par\smallskip Esta es la ecuación antes de multiplicar:
\[x + y - 2z = 5\]
\par\smallskip {\small\renewcommand{\arraystretch}{1.3}\begin{tabular}{>{\raggedright\arraybackslash}p{2.80cm}>{\raggedright\arraybackslash}p{1.91cm}>{\raggedright\arraybackslash}p{6.50cm}}\hline\rowcolor{black!7} \textbf{Parte de la ecuación} & \textbf{Antes} & \textbf{La cuenta y su resultado}\\\hline Término con x & $x$ & $2\cdot\left(x\right)=2x$\\ Término con y & $y$ & $2\cdot\left(y\right)=2y$\\ Término con z & $-2z$ & $2\cdot\left(-2z\right)=-4z$\\ Número a la derecha del igual & $5$ & $2\cdot\left(5\right)=10$\\ \hline\end{tabular}}\par\smallskip
\par\smallskip Después de hacer todas las multiplicaciones queda:
\[2x + 2y - 4z = 10\]
\par\smallskip\textbf{Qué aprendimos.} La misma condición ahora está escrita con los términos multiplicados. Ya tenemos el tamaño que buscábamos para x.
\par\smallskip\textbf{Comprueba que lo entiendes.} ¿Podríamos multiplicar solamente el término con x?\par\textit{Respuesta:} No. Eso cambiaría la condición. Debemos multiplicar todos los términos, también el número que está después del igual.
\end{minipage}\par\medskip
\noindent\begin{minipage}{\linewidth}\textbf{5. Restamos las dos ecuaciones preparadas}\\[2pt]Restamos la ecuación 1 preparada a la ecuación 2 preparada. Juntamos cada término solo con otro que tenga la misma letra. Los números sin letra se calculan por separado.
\par\smallskip {\small\renewcommand{\arraystretch}{1.3}\begin{tabular}{>{\raggedright\arraybackslash}p{2.80cm}>{\raggedright\arraybackslash}p{5.38cm}>{\raggedright\arraybackslash}p{3.03cm}}\hline\rowcolor{black!7} \textbf{Qué agrupamos} & \textbf{Cuenta} & \textbf{Qué significa}\\\hline Los términos con x & $2x - \left(2x\right)=0$ & Se cancela esta letra.\\ Los términos con y & $3y - \left(2y\right)=y$ & Quedan y.\\ Los términos con z & $4z - \left(-4z\right)=8z$ & Quedan 8z.\\ Los números a la derecha del igual & $2 - \left(10\right)=-8$ & Este será el número de la nueva igualdad.\\ \hline\end{tabular}}\par\smallskip
\par\smallskip Primera igualdad que usaremos:
\[2x + 3y + 4z = 2\]
\par\smallskip Le restamos esta igualdad:
\[2x + 2y - 4z = 10\]
\par\smallskip\textbf{Qué aprendimos.} Los términos con x dan cero. Los términos con letras distintas no se mezclan.
\end{minipage}\par\medskip
\noindent\begin{minipage}{\linewidth}\textbf{6. Mira lo que hemos conseguido}\\[2pt]La nueva ecuación es el resultado de la cuenta anterior. Todavía contiene y, z. Seguiremos reduciendo las condiciones.
\par\smallskip Escribimos la nueva condición:
\[y + 8z = -8\]
\[\left\{\begin{aligned}x + y - 2z &= 5\\y + 8z &= -8\end{aligned}\right.\]
\par\smallskip\textbf{Qué aprendimos.} Conservamos la ecuación 1 y reemplazamos la 2 por esta nueva condición. De esa manera conservamos toda la información necesaria para resolver el sistema.
\end{minipage}\par\medskip
\noindent\begin{minipage}{\linewidth}\textbf{7. Una letra puede elegirse libremente}\\[2pt]Las condiciones que quedan no fijan un único valor para z. Podemos elegir ese valor y calcular los demás. Para representar todas las elecciones usamos otra letra, llamada parámetro.
\par\smallskip La letra z será el número que elijamos. Lo llamaremos t:
\[z=t\]
\par\smallskip\textbf{Qué aprendimos.} Un parámetro no es una respuesta que falte: es una forma de describir muchas respuestas a la vez. Primero lo dejamos libre; en un problema real también revisaremos si hay límites como cantidades no negativas.
\end{minipage}\par\medskip
\noindent\begin{minipage}{\linewidth}\textbf{8. Usamos lo que ya sabemos para hallar y}\\[2pt]Ya sabemos z = t. Sustituir significa escribir ese valor cada vez que aparezca su letra. Usamos paréntesis para que se vea qué número estamos multiplicando.
\par\smallskip\textbf{Qué buscamos.} Dejar solamente y como letra desconocida en esta ecuación.
\par\smallskip Tomamos esta condición que obtuvimos antes:
\[y + 8z = -8\]
\par\smallskip Reemplazamos las letras conocidas por sus valores:
\[y + 8\left(t\right)=-8\]
\end{minipage}\par\medskip
\noindent\begin{minipage}{\linewidth}\textbf{9. Calculamos las multiplicaciones que ya se pueden hacer}\\[2pt]La letra y sigue sin conocerse. Las otras partes ya tienen un valor o una expresión conocida y podemos efectuar sus productos.
\par\smallskip {\small\renewcommand{\arraystretch}{1.3}\begin{tabular}{>{\raggedright\arraybackslash}p{3.36cm}>{\raggedright\arraybackslash}p{4.48cm}>{\raggedright\arraybackslash}p{3.36cm}}\hline\rowcolor{black!7} \textbf{Parte conocida} & \textbf{La multiplicación} & \textbf{Resultado}\\\hline El término 8z & $8\cdot\left(t\right)$ & $8 t$\\ \hline\end{tabular}}\par\smallskip
\par\smallskip Al reunir esas partes conocidas, la ecuación queda:
\[y + 8 t=-8\]
\end{minipage}\par\medskip
\noindent\begin{minipage}{\linewidth}\textbf{10. Dejamos el término con y solo a un lado}\\[2pt]Para quitar la parte conocida, quitamos la expresión 8*t en AMBOS lados. Piensa en una balanza: si hacemos lo mismo a ambos lados, la igualdad se mantiene.
\par\smallskip Hacemos la misma operación a la izquierda y a la derecha:
\[y=-8-\left(8 t\right)\]
\par\smallskip Las partes que se compensan dan cero; queda:
\[y=- 8 t - 8\]
\par\smallskip\textbf{Qué aprendimos.} «Pasar al otro lado» es una forma abreviada de decir que hicimos la misma suma o resta en los dos lados.
\end{minipage}\par\medskip
\noindent\begin{minipage}{\linewidth}\textbf{11. Leemos el valor de y}\\[2pt]La letra ya está sola: su número multiplicador es 1, que no se escribe. La igualdad nos dice directamente cuánto vale.
\par\smallskip Anotamos el valor encontrado:
\[y=- 8 t - 8\]
\end{minipage}\par\medskip
\noindent\begin{minipage}{\linewidth}\textbf{12. Usamos lo que ya sabemos para hallar x}\\[2pt]Ya sabemos y = -8*t - 8; z = t. Sustituir significa escribir ese valor cada vez que aparezca su letra. Usamos paréntesis para que se vea qué número estamos multiplicando.
\par\smallskip\textbf{Qué buscamos.} Dejar solamente x como letra desconocida en esta ecuación.
\par\smallskip Tomamos esta condición que obtuvimos antes:
\[x + y - 2z = 5\]
\par\smallskip Reemplazamos las letras conocidas por sus valores:
\[x + \left(- 8 t - 8\right) - 2\left(t\right)=5\]
\end{minipage}\par\medskip
\noindent\begin{minipage}{\linewidth}\textbf{13. Calculamos las multiplicaciones que ya se pueden hacer}\\[2pt]La letra x sigue sin conocerse. Las otras partes ya tienen un valor o una expresión conocida y podemos efectuar sus productos.
\par\smallskip {\small\renewcommand{\arraystretch}{1.3}\begin{tabular}{>{\raggedright\arraybackslash}p{3.36cm}>{\raggedright\arraybackslash}p{4.48cm}>{\raggedright\arraybackslash}p{3.36cm}}\hline\rowcolor{black!7} \textbf{Parte conocida} & \textbf{La multiplicación} & \textbf{Resultado}\\\hline El término y & $1\cdot\left(- 8 t - 8\right)$ & $- 8 t - 8$\\ El término -2z & $-2\cdot\left(t\right)$ & $- 2 t$\\ \hline\end{tabular}}\par\smallskip
\par\smallskip Al reunir esas partes conocidas, la ecuación queda:
\[x + \left(- 10 t - 8\right)=5\]
\end{minipage}\par\medskip
\noindent\begin{minipage}{\linewidth}\textbf{14. Dejamos el término con x solo a un lado}\\[2pt]Para quitar la parte conocida, quitamos la expresión -10*t - 8 en AMBOS lados. Piensa en una balanza: si hacemos lo mismo a ambos lados, la igualdad se mantiene.
\par\smallskip Hacemos la misma operación a la izquierda y a la derecha:
\[x=5-\left(- 10 t - 8\right)\]
\par\smallskip Las partes que se compensan dan cero; queda:
\[x=10 t + 13\]
\par\smallskip\textbf{Qué aprendimos.} «Pasar al otro lado» es una forma abreviada de decir que hicimos la misma suma o resta en los dos lados.
\end{minipage}\par\medskip
\noindent\begin{minipage}{\linewidth}\textbf{15. Leemos el valor de x}\\[2pt]La letra ya está sola: su número multiplicador es 1, que no se escribe. La igualdad nos dice directamente cuánto vale.
\par\smallskip Anotamos el valor encontrado:
\[x=10 t + 13\]
\end{minipage}\par\medskip
\noindent\begin{minipage}{\linewidth}\textbf{16. Reunimos la respuesta; todavía falta comprobarla}\\[2pt]Escribimos todos los valores en el orden x, y, z. Dependen del número que elijamos para el parámetro, por eso describen una familia de respuestas.
\[\left\{\begin{aligned}x&=10 t + 13\\y&=- 8 t - 8\\z&=t\end{aligned}\right.\]
\par\smallskip\textbf{Qué aprendimos.} El siguiente paso es volver a las ecuaciones del enunciado y comprobar cada condición.
\end{minipage}\par\medskip
\noindent\begin{minipage}{\linewidth}\textbf{17. Dos elecciones distintas producen dos soluciones}\\[2pt]Esta tabla muestra qué pasa al elegir 0 o 1 para el parámetro. Sustituimos la elección en las expresiones de la respuesta. Cada fila es un grupo completo de valores.
\par\smallskip {\small\renewcommand{\arraystretch}{1.3}\begin{tabular}{>{\raggedright\arraybackslash}p{2.68cm}>{\raggedright\arraybackslash}p{2.68cm}>{\raggedright\arraybackslash}p{2.68cm}>{\raggedright\arraybackslash}p{2.68cm}}\hline\rowcolor{black!7} \textbf{Elección de t} & \textbf{x} & \textbf{y} & \textbf{z}\\\hline 0 & 13 & -8 & 0\\ 1 & 23 & -16 & 1\\ \hline\end{tabular}}\par\smallskip
\par\smallskip\textbf{Qué aprendimos.} No hay una sola solución. Hay una para cada elección permitida del parámetro. El deslizador de la web permite comprobar más ejemplos.
\end{minipage}\par\medskip
\noindent\begin{minipage}{\linewidth}\textbf{18. ¿La respuesta cumple la ecuación 1?}\\[2pt]Comprobar significa volver a la condición original y calcularla con nuestra respuesta. Si los dos lados dan el mismo número, esta condición sí se cumple. Deben cumplirse todas.
\par\smallskip La condición del enunciado era:
\[x + y - 2z = 5\]
\par\smallskip Ponemos nuestros valores donde estaban las letras:
\[\left(10 t + 13\right) + \left(- 8 t - 8\right) - 2\left(t\right)=5\]
\par\smallskip Hacemos las multiplicaciones y luego la suma:
\[\left(10 t + 13\right) + \left(- 8 t - 8\right) + - 2 t=5\]
\par\smallskip Comparamos los dos lados:
\[5=5\]
\par\smallskip\textbf{Qué aprendimos.} La igualdad final es verdadera. Al simplificar, los términos del parámetro se compensan: la condición vale para cualquier elección permitida.
\end{minipage}\par\medskip
\noindent\begin{minipage}{\linewidth}\textbf{19. ¿La respuesta cumple la ecuación 2?}\\[2pt]Comprobar significa volver a la condición original y calcularla con nuestra respuesta. Si los dos lados dan el mismo número, esta condición sí se cumple. Deben cumplirse todas.
\par\smallskip La condición del enunciado era:
\[2x + 3y + 4z = 2\]
\par\smallskip Ponemos nuestros valores donde estaban las letras:
\[2\left(10 t + 13\right) + 3\left(- 8 t - 8\right) + 4\left(t\right)=2\]
\par\smallskip Hacemos las multiplicaciones y luego la suma:
\[\left(20 t + 26\right) + \left(- 24 t - 24\right) + 4 t=2\]
\par\smallskip Comparamos los dos lados:
\[2=2\]
\par\smallskip\textbf{Qué aprendimos.} La igualdad final es verdadera. Al simplificar, los términos del parámetro se compensan: la condición vale para cualquier elección permitida.
\end{minipage}\par\medskip
\respuesta{x = 10*t + 13; y = -8*t - 8; z = t. Parámetro real.}
\noindent\begin{minipage}{\linewidth}\centering\includegraphics[width=.96\linewidth,height=6.8cm,keepaspectratio]{../recursos/grafico_1_1-6.png}\\{\small Lectura visual complementaria; la demostración está en las operaciones.}\end{minipage}
\end{multicols*}
\newpage\fancyfoot[L]{\small Sección 1.1 | Ejercicio 7 | Paso a paso}\hypertarget{1.1-7}{}
\pdfbookmark[1]{1.1 / 7}{bm1.1-7}
\titulo{7 | Sistema lineal por eliminación}
\textbf{Enunciado}\par Resuelva el sistema lineal dado por medio del método de eliminación.
\[\left\{\begin{aligned}x + 4y - z &= 12\\3x + 8y - 2z &= 4\end{aligned}\right.\]
\begin{multicols*}{2}
\noindent\begin{minipage}{\linewidth}\textbf{1. Primero: ¿qué significa resolver este sistema?}\\[2pt]Las letras x, y, z representan números que todavía no conocemos. Una solución es una elección de esos números que hace verdaderas TODAS las ecuaciones al mismo tiempo. La llave de la izquierda significa que las condiciones se deben cumplir juntas.
\par\smallskip\textbf{Qué buscamos.} Encontrar un mismo grupo de valores para las 3 letras.
\par\smallskip El número 3 delante de x indica tres veces ese mismo valor; es decir, sumar tres copias:
\[3x=x+x+x\]
\par\smallskip Ejemplo numérico de la notación: si x valiera 2, tres veces ese valor sería 6. Este 2 es solo un ejemplo; todavía no es la respuesta del ejercicio. El punto centrado indica multiplicación:
\[3\cdot 2=2+2+2=6\]
\par\smallskip\textbf{Comprueba que lo entiendes.} ¿Basta con que los valores funcionen en una sola ecuación?\par\textit{Respuesta:} No. Deben funcionar en todas las ecuaciones del sistema, usando los mismos valores.
\end{minipage}\par\medskip
\noindent\begin{minipage}{\linewidth}\textbf{2. Nuestro plan: ir quitando letras}\\[2pt]Una ecuación con una sola letra es más fácil de resolver. Vamos a combinar ecuaciones para que una letra desaparezca. Cuando sepamos un valor, lo pondremos en las ecuaciones que faltan.
\par\smallskip\textbf{Qué buscamos.} Primero quitar x de las otras ecuaciones. Después seguir con las letras restantes.
\par\smallskip\textbf{Por qué.} Si una cantidad se resta de sí misma, queda cero. Usaremos esta idea con términos que tengan la misma letra.
\par\smallskip Ejemplo de la idea que utilizaremos:
\[3x-3x=0\]
\par\smallskip\textbf{Qué aprendimos.} Eliminar una letra de una ecuación NO significa que su valor sea cero; significa que sus términos se cancelaron en esa cuenta.
\end{minipage}\par\medskip
\noindent\begin{minipage}{\linewidth}\textbf{3. ¿Cómo haremos desaparecer x?}\\[2pt]Trabajaremos con las ecuaciones 1 y 2. Queremos que los términos con x tengan el mismo tamaño. Así podremos restarlos.
\par\smallskip\textbf{Qué buscamos.} Conseguir términos con 3 veces x en las dos ecuaciones.
\par\smallskip\textbf{Por qué.} Tener 3 veces la misma letra en ambos lados de la resta nos permitirá cancelarla.
\par\smallskip {\small\renewcommand{\arraystretch}{1.3}\begin{tabular}{>{\raggedright\arraybackslash}p{2.58cm}>{\raggedright\arraybackslash}p{2.47cm}>{\raggedright\arraybackslash}p{6.17cm}}\hline\rowcolor{black!7} \textbf{Dónde miramos} & \textbf{Qué término hay} & \textbf{Qué necesitamos hacer}\\\hline Ecuación 1 & $x$ & Multiplicar TODA esta ecuación por 3\\ Ecuación 2 & $3x$ & La dejamos igual.\\ \hline\end{tabular}}\par\smallskip
\par\smallskip\textbf{Qué aprendimos.} Estos multiplicadores se eligen para igualar los términos con x. Los demás términos también se multiplican para conservar la igualdad.
\end{minipage}\par\medskip
\noindent\begin{minipage}{\linewidth}\textbf{4. Multiplicamos la ecuación 1 por 3}\\[2pt]Multiplicar una ecuación por 3 significa multiplicar CADA término de la izquierda y también el número de la derecha. La igualdad se conserva porque hacemos lo mismo en ambos lados.
\par\smallskip Esta es la ecuación antes de multiplicar:
\[x + 4y - z = 12\]
\par\smallskip {\small\renewcommand{\arraystretch}{1.3}\begin{tabular}{>{\raggedright\arraybackslash}p{2.80cm}>{\raggedright\arraybackslash}p{1.91cm}>{\raggedright\arraybackslash}p{6.50cm}}\hline\rowcolor{black!7} \textbf{Parte de la ecuación} & \textbf{Antes} & \textbf{La cuenta y su resultado}\\\hline Término con x & $x$ & $3\cdot\left(x\right)=3x$\\ Término con y & $4y$ & $3\cdot\left(4y\right)=12y$\\ Término con z & $-z$ & $3\cdot\left(-z\right)=-3z$\\ Número a la derecha del igual & $12$ & $3\cdot\left(12\right)=36$\\ \hline\end{tabular}}\par\smallskip
\par\smallskip Después de hacer todas las multiplicaciones queda:
\[3x + 12y - 3z = 36\]
\par\smallskip\textbf{Qué aprendimos.} La misma condición ahora está escrita con los términos multiplicados. Ya tenemos el tamaño que buscábamos para x.
\par\smallskip\textbf{Comprueba que lo entiendes.} ¿Podríamos multiplicar solamente el término con x?\par\textit{Respuesta:} No. Eso cambiaría la condición. Debemos multiplicar todos los términos, también el número que está después del igual.
\end{minipage}\par\medskip
\noindent\begin{minipage}{\linewidth}\textbf{5. Restamos las dos ecuaciones preparadas}\\[2pt]Restamos la ecuación 1 preparada a la ecuación 2 preparada. Juntamos cada término solo con otro que tenga la misma letra. Los números sin letra se calculan por separado.
\par\smallskip {\small\renewcommand{\arraystretch}{1.3}\begin{tabular}{>{\raggedright\arraybackslash}p{2.80cm}>{\raggedright\arraybackslash}p{5.38cm}>{\raggedright\arraybackslash}p{3.03cm}}\hline\rowcolor{black!7} \textbf{Qué agrupamos} & \textbf{Cuenta} & \textbf{Qué significa}\\\hline Los términos con x & $3x - \left(3x\right)=0$ & Se cancela esta letra.\\ Los términos con y & $8y - \left(12y\right)=-4y$ & Quedan -4y.\\ Los términos con z & $-2z - \left(-3z\right)=z$ & Quedan z.\\ Los números a la derecha del igual & $4 - \left(36\right)=-32$ & Este será el número de la nueva igualdad.\\ \hline\end{tabular}}\par\smallskip
\par\smallskip Primera igualdad que usaremos:
\[3x + 8y - 2z = 4\]
\par\smallskip Le restamos esta igualdad:
\[3x + 12y - 3z = 36\]
\par\smallskip\textbf{Qué aprendimos.} Los términos con x dan cero. Los términos con letras distintas no se mezclan.
\end{minipage}\par\medskip
\noindent\begin{minipage}{\linewidth}\textbf{6. Mira lo que hemos conseguido}\\[2pt]La nueva ecuación es el resultado de la cuenta anterior. Todavía contiene y, z. Seguiremos reduciendo las condiciones.
\par\smallskip Escribimos la nueva condición:
\[-4y + z = -32\]
\[\left\{\begin{aligned}x + 4y - z &= 12\\-4y + z &= -32\end{aligned}\right.\]
\par\smallskip\textbf{Qué aprendimos.} Conservamos la ecuación 1 y reemplazamos la 2 por esta nueva condición. De esa manera conservamos toda la información necesaria para resolver el sistema.
\end{minipage}\par\medskip
\noindent\begin{minipage}{\linewidth}\textbf{7. Una letra puede elegirse libremente}\\[2pt]Las condiciones que quedan no fijan un único valor para z. Podemos elegir ese valor y calcular los demás. Para representar todas las elecciones usamos otra letra, llamada parámetro.
\par\smallskip La letra z será el número que elijamos. Lo llamaremos t:
\[z=t\]
\par\smallskip\textbf{Qué aprendimos.} Un parámetro no es una respuesta que falte: es una forma de describir muchas respuestas a la vez. Primero lo dejamos libre; en un problema real también revisaremos si hay límites como cantidades no negativas.
\end{minipage}\par\medskip
\noindent\begin{minipage}{\linewidth}\textbf{8. Usamos lo que ya sabemos para hallar y}\\[2pt]Ya sabemos z = t. Sustituir significa escribir ese valor cada vez que aparezca su letra. Usamos paréntesis para que se vea qué número estamos multiplicando.
\par\smallskip\textbf{Qué buscamos.} Dejar solamente y como letra desconocida en esta ecuación.
\par\smallskip Tomamos esta condición que obtuvimos antes:
\[-4y + z = -32\]
\par\smallskip Reemplazamos las letras conocidas por sus valores:
\[-4y + \left(t\right)=-32\]
\end{minipage}\par\medskip
\noindent\begin{minipage}{\linewidth}\textbf{9. Calculamos las multiplicaciones que ya se pueden hacer}\\[2pt]La letra y sigue sin conocerse. Las otras partes ya tienen un valor o una expresión conocida y podemos efectuar sus productos.
\par\smallskip {\small\renewcommand{\arraystretch}{1.3}\begin{tabular}{>{\raggedright\arraybackslash}p{3.36cm}>{\raggedright\arraybackslash}p{4.48cm}>{\raggedright\arraybackslash}p{3.36cm}}\hline\rowcolor{black!7} \textbf{Parte conocida} & \textbf{La multiplicación} & \textbf{Resultado}\\\hline El término z & $1\cdot\left(t\right)$ & $t$\\ \hline\end{tabular}}\par\smallskip
\par\smallskip Al reunir esas partes conocidas, la ecuación queda:
\[-4y + t=-32\]
\end{minipage}\par\medskip
\noindent\begin{minipage}{\linewidth}\textbf{10. Dejamos el término con y solo a un lado}\\[2pt]Para quitar la parte conocida, quitamos la expresión t en AMBOS lados. Piensa en una balanza: si hacemos lo mismo a ambos lados, la igualdad se mantiene.
\par\smallskip Hacemos la misma operación a la izquierda y a la derecha:
\[-4y=-32-\left(t\right)\]
\par\smallskip Las partes que se compensan dan cero; queda:
\[-4y=- t - 32\]
\par\smallskip\textbf{Qué aprendimos.} «Pasar al otro lado» es una forma abreviada de decir que hicimos la misma suma o resta en los dos lados.
\end{minipage}\par\medskip
\noindent\begin{minipage}{\linewidth}\textbf{11. Calculamos cuánto vale y}\\[2pt]Para deshacer una multiplicación por -4, dividimos AMBOS lados entre -4. En el lado izquierdo queda una sola y. En el derecho hacemos la división.
\par\smallskip Dividimos los dos lados entre el mismo número:
\[\frac{-4y}{-4}=\frac{- t - 32}{-4}\]
\par\smallskip Ya podemos leer el valor de la letra:
\[y=\frac{- t - 32}{-4}=\frac{t}{4} + 8\]
\par\smallskip\textbf{Qué aprendimos.} Este valor se usará en las ecuaciones donde todavía falten letras.
\end{minipage}\par\medskip
\noindent\begin{minipage}{\linewidth}\textbf{12. Usamos lo que ya sabemos para hallar x}\\[2pt]Ya sabemos y = t/4 + 8; z = t. Sustituir significa escribir ese valor cada vez que aparezca su letra. Usamos paréntesis para que se vea qué número estamos multiplicando.
\par\smallskip\textbf{Qué buscamos.} Dejar solamente x como letra desconocida en esta ecuación.
\par\smallskip Tomamos esta condición que obtuvimos antes:
\[x + 4y - z = 12\]
\par\smallskip Reemplazamos las letras conocidas por sus valores:
\[x + 4\left(\frac{t}{4} + 8\right) - \left(t\right)=12\]
\end{minipage}\par\medskip
\noindent\begin{minipage}{\linewidth}\textbf{13. Calculamos las multiplicaciones que ya se pueden hacer}\\[2pt]La letra x sigue sin conocerse. Las otras partes ya tienen un valor o una expresión conocida y podemos efectuar sus productos.
\par\smallskip {\small\renewcommand{\arraystretch}{1.3}\begin{tabular}{>{\raggedright\arraybackslash}p{3.36cm}>{\raggedright\arraybackslash}p{4.48cm}>{\raggedright\arraybackslash}p{3.36cm}}\hline\rowcolor{black!7} \textbf{Parte conocida} & \textbf{La multiplicación} & \textbf{Resultado}\\\hline El término 4y & $4\cdot\left(\frac{t}{4} + 8\right)$ & $t + 32$\\ El término -z & $-1\cdot\left(t\right)$ & $- t$\\ \hline\end{tabular}}\par\smallskip
\par\smallskip Al reunir esas partes conocidas, la ecuación queda:
\[x + 32=12\]
\end{minipage}\par\medskip
\noindent\begin{minipage}{\linewidth}\textbf{14. Dejamos el término con x solo a un lado}\\[2pt]Para quitar la parte conocida, restamos 32 en AMBOS lados. Piensa en una balanza: si hacemos lo mismo a ambos lados, la igualdad se mantiene.
\par\smallskip Hacemos la misma operación a la izquierda y a la derecha:
\[x + 32 - 32=12 - 32\]
\par\smallskip Las partes que se compensan dan cero; queda:
\[x=-20\]
\par\smallskip\textbf{Qué aprendimos.} «Pasar al otro lado» es una forma abreviada de decir que hicimos la misma suma o resta en los dos lados.
\end{minipage}\par\medskip
\noindent\begin{minipage}{\linewidth}\textbf{15. Leemos el valor de x}\\[2pt]La letra ya está sola: su número multiplicador es 1, que no se escribe. La igualdad nos dice directamente cuánto vale.
\par\smallskip Anotamos el valor encontrado:
\[x=-20\]
\end{minipage}\par\medskip
\noindent\begin{minipage}{\linewidth}\textbf{16. Reunimos la respuesta; todavía falta comprobarla}\\[2pt]Escribimos todos los valores en el orden x, y, z. Dependen del número que elijamos para el parámetro, por eso describen una familia de respuestas.
\[\left\{\begin{aligned}x&=-20\\y&=\frac{t}{4} + 8\\z&=t\end{aligned}\right.\]
\par\smallskip\textbf{Qué aprendimos.} El siguiente paso es volver a las ecuaciones del enunciado y comprobar cada condición.
\end{minipage}\par\medskip
\noindent\begin{minipage}{\linewidth}\textbf{17. Dos elecciones distintas producen dos soluciones}\\[2pt]Esta tabla muestra qué pasa al elegir 0 o 1 para el parámetro. Sustituimos la elección en las expresiones de la respuesta. Cada fila es un grupo completo de valores.
\par\smallskip {\small\renewcommand{\arraystretch}{1.3}\begin{tabular}{>{\raggedright\arraybackslash}p{2.68cm}>{\raggedright\arraybackslash}p{2.68cm}>{\raggedright\arraybackslash}p{2.68cm}>{\raggedright\arraybackslash}p{2.68cm}}\hline\rowcolor{black!7} \textbf{Elección de t} & \textbf{x} & \textbf{y} & \textbf{z}\\\hline 0 & -20 & 8 & 0\\ 1 & -20 & 33/4 & 1\\ \hline\end{tabular}}\par\smallskip
\par\smallskip\textbf{Qué aprendimos.} No hay una sola solución. Hay una para cada elección permitida del parámetro. El deslizador de la web permite comprobar más ejemplos.
\end{minipage}\par\medskip
\noindent\begin{minipage}{\linewidth}\textbf{18. ¿La respuesta cumple la ecuación 1?}\\[2pt]Comprobar significa volver a la condición original y calcularla con nuestra respuesta. Si los dos lados dan el mismo número, esta condición sí se cumple. Deben cumplirse todas.
\par\smallskip La condición del enunciado era:
\[x + 4y - z = 12\]
\par\smallskip Ponemos nuestros valores donde estaban las letras:
\[\left(-20\right) + 4\left(\frac{t}{4} + 8\right) - \left(t\right)=12\]
\par\smallskip Hacemos las multiplicaciones y luego la suma:
\[-20 + \left(t + 32\right) + - t=12\]
\par\smallskip Comparamos los dos lados:
\[12=12\]
\par\smallskip\textbf{Qué aprendimos.} La igualdad final es verdadera. Al simplificar, los términos del parámetro se compensan: la condición vale para cualquier elección permitida.
\end{minipage}\par\medskip
\noindent\begin{minipage}{\linewidth}\textbf{19. ¿La respuesta cumple la ecuación 2?}\\[2pt]Comprobar significa volver a la condición original y calcularla con nuestra respuesta. Si los dos lados dan el mismo número, esta condición sí se cumple. Deben cumplirse todas.
\par\smallskip La condición del enunciado era:
\[3x + 8y - 2z = 4\]
\par\smallskip Ponemos nuestros valores donde estaban las letras:
\[3\left(-20\right) + 8\left(\frac{t}{4} + 8\right) - 2\left(t\right)=4\]
\par\smallskip Hacemos las multiplicaciones y luego la suma:
\[-60 + \left(2 t + 64\right) + - 2 t=4\]
\par\smallskip Comparamos los dos lados:
\[4=4\]
\par\smallskip\textbf{Qué aprendimos.} La igualdad final es verdadera. Al simplificar, los términos del parámetro se compensan: la condición vale para cualquier elección permitida.
\end{minipage}\par\medskip
\respuesta{x = -20; y = t/4 + 8; z = t. Parámetro real.}
\noindent\begin{minipage}{\linewidth}\centering\includegraphics[width=.96\linewidth,height=6.8cm,keepaspectratio]{../recursos/grafico_1_1-7.png}\\{\small Lectura visual complementaria; la demostración está en las operaciones.}\end{minipage}
\end{multicols*}
\newpage\fancyfoot[L]{\small Sección 1.1 | Ejercicio 8 | Paso a paso}\hypertarget{1.1-8}{}
\pdfbookmark[1]{1.1 / 8}{bm1.1-8}
\titulo{8 | Sistema lineal por eliminación}
\textbf{Enunciado}\par Resuelva el sistema lineal dado por medio del método de eliminación.
\[\left\{\begin{aligned}3x + 4y - z &= 8\\6x + 8y - 2z &= 3\end{aligned}\right.\]
\begin{multicols*}{2}
\noindent\begin{minipage}{\linewidth}\textbf{1. Primero: ¿qué significa resolver este sistema?}\\[2pt]Las letras x, y, z representan números que todavía no conocemos. Una solución es una elección de esos números que hace verdaderas TODAS las ecuaciones al mismo tiempo. La llave de la izquierda significa que las condiciones se deben cumplir juntas.
\par\smallskip\textbf{Qué buscamos.} Encontrar un mismo grupo de valores para las 3 letras.
\par\smallskip El número 3 delante de x indica tres veces ese mismo valor; es decir, sumar tres copias:
\[3x=x+x+x\]
\par\smallskip Ejemplo numérico de la notación: si x valiera 2, tres veces ese valor sería 6. Este 2 es solo un ejemplo; todavía no es la respuesta del ejercicio. El punto centrado indica multiplicación:
\[3\cdot 2=2+2+2=6\]
\par\smallskip\textbf{Comprueba que lo entiendes.} ¿Basta con que los valores funcionen en una sola ecuación?\par\textit{Respuesta:} No. Deben funcionar en todas las ecuaciones del sistema, usando los mismos valores.
\end{minipage}\par\medskip
\noindent\begin{minipage}{\linewidth}\textbf{2. Nuestro plan: ir quitando letras}\\[2pt]Una ecuación con una sola letra es más fácil de resolver. Vamos a combinar ecuaciones para que una letra desaparezca. Cuando sepamos un valor, lo pondremos en las ecuaciones que faltan.
\par\smallskip\textbf{Qué buscamos.} Primero quitar x de las otras ecuaciones. Después seguir con las letras restantes.
\par\smallskip\textbf{Por qué.} Si una cantidad se resta de sí misma, queda cero. Usaremos esta idea con términos que tengan la misma letra.
\par\smallskip Ejemplo de la idea que utilizaremos:
\[3x-3x=0\]
\par\smallskip\textbf{Qué aprendimos.} Eliminar una letra de una ecuación NO significa que su valor sea cero; significa que sus términos se cancelaron en esa cuenta.
\end{minipage}\par\medskip
\noindent\begin{minipage}{\linewidth}\textbf{3. ¿Cómo haremos desaparecer x?}\\[2pt]Trabajaremos con las ecuaciones 1 y 2. Queremos que los términos con x tengan el mismo tamaño. Así podremos restarlos.
\par\smallskip\textbf{Qué buscamos.} Conseguir términos con 6 veces x en las dos ecuaciones.
\par\smallskip\textbf{Por qué.} Tener 6 veces la misma letra en ambos lados de la resta nos permitirá cancelarla.
\par\smallskip {\small\renewcommand{\arraystretch}{1.3}\begin{tabular}{>{\raggedright\arraybackslash}p{2.58cm}>{\raggedright\arraybackslash}p{2.47cm}>{\raggedright\arraybackslash}p{6.17cm}}\hline\rowcolor{black!7} \textbf{Dónde miramos} & \textbf{Qué término hay} & \textbf{Qué necesitamos hacer}\\\hline Ecuación 1 & $3x$ & Multiplicar TODA esta ecuación por 2\\ Ecuación 2 & $6x$ & La dejamos igual.\\ \hline\end{tabular}}\par\smallskip
\par\smallskip\textbf{Qué aprendimos.} Estos multiplicadores se eligen para igualar los términos con x. Los demás términos también se multiplican para conservar la igualdad.
\end{minipage}\par\medskip
\noindent\begin{minipage}{\linewidth}\textbf{4. Multiplicamos la ecuación 1 por 2}\\[2pt]Multiplicar una ecuación por 2 significa multiplicar CADA término de la izquierda y también el número de la derecha. La igualdad se conserva porque hacemos lo mismo en ambos lados.
\par\smallskip Esta es la ecuación antes de multiplicar:
\[3x + 4y - z = 8\]
\par\smallskip {\small\renewcommand{\arraystretch}{1.3}\begin{tabular}{>{\raggedright\arraybackslash}p{2.80cm}>{\raggedright\arraybackslash}p{1.91cm}>{\raggedright\arraybackslash}p{6.50cm}}\hline\rowcolor{black!7} \textbf{Parte de la ecuación} & \textbf{Antes} & \textbf{La cuenta y su resultado}\\\hline Término con x & $3x$ & $2\cdot\left(3x\right)=6x$\\ Término con y & $4y$ & $2\cdot\left(4y\right)=8y$\\ Término con z & $-z$ & $2\cdot\left(-z\right)=-2z$\\ Número a la derecha del igual & $8$ & $2\cdot\left(8\right)=16$\\ \hline\end{tabular}}\par\smallskip
\par\smallskip Después de hacer todas las multiplicaciones queda:
\[6x + 8y - 2z = 16\]
\par\smallskip\textbf{Qué aprendimos.} La misma condición ahora está escrita con los términos multiplicados. Ya tenemos el tamaño que buscábamos para x.
\par\smallskip\textbf{Comprueba que lo entiendes.} ¿Podríamos multiplicar solamente el término con x?\par\textit{Respuesta:} No. Eso cambiaría la condición. Debemos multiplicar todos los términos, también el número que está después del igual.
\end{minipage}\par\medskip
\noindent\begin{minipage}{\linewidth}\textbf{5. Restamos las dos ecuaciones preparadas}\\[2pt]Restamos la ecuación 1 preparada a la ecuación 2 preparada. Juntamos cada término solo con otro que tenga la misma letra. Los números sin letra se calculan por separado.
\par\smallskip {\small\renewcommand{\arraystretch}{1.3}\begin{tabular}{>{\raggedright\arraybackslash}p{2.80cm}>{\raggedright\arraybackslash}p{5.38cm}>{\raggedright\arraybackslash}p{3.03cm}}\hline\rowcolor{black!7} \textbf{Qué agrupamos} & \textbf{Cuenta} & \textbf{Qué significa}\\\hline Los términos con x & $6x - \left(6x\right)=0$ & Se cancela esta letra.\\ Los términos con y & $8y - \left(8y\right)=0$ & Se cancela esta letra.\\ Los términos con z & $-2z - \left(-2z\right)=0$ & Se cancela esta letra.\\ Los números a la derecha del igual & $3 - \left(16\right)=-13$ & Este será el número de la nueva igualdad.\\ \hline\end{tabular}}\par\smallskip
\par\smallskip Primera igualdad que usaremos:
\[6x + 8y - 2z = 3\]
\par\smallskip Le restamos esta igualdad:
\[6x + 8y - 2z = 16\]
\par\smallskip\textbf{Qué aprendimos.} Los términos con x dan cero. Los términos con letras distintas no se mezclan.
\end{minipage}\par\medskip
\noindent\begin{minipage}{\linewidth}\textbf{6. Mira lo que hemos conseguido}\\[2pt]La nueva ecuación es el resultado de la cuenta anterior. Ya no contiene letras. Tenemos que decidir si la igualdad numérica es verdadera o imposible.
\par\smallskip Escribimos la nueva condición:
\[0 = -13\]
\[\left\{\begin{aligned}3x + 4y - z &= 8\\0 &= -13\end{aligned}\right.\]
\par\smallskip\textbf{Qué aprendimos.} Conservamos la ecuación 1 y reemplazamos la 2 por esta nueva condición. De esa manera conservamos toda la información necesaria para resolver el sistema.
\end{minipage}\par\medskip
\noindent\begin{minipage}{\linewidth}\textbf{7. Esta igualdad nos dice que no hay solución}\\[2pt]Al hacer las cuentas, todas las letras desaparecieron y quedó la igualdad que ves abajo. Cero no puede ser igual a -13. Como esa condición tendría que cumplirse junto con las demás, ningún grupo de valores puede resolver el sistema.
\par\smallskip Esta es la contradicción encontrada:
\[0 = -13\]
\par\smallskip\textbf{Qué aprendimos.} El resultado es «sin solución». No significa que debamos responder con ceros: significa que no existe una respuesta que cumpla todas las condiciones.
\par\smallskip\textbf{Comprueba que lo entiendes.} ¿Responder x = 0, y = 0, z = 0 resolvería esta contradicción?\par\textit{Respuesta:} No. La contradicción ya no depende de las letras. Ningún valor puede convertir cero en otro número.
\end{minipage}\par\medskip
\respuesta{Sin solución.}
\noindent\begin{minipage}{\linewidth}\centering\includegraphics[width=.96\linewidth,height=6.8cm,keepaspectratio]{../recursos/grafico_1_1-8.png}\\{\small Lectura visual complementaria; la demostración está en las operaciones.}\end{minipage}
\end{multicols*}
\newpage\fancyfoot[L]{\small Sección 1.1 | Ejercicio 9 | Paso a paso}\hypertarget{1.1-9}{}
\pdfbookmark[1]{1.1 / 9}{bm1.1-9}
\titulo{9 | Sistema lineal por eliminación}
\textbf{Enunciado}\par Resuelva el sistema lineal dado por medio del método de eliminación.
\[\left\{\begin{aligned}x + y + 3z &= 12\\2x + 2y + 6z &= 6\end{aligned}\right.\]
\begin{multicols*}{2}
\noindent\begin{minipage}{\linewidth}\textbf{1. Primero: ¿qué significa resolver este sistema?}\\[2pt]Las letras x, y, z representan números que todavía no conocemos. Una solución es una elección de esos números que hace verdaderas TODAS las ecuaciones al mismo tiempo. La llave de la izquierda significa que las condiciones se deben cumplir juntas.
\par\smallskip\textbf{Qué buscamos.} Encontrar un mismo grupo de valores para las 3 letras.
\par\smallskip El número 3 delante de x indica tres veces ese mismo valor; es decir, sumar tres copias:
\[3x=x+x+x\]
\par\smallskip Ejemplo numérico de la notación: si x valiera 2, tres veces ese valor sería 6. Este 2 es solo un ejemplo; todavía no es la respuesta del ejercicio. El punto centrado indica multiplicación:
\[3\cdot 2=2+2+2=6\]
\par\smallskip\textbf{Comprueba que lo entiendes.} ¿Basta con que los valores funcionen en una sola ecuación?\par\textit{Respuesta:} No. Deben funcionar en todas las ecuaciones del sistema, usando los mismos valores.
\end{minipage}\par\medskip
\noindent\begin{minipage}{\linewidth}\textbf{2. Nuestro plan: ir quitando letras}\\[2pt]Una ecuación con una sola letra es más fácil de resolver. Vamos a combinar ecuaciones para que una letra desaparezca. Cuando sepamos un valor, lo pondremos en las ecuaciones que faltan.
\par\smallskip\textbf{Qué buscamos.} Primero quitar x de las otras ecuaciones. Después seguir con las letras restantes.
\par\smallskip\textbf{Por qué.} Si una cantidad se resta de sí misma, queda cero. Usaremos esta idea con términos que tengan la misma letra.
\par\smallskip Ejemplo de la idea que utilizaremos:
\[3x-3x=0\]
\par\smallskip\textbf{Qué aprendimos.} Eliminar una letra de una ecuación NO significa que su valor sea cero; significa que sus términos se cancelaron en esa cuenta.
\end{minipage}\par\medskip
\noindent\begin{minipage}{\linewidth}\textbf{3. ¿Cómo haremos desaparecer x?}\\[2pt]Trabajaremos con las ecuaciones 1 y 2. Queremos que los términos con x tengan el mismo tamaño. Así podremos restarlos.
\par\smallskip\textbf{Qué buscamos.} Conseguir términos con 2 veces x en las dos ecuaciones.
\par\smallskip\textbf{Por qué.} Tener 2 veces la misma letra en ambos lados de la resta nos permitirá cancelarla.
\par\smallskip {\small\renewcommand{\arraystretch}{1.3}\begin{tabular}{>{\raggedright\arraybackslash}p{2.58cm}>{\raggedright\arraybackslash}p{2.47cm}>{\raggedright\arraybackslash}p{6.17cm}}\hline\rowcolor{black!7} \textbf{Dónde miramos} & \textbf{Qué término hay} & \textbf{Qué necesitamos hacer}\\\hline Ecuación 1 & $x$ & Multiplicar TODA esta ecuación por 2\\ Ecuación 2 & $2x$ & La dejamos igual.\\ \hline\end{tabular}}\par\smallskip
\par\smallskip\textbf{Qué aprendimos.} Estos multiplicadores se eligen para igualar los términos con x. Los demás términos también se multiplican para conservar la igualdad.
\end{minipage}\par\medskip
\noindent\begin{minipage}{\linewidth}\textbf{4. Multiplicamos la ecuación 1 por 2}\\[2pt]Multiplicar una ecuación por 2 significa multiplicar CADA término de la izquierda y también el número de la derecha. La igualdad se conserva porque hacemos lo mismo en ambos lados.
\par\smallskip Esta es la ecuación antes de multiplicar:
\[x + y + 3z = 12\]
\par\smallskip {\small\renewcommand{\arraystretch}{1.3}\begin{tabular}{>{\raggedright\arraybackslash}p{2.80cm}>{\raggedright\arraybackslash}p{1.91cm}>{\raggedright\arraybackslash}p{6.50cm}}\hline\rowcolor{black!7} \textbf{Parte de la ecuación} & \textbf{Antes} & \textbf{La cuenta y su resultado}\\\hline Término con x & $x$ & $2\cdot\left(x\right)=2x$\\ Término con y & $y$ & $2\cdot\left(y\right)=2y$\\ Término con z & $3z$ & $2\cdot\left(3z\right)=6z$\\ Número a la derecha del igual & $12$ & $2\cdot\left(12\right)=24$\\ \hline\end{tabular}}\par\smallskip
\par\smallskip Después de hacer todas las multiplicaciones queda:
\[2x + 2y + 6z = 24\]
\par\smallskip\textbf{Qué aprendimos.} La misma condición ahora está escrita con los términos multiplicados. Ya tenemos el tamaño que buscábamos para x.
\par\smallskip\textbf{Comprueba que lo entiendes.} ¿Podríamos multiplicar solamente el término con x?\par\textit{Respuesta:} No. Eso cambiaría la condición. Debemos multiplicar todos los términos, también el número que está después del igual.
\end{minipage}\par\medskip
\noindent\begin{minipage}{\linewidth}\textbf{5. Restamos las dos ecuaciones preparadas}\\[2pt]Restamos la ecuación 1 preparada a la ecuación 2 preparada. Juntamos cada término solo con otro que tenga la misma letra. Los números sin letra se calculan por separado.
\par\smallskip {\small\renewcommand{\arraystretch}{1.3}\begin{tabular}{>{\raggedright\arraybackslash}p{2.80cm}>{\raggedright\arraybackslash}p{5.38cm}>{\raggedright\arraybackslash}p{3.03cm}}\hline\rowcolor{black!7} \textbf{Qué agrupamos} & \textbf{Cuenta} & \textbf{Qué significa}\\\hline Los términos con x & $2x - \left(2x\right)=0$ & Se cancela esta letra.\\ Los términos con y & $2y - \left(2y\right)=0$ & Se cancela esta letra.\\ Los términos con z & $6z - \left(6z\right)=0$ & Se cancela esta letra.\\ Los números a la derecha del igual & $6 - \left(24\right)=-18$ & Este será el número de la nueva igualdad.\\ \hline\end{tabular}}\par\smallskip
\par\smallskip Primera igualdad que usaremos:
\[2x + 2y + 6z = 6\]
\par\smallskip Le restamos esta igualdad:
\[2x + 2y + 6z = 24\]
\par\smallskip\textbf{Qué aprendimos.} Los términos con x dan cero. Los términos con letras distintas no se mezclan.
\end{minipage}\par\medskip
\noindent\begin{minipage}{\linewidth}\textbf{6. Mira lo que hemos conseguido}\\[2pt]La nueva ecuación es el resultado de la cuenta anterior. Ya no contiene letras. Tenemos que decidir si la igualdad numérica es verdadera o imposible.
\par\smallskip Escribimos la nueva condición:
\[0 = -18\]
\[\left\{\begin{aligned}x + y + 3z &= 12\\0 &= -18\end{aligned}\right.\]
\par\smallskip\textbf{Qué aprendimos.} Conservamos la ecuación 1 y reemplazamos la 2 por esta nueva condición. De esa manera conservamos toda la información necesaria para resolver el sistema.
\end{minipage}\par\medskip
\noindent\begin{minipage}{\linewidth}\textbf{7. Esta igualdad nos dice que no hay solución}\\[2pt]Al hacer las cuentas, todas las letras desaparecieron y quedó la igualdad que ves abajo. Cero no puede ser igual a -18. Como esa condición tendría que cumplirse junto con las demás, ningún grupo de valores puede resolver el sistema.
\par\smallskip Esta es la contradicción encontrada:
\[0 = -18\]
\par\smallskip\textbf{Qué aprendimos.} El resultado es «sin solución». No significa que debamos responder con ceros: significa que no existe una respuesta que cumpla todas las condiciones.
\par\smallskip\textbf{Comprueba que lo entiendes.} ¿Responder x = 0, y = 0, z = 0 resolvería esta contradicción?\par\textit{Respuesta:} No. La contradicción ya no depende de las letras. Ningún valor puede convertir cero en otro número.
\end{minipage}\par\medskip
\respuesta{Sin solución.}
\noindent\begin{minipage}{\linewidth}\centering\includegraphics[width=.96\linewidth,height=6.8cm,keepaspectratio]{../recursos/grafico_1_1-9.png}\\{\small Lectura visual complementaria; la demostración está en las operaciones.}\end{minipage}
\end{multicols*}
\newpage\fancyfoot[L]{\small Sección 1.1 | Ejercicio 10 | Paso a paso}\hypertarget{1.1-10}{}
\pdfbookmark[1]{1.1 / 10}{bm1.1-10}
\titulo{10 | Sistema lineal por eliminación}
\textbf{Enunciado}\par Resuelva el sistema lineal dado por medio del método de eliminación.
\[\left\{\begin{aligned}x + y &= 1\\2x - y &= 5\\3x + 4y &= 2\end{aligned}\right.\]
\begin{multicols*}{2}
\noindent\begin{minipage}{\linewidth}\textbf{1. Primero: ¿qué significa resolver este sistema?}\\[2pt]Las letras x, y representan números que todavía no conocemos. Una solución es una elección de esos números que hace verdaderas TODAS las ecuaciones al mismo tiempo. La llave de la izquierda significa que las condiciones se deben cumplir juntas.
\par\smallskip\textbf{Qué buscamos.} Encontrar una misma pareja de valores.
\par\smallskip El número 3 delante de x indica tres veces ese mismo valor; es decir, sumar tres copias:
\[3x=x+x+x\]
\par\smallskip Ejemplo numérico de la notación: si x valiera 2, tres veces ese valor sería 6. Este 2 es solo un ejemplo; todavía no es la respuesta del ejercicio. El punto centrado indica multiplicación:
\[3\cdot 2=2+2+2=6\]
\par\smallskip\textbf{Comprueba que lo entiendes.} ¿Basta con que los valores funcionen en una sola ecuación?\par\textit{Respuesta:} No. Deben funcionar en todas las ecuaciones del sistema, usando los mismos valores.
\end{minipage}\par\medskip
\noindent\begin{minipage}{\linewidth}\textbf{2. Nuestro plan: ir quitando letras}\\[2pt]Una ecuación con una sola letra es más fácil de resolver. Vamos a combinar ecuaciones para que una letra desaparezca. Cuando sepamos un valor, lo pondremos en las ecuaciones que faltan.
\par\smallskip\textbf{Qué buscamos.} Primero quitar x para calcular y. Después calcular x.
\par\smallskip\textbf{Por qué.} Si una cantidad se resta de sí misma, queda cero. Usaremos esta idea con términos que tengan la misma letra.
\par\smallskip Ejemplo de la idea que utilizaremos:
\[3x-3x=0\]
\par\smallskip\textbf{Qué aprendimos.} Eliminar una letra de una ecuación NO significa que su valor sea cero; significa que sus términos se cancelaron en esa cuenta.
\end{minipage}\par\medskip
\noindent\begin{minipage}{\linewidth}\textbf{3. ¿Cómo haremos desaparecer x?}\\[2pt]Trabajaremos con las ecuaciones 1 y 2. Queremos que los términos con x tengan el mismo tamaño. Así podremos restarlos.
\par\smallskip\textbf{Qué buscamos.} Conseguir términos con 2 veces x en las dos ecuaciones.
\par\smallskip\textbf{Por qué.} Tener 2 veces la misma letra en ambos lados de la resta nos permitirá cancelarla.
\par\smallskip {\small\renewcommand{\arraystretch}{1.3}\begin{tabular}{>{\raggedright\arraybackslash}p{2.58cm}>{\raggedright\arraybackslash}p{2.47cm}>{\raggedright\arraybackslash}p{6.17cm}}\hline\rowcolor{black!7} \textbf{Dónde miramos} & \textbf{Qué término hay} & \textbf{Qué necesitamos hacer}\\\hline Ecuación 1 & $x$ & Multiplicar TODA esta ecuación por 2\\ Ecuación 2 & $2x$ & La dejamos igual.\\ \hline\end{tabular}}\par\smallskip
\par\smallskip\textbf{Qué aprendimos.} Estos multiplicadores se eligen para igualar los términos con x. Los demás términos también se multiplican para conservar la igualdad.
\end{minipage}\par\medskip
\noindent\begin{minipage}{\linewidth}\textbf{4. Multiplicamos la ecuación 1 por 2}\\[2pt]Multiplicar una ecuación por 2 significa multiplicar CADA término de la izquierda y también el número de la derecha. La igualdad se conserva porque hacemos lo mismo en ambos lados.
\par\smallskip Esta es la ecuación antes de multiplicar:
\[x + y = 1\]
\par\smallskip {\small\renewcommand{\arraystretch}{1.3}\begin{tabular}{>{\raggedright\arraybackslash}p{2.80cm}>{\raggedright\arraybackslash}p{1.91cm}>{\raggedright\arraybackslash}p{6.50cm}}\hline\rowcolor{black!7} \textbf{Parte de la ecuación} & \textbf{Antes} & \textbf{La cuenta y su resultado}\\\hline Término con x & $x$ & $2\cdot\left(x\right)=2x$\\ Término con y & $y$ & $2\cdot\left(y\right)=2y$\\ Número a la derecha del igual & $1$ & $2\cdot\left(1\right)=2$\\ \hline\end{tabular}}\par\smallskip
\par\smallskip Después de hacer todas las multiplicaciones queda:
\[2x + 2y = 2\]
\par\smallskip\textbf{Qué aprendimos.} La misma condición ahora está escrita con los términos multiplicados. Ya tenemos el tamaño que buscábamos para x.
\par\smallskip\textbf{Comprueba que lo entiendes.} ¿Podríamos multiplicar solamente el término con x?\par\textit{Respuesta:} No. Eso cambiaría la condición. Debemos multiplicar todos los términos, también el número que está después del igual.
\end{minipage}\par\medskip
\noindent\begin{minipage}{\linewidth}\textbf{5. Restamos las dos ecuaciones preparadas}\\[2pt]Restamos la ecuación 1 preparada a la ecuación 2 preparada. Juntamos cada término solo con otro que tenga la misma letra. Los números sin letra se calculan por separado.
\par\smallskip {\small\renewcommand{\arraystretch}{1.3}\begin{tabular}{>{\raggedright\arraybackslash}p{2.80cm}>{\raggedright\arraybackslash}p{5.38cm}>{\raggedright\arraybackslash}p{3.03cm}}\hline\rowcolor{black!7} \textbf{Qué agrupamos} & \textbf{Cuenta} & \textbf{Qué significa}\\\hline Los términos con x & $2x - \left(2x\right)=0$ & Se cancela esta letra.\\ Los términos con y & $-y - \left(2y\right)=-3y$ & Quedan -3y.\\ Los números a la derecha del igual & $5 - \left(2\right)=3$ & Este será el número de la nueva igualdad.\\ \hline\end{tabular}}\par\smallskip
\par\smallskip Primera igualdad que usaremos:
\[2x - y = 5\]
\par\smallskip Le restamos esta igualdad:
\[2x + 2y = 2\]
\par\smallskip\textbf{Qué aprendimos.} Los términos con x dan cero. Los términos con letras distintas no se mezclan.
\end{minipage}\par\medskip
\noindent\begin{minipage}{\linewidth}\textbf{6. Mira lo que hemos conseguido}\\[2pt]La nueva ecuación es el resultado de la cuenta anterior. Solo contiene y: después podremos despejarla.
\par\smallskip Escribimos la nueva condición:
\[-3y = 3\]
\[\left\{\begin{aligned}x + y &= 1\\-3y &= 3\\3x + 4y &= 2\end{aligned}\right.\]
\par\smallskip\textbf{Qué aprendimos.} Conservamos la ecuación 1 y reemplazamos la 2 por esta nueva condición. De esa manera conservamos toda la información necesaria para resolver el sistema.
\end{minipage}\par\medskip
\noindent\begin{minipage}{\linewidth}\textbf{7. ¿Cómo haremos desaparecer x?}\\[2pt]Trabajaremos con las ecuaciones 1 y 3. Queremos que los términos con x tengan el mismo tamaño. Así podremos restarlos.
\par\smallskip\textbf{Qué buscamos.} Conseguir términos con 3 veces x en las dos ecuaciones.
\par\smallskip\textbf{Por qué.} Tener 3 veces la misma letra en ambos lados de la resta nos permitirá cancelarla.
\par\smallskip {\small\renewcommand{\arraystretch}{1.3}\begin{tabular}{>{\raggedright\arraybackslash}p{2.58cm}>{\raggedright\arraybackslash}p{2.47cm}>{\raggedright\arraybackslash}p{6.17cm}}\hline\rowcolor{black!7} \textbf{Dónde miramos} & \textbf{Qué término hay} & \textbf{Qué necesitamos hacer}\\\hline Ecuación 1 & $x$ & Multiplicar TODA esta ecuación por 3\\ Ecuación 3 & $3x$ & La dejamos igual.\\ \hline\end{tabular}}\par\smallskip
\par\smallskip\textbf{Qué aprendimos.} Estos multiplicadores se eligen para igualar los términos con x. Los demás términos también se multiplican para conservar la igualdad.
\end{minipage}\par\medskip
\noindent\begin{minipage}{\linewidth}\textbf{8. Multiplicamos la ecuación 1 por 3}\\[2pt]Multiplicar una ecuación por 3 significa multiplicar CADA término de la izquierda y también el número de la derecha. La igualdad se conserva porque hacemos lo mismo en ambos lados.
\par\smallskip Esta es la ecuación antes de multiplicar:
\[x + y = 1\]
\par\smallskip {\small\renewcommand{\arraystretch}{1.3}\begin{tabular}{>{\raggedright\arraybackslash}p{2.80cm}>{\raggedright\arraybackslash}p{1.91cm}>{\raggedright\arraybackslash}p{6.50cm}}\hline\rowcolor{black!7} \textbf{Parte de la ecuación} & \textbf{Antes} & \textbf{La cuenta y su resultado}\\\hline Término con x & $x$ & $3\cdot\left(x\right)=3x$\\ Término con y & $y$ & $3\cdot\left(y\right)=3y$\\ Número a la derecha del igual & $1$ & $3\cdot\left(1\right)=3$\\ \hline\end{tabular}}\par\smallskip
\par\smallskip Después de hacer todas las multiplicaciones queda:
\[3x + 3y = 3\]
\par\smallskip\textbf{Qué aprendimos.} La misma condición ahora está escrita con los términos multiplicados. Ya tenemos el tamaño que buscábamos para x.
\par\smallskip\textbf{Comprueba que lo entiendes.} ¿Podríamos multiplicar solamente el término con x?\par\textit{Respuesta:} No. Eso cambiaría la condición. Debemos multiplicar todos los términos, también el número que está después del igual.
\end{minipage}\par\medskip
\noindent\begin{minipage}{\linewidth}\textbf{9. Restamos las dos ecuaciones preparadas}\\[2pt]Restamos la ecuación 1 preparada a la ecuación 3 preparada. Juntamos cada término solo con otro que tenga la misma letra. Los números sin letra se calculan por separado.
\par\smallskip {\small\renewcommand{\arraystretch}{1.3}\begin{tabular}{>{\raggedright\arraybackslash}p{2.80cm}>{\raggedright\arraybackslash}p{5.38cm}>{\raggedright\arraybackslash}p{3.03cm}}\hline\rowcolor{black!7} \textbf{Qué agrupamos} & \textbf{Cuenta} & \textbf{Qué significa}\\\hline Los términos con x & $3x - \left(3x\right)=0$ & Se cancela esta letra.\\ Los términos con y & $4y - \left(3y\right)=y$ & Quedan y.\\ Los números a la derecha del igual & $2 - \left(3\right)=-1$ & Este será el número de la nueva igualdad.\\ \hline\end{tabular}}\par\smallskip
\par\smallskip Primera igualdad que usaremos:
\[3x + 4y = 2\]
\par\smallskip Le restamos esta igualdad:
\[3x + 3y = 3\]
\par\smallskip\textbf{Qué aprendimos.} Los términos con x dan cero. Los términos con letras distintas no se mezclan.
\end{minipage}\par\medskip
\noindent\begin{minipage}{\linewidth}\textbf{10. Mira lo que hemos conseguido}\\[2pt]La nueva ecuación es el resultado de la cuenta anterior. Solo contiene y: después podremos despejarla.
\par\smallskip Escribimos la nueva condición:
\[y = -1\]
\[\left\{\begin{aligned}x + y &= 1\\-3y &= 3\\y &= -1\end{aligned}\right.\]
\par\smallskip\textbf{Qué aprendimos.} Conservamos la ecuación 1 y reemplazamos la 3 por esta nueva condición. De esa manera conservamos toda la información necesaria para resolver el sistema.
\end{minipage}\par\medskip
\noindent\begin{minipage}{\linewidth}\textbf{11. Ponemos una ecuación más cómoda arriba}\\[2pt]Cambiamos de lugar las ecuaciones 2 y 3. Las condiciones siguen siendo las mismas; solamente cambia el orden en que las leemos.
\par\smallskip\textbf{Por qué.} La ecuación que subimos tiene un número más pequeño delante de y. Eso facilita las multiplicaciones que haremos a continuación.
\[\left\{\begin{aligned}x + y &= 1\\y &= -1\\-3y &= 3\end{aligned}\right.\]
\end{minipage}\par\medskip
\noindent\begin{minipage}{\linewidth}\textbf{12. ¿Cómo haremos desaparecer y?}\\[2pt]Trabajaremos con las ecuaciones 2 y 3. Queremos que los términos con y tengan el mismo tamaño. Así podremos sumarlos, porque tienen signos opuestos.
\par\smallskip\textbf{Qué buscamos.} Conseguir términos con 3 veces y en las dos ecuaciones.
\par\smallskip\textbf{Por qué.} Al sumar la misma cantidad con signos opuestos obtenemos cero.
\par\smallskip {\small\renewcommand{\arraystretch}{1.3}\begin{tabular}{>{\raggedright\arraybackslash}p{2.58cm}>{\raggedright\arraybackslash}p{2.47cm}>{\raggedright\arraybackslash}p{6.17cm}}\hline\rowcolor{black!7} \textbf{Dónde miramos} & \textbf{Qué término hay} & \textbf{Qué necesitamos hacer}\\\hline Ecuación 2 & $y$ & Multiplicar TODA esta ecuación por 3\\ Ecuación 3 & $-3y$ & La dejamos igual.\\ \hline\end{tabular}}\par\smallskip
\par\smallskip\textbf{Qué aprendimos.} Estos multiplicadores se eligen para igualar los términos con y. Los demás términos también se multiplican para conservar la igualdad.
\end{minipage}\par\medskip
\noindent\begin{minipage}{\linewidth}\textbf{13. Multiplicamos la ecuación 2 por 3}\\[2pt]Multiplicar una ecuación por 3 significa multiplicar CADA término de la izquierda y también el número de la derecha. La igualdad se conserva porque hacemos lo mismo en ambos lados.
\par\smallskip Esta es la ecuación antes de multiplicar:
\[y = -1\]
\par\smallskip {\small\renewcommand{\arraystretch}{1.3}\begin{tabular}{>{\raggedright\arraybackslash}p{2.80cm}>{\raggedright\arraybackslash}p{1.91cm}>{\raggedright\arraybackslash}p{6.50cm}}\hline\rowcolor{black!7} \textbf{Parte de la ecuación} & \textbf{Antes} & \textbf{La cuenta y su resultado}\\\hline Término con y & $y$ & $3\cdot\left(y\right)=3y$\\ Número a la derecha del igual & $-1$ & $3\cdot\left(-1\right)=-3$\\ \hline\end{tabular}}\par\smallskip
\par\smallskip Después de hacer todas las multiplicaciones queda:
\[3y = -3\]
\par\smallskip\textbf{Qué aprendimos.} La misma condición ahora está escrita con los términos multiplicados. Ya tenemos el tamaño que buscábamos para y.
\par\smallskip\textbf{Comprueba que lo entiendes.} ¿Podríamos multiplicar solamente el término con y?\par\textit{Respuesta:} No. Eso cambiaría la condición. Debemos multiplicar todos los términos, también el número que está después del igual.
\end{minipage}\par\medskip
\noindent\begin{minipage}{\linewidth}\textbf{14. Sumamos las dos ecuaciones preparadas}\\[2pt]Sumamos las dos ecuaciones preparadas. Juntamos cada término solo con otro que tenga la misma letra. Los números sin letra se calculan por separado.
\par\smallskip {\small\renewcommand{\arraystretch}{1.3}\begin{tabular}{>{\raggedright\arraybackslash}p{2.80cm}>{\raggedright\arraybackslash}p{5.38cm}>{\raggedright\arraybackslash}p{3.03cm}}\hline\rowcolor{black!7} \textbf{Qué agrupamos} & \textbf{Cuenta} & \textbf{Qué significa}\\\hline Los términos con y & $-3y + \left(3y\right)=0$ & Se cancela esta letra.\\ Los números a la derecha del igual & $3 + \left(-3\right)=0$ & Este será el número de la nueva igualdad.\\ \hline\end{tabular}}\par\smallskip
\par\smallskip Primera igualdad que usaremos:
\[-3y = 3\]
\par\smallskip Le sumamos esta igualdad:
\[3y = -3\]
\par\smallskip\textbf{Qué aprendimos.} Los términos con y dan cero. Los términos con letras distintas no se mezclan.
\end{minipage}\par\medskip
\noindent\begin{minipage}{\linewidth}\textbf{15. Mira lo que hemos conseguido}\\[2pt]La nueva ecuación es el resultado de la cuenta anterior. Ya no contiene letras. Tenemos que decidir si la igualdad numérica es verdadera o imposible.
\par\smallskip Escribimos la nueva condición:
\[0 = 0\]
\[\left\{\begin{aligned}x + y &= 1\\y &= -1\\0 &= 0\end{aligned}\right.\]
\par\smallskip\textbf{Qué aprendimos.} Conservamos la ecuación 2 y reemplazamos la 3 por esta nueva condición. De esa manera conservamos toda la información necesaria para resolver el sistema.
\end{minipage}\par\medskip
\noindent\begin{minipage}{\linewidth}\textbf{16. Ya tenemos una ecuación con una sola letra}\\[2pt]En esta condición solamente falta conocer y. El número 1 está multiplicando a esa letra. Ahora podemos deshacer esa multiplicación.
\par\smallskip Esta es la condición que vamos a resolver:
\[y = -1\]
\end{minipage}\par\medskip
\noindent\begin{minipage}{\linewidth}\textbf{17. Leemos el valor de y}\\[2pt]La letra ya está sola: su número multiplicador es 1, que no se escribe. La igualdad nos dice directamente cuánto vale.
\par\smallskip Anotamos el valor encontrado:
\[y=-1\]
\end{minipage}\par\medskip
\noindent\begin{minipage}{\linewidth}\textbf{18. Usamos lo que ya sabemos para hallar x}\\[2pt]Ya sabemos y = -1. Sustituir significa escribir ese valor cada vez que aparezca su letra. Usamos paréntesis para que se vea qué número estamos multiplicando.
\par\smallskip\textbf{Qué buscamos.} Dejar solamente x como letra desconocida en esta ecuación.
\par\smallskip Tomamos esta condición que obtuvimos antes:
\[x + y = 1\]
\par\smallskip Reemplazamos las letras conocidas por sus valores:
\[x + \left(-1\right)=1\]
\end{minipage}\par\medskip
\noindent\begin{minipage}{\linewidth}\textbf{19. Calculamos las multiplicaciones que ya se pueden hacer}\\[2pt]La letra x sigue sin conocerse. Las otras partes ya tienen un valor o una expresión conocida y podemos efectuar sus productos.
\par\smallskip {\small\renewcommand{\arraystretch}{1.3}\begin{tabular}{>{\raggedright\arraybackslash}p{3.36cm}>{\raggedright\arraybackslash}p{4.48cm}>{\raggedright\arraybackslash}p{3.36cm}}\hline\rowcolor{black!7} \textbf{Parte conocida} & \textbf{La multiplicación} & \textbf{Resultado}\\\hline El término y & $1\cdot\left(-1\right)$ & $-1$\\ \hline\end{tabular}}\par\smallskip
\par\smallskip Al reunir esas partes conocidas, la ecuación queda:
\[x - 1=1\]
\end{minipage}\par\medskip
\noindent\begin{minipage}{\linewidth}\textbf{20. Dejamos el término con x solo a un lado}\\[2pt]Para quitar la parte conocida, sumamos 1 en AMBOS lados. Piensa en una balanza: si hacemos lo mismo a ambos lados, la igualdad se mantiene.
\par\smallskip Hacemos la misma operación a la izquierda y a la derecha:
\[x - 1 + 1=1 + 1\]
\par\smallskip Las partes que se compensan dan cero; queda:
\[x=2\]
\par\smallskip\textbf{Qué aprendimos.} «Pasar al otro lado» es una forma abreviada de decir que hicimos la misma suma o resta en los dos lados.
\end{minipage}\par\medskip
\noindent\begin{minipage}{\linewidth}\textbf{21. Leemos el valor de x}\\[2pt]La letra ya está sola: su número multiplicador es 1, que no se escribe. La igualdad nos dice directamente cuánto vale.
\par\smallskip Anotamos el valor encontrado:
\[x=2\]
\end{minipage}\par\medskip
\noindent\begin{minipage}{\linewidth}\textbf{22. Reunimos la respuesta; todavía falta comprobarla}\\[2pt]Escribimos todos los valores en el orden x, y. Estos valores deben funcionar juntos, no cada uno en una ecuación diferente.
\[\left\{\begin{aligned}x&=2\\y&=-1\end{aligned}\right.\]
\par\smallskip\textbf{Qué aprendimos.} El siguiente paso es volver a las ecuaciones del enunciado y comprobar cada condición.
\end{minipage}\par\medskip
\noindent\begin{minipage}{\linewidth}\textbf{23. ¿La respuesta cumple la ecuación 1?}\\[2pt]Comprobar significa volver a la condición original y calcularla con nuestra respuesta. Si los dos lados dan el mismo número, esta condición sí se cumple. Deben cumplirse todas.
\par\smallskip La condición del enunciado era:
\[x + y = 1\]
\par\smallskip Ponemos nuestros valores donde estaban las letras:
\[\left(2\right) + \left(-1\right)=1\]
\par\smallskip Hacemos las multiplicaciones y luego la suma:
\[2 - 1=1\]
\par\smallskip Comparamos los dos lados:
\[1=1\]
\par\smallskip\textbf{Qué aprendimos.} La igualdad final es verdadera. Nuestra respuesta sí cumple la ecuación 1.
\end{minipage}\par\medskip
\noindent\begin{minipage}{\linewidth}\textbf{24. ¿La respuesta cumple la ecuación 2?}\\[2pt]Comprobar significa volver a la condición original y calcularla con nuestra respuesta. Si los dos lados dan el mismo número, esta condición sí se cumple. Deben cumplirse todas.
\par\smallskip La condición del enunciado era:
\[2x - y = 5\]
\par\smallskip Ponemos nuestros valores donde estaban las letras:
\[2\left(2\right) - \left(-1\right)=5\]
\par\smallskip Hacemos las multiplicaciones y luego la suma:
\[4 + 1=5\]
\par\smallskip Comparamos los dos lados:
\[5=5\]
\par\smallskip\textbf{Qué aprendimos.} La igualdad final es verdadera. Nuestra respuesta sí cumple la ecuación 2.
\end{minipage}\par\medskip
\noindent\begin{minipage}{\linewidth}\textbf{25. ¿La respuesta cumple la ecuación 3?}\\[2pt]Comprobar significa volver a la condición original y calcularla con nuestra respuesta. Si los dos lados dan el mismo número, esta condición sí se cumple. Deben cumplirse todas.
\par\smallskip La condición del enunciado era:
\[3x + 4y = 2\]
\par\smallskip Ponemos nuestros valores donde estaban las letras:
\[3\left(2\right) + 4\left(-1\right)=2\]
\par\smallskip Hacemos las multiplicaciones y luego la suma:
\[6 - 4=2\]
\par\smallskip Comparamos los dos lados:
\[2=2\]
\par\smallskip\textbf{Qué aprendimos.} La igualdad final es verdadera. Nuestra respuesta sí cumple la ecuación 3.
\end{minipage}\par\medskip
\respuesta{x = 2; y = -1.}
\noindent\begin{minipage}{\linewidth}\centering\includegraphics[width=.96\linewidth,height=6.8cm,keepaspectratio]{../recursos/grafico_1_1-10.png}\\{\small Lectura visual complementaria; la demostración está en las operaciones.}\end{minipage}
\end{multicols*}
\newpage\fancyfoot[L]{\small Sección 1.1 | Ejercicio 11 | Paso a paso}\hypertarget{1.1-11}{}
\pdfbookmark[1]{1.1 / 11}{bm1.1-11}
\titulo{11 | Sistema lineal por eliminación}
\textbf{Enunciado}\par Resuelva el sistema lineal dado por medio del método de eliminación.
\[\left\{\begin{aligned}2x + 3y &= 13\\x - 2y &= 3\\5x + 2y &= 27\end{aligned}\right.\]
\begin{multicols*}{2}
\noindent\begin{minipage}{\linewidth}\textbf{1. Primero: ¿qué significa resolver este sistema?}\\[2pt]Las letras x, y representan números que todavía no conocemos. Una solución es una elección de esos números que hace verdaderas TODAS las ecuaciones al mismo tiempo. La llave de la izquierda significa que las condiciones se deben cumplir juntas.
\par\smallskip\textbf{Qué buscamos.} Encontrar una misma pareja de valores.
\par\smallskip El número 3 delante de x indica tres veces ese mismo valor; es decir, sumar tres copias:
\[3x=x+x+x\]
\par\smallskip Ejemplo numérico de la notación: si x valiera 2, tres veces ese valor sería 6. Este 2 es solo un ejemplo; todavía no es la respuesta del ejercicio. El punto centrado indica multiplicación:
\[3\cdot 2=2+2+2=6\]
\par\smallskip\textbf{Comprueba que lo entiendes.} ¿Basta con que los valores funcionen en una sola ecuación?\par\textit{Respuesta:} No. Deben funcionar en todas las ecuaciones del sistema, usando los mismos valores.
\end{minipage}\par\medskip
\noindent\begin{minipage}{\linewidth}\textbf{2. Nuestro plan: ir quitando letras}\\[2pt]Una ecuación con una sola letra es más fácil de resolver. Vamos a combinar ecuaciones para que una letra desaparezca. Cuando sepamos un valor, lo pondremos en las ecuaciones que faltan.
\par\smallskip\textbf{Qué buscamos.} Primero quitar x para calcular y. Después calcular x.
\par\smallskip\textbf{Por qué.} Si una cantidad se resta de sí misma, queda cero. Usaremos esta idea con términos que tengan la misma letra.
\par\smallskip Ejemplo de la idea que utilizaremos:
\[3x-3x=0\]
\par\smallskip\textbf{Qué aprendimos.} Eliminar una letra de una ecuación NO significa que su valor sea cero; significa que sus términos se cancelaron en esa cuenta.
\end{minipage}\par\medskip
\noindent\begin{minipage}{\linewidth}\textbf{3. Ponemos una ecuación más cómoda arriba}\\[2pt]Cambiamos de lugar las ecuaciones 1 y 2. Las condiciones siguen siendo las mismas; solamente cambia el orden en que las leemos.
\par\smallskip\textbf{Por qué.} La ecuación que subimos tiene un número más pequeño delante de x. Eso facilita las multiplicaciones que haremos a continuación.
\[\left\{\begin{aligned}x - 2y &= 3\\2x + 3y &= 13\\5x + 2y &= 27\end{aligned}\right.\]
\end{minipage}\par\medskip
\noindent\begin{minipage}{\linewidth}\textbf{4. ¿Cómo haremos desaparecer x?}\\[2pt]Trabajaremos con las ecuaciones 1 y 2. Queremos que los términos con x tengan el mismo tamaño. Así podremos restarlos.
\par\smallskip\textbf{Qué buscamos.} Conseguir términos con 2 veces x en las dos ecuaciones.
\par\smallskip\textbf{Por qué.} Tener 2 veces la misma letra en ambos lados de la resta nos permitirá cancelarla.
\par\smallskip {\small\renewcommand{\arraystretch}{1.3}\begin{tabular}{>{\raggedright\arraybackslash}p{2.58cm}>{\raggedright\arraybackslash}p{2.47cm}>{\raggedright\arraybackslash}p{6.17cm}}\hline\rowcolor{black!7} \textbf{Dónde miramos} & \textbf{Qué término hay} & \textbf{Qué necesitamos hacer}\\\hline Ecuación 1 & $x$ & Multiplicar TODA esta ecuación por 2\\ Ecuación 2 & $2x$ & La dejamos igual.\\ \hline\end{tabular}}\par\smallskip
\par\smallskip\textbf{Qué aprendimos.} Estos multiplicadores se eligen para igualar los términos con x. Los demás términos también se multiplican para conservar la igualdad.
\end{minipage}\par\medskip
\noindent\begin{minipage}{\linewidth}\textbf{5. Multiplicamos la ecuación 1 por 2}\\[2pt]Multiplicar una ecuación por 2 significa multiplicar CADA término de la izquierda y también el número de la derecha. La igualdad se conserva porque hacemos lo mismo en ambos lados.
\par\smallskip Esta es la ecuación antes de multiplicar:
\[x - 2y = 3\]
\par\smallskip {\small\renewcommand{\arraystretch}{1.3}\begin{tabular}{>{\raggedright\arraybackslash}p{2.80cm}>{\raggedright\arraybackslash}p{1.91cm}>{\raggedright\arraybackslash}p{6.50cm}}\hline\rowcolor{black!7} \textbf{Parte de la ecuación} & \textbf{Antes} & \textbf{La cuenta y su resultado}\\\hline Término con x & $x$ & $2\cdot\left(x\right)=2x$\\ Término con y & $-2y$ & $2\cdot\left(-2y\right)=-4y$\\ Número a la derecha del igual & $3$ & $2\cdot\left(3\right)=6$\\ \hline\end{tabular}}\par\smallskip
\par\smallskip Después de hacer todas las multiplicaciones queda:
\[2x - 4y = 6\]
\par\smallskip\textbf{Qué aprendimos.} La misma condición ahora está escrita con los términos multiplicados. Ya tenemos el tamaño que buscábamos para x.
\par\smallskip\textbf{Comprueba que lo entiendes.} ¿Podríamos multiplicar solamente el término con x?\par\textit{Respuesta:} No. Eso cambiaría la condición. Debemos multiplicar todos los términos, también el número que está después del igual.
\end{minipage}\par\medskip
\noindent\begin{minipage}{\linewidth}\textbf{6. Restamos las dos ecuaciones preparadas}\\[2pt]Restamos la ecuación 1 preparada a la ecuación 2 preparada. Juntamos cada término solo con otro que tenga la misma letra. Los números sin letra se calculan por separado.
\par\smallskip {\small\renewcommand{\arraystretch}{1.3}\begin{tabular}{>{\raggedright\arraybackslash}p{2.80cm}>{\raggedright\arraybackslash}p{5.38cm}>{\raggedright\arraybackslash}p{3.03cm}}\hline\rowcolor{black!7} \textbf{Qué agrupamos} & \textbf{Cuenta} & \textbf{Qué significa}\\\hline Los términos con x & $2x - \left(2x\right)=0$ & Se cancela esta letra.\\ Los términos con y & $3y - \left(-4y\right)=7y$ & Quedan 7y.\\ Los números a la derecha del igual & $13 - \left(6\right)=7$ & Este será el número de la nueva igualdad.\\ \hline\end{tabular}}\par\smallskip
\par\smallskip Primera igualdad que usaremos:
\[2x + 3y = 13\]
\par\smallskip Le restamos esta igualdad:
\[2x - 4y = 6\]
\par\smallskip\textbf{Qué aprendimos.} Los términos con x dan cero. Los términos con letras distintas no se mezclan.
\end{minipage}\par\medskip
\noindent\begin{minipage}{\linewidth}\textbf{7. Mira lo que hemos conseguido}\\[2pt]La nueva ecuación es el resultado de la cuenta anterior. Solo contiene y: después podremos despejarla.
\par\smallskip Escribimos la nueva condición:
\[7y = 7\]
\[\left\{\begin{aligned}x - 2y &= 3\\7y &= 7\\5x + 2y &= 27\end{aligned}\right.\]
\par\smallskip\textbf{Qué aprendimos.} Conservamos la ecuación 1 y reemplazamos la 2 por esta nueva condición. De esa manera conservamos toda la información necesaria para resolver el sistema.
\end{minipage}\par\medskip
\noindent\begin{minipage}{\linewidth}\textbf{8. ¿Cómo haremos desaparecer x?}\\[2pt]Trabajaremos con las ecuaciones 1 y 3. Queremos que los términos con x tengan el mismo tamaño. Así podremos restarlos.
\par\smallskip\textbf{Qué buscamos.} Conseguir términos con 5 veces x en las dos ecuaciones.
\par\smallskip\textbf{Por qué.} Tener 5 veces la misma letra en ambos lados de la resta nos permitirá cancelarla.
\par\smallskip {\small\renewcommand{\arraystretch}{1.3}\begin{tabular}{>{\raggedright\arraybackslash}p{2.58cm}>{\raggedright\arraybackslash}p{2.47cm}>{\raggedright\arraybackslash}p{6.17cm}}\hline\rowcolor{black!7} \textbf{Dónde miramos} & \textbf{Qué término hay} & \textbf{Qué necesitamos hacer}\\\hline Ecuación 1 & $x$ & Multiplicar TODA esta ecuación por 5\\ Ecuación 3 & $5x$ & La dejamos igual.\\ \hline\end{tabular}}\par\smallskip
\par\smallskip\textbf{Qué aprendimos.} Estos multiplicadores se eligen para igualar los términos con x. Los demás términos también se multiplican para conservar la igualdad.
\end{minipage}\par\medskip
\noindent\begin{minipage}{\linewidth}\textbf{9. Multiplicamos la ecuación 1 por 5}\\[2pt]Multiplicar una ecuación por 5 significa multiplicar CADA término de la izquierda y también el número de la derecha. La igualdad se conserva porque hacemos lo mismo en ambos lados.
\par\smallskip Esta es la ecuación antes de multiplicar:
\[x - 2y = 3\]
\par\smallskip {\small\renewcommand{\arraystretch}{1.3}\begin{tabular}{>{\raggedright\arraybackslash}p{2.80cm}>{\raggedright\arraybackslash}p{1.91cm}>{\raggedright\arraybackslash}p{6.50cm}}\hline\rowcolor{black!7} \textbf{Parte de la ecuación} & \textbf{Antes} & \textbf{La cuenta y su resultado}\\\hline Término con x & $x$ & $5\cdot\left(x\right)=5x$\\ Término con y & $-2y$ & $5\cdot\left(-2y\right)=-10y$\\ Número a la derecha del igual & $3$ & $5\cdot\left(3\right)=15$\\ \hline\end{tabular}}\par\smallskip
\par\smallskip Después de hacer todas las multiplicaciones queda:
\[5x - 10y = 15\]
\par\smallskip\textbf{Qué aprendimos.} La misma condición ahora está escrita con los términos multiplicados. Ya tenemos el tamaño que buscábamos para x.
\par\smallskip\textbf{Comprueba que lo entiendes.} ¿Podríamos multiplicar solamente el término con x?\par\textit{Respuesta:} No. Eso cambiaría la condición. Debemos multiplicar todos los términos, también el número que está después del igual.
\end{minipage}\par\medskip
\noindent\begin{minipage}{\linewidth}\textbf{10. Restamos las dos ecuaciones preparadas}\\[2pt]Restamos la ecuación 1 preparada a la ecuación 3 preparada. Juntamos cada término solo con otro que tenga la misma letra. Los números sin letra se calculan por separado.
\par\smallskip {\small\renewcommand{\arraystretch}{1.3}\begin{tabular}{>{\raggedright\arraybackslash}p{2.80cm}>{\raggedright\arraybackslash}p{5.38cm}>{\raggedright\arraybackslash}p{3.03cm}}\hline\rowcolor{black!7} \textbf{Qué agrupamos} & \textbf{Cuenta} & \textbf{Qué significa}\\\hline Los términos con x & $5x - \left(5x\right)=0$ & Se cancela esta letra.\\ Los términos con y & $2y - \left(-10y\right)=12y$ & Quedan 12y.\\ Los números a la derecha del igual & $27 - \left(15\right)=12$ & Este será el número de la nueva igualdad.\\ \hline\end{tabular}}\par\smallskip
\par\smallskip Primera igualdad que usaremos:
\[5x + 2y = 27\]
\par\smallskip Le restamos esta igualdad:
\[5x - 10y = 15\]
\par\smallskip\textbf{Qué aprendimos.} Los términos con x dan cero. Los términos con letras distintas no se mezclan.
\end{minipage}\par\medskip
\noindent\begin{minipage}{\linewidth}\textbf{11. Mira lo que hemos conseguido}\\[2pt]La nueva ecuación es el resultado de la cuenta anterior. Solo contiene y: después podremos despejarla.
\par\smallskip Escribimos la nueva condición:
\[12y = 12\]
\[\left\{\begin{aligned}x - 2y &= 3\\7y &= 7\\12y &= 12\end{aligned}\right.\]
\par\smallskip\textbf{Qué aprendimos.} Conservamos la ecuación 1 y reemplazamos la 3 por esta nueva condición. De esa manera conservamos toda la información necesaria para resolver el sistema.
\end{minipage}\par\medskip
\noindent\begin{minipage}{\linewidth}\textbf{12. ¿Cómo haremos desaparecer y?}\\[2pt]Trabajaremos con las ecuaciones 2 y 3. Queremos que los términos con y tengan el mismo tamaño. Así podremos restarlos.
\par\smallskip\textbf{Qué buscamos.} Conseguir términos con 84 veces y en las dos ecuaciones.
\par\smallskip\textbf{Por qué.} Tener 84 veces la misma letra en ambos lados de la resta nos permitirá cancelarla.
\par\smallskip {\small\renewcommand{\arraystretch}{1.3}\begin{tabular}{>{\raggedright\arraybackslash}p{2.58cm}>{\raggedright\arraybackslash}p{2.47cm}>{\raggedright\arraybackslash}p{6.17cm}}\hline\rowcolor{black!7} \textbf{Dónde miramos} & \textbf{Qué término hay} & \textbf{Qué necesitamos hacer}\\\hline Ecuación 2 & $7y$ & Multiplicar TODA esta ecuación por 12\\ Ecuación 3 & $12y$ & Multiplicar TODA esta ecuación por 7\\ \hline\end{tabular}}\par\smallskip
\par\smallskip\textbf{Qué aprendimos.} Estos multiplicadores se eligen para igualar los términos con y. Los demás términos también se multiplican para conservar la igualdad.
\end{minipage}\par\medskip
\noindent\begin{minipage}{\linewidth}\textbf{13. Multiplicamos la ecuación 2 por 12}\\[2pt]Multiplicar una ecuación por 12 significa multiplicar CADA término de la izquierda y también el número de la derecha. La igualdad se conserva porque hacemos lo mismo en ambos lados.
\par\smallskip Esta es la ecuación antes de multiplicar:
\[7y = 7\]
\par\smallskip {\small\renewcommand{\arraystretch}{1.3}\begin{tabular}{>{\raggedright\arraybackslash}p{2.80cm}>{\raggedright\arraybackslash}p{1.91cm}>{\raggedright\arraybackslash}p{6.50cm}}\hline\rowcolor{black!7} \textbf{Parte de la ecuación} & \textbf{Antes} & \textbf{La cuenta y su resultado}\\\hline Término con y & $7y$ & $12\cdot\left(7y\right)=84y$\\ Número a la derecha del igual & $7$ & $12\cdot\left(7\right)=84$\\ \hline\end{tabular}}\par\smallskip
\par\smallskip Después de hacer todas las multiplicaciones queda:
\[84y = 84\]
\par\smallskip\textbf{Qué aprendimos.} La misma condición ahora está escrita con los términos multiplicados. Ya tenemos el tamaño que buscábamos para y.
\par\smallskip\textbf{Comprueba que lo entiendes.} ¿Podríamos multiplicar solamente el término con y?\par\textit{Respuesta:} No. Eso cambiaría la condición. Debemos multiplicar todos los términos, también el número que está después del igual.
\end{minipage}\par\medskip
\noindent\begin{minipage}{\linewidth}\textbf{14. Multiplicamos la ecuación 3 por 7}\\[2pt]Multiplicar una ecuación por 7 significa multiplicar CADA término de la izquierda y también el número de la derecha. La igualdad se conserva porque hacemos lo mismo en ambos lados.
\par\smallskip Esta es la ecuación antes de multiplicar:
\[12y = 12\]
\par\smallskip {\small\renewcommand{\arraystretch}{1.3}\begin{tabular}{>{\raggedright\arraybackslash}p{2.80cm}>{\raggedright\arraybackslash}p{1.91cm}>{\raggedright\arraybackslash}p{6.50cm}}\hline\rowcolor{black!7} \textbf{Parte de la ecuación} & \textbf{Antes} & \textbf{La cuenta y su resultado}\\\hline Término con y & $12y$ & $7\cdot\left(12y\right)=84y$\\ Número a la derecha del igual & $12$ & $7\cdot\left(12\right)=84$\\ \hline\end{tabular}}\par\smallskip
\par\smallskip Después de hacer todas las multiplicaciones queda:
\[84y = 84\]
\par\smallskip\textbf{Qué aprendimos.} La misma condición ahora está escrita con los términos multiplicados. Ya tenemos el tamaño que buscábamos para y.
\par\smallskip\textbf{Comprueba que lo entiendes.} ¿Podríamos multiplicar solamente el término con y?\par\textit{Respuesta:} No. Eso cambiaría la condición. Debemos multiplicar todos los términos, también el número que está después del igual.
\end{minipage}\par\medskip
\noindent\begin{minipage}{\linewidth}\textbf{15. Restamos las dos ecuaciones preparadas}\\[2pt]Restamos la ecuación 2 preparada a la ecuación 3 preparada. Juntamos cada término solo con otro que tenga la misma letra. Los números sin letra se calculan por separado.
\par\smallskip {\small\renewcommand{\arraystretch}{1.3}\begin{tabular}{>{\raggedright\arraybackslash}p{2.80cm}>{\raggedright\arraybackslash}p{5.38cm}>{\raggedright\arraybackslash}p{3.03cm}}\hline\rowcolor{black!7} \textbf{Qué agrupamos} & \textbf{Cuenta} & \textbf{Qué significa}\\\hline Los términos con y & $84y - \left(84y\right)=0$ & Se cancela esta letra.\\ Los números a la derecha del igual & $84 - \left(84\right)=0$ & Este será el número de la nueva igualdad.\\ \hline\end{tabular}}\par\smallskip
\par\smallskip Primera igualdad que usaremos:
\[84y = 84\]
\par\smallskip Le restamos esta igualdad:
\[84y = 84\]
\par\smallskip\textbf{Qué aprendimos.} Los términos con y dan cero. Los términos con letras distintas no se mezclan.
\end{minipage}\par\medskip
\noindent\begin{minipage}{\linewidth}\textbf{16. Mira lo que hemos conseguido}\\[2pt]La nueva ecuación es el resultado de la cuenta anterior. Ya no contiene letras. Tenemos que decidir si la igualdad numérica es verdadera o imposible.
\par\smallskip Escribimos la nueva condición:
\[0 = 0\]
\[\left\{\begin{aligned}x - 2y &= 3\\7y &= 7\\0 &= 0\end{aligned}\right.\]
\par\smallskip\textbf{Qué aprendimos.} Conservamos la ecuación 2 y reemplazamos la 3 por esta nueva condición. De esa manera conservamos toda la información necesaria para resolver el sistema.
\end{minipage}\par\medskip
\noindent\begin{minipage}{\linewidth}\textbf{17. Ya tenemos una ecuación con una sola letra}\\[2pt]En esta condición solamente falta conocer y. El número 7 está multiplicando a esa letra. Ahora podemos deshacer esa multiplicación.
\par\smallskip Esta es la condición que vamos a resolver:
\[7y = 7\]
\end{minipage}\par\medskip
\noindent\begin{minipage}{\linewidth}\textbf{18. Calculamos cuánto vale y}\\[2pt]Para deshacer una multiplicación por 7, dividimos AMBOS lados entre 7. En el lado izquierdo queda una sola y. En el derecho hacemos la división.
\par\smallskip Dividimos los dos lados entre el mismo número:
\[\frac{7y}{7}=\frac{7}{7}\]
\par\smallskip Ya podemos leer el valor de la letra:
\[y=\frac{7}{7}=1\]
\par\smallskip\textbf{Qué aprendimos.} Este valor se usará en las ecuaciones donde todavía falten letras.
\end{minipage}\par\medskip
\noindent\begin{minipage}{\linewidth}\textbf{19. Usamos lo que ya sabemos para hallar x}\\[2pt]Ya sabemos y = 1. Sustituir significa escribir ese valor cada vez que aparezca su letra. Usamos paréntesis para que se vea qué número estamos multiplicando.
\par\smallskip\textbf{Qué buscamos.} Dejar solamente x como letra desconocida en esta ecuación.
\par\smallskip Tomamos esta condición que obtuvimos antes:
\[x - 2y = 3\]
\par\smallskip Reemplazamos las letras conocidas por sus valores:
\[x - 2\left(1\right)=3\]
\end{minipage}\par\medskip
\noindent\begin{minipage}{\linewidth}\textbf{20. Calculamos las multiplicaciones que ya se pueden hacer}\\[2pt]La letra x sigue sin conocerse. Las otras partes ya tienen un valor o una expresión conocida y podemos efectuar sus productos.
\par\smallskip {\small\renewcommand{\arraystretch}{1.3}\begin{tabular}{>{\raggedright\arraybackslash}p{3.36cm}>{\raggedright\arraybackslash}p{4.48cm}>{\raggedright\arraybackslash}p{3.36cm}}\hline\rowcolor{black!7} \textbf{Parte conocida} & \textbf{La multiplicación} & \textbf{Resultado}\\\hline El término -2y & $-2\cdot\left(1\right)$ & $-2$\\ \hline\end{tabular}}\par\smallskip
\par\smallskip Al reunir esas partes conocidas, la ecuación queda:
\[x - 2=3\]
\end{minipage}\par\medskip
\noindent\begin{minipage}{\linewidth}\textbf{21. Dejamos el término con x solo a un lado}\\[2pt]Para quitar la parte conocida, sumamos 2 en AMBOS lados. Piensa en una balanza: si hacemos lo mismo a ambos lados, la igualdad se mantiene.
\par\smallskip Hacemos la misma operación a la izquierda y a la derecha:
\[x - 2 + 2=3 + 2\]
\par\smallskip Las partes que se compensan dan cero; queda:
\[x=5\]
\par\smallskip\textbf{Qué aprendimos.} «Pasar al otro lado» es una forma abreviada de decir que hicimos la misma suma o resta en los dos lados.
\end{minipage}\par\medskip
\noindent\begin{minipage}{\linewidth}\textbf{22. Leemos el valor de x}\\[2pt]La letra ya está sola: su número multiplicador es 1, que no se escribe. La igualdad nos dice directamente cuánto vale.
\par\smallskip Anotamos el valor encontrado:
\[x=5\]
\end{minipage}\par\medskip
\noindent\begin{minipage}{\linewidth}\textbf{23. Reunimos la respuesta; todavía falta comprobarla}\\[2pt]Escribimos todos los valores en el orden x, y. Estos valores deben funcionar juntos, no cada uno en una ecuación diferente.
\[\left\{\begin{aligned}x&=5\\y&=1\end{aligned}\right.\]
\par\smallskip\textbf{Qué aprendimos.} El siguiente paso es volver a las ecuaciones del enunciado y comprobar cada condición.
\end{minipage}\par\medskip
\noindent\begin{minipage}{\linewidth}\textbf{24. ¿La respuesta cumple la ecuación 1?}\\[2pt]Comprobar significa volver a la condición original y calcularla con nuestra respuesta. Si los dos lados dan el mismo número, esta condición sí se cumple. Deben cumplirse todas.
\par\smallskip La condición del enunciado era:
\[2x + 3y = 13\]
\par\smallskip Ponemos nuestros valores donde estaban las letras:
\[2\left(5\right) + 3\left(1\right)=13\]
\par\smallskip Hacemos las multiplicaciones y luego la suma:
\[10 + 3=13\]
\par\smallskip Comparamos los dos lados:
\[13=13\]
\par\smallskip\textbf{Qué aprendimos.} La igualdad final es verdadera. Nuestra respuesta sí cumple la ecuación 1.
\end{minipage}\par\medskip
\noindent\begin{minipage}{\linewidth}\textbf{25. ¿La respuesta cumple la ecuación 2?}\\[2pt]Comprobar significa volver a la condición original y calcularla con nuestra respuesta. Si los dos lados dan el mismo número, esta condición sí se cumple. Deben cumplirse todas.
\par\smallskip La condición del enunciado era:
\[x - 2y = 3\]
\par\smallskip Ponemos nuestros valores donde estaban las letras:
\[\left(5\right) - 2\left(1\right)=3\]
\par\smallskip Hacemos las multiplicaciones y luego la suma:
\[5 - 2=3\]
\par\smallskip Comparamos los dos lados:
\[3=3\]
\par\smallskip\textbf{Qué aprendimos.} La igualdad final es verdadera. Nuestra respuesta sí cumple la ecuación 2.
\end{minipage}\par\medskip
\noindent\begin{minipage}{\linewidth}\textbf{26. ¿La respuesta cumple la ecuación 3?}\\[2pt]Comprobar significa volver a la condición original y calcularla con nuestra respuesta. Si los dos lados dan el mismo número, esta condición sí se cumple. Deben cumplirse todas.
\par\smallskip La condición del enunciado era:
\[5x + 2y = 27\]
\par\smallskip Ponemos nuestros valores donde estaban las letras:
\[5\left(5\right) + 2\left(1\right)=27\]
\par\smallskip Hacemos las multiplicaciones y luego la suma:
\[25 + 2=27\]
\par\smallskip Comparamos los dos lados:
\[27=27\]
\par\smallskip\textbf{Qué aprendimos.} La igualdad final es verdadera. Nuestra respuesta sí cumple la ecuación 3.
\end{minipage}\par\medskip
\respuesta{x = 5; y = 1.}
\noindent\begin{minipage}{\linewidth}\centering\includegraphics[width=.96\linewidth,height=6.8cm,keepaspectratio]{../recursos/grafico_1_1-11.png}\\{\small Lectura visual complementaria; la demostración está en las operaciones.}\end{minipage}
\end{multicols*}
\newpage\fancyfoot[L]{\small Sección 1.1 | Ejercicio 12 | Paso a paso}\hypertarget{1.1-12}{}
\pdfbookmark[1]{1.1 / 12}{bm1.1-12}
\titulo{12 | Sistema lineal por eliminación}
\textbf{Enunciado}\par Resuelva el sistema lineal dado por medio del método de eliminación.
\[\left\{\begin{aligned}x - 5y &= 6\\3x + 2y &= 1\\5x + 2y &= 1\end{aligned}\right.\]
\begin{multicols*}{2}
\noindent\begin{minipage}{\linewidth}\textbf{1. Primero: ¿qué significa resolver este sistema?}\\[2pt]Las letras x, y representan números que todavía no conocemos. Una solución es una elección de esos números que hace verdaderas TODAS las ecuaciones al mismo tiempo. La llave de la izquierda significa que las condiciones se deben cumplir juntas.
\par\smallskip\textbf{Qué buscamos.} Encontrar una misma pareja de valores.
\par\smallskip El número 3 delante de x indica tres veces ese mismo valor; es decir, sumar tres copias:
\[3x=x+x+x\]
\par\smallskip Ejemplo numérico de la notación: si x valiera 2, tres veces ese valor sería 6. Este 2 es solo un ejemplo; todavía no es la respuesta del ejercicio. El punto centrado indica multiplicación:
\[3\cdot 2=2+2+2=6\]
\par\smallskip\textbf{Comprueba que lo entiendes.} ¿Basta con que los valores funcionen en una sola ecuación?\par\textit{Respuesta:} No. Deben funcionar en todas las ecuaciones del sistema, usando los mismos valores.
\end{minipage}\par\medskip
\noindent\begin{minipage}{\linewidth}\textbf{2. Nuestro plan: ir quitando letras}\\[2pt]Una ecuación con una sola letra es más fácil de resolver. Vamos a combinar ecuaciones para que una letra desaparezca. Cuando sepamos un valor, lo pondremos en las ecuaciones que faltan.
\par\smallskip\textbf{Qué buscamos.} Primero quitar x para calcular y. Después calcular x.
\par\smallskip\textbf{Por qué.} Si una cantidad se resta de sí misma, queda cero. Usaremos esta idea con términos que tengan la misma letra.
\par\smallskip Ejemplo de la idea que utilizaremos:
\[3x-3x=0\]
\par\smallskip\textbf{Qué aprendimos.} Eliminar una letra de una ecuación NO significa que su valor sea cero; significa que sus términos se cancelaron en esa cuenta.
\end{minipage}\par\medskip
\noindent\begin{minipage}{\linewidth}\textbf{3. ¿Cómo haremos desaparecer x?}\\[2pt]Trabajaremos con las ecuaciones 1 y 2. Queremos que los términos con x tengan el mismo tamaño. Así podremos restarlos.
\par\smallskip\textbf{Qué buscamos.} Conseguir términos con 3 veces x en las dos ecuaciones.
\par\smallskip\textbf{Por qué.} Tener 3 veces la misma letra en ambos lados de la resta nos permitirá cancelarla.
\par\smallskip {\small\renewcommand{\arraystretch}{1.3}\begin{tabular}{>{\raggedright\arraybackslash}p{2.58cm}>{\raggedright\arraybackslash}p{2.47cm}>{\raggedright\arraybackslash}p{6.17cm}}\hline\rowcolor{black!7} \textbf{Dónde miramos} & \textbf{Qué término hay} & \textbf{Qué necesitamos hacer}\\\hline Ecuación 1 & $x$ & Multiplicar TODA esta ecuación por 3\\ Ecuación 2 & $3x$ & La dejamos igual.\\ \hline\end{tabular}}\par\smallskip
\par\smallskip\textbf{Qué aprendimos.} Estos multiplicadores se eligen para igualar los términos con x. Los demás términos también se multiplican para conservar la igualdad.
\end{minipage}\par\medskip
\noindent\begin{minipage}{\linewidth}\textbf{4. Multiplicamos la ecuación 1 por 3}\\[2pt]Multiplicar una ecuación por 3 significa multiplicar CADA término de la izquierda y también el número de la derecha. La igualdad se conserva porque hacemos lo mismo en ambos lados.
\par\smallskip Esta es la ecuación antes de multiplicar:
\[x - 5y = 6\]
\par\smallskip {\small\renewcommand{\arraystretch}{1.3}\begin{tabular}{>{\raggedright\arraybackslash}p{2.80cm}>{\raggedright\arraybackslash}p{1.91cm}>{\raggedright\arraybackslash}p{6.50cm}}\hline\rowcolor{black!7} \textbf{Parte de la ecuación} & \textbf{Antes} & \textbf{La cuenta y su resultado}\\\hline Término con x & $x$ & $3\cdot\left(x\right)=3x$\\ Término con y & $-5y$ & $3\cdot\left(-5y\right)=-15y$\\ Número a la derecha del igual & $6$ & $3\cdot\left(6\right)=18$\\ \hline\end{tabular}}\par\smallskip
\par\smallskip Después de hacer todas las multiplicaciones queda:
\[3x - 15y = 18\]
\par\smallskip\textbf{Qué aprendimos.} La misma condición ahora está escrita con los términos multiplicados. Ya tenemos el tamaño que buscábamos para x.
\par\smallskip\textbf{Comprueba que lo entiendes.} ¿Podríamos multiplicar solamente el término con x?\par\textit{Respuesta:} No. Eso cambiaría la condición. Debemos multiplicar todos los términos, también el número que está después del igual.
\end{minipage}\par\medskip
\noindent\begin{minipage}{\linewidth}\textbf{5. Restamos las dos ecuaciones preparadas}\\[2pt]Restamos la ecuación 1 preparada a la ecuación 2 preparada. Juntamos cada término solo con otro que tenga la misma letra. Los números sin letra se calculan por separado.
\par\smallskip {\small\renewcommand{\arraystretch}{1.3}\begin{tabular}{>{\raggedright\arraybackslash}p{2.80cm}>{\raggedright\arraybackslash}p{5.38cm}>{\raggedright\arraybackslash}p{3.03cm}}\hline\rowcolor{black!7} \textbf{Qué agrupamos} & \textbf{Cuenta} & \textbf{Qué significa}\\\hline Los términos con x & $3x - \left(3x\right)=0$ & Se cancela esta letra.\\ Los términos con y & $2y - \left(-15y\right)=17y$ & Quedan 17y.\\ Los números a la derecha del igual & $1 - \left(18\right)=-17$ & Este será el número de la nueva igualdad.\\ \hline\end{tabular}}\par\smallskip
\par\smallskip Primera igualdad que usaremos:
\[3x + 2y = 1\]
\par\smallskip Le restamos esta igualdad:
\[3x - 15y = 18\]
\par\smallskip\textbf{Qué aprendimos.} Los términos con x dan cero. Los términos con letras distintas no se mezclan.
\end{minipage}\par\medskip
\noindent\begin{minipage}{\linewidth}\textbf{6. Mira lo que hemos conseguido}\\[2pt]La nueva ecuación es el resultado de la cuenta anterior. Solo contiene y: después podremos despejarla.
\par\smallskip Escribimos la nueva condición:
\[17y = -17\]
\[\left\{\begin{aligned}x - 5y &= 6\\17y &= -17\\5x + 2y &= 1\end{aligned}\right.\]
\par\smallskip\textbf{Qué aprendimos.} Conservamos la ecuación 1 y reemplazamos la 2 por esta nueva condición. De esa manera conservamos toda la información necesaria para resolver el sistema.
\end{minipage}\par\medskip
\noindent\begin{minipage}{\linewidth}\textbf{7. ¿Cómo haremos desaparecer x?}\\[2pt]Trabajaremos con las ecuaciones 1 y 3. Queremos que los términos con x tengan el mismo tamaño. Así podremos restarlos.
\par\smallskip\textbf{Qué buscamos.} Conseguir términos con 5 veces x en las dos ecuaciones.
\par\smallskip\textbf{Por qué.} Tener 5 veces la misma letra en ambos lados de la resta nos permitirá cancelarla.
\par\smallskip {\small\renewcommand{\arraystretch}{1.3}\begin{tabular}{>{\raggedright\arraybackslash}p{2.58cm}>{\raggedright\arraybackslash}p{2.47cm}>{\raggedright\arraybackslash}p{6.17cm}}\hline\rowcolor{black!7} \textbf{Dónde miramos} & \textbf{Qué término hay} & \textbf{Qué necesitamos hacer}\\\hline Ecuación 1 & $x$ & Multiplicar TODA esta ecuación por 5\\ Ecuación 3 & $5x$ & La dejamos igual.\\ \hline\end{tabular}}\par\smallskip
\par\smallskip\textbf{Qué aprendimos.} Estos multiplicadores se eligen para igualar los términos con x. Los demás términos también se multiplican para conservar la igualdad.
\end{minipage}\par\medskip
\noindent\begin{minipage}{\linewidth}\textbf{8. Multiplicamos la ecuación 1 por 5}\\[2pt]Multiplicar una ecuación por 5 significa multiplicar CADA término de la izquierda y también el número de la derecha. La igualdad se conserva porque hacemos lo mismo en ambos lados.
\par\smallskip Esta es la ecuación antes de multiplicar:
\[x - 5y = 6\]
\par\smallskip {\small\renewcommand{\arraystretch}{1.3}\begin{tabular}{>{\raggedright\arraybackslash}p{2.80cm}>{\raggedright\arraybackslash}p{1.91cm}>{\raggedright\arraybackslash}p{6.50cm}}\hline\rowcolor{black!7} \textbf{Parte de la ecuación} & \textbf{Antes} & \textbf{La cuenta y su resultado}\\\hline Término con x & $x$ & $5\cdot\left(x\right)=5x$\\ Término con y & $-5y$ & $5\cdot\left(-5y\right)=-25y$\\ Número a la derecha del igual & $6$ & $5\cdot\left(6\right)=30$\\ \hline\end{tabular}}\par\smallskip
\par\smallskip Después de hacer todas las multiplicaciones queda:
\[5x - 25y = 30\]
\par\smallskip\textbf{Qué aprendimos.} La misma condición ahora está escrita con los términos multiplicados. Ya tenemos el tamaño que buscábamos para x.
\par\smallskip\textbf{Comprueba que lo entiendes.} ¿Podríamos multiplicar solamente el término con x?\par\textit{Respuesta:} No. Eso cambiaría la condición. Debemos multiplicar todos los términos, también el número que está después del igual.
\end{minipage}\par\medskip
\noindent\begin{minipage}{\linewidth}\textbf{9. Restamos las dos ecuaciones preparadas}\\[2pt]Restamos la ecuación 1 preparada a la ecuación 3 preparada. Juntamos cada término solo con otro que tenga la misma letra. Los números sin letra se calculan por separado.
\par\smallskip {\small\renewcommand{\arraystretch}{1.3}\begin{tabular}{>{\raggedright\arraybackslash}p{2.80cm}>{\raggedright\arraybackslash}p{5.38cm}>{\raggedright\arraybackslash}p{3.03cm}}\hline\rowcolor{black!7} \textbf{Qué agrupamos} & \textbf{Cuenta} & \textbf{Qué significa}\\\hline Los términos con x & $5x - \left(5x\right)=0$ & Se cancela esta letra.\\ Los términos con y & $2y - \left(-25y\right)=27y$ & Quedan 27y.\\ Los números a la derecha del igual & $1 - \left(30\right)=-29$ & Este será el número de la nueva igualdad.\\ \hline\end{tabular}}\par\smallskip
\par\smallskip Primera igualdad que usaremos:
\[5x + 2y = 1\]
\par\smallskip Le restamos esta igualdad:
\[5x - 25y = 30\]
\par\smallskip\textbf{Qué aprendimos.} Los términos con x dan cero. Los términos con letras distintas no se mezclan.
\end{minipage}\par\medskip
\noindent\begin{minipage}{\linewidth}\textbf{10. Mira lo que hemos conseguido}\\[2pt]La nueva ecuación es el resultado de la cuenta anterior. Solo contiene y: después podremos despejarla.
\par\smallskip Escribimos la nueva condición:
\[27y = -29\]
\[\left\{\begin{aligned}x - 5y &= 6\\17y &= -17\\27y &= -29\end{aligned}\right.\]
\par\smallskip\textbf{Qué aprendimos.} Conservamos la ecuación 1 y reemplazamos la 3 por esta nueva condición. De esa manera conservamos toda la información necesaria para resolver el sistema.
\end{minipage}\par\medskip
\noindent\begin{minipage}{\linewidth}\textbf{11. ¿Cómo haremos desaparecer y?}\\[2pt]Trabajaremos con las ecuaciones 2 y 3. Queremos que los términos con y tengan el mismo tamaño. Así podremos restarlos.
\par\smallskip\textbf{Qué buscamos.} Conseguir términos con 459 veces y en las dos ecuaciones.
\par\smallskip\textbf{Por qué.} Tener 459 veces la misma letra en ambos lados de la resta nos permitirá cancelarla.
\par\smallskip {\small\renewcommand{\arraystretch}{1.3}\begin{tabular}{>{\raggedright\arraybackslash}p{2.58cm}>{\raggedright\arraybackslash}p{2.47cm}>{\raggedright\arraybackslash}p{6.17cm}}\hline\rowcolor{black!7} \textbf{Dónde miramos} & \textbf{Qué término hay} & \textbf{Qué necesitamos hacer}\\\hline Ecuación 2 & $17y$ & Multiplicar TODA esta ecuación por 27\\ Ecuación 3 & $27y$ & Multiplicar TODA esta ecuación por 17\\ \hline\end{tabular}}\par\smallskip
\par\smallskip\textbf{Qué aprendimos.} Estos multiplicadores se eligen para igualar los términos con y. Los demás términos también se multiplican para conservar la igualdad.
\end{minipage}\par\medskip
\noindent\begin{minipage}{\linewidth}\textbf{12. Multiplicamos la ecuación 2 por 27}\\[2pt]Multiplicar una ecuación por 27 significa multiplicar CADA término de la izquierda y también el número de la derecha. La igualdad se conserva porque hacemos lo mismo en ambos lados.
\par\smallskip Esta es la ecuación antes de multiplicar:
\[17y = -17\]
\par\smallskip {\small\renewcommand{\arraystretch}{1.3}\begin{tabular}{>{\raggedright\arraybackslash}p{2.80cm}>{\raggedright\arraybackslash}p{1.91cm}>{\raggedright\arraybackslash}p{6.50cm}}\hline\rowcolor{black!7} \textbf{Parte de la ecuación} & \textbf{Antes} & \textbf{La cuenta y su resultado}\\\hline Término con y & $17y$ & $27\cdot\left(17y\right)=459y$\\ Número a la derecha del igual & $-17$ & $27\cdot\left(-17\right)=-459$\\ \hline\end{tabular}}\par\smallskip
\par\smallskip Después de hacer todas las multiplicaciones queda:
\[459y = -459\]
\par\smallskip\textbf{Qué aprendimos.} La misma condición ahora está escrita con los términos multiplicados. Ya tenemos el tamaño que buscábamos para y.
\par\smallskip\textbf{Comprueba que lo entiendes.} ¿Podríamos multiplicar solamente el término con y?\par\textit{Respuesta:} No. Eso cambiaría la condición. Debemos multiplicar todos los términos, también el número que está después del igual.
\end{minipage}\par\medskip
\noindent\begin{minipage}{\linewidth}\textbf{13. Multiplicamos la ecuación 3 por 17}\\[2pt]Multiplicar una ecuación por 17 significa multiplicar CADA término de la izquierda y también el número de la derecha. La igualdad se conserva porque hacemos lo mismo en ambos lados.
\par\smallskip Esta es la ecuación antes de multiplicar:
\[27y = -29\]
\par\smallskip {\small\renewcommand{\arraystretch}{1.3}\begin{tabular}{>{\raggedright\arraybackslash}p{2.80cm}>{\raggedright\arraybackslash}p{1.91cm}>{\raggedright\arraybackslash}p{6.50cm}}\hline\rowcolor{black!7} \textbf{Parte de la ecuación} & \textbf{Antes} & \textbf{La cuenta y su resultado}\\\hline Término con y & $27y$ & $17\cdot\left(27y\right)=459y$\\ Número a la derecha del igual & $-29$ & $17\cdot\left(-29\right)=-493$\\ \hline\end{tabular}}\par\smallskip
\par\smallskip Después de hacer todas las multiplicaciones queda:
\[459y = -493\]
\par\smallskip\textbf{Qué aprendimos.} La misma condición ahora está escrita con los términos multiplicados. Ya tenemos el tamaño que buscábamos para y.
\par\smallskip\textbf{Comprueba que lo entiendes.} ¿Podríamos multiplicar solamente el término con y?\par\textit{Respuesta:} No. Eso cambiaría la condición. Debemos multiplicar todos los términos, también el número que está después del igual.
\end{minipage}\par\medskip
\noindent\begin{minipage}{\linewidth}\textbf{14. Restamos las dos ecuaciones preparadas}\\[2pt]Restamos la ecuación 2 preparada a la ecuación 3 preparada. Juntamos cada término solo con otro que tenga la misma letra. Los números sin letra se calculan por separado.
\par\smallskip {\small\renewcommand{\arraystretch}{1.3}\begin{tabular}{>{\raggedright\arraybackslash}p{2.80cm}>{\raggedright\arraybackslash}p{5.38cm}>{\raggedright\arraybackslash}p{3.03cm}}\hline\rowcolor{black!7} \textbf{Qué agrupamos} & \textbf{Cuenta} & \textbf{Qué significa}\\\hline Los términos con y & $459y - \left(459y\right)=0$ & Se cancela esta letra.\\ Los números a la derecha del igual & $-493 - \left(-459\right)=-493+459=-34$ & Este será el número de la nueva igualdad.\\ \hline\end{tabular}}\par\smallskip
\par\smallskip Primera igualdad que usaremos:
\[459y = -493\]
\par\smallskip Le restamos esta igualdad:
\[459y = -459\]
\par\smallskip\textbf{Qué aprendimos.} Los términos con y dan cero. Recuerda: restar un número negativo equivale a sumar su opuesto.
\end{minipage}\par\medskip
\noindent\begin{minipage}{\linewidth}\textbf{15. Mira lo que hemos conseguido}\\[2pt]La nueva ecuación es el resultado de la cuenta anterior. Ya no contiene letras. Tenemos que decidir si la igualdad numérica es verdadera o imposible.
\par\smallskip Escribimos la nueva condición:
\[0 = -34\]
\[\left\{\begin{aligned}x - 5y &= 6\\17y &= -17\\0 &= -34\end{aligned}\right.\]
\par\smallskip\textbf{Qué aprendimos.} Conservamos la ecuación 2 y reemplazamos la 3 por esta nueva condición. De esa manera conservamos toda la información necesaria para resolver el sistema.
\end{minipage}\par\medskip
\noindent\begin{minipage}{\linewidth}\textbf{16. Esta igualdad nos dice que no hay solución}\\[2pt]Al hacer las cuentas, todas las letras desaparecieron y quedó la igualdad que ves abajo. Cero no puede ser igual a -34. Como esa condición tendría que cumplirse junto con las demás, ningún grupo de valores puede resolver el sistema.
\par\smallskip Esta es la contradicción encontrada:
\[0 = -34\]
\par\smallskip\textbf{Qué aprendimos.} El resultado es «sin solución». No significa que debamos responder con ceros: significa que no existe una respuesta que cumpla todas las condiciones.
\par\smallskip\textbf{Comprueba que lo entiendes.} ¿Responder x = 0, y = 0 resolvería esta contradicción?\par\textit{Respuesta:} No. La contradicción ya no depende de las letras. Ningún valor puede convertir cero en otro número.
\end{minipage}\par\medskip
\respuesta{Sin solución.}
\noindent\begin{minipage}{\linewidth}\centering\includegraphics[width=.96\linewidth,height=6.8cm,keepaspectratio]{../recursos/grafico_1_1-12.png}\\{\small Lectura visual complementaria; la demostración está en las operaciones.}\end{minipage}
\end{multicols*}
\newpage\fancyfoot[L]{\small Sección 1.1 | Ejercicio 13 | Paso a paso}\hypertarget{1.1-13}{}
\pdfbookmark[1]{1.1 / 13}{bm1.1-13}
\titulo{13 | Sistema lineal por eliminación}
\textbf{Enunciado}\par Resuelva el sistema lineal dado por medio del método de eliminación.
\[\left\{\begin{aligned}x + 3y &= -4\\2x + 5y &= -8\\x + 3y &= -5\end{aligned}\right.\]
\begin{multicols*}{2}
\noindent\begin{minipage}{\linewidth}\textbf{1. Primero: ¿qué significa resolver este sistema?}\\[2pt]Las letras x, y representan números que todavía no conocemos. Una solución es una elección de esos números que hace verdaderas TODAS las ecuaciones al mismo tiempo. La llave de la izquierda significa que las condiciones se deben cumplir juntas.
\par\smallskip\textbf{Qué buscamos.} Encontrar una misma pareja de valores.
\par\smallskip El número 3 delante de x indica tres veces ese mismo valor; es decir, sumar tres copias:
\[3x=x+x+x\]
\par\smallskip Ejemplo numérico de la notación: si x valiera 2, tres veces ese valor sería 6. Este 2 es solo un ejemplo; todavía no es la respuesta del ejercicio. El punto centrado indica multiplicación:
\[3\cdot 2=2+2+2=6\]
\par\smallskip\textbf{Comprueba que lo entiendes.} ¿Basta con que los valores funcionen en una sola ecuación?\par\textit{Respuesta:} No. Deben funcionar en todas las ecuaciones del sistema, usando los mismos valores.
\end{minipage}\par\medskip
\noindent\begin{minipage}{\linewidth}\textbf{2. Nuestro plan: ir quitando letras}\\[2pt]Una ecuación con una sola letra es más fácil de resolver. Vamos a combinar ecuaciones para que una letra desaparezca. Cuando sepamos un valor, lo pondremos en las ecuaciones que faltan.
\par\smallskip\textbf{Qué buscamos.} Primero quitar x para calcular y. Después calcular x.
\par\smallskip\textbf{Por qué.} Si una cantidad se resta de sí misma, queda cero. Usaremos esta idea con términos que tengan la misma letra.
\par\smallskip Ejemplo de la idea que utilizaremos:
\[3x-3x=0\]
\par\smallskip\textbf{Qué aprendimos.} Eliminar una letra de una ecuación NO significa que su valor sea cero; significa que sus términos se cancelaron en esa cuenta.
\end{minipage}\par\medskip
\noindent\begin{minipage}{\linewidth}\textbf{3. ¿Cómo haremos desaparecer x?}\\[2pt]Trabajaremos con las ecuaciones 1 y 2. Queremos que los términos con x tengan el mismo tamaño. Así podremos restarlos.
\par\smallskip\textbf{Qué buscamos.} Conseguir términos con 2 veces x en las dos ecuaciones.
\par\smallskip\textbf{Por qué.} Tener 2 veces la misma letra en ambos lados de la resta nos permitirá cancelarla.
\par\smallskip {\small\renewcommand{\arraystretch}{1.3}\begin{tabular}{>{\raggedright\arraybackslash}p{2.58cm}>{\raggedright\arraybackslash}p{2.47cm}>{\raggedright\arraybackslash}p{6.17cm}}\hline\rowcolor{black!7} \textbf{Dónde miramos} & \textbf{Qué término hay} & \textbf{Qué necesitamos hacer}\\\hline Ecuación 1 & $x$ & Multiplicar TODA esta ecuación por 2\\ Ecuación 2 & $2x$ & La dejamos igual.\\ \hline\end{tabular}}\par\smallskip
\par\smallskip\textbf{Qué aprendimos.} Estos multiplicadores se eligen para igualar los términos con x. Los demás términos también se multiplican para conservar la igualdad.
\end{minipage}\par\medskip
\noindent\begin{minipage}{\linewidth}\textbf{4. Multiplicamos la ecuación 1 por 2}\\[2pt]Multiplicar una ecuación por 2 significa multiplicar CADA término de la izquierda y también el número de la derecha. La igualdad se conserva porque hacemos lo mismo en ambos lados.
\par\smallskip Esta es la ecuación antes de multiplicar:
\[x + 3y = -4\]
\par\smallskip {\small\renewcommand{\arraystretch}{1.3}\begin{tabular}{>{\raggedright\arraybackslash}p{2.80cm}>{\raggedright\arraybackslash}p{1.91cm}>{\raggedright\arraybackslash}p{6.50cm}}\hline\rowcolor{black!7} \textbf{Parte de la ecuación} & \textbf{Antes} & \textbf{La cuenta y su resultado}\\\hline Término con x & $x$ & $2\cdot\left(x\right)=2x$\\ Término con y & $3y$ & $2\cdot\left(3y\right)=6y$\\ Número a la derecha del igual & $-4$ & $2\cdot\left(-4\right)=-8$\\ \hline\end{tabular}}\par\smallskip
\par\smallskip Después de hacer todas las multiplicaciones queda:
\[2x + 6y = -8\]
\par\smallskip\textbf{Qué aprendimos.} La misma condición ahora está escrita con los términos multiplicados. Ya tenemos el tamaño que buscábamos para x.
\par\smallskip\textbf{Comprueba que lo entiendes.} ¿Podríamos multiplicar solamente el término con x?\par\textit{Respuesta:} No. Eso cambiaría la condición. Debemos multiplicar todos los términos, también el número que está después del igual.
\end{minipage}\par\medskip
\noindent\begin{minipage}{\linewidth}\textbf{5. Restamos las dos ecuaciones preparadas}\\[2pt]Restamos la ecuación 1 preparada a la ecuación 2 preparada. Juntamos cada término solo con otro que tenga la misma letra. Los números sin letra se calculan por separado.
\par\smallskip {\small\renewcommand{\arraystretch}{1.3}\begin{tabular}{>{\raggedright\arraybackslash}p{2.80cm}>{\raggedright\arraybackslash}p{5.38cm}>{\raggedright\arraybackslash}p{3.03cm}}\hline\rowcolor{black!7} \textbf{Qué agrupamos} & \textbf{Cuenta} & \textbf{Qué significa}\\\hline Los términos con x & $2x - \left(2x\right)=0$ & Se cancela esta letra.\\ Los términos con y & $5y - \left(6y\right)=-y$ & Quedan -y.\\ Los números a la derecha del igual & $-8 - \left(-8\right)=-8+8=0$ & Este será el número de la nueva igualdad.\\ \hline\end{tabular}}\par\smallskip
\par\smallskip Primera igualdad que usaremos:
\[2x + 5y = -8\]
\par\smallskip Le restamos esta igualdad:
\[2x + 6y = -8\]
\par\smallskip\textbf{Qué aprendimos.} Los términos con x dan cero. Recuerda: restar un número negativo equivale a sumar su opuesto.
\end{minipage}\par\medskip
\noindent\begin{minipage}{\linewidth}\textbf{6. Mira lo que hemos conseguido}\\[2pt]La nueva ecuación es el resultado de la cuenta anterior. Solo contiene y: después podremos despejarla.
\par\smallskip Escribimos la nueva condición:
\[-y = 0\]
\[\left\{\begin{aligned}x + 3y &= -4\\-y &= 0\\x + 3y &= -5\end{aligned}\right.\]
\par\smallskip\textbf{Qué aprendimos.} Conservamos la ecuación 1 y reemplazamos la 2 por esta nueva condición. De esa manera conservamos toda la información necesaria para resolver el sistema.
\end{minipage}\par\medskip
\noindent\begin{minipage}{\linewidth}\textbf{7. ¿Cómo haremos desaparecer x?}\\[2pt]Trabajaremos con las ecuaciones 1 y 3. Queremos que los términos con x tengan el mismo tamaño. Así podremos restarlos.
\par\smallskip\textbf{Qué buscamos.} Conseguir términos con 1 veces x en las dos ecuaciones.
\par\smallskip\textbf{Por qué.} Tener 1 veces la misma letra en ambos lados de la resta nos permitirá cancelarla.
\par\smallskip {\small\renewcommand{\arraystretch}{1.3}\begin{tabular}{>{\raggedright\arraybackslash}p{2.58cm}>{\raggedright\arraybackslash}p{2.47cm}>{\raggedright\arraybackslash}p{6.17cm}}\hline\rowcolor{black!7} \textbf{Dónde miramos} & \textbf{Qué término hay} & \textbf{Qué necesitamos hacer}\\\hline Ecuación 1 & $x$ & La dejamos igual.\\ Ecuación 3 & $x$ & La dejamos igual.\\ \hline\end{tabular}}\par\smallskip
\par\smallskip\textbf{Qué aprendimos.} Estos multiplicadores se eligen para igualar los términos con x. Los demás términos también se multiplican para conservar la igualdad.
\end{minipage}\par\medskip
\noindent\begin{minipage}{\linewidth}\textbf{8. Restamos las dos ecuaciones preparadas}\\[2pt]Restamos la ecuación 1 preparada a la ecuación 3 preparada. Juntamos cada término solo con otro que tenga la misma letra. Los números sin letra se calculan por separado.
\par\smallskip {\small\renewcommand{\arraystretch}{1.3}\begin{tabular}{>{\raggedright\arraybackslash}p{2.80cm}>{\raggedright\arraybackslash}p{5.38cm}>{\raggedright\arraybackslash}p{3.03cm}}\hline\rowcolor{black!7} \textbf{Qué agrupamos} & \textbf{Cuenta} & \textbf{Qué significa}\\\hline Los términos con x & $x - \left(x\right)=0$ & Se cancela esta letra.\\ Los términos con y & $3y - \left(3y\right)=0$ & Se cancela esta letra.\\ Los números a la derecha del igual & $-5 - \left(-4\right)=-5+4=-1$ & Este será el número de la nueva igualdad.\\ \hline\end{tabular}}\par\smallskip
\par\smallskip Primera igualdad que usaremos:
\[x + 3y = -5\]
\par\smallskip Le restamos esta igualdad:
\[x + 3y = -4\]
\par\smallskip\textbf{Qué aprendimos.} Los términos con x dan cero. Recuerda: restar un número negativo equivale a sumar su opuesto.
\end{minipage}\par\medskip
\noindent\begin{minipage}{\linewidth}\textbf{9. Mira lo que hemos conseguido}\\[2pt]La nueva ecuación es el resultado de la cuenta anterior. Ya no contiene letras. Tenemos que decidir si la igualdad numérica es verdadera o imposible.
\par\smallskip Escribimos la nueva condición:
\[0 = -1\]
\[\left\{\begin{aligned}x + 3y &= -4\\-y &= 0\\0 &= -1\end{aligned}\right.\]
\par\smallskip\textbf{Qué aprendimos.} Conservamos la ecuación 1 y reemplazamos la 3 por esta nueva condición. De esa manera conservamos toda la información necesaria para resolver el sistema.
\end{minipage}\par\medskip
\noindent\begin{minipage}{\linewidth}\textbf{10. Esta igualdad nos dice que no hay solución}\\[2pt]Al hacer las cuentas, todas las letras desaparecieron y quedó la igualdad que ves abajo. Cero no puede ser igual a -1. Como esa condición tendría que cumplirse junto con las demás, ningún grupo de valores puede resolver el sistema.
\par\smallskip Esta es la contradicción encontrada:
\[0 = -1\]
\par\smallskip\textbf{Qué aprendimos.} El resultado es «sin solución». No significa que debamos responder con ceros: significa que no existe una respuesta que cumpla todas las condiciones.
\par\smallskip\textbf{Comprueba que lo entiendes.} ¿Responder x = 0, y = 0 resolvería esta contradicción?\par\textit{Respuesta:} No. La contradicción ya no depende de las letras. Ningún valor puede convertir cero en otro número.
\end{minipage}\par\medskip
\respuesta{Sin solución.}
\noindent\begin{minipage}{\linewidth}\centering\includegraphics[width=.96\linewidth,height=6.8cm,keepaspectratio]{../recursos/grafico_1_1-13.png}\\{\small Lectura visual complementaria; la demostración está en las operaciones.}\end{minipage}
\end{multicols*}
\newpage\fancyfoot[L]{\small Sección 1.1 | Ejercicio 14 | Paso a paso}\hypertarget{1.1-14}{}
\pdfbookmark[1]{1.1 / 14}{bm1.1-14}
\titulo{14 | Sistema lineal por eliminación}
\textbf{Enunciado}\par Resuelva el sistema lineal dado por medio del método de eliminación.
\[\left\{\begin{aligned}2x + 3y - z &= 6\\2x - y + 2z &= -8\\3x - y + z &= -7\end{aligned}\right.\]
\begin{multicols*}{2}
\noindent\begin{minipage}{\linewidth}\textbf{1. Primero: ¿qué significa resolver este sistema?}\\[2pt]Las letras x, y, z representan números que todavía no conocemos. Una solución es una elección de esos números que hace verdaderas TODAS las ecuaciones al mismo tiempo. La llave de la izquierda significa que las condiciones se deben cumplir juntas.
\par\smallskip\textbf{Qué buscamos.} Encontrar un mismo grupo de valores para las 3 letras.
\par\smallskip El número 3 delante de x indica tres veces ese mismo valor; es decir, sumar tres copias:
\[3x=x+x+x\]
\par\smallskip Ejemplo numérico de la notación: si x valiera 2, tres veces ese valor sería 6. Este 2 es solo un ejemplo; todavía no es la respuesta del ejercicio. El punto centrado indica multiplicación:
\[3\cdot 2=2+2+2=6\]
\par\smallskip\textbf{Comprueba que lo entiendes.} ¿Basta con que los valores funcionen en una sola ecuación?\par\textit{Respuesta:} No. Deben funcionar en todas las ecuaciones del sistema, usando los mismos valores.
\end{minipage}\par\medskip
\noindent\begin{minipage}{\linewidth}\textbf{2. Nuestro plan: ir quitando letras}\\[2pt]Una ecuación con una sola letra es más fácil de resolver. Vamos a combinar ecuaciones para que una letra desaparezca. Cuando sepamos un valor, lo pondremos en las ecuaciones que faltan.
\par\smallskip\textbf{Qué buscamos.} Primero quitar x de las otras ecuaciones. Después seguir con las letras restantes.
\par\smallskip\textbf{Por qué.} Si una cantidad se resta de sí misma, queda cero. Usaremos esta idea con términos que tengan la misma letra.
\par\smallskip Ejemplo de la idea que utilizaremos:
\[3x-3x=0\]
\par\smallskip\textbf{Qué aprendimos.} Eliminar una letra de una ecuación NO significa que su valor sea cero; significa que sus términos se cancelaron en esa cuenta.
\end{minipage}\par\medskip
\noindent\begin{minipage}{\linewidth}\textbf{3. ¿Cómo haremos desaparecer x?}\\[2pt]Trabajaremos con las ecuaciones 1 y 2. Queremos que los términos con x tengan el mismo tamaño. Así podremos restarlos.
\par\smallskip\textbf{Qué buscamos.} Conseguir términos con 2 veces x en las dos ecuaciones.
\par\smallskip\textbf{Por qué.} Tener 2 veces la misma letra en ambos lados de la resta nos permitirá cancelarla.
\par\smallskip {\small\renewcommand{\arraystretch}{1.3}\begin{tabular}{>{\raggedright\arraybackslash}p{2.58cm}>{\raggedright\arraybackslash}p{2.47cm}>{\raggedright\arraybackslash}p{6.17cm}}\hline\rowcolor{black!7} \textbf{Dónde miramos} & \textbf{Qué término hay} & \textbf{Qué necesitamos hacer}\\\hline Ecuación 1 & $2x$ & La dejamos igual.\\ Ecuación 2 & $2x$ & La dejamos igual.\\ \hline\end{tabular}}\par\smallskip
\par\smallskip\textbf{Qué aprendimos.} Estos multiplicadores se eligen para igualar los términos con x. Los demás términos también se multiplican para conservar la igualdad.
\end{minipage}\par\medskip
\noindent\begin{minipage}{\linewidth}\textbf{4. Restamos las dos ecuaciones preparadas}\\[2pt]Restamos la ecuación 1 preparada a la ecuación 2 preparada. Juntamos cada término solo con otro que tenga la misma letra. Los números sin letra se calculan por separado.
\par\smallskip {\small\renewcommand{\arraystretch}{1.3}\begin{tabular}{>{\raggedright\arraybackslash}p{2.80cm}>{\raggedright\arraybackslash}p{5.38cm}>{\raggedright\arraybackslash}p{3.03cm}}\hline\rowcolor{black!7} \textbf{Qué agrupamos} & \textbf{Cuenta} & \textbf{Qué significa}\\\hline Los términos con x & $2x - \left(2x\right)=0$ & Se cancela esta letra.\\ Los términos con y & $-y - \left(3y\right)=-4y$ & Quedan -4y.\\ Los términos con z & $2z - \left(-z\right)=3z$ & Quedan 3z.\\ Los números a la derecha del igual & $-8 - \left(6\right)=-14$ & Este será el número de la nueva igualdad.\\ \hline\end{tabular}}\par\smallskip
\par\smallskip Primera igualdad que usaremos:
\[2x - y + 2z = -8\]
\par\smallskip Le restamos esta igualdad:
\[2x + 3y - z = 6\]
\par\smallskip\textbf{Qué aprendimos.} Los términos con x dan cero. Los términos con letras distintas no se mezclan.
\end{minipage}\par\medskip
\noindent\begin{minipage}{\linewidth}\textbf{5. Mira lo que hemos conseguido}\\[2pt]La nueva ecuación es el resultado de la cuenta anterior. Todavía contiene y, z. Seguiremos reduciendo las condiciones.
\par\smallskip Escribimos la nueva condición:
\[-4y + 3z = -14\]
\[\left\{\begin{aligned}2x + 3y - z &= 6\\-4y + 3z &= -14\\3x - y + z &= -7\end{aligned}\right.\]
\par\smallskip\textbf{Qué aprendimos.} Conservamos la ecuación 1 y reemplazamos la 2 por esta nueva condición. De esa manera conservamos toda la información necesaria para resolver el sistema.
\end{minipage}\par\medskip
\noindent\begin{minipage}{\linewidth}\textbf{6. ¿Cómo haremos desaparecer x?}\\[2pt]Trabajaremos con las ecuaciones 1 y 3. Queremos que los términos con x tengan el mismo tamaño. Así podremos restarlos.
\par\smallskip\textbf{Qué buscamos.} Conseguir términos con 6 veces x en las dos ecuaciones.
\par\smallskip\textbf{Por qué.} Tener 6 veces la misma letra en ambos lados de la resta nos permitirá cancelarla.
\par\smallskip {\small\renewcommand{\arraystretch}{1.3}\begin{tabular}{>{\raggedright\arraybackslash}p{2.58cm}>{\raggedright\arraybackslash}p{2.47cm}>{\raggedright\arraybackslash}p{6.17cm}}\hline\rowcolor{black!7} \textbf{Dónde miramos} & \textbf{Qué término hay} & \textbf{Qué necesitamos hacer}\\\hline Ecuación 1 & $2x$ & Multiplicar TODA esta ecuación por 3\\ Ecuación 3 & $3x$ & Multiplicar TODA esta ecuación por 2\\ \hline\end{tabular}}\par\smallskip
\par\smallskip\textbf{Qué aprendimos.} Estos multiplicadores se eligen para igualar los términos con x. Los demás términos también se multiplican para conservar la igualdad.
\end{minipage}\par\medskip
\noindent\begin{minipage}{\linewidth}\textbf{7. Multiplicamos la ecuación 1 por 3}\\[2pt]Multiplicar una ecuación por 3 significa multiplicar CADA término de la izquierda y también el número de la derecha. La igualdad se conserva porque hacemos lo mismo en ambos lados.
\par\smallskip Esta es la ecuación antes de multiplicar:
\[2x + 3y - z = 6\]
\par\smallskip {\small\renewcommand{\arraystretch}{1.3}\begin{tabular}{>{\raggedright\arraybackslash}p{2.80cm}>{\raggedright\arraybackslash}p{1.91cm}>{\raggedright\arraybackslash}p{6.50cm}}\hline\rowcolor{black!7} \textbf{Parte de la ecuación} & \textbf{Antes} & \textbf{La cuenta y su resultado}\\\hline Término con x & $2x$ & $3\cdot\left(2x\right)=6x$\\ Término con y & $3y$ & $3\cdot\left(3y\right)=9y$\\ Término con z & $-z$ & $3\cdot\left(-z\right)=-3z$\\ Número a la derecha del igual & $6$ & $3\cdot\left(6\right)=18$\\ \hline\end{tabular}}\par\smallskip
\par\smallskip Después de hacer todas las multiplicaciones queda:
\[6x + 9y - 3z = 18\]
\par\smallskip\textbf{Qué aprendimos.} La misma condición ahora está escrita con los términos multiplicados. Ya tenemos el tamaño que buscábamos para x.
\par\smallskip\textbf{Comprueba que lo entiendes.} ¿Podríamos multiplicar solamente el término con x?\par\textit{Respuesta:} No. Eso cambiaría la condición. Debemos multiplicar todos los términos, también el número que está después del igual.
\end{minipage}\par\medskip
\noindent\begin{minipage}{\linewidth}\textbf{8. Multiplicamos la ecuación 3 por 2}\\[2pt]Multiplicar una ecuación por 2 significa multiplicar CADA término de la izquierda y también el número de la derecha. La igualdad se conserva porque hacemos lo mismo en ambos lados.
\par\smallskip Esta es la ecuación antes de multiplicar:
\[3x - y + z = -7\]
\par\smallskip {\small\renewcommand{\arraystretch}{1.3}\begin{tabular}{>{\raggedright\arraybackslash}p{2.80cm}>{\raggedright\arraybackslash}p{1.91cm}>{\raggedright\arraybackslash}p{6.50cm}}\hline\rowcolor{black!7} \textbf{Parte de la ecuación} & \textbf{Antes} & \textbf{La cuenta y su resultado}\\\hline Término con x & $3x$ & $2\cdot\left(3x\right)=6x$\\ Término con y & $-y$ & $2\cdot\left(-y\right)=-2y$\\ Término con z & $z$ & $2\cdot\left(z\right)=2z$\\ Número a la derecha del igual & $-7$ & $2\cdot\left(-7\right)=-14$\\ \hline\end{tabular}}\par\smallskip
\par\smallskip Después de hacer todas las multiplicaciones queda:
\[6x - 2y + 2z = -14\]
\par\smallskip\textbf{Qué aprendimos.} La misma condición ahora está escrita con los términos multiplicados. Ya tenemos el tamaño que buscábamos para x.
\par\smallskip\textbf{Comprueba que lo entiendes.} ¿Podríamos multiplicar solamente el término con x?\par\textit{Respuesta:} No. Eso cambiaría la condición. Debemos multiplicar todos los términos, también el número que está después del igual.
\end{minipage}\par\medskip
\noindent\begin{minipage}{\linewidth}\textbf{9. Restamos las dos ecuaciones preparadas}\\[2pt]Restamos la ecuación 1 preparada a la ecuación 3 preparada. Juntamos cada término solo con otro que tenga la misma letra. Los números sin letra se calculan por separado.
\par\smallskip {\small\renewcommand{\arraystretch}{1.3}\begin{tabular}{>{\raggedright\arraybackslash}p{2.80cm}>{\raggedright\arraybackslash}p{5.38cm}>{\raggedright\arraybackslash}p{3.03cm}}\hline\rowcolor{black!7} \textbf{Qué agrupamos} & \textbf{Cuenta} & \textbf{Qué significa}\\\hline Los términos con x & $6x - \left(6x\right)=0$ & Se cancela esta letra.\\ Los términos con y & $-2y - \left(9y\right)=-11y$ & Quedan -11y.\\ Los términos con z & $2z - \left(-3z\right)=5z$ & Quedan 5z.\\ Los números a la derecha del igual & $-14 - \left(18\right)=-32$ & Este será el número de la nueva igualdad.\\ \hline\end{tabular}}\par\smallskip
\par\smallskip Primera igualdad que usaremos:
\[6x - 2y + 2z = -14\]
\par\smallskip Le restamos esta igualdad:
\[6x + 9y - 3z = 18\]
\par\smallskip\textbf{Qué aprendimos.} Los términos con x dan cero. Los términos con letras distintas no se mezclan.
\end{minipage}\par\medskip
\noindent\begin{minipage}{\linewidth}\textbf{10. Mira lo que hemos conseguido}\\[2pt]La nueva ecuación es el resultado de la cuenta anterior. Todavía contiene y, z. Seguiremos reduciendo las condiciones.
\par\smallskip Escribimos la nueva condición:
\[-11y + 5z = -32\]
\[\left\{\begin{aligned}2x + 3y - z &= 6\\-4y + 3z &= -14\\-11y + 5z &= -32\end{aligned}\right.\]
\par\smallskip\textbf{Qué aprendimos.} Conservamos la ecuación 1 y reemplazamos la 3 por esta nueva condición. De esa manera conservamos toda la información necesaria para resolver el sistema.
\end{minipage}\par\medskip
\noindent\begin{minipage}{\linewidth}\textbf{11. ¿Cómo haremos desaparecer y?}\\[2pt]Trabajaremos con las ecuaciones 2 y 3. Queremos que los términos con y tengan el mismo tamaño. Así podremos restarlos.
\par\smallskip\textbf{Qué buscamos.} Conseguir términos con 44 veces y en las dos ecuaciones.
\par\smallskip\textbf{Por qué.} Tener 44 veces la misma letra en ambos lados de la resta nos permitirá cancelarla.
\par\smallskip {\small\renewcommand{\arraystretch}{1.3}\begin{tabular}{>{\raggedright\arraybackslash}p{2.58cm}>{\raggedright\arraybackslash}p{2.47cm}>{\raggedright\arraybackslash}p{6.17cm}}\hline\rowcolor{black!7} \textbf{Dónde miramos} & \textbf{Qué término hay} & \textbf{Qué necesitamos hacer}\\\hline Ecuación 2 & $-4y$ & Multiplicar TODA esta ecuación por 11\\ Ecuación 3 & $-11y$ & Multiplicar TODA esta ecuación por 4\\ \hline\end{tabular}}\par\smallskip
\par\smallskip\textbf{Qué aprendimos.} Estos multiplicadores se eligen para igualar los términos con y. Los demás términos también se multiplican para conservar la igualdad.
\end{minipage}\par\medskip
\noindent\begin{minipage}{\linewidth}\textbf{12. Multiplicamos la ecuación 2 por 11}\\[2pt]Multiplicar una ecuación por 11 significa multiplicar CADA término de la izquierda y también el número de la derecha. La igualdad se conserva porque hacemos lo mismo en ambos lados.
\par\smallskip Esta es la ecuación antes de multiplicar:
\[-4y + 3z = -14\]
\par\smallskip {\small\renewcommand{\arraystretch}{1.3}\begin{tabular}{>{\raggedright\arraybackslash}p{2.80cm}>{\raggedright\arraybackslash}p{1.91cm}>{\raggedright\arraybackslash}p{6.50cm}}\hline\rowcolor{black!7} \textbf{Parte de la ecuación} & \textbf{Antes} & \textbf{La cuenta y su resultado}\\\hline Término con y & $-4y$ & $11\cdot\left(-4y\right)=-44y$\\ Término con z & $3z$ & $11\cdot\left(3z\right)=33z$\\ Número a la derecha del igual & $-14$ & $11\cdot\left(-14\right)=-154$\\ \hline\end{tabular}}\par\smallskip
\par\smallskip Después de hacer todas las multiplicaciones queda:
\[-44y + 33z = -154\]
\par\smallskip\textbf{Qué aprendimos.} La misma condición ahora está escrita con los términos multiplicados. Ya tenemos el tamaño que buscábamos para y.
\par\smallskip\textbf{Comprueba que lo entiendes.} ¿Podríamos multiplicar solamente el término con y?\par\textit{Respuesta:} No. Eso cambiaría la condición. Debemos multiplicar todos los términos, también el número que está después del igual.
\end{minipage}\par\medskip
\noindent\begin{minipage}{\linewidth}\textbf{13. Multiplicamos la ecuación 3 por 4}\\[2pt]Multiplicar una ecuación por 4 significa multiplicar CADA término de la izquierda y también el número de la derecha. La igualdad se conserva porque hacemos lo mismo en ambos lados.
\par\smallskip Esta es la ecuación antes de multiplicar:
\[-11y + 5z = -32\]
\par\smallskip {\small\renewcommand{\arraystretch}{1.3}\begin{tabular}{>{\raggedright\arraybackslash}p{2.80cm}>{\raggedright\arraybackslash}p{1.91cm}>{\raggedright\arraybackslash}p{6.50cm}}\hline\rowcolor{black!7} \textbf{Parte de la ecuación} & \textbf{Antes} & \textbf{La cuenta y su resultado}\\\hline Término con y & $-11y$ & $4\cdot\left(-11y\right)=-44y$\\ Término con z & $5z$ & $4\cdot\left(5z\right)=20z$\\ Número a la derecha del igual & $-32$ & $4\cdot\left(-32\right)=-128$\\ \hline\end{tabular}}\par\smallskip
\par\smallskip Después de hacer todas las multiplicaciones queda:
\[-44y + 20z = -128\]
\par\smallskip\textbf{Qué aprendimos.} La misma condición ahora está escrita con los términos multiplicados. Ya tenemos el tamaño que buscábamos para y.
\par\smallskip\textbf{Comprueba que lo entiendes.} ¿Podríamos multiplicar solamente el término con y?\par\textit{Respuesta:} No. Eso cambiaría la condición. Debemos multiplicar todos los términos, también el número que está después del igual.
\end{minipage}\par\medskip
\noindent\begin{minipage}{\linewidth}\textbf{14. Restamos las dos ecuaciones preparadas}\\[2pt]Restamos la ecuación 2 preparada a la ecuación 3 preparada. Juntamos cada término solo con otro que tenga la misma letra. Los números sin letra se calculan por separado.
\par\smallskip {\small\renewcommand{\arraystretch}{1.3}\begin{tabular}{>{\raggedright\arraybackslash}p{2.80cm}>{\raggedright\arraybackslash}p{5.38cm}>{\raggedright\arraybackslash}p{3.03cm}}\hline\rowcolor{black!7} \textbf{Qué agrupamos} & \textbf{Cuenta} & \textbf{Qué significa}\\\hline Los términos con y & $-44y - \left(-44y\right)=0$ & Se cancela esta letra.\\ Los términos con z & $20z - \left(33z\right)=-13z$ & Quedan -13z.\\ Los números a la derecha del igual & $-128 - \left(-154\right)=-128+154=26$ & Este será el número de la nueva igualdad.\\ \hline\end{tabular}}\par\smallskip
\par\smallskip Primera igualdad que usaremos:
\[-44y + 20z = -128\]
\par\smallskip Le restamos esta igualdad:
\[-44y + 33z = -154\]
\par\smallskip\textbf{Qué aprendimos.} Los términos con y dan cero. Recuerda: restar un número negativo equivale a sumar su opuesto.
\end{minipage}\par\medskip
\noindent\begin{minipage}{\linewidth}\textbf{15. Mira lo que hemos conseguido}\\[2pt]La nueva ecuación es el resultado de la cuenta anterior. Solo contiene z: después podremos despejarla.
\par\smallskip Escribimos la nueva condición:
\[-13z = 26\]
\[\left\{\begin{aligned}2x + 3y - z &= 6\\-4y + 3z &= -14\\-13z &= 26\end{aligned}\right.\]
\par\smallskip\textbf{Qué aprendimos.} Conservamos la ecuación 2 y reemplazamos la 3 por esta nueva condición. De esa manera conservamos toda la información necesaria para resolver el sistema.
\end{minipage}\par\medskip
\noindent\begin{minipage}{\linewidth}\textbf{16. Ya tenemos una ecuación con una sola letra}\\[2pt]En esta condición solamente falta conocer z. El número -13 está multiplicando a esa letra. Ahora podemos deshacer esa multiplicación.
\par\smallskip Esta es la condición que vamos a resolver:
\[-13z = 26\]
\end{minipage}\par\medskip
\noindent\begin{minipage}{\linewidth}\textbf{17. Calculamos cuánto vale z}\\[2pt]Para deshacer una multiplicación por -13, dividimos AMBOS lados entre -13. En el lado izquierdo queda una sola z. En el derecho hacemos la división.
\par\smallskip Dividimos los dos lados entre el mismo número:
\[\frac{-13z}{-13}=\frac{26}{-13}\]
\par\smallskip Ya podemos leer el valor de la letra:
\[z=\frac{26}{-13}=-2\]
\par\smallskip\textbf{Qué aprendimos.} Este valor se usará en las ecuaciones donde todavía falten letras.
\end{minipage}\par\medskip
\noindent\begin{minipage}{\linewidth}\textbf{18. Usamos lo que ya sabemos para hallar y}\\[2pt]Ya sabemos z = -2. Sustituir significa escribir ese valor cada vez que aparezca su letra. Usamos paréntesis para que se vea qué número estamos multiplicando.
\par\smallskip\textbf{Qué buscamos.} Dejar solamente y como letra desconocida en esta ecuación.
\par\smallskip Tomamos esta condición que obtuvimos antes:
\[-4y + 3z = -14\]
\par\smallskip Reemplazamos las letras conocidas por sus valores:
\[-4y + 3\left(-2\right)=-14\]
\end{minipage}\par\medskip
\noindent\begin{minipage}{\linewidth}\textbf{19. Calculamos las multiplicaciones que ya se pueden hacer}\\[2pt]La letra y sigue sin conocerse. Las otras partes ya tienen un valor o una expresión conocida y podemos efectuar sus productos.
\par\smallskip {\small\renewcommand{\arraystretch}{1.3}\begin{tabular}{>{\raggedright\arraybackslash}p{3.36cm}>{\raggedright\arraybackslash}p{4.48cm}>{\raggedright\arraybackslash}p{3.36cm}}\hline\rowcolor{black!7} \textbf{Parte conocida} & \textbf{La multiplicación} & \textbf{Resultado}\\\hline El término 3z & $3\cdot\left(-2\right)$ & $-6$\\ \hline\end{tabular}}\par\smallskip
\par\smallskip Al reunir esas partes conocidas, la ecuación queda:
\[-4y - 6=-14\]
\end{minipage}\par\medskip
\noindent\begin{minipage}{\linewidth}\textbf{20. Dejamos el término con y solo a un lado}\\[2pt]Para quitar la parte conocida, sumamos 6 en AMBOS lados. Piensa en una balanza: si hacemos lo mismo a ambos lados, la igualdad se mantiene.
\par\smallskip Hacemos la misma operación a la izquierda y a la derecha:
\[-4y - 6 + 6=-14 + 6\]
\par\smallskip Las partes que se compensan dan cero; queda:
\[-4y=-8\]
\par\smallskip\textbf{Qué aprendimos.} «Pasar al otro lado» es una forma abreviada de decir que hicimos la misma suma o resta en los dos lados.
\end{minipage}\par\medskip
\noindent\begin{minipage}{\linewidth}\textbf{21. Calculamos cuánto vale y}\\[2pt]Para deshacer una multiplicación por -4, dividimos AMBOS lados entre -4. En el lado izquierdo queda una sola y. En el derecho hacemos la división.
\par\smallskip Dividimos los dos lados entre el mismo número:
\[\frac{-4y}{-4}=\frac{-8}{-4}\]
\par\smallskip Ya podemos leer el valor de la letra:
\[y=\frac{-8}{-4}=2\]
\par\smallskip\textbf{Qué aprendimos.} Un número negativo dividido entre otro negativo da positivo.
\end{minipage}\par\medskip
\noindent\begin{minipage}{\linewidth}\textbf{22. Usamos lo que ya sabemos para hallar x}\\[2pt]Ya sabemos y = 2; z = -2. Sustituir significa escribir ese valor cada vez que aparezca su letra. Usamos paréntesis para que se vea qué número estamos multiplicando.
\par\smallskip\textbf{Qué buscamos.} Dejar solamente x como letra desconocida en esta ecuación.
\par\smallskip Tomamos esta condición que obtuvimos antes:
\[2x + 3y - z = 6\]
\par\smallskip Reemplazamos las letras conocidas por sus valores:
\[2x + 3\left(2\right) - \left(-2\right)=6\]
\end{minipage}\par\medskip
\noindent\begin{minipage}{\linewidth}\textbf{23. Calculamos las multiplicaciones que ya se pueden hacer}\\[2pt]La letra x sigue sin conocerse. Las otras partes ya tienen un valor o una expresión conocida y podemos efectuar sus productos.
\par\smallskip {\small\renewcommand{\arraystretch}{1.3}\begin{tabular}{>{\raggedright\arraybackslash}p{3.36cm}>{\raggedright\arraybackslash}p{4.48cm}>{\raggedright\arraybackslash}p{3.36cm}}\hline\rowcolor{black!7} \textbf{Parte conocida} & \textbf{La multiplicación} & \textbf{Resultado}\\\hline El término 3y & $3\cdot\left(2\right)$ & $6$\\ El término -z & $-1\cdot\left(-2\right)$ & $2$\\ \hline\end{tabular}}\par\smallskip
\par\smallskip Al reunir esas partes conocidas, la ecuación queda:
\[2x + 8=6\]
\end{minipage}\par\medskip
\noindent\begin{minipage}{\linewidth}\textbf{24. Dejamos el término con x solo a un lado}\\[2pt]Para quitar la parte conocida, restamos 8 en AMBOS lados. Piensa en una balanza: si hacemos lo mismo a ambos lados, la igualdad se mantiene.
\par\smallskip Hacemos la misma operación a la izquierda y a la derecha:
\[2x + 8 - 8=6 - 8\]
\par\smallskip Las partes que se compensan dan cero; queda:
\[2x=-2\]
\par\smallskip\textbf{Qué aprendimos.} «Pasar al otro lado» es una forma abreviada de decir que hicimos la misma suma o resta en los dos lados.
\end{minipage}\par\medskip
\noindent\begin{minipage}{\linewidth}\textbf{25. Calculamos cuánto vale x}\\[2pt]Para deshacer una multiplicación por 2, dividimos AMBOS lados entre 2. En el lado izquierdo queda una sola x. En el derecho hacemos la división.
\par\smallskip Dividimos los dos lados entre el mismo número:
\[\frac{2x}{2}=\frac{-2}{2}\]
\par\smallskip Ya podemos leer el valor de la letra:
\[x=\frac{-2}{2}=-1\]
\par\smallskip\textbf{Qué aprendimos.} Este valor se usará en las ecuaciones donde todavía falten letras.
\end{minipage}\par\medskip
\noindent\begin{minipage}{\linewidth}\textbf{26. Reunimos la respuesta; todavía falta comprobarla}\\[2pt]Escribimos todos los valores en el orden x, y, z. Estos valores deben funcionar juntos, no cada uno en una ecuación diferente.
\[\left\{\begin{aligned}x&=-1\\y&=2\\z&=-2\end{aligned}\right.\]
\par\smallskip\textbf{Qué aprendimos.} El siguiente paso es volver a las ecuaciones del enunciado y comprobar cada condición.
\end{minipage}\par\medskip
\noindent\begin{minipage}{\linewidth}\textbf{27. ¿La respuesta cumple la ecuación 1?}\\[2pt]Comprobar significa volver a la condición original y calcularla con nuestra respuesta. Si los dos lados dan el mismo número, esta condición sí se cumple. Deben cumplirse todas.
\par\smallskip La condición del enunciado era:
\[2x + 3y - z = 6\]
\par\smallskip Ponemos nuestros valores donde estaban las letras:
\[2\left(-1\right) + 3\left(2\right) - \left(-2\right)=6\]
\par\smallskip Hacemos las multiplicaciones y luego la suma:
\[-2 + 6 + 2=6\]
\par\smallskip Comparamos los dos lados:
\[6=6\]
\par\smallskip\textbf{Qué aprendimos.} La igualdad final es verdadera. Nuestra respuesta sí cumple la ecuación 1.
\end{minipage}\par\medskip
\noindent\begin{minipage}{\linewidth}\textbf{28. ¿La respuesta cumple la ecuación 2?}\\[2pt]Comprobar significa volver a la condición original y calcularla con nuestra respuesta. Si los dos lados dan el mismo número, esta condición sí se cumple. Deben cumplirse todas.
\par\smallskip La condición del enunciado era:
\[2x - y + 2z = -8\]
\par\smallskip Ponemos nuestros valores donde estaban las letras:
\[2\left(-1\right) - \left(2\right) + 2\left(-2\right)=-8\]
\par\smallskip Hacemos las multiplicaciones y luego la suma:
\[-2 - 2 - 4=-8\]
\par\smallskip Comparamos los dos lados:
\[-8=-8\]
\par\smallskip\textbf{Qué aprendimos.} La igualdad final es verdadera. Nuestra respuesta sí cumple la ecuación 2.
\end{minipage}\par\medskip
\noindent\begin{minipage}{\linewidth}\textbf{29. ¿La respuesta cumple la ecuación 3?}\\[2pt]Comprobar significa volver a la condición original y calcularla con nuestra respuesta. Si los dos lados dan el mismo número, esta condición sí se cumple. Deben cumplirse todas.
\par\smallskip La condición del enunciado era:
\[3x - y + z = -7\]
\par\smallskip Ponemos nuestros valores donde estaban las letras:
\[3\left(-1\right) - \left(2\right) + \left(-2\right)=-7\]
\par\smallskip Hacemos las multiplicaciones y luego la suma:
\[-3 - 2 - 2=-7\]
\par\smallskip Comparamos los dos lados:
\[-7=-7\]
\par\smallskip\textbf{Qué aprendimos.} La igualdad final es verdadera. Nuestra respuesta sí cumple la ecuación 3.
\end{minipage}\par\medskip
\respuesta{x = -1; y = 2; z = -2.}
\noindent\begin{minipage}{\linewidth}\centering\includegraphics[width=.96\linewidth,height=6.8cm,keepaspectratio]{../recursos/grafico_1_1-14.png}\\{\small Lectura visual complementaria; la demostración está en las operaciones.}\end{minipage}
\end{multicols*}
\newpage\fancyfoot[L]{\small Sección 1.1 | Ejercicio 15 | Paso a paso}\hypertarget{1.1-15}{}
\pdfbookmark[1]{1.1 / 15}{bm1.1-15}
\titulo{15 | Un parámetro que decide la compatibilidad}
\textbf{Enunciado}\par Dado el sistema lineal: (a) determine un valor de t para que el sistema tenga una solución. (b) determine un valor de t para que el sistema no tenga solución. (c) ¿Cuántos valores diferentes de t pueden seleccionarse en la parte (b)?
\[\left\{\begin{aligned}2x-y&=5\\4x-2y&=t\end{aligned}\right.\]
\begin{multicols*}{2}
\noindent\begin{minipage}{\linewidth}\textbf{1. Reconocemos los múltiplos}\\[2pt]La segunda ecuación tiene el doble de los coeficientes de la primera. Multiplicamos AMBOS lados de la primera por 2.
\par\smallskip\textbf{Qué buscamos.} Encontrar cuándo dos condiciones pueden cumplirse juntas al cambiar t. Aquí t es un número que el enunciado deja por elegir.
\[2(2x-y)=2(5)\]
\[4x-2y=10\]
\end{minipage}\par\medskip
\noindent\begin{minipage}{\linewidth}\textbf{2. Restamos las ecuaciones término por término}\\[2pt]Restamos la primera ecuación multiplicada por 2 a la segunda original. Incluimos los términos independientes.
\[(4x-2y)-(4x-2y)=t-10\]
\[(4-4)x+(-2+2)y=t-10\]
\[0=t-10\]
\end{minipage}\par\medskip
\noindent\begin{minipage}{\linewidth}\textbf{3. Parte (a): buscamos cuándo la igualdad es verdadera}\\[2pt]Para que cero sea igual a t menos 10, sumamos 10 a ambos lados. Con t = 10 las dos ecuaciones coinciden.
\[0=t-10\]
\[0+10=t-10+10\]
\[t=10\]
\end{minipage}\par\medskip
\noindent\begin{minipage}{\linewidth}\textbf{4. Describimos las soluciones cuando t = 10}\\[2pt]De 2x - y = 5 restamos 2x a ambos lados y luego multiplicamos por -1. Elegimos x libremente.
\[-y=5-2x\]
\[y=2x-5\]
\[x=s,\quad y=2s-5,\quad s\in\mathbb{R}\]
\end{minipage}\par\medskip
\noindent\begin{minipage}{\linewidth}\textbf{5. Parte (b): elegimos un valor que produzca contradicción}\\[2pt]Por ejemplo, con t = 0 la última ecuación queda imposible. No hay solución.
\[t=0\]
\[0=0-10\]
\[0=-10\]
\end{minipage}\par\medskip
\noindent\begin{minipage}{\linewidth}\textbf{6. Parte (c): contamos los valores posibles}\\[2pt]Cualquier real distinto de 10 hace falsa la igualdad 0 = t - 10. Hay infinitos valores para la parte (b). Ningún valor de t produce una solución única.
\[t\in\mathbb{R}\setminus\{10\}\]
\end{minipage}\par\medskip
\respuesta{(a) t = 10: infinitas soluciones. (b) Cualquier t distinto de 10; por ejemplo 0. (c) Infinitos valores.}
\noindent\begin{minipage}{\linewidth}\centering\includegraphics[width=.96\linewidth,height=6.8cm,keepaspectratio]{../recursos/grafico_1_1-15.png}\\{\small Lectura visual complementaria; la demostración está en las operaciones.}\end{minipage}
\end{multicols*}
\newpage\fancyfoot[L]{\small Sección 1.1 | Ejercicio 16 | Paso a paso}\hypertarget{1.1-16}{}
\pdfbookmark[1]{1.1 / 16}{bm1.1-16}
\titulo{16 | Sumar y multiplicar soluciones}
\textbf{Enunciado}\par Dado el sistema lineal: (a) verifique que (1, -1, -1) es una solución. (b) verifique que (-2, 2, 2) es una solución. (c) ¿Su suma (-1, 1, 1) es una solución? (d) ¿El triple de la suma es una solución?
\[\left\{\begin{aligned}2x+3y-z&=0\\x-4y+5z&=0\end{aligned}\right.\]
\begin{multicols*}{2}
\noindent\begin{minipage}{\linewidth}\textbf{1. Parte (a): comprobamos (1, -1, -1)}\\[2pt]Un punto es solución si satisface simultáneamente las dos ecuaciones. Cuidado: restar un número negativo cambia el signo.
\par\smallskip\textbf{Qué buscamos.} Revisar si cada grupo de números propuesto cumple las dos ecuaciones.
\[2(1)+3(-1)-(-1)=2-3+1=0\]
\[1-4(-1)+5(-1)=1+4-5=0\]
\end{minipage}\par\medskip
\noindent\begin{minipage}{\linewidth}\textbf{2. Parte (b): comprobamos (-2, 2, 2)}\\[2pt]Sustituimos de nuevo, conservando los signos.
\[2(-2)+3(2)-2=-4+6-2=0\]
\[-2-4(2)+5(2)=-2-8+10=0\]
\end{minipage}\par\medskip
\noindent\begin{minipage}{\linewidth}\textbf{3. Parte (c): sumamos componente a componente}\\[2pt]La suma es (-1, 1, 1). También satisface las dos ecuaciones.
\[(1,-1,-1)+(-2,2,2)=(-1,1,1)\]
\[2(-1)+3(1)-1=0\]
\[-1-4(1)+5(1)=0\]
\end{minipage}\par\medskip
\noindent\begin{minipage}{\linewidth}\textbf{4. Parte (d): multiplicamos por 3}\\[2pt]El triple es (-3, 3, 3). Verificamos ambas igualdades.
\[3(-1,1,1)=(-3,3,3)\]
\[2(-3)+3(3)-3=0\]
\[-3-4(3)+5(3)=0\]
\end{minipage}\par\medskip
\noindent\begin{minipage}{\linewidth}\textbf{5. ¿Por qué la suma y el triple también funcionaron?}\\[2pt]Las ecuaciones de este ejercicio terminan en cero. Al sumar dos soluciones estamos sumando también sus lados derechos: cero más cero sigue siendo cero. Al triplicar una solución, el lado derecho es tres veces cero y sigue siendo cero. Esta explicación se apoya en que las letras aparecen solo multiplicadas por números y sumadas.
\par\smallskip Lo que ocurre a la derecha cuando sumamos soluciones:
\[0+0=0\]
\par\smallskip Lo que ocurre cuando triplicamos una solución:
\[3\cdot0=0\]
\par\smallskip\textbf{Qué aprendimos.} Un sistema cuyas ecuaciones tienen cero a la derecha se llama homogéneo. Aquí el nombre se refiere a esos ceros.
\end{minipage}\par\medskip
\respuesta{(a), (b), (c) y (d): sí son soluciones; la suma es (-1, 1, 1) y su triple es (-3, 3, 3).}
\noindent\begin{minipage}{\linewidth}\centering\includegraphics[width=.96\linewidth,height=6.8cm,keepaspectratio]{../recursos/grafico_1_1-16.png}\\{\small Lectura visual complementaria; la demostración está en las operaciones.}\end{minipage}
\end{multicols*}
\newpage\fancyfoot[L]{\small Sección 1.1 | Ejercicio 17 | Paso a paso}\hypertarget{1.1-17}{}
\pdfbookmark[1]{1.1 / 17}{bm1.1-17}
\titulo{17 | Sustitución hacia atrás}
\textbf{Enunciado}\par Resuelva el sistema lineal siguiente sin utilizar el método de eliminación.
\[\left\{\begin{aligned}2x+y-2z&=-5\\3y+z&=7\\z&=4\end{aligned}\right.\]
\begin{multicols*}{2}
\noindent\begin{minipage}{\linewidth}\textbf{1. La última ecuación ya da z}\\[2pt]No hacemos eliminación: este sistema ya está escalonado. Leemos la tercera ecuación.
\par\smallskip\textbf{Qué buscamos.} Encontrar z, después y y al final x, usando sustitución: reemplazar una letra por el valor que ya conocemos.
\[z=4\]
\end{minipage}\par\medskip
\noindent\begin{minipage}{\linewidth}\textbf{2. Sustituimos z en la segunda ecuación}\\[2pt]Reemplazamos z por 4 en 3y + z = 7.
\[3y+z=7\]
\[3y+4=7\]
\end{minipage}\par\medskip
\noindent\begin{minipage}{\linewidth}\textbf{3. Despejamos y}\\[2pt]Restamos 4 a ambos lados y dividimos entre 3.
\[3y+4-4=7-4\]
\[3y=3\]
\[y=\frac{3}{3}=1\]
\end{minipage}\par\medskip
\noindent\begin{minipage}{\linewidth}\textbf{4. Sustituimos en la primera ecuación}\\[2pt]Ahora conocemos y = 1 y z = 4. Primero efectuamos la multiplicación por 2.
\[2x+y-2z=-5\]
\[2x+1-2(4)=-5\]
\[2x+1-8=-5\]
\[2x-7=-5\]
\end{minipage}\par\medskip
\noindent\begin{minipage}{\linewidth}\textbf{5. Despejamos x}\\[2pt]Sumamos 7 a ambos lados; después dividimos entre 2.
\[2x-7+7=-5+7\]
\[2x=2\]
\[x=\frac{2}{2}=1\]
\end{minipage}\par\medskip
\noindent\begin{minipage}{\linewidth}\textbf{6. Reunimos la respuesta}\\[2pt]El orden de las coordenadas es x, y, z.
\[\left\{\begin{aligned}x&=1\\y&=1\\z&=4\end{aligned}\right.\]
\end{minipage}\par\medskip
\noindent\begin{minipage}{\linewidth}\textbf{7. Comprobamos en las ecuaciones originales}\\[2pt]Sustituimos en las tres ecuaciones; las tres igualdades son verdaderas.
\[2(1)+1-2(4)=2+1-8=-5\]
\[3(1)+4=3+4=7\]
\[z=4\]
\end{minipage}\par\medskip
\respuesta{(x, y, z) = (1, 1, 4).}
\noindent\begin{minipage}{\linewidth}\centering\includegraphics[width=.96\linewidth,height=6.8cm,keepaspectratio]{../recursos/grafico_1_1-17.png}\\{\small Lectura visual complementaria; la demostración está en las operaciones.}\end{minipage}
\end{multicols*}
\newpage\fancyfoot[L]{\small Sección 1.1 | Ejercicio 18 | Paso a paso}\hypertarget{1.1-18}{}
\pdfbookmark[1]{1.1 / 18}{bm1.1-18}
\titulo{18 | Sustitución hacia adelante}
\textbf{Enunciado}\par Resuelva el sistema lineal siguiente sin utilizar el método de eliminación.
\[\left\{\begin{aligned}4x&=8\\-2x+3y&=-1\\3x+5y-2z&=11\end{aligned}\right.\]
\begin{multicols*}{2}
\noindent\begin{minipage}{\linewidth}\textbf{1. Despejamos x}\\[2pt]Dividimos ambos lados de 4x = 8 entre 4.
\par\smallskip\textbf{Qué buscamos.} Encontrar x, después y y al final z. Empezamos por la ecuación que ya contiene una sola letra.
\[4x=8\]
\[\frac{4x}{4}=\frac{8}{4}\]
\[x=2\]
\end{minipage}\par\medskip
\noindent\begin{minipage}{\linewidth}\textbf{2. Sustituimos x en la segunda ecuación}\\[2pt]Reemplazamos x por 2 y efectuamos el producto.
\[-2x+3y=-1\]
\[-2(2)+3y=-1\]
\[-4+3y=-1\]
\end{minipage}\par\medskip
\noindent\begin{minipage}{\linewidth}\textbf{3. Despejamos y}\\[2pt]Sumamos 4 a ambos lados y luego dividimos entre 3.
\[-4+3y+4=-1+4\]
\[3y=3\]
\[y=\frac{3}{3}=1\]
\end{minipage}\par\medskip
\noindent\begin{minipage}{\linewidth}\textbf{4. Sustituimos en la tercera ecuación}\\[2pt]Usamos x = 2 y y = 1. Multiplicamos y sumamos los términos conocidos.
\[3x+5y-2z=11\]
\[3(2)+5(1)-2z=11\]
\[6+5-2z=11\]
\[11-2z=11\]
\end{minipage}\par\medskip
\noindent\begin{minipage}{\linewidth}\textbf{5. Despejamos z}\\[2pt]Restamos 11 a ambos lados y dividimos entre -2.
\[11-2z-11=11-11\]
\[-2z=0\]
\[z=\frac{0}{-2}=0\]
\end{minipage}\par\medskip
\noindent\begin{minipage}{\linewidth}\textbf{6. Respuesta}\\[2pt]Reunimos los valores calculados.
\[\left\{\begin{aligned}x&=2\\y&=1\\z&=0\end{aligned}\right.\]
\end{minipage}\par\medskip
\noindent\begin{minipage}{\linewidth}\textbf{7. Comprobación completa}\\[2pt]Verificamos las tres condiciones.
\[4(2)=8\]
\[-2(2)+3(1)=-4+3=-1\]
\[3(2)+5(1)-2(0)=6+5-0=11\]
\end{minipage}\par\medskip
\respuesta{(x, y, z) = (2, 1, 0).}
\noindent\begin{minipage}{\linewidth}\centering\includegraphics[width=.96\linewidth,height=6.8cm,keepaspectratio]{../recursos/grafico_1_1-18.png}\\{\small Lectura visual complementaria; la demostración está en las operaciones.}\end{minipage}
\end{multicols*}
\newpage\fancyfoot[L]{\small Sección 1.1 | Ejercicio 19 | Paso a paso}\hypertarget{1.1-19}{}
\pdfbookmark[1]{1.1 / 19}{bm1.1-19}
\titulo{19 | Un mismo valor debe cumplir todo el sistema}
\textbf{Enunciado}\par ¿Existe un valor de r tal que x = 1, y = 2, z = r sea una solución del siguiente sistema lineal? De ser así, determínelo.
\[\left\{\begin{aligned}2x+3y-z&=11\\x-y+2z&=-7\\4x+y-2z&=12\end{aligned}\right.\]
\begin{multicols*}{2}
\noindent\begin{minipage}{\linewidth}\textbf{1. Sustituimos el candidato en la primera ecuación}\\[2pt]Ponemos x = 1, y = 2 y z = r. Multiplicamos los valores conocidos y sumamos.
\par\smallskip\textbf{Qué buscamos.} Encontrar un solo valor de r que funcione en las tres condiciones.
\[2(1)+3(2)-r=11\]
\[2+6-r=11\]
\[8-r=11\]
\end{minipage}\par\medskip
\noindent\begin{minipage}{\linewidth}\textbf{2. Despejamos r en la primera ecuación}\\[2pt]Restamos 8 a ambos lados y multiplicamos por -1.
\[8-r-8=11-8\]
\[-r=3\]
\[r=-3\]
\end{minipage}\par\medskip
\noindent\begin{minipage}{\linewidth}\textbf{3. Usamos la segunda ecuación}\\[2pt]Sustituimos y despejamos. Esta condición debe exigir el mismo r.
\[1-2+2r=-7\]
\[-1+2r=-7\]
\[2r=-7+1=-6\]
\[r=\frac{-6}{2}=-3\]
\end{minipage}\par\medskip
\noindent\begin{minipage}{\linewidth}\textbf{4. Usamos la tercera ecuación}\\[2pt]Restamos 6 a ambos lados y dividimos entre -2.
\[4(1)+2-2r=12\]
\[6-2r=12\]
\[-2r=12-6=6\]
\[r=\frac{6}{-2}=-3\]
\end{minipage}\par\medskip
\noindent\begin{minipage}{\linewidth}\textbf{5. Verificamos el valor común}\\[2pt]Las tres ecuaciones requieren r = -3 y quedan satisfechas.
\[2+6-(-3)=8+3=11\]
\[1-2+2(-3)=-1-6=-7\]
\[4+2-2(-3)=6+6=12\]
\end{minipage}\par\medskip
\respuesta{Sí: r = -3.}
\noindent\begin{minipage}{\linewidth}\centering\includegraphics[width=.96\linewidth,height=6.8cm,keepaspectratio]{../recursos/grafico_1_1-19.png}\\{\small Lectura visual complementaria; la demostración está en las operaciones.}\end{minipage}
\end{multicols*}
\newpage\fancyfoot[L]{\small Sección 1.1 | Ejercicio 20 | Paso a paso}\hypertarget{1.1-20}{}
\pdfbookmark[1]{1.1 / 20}{bm1.1-20}
\titulo{20 | Comprobamos antes de concluir}
\textbf{Enunciado}\par ¿Existe un valor de r tal que x = r, y = 2, z = 1 sea una solución del siguiente sistema lineal? De ser así, determínelo.
\[\left\{\begin{aligned}3x-2z&=4\\x-4y+z&=-5\\-2x+3y+2z&=9\end{aligned}\right.\]
\begin{multicols*}{2}
\noindent\begin{minipage}{\linewidth}\textbf{1. Primera ecuación}\\[2pt]Sustituimos x = r, y = 2, z = 1. Sumamos 2 a ambos lados y dividimos entre 3.
\par\smallskip\textbf{Qué buscamos.} Decidir si un mismo r puede satisfacer las tres ecuaciones. Si piden valores diferentes, no existe ese r.
\[3r-2(1)=4\]
\[3r-2+2=4+2\]
\[3r=6\]
\[r=\frac{6}{3}=2\]
\end{minipage}\par\medskip
\noindent\begin{minipage}{\linewidth}\textbf{2. Segunda ecuación}\\[2pt]Primero multiplicamos; después sumamos 7 a ambos lados.
\[r-4(2)+1=-5\]
\[r-8+1=-5\]
\[r-7=-5\]
\[r=-5+7=2\]
\end{minipage}\par\medskip
\noindent\begin{minipage}{\linewidth}\textbf{3. Tercera ecuación}\\[2pt]Primero hacemos las multiplicaciones y luego restamos 8 a ambos lados.
\[-2r+3(2)+2(1)=9\]
\[-2r+6+2=9\]
\[-2r=9-8=1\]
\[r=\frac{1}{-2}=-\frac12\]
\end{minipage}\par\medskip
\noindent\begin{minipage}{\linewidth}\textbf{4. Comparamos lo que exige cada ecuación}\\[2pt]Las dos primeras necesitan r = 2; la tercera necesita r = -1/2. Un mismo valor no puede cumplir ambas exigencias. No existe el valor pedido.
\[2\neq-\frac12\]
\[r=2\quad\Longrightarrow\quad -2(2)+6+2=4\neq9\]
\end{minipage}\par\medskip
\respuesta{No existe tal valor de r.}
\noindent\begin{minipage}{\linewidth}\centering\includegraphics[width=.96\linewidth,height=6.8cm,keepaspectratio]{../recursos/grafico_1_1-20.png}\\{\small Lectura visual complementaria; la demostración está en las operaciones.}\end{minipage}
\end{multicols*}
\newpage\fancyfoot[L]{\small Sección 1.1 | Ejercicio 21 | Paso a paso}\hypertarget{1.1-21}{}
\pdfbookmark[1]{1.1 / 21}{bm1.1-21}
\titulo{21 | Tres planos: punto, vacío o recta}
\textbf{Enunciado}\par Diga cuál es el número de puntos que están simultáneamente en los tres planos que se muestran en cada inciso de la figura 1.2.
\begin{multicols*}{2}
\noindent\begin{minipage}{\linewidth}\textbf{1. Antes de contar: ¿qué representa un plano?}\\[2pt]Un plano es una superficie plana que se extiende sin terminar, como una hoja imaginaria infinita. Cada ecuación con tres letras puede representar un plano. Resolver las tres juntas significa buscar puntos que estén en los TRES planos al mismo tiempo.
\par\smallskip\textbf{Qué buscamos.} Identificar la parte que los tres planos comparten, no solo dónde se cruza una pareja.
\par\smallskip\textbf{Qué aprendimos.} Un punto común da una solución; una recta común contiene infinitos puntos y da infinitas soluciones; si no hay nada común a los tres, no hay solución.
\end{minipage}\par\medskip
\noindent\begin{minipage}{\linewidth}\textbf{2. Buscamos puntos comunes a los tres planos}\\[2pt]La figura 1.2 compara tres configuraciones. Una solución debe pertenecer simultáneamente a P1, P2 y P3. No basta con estar en dos de ellos.
\end{minipage}\par\medskip
\noindent\begin{minipage}{\linewidth}\textbf{3. Inciso (a): una única solución}\\[2pt]Los tres planos se encuentran en el punto negro señalado. Aunque dos planos se corten en una recta, al exigir que el punto esté también en el tercero solo queda ese punto.
\[P_1\cap P_2\cap P_3=\{Q\}\]
\end{minipage}\par\medskip
\noindent\begin{minipage}{\linewidth}\textbf{4. Inciso (b): ninguna solución}\\[2pt]P2 y P3 son planos paralelos distintos. No comparten puntos, de modo que tampoco puede existir uno que pertenezca a ambos y a P1. Que P1 corte a cada uno por separado no cambia esta conclusión.
\[P_2\cap P_3=\varnothing\]
\[P_1\cap P_2\cap P_3=\varnothing\]
\end{minipage}\par\medskip
\noindent\begin{minipage}{\linewidth}\textbf{5. Inciso (c): infinitas soluciones}\\[2pt]Los tres planos comparten la misma recta dibujada como arista. Cada punto de esa recta satisface las tres ecuaciones: hay infinitos puntos comunes.
\[P_1\cap P_2\cap P_3=L\]
\end{minipage}\par\medskip
\noindent\begin{minipage}{\linewidth}\textbf{6. Relacionamos la geometría con el álgebra}\\[2pt]Un punto corresponde a una solución única; ninguna intersección común corresponde a un sistema incompatible; una recta común corresponde a infinitas soluciones con un parámetro libre.
\end{minipage}\par\medskip
\respuesta{(a) Un punto. (b) Ningún punto. (c) Infinitos puntos sobre una recta.}
\noindent\begin{minipage}{\linewidth}\centering\includegraphics[width=.96\linewidth,height=6.8cm,keepaspectratio]{../recursos/figura_1_2.png}\\{\small Figura 1.2 suministrada en el enunciado.}\end{minipage}
\par\medskip\noindent\begin{minipage}{\linewidth}\centering\includegraphics[width=.96\linewidth,height=7cm,keepaspectratio]{../recursos/grafico_1_1-21.png}\\{\small Modelos cualitativos de los tres incisos: las coordenadas son ilustrativas.}\end{minipage}
\end{multicols*}
\newpage\fancyfoot[L]{\small Sección 1.1 | Ejercicio 22 | Paso a paso}\hypertarget{1.1-22}{}
\pdfbookmark[1]{1.1 / 22}{bm1.1-22}
\titulo{22 | ¿Qué comparten los tres planos?}
\textbf{Enunciado}\par Diga cuál es el número de puntos que están simultáneamente en los tres planos que se muestran en cada inciso de la figura 1.3.
\begin{multicols*}{2}
\noindent\begin{minipage}{\linewidth}\textbf{1. Antes de contar: ¿qué representa un plano?}\\[2pt]Un plano es una superficie plana que se extiende sin terminar, como una hoja imaginaria infinita. Cada ecuación con tres letras puede representar un plano. Resolver las tres juntas significa buscar puntos que estén en los TRES planos al mismo tiempo.
\par\smallskip\textbf{Qué buscamos.} Identificar la parte que los tres planos comparten, no solo dónde se cruza una pareja.
\par\smallskip\textbf{Qué aprendimos.} Un punto común da una solución; una recta común contiene infinitos puntos y da infinitas soluciones; si no hay nada común a los tres, no hay solución.
\end{minipage}\par\medskip
\noindent\begin{minipage}{\linewidth}\textbf{2. Leemos la figura suministrada}\\[2pt]Buscamos la intersección de los tres planos a la vez. Las zonas dibujadas son partes visibles de planos infinitos; sus bordes no limitan el plano.
\end{minipage}\par\medskip
\noindent\begin{minipage}{\linewidth}\textbf{3. Inciso (a): planos paralelos distintos}\\[2pt]Los tres planos están a diferentes alturas. No hay un punto que pertenezca a todos ellos. Número de puntos comunes: cero.
\[P_1\cap P_2\cap P_3=\varnothing\]
\end{minipage}\par\medskip
\noindent\begin{minipage}{\linewidth}\textbf{4. Inciso (b): una recta común}\\[2pt]Los tres planos comparten la recta que aparece como arista central. Cada punto de esa recta pertenece a los tres: hay infinitos puntos comunes.
\[P_1\cap P_2\cap P_3=L\]
\end{minipage}\par\medskip
\noindent\begin{minipage}{\linewidth}\textbf{5. Inciso (c): sin punto común a los tres}\\[2pt]Los dos planos inclinados se cortan en la recta superior. Esa recta es paralela al plano inferior y no está contenida en él. Aunque cada pareja se intersecta, no hay un punto común a los tres.
\[P_1\cap P_2\cap P_3=\varnothing\]
\end{minipage}\par\medskip
\noindent\begin{minipage}{\linewidth}\textbf{6. Recordatorio}\\[2pt]Para responder no contamos cuántas rectas de intersección de parejas vemos, sino cuántos puntos están en los tres planos simultáneamente.
\end{minipage}\par\medskip
\respuesta{(a) Cero. (b) Infinitos, sobre una recta. (c) Cero.}
\noindent\begin{minipage}{\linewidth}\centering\includegraphics[width=.96\linewidth,height=6.8cm,keepaspectratio]{../recursos/figura_1_3.png}\\{\small Figura 1.3 suministrada en el enunciado.}\end{minipage}
\par\medskip\noindent\begin{minipage}{\linewidth}\centering\includegraphics[width=.96\linewidth,height=7cm,keepaspectratio]{../recursos/grafico_1_1-22.png}\\{\small Modelos cualitativos de los tres incisos: las coordenadas son ilustrativas.}\end{minipage}
\end{multicols*}
\newpage\fancyfoot[L]{\small Sección 1.1 | Ejercicio 23 | Paso a paso}\hypertarget{1.1-23}{}
\pdfbookmark[1]{1.1 / 23}{bm1.1-23}
\titulo{23 | Producción de gasolina}
\textbf{Enunciado}\par Una refinería produce gasolina con azufre y sin azufre. Para producir cada tonelada de gasolina sin azufre se requieren 5 minutos en la planta mezcladora y 4 minutos en la planta de refinación, mientras que cada tonelada de gasolina con azufre requiere 4 minutos en la planta mezcladora y 2 minutos en la planta de refinación. Si la planta mezcladora está disponible 3 horas y la de refinación 2 horas, ¿cuántas toneladas de cada tipo de gasolina deben producirse de modo que las plantas operen a toda su capacidad?
\[\left\{\begin{aligned}5x + 4y &= 180\\4x + 2y &= 120\end{aligned}\right.\]
\begin{multicols*}{2}
\noindent\begin{minipage}{\linewidth}\textbf{1. Definimos las cantidades}\\[2pt]x = toneladas de gasolina sin azufre; y = toneladas de gasolina con azufre. Las incógnitas son cantidades de producto, no tiempos.
\par\smallskip\textbf{Por qué.} Una letra representa una cantidad de producto que debemos decidir. Los minutos por tonelada son datos del problema.
\par\smallskip {\small\renewcommand{\arraystretch}{1.3}\begin{tabular}{>{\raggedright\arraybackslash}p{2.68cm}>{\raggedright\arraybackslash}p{2.68cm}>{\raggedright\arraybackslash}p{2.68cm}>{\raggedright\arraybackslash}p{2.68cm}}\hline\rowcolor{black!7} \textbf{Gasolina} & \textbf{Letra: toneladas} & \textbf{Minutos mezcladora} & \textbf{Minutos refinación}\\\hline Sin azufre & x & 5 por tonelada & 4 por tonelada\\ Con azufre & y & 4 por tonelada & 2 por tonelada\\ \hline\end{tabular}}\par\smallskip
\end{minipage}\par\medskip
\noindent\begin{minipage}{\linewidth}\textbf{2. Unificamos las unidades}\\[2pt]Los consumos están dados en minutos por tonelada. Convertimos las horas disponibles a minutos antes de igualar.
\[3\,\mathrm{h}=180\,\mathrm{min}\]
\[2\,\mathrm{h}=120\,\mathrm{min}\]
\end{minipage}\par\medskip
\noindent\begin{minipage}{\linewidth}\textbf{3. Construimos las ecuaciones}\\[2pt]En la mezcladora, una tonelada sin azufre usa 5 minutos; x toneladas usan 5x minutos. La gasolina con azufre aporta 4y minutos. Sumamos ambos tiempos y los igualamos a los 180 minutos disponibles. Hacemos el mismo razonamiento en refinación.
\[5x+4y=180\quad(\mathrm{mezcladora})\]
\[4x+2y=120\quad(\mathrm{refinacion})\]
\par\smallskip\textbf{Qué aprendimos.} Usamos igualdades porque el enunciado exige ocupar todo el tiempo disponible en ambas plantas.
\end{minipage}\par\medskip
\noindent\begin{minipage}{\linewidth}\textbf{4. Primero: ¿qué significa resolver este sistema?}\\[2pt]Las letras x, y representan números que todavía no conocemos. Una solución es una elección de esos números que hace verdaderas TODAS las ecuaciones al mismo tiempo. La llave de la izquierda significa que las condiciones se deben cumplir juntas.
\par\smallskip\textbf{Qué buscamos.} Encontrar una misma pareja de valores.
\par\smallskip El número 3 delante de x indica tres veces ese mismo valor; es decir, sumar tres copias:
\[3x=x+x+x\]
\par\smallskip Ejemplo numérico de la notación: si x valiera 2, tres veces ese valor sería 6. Este 2 es solo un ejemplo; todavía no es la respuesta del ejercicio. El punto centrado indica multiplicación:
\[3\cdot 2=2+2+2=6\]
\par\smallskip\textbf{Comprueba que lo entiendes.} ¿Basta con que los valores funcionen en una sola ecuación?\par\textit{Respuesta:} No. Deben funcionar en todas las ecuaciones del sistema, usando los mismos valores.
\end{minipage}\par\medskip
\noindent\begin{minipage}{\linewidth}\textbf{5. Nuestro plan: ir quitando letras}\\[2pt]Una ecuación con una sola letra es más fácil de resolver. Vamos a combinar ecuaciones para que una letra desaparezca. Cuando sepamos un valor, lo pondremos en las ecuaciones que faltan.
\par\smallskip\textbf{Qué buscamos.} Primero quitar x para calcular y. Después calcular x.
\par\smallskip\textbf{Por qué.} Si una cantidad se resta de sí misma, queda cero. Usaremos esta idea con términos que tengan la misma letra.
\par\smallskip Ejemplo de la idea que utilizaremos:
\[3x-3x=0\]
\par\smallskip\textbf{Qué aprendimos.} Eliminar una letra de una ecuación NO significa que su valor sea cero; significa que sus términos se cancelaron en esa cuenta.
\end{minipage}\par\medskip
\noindent\begin{minipage}{\linewidth}\textbf{6. Ponemos una ecuación más cómoda arriba}\\[2pt]Cambiamos de lugar las ecuaciones 1 y 2. Las condiciones siguen siendo las mismas; solamente cambia el orden en que las leemos.
\par\smallskip\textbf{Por qué.} La ecuación que subimos tiene un número más pequeño delante de x. Eso facilita las multiplicaciones que haremos a continuación.
\[\left\{\begin{aligned}4x + 2y &= 120\\5x + 4y &= 180\end{aligned}\right.\]
\end{minipage}\par\medskip
\noindent\begin{minipage}{\linewidth}\textbf{7. ¿Cómo haremos desaparecer x?}\\[2pt]Trabajaremos con las ecuaciones 1 y 2. Queremos que los términos con x tengan el mismo tamaño. Así podremos restarlos.
\par\smallskip\textbf{Qué buscamos.} Conseguir términos con 20 veces x en las dos ecuaciones.
\par\smallskip\textbf{Por qué.} Tener 20 veces la misma letra en ambos lados de la resta nos permitirá cancelarla.
\par\smallskip {\small\renewcommand{\arraystretch}{1.3}\begin{tabular}{>{\raggedright\arraybackslash}p{2.58cm}>{\raggedright\arraybackslash}p{2.47cm}>{\raggedright\arraybackslash}p{6.17cm}}\hline\rowcolor{black!7} \textbf{Dónde miramos} & \textbf{Qué término hay} & \textbf{Qué necesitamos hacer}\\\hline Ecuación 1 & $4x$ & Multiplicar TODA esta ecuación por 5\\ Ecuación 2 & $5x$ & Multiplicar TODA esta ecuación por 4\\ \hline\end{tabular}}\par\smallskip
\par\smallskip\textbf{Qué aprendimos.} Estos multiplicadores se eligen para igualar los términos con x. Los demás términos también se multiplican para conservar la igualdad.
\end{minipage}\par\medskip
\noindent\begin{minipage}{\linewidth}\textbf{8. Multiplicamos la ecuación 1 por 5}\\[2pt]Multiplicar una ecuación por 5 significa multiplicar CADA término de la izquierda y también el número de la derecha. La igualdad se conserva porque hacemos lo mismo en ambos lados.
\par\smallskip Esta es la ecuación antes de multiplicar:
\[4x + 2y = 120\]
\par\smallskip {\small\renewcommand{\arraystretch}{1.3}\begin{tabular}{>{\raggedright\arraybackslash}p{2.80cm}>{\raggedright\arraybackslash}p{1.91cm}>{\raggedright\arraybackslash}p{6.50cm}}\hline\rowcolor{black!7} \textbf{Parte de la ecuación} & \textbf{Antes} & \textbf{La cuenta y su resultado}\\\hline Término con x & $4x$ & $5\cdot\left(4x\right)=20x$\\ Término con y & $2y$ & $5\cdot\left(2y\right)=10y$\\ Número a la derecha del igual & $120$ & $5\cdot\left(120\right)=600$\\ \hline\end{tabular}}\par\smallskip
\par\smallskip Después de hacer todas las multiplicaciones queda:
\[20x + 10y = 600\]
\par\smallskip\textbf{Qué aprendimos.} La misma condición ahora está escrita con los términos multiplicados. Ya tenemos el tamaño que buscábamos para x.
\par\smallskip\textbf{Comprueba que lo entiendes.} ¿Podríamos multiplicar solamente el término con x?\par\textit{Respuesta:} No. Eso cambiaría la condición. Debemos multiplicar todos los términos, también el número que está después del igual.
\end{minipage}\par\medskip
\noindent\begin{minipage}{\linewidth}\textbf{9. Multiplicamos la ecuación 2 por 4}\\[2pt]Multiplicar una ecuación por 4 significa multiplicar CADA término de la izquierda y también el número de la derecha. La igualdad se conserva porque hacemos lo mismo en ambos lados.
\par\smallskip Esta es la ecuación antes de multiplicar:
\[5x + 4y = 180\]
\par\smallskip {\small\renewcommand{\arraystretch}{1.3}\begin{tabular}{>{\raggedright\arraybackslash}p{2.80cm}>{\raggedright\arraybackslash}p{1.91cm}>{\raggedright\arraybackslash}p{6.50cm}}\hline\rowcolor{black!7} \textbf{Parte de la ecuación} & \textbf{Antes} & \textbf{La cuenta y su resultado}\\\hline Término con x & $5x$ & $4\cdot\left(5x\right)=20x$\\ Término con y & $4y$ & $4\cdot\left(4y\right)=16y$\\ Número a la derecha del igual & $180$ & $4\cdot\left(180\right)=720$\\ \hline\end{tabular}}\par\smallskip
\par\smallskip Después de hacer todas las multiplicaciones queda:
\[20x + 16y = 720\]
\par\smallskip\textbf{Qué aprendimos.} La misma condición ahora está escrita con los términos multiplicados. Ya tenemos el tamaño que buscábamos para x.
\par\smallskip\textbf{Comprueba que lo entiendes.} ¿Podríamos multiplicar solamente el término con x?\par\textit{Respuesta:} No. Eso cambiaría la condición. Debemos multiplicar todos los términos, también el número que está después del igual.
\end{minipage}\par\medskip
\noindent\begin{minipage}{\linewidth}\textbf{10. Restamos las dos ecuaciones preparadas}\\[2pt]Restamos la ecuación 1 preparada a la ecuación 2 preparada. Juntamos cada término solo con otro que tenga la misma letra. Los números sin letra se calculan por separado.
\par\smallskip {\small\renewcommand{\arraystretch}{1.3}\begin{tabular}{>{\raggedright\arraybackslash}p{2.80cm}>{\raggedright\arraybackslash}p{5.38cm}>{\raggedright\arraybackslash}p{3.03cm}}\hline\rowcolor{black!7} \textbf{Qué agrupamos} & \textbf{Cuenta} & \textbf{Qué significa}\\\hline Los términos con x & $20x - \left(20x\right)=0$ & Se cancela esta letra.\\ Los términos con y & $16y - \left(10y\right)=6y$ & Quedan 6y.\\ Los números a la derecha del igual & $720 - \left(600\right)=120$ & Este será el número de la nueva igualdad.\\ \hline\end{tabular}}\par\smallskip
\par\smallskip Primera igualdad que usaremos:
\[20x + 16y = 720\]
\par\smallskip Le restamos esta igualdad:
\[20x + 10y = 600\]
\par\smallskip\textbf{Qué aprendimos.} Los términos con x dan cero. Los términos con letras distintas no se mezclan.
\end{minipage}\par\medskip
\noindent\begin{minipage}{\linewidth}\textbf{11. Mira lo que hemos conseguido}\\[2pt]La nueva ecuación es el resultado de la cuenta anterior. Solo contiene y: después podremos despejarla.
\par\smallskip Escribimos la nueva condición:
\[6y = 120\]
\[\left\{\begin{aligned}4x + 2y &= 120\\6y &= 120\end{aligned}\right.\]
\par\smallskip\textbf{Qué aprendimos.} Conservamos la ecuación 1 y reemplazamos la 2 por esta nueva condición. De esa manera conservamos toda la información necesaria para resolver el sistema.
\end{minipage}\par\medskip
\noindent\begin{minipage}{\linewidth}\textbf{12. Ya tenemos una ecuación con una sola letra}\\[2pt]En esta condición solamente falta conocer y. El número 6 está multiplicando a esa letra. Ahora podemos deshacer esa multiplicación.
\par\smallskip Esta es la condición que vamos a resolver:
\[6y = 120\]
\end{minipage}\par\medskip
\noindent\begin{minipage}{\linewidth}\textbf{13. Calculamos cuánto vale y}\\[2pt]Para deshacer una multiplicación por 6, dividimos AMBOS lados entre 6. En el lado izquierdo queda una sola y. En el derecho hacemos la división.
\par\smallskip Dividimos los dos lados entre el mismo número:
\[\frac{6y}{6}=\frac{120}{6}\]
\par\smallskip Ya podemos leer el valor de la letra:
\[y=\frac{120}{6}=20\]
\par\smallskip\textbf{Qué aprendimos.} Este valor se usará en las ecuaciones donde todavía falten letras.
\end{minipage}\par\medskip
\noindent\begin{minipage}{\linewidth}\textbf{14. Usamos lo que ya sabemos para hallar x}\\[2pt]Ya sabemos y = 20. Sustituir significa escribir ese valor cada vez que aparezca su letra. Usamos paréntesis para que se vea qué número estamos multiplicando.
\par\smallskip\textbf{Qué buscamos.} Dejar solamente x como letra desconocida en esta ecuación.
\par\smallskip Tomamos esta condición que obtuvimos antes:
\[4x + 2y = 120\]
\par\smallskip Reemplazamos las letras conocidas por sus valores:
\[4x + 2\left(20\right)=120\]
\end{minipage}\par\medskip
\noindent\begin{minipage}{\linewidth}\textbf{15. Calculamos las multiplicaciones que ya se pueden hacer}\\[2pt]La letra x sigue sin conocerse. Las otras partes ya tienen un valor o una expresión conocida y podemos efectuar sus productos.
\par\smallskip {\small\renewcommand{\arraystretch}{1.3}\begin{tabular}{>{\raggedright\arraybackslash}p{3.36cm}>{\raggedright\arraybackslash}p{4.48cm}>{\raggedright\arraybackslash}p{3.36cm}}\hline\rowcolor{black!7} \textbf{Parte conocida} & \textbf{La multiplicación} & \textbf{Resultado}\\\hline El término 2y & $2\cdot\left(20\right)$ & $40$\\ \hline\end{tabular}}\par\smallskip
\par\smallskip Al reunir esas partes conocidas, la ecuación queda:
\[4x + 40=120\]
\end{minipage}\par\medskip
\noindent\begin{minipage}{\linewidth}\textbf{16. Dejamos el término con x solo a un lado}\\[2pt]Para quitar la parte conocida, restamos 40 en AMBOS lados. Piensa en una balanza: si hacemos lo mismo a ambos lados, la igualdad se mantiene.
\par\smallskip Hacemos la misma operación a la izquierda y a la derecha:
\[4x + 40 - 40=120 - 40\]
\par\smallskip Las partes que se compensan dan cero; queda:
\[4x=80\]
\par\smallskip\textbf{Qué aprendimos.} «Pasar al otro lado» es una forma abreviada de decir que hicimos la misma suma o resta en los dos lados.
\end{minipage}\par\medskip
\noindent\begin{minipage}{\linewidth}\textbf{17. Calculamos cuánto vale x}\\[2pt]Para deshacer una multiplicación por 4, dividimos AMBOS lados entre 4. En el lado izquierdo queda una sola x. En el derecho hacemos la división.
\par\smallskip Dividimos los dos lados entre el mismo número:
\[\frac{4x}{4}=\frac{80}{4}\]
\par\smallskip Ya podemos leer el valor de la letra:
\[x=\frac{80}{4}=20\]
\par\smallskip\textbf{Qué aprendimos.} Este valor se usará en las ecuaciones donde todavía falten letras.
\end{minipage}\par\medskip
\noindent\begin{minipage}{\linewidth}\textbf{18. Reunimos la respuesta; todavía falta comprobarla}\\[2pt]Escribimos todos los valores en el orden x, y. Estos valores deben funcionar juntos, no cada uno en una ecuación diferente.
\[\left\{\begin{aligned}x&=20\\y&=20\end{aligned}\right.\]
\par\smallskip\textbf{Qué aprendimos.} El siguiente paso es volver a las ecuaciones del enunciado y comprobar cada condición.
\end{minipage}\par\medskip
\noindent\begin{minipage}{\linewidth}\textbf{19. ¿La respuesta cumple la ecuación 1?}\\[2pt]Comprobar significa volver a la condición original y calcularla con nuestra respuesta. Si los dos lados dan el mismo número, esta condición sí se cumple. Deben cumplirse todas.
\par\smallskip La condición del enunciado era:
\[5x + 4y = 180\]
\par\smallskip Ponemos nuestros valores donde estaban las letras:
\[5\left(20\right) + 4\left(20\right)=180\]
\par\smallskip Hacemos las multiplicaciones y luego la suma:
\[100 + 80=180\]
\par\smallskip Comparamos los dos lados:
\[180=180\]
\par\smallskip\textbf{Qué aprendimos.} La igualdad final es verdadera. Nuestra respuesta sí cumple la ecuación 1.
\end{minipage}\par\medskip
\noindent\begin{minipage}{\linewidth}\textbf{20. ¿La respuesta cumple la ecuación 2?}\\[2pt]Comprobar significa volver a la condición original y calcularla con nuestra respuesta. Si los dos lados dan el mismo número, esta condición sí se cumple. Deben cumplirse todas.
\par\smallskip La condición del enunciado era:
\[4x + 2y = 120\]
\par\smallskip Ponemos nuestros valores donde estaban las letras:
\[4\left(20\right) + 2\left(20\right)=120\]
\par\smallskip Hacemos las multiplicaciones y luego la suma:
\[80 + 40=120\]
\par\smallskip Comparamos los dos lados:
\[120=120\]
\par\smallskip\textbf{Qué aprendimos.} La igualdad final es verdadera. Nuestra respuesta sí cumple la ecuación 2.
\end{minipage}\par\medskip
\noindent\begin{minipage}{\linewidth}\textbf{21. Interpretamos y verificamos las unidades}\\[2pt]Se producen 20 toneladas sin azufre y 20 con azufre. Mezcladora: 100 + 80 = 180 min; refinación: 80 + 40 = 120 min.
\end{minipage}\par\medskip
\respuesta{x = 20; y = 20. Se producen 20 toneladas sin azufre y 20 con azufre. Mezcladora: 100 + 80 = 180 min; refinación: 80 + 40 = 120 min.}
\noindent\begin{minipage}{\linewidth}\centering\includegraphics[width=.96\linewidth,height=6.8cm,keepaspectratio]{../recursos/grafico_1_1-23.png}\\{\small Lectura visual complementaria; la demostración está en las operaciones.}\end{minipage}
\end{multicols*}
\newpage\fancyfoot[L]{\small Sección 1.1 | Ejercicio 24 | Paso a paso}\hypertarget{1.1-24}{}
\pdfbookmark[1]{1.1 / 24}{bm1.1-24}
\titulo{24 | Dos tipos de plástico}
\textbf{Enunciado}\par Un fabricante produce dos tipos de plásticos: regular y especial. La producción de cada tonelada de plástico regular requiere dos horas en la planta A y 5 horas en la planta B; para producir cada tonelada de plástico especial se necesitan 2 horas en la planta A y 3 horas en la planta B. Si la planta A está disponible 8 horas diarias y la planta B 15 horas al día, ¿cuántas toneladas de cada tipo de plástico pueden producirse diariamente de modo que ambas plantas se utilicen al máximo de su capacidad?
\[\left\{\begin{aligned}2x + 2y &= 8\\5x + 3y &= 15\end{aligned}\right.\]
\begin{multicols*}{2}
\noindent\begin{minipage}{\linewidth}\textbf{1. Identificamos las incógnitas}\\[2pt]x = toneladas de plástico regular por día; y = toneladas de plástico especial por día.
\par\smallskip {\small\renewcommand{\arraystretch}{1.3}\begin{tabular}{>{\raggedright\arraybackslash}p{2.68cm}>{\raggedright\arraybackslash}p{2.68cm}>{\raggedright\arraybackslash}p{2.68cm}>{\raggedright\arraybackslash}p{2.68cm}}\hline\rowcolor{black!7} \textbf{Producto} & \textbf{Cantidad por día} & \textbf{Horas en A por tonelada} & \textbf{Horas en B por tonelada}\\\hline Plástico regular & x toneladas & 2 & 5\\ Plástico especial & y toneladas & 2 & 3\\ \hline\end{tabular}}\par\smallskip
\end{minipage}\par\medskip
\noindent\begin{minipage}{\linewidth}\textbf{2. Escribimos un balance por planta}\\[2pt]En A, ambos tipos consumen 2 h por tonelada. En B, regular consume 5 h y especial 3 h. Todos los datos ya usan horas.
\par\smallskip\textbf{Por qué.} Si una tonelada de regular ocupa 5 horas en B, x toneladas ocupan 5x horas. La producción especial ocupa 3y horas en B. La suma debe llenar las 15 horas disponibles.
\[2x+2y=8\quad(\mathrm{A})\]
\[5x+3y=15\quad(\mathrm{B})\]
\end{minipage}\par\medskip
\noindent\begin{minipage}{\linewidth}\textbf{3. Primero: ¿qué significa resolver este sistema?}\\[2pt]Las letras x, y representan números que todavía no conocemos. Una solución es una elección de esos números que hace verdaderas TODAS las ecuaciones al mismo tiempo. La llave de la izquierda significa que las condiciones se deben cumplir juntas.
\par\smallskip\textbf{Qué buscamos.} Encontrar una misma pareja de valores.
\par\smallskip El número 3 delante de x indica tres veces ese mismo valor; es decir, sumar tres copias:
\[3x=x+x+x\]
\par\smallskip Ejemplo numérico de la notación: si x valiera 2, tres veces ese valor sería 6. Este 2 es solo un ejemplo; todavía no es la respuesta del ejercicio. El punto centrado indica multiplicación:
\[3\cdot 2=2+2+2=6\]
\par\smallskip\textbf{Comprueba que lo entiendes.} ¿Basta con que los valores funcionen en una sola ecuación?\par\textit{Respuesta:} No. Deben funcionar en todas las ecuaciones del sistema, usando los mismos valores.
\end{minipage}\par\medskip
\noindent\begin{minipage}{\linewidth}\textbf{4. Nuestro plan: ir quitando letras}\\[2pt]Una ecuación con una sola letra es más fácil de resolver. Vamos a combinar ecuaciones para que una letra desaparezca. Cuando sepamos un valor, lo pondremos en las ecuaciones que faltan.
\par\smallskip\textbf{Qué buscamos.} Primero quitar x para calcular y. Después calcular x.
\par\smallskip\textbf{Por qué.} Si una cantidad se resta de sí misma, queda cero. Usaremos esta idea con términos que tengan la misma letra.
\par\smallskip Ejemplo de la idea que utilizaremos:
\[3x-3x=0\]
\par\smallskip\textbf{Qué aprendimos.} Eliminar una letra de una ecuación NO significa que su valor sea cero; significa que sus términos se cancelaron en esa cuenta.
\end{minipage}\par\medskip
\noindent\begin{minipage}{\linewidth}\textbf{5. ¿Cómo haremos desaparecer x?}\\[2pt]Trabajaremos con las ecuaciones 1 y 2. Queremos que los términos con x tengan el mismo tamaño. Así podremos restarlos.
\par\smallskip\textbf{Qué buscamos.} Conseguir términos con 10 veces x en las dos ecuaciones.
\par\smallskip\textbf{Por qué.} Tener 10 veces la misma letra en ambos lados de la resta nos permitirá cancelarla.
\par\smallskip {\small\renewcommand{\arraystretch}{1.3}\begin{tabular}{>{\raggedright\arraybackslash}p{2.58cm}>{\raggedright\arraybackslash}p{2.47cm}>{\raggedright\arraybackslash}p{6.17cm}}\hline\rowcolor{black!7} \textbf{Dónde miramos} & \textbf{Qué término hay} & \textbf{Qué necesitamos hacer}\\\hline Ecuación 1 & $2x$ & Multiplicar TODA esta ecuación por 5\\ Ecuación 2 & $5x$ & Multiplicar TODA esta ecuación por 2\\ \hline\end{tabular}}\par\smallskip
\par\smallskip\textbf{Qué aprendimos.} Estos multiplicadores se eligen para igualar los términos con x. Los demás términos también se multiplican para conservar la igualdad.
\end{minipage}\par\medskip
\noindent\begin{minipage}{\linewidth}\textbf{6. Multiplicamos la ecuación 1 por 5}\\[2pt]Multiplicar una ecuación por 5 significa multiplicar CADA término de la izquierda y también el número de la derecha. La igualdad se conserva porque hacemos lo mismo en ambos lados.
\par\smallskip Esta es la ecuación antes de multiplicar:
\[2x + 2y = 8\]
\par\smallskip {\small\renewcommand{\arraystretch}{1.3}\begin{tabular}{>{\raggedright\arraybackslash}p{2.80cm}>{\raggedright\arraybackslash}p{1.91cm}>{\raggedright\arraybackslash}p{6.50cm}}\hline\rowcolor{black!7} \textbf{Parte de la ecuación} & \textbf{Antes} & \textbf{La cuenta y su resultado}\\\hline Término con x & $2x$ & $5\cdot\left(2x\right)=10x$\\ Término con y & $2y$ & $5\cdot\left(2y\right)=10y$\\ Número a la derecha del igual & $8$ & $5\cdot\left(8\right)=40$\\ \hline\end{tabular}}\par\smallskip
\par\smallskip Después de hacer todas las multiplicaciones queda:
\[10x + 10y = 40\]
\par\smallskip\textbf{Qué aprendimos.} La misma condición ahora está escrita con los términos multiplicados. Ya tenemos el tamaño que buscábamos para x.
\par\smallskip\textbf{Comprueba que lo entiendes.} ¿Podríamos multiplicar solamente el término con x?\par\textit{Respuesta:} No. Eso cambiaría la condición. Debemos multiplicar todos los términos, también el número que está después del igual.
\end{minipage}\par\medskip
\noindent\begin{minipage}{\linewidth}\textbf{7. Multiplicamos la ecuación 2 por 2}\\[2pt]Multiplicar una ecuación por 2 significa multiplicar CADA término de la izquierda y también el número de la derecha. La igualdad se conserva porque hacemos lo mismo en ambos lados.
\par\smallskip Esta es la ecuación antes de multiplicar:
\[5x + 3y = 15\]
\par\smallskip {\small\renewcommand{\arraystretch}{1.3}\begin{tabular}{>{\raggedright\arraybackslash}p{2.80cm}>{\raggedright\arraybackslash}p{1.91cm}>{\raggedright\arraybackslash}p{6.50cm}}\hline\rowcolor{black!7} \textbf{Parte de la ecuación} & \textbf{Antes} & \textbf{La cuenta y su resultado}\\\hline Término con x & $5x$ & $2\cdot\left(5x\right)=10x$\\ Término con y & $3y$ & $2\cdot\left(3y\right)=6y$\\ Número a la derecha del igual & $15$ & $2\cdot\left(15\right)=30$\\ \hline\end{tabular}}\par\smallskip
\par\smallskip Después de hacer todas las multiplicaciones queda:
\[10x + 6y = 30\]
\par\smallskip\textbf{Qué aprendimos.} La misma condición ahora está escrita con los términos multiplicados. Ya tenemos el tamaño que buscábamos para x.
\par\smallskip\textbf{Comprueba que lo entiendes.} ¿Podríamos multiplicar solamente el término con x?\par\textit{Respuesta:} No. Eso cambiaría la condición. Debemos multiplicar todos los términos, también el número que está después del igual.
\end{minipage}\par\medskip
\noindent\begin{minipage}{\linewidth}\textbf{8. Restamos las dos ecuaciones preparadas}\\[2pt]Restamos la ecuación 1 preparada a la ecuación 2 preparada. Juntamos cada término solo con otro que tenga la misma letra. Los números sin letra se calculan por separado.
\par\smallskip {\small\renewcommand{\arraystretch}{1.3}\begin{tabular}{>{\raggedright\arraybackslash}p{2.80cm}>{\raggedright\arraybackslash}p{5.38cm}>{\raggedright\arraybackslash}p{3.03cm}}\hline\rowcolor{black!7} \textbf{Qué agrupamos} & \textbf{Cuenta} & \textbf{Qué significa}\\\hline Los términos con x & $10x - \left(10x\right)=0$ & Se cancela esta letra.\\ Los términos con y & $6y - \left(10y\right)=-4y$ & Quedan -4y.\\ Los números a la derecha del igual & $30 - \left(40\right)=-10$ & Este será el número de la nueva igualdad.\\ \hline\end{tabular}}\par\smallskip
\par\smallskip Primera igualdad que usaremos:
\[10x + 6y = 30\]
\par\smallskip Le restamos esta igualdad:
\[10x + 10y = 40\]
\par\smallskip\textbf{Qué aprendimos.} Los términos con x dan cero. Los términos con letras distintas no se mezclan.
\end{minipage}\par\medskip
\noindent\begin{minipage}{\linewidth}\textbf{9. Mira lo que hemos conseguido}\\[2pt]La nueva ecuación es el resultado de la cuenta anterior. Solo contiene y: después podremos despejarla.
\par\smallskip Escribimos la nueva condición:
\[-4y = -10\]
\[\left\{\begin{aligned}2x + 2y &= 8\\-4y &= -10\end{aligned}\right.\]
\par\smallskip\textbf{Qué aprendimos.} Conservamos la ecuación 1 y reemplazamos la 2 por esta nueva condición. De esa manera conservamos toda la información necesaria para resolver el sistema.
\end{minipage}\par\medskip
\noindent\begin{minipage}{\linewidth}\textbf{10. Ya tenemos una ecuación con una sola letra}\\[2pt]En esta condición solamente falta conocer y. El número -4 está multiplicando a esa letra. Ahora podemos deshacer esa multiplicación.
\par\smallskip Esta es la condición que vamos a resolver:
\[-4y = -10\]
\end{minipage}\par\medskip
\noindent\begin{minipage}{\linewidth}\textbf{11. Calculamos cuánto vale y}\\[2pt]Para deshacer una multiplicación por -4, dividimos AMBOS lados entre -4. En el lado izquierdo queda una sola y. En el derecho hacemos la división.
\par\smallskip Dividimos los dos lados entre el mismo número:
\[\frac{-4y}{-4}=\frac{-10}{-4}\]
\par\smallskip Ya podemos leer el valor de la letra:
\[y=\frac{-10}{-4}=\frac{5}{2}\]
\par\smallskip\textbf{Qué aprendimos.} Un número negativo dividido entre otro negativo da positivo.
\end{minipage}\par\medskip
\noindent\begin{minipage}{\linewidth}\textbf{12. Usamos lo que ya sabemos para hallar x}\\[2pt]Ya sabemos y = 5/2. Sustituir significa escribir ese valor cada vez que aparezca su letra. Usamos paréntesis para que se vea qué número estamos multiplicando.
\par\smallskip\textbf{Qué buscamos.} Dejar solamente x como letra desconocida en esta ecuación.
\par\smallskip Tomamos esta condición que obtuvimos antes:
\[2x + 2y = 8\]
\par\smallskip Reemplazamos las letras conocidas por sus valores:
\[2x + 2\left(\frac{5}{2}\right)=8\]
\end{minipage}\par\medskip
\noindent\begin{minipage}{\linewidth}\textbf{13. Calculamos las multiplicaciones que ya se pueden hacer}\\[2pt]La letra x sigue sin conocerse. Las otras partes ya tienen un valor o una expresión conocida y podemos efectuar sus productos.
\par\smallskip {\small\renewcommand{\arraystretch}{1.3}\begin{tabular}{>{\raggedright\arraybackslash}p{3.36cm}>{\raggedright\arraybackslash}p{4.48cm}>{\raggedright\arraybackslash}p{3.36cm}}\hline\rowcolor{black!7} \textbf{Parte conocida} & \textbf{La multiplicación} & \textbf{Resultado}\\\hline El término 2y & $2\cdot\left(\frac{5}{2}\right)$ & $5$\\ \hline\end{tabular}}\par\smallskip
\par\smallskip Al reunir esas partes conocidas, la ecuación queda:
\[2x + 5=8\]
\end{minipage}\par\medskip
\noindent\begin{minipage}{\linewidth}\textbf{14. Dejamos el término con x solo a un lado}\\[2pt]Para quitar la parte conocida, restamos 5 en AMBOS lados. Piensa en una balanza: si hacemos lo mismo a ambos lados, la igualdad se mantiene.
\par\smallskip Hacemos la misma operación a la izquierda y a la derecha:
\[2x + 5 - 5=8 - 5\]
\par\smallskip Las partes que se compensan dan cero; queda:
\[2x=3\]
\par\smallskip\textbf{Qué aprendimos.} «Pasar al otro lado» es una forma abreviada de decir que hicimos la misma suma o resta en los dos lados.
\end{minipage}\par\medskip
\noindent\begin{minipage}{\linewidth}\textbf{15. Calculamos cuánto vale x}\\[2pt]Para deshacer una multiplicación por 2, dividimos AMBOS lados entre 2. En el lado izquierdo queda una sola x. En el derecho hacemos la división.
\par\smallskip Dividimos los dos lados entre el mismo número:
\[\frac{2x}{2}=\frac{3}{2}\]
\par\smallskip Ya podemos leer el valor de la letra:
\[x=\frac{3}{2}=\frac{3}{2}\]
\par\smallskip\textbf{Qué aprendimos.} Este valor se usará en las ecuaciones donde todavía falten letras.
\end{minipage}\par\medskip
\noindent\begin{minipage}{\linewidth}\textbf{16. Reunimos la respuesta; todavía falta comprobarla}\\[2pt]Escribimos todos los valores en el orden x, y. Estos valores deben funcionar juntos, no cada uno en una ecuación diferente.
\[\left\{\begin{aligned}x&=\frac{3}{2}\\y&=\frac{5}{2}\end{aligned}\right.\]
\par\smallskip\textbf{Qué aprendimos.} El siguiente paso es volver a las ecuaciones del enunciado y comprobar cada condición.
\end{minipage}\par\medskip
\noindent\begin{minipage}{\linewidth}\textbf{17. ¿La respuesta cumple la ecuación 1?}\\[2pt]Comprobar significa volver a la condición original y calcularla con nuestra respuesta. Si los dos lados dan el mismo número, esta condición sí se cumple. Deben cumplirse todas.
\par\smallskip La condición del enunciado era:
\[2x + 2y = 8\]
\par\smallskip Ponemos nuestros valores donde estaban las letras:
\[2\left(\frac{3}{2}\right) + 2\left(\frac{5}{2}\right)=8\]
\par\smallskip Hacemos las multiplicaciones y luego la suma:
\[3 + 5=8\]
\par\smallskip Comparamos los dos lados:
\[8=8\]
\par\smallskip\textbf{Qué aprendimos.} La igualdad final es verdadera. Nuestra respuesta sí cumple la ecuación 1.
\end{minipage}\par\medskip
\noindent\begin{minipage}{\linewidth}\textbf{18. ¿La respuesta cumple la ecuación 2?}\\[2pt]Comprobar significa volver a la condición original y calcularla con nuestra respuesta. Si los dos lados dan el mismo número, esta condición sí se cumple. Deben cumplirse todas.
\par\smallskip La condición del enunciado era:
\[5x + 3y = 15\]
\par\smallskip Ponemos nuestros valores donde estaban las letras:
\[5\left(\frac{3}{2}\right) + 3\left(\frac{5}{2}\right)=15\]
\par\smallskip Hacemos las multiplicaciones y luego la suma:
\[\frac{15}{2} + \frac{15}{2}=15\]
\par\smallskip Comparamos los dos lados:
\[15=15\]
\par\smallskip\textbf{Qué aprendimos.} La igualdad final es verdadera. Nuestra respuesta sí cumple la ecuación 2.
\end{minipage}\par\medskip
\noindent\begin{minipage}{\linewidth}\textbf{19. Interpretamos y verificamos las unidades}\\[2pt]Se producen 1.5 toneladas de regular y 2.5 toneladas de especial al día. En A: 3 + 5 = 8 h. En B: 7.5 + 7.5 = 15 h. Las toneladas pueden ser fraccionarias.
\end{minipage}\par\medskip
\respuesta{x = 3/2; y = 5/2. Se producen 1.5 toneladas de regular y 2.5 toneladas de especial al día. En A: 3 + 5 = 8 h. En B: 7.5 + 7.5 = 15 h. Las toneladas pueden ser fraccionarias.}
\noindent\begin{minipage}{\linewidth}\centering\includegraphics[width=.96\linewidth,height=6.8cm,keepaspectratio]{../recursos/grafico_1_1-24.png}\\{\small Lectura visual complementaria; la demostración está en las operaciones.}\end{minipage}
\end{multicols*}
\newpage\fancyfoot[L]{\small Sección 1.1 | Ejercicio 25 | Paso a paso}\hypertarget{1.1-25}{}
\pdfbookmark[1]{1.1 / 25}{bm1.1-25}
\titulo{25 | Una dieta con tres alimentos}
\textbf{Enunciado}\par Un nutriólogo prepara una dieta que consiste en los alimentos A, B y C. Cada onza del alimento A contiene 2 unidades de proteína, 3 unidades de grasa y 4 unidades de carbohidratos. Cada onza del alimento B contiene 3 unidades de proteínas, 2 unidades de grasa y 1 unidad de carbohidratos. Por su parte, cada onza del alimento C contiene 3 unidades de proteínas, 3 unidades de grasa y 2 unidades de carbohidratos. Si la dieta debe proporcionar exactamente 25 unidades de proteínas, 24 unidades de grasa y 21 unidades de carbohidratos, ¿cuántas onzas de cada tipo de alimento deben utilizarse?
\[\left\{\begin{aligned}2x + 3y + 3z &= 25\\3x + 2y + 3z &= 24\\4x + y + 2z &= 21\end{aligned}\right.\]
\begin{multicols*}{2}
\noindent\begin{minipage}{\linewidth}\textbf{1. Definimos las variables}\\[2pt]x, y, z representan las onzas de alimentos A, B y C, respectivamente. Los nutrientes se suman por separado.
\par\smallskip {\small\renewcommand{\arraystretch}{1.3}\begin{tabular}{>{\raggedright\arraybackslash}p{2.68cm}>{\raggedright\arraybackslash}p{2.68cm}>{\raggedright\arraybackslash}p{2.68cm}>{\raggedright\arraybackslash}p{2.68cm}}\hline\rowcolor{black!7} \textbf{Alimento / cantidad} & \textbf{Proteína por onza} & \textbf{Grasa por onza} & \textbf{Carbohidratos por onza}\\\hline A: x onzas & 2 & 3 & 4\\ B: y onzas & 3 & 2 & 1\\ C: z onzas & 3 & 3 & 2\\ \hline\end{tabular}}\par\smallskip
\end{minipage}\par\medskip
\noindent\begin{minipage}{\linewidth}\textbf{2. Organizamos los tres balances}\\[2pt]Para cada nutriente multiplicamos «lo que aporta una onza» por «las onzas elegidas». Por ejemplo, A aporta 2x unidades de proteína, B aporta 3y y C aporta 3z. La dieta necesita exactamente 25, así que su suma debe ser 25. Repetimos el razonamiento con grasa y carbohidratos.
\par\smallskip\textbf{Qué buscamos.} Escribir una ecuación para proteínas, otra para grasa y otra para carbohidratos. No mezclamos nutrientes distintos en una misma suma.
\[2x+3y+3z=25\quad(\mathrm{proteinas})\]
\[3x+2y+3z=24\quad(\mathrm{grasa})\]
\[4x+y+2z=21\quad(\mathrm{carbohidratos})\]
\end{minipage}\par\medskip
\noindent\begin{minipage}{\linewidth}\textbf{3. Primero: ¿qué significa resolver este sistema?}\\[2pt]Las letras x, y, z representan números que todavía no conocemos. Una solución es una elección de esos números que hace verdaderas TODAS las ecuaciones al mismo tiempo. La llave de la izquierda significa que las condiciones se deben cumplir juntas.
\par\smallskip\textbf{Qué buscamos.} Encontrar un mismo grupo de valores para las 3 letras.
\par\smallskip El número 3 delante de x indica tres veces ese mismo valor; es decir, sumar tres copias:
\[3x=x+x+x\]
\par\smallskip Ejemplo numérico de la notación: si x valiera 2, tres veces ese valor sería 6. Este 2 es solo un ejemplo; todavía no es la respuesta del ejercicio. El punto centrado indica multiplicación:
\[3\cdot 2=2+2+2=6\]
\par\smallskip\textbf{Comprueba que lo entiendes.} ¿Basta con que los valores funcionen en una sola ecuación?\par\textit{Respuesta:} No. Deben funcionar en todas las ecuaciones del sistema, usando los mismos valores.
\end{minipage}\par\medskip
\noindent\begin{minipage}{\linewidth}\textbf{4. Nuestro plan: ir quitando letras}\\[2pt]Una ecuación con una sola letra es más fácil de resolver. Vamos a combinar ecuaciones para que una letra desaparezca. Cuando sepamos un valor, lo pondremos en las ecuaciones que faltan.
\par\smallskip\textbf{Qué buscamos.} Primero quitar x de las otras ecuaciones. Después seguir con las letras restantes.
\par\smallskip\textbf{Por qué.} Si una cantidad se resta de sí misma, queda cero. Usaremos esta idea con términos que tengan la misma letra.
\par\smallskip Ejemplo de la idea que utilizaremos:
\[3x-3x=0\]
\par\smallskip\textbf{Qué aprendimos.} Eliminar una letra de una ecuación NO significa que su valor sea cero; significa que sus términos se cancelaron en esa cuenta.
\end{minipage}\par\medskip
\noindent\begin{minipage}{\linewidth}\textbf{5. ¿Cómo haremos desaparecer x?}\\[2pt]Trabajaremos con las ecuaciones 1 y 2. Queremos que los términos con x tengan el mismo tamaño. Así podremos restarlos.
\par\smallskip\textbf{Qué buscamos.} Conseguir términos con 6 veces x en las dos ecuaciones.
\par\smallskip\textbf{Por qué.} Tener 6 veces la misma letra en ambos lados de la resta nos permitirá cancelarla.
\par\smallskip {\small\renewcommand{\arraystretch}{1.3}\begin{tabular}{>{\raggedright\arraybackslash}p{2.58cm}>{\raggedright\arraybackslash}p{2.47cm}>{\raggedright\arraybackslash}p{6.17cm}}\hline\rowcolor{black!7} \textbf{Dónde miramos} & \textbf{Qué término hay} & \textbf{Qué necesitamos hacer}\\\hline Ecuación 1 & $2x$ & Multiplicar TODA esta ecuación por 3\\ Ecuación 2 & $3x$ & Multiplicar TODA esta ecuación por 2\\ \hline\end{tabular}}\par\smallskip
\par\smallskip\textbf{Qué aprendimos.} Estos multiplicadores se eligen para igualar los términos con x. Los demás términos también se multiplican para conservar la igualdad.
\end{minipage}\par\medskip
\noindent\begin{minipage}{\linewidth}\textbf{6. Multiplicamos la ecuación 1 por 3}\\[2pt]Multiplicar una ecuación por 3 significa multiplicar CADA término de la izquierda y también el número de la derecha. La igualdad se conserva porque hacemos lo mismo en ambos lados.
\par\smallskip Esta es la ecuación antes de multiplicar:
\[2x + 3y + 3z = 25\]
\par\smallskip {\small\renewcommand{\arraystretch}{1.3}\begin{tabular}{>{\raggedright\arraybackslash}p{2.80cm}>{\raggedright\arraybackslash}p{1.91cm}>{\raggedright\arraybackslash}p{6.50cm}}\hline\rowcolor{black!7} \textbf{Parte de la ecuación} & \textbf{Antes} & \textbf{La cuenta y su resultado}\\\hline Término con x & $2x$ & $3\cdot\left(2x\right)=6x$\\ Término con y & $3y$ & $3\cdot\left(3y\right)=9y$\\ Término con z & $3z$ & $3\cdot\left(3z\right)=9z$\\ Número a la derecha del igual & $25$ & $3\cdot\left(25\right)=75$\\ \hline\end{tabular}}\par\smallskip
\par\smallskip Después de hacer todas las multiplicaciones queda:
\[6x + 9y + 9z = 75\]
\par\smallskip\textbf{Qué aprendimos.} La misma condición ahora está escrita con los términos multiplicados. Ya tenemos el tamaño que buscábamos para x.
\par\smallskip\textbf{Comprueba que lo entiendes.} ¿Podríamos multiplicar solamente el término con x?\par\textit{Respuesta:} No. Eso cambiaría la condición. Debemos multiplicar todos los términos, también el número que está después del igual.
\end{minipage}\par\medskip
\noindent\begin{minipage}{\linewidth}\textbf{7. Multiplicamos la ecuación 2 por 2}\\[2pt]Multiplicar una ecuación por 2 significa multiplicar CADA término de la izquierda y también el número de la derecha. La igualdad se conserva porque hacemos lo mismo en ambos lados.
\par\smallskip Esta es la ecuación antes de multiplicar:
\[3x + 2y + 3z = 24\]
\par\smallskip {\small\renewcommand{\arraystretch}{1.3}\begin{tabular}{>{\raggedright\arraybackslash}p{2.80cm}>{\raggedright\arraybackslash}p{1.91cm}>{\raggedright\arraybackslash}p{6.50cm}}\hline\rowcolor{black!7} \textbf{Parte de la ecuación} & \textbf{Antes} & \textbf{La cuenta y su resultado}\\\hline Término con x & $3x$ & $2\cdot\left(3x\right)=6x$\\ Término con y & $2y$ & $2\cdot\left(2y\right)=4y$\\ Término con z & $3z$ & $2\cdot\left(3z\right)=6z$\\ Número a la derecha del igual & $24$ & $2\cdot\left(24\right)=48$\\ \hline\end{tabular}}\par\smallskip
\par\smallskip Después de hacer todas las multiplicaciones queda:
\[6x + 4y + 6z = 48\]
\par\smallskip\textbf{Qué aprendimos.} La misma condición ahora está escrita con los términos multiplicados. Ya tenemos el tamaño que buscábamos para x.
\par\smallskip\textbf{Comprueba que lo entiendes.} ¿Podríamos multiplicar solamente el término con x?\par\textit{Respuesta:} No. Eso cambiaría la condición. Debemos multiplicar todos los términos, también el número que está después del igual.
\end{minipage}\par\medskip
\noindent\begin{minipage}{\linewidth}\textbf{8. Restamos las dos ecuaciones preparadas}\\[2pt]Restamos la ecuación 1 preparada a la ecuación 2 preparada. Juntamos cada término solo con otro que tenga la misma letra. Los números sin letra se calculan por separado.
\par\smallskip {\small\renewcommand{\arraystretch}{1.3}\begin{tabular}{>{\raggedright\arraybackslash}p{2.80cm}>{\raggedright\arraybackslash}p{5.38cm}>{\raggedright\arraybackslash}p{3.03cm}}\hline\rowcolor{black!7} \textbf{Qué agrupamos} & \textbf{Cuenta} & \textbf{Qué significa}\\\hline Los términos con x & $6x - \left(6x\right)=0$ & Se cancela esta letra.\\ Los términos con y & $4y - \left(9y\right)=-5y$ & Quedan -5y.\\ Los términos con z & $6z - \left(9z\right)=-3z$ & Quedan -3z.\\ Los números a la derecha del igual & $48 - \left(75\right)=-27$ & Este será el número de la nueva igualdad.\\ \hline\end{tabular}}\par\smallskip
\par\smallskip Primera igualdad que usaremos:
\[6x + 4y + 6z = 48\]
\par\smallskip Le restamos esta igualdad:
\[6x + 9y + 9z = 75\]
\par\smallskip\textbf{Qué aprendimos.} Los términos con x dan cero. Los términos con letras distintas no se mezclan.
\end{minipage}\par\medskip
\noindent\begin{minipage}{\linewidth}\textbf{9. Mira lo que hemos conseguido}\\[2pt]La nueva ecuación es el resultado de la cuenta anterior. Todavía contiene y, z. Seguiremos reduciendo las condiciones.
\par\smallskip Escribimos la nueva condición:
\[-5y - 3z = -27\]
\[\left\{\begin{aligned}2x + 3y + 3z &= 25\\-5y - 3z &= -27\\4x + y + 2z &= 21\end{aligned}\right.\]
\par\smallskip\textbf{Qué aprendimos.} Conservamos la ecuación 1 y reemplazamos la 2 por esta nueva condición. De esa manera conservamos toda la información necesaria para resolver el sistema.
\end{minipage}\par\medskip
\noindent\begin{minipage}{\linewidth}\textbf{10. ¿Cómo haremos desaparecer x?}\\[2pt]Trabajaremos con las ecuaciones 1 y 3. Queremos que los términos con x tengan el mismo tamaño. Así podremos restarlos.
\par\smallskip\textbf{Qué buscamos.} Conseguir términos con 4 veces x en las dos ecuaciones.
\par\smallskip\textbf{Por qué.} Tener 4 veces la misma letra en ambos lados de la resta nos permitirá cancelarla.
\par\smallskip {\small\renewcommand{\arraystretch}{1.3}\begin{tabular}{>{\raggedright\arraybackslash}p{2.58cm}>{\raggedright\arraybackslash}p{2.47cm}>{\raggedright\arraybackslash}p{6.17cm}}\hline\rowcolor{black!7} \textbf{Dónde miramos} & \textbf{Qué término hay} & \textbf{Qué necesitamos hacer}\\\hline Ecuación 1 & $2x$ & Multiplicar TODA esta ecuación por 2\\ Ecuación 3 & $4x$ & La dejamos igual.\\ \hline\end{tabular}}\par\smallskip
\par\smallskip\textbf{Qué aprendimos.} Estos multiplicadores se eligen para igualar los términos con x. Los demás términos también se multiplican para conservar la igualdad.
\end{minipage}\par\medskip
\noindent\begin{minipage}{\linewidth}\textbf{11. Multiplicamos la ecuación 1 por 2}\\[2pt]Multiplicar una ecuación por 2 significa multiplicar CADA término de la izquierda y también el número de la derecha. La igualdad se conserva porque hacemos lo mismo en ambos lados.
\par\smallskip Esta es la ecuación antes de multiplicar:
\[2x + 3y + 3z = 25\]
\par\smallskip {\small\renewcommand{\arraystretch}{1.3}\begin{tabular}{>{\raggedright\arraybackslash}p{2.80cm}>{\raggedright\arraybackslash}p{1.91cm}>{\raggedright\arraybackslash}p{6.50cm}}\hline\rowcolor{black!7} \textbf{Parte de la ecuación} & \textbf{Antes} & \textbf{La cuenta y su resultado}\\\hline Término con x & $2x$ & $2\cdot\left(2x\right)=4x$\\ Término con y & $3y$ & $2\cdot\left(3y\right)=6y$\\ Término con z & $3z$ & $2\cdot\left(3z\right)=6z$\\ Número a la derecha del igual & $25$ & $2\cdot\left(25\right)=50$\\ \hline\end{tabular}}\par\smallskip
\par\smallskip Después de hacer todas las multiplicaciones queda:
\[4x + 6y + 6z = 50\]
\par\smallskip\textbf{Qué aprendimos.} La misma condición ahora está escrita con los términos multiplicados. Ya tenemos el tamaño que buscábamos para x.
\par\smallskip\textbf{Comprueba que lo entiendes.} ¿Podríamos multiplicar solamente el término con x?\par\textit{Respuesta:} No. Eso cambiaría la condición. Debemos multiplicar todos los términos, también el número que está después del igual.
\end{minipage}\par\medskip
\noindent\begin{minipage}{\linewidth}\textbf{12. Restamos las dos ecuaciones preparadas}\\[2pt]Restamos la ecuación 1 preparada a la ecuación 3 preparada. Juntamos cada término solo con otro que tenga la misma letra. Los números sin letra se calculan por separado.
\par\smallskip {\small\renewcommand{\arraystretch}{1.3}\begin{tabular}{>{\raggedright\arraybackslash}p{2.80cm}>{\raggedright\arraybackslash}p{5.38cm}>{\raggedright\arraybackslash}p{3.03cm}}\hline\rowcolor{black!7} \textbf{Qué agrupamos} & \textbf{Cuenta} & \textbf{Qué significa}\\\hline Los términos con x & $4x - \left(4x\right)=0$ & Se cancela esta letra.\\ Los términos con y & $y - \left(6y\right)=-5y$ & Quedan -5y.\\ Los términos con z & $2z - \left(6z\right)=-4z$ & Quedan -4z.\\ Los números a la derecha del igual & $21 - \left(50\right)=-29$ & Este será el número de la nueva igualdad.\\ \hline\end{tabular}}\par\smallskip
\par\smallskip Primera igualdad que usaremos:
\[4x + y + 2z = 21\]
\par\smallskip Le restamos esta igualdad:
\[4x + 6y + 6z = 50\]
\par\smallskip\textbf{Qué aprendimos.} Los términos con x dan cero. Los términos con letras distintas no se mezclan.
\end{minipage}\par\medskip
\noindent\begin{minipage}{\linewidth}\textbf{13. Mira lo que hemos conseguido}\\[2pt]La nueva ecuación es el resultado de la cuenta anterior. Todavía contiene y, z. Seguiremos reduciendo las condiciones.
\par\smallskip Escribimos la nueva condición:
\[-5y - 4z = -29\]
\[\left\{\begin{aligned}2x + 3y + 3z &= 25\\-5y - 3z &= -27\\-5y - 4z &= -29\end{aligned}\right.\]
\par\smallskip\textbf{Qué aprendimos.} Conservamos la ecuación 1 y reemplazamos la 3 por esta nueva condición. De esa manera conservamos toda la información necesaria para resolver el sistema.
\end{minipage}\par\medskip
\noindent\begin{minipage}{\linewidth}\textbf{14. ¿Cómo haremos desaparecer y?}\\[2pt]Trabajaremos con las ecuaciones 2 y 3. Queremos que los términos con y tengan el mismo tamaño. Así podremos restarlos.
\par\smallskip\textbf{Qué buscamos.} Conseguir términos con 5 veces y en las dos ecuaciones.
\par\smallskip\textbf{Por qué.} Tener 5 veces la misma letra en ambos lados de la resta nos permitirá cancelarla.
\par\smallskip {\small\renewcommand{\arraystretch}{1.3}\begin{tabular}{>{\raggedright\arraybackslash}p{2.58cm}>{\raggedright\arraybackslash}p{2.47cm}>{\raggedright\arraybackslash}p{6.17cm}}\hline\rowcolor{black!7} \textbf{Dónde miramos} & \textbf{Qué término hay} & \textbf{Qué necesitamos hacer}\\\hline Ecuación 2 & $-5y$ & La dejamos igual.\\ Ecuación 3 & $-5y$ & La dejamos igual.\\ \hline\end{tabular}}\par\smallskip
\par\smallskip\textbf{Qué aprendimos.} Estos multiplicadores se eligen para igualar los términos con y. Los demás términos también se multiplican para conservar la igualdad.
\end{minipage}\par\medskip
\noindent\begin{minipage}{\linewidth}\textbf{15. Restamos las dos ecuaciones preparadas}\\[2pt]Restamos la ecuación 2 preparada a la ecuación 3 preparada. Juntamos cada término solo con otro que tenga la misma letra. Los números sin letra se calculan por separado.
\par\smallskip {\small\renewcommand{\arraystretch}{1.3}\begin{tabular}{>{\raggedright\arraybackslash}p{2.80cm}>{\raggedright\arraybackslash}p{5.38cm}>{\raggedright\arraybackslash}p{3.03cm}}\hline\rowcolor{black!7} \textbf{Qué agrupamos} & \textbf{Cuenta} & \textbf{Qué significa}\\\hline Los términos con y & $-5y - \left(-5y\right)=0$ & Se cancela esta letra.\\ Los términos con z & $-4z - \left(-3z\right)=-z$ & Quedan -z.\\ Los números a la derecha del igual & $-29 - \left(-27\right)=-29+27=-2$ & Este será el número de la nueva igualdad.\\ \hline\end{tabular}}\par\smallskip
\par\smallskip Primera igualdad que usaremos:
\[-5y - 4z = -29\]
\par\smallskip Le restamos esta igualdad:
\[-5y - 3z = -27\]
\par\smallskip\textbf{Qué aprendimos.} Los términos con y dan cero. Recuerda: restar un número negativo equivale a sumar su opuesto.
\end{minipage}\par\medskip
\noindent\begin{minipage}{\linewidth}\textbf{16. Mira lo que hemos conseguido}\\[2pt]La nueva ecuación es el resultado de la cuenta anterior. Solo contiene z: después podremos despejarla.
\par\smallskip Escribimos la nueva condición:
\[-z = -2\]
\[\left\{\begin{aligned}2x + 3y + 3z &= 25\\-5y - 3z &= -27\\-z &= -2\end{aligned}\right.\]
\par\smallskip\textbf{Qué aprendimos.} Conservamos la ecuación 2 y reemplazamos la 3 por esta nueva condición. De esa manera conservamos toda la información necesaria para resolver el sistema.
\end{minipage}\par\medskip
\noindent\begin{minipage}{\linewidth}\textbf{17. Ya tenemos una ecuación con una sola letra}\\[2pt]En esta condición solamente falta conocer z. El número -1 está multiplicando a esa letra. Ahora podemos deshacer esa multiplicación.
\par\smallskip Esta es la condición que vamos a resolver:
\[-z = -2\]
\end{minipage}\par\medskip
\noindent\begin{minipage}{\linewidth}\textbf{18. Calculamos cuánto vale z}\\[2pt]Para deshacer una multiplicación por -1, dividimos AMBOS lados entre -1. En el lado izquierdo queda una sola z. En el derecho hacemos la división.
\par\smallskip Dividimos los dos lados entre el mismo número:
\[\frac{-z}{-1}=\frac{-2}{-1}\]
\par\smallskip Ya podemos leer el valor de la letra:
\[z=\frac{-2}{-1}=2\]
\par\smallskip\textbf{Qué aprendimos.} Un número negativo dividido entre otro negativo da positivo.
\end{minipage}\par\medskip
\noindent\begin{minipage}{\linewidth}\textbf{19. Usamos lo que ya sabemos para hallar y}\\[2pt]Ya sabemos z = 2. Sustituir significa escribir ese valor cada vez que aparezca su letra. Usamos paréntesis para que se vea qué número estamos multiplicando.
\par\smallskip\textbf{Qué buscamos.} Dejar solamente y como letra desconocida en esta ecuación.
\par\smallskip Tomamos esta condición que obtuvimos antes:
\[-5y - 3z = -27\]
\par\smallskip Reemplazamos las letras conocidas por sus valores:
\[-5y - 3\left(2\right)=-27\]
\end{minipage}\par\medskip
\noindent\begin{minipage}{\linewidth}\textbf{20. Calculamos las multiplicaciones que ya se pueden hacer}\\[2pt]La letra y sigue sin conocerse. Las otras partes ya tienen un valor o una expresión conocida y podemos efectuar sus productos.
\par\smallskip {\small\renewcommand{\arraystretch}{1.3}\begin{tabular}{>{\raggedright\arraybackslash}p{3.36cm}>{\raggedright\arraybackslash}p{4.48cm}>{\raggedright\arraybackslash}p{3.36cm}}\hline\rowcolor{black!7} \textbf{Parte conocida} & \textbf{La multiplicación} & \textbf{Resultado}\\\hline El término -3z & $-3\cdot\left(2\right)$ & $-6$\\ \hline\end{tabular}}\par\smallskip
\par\smallskip Al reunir esas partes conocidas, la ecuación queda:
\[-5y - 6=-27\]
\end{minipage}\par\medskip
\noindent\begin{minipage}{\linewidth}\textbf{21. Dejamos el término con y solo a un lado}\\[2pt]Para quitar la parte conocida, sumamos 6 en AMBOS lados. Piensa en una balanza: si hacemos lo mismo a ambos lados, la igualdad se mantiene.
\par\smallskip Hacemos la misma operación a la izquierda y a la derecha:
\[-5y - 6 + 6=-27 + 6\]
\par\smallskip Las partes que se compensan dan cero; queda:
\[-5y=-21\]
\par\smallskip\textbf{Qué aprendimos.} «Pasar al otro lado» es una forma abreviada de decir que hicimos la misma suma o resta en los dos lados.
\end{minipage}\par\medskip
\noindent\begin{minipage}{\linewidth}\textbf{22. Calculamos cuánto vale y}\\[2pt]Para deshacer una multiplicación por -5, dividimos AMBOS lados entre -5. En el lado izquierdo queda una sola y. En el derecho hacemos la división.
\par\smallskip Dividimos los dos lados entre el mismo número:
\[\frac{-5y}{-5}=\frac{-21}{-5}\]
\par\smallskip Ya podemos leer el valor de la letra:
\[y=\frac{-21}{-5}=\frac{21}{5}\]
\par\smallskip\textbf{Qué aprendimos.} Un número negativo dividido entre otro negativo da positivo.
\end{minipage}\par\medskip
\noindent\begin{minipage}{\linewidth}\textbf{23. Usamos lo que ya sabemos para hallar x}\\[2pt]Ya sabemos y = 21/5; z = 2. Sustituir significa escribir ese valor cada vez que aparezca su letra. Usamos paréntesis para que se vea qué número estamos multiplicando.
\par\smallskip\textbf{Qué buscamos.} Dejar solamente x como letra desconocida en esta ecuación.
\par\smallskip Tomamos esta condición que obtuvimos antes:
\[2x + 3y + 3z = 25\]
\par\smallskip Reemplazamos las letras conocidas por sus valores:
\[2x + 3\left(\frac{21}{5}\right) + 3\left(2\right)=25\]
\end{minipage}\par\medskip
\noindent\begin{minipage}{\linewidth}\textbf{24. Calculamos las multiplicaciones que ya se pueden hacer}\\[2pt]La letra x sigue sin conocerse. Las otras partes ya tienen un valor o una expresión conocida y podemos efectuar sus productos.
\par\smallskip {\small\renewcommand{\arraystretch}{1.3}\begin{tabular}{>{\raggedright\arraybackslash}p{3.36cm}>{\raggedright\arraybackslash}p{4.48cm}>{\raggedright\arraybackslash}p{3.36cm}}\hline\rowcolor{black!7} \textbf{Parte conocida} & \textbf{La multiplicación} & \textbf{Resultado}\\\hline El término 3y & $3\cdot\left(\frac{21}{5}\right)$ & $\frac{63}{5}$\\ El término 3z & $3\cdot\left(2\right)$ & $6$\\ \hline\end{tabular}}\par\smallskip
\par\smallskip Al reunir esas partes conocidas, la ecuación queda:
\[2x + \frac{93}{5}=25\]
\end{minipage}\par\medskip
\noindent\begin{minipage}{\linewidth}\textbf{25. Dejamos el término con x solo a un lado}\\[2pt]Para quitar la parte conocida, restamos 93/5 en AMBOS lados. Piensa en una balanza: si hacemos lo mismo a ambos lados, la igualdad se mantiene.
\par\smallskip Hacemos la misma operación a la izquierda y a la derecha:
\[2x + \frac{93}{5} - \frac{93}{5}=25 - \frac{93}{5}\]
\par\smallskip Las partes que se compensan dan cero; queda:
\[2x=\frac{32}{5}\]
\par\smallskip\textbf{Qué aprendimos.} «Pasar al otro lado» es una forma abreviada de decir que hicimos la misma suma o resta en los dos lados.
\end{minipage}\par\medskip
\noindent\begin{minipage}{\linewidth}\textbf{26. Calculamos cuánto vale x}\\[2pt]Para deshacer una multiplicación por 2, dividimos AMBOS lados entre 2. En el lado izquierdo queda una sola x. En el derecho hacemos la división.
\par\smallskip Dividimos los dos lados entre el mismo número:
\[\frac{2x}{2}=\frac{\frac{32}{5}}{2}\]
\par\smallskip Ya podemos leer el valor de la letra:
\[x=\frac{\frac{32}{5}}{2}=\frac{16}{5}\]
\par\smallskip\textbf{Qué aprendimos.} Este valor se usará en las ecuaciones donde todavía falten letras.
\end{minipage}\par\medskip
\noindent\begin{minipage}{\linewidth}\textbf{27. Reunimos la respuesta; todavía falta comprobarla}\\[2pt]Escribimos todos los valores en el orden x, y, z. Estos valores deben funcionar juntos, no cada uno en una ecuación diferente.
\[\left\{\begin{aligned}x&=\frac{16}{5}\\y&=\frac{21}{5}\\z&=2\end{aligned}\right.\]
\par\smallskip\textbf{Qué aprendimos.} El siguiente paso es volver a las ecuaciones del enunciado y comprobar cada condición.
\end{minipage}\par\medskip
\noindent\begin{minipage}{\linewidth}\textbf{28. ¿La respuesta cumple la ecuación 1?}\\[2pt]Comprobar significa volver a la condición original y calcularla con nuestra respuesta. Si los dos lados dan el mismo número, esta condición sí se cumple. Deben cumplirse todas.
\par\smallskip La condición del enunciado era:
\[2x + 3y + 3z = 25\]
\par\smallskip Ponemos nuestros valores donde estaban las letras:
\[2\left(\frac{16}{5}\right) + 3\left(\frac{21}{5}\right) + 3\left(2\right)=25\]
\par\smallskip Hacemos las multiplicaciones y luego la suma:
\[\frac{32}{5} + \frac{63}{5} + 6=25\]
\par\smallskip Comparamos los dos lados:
\[25=25\]
\par\smallskip\textbf{Qué aprendimos.} La igualdad final es verdadera. Nuestra respuesta sí cumple la ecuación 1.
\end{minipage}\par\medskip
\noindent\begin{minipage}{\linewidth}\textbf{29. ¿La respuesta cumple la ecuación 2?}\\[2pt]Comprobar significa volver a la condición original y calcularla con nuestra respuesta. Si los dos lados dan el mismo número, esta condición sí se cumple. Deben cumplirse todas.
\par\smallskip La condición del enunciado era:
\[3x + 2y + 3z = 24\]
\par\smallskip Ponemos nuestros valores donde estaban las letras:
\[3\left(\frac{16}{5}\right) + 2\left(\frac{21}{5}\right) + 3\left(2\right)=24\]
\par\smallskip Hacemos las multiplicaciones y luego la suma:
\[\frac{48}{5} + \frac{42}{5} + 6=24\]
\par\smallskip Comparamos los dos lados:
\[24=24\]
\par\smallskip\textbf{Qué aprendimos.} La igualdad final es verdadera. Nuestra respuesta sí cumple la ecuación 2.
\end{minipage}\par\medskip
\noindent\begin{minipage}{\linewidth}\textbf{30. ¿La respuesta cumple la ecuación 3?}\\[2pt]Comprobar significa volver a la condición original y calcularla con nuestra respuesta. Si los dos lados dan el mismo número, esta condición sí se cumple. Deben cumplirse todas.
\par\smallskip La condición del enunciado era:
\[4x + y + 2z = 21\]
\par\smallskip Ponemos nuestros valores donde estaban las letras:
\[4\left(\frac{16}{5}\right) + \left(\frac{21}{5}\right) + 2\left(2\right)=21\]
\par\smallskip Hacemos las multiplicaciones y luego la suma:
\[\frac{64}{5} + \frac{21}{5} + 4=21\]
\par\smallskip Comparamos los dos lados:
\[21=21\]
\par\smallskip\textbf{Qué aprendimos.} La igualdad final es verdadera. Nuestra respuesta sí cumple la ecuación 3.
\end{minipage}\par\medskip
\noindent\begin{minipage}{\linewidth}\textbf{31. Interpretamos y verificamos las unidades}\\[2pt]Se utilizan 16/5 onzas de A, 21/5 onzas de B y 2 onzas de C. Las tres cantidades son no negativas y los totales de nutrientes son exactamente 25, 24 y 21.
\end{minipage}\par\medskip
\respuesta{Alimento A: 16/5 = 3.2 onzas. Alimento B: 21/5 = 4.2 onzas. Alimento C: 2 onzas.}
\noindent\begin{minipage}{\linewidth}\centering\includegraphics[width=.96\linewidth,height=6.8cm,keepaspectratio]{../recursos/grafico_1_1-25.png}\\{\small Lectura visual complementaria; la demostración está en las operaciones.}\end{minipage}
\end{multicols*}
\newpage\fancyfoot[L]{\small Sección 1.1 | Ejercicio 26 | Paso a paso}\hypertarget{1.1-26}{}
\pdfbookmark[1]{1.1 / 26}{bm1.1-26}
\titulo{26 | Reveladores: una familia de planes de producción}
\textbf{Enunciado}\par Un fabricante produce reveladores de película de 2, 6 y 9 minutos. La fabricación de cada tonelada del revelador de 2 minutos requiere 6 minutos en la planta A y 24 minutos en la planta B. Para manufacturar cada tonelada del revelador de 6 minutos son necesarios 12 minutos en la planta A y 12 minutos en la planta B. Por último, para producir cada tonelada del revelador de 9 minutos se utiliza 12 minutos en la planta A y 12 minutos en la planta B. Si la planta A está disponible 10 horas al día y la planta B 16 horas diarias, ¿cuántas toneladas de cada tipo de revelador de película pueden producirse de modo que las plantas operen a toda su capacidad?
\[\left\{\begin{aligned}6x + 12y + 12z &= 600\\24x + 12y + 12z &= 960\end{aligned}\right.\]
\begin{multicols*}{2}
\noindent\begin{minipage}{\linewidth}\textbf{1. Variables y unidades}\\[2pt]x, y, z son toneladas por día de los reveladores de 2, 6 y 9 minutos. Esos nombres identifican productos; sus consumos de planta son los datos de cada coeficiente.
\par\smallskip {\small\renewcommand{\arraystretch}{1.3}\begin{tabular}{>{\raggedright\arraybackslash}p{3.74cm}>{\raggedright\arraybackslash}p{3.74cm}>{\raggedright\arraybackslash}p{3.74cm}}\hline\rowcolor{black!7} \textbf{Revelador / cantidad} & \textbf{Minutos en A por tonelada} & \textbf{Minutos en B por tonelada}\\\hline De 2 minutos: x toneladas & 6 & 24\\ De 6 minutos: y toneladas & 12 & 12\\ De 9 minutos: z toneladas & 12 & 12\\ \hline\end{tabular}}\par\smallskip
\par\smallskip\textbf{Qué aprendimos.} «De 2, 6 y 9 minutos» son los nombres de los productos. Para calcular el uso de las plantas usamos los tiempos de la tabla, no esos nombres.
\end{minipage}\par\medskip
\noindent\begin{minipage}{\linewidth}\textbf{2. Convertimos la disponibilidad a minutos}\\[2pt]La planta A aporta 600 min y la B aporta 960 min al día.
\[10(60)=600\]
\[16(60)=960\]
\end{minipage}\par\medskip
\noindent\begin{minipage}{\linewidth}\textbf{3. Planteamos los dos balances}\\[2pt]Los reveladores de 6 y 9 minutos tienen el mismo consumo en A y en B. Por eso podemos cambiar uno por otro sin alterar los totales.
\par\smallskip\textbf{Por qué.} El consumo total de una planta se obtiene sumando el tiempo que necesita cada producto. Los productos de 6 y 9 minutos tienen consumos iguales: esto permitirá varias mezclas posibles.
\[6x+12y+12z=600\]
\[24x+12y+12z=960\]
\end{minipage}\par\medskip
\noindent\begin{minipage}{\linewidth}\textbf{4. Primero: ¿qué significa resolver este sistema?}\\[2pt]Las letras x, y, z representan números que todavía no conocemos. Una solución es una elección de esos números que hace verdaderas TODAS las ecuaciones al mismo tiempo. La llave de la izquierda significa que las condiciones se deben cumplir juntas.
\par\smallskip\textbf{Qué buscamos.} Encontrar un mismo grupo de valores para las 3 letras.
\par\smallskip El número 3 delante de x indica tres veces ese mismo valor; es decir, sumar tres copias:
\[3x=x+x+x\]
\par\smallskip Ejemplo numérico de la notación: si x valiera 2, tres veces ese valor sería 6. Este 2 es solo un ejemplo; todavía no es la respuesta del ejercicio. El punto centrado indica multiplicación:
\[3\cdot 2=2+2+2=6\]
\par\smallskip\textbf{Comprueba que lo entiendes.} ¿Basta con que los valores funcionen en una sola ecuación?\par\textit{Respuesta:} No. Deben funcionar en todas las ecuaciones del sistema, usando los mismos valores.
\end{minipage}\par\medskip
\noindent\begin{minipage}{\linewidth}\textbf{5. Nuestro plan: ir quitando letras}\\[2pt]Una ecuación con una sola letra es más fácil de resolver. Vamos a combinar ecuaciones para que una letra desaparezca. Cuando sepamos un valor, lo pondremos en las ecuaciones que faltan.
\par\smallskip\textbf{Qué buscamos.} Primero quitar x de las otras ecuaciones. Después seguir con las letras restantes.
\par\smallskip\textbf{Por qué.} Si una cantidad se resta de sí misma, queda cero. Usaremos esta idea con términos que tengan la misma letra.
\par\smallskip Ejemplo de la idea que utilizaremos:
\[3x-3x=0\]
\par\smallskip\textbf{Qué aprendimos.} Eliminar una letra de una ecuación NO significa que su valor sea cero; significa que sus términos se cancelaron en esa cuenta.
\end{minipage}\par\medskip
\noindent\begin{minipage}{\linewidth}\textbf{6. ¿Cómo haremos desaparecer x?}\\[2pt]Trabajaremos con las ecuaciones 1 y 2. Queremos que los términos con x tengan el mismo tamaño. Así podremos restarlos.
\par\smallskip\textbf{Qué buscamos.} Conseguir términos con 24 veces x en las dos ecuaciones.
\par\smallskip\textbf{Por qué.} Tener 24 veces la misma letra en ambos lados de la resta nos permitirá cancelarla.
\par\smallskip {\small\renewcommand{\arraystretch}{1.3}\begin{tabular}{>{\raggedright\arraybackslash}p{2.58cm}>{\raggedright\arraybackslash}p{2.47cm}>{\raggedright\arraybackslash}p{6.17cm}}\hline\rowcolor{black!7} \textbf{Dónde miramos} & \textbf{Qué término hay} & \textbf{Qué necesitamos hacer}\\\hline Ecuación 1 & $6x$ & Multiplicar TODA esta ecuación por 4\\ Ecuación 2 & $24x$ & La dejamos igual.\\ \hline\end{tabular}}\par\smallskip
\par\smallskip\textbf{Qué aprendimos.} Estos multiplicadores se eligen para igualar los términos con x. Los demás términos también se multiplican para conservar la igualdad.
\end{minipage}\par\medskip
\noindent\begin{minipage}{\linewidth}\textbf{7. Multiplicamos la ecuación 1 por 4}\\[2pt]Multiplicar una ecuación por 4 significa multiplicar CADA término de la izquierda y también el número de la derecha. La igualdad se conserva porque hacemos lo mismo en ambos lados.
\par\smallskip Esta es la ecuación antes de multiplicar:
\[6x + 12y + 12z = 600\]
\par\smallskip {\small\renewcommand{\arraystretch}{1.3}\begin{tabular}{>{\raggedright\arraybackslash}p{2.80cm}>{\raggedright\arraybackslash}p{1.91cm}>{\raggedright\arraybackslash}p{6.50cm}}\hline\rowcolor{black!7} \textbf{Parte de la ecuación} & \textbf{Antes} & \textbf{La cuenta y su resultado}\\\hline Término con x & $6x$ & $4\cdot\left(6x\right)=24x$\\ Término con y & $12y$ & $4\cdot\left(12y\right)=48y$\\ Término con z & $12z$ & $4\cdot\left(12z\right)=48z$\\ Número a la derecha del igual & $600$ & $4\cdot\left(600\right)=2400$\\ \hline\end{tabular}}\par\smallskip
\par\smallskip Después de hacer todas las multiplicaciones queda:
\[24x + 48y + 48z = 2400\]
\par\smallskip\textbf{Qué aprendimos.} La misma condición ahora está escrita con los términos multiplicados. Ya tenemos el tamaño que buscábamos para x.
\par\smallskip\textbf{Comprueba que lo entiendes.} ¿Podríamos multiplicar solamente el término con x?\par\textit{Respuesta:} No. Eso cambiaría la condición. Debemos multiplicar todos los términos, también el número que está después del igual.
\end{minipage}\par\medskip
\noindent\begin{minipage}{\linewidth}\textbf{8. Restamos las dos ecuaciones preparadas}\\[2pt]Restamos la ecuación 1 preparada a la ecuación 2 preparada. Juntamos cada término solo con otro que tenga la misma letra. Los números sin letra se calculan por separado.
\par\smallskip {\small\renewcommand{\arraystretch}{1.3}\begin{tabular}{>{\raggedright\arraybackslash}p{2.80cm}>{\raggedright\arraybackslash}p{5.38cm}>{\raggedright\arraybackslash}p{3.03cm}}\hline\rowcolor{black!7} \textbf{Qué agrupamos} & \textbf{Cuenta} & \textbf{Qué significa}\\\hline Los términos con x & $24x - \left(24x\right)=0$ & Se cancela esta letra.\\ Los términos con y & $12y - \left(48y\right)=-36y$ & Quedan -36y.\\ Los términos con z & $12z - \left(48z\right)=-36z$ & Quedan -36z.\\ Los números a la derecha del igual & $960 - \left(2400\right)=-1440$ & Este será el número de la nueva igualdad.\\ \hline\end{tabular}}\par\smallskip
\par\smallskip Primera igualdad que usaremos:
\[24x + 12y + 12z = 960\]
\par\smallskip Le restamos esta igualdad:
\[24x + 48y + 48z = 2400\]
\par\smallskip\textbf{Qué aprendimos.} Los términos con x dan cero. Los términos con letras distintas no se mezclan.
\end{minipage}\par\medskip
\noindent\begin{minipage}{\linewidth}\textbf{9. Mira lo que hemos conseguido}\\[2pt]La nueva ecuación es el resultado de la cuenta anterior. Todavía contiene y, z. Seguiremos reduciendo las condiciones.
\par\smallskip Escribimos la nueva condición:
\[-36y - 36z = -1440\]
\[\left\{\begin{aligned}6x + 12y + 12z &= 600\\-36y - 36z &= -1440\end{aligned}\right.\]
\par\smallskip\textbf{Qué aprendimos.} Conservamos la ecuación 1 y reemplazamos la 2 por esta nueva condición. De esa manera conservamos toda la información necesaria para resolver el sistema.
\end{minipage}\par\medskip
\noindent\begin{minipage}{\linewidth}\textbf{10. Escribimos la misma condición con números más pequeños}\\[2pt]Todos los números de esta ecuación se pueden dividir entre 36 sin dejar restos. Dividimos ambos lados entre ese mismo número. Así la condición no cambia y las siguientes cuentas serán más sencillas.
\par\smallskip Antes de simplificar:
\[-36y - 36z = -1440\]
\par\smallskip Después de dividir todos los términos entre 36:
\[-y - z = -40\]
\end{minipage}\par\medskip
\noindent\begin{minipage}{\linewidth}\textbf{11. Una letra puede elegirse libremente}\\[2pt]Las condiciones que quedan no fijan un único valor para z. Podemos elegir ese valor y calcular los demás. Para representar todas las elecciones usamos otra letra, llamada parámetro.
\par\smallskip La letra z será el número que elijamos. Lo llamaremos t:
\[z=t\]
\par\smallskip\textbf{Qué aprendimos.} Un parámetro no es una respuesta que falte: es una forma de describir muchas respuestas a la vez. Primero lo dejamos libre; en un problema real también revisaremos si hay límites como cantidades no negativas.
\end{minipage}\par\medskip
\noindent\begin{minipage}{\linewidth}\textbf{12. Usamos lo que ya sabemos para hallar y}\\[2pt]Ya sabemos z = t. Sustituir significa escribir ese valor cada vez que aparezca su letra. Usamos paréntesis para que se vea qué número estamos multiplicando.
\par\smallskip\textbf{Qué buscamos.} Dejar solamente y como letra desconocida en esta ecuación.
\par\smallskip Tomamos esta condición que obtuvimos antes:
\[-y - z = -40\]
\par\smallskip Reemplazamos las letras conocidas por sus valores:
\[-y - \left(t\right)=-40\]
\end{minipage}\par\medskip
\noindent\begin{minipage}{\linewidth}\textbf{13. Calculamos las multiplicaciones que ya se pueden hacer}\\[2pt]La letra y sigue sin conocerse. Las otras partes ya tienen un valor o una expresión conocida y podemos efectuar sus productos.
\par\smallskip {\small\renewcommand{\arraystretch}{1.3}\begin{tabular}{>{\raggedright\arraybackslash}p{3.36cm}>{\raggedright\arraybackslash}p{4.48cm}>{\raggedright\arraybackslash}p{3.36cm}}\hline\rowcolor{black!7} \textbf{Parte conocida} & \textbf{La multiplicación} & \textbf{Resultado}\\\hline El término -z & $-1\cdot\left(t\right)$ & $- t$\\ \hline\end{tabular}}\par\smallskip
\par\smallskip Al reunir esas partes conocidas, la ecuación queda:
\[-y + - t=-40\]
\end{minipage}\par\medskip
\noindent\begin{minipage}{\linewidth}\textbf{14. Dejamos el término con y solo a un lado}\\[2pt]Para quitar la parte conocida, quitamos la expresión -t en AMBOS lados. Piensa en una balanza: si hacemos lo mismo a ambos lados, la igualdad se mantiene.
\par\smallskip Hacemos la misma operación a la izquierda y a la derecha:
\[-y=-40-\left(- t\right)\]
\par\smallskip Las partes que se compensan dan cero; queda:
\[-y=t - 40\]
\par\smallskip\textbf{Qué aprendimos.} «Pasar al otro lado» es una forma abreviada de decir que hicimos la misma suma o resta en los dos lados.
\end{minipage}\par\medskip
\noindent\begin{minipage}{\linewidth}\textbf{15. Calculamos cuánto vale y}\\[2pt]Para deshacer una multiplicación por -1, dividimos AMBOS lados entre -1. En el lado izquierdo queda una sola y. En el derecho hacemos la división.
\par\smallskip Dividimos los dos lados entre el mismo número:
\[\frac{-y}{-1}=\frac{t - 40}{-1}\]
\par\smallskip Ya podemos leer el valor de la letra:
\[y=\frac{t - 40}{-1}=40 - t\]
\par\smallskip\textbf{Qué aprendimos.} Este valor se usará en las ecuaciones donde todavía falten letras.
\end{minipage}\par\medskip
\noindent\begin{minipage}{\linewidth}\textbf{16. Usamos lo que ya sabemos para hallar x}\\[2pt]Ya sabemos y = 40 - t; z = t. Sustituir significa escribir ese valor cada vez que aparezca su letra. Usamos paréntesis para que se vea qué número estamos multiplicando.
\par\smallskip\textbf{Qué buscamos.} Dejar solamente x como letra desconocida en esta ecuación.
\par\smallskip Tomamos esta condición que obtuvimos antes:
\[6x + 12y + 12z = 600\]
\par\smallskip Reemplazamos las letras conocidas por sus valores:
\[6x + 12\left(40 - t\right) + 12\left(t\right)=600\]
\end{minipage}\par\medskip
\noindent\begin{minipage}{\linewidth}\textbf{17. Calculamos las multiplicaciones que ya se pueden hacer}\\[2pt]La letra x sigue sin conocerse. Las otras partes ya tienen un valor o una expresión conocida y podemos efectuar sus productos.
\par\smallskip {\small\renewcommand{\arraystretch}{1.3}\begin{tabular}{>{\raggedright\arraybackslash}p{3.36cm}>{\raggedright\arraybackslash}p{4.48cm}>{\raggedright\arraybackslash}p{3.36cm}}\hline\rowcolor{black!7} \textbf{Parte conocida} & \textbf{La multiplicación} & \textbf{Resultado}\\\hline El término 12y & $12\cdot\left(40 - t\right)$ & $480 - 12 t$\\ El término 12z & $12\cdot\left(t\right)$ & $12 t$\\ \hline\end{tabular}}\par\smallskip
\par\smallskip Al reunir esas partes conocidas, la ecuación queda:
\[6x + 480=600\]
\end{minipage}\par\medskip
\noindent\begin{minipage}{\linewidth}\textbf{18. Dejamos el término con x solo a un lado}\\[2pt]Para quitar la parte conocida, restamos 480 en AMBOS lados. Piensa en una balanza: si hacemos lo mismo a ambos lados, la igualdad se mantiene.
\par\smallskip Hacemos la misma operación a la izquierda y a la derecha:
\[6x + 480 - 480=600 - 480\]
\par\smallskip Las partes que se compensan dan cero; queda:
\[6x=120\]
\par\smallskip\textbf{Qué aprendimos.} «Pasar al otro lado» es una forma abreviada de decir que hicimos la misma suma o resta en los dos lados.
\end{minipage}\par\medskip
\noindent\begin{minipage}{\linewidth}\textbf{19. Calculamos cuánto vale x}\\[2pt]Para deshacer una multiplicación por 6, dividimos AMBOS lados entre 6. En el lado izquierdo queda una sola x. En el derecho hacemos la división.
\par\smallskip Dividimos los dos lados entre el mismo número:
\[\frac{6x}{6}=\frac{120}{6}\]
\par\smallskip Ya podemos leer el valor de la letra:
\[x=\frac{120}{6}=20\]
\par\smallskip\textbf{Qué aprendimos.} Este valor se usará en las ecuaciones donde todavía falten letras.
\end{minipage}\par\medskip
\noindent\begin{minipage}{\linewidth}\textbf{20. Reunimos la respuesta; todavía falta comprobarla}\\[2pt]Escribimos todos los valores en el orden x, y, z. Dependen del número que elijamos para el parámetro, por eso describen una familia de respuestas.
\[\left\{\begin{aligned}x&=20\\y&=40 - t\\z&=t\end{aligned}\right.\]
\par\smallskip\textbf{Qué aprendimos.} El siguiente paso es volver a las ecuaciones del enunciado y comprobar cada condición.
\end{minipage}\par\medskip
\noindent\begin{minipage}{\linewidth}\textbf{21. Dos elecciones distintas producen dos soluciones}\\[2pt]Esta tabla muestra qué pasa al elegir 0 o 1 para el parámetro. Sustituimos la elección en las expresiones de la respuesta. Cada fila es un grupo completo de valores.
\par\smallskip {\small\renewcommand{\arraystretch}{1.3}\begin{tabular}{>{\raggedright\arraybackslash}p{2.68cm}>{\raggedright\arraybackslash}p{2.68cm}>{\raggedright\arraybackslash}p{2.68cm}>{\raggedright\arraybackslash}p{2.68cm}}\hline\rowcolor{black!7} \textbf{Elección de t} & \textbf{x} & \textbf{y} & \textbf{z}\\\hline 0 & 20 & 40 & 0\\ 1 & 20 & 39 & 1\\ \hline\end{tabular}}\par\smallskip
\par\smallskip\textbf{Qué aprendimos.} No hay una sola solución. Hay una para cada elección permitida del parámetro. El deslizador de la web permite comprobar más ejemplos.
\end{minipage}\par\medskip
\noindent\begin{minipage}{\linewidth}\textbf{22. ¿La respuesta cumple la ecuación 1?}\\[2pt]Comprobar significa volver a la condición original y calcularla con nuestra respuesta. Si los dos lados dan el mismo número, esta condición sí se cumple. Deben cumplirse todas.
\par\smallskip La condición del enunciado era:
\[6x + 12y + 12z = 600\]
\par\smallskip Ponemos nuestros valores donde estaban las letras:
\[6\left(20\right) + 12\left(40 - t\right) + 12\left(t\right)=600\]
\par\smallskip Hacemos las multiplicaciones y luego la suma:
\[120 + \left(480 - 12 t\right) + 12 t=600\]
\par\smallskip Comparamos los dos lados:
\[600=600\]
\par\smallskip\textbf{Qué aprendimos.} La igualdad final es verdadera. Al simplificar, los términos del parámetro se compensan: la condición vale para cualquier elección permitida.
\end{minipage}\par\medskip
\noindent\begin{minipage}{\linewidth}\textbf{23. ¿La respuesta cumple la ecuación 2?}\\[2pt]Comprobar significa volver a la condición original y calcularla con nuestra respuesta. Si los dos lados dan el mismo número, esta condición sí se cumple. Deben cumplirse todas.
\par\smallskip La condición del enunciado era:
\[24x + 12y + 12z = 960\]
\par\smallskip Ponemos nuestros valores donde estaban las letras:
\[24\left(20\right) + 12\left(40 - t\right) + 12\left(t\right)=960\]
\par\smallskip Hacemos las multiplicaciones y luego la suma:
\[480 + \left(480 - 12 t\right) + 12 t=960\]
\par\smallskip Comparamos los dos lados:
\[960=960\]
\par\smallskip\textbf{Qué aprendimos.} La igualdad final es verdadera. Al simplificar, los términos del parámetro se compensan: la condición vale para cualquier elección permitida.
\end{minipage}\par\medskip
\noindent\begin{minipage}{\linewidth}\textbf{24. Interpretamos y verificamos las unidades}\\[2pt]El revelador de 2 minutos debe producirse en 20 toneladas al día. Los otros dos deben sumar 40 toneladas: y = 40 - t, z = t, con 0 <= t <= 40. Hay infinitos planes no negativos.
\end{minipage}\par\medskip
\noindent\begin{minipage}{\linewidth}\textbf{25. Restringimos el parámetro al contexto}\\[2pt]El álgebra permite cualquier t real, pero las toneladas producidas no pueden ser negativas. Aplicamos y >= 0 y z >= 0.
\[y=40-t\geq0,\quad z=t\geq0\]
\[0\leq t\leq40\]
\end{minipage}\par\medskip
\noindent\begin{minipage}{\linewidth}\textbf{26. Ejemplos y comprobación}\\[2pt]Sirven (20, 40, 0), (20, 20, 20) y (20, 0, 40). No se exige que las toneladas sean números enteros. Para cualquier t permitido, las dos plantas se usan por completo.
\[6(20)+12(40-t)+12t=600\]
\[24(20)+12(40-t)+12t=960\]
\end{minipage}\par\medskip
\respuesta{x = 20; y = 40 - t; z = t. Parámetro real. El revelador de 2 minutos debe producirse en 20 toneladas al día. Los otros dos deben sumar 40 toneladas: y = 40 - t, z = t, con 0 <= t <= 40. Hay infinitos planes no negativos.}
\noindent\begin{minipage}{\linewidth}\centering\includegraphics[width=.96\linewidth,height=6.8cm,keepaspectratio]{../recursos/grafico_1_1-26.png}\\{\small Lectura visual complementaria; la demostración está en las operaciones.}\end{minipage}
\end{multicols*}
\newpage\fancyfoot[L]{\small Sección 1.1 | Ejercicio 27 | Paso a paso}\hypertarget{1.1-27}{}
\pdfbookmark[1]{1.1 / 27}{bm1.1-27}
\titulo{27 | Una parábola determinada por tres puntos}
\textbf{Enunciado}\par Suponga que los tres puntos (1, -5), (-1, 1) y (2, 7) están en la parábola p(x) = ax\textsuperscript{2} + bx + c. (a) Determine un sistema lineal de tres ecuaciones con tres incógnitas que deba resolverse para determinar a, b y c. (b) Resuelva el sistema lineal que obtuvo en la parte (a) para a, b y c.
\[\left\{\begin{aligned}a + b + c &= -5\\a - b + c &= 1\\4a + 2b + c &= 7\end{aligned}\right.\]
\begin{multicols*}{2}
\noindent\begin{minipage}{\linewidth}\textbf{1. ¿Cuáles son las incógnitas?}\\[2pt]En p(x) = ax\textsuperscript{2} + bx + c, la letra x es la entrada de la función y p(x) es su altura. Los puntos del enunciado nos dan entradas y alturas conocidas. Lo que falta encontrar son a, b y c: los números que definen la parábola.
\end{minipage}\par\medskip
\noindent\begin{minipage}{\linewidth}\textbf{2. Parte (a): usamos cada punto}\\[2pt]Cada punto (x, y) exige p(x) = y. Sustituimos uno por uno, incluyendo el cuadrado de -1.
\par\smallskip {\small\renewcommand{\arraystretch}{1.3}\begin{tabular}{>{\raggedright\arraybackslash}p{3.74cm}>{\raggedright\arraybackslash}p{3.74cm}>{\raggedright\arraybackslash}p{3.74cm}}\hline\rowcolor{black!7} \textbf{Punto conocido} & \textbf{Qué nos dice} & \textbf{Ecuación que obtenemos}\\\hline (1, -5) & Con x = 1, la altura es -5. & $a+b+c=-5$\\ (-1, 1) & Con x = -1, la altura es 1. & $a-b+c=1$\\ (2, 7) & Con x = 2, la altura es 7. & $4a+2b+c=7$\\ \hline\end{tabular}}\par\smallskip
\[p(1)=a+b+c=-5\]
\[p(-1)=a-b+c=1\]
\[p(2)=4a+2b+c=7\]
\par\smallskip\textbf{Qué aprendimos.} Al elevar -1 al cuadrado obtenemos 1. Por eso el término de a sigue siendo +a en la segunda ecuación. Estamos resolviendo para a, b y c, no para las coordenadas x de los puntos.
\end{minipage}\par\medskip
\noindent\begin{minipage}{\linewidth}\textbf{3. Primero: ¿qué significa resolver este sistema?}\\[2pt]Las letras a, b, c representan números que todavía no conocemos. Una solución es una elección de esos números que hace verdaderas TODAS las ecuaciones al mismo tiempo. La llave de la izquierda significa que las condiciones se deben cumplir juntas.
\par\smallskip\textbf{Qué buscamos.} Encontrar un mismo grupo de valores para las 3 letras.
\par\smallskip El número 3 delante de a indica tres veces ese mismo valor; es decir, sumar tres copias:
\[3a=a+a+a\]
\par\smallskip Ejemplo numérico de la notación: si a valiera 2, tres veces ese valor sería 6. Este 2 es solo un ejemplo; todavía no es la respuesta del ejercicio. El punto centrado indica multiplicación:
\[3\cdot 2=2+2+2=6\]
\par\smallskip\textbf{Comprueba que lo entiendes.} ¿Basta con que los valores funcionen en una sola ecuación?\par\textit{Respuesta:} No. Deben funcionar en todas las ecuaciones del sistema, usando los mismos valores.
\end{minipage}\par\medskip
\noindent\begin{minipage}{\linewidth}\textbf{4. Nuestro plan: ir quitando letras}\\[2pt]Una ecuación con una sola letra es más fácil de resolver. Vamos a combinar ecuaciones para que una letra desaparezca. Cuando sepamos un valor, lo pondremos en las ecuaciones que faltan.
\par\smallskip\textbf{Qué buscamos.} Primero quitar a de las otras ecuaciones. Después seguir con las letras restantes.
\par\smallskip\textbf{Por qué.} Si una cantidad se resta de sí misma, queda cero. Usaremos esta idea con términos que tengan la misma letra.
\par\smallskip Ejemplo de la idea que utilizaremos:
\[3x-3x=0\]
\par\smallskip\textbf{Qué aprendimos.} Eliminar una letra de una ecuación NO significa que su valor sea cero; significa que sus términos se cancelaron en esa cuenta.
\end{minipage}\par\medskip
\noindent\begin{minipage}{\linewidth}\textbf{5. ¿Cómo haremos desaparecer a?}\\[2pt]Trabajaremos con las ecuaciones 1 y 2. Queremos que los términos con a tengan el mismo tamaño. Así podremos restarlos.
\par\smallskip\textbf{Qué buscamos.} Conseguir términos con 1 veces a en las dos ecuaciones.
\par\smallskip\textbf{Por qué.} Tener 1 veces la misma letra en ambos lados de la resta nos permitirá cancelarla.
\par\smallskip {\small\renewcommand{\arraystretch}{1.3}\begin{tabular}{>{\raggedright\arraybackslash}p{2.58cm}>{\raggedright\arraybackslash}p{2.47cm}>{\raggedright\arraybackslash}p{6.17cm}}\hline\rowcolor{black!7} \textbf{Dónde miramos} & \textbf{Qué término hay} & \textbf{Qué necesitamos hacer}\\\hline Ecuación 1 & $a$ & La dejamos igual.\\ Ecuación 2 & $a$ & La dejamos igual.\\ \hline\end{tabular}}\par\smallskip
\par\smallskip\textbf{Qué aprendimos.} Estos multiplicadores se eligen para igualar los términos con a. Los demás términos también se multiplican para conservar la igualdad.
\end{minipage}\par\medskip
\noindent\begin{minipage}{\linewidth}\textbf{6. Restamos las dos ecuaciones preparadas}\\[2pt]Restamos la ecuación 1 preparada a la ecuación 2 preparada. Juntamos cada término solo con otro que tenga la misma letra. Los números sin letra se calculan por separado.
\par\smallskip {\small\renewcommand{\arraystretch}{1.3}\begin{tabular}{>{\raggedright\arraybackslash}p{2.80cm}>{\raggedright\arraybackslash}p{5.38cm}>{\raggedright\arraybackslash}p{3.03cm}}\hline\rowcolor{black!7} \textbf{Qué agrupamos} & \textbf{Cuenta} & \textbf{Qué significa}\\\hline Los términos con a & $a - \left(a\right)=0$ & Se cancela esta letra.\\ Los términos con b & $-b - \left(b\right)=-2b$ & Quedan -2b.\\ Los términos con c & $c - \left(c\right)=0$ & Se cancela esta letra.\\ Los números a la derecha del igual & $1 - \left(-5\right)=1+5=6$ & Este será el número de la nueva igualdad.\\ \hline\end{tabular}}\par\smallskip
\par\smallskip Primera igualdad que usaremos:
\[a - b + c = 1\]
\par\smallskip Le restamos esta igualdad:
\[a + b + c = -5\]
\par\smallskip\textbf{Qué aprendimos.} Los términos con a dan cero. Recuerda: restar un número negativo equivale a sumar su opuesto.
\end{minipage}\par\medskip
\noindent\begin{minipage}{\linewidth}\textbf{7. Mira lo que hemos conseguido}\\[2pt]La nueva ecuación es el resultado de la cuenta anterior. Solo contiene b: después podremos despejarla.
\par\smallskip Escribimos la nueva condición:
\[-2b = 6\]
\[\left\{\begin{aligned}a + b + c &= -5\\-2b &= 6\\4a + 2b + c &= 7\end{aligned}\right.\]
\par\smallskip\textbf{Qué aprendimos.} Conservamos la ecuación 1 y reemplazamos la 2 por esta nueva condición. De esa manera conservamos toda la información necesaria para resolver el sistema.
\end{minipage}\par\medskip
\noindent\begin{minipage}{\linewidth}\textbf{8. ¿Cómo haremos desaparecer a?}\\[2pt]Trabajaremos con las ecuaciones 1 y 3. Queremos que los términos con a tengan el mismo tamaño. Así podremos restarlos.
\par\smallskip\textbf{Qué buscamos.} Conseguir términos con 4 veces a en las dos ecuaciones.
\par\smallskip\textbf{Por qué.} Tener 4 veces la misma letra en ambos lados de la resta nos permitirá cancelarla.
\par\smallskip {\small\renewcommand{\arraystretch}{1.3}\begin{tabular}{>{\raggedright\arraybackslash}p{2.58cm}>{\raggedright\arraybackslash}p{2.47cm}>{\raggedright\arraybackslash}p{6.17cm}}\hline\rowcolor{black!7} \textbf{Dónde miramos} & \textbf{Qué término hay} & \textbf{Qué necesitamos hacer}\\\hline Ecuación 1 & $a$ & Multiplicar TODA esta ecuación por 4\\ Ecuación 3 & $4a$ & La dejamos igual.\\ \hline\end{tabular}}\par\smallskip
\par\smallskip\textbf{Qué aprendimos.} Estos multiplicadores se eligen para igualar los términos con a. Los demás términos también se multiplican para conservar la igualdad.
\end{minipage}\par\medskip
\noindent\begin{minipage}{\linewidth}\textbf{9. Multiplicamos la ecuación 1 por 4}\\[2pt]Multiplicar una ecuación por 4 significa multiplicar CADA término de la izquierda y también el número de la derecha. La igualdad se conserva porque hacemos lo mismo en ambos lados.
\par\smallskip Esta es la ecuación antes de multiplicar:
\[a + b + c = -5\]
\par\smallskip {\small\renewcommand{\arraystretch}{1.3}\begin{tabular}{>{\raggedright\arraybackslash}p{2.80cm}>{\raggedright\arraybackslash}p{1.91cm}>{\raggedright\arraybackslash}p{6.50cm}}\hline\rowcolor{black!7} \textbf{Parte de la ecuación} & \textbf{Antes} & \textbf{La cuenta y su resultado}\\\hline Término con a & $a$ & $4\cdot\left(a\right)=4a$\\ Término con b & $b$ & $4\cdot\left(b\right)=4b$\\ Término con c & $c$ & $4\cdot\left(c\right)=4c$\\ Número a la derecha del igual & $-5$ & $4\cdot\left(-5\right)=-20$\\ \hline\end{tabular}}\par\smallskip
\par\smallskip Después de hacer todas las multiplicaciones queda:
\[4a + 4b + 4c = -20\]
\par\smallskip\textbf{Qué aprendimos.} La misma condición ahora está escrita con los términos multiplicados. Ya tenemos el tamaño que buscábamos para a.
\par\smallskip\textbf{Comprueba que lo entiendes.} ¿Podríamos multiplicar solamente el término con a?\par\textit{Respuesta:} No. Eso cambiaría la condición. Debemos multiplicar todos los términos, también el número que está después del igual.
\end{minipage}\par\medskip
\noindent\begin{minipage}{\linewidth}\textbf{10. Restamos las dos ecuaciones preparadas}\\[2pt]Restamos la ecuación 1 preparada a la ecuación 3 preparada. Juntamos cada término solo con otro que tenga la misma letra. Los números sin letra se calculan por separado.
\par\smallskip {\small\renewcommand{\arraystretch}{1.3}\begin{tabular}{>{\raggedright\arraybackslash}p{2.80cm}>{\raggedright\arraybackslash}p{5.38cm}>{\raggedright\arraybackslash}p{3.03cm}}\hline\rowcolor{black!7} \textbf{Qué agrupamos} & \textbf{Cuenta} & \textbf{Qué significa}\\\hline Los términos con a & $4a - \left(4a\right)=0$ & Se cancela esta letra.\\ Los términos con b & $2b - \left(4b\right)=-2b$ & Quedan -2b.\\ Los términos con c & $c - \left(4c\right)=-3c$ & Quedan -3c.\\ Los números a la derecha del igual & $7 - \left(-20\right)=7+20=27$ & Este será el número de la nueva igualdad.\\ \hline\end{tabular}}\par\smallskip
\par\smallskip Primera igualdad que usaremos:
\[4a + 2b + c = 7\]
\par\smallskip Le restamos esta igualdad:
\[4a + 4b + 4c = -20\]
\par\smallskip\textbf{Qué aprendimos.} Los términos con a dan cero. Recuerda: restar un número negativo equivale a sumar su opuesto.
\end{minipage}\par\medskip
\noindent\begin{minipage}{\linewidth}\textbf{11. Mira lo que hemos conseguido}\\[2pt]La nueva ecuación es el resultado de la cuenta anterior. Todavía contiene b, c. Seguiremos reduciendo las condiciones.
\par\smallskip Escribimos la nueva condición:
\[-2b - 3c = 27\]
\[\left\{\begin{aligned}a + b + c &= -5\\-2b &= 6\\-2b - 3c &= 27\end{aligned}\right.\]
\par\smallskip\textbf{Qué aprendimos.} Conservamos la ecuación 1 y reemplazamos la 3 por esta nueva condición. De esa manera conservamos toda la información necesaria para resolver el sistema.
\end{minipage}\par\medskip
\noindent\begin{minipage}{\linewidth}\textbf{12. ¿Cómo haremos desaparecer b?}\\[2pt]Trabajaremos con las ecuaciones 2 y 3. Queremos que los términos con b tengan el mismo tamaño. Así podremos restarlos.
\par\smallskip\textbf{Qué buscamos.} Conseguir términos con 2 veces b en las dos ecuaciones.
\par\smallskip\textbf{Por qué.} Tener 2 veces la misma letra en ambos lados de la resta nos permitirá cancelarla.
\par\smallskip {\small\renewcommand{\arraystretch}{1.3}\begin{tabular}{>{\raggedright\arraybackslash}p{2.58cm}>{\raggedright\arraybackslash}p{2.47cm}>{\raggedright\arraybackslash}p{6.17cm}}\hline\rowcolor{black!7} \textbf{Dónde miramos} & \textbf{Qué término hay} & \textbf{Qué necesitamos hacer}\\\hline Ecuación 2 & $-2b$ & La dejamos igual.\\ Ecuación 3 & $-2b$ & La dejamos igual.\\ \hline\end{tabular}}\par\smallskip
\par\smallskip\textbf{Qué aprendimos.} Estos multiplicadores se eligen para igualar los términos con b. Los demás términos también se multiplican para conservar la igualdad.
\end{minipage}\par\medskip
\noindent\begin{minipage}{\linewidth}\textbf{13. Restamos las dos ecuaciones preparadas}\\[2pt]Restamos la ecuación 2 preparada a la ecuación 3 preparada. Juntamos cada término solo con otro que tenga la misma letra. Los números sin letra se calculan por separado.
\par\smallskip {\small\renewcommand{\arraystretch}{1.3}\begin{tabular}{>{\raggedright\arraybackslash}p{2.80cm}>{\raggedright\arraybackslash}p{5.38cm}>{\raggedright\arraybackslash}p{3.03cm}}\hline\rowcolor{black!7} \textbf{Qué agrupamos} & \textbf{Cuenta} & \textbf{Qué significa}\\\hline Los términos con b & $-2b - \left(-2b\right)=0$ & Se cancela esta letra.\\ Los términos con c & $-3c - \left(0\right)=-3c$ & Quedan -3c.\\ Los números a la derecha del igual & $27 - \left(6\right)=21$ & Este será el número de la nueva igualdad.\\ \hline\end{tabular}}\par\smallskip
\par\smallskip Primera igualdad que usaremos:
\[-2b - 3c = 27\]
\par\smallskip Le restamos esta igualdad:
\[-2b = 6\]
\par\smallskip\textbf{Qué aprendimos.} Los términos con b dan cero. Los términos con letras distintas no se mezclan.
\end{minipage}\par\medskip
\noindent\begin{minipage}{\linewidth}\textbf{14. Mira lo que hemos conseguido}\\[2pt]La nueva ecuación es el resultado de la cuenta anterior. Solo contiene c: después podremos despejarla.
\par\smallskip Escribimos la nueva condición:
\[-3c = 21\]
\[\left\{\begin{aligned}a + b + c &= -5\\-2b &= 6\\-3c &= 21\end{aligned}\right.\]
\par\smallskip\textbf{Qué aprendimos.} Conservamos la ecuación 2 y reemplazamos la 3 por esta nueva condición. De esa manera conservamos toda la información necesaria para resolver el sistema.
\end{minipage}\par\medskip
\noindent\begin{minipage}{\linewidth}\textbf{15. Ya tenemos una ecuación con una sola letra}\\[2pt]En esta condición solamente falta conocer c. El número -3 está multiplicando a esa letra. Ahora podemos deshacer esa multiplicación.
\par\smallskip Esta es la condición que vamos a resolver:
\[-3c = 21\]
\end{minipage}\par\medskip
\noindent\begin{minipage}{\linewidth}\textbf{16. Calculamos cuánto vale c}\\[2pt]Para deshacer una multiplicación por -3, dividimos AMBOS lados entre -3. En el lado izquierdo queda una sola c. En el derecho hacemos la división.
\par\smallskip Dividimos los dos lados entre el mismo número:
\[\frac{-3c}{-3}=\frac{21}{-3}\]
\par\smallskip Ya podemos leer el valor de la letra:
\[c=\frac{21}{-3}=-7\]
\par\smallskip\textbf{Qué aprendimos.} Este valor se usará en las ecuaciones donde todavía falten letras.
\end{minipage}\par\medskip
\noindent\begin{minipage}{\linewidth}\textbf{17. Ya tenemos una ecuación con una sola letra}\\[2pt]En esta condición solamente falta conocer b. El número -2 está multiplicando a esa letra. Ahora podemos deshacer esa multiplicación.
\par\smallskip Esta es la condición que vamos a resolver:
\[-2b = 6\]
\end{minipage}\par\medskip
\noindent\begin{minipage}{\linewidth}\textbf{18. Calculamos cuánto vale b}\\[2pt]Para deshacer una multiplicación por -2, dividimos AMBOS lados entre -2. En el lado izquierdo queda una sola b. En el derecho hacemos la división.
\par\smallskip Dividimos los dos lados entre el mismo número:
\[\frac{-2b}{-2}=\frac{6}{-2}\]
\par\smallskip Ya podemos leer el valor de la letra:
\[b=\frac{6}{-2}=-3\]
\par\smallskip\textbf{Qué aprendimos.} Este valor se usará en las ecuaciones donde todavía falten letras.
\end{minipage}\par\medskip
\noindent\begin{minipage}{\linewidth}\textbf{19. Usamos lo que ya sabemos para hallar a}\\[2pt]Ya sabemos b = -3; c = -7. Sustituir significa escribir ese valor cada vez que aparezca su letra. Usamos paréntesis para que se vea qué número estamos multiplicando.
\par\smallskip\textbf{Qué buscamos.} Dejar solamente a como letra desconocida en esta ecuación.
\par\smallskip Tomamos esta condición que obtuvimos antes:
\[a + b + c = -5\]
\par\smallskip Reemplazamos las letras conocidas por sus valores:
\[a + \left(-3\right) + \left(-7\right)=-5\]
\end{minipage}\par\medskip
\noindent\begin{minipage}{\linewidth}\textbf{20. Calculamos las multiplicaciones que ya se pueden hacer}\\[2pt]La letra a sigue sin conocerse. Las otras partes ya tienen un valor o una expresión conocida y podemos efectuar sus productos.
\par\smallskip {\small\renewcommand{\arraystretch}{1.3}\begin{tabular}{>{\raggedright\arraybackslash}p{3.36cm}>{\raggedright\arraybackslash}p{4.48cm}>{\raggedright\arraybackslash}p{3.36cm}}\hline\rowcolor{black!7} \textbf{Parte conocida} & \textbf{La multiplicación} & \textbf{Resultado}\\\hline El término b & $1\cdot\left(-3\right)$ & $-3$\\ El término c & $1\cdot\left(-7\right)$ & $-7$\\ \hline\end{tabular}}\par\smallskip
\par\smallskip Al reunir esas partes conocidas, la ecuación queda:
\[a - 10=-5\]
\end{minipage}\par\medskip
\noindent\begin{minipage}{\linewidth}\textbf{21. Dejamos el término con a solo a un lado}\\[2pt]Para quitar la parte conocida, sumamos 10 en AMBOS lados. Piensa en una balanza: si hacemos lo mismo a ambos lados, la igualdad se mantiene.
\par\smallskip Hacemos la misma operación a la izquierda y a la derecha:
\[a - 10 + 10=-5 + 10\]
\par\smallskip Las partes que se compensan dan cero; queda:
\[a=5\]
\par\smallskip\textbf{Qué aprendimos.} «Pasar al otro lado» es una forma abreviada de decir que hicimos la misma suma o resta en los dos lados.
\end{minipage}\par\medskip
\noindent\begin{minipage}{\linewidth}\textbf{22. Leemos el valor de a}\\[2pt]La letra ya está sola: su número multiplicador es 1, que no se escribe. La igualdad nos dice directamente cuánto vale.
\par\smallskip Anotamos el valor encontrado:
\[a=5\]
\end{minipage}\par\medskip
\noindent\begin{minipage}{\linewidth}\textbf{23. Reunimos la respuesta; todavía falta comprobarla}\\[2pt]Escribimos todos los valores en el orden a, b, c. Estos valores deben funcionar juntos, no cada uno en una ecuación diferente.
\[\left\{\begin{aligned}a&=5\\b&=-3\\c&=-7\end{aligned}\right.\]
\par\smallskip\textbf{Qué aprendimos.} El siguiente paso es volver a las ecuaciones del enunciado y comprobar cada condición.
\end{minipage}\par\medskip
\noindent\begin{minipage}{\linewidth}\textbf{24. ¿La respuesta cumple la ecuación 1?}\\[2pt]Comprobar significa volver a la condición original y calcularla con nuestra respuesta. Si los dos lados dan el mismo número, esta condición sí se cumple. Deben cumplirse todas.
\par\smallskip La condición del enunciado era:
\[a + b + c = -5\]
\par\smallskip Ponemos nuestros valores donde estaban las letras:
\[\left(5\right) + \left(-3\right) + \left(-7\right)=-5\]
\par\smallskip Hacemos las multiplicaciones y luego la suma:
\[5 - 3 - 7=-5\]
\par\smallskip Comparamos los dos lados:
\[-5=-5\]
\par\smallskip\textbf{Qué aprendimos.} La igualdad final es verdadera. Nuestra respuesta sí cumple la ecuación 1.
\end{minipage}\par\medskip
\noindent\begin{minipage}{\linewidth}\textbf{25. ¿La respuesta cumple la ecuación 2?}\\[2pt]Comprobar significa volver a la condición original y calcularla con nuestra respuesta. Si los dos lados dan el mismo número, esta condición sí se cumple. Deben cumplirse todas.
\par\smallskip La condición del enunciado era:
\[a - b + c = 1\]
\par\smallskip Ponemos nuestros valores donde estaban las letras:
\[\left(5\right) - \left(-3\right) + \left(-7\right)=1\]
\par\smallskip Hacemos las multiplicaciones y luego la suma:
\[5 + 3 - 7=1\]
\par\smallskip Comparamos los dos lados:
\[1=1\]
\par\smallskip\textbf{Qué aprendimos.} La igualdad final es verdadera. Nuestra respuesta sí cumple la ecuación 2.
\end{minipage}\par\medskip
\noindent\begin{minipage}{\linewidth}\textbf{26. ¿La respuesta cumple la ecuación 3?}\\[2pt]Comprobar significa volver a la condición original y calcularla con nuestra respuesta. Si los dos lados dan el mismo número, esta condición sí se cumple. Deben cumplirse todas.
\par\smallskip La condición del enunciado era:
\[4a + 2b + c = 7\]
\par\smallskip Ponemos nuestros valores donde estaban las letras:
\[4\left(5\right) + 2\left(-3\right) + \left(-7\right)=7\]
\par\smallskip Hacemos las multiplicaciones y luego la suma:
\[20 - 6 - 7=7\]
\par\smallskip Comparamos los dos lados:
\[7=7\]
\par\smallskip\textbf{Qué aprendimos.} La igualdad final es verdadera. Nuestra respuesta sí cumple la ecuación 3.
\end{minipage}\par\medskip
\noindent\begin{minipage}{\linewidth}\textbf{27. Interpretamos y verificamos las unidades}\\[2pt]a = 5, b = -3 y c = -7. La parábola es p(x) = 5x\textsuperscript{2} - 3x - 7.
\end{minipage}\par\medskip
\noindent\begin{minipage}{\linewidth}\textbf{28. Comprobamos los tres puntos}\\[2pt]La fórmula recupera exactamente las tres alturas dadas.
\[p(1)=5-3-7=-5\]
\[p(-1)=5+3-7=1\]
\[p(2)=20-6-7=7\]
\end{minipage}\par\medskip
\respuesta{a = 5, b = -3 y c = -7. La parábola es p(x) = 5x\textsuperscript{2} - 3x - 7.}
\noindent\begin{minipage}{\linewidth}\centering\includegraphics[width=.96\linewidth,height=6.8cm,keepaspectratio]{../recursos/grafico_1_1-27.png}\\{\small Lectura visual complementaria; la demostración está en las operaciones.}\end{minipage}
\end{multicols*}
\newpage\fancyfoot[L]{\small Sección 1.1 | Ejercicio 28 | Paso a paso}\hypertarget{1.1-28}{}
\pdfbookmark[1]{1.1 / 28}{bm1.1-28}
\titulo{28 | Reparto de una herencia e intereses}
\textbf{Enunciado}\par Una herencia de \$24,000 se dividió en tres fideicomisos; el segundo fideicomiso recibió el doble del primero. Los tres fideicomisos pagan una tasa de interés de 9, 10 y 6\% anual, respectivamente; al final del primer año, el rendimiento total fue de \$2,210. ¿Cuánto se invirtió en cada fideicomiso?
\[\left\{\begin{aligned}x + y + z &= 24000\\ - 2x + y &= 0\\9x + 10y + 6z &= 221000\end{aligned}\right.\]
\begin{multicols*}{2}
\noindent\begin{minipage}{\linewidth}\textbf{1. Definimos los capitales}\\[2pt]Un fideicomiso es aquí una inversión que paga intereses. Llamaremos x, y, z a los dólares invertidos en el primero, segundo y tercero. Capital es el dinero colocado; interés es lo que ese dinero produce. Los \$24,000 son capital y los \$2,210 son intereses: son dos cantidades distintas.
\par\smallskip {\small\renewcommand{\arraystretch}{1.3}\begin{tabular}{>{\raggedright\arraybackslash}p{2.68cm}>{\raggedright\arraybackslash}p{2.68cm}>{\raggedright\arraybackslash}p{2.68cm}>{\raggedright\arraybackslash}p{2.68cm}}\hline\rowcolor{black!7} \textbf{Inversión} & \textbf{Capital} & \textbf{Tasa anual como decimal} & \textbf{Interés del primer año}\\\hline Primera & x & 9\% = 0.09 & $0.09x$\\ Segunda & y & 10\% = 0.10 & $0.10y$\\ Tercera & z & 6\% = 0.06 & $0.06z$\\ \hline\end{tabular}}\par\smallskip
\end{minipage}\par\medskip
\noindent\begin{minipage}{\linewidth}\textbf{2. Escribimos las tres condiciones}\\[2pt]La suma de capitales es \$24,000; el segundo duplica al primero; cada interés anual es tasa por capital.
\par\smallskip\textbf{Por qué.} Decir que la segunda inversión recibió el doble significa y = 2x. La suma de los tres capitales es 24000. La suma de los intereses, calculados como tasa por capital, es 2210.
\[x+y+z=24000\]
\[y=2x\]
\[0.09x+0.10y+0.06z=2210\]
\end{minipage}\par\medskip
\noindent\begin{minipage}{\linewidth}\textbf{3. Quitamos decimales de la ecuación de intereses}\\[2pt]Multiplicamos TODA la tercera ecuación por 100. El término de la derecha pasa a 221000; no a 2210.
\[9x+10y+6z=221000\]
\par\smallskip\textbf{Qué aprendimos.} Multiplicar por 100 sirve para trabajar sin decimales; no significa que ahora tengamos más dinero. Solo reescribimos la misma igualdad.
\end{minipage}\par\medskip
\noindent\begin{minipage}{\linewidth}\textbf{4. Primero: ¿qué significa resolver este sistema?}\\[2pt]Las letras x, y, z representan números que todavía no conocemos. Una solución es una elección de esos números que hace verdaderas TODAS las ecuaciones al mismo tiempo. La llave de la izquierda significa que las condiciones se deben cumplir juntas.
\par\smallskip\textbf{Qué buscamos.} Encontrar un mismo grupo de valores para las 3 letras.
\par\smallskip El número 3 delante de x indica tres veces ese mismo valor; es decir, sumar tres copias:
\[3x=x+x+x\]
\par\smallskip Ejemplo numérico de la notación: si x valiera 2, tres veces ese valor sería 6. Este 2 es solo un ejemplo; todavía no es la respuesta del ejercicio. El punto centrado indica multiplicación:
\[3\cdot 2=2+2+2=6\]
\par\smallskip\textbf{Comprueba que lo entiendes.} ¿Basta con que los valores funcionen en una sola ecuación?\par\textit{Respuesta:} No. Deben funcionar en todas las ecuaciones del sistema, usando los mismos valores.
\end{minipage}\par\medskip
\noindent\begin{minipage}{\linewidth}\textbf{5. Nuestro plan: ir quitando letras}\\[2pt]Una ecuación con una sola letra es más fácil de resolver. Vamos a combinar ecuaciones para que una letra desaparezca. Cuando sepamos un valor, lo pondremos en las ecuaciones que faltan.
\par\smallskip\textbf{Qué buscamos.} Primero quitar x de las otras ecuaciones. Después seguir con las letras restantes.
\par\smallskip\textbf{Por qué.} Si una cantidad se resta de sí misma, queda cero. Usaremos esta idea con términos que tengan la misma letra.
\par\smallskip Ejemplo de la idea que utilizaremos:
\[3x-3x=0\]
\par\smallskip\textbf{Qué aprendimos.} Eliminar una letra de una ecuación NO significa que su valor sea cero; significa que sus términos se cancelaron en esa cuenta.
\end{minipage}\par\medskip
\noindent\begin{minipage}{\linewidth}\textbf{6. ¿Cómo haremos desaparecer x?}\\[2pt]Trabajaremos con las ecuaciones 1 y 2. Queremos que los términos con x tengan el mismo tamaño. Así podremos sumarlos, porque tienen signos opuestos.
\par\smallskip\textbf{Qué buscamos.} Conseguir términos con 2 veces x en las dos ecuaciones.
\par\smallskip\textbf{Por qué.} Al sumar la misma cantidad con signos opuestos obtenemos cero.
\par\smallskip {\small\renewcommand{\arraystretch}{1.3}\begin{tabular}{>{\raggedright\arraybackslash}p{2.58cm}>{\raggedright\arraybackslash}p{2.47cm}>{\raggedright\arraybackslash}p{6.17cm}}\hline\rowcolor{black!7} \textbf{Dónde miramos} & \textbf{Qué término hay} & \textbf{Qué necesitamos hacer}\\\hline Ecuación 1 & $x$ & Multiplicar TODA esta ecuación por 2\\ Ecuación 2 & $-2x$ & La dejamos igual.\\ \hline\end{tabular}}\par\smallskip
\par\smallskip\textbf{Qué aprendimos.} Estos multiplicadores se eligen para igualar los términos con x. Los demás términos también se multiplican para conservar la igualdad.
\end{minipage}\par\medskip
\noindent\begin{minipage}{\linewidth}\textbf{7. Multiplicamos la ecuación 1 por 2}\\[2pt]Multiplicar una ecuación por 2 significa multiplicar CADA término de la izquierda y también el número de la derecha. La igualdad se conserva porque hacemos lo mismo en ambos lados.
\par\smallskip Esta es la ecuación antes de multiplicar:
\[x + y + z = 24000\]
\par\smallskip {\small\renewcommand{\arraystretch}{1.3}\begin{tabular}{>{\raggedright\arraybackslash}p{2.80cm}>{\raggedright\arraybackslash}p{1.91cm}>{\raggedright\arraybackslash}p{6.50cm}}\hline\rowcolor{black!7} \textbf{Parte de la ecuación} & \textbf{Antes} & \textbf{La cuenta y su resultado}\\\hline Término con x & $x$ & $2\cdot\left(x\right)=2x$\\ Término con y & $y$ & $2\cdot\left(y\right)=2y$\\ Término con z & $z$ & $2\cdot\left(z\right)=2z$\\ Número a la derecha del igual & $24000$ & $2\cdot\left(24000\right)=48000$\\ \hline\end{tabular}}\par\smallskip
\par\smallskip Después de hacer todas las multiplicaciones queda:
\[2x + 2y + 2z = 48000\]
\par\smallskip\textbf{Qué aprendimos.} La misma condición ahora está escrita con los términos multiplicados. Ya tenemos el tamaño que buscábamos para x.
\par\smallskip\textbf{Comprueba que lo entiendes.} ¿Podríamos multiplicar solamente el término con x?\par\textit{Respuesta:} No. Eso cambiaría la condición. Debemos multiplicar todos los términos, también el número que está después del igual.
\end{minipage}\par\medskip
\noindent\begin{minipage}{\linewidth}\textbf{8. Sumamos las dos ecuaciones preparadas}\\[2pt]Sumamos las dos ecuaciones preparadas. Juntamos cada término solo con otro que tenga la misma letra. Los números sin letra se calculan por separado.
\par\smallskip {\small\renewcommand{\arraystretch}{1.3}\begin{tabular}{>{\raggedright\arraybackslash}p{2.80cm}>{\raggedright\arraybackslash}p{5.38cm}>{\raggedright\arraybackslash}p{3.03cm}}\hline\rowcolor{black!7} \textbf{Qué agrupamos} & \textbf{Cuenta} & \textbf{Qué significa}\\\hline Los términos con x & $-2x + \left(2x\right)=0$ & Se cancela esta letra.\\ Los términos con y & $y + \left(2y\right)=3y$ & Quedan 3y.\\ Los términos con z & $0 + \left(2z\right)=2z$ & Quedan 2z.\\ Los números a la derecha del igual & $0 + \left(48000\right)=48000$ & Este será el número de la nueva igualdad.\\ \hline\end{tabular}}\par\smallskip
\par\smallskip Primera igualdad que usaremos:
\[-2x + y = 0\]
\par\smallskip Le sumamos esta igualdad:
\[2x + 2y + 2z = 48000\]
\par\smallskip\textbf{Qué aprendimos.} Los términos con x dan cero. Los términos con letras distintas no se mezclan.
\end{minipage}\par\medskip
\noindent\begin{minipage}{\linewidth}\textbf{9. Mira lo que hemos conseguido}\\[2pt]La nueva ecuación es el resultado de la cuenta anterior. Todavía contiene y, z. Seguiremos reduciendo las condiciones.
\par\smallskip Escribimos la nueva condición:
\[3y + 2z = 48000\]
\[\left\{\begin{aligned}x + y + z &= 24000\\3y + 2z &= 48000\\9x + 10y + 6z &= 221000\end{aligned}\right.\]
\par\smallskip\textbf{Qué aprendimos.} Conservamos la ecuación 1 y reemplazamos la 2 por esta nueva condición. De esa manera conservamos toda la información necesaria para resolver el sistema.
\end{minipage}\par\medskip
\noindent\begin{minipage}{\linewidth}\textbf{10. ¿Cómo haremos desaparecer x?}\\[2pt]Trabajaremos con las ecuaciones 1 y 3. Queremos que los términos con x tengan el mismo tamaño. Así podremos restarlos.
\par\smallskip\textbf{Qué buscamos.} Conseguir términos con 9 veces x en las dos ecuaciones.
\par\smallskip\textbf{Por qué.} Tener 9 veces la misma letra en ambos lados de la resta nos permitirá cancelarla.
\par\smallskip {\small\renewcommand{\arraystretch}{1.3}\begin{tabular}{>{\raggedright\arraybackslash}p{2.58cm}>{\raggedright\arraybackslash}p{2.47cm}>{\raggedright\arraybackslash}p{6.17cm}}\hline\rowcolor{black!7} \textbf{Dónde miramos} & \textbf{Qué término hay} & \textbf{Qué necesitamos hacer}\\\hline Ecuación 1 & $x$ & Multiplicar TODA esta ecuación por 9\\ Ecuación 3 & $9x$ & La dejamos igual.\\ \hline\end{tabular}}\par\smallskip
\par\smallskip\textbf{Qué aprendimos.} Estos multiplicadores se eligen para igualar los términos con x. Los demás términos también se multiplican para conservar la igualdad.
\end{minipage}\par\medskip
\noindent\begin{minipage}{\linewidth}\textbf{11. Multiplicamos la ecuación 1 por 9}\\[2pt]Multiplicar una ecuación por 9 significa multiplicar CADA término de la izquierda y también el número de la derecha. La igualdad se conserva porque hacemos lo mismo en ambos lados.
\par\smallskip Esta es la ecuación antes de multiplicar:
\[x + y + z = 24000\]
\par\smallskip {\small\renewcommand{\arraystretch}{1.3}\begin{tabular}{>{\raggedright\arraybackslash}p{2.80cm}>{\raggedright\arraybackslash}p{1.91cm}>{\raggedright\arraybackslash}p{6.50cm}}\hline\rowcolor{black!7} \textbf{Parte de la ecuación} & \textbf{Antes} & \textbf{La cuenta y su resultado}\\\hline Término con x & $x$ & $9\cdot\left(x\right)=9x$\\ Término con y & $y$ & $9\cdot\left(y\right)=9y$\\ Término con z & $z$ & $9\cdot\left(z\right)=9z$\\ Número a la derecha del igual & $24000$ & $9\cdot\left(24000\right)=216000$\\ \hline\end{tabular}}\par\smallskip
\par\smallskip Después de hacer todas las multiplicaciones queda:
\[9x + 9y + 9z = 216000\]
\par\smallskip\textbf{Qué aprendimos.} La misma condición ahora está escrita con los términos multiplicados. Ya tenemos el tamaño que buscábamos para x.
\par\smallskip\textbf{Comprueba que lo entiendes.} ¿Podríamos multiplicar solamente el término con x?\par\textit{Respuesta:} No. Eso cambiaría la condición. Debemos multiplicar todos los términos, también el número que está después del igual.
\end{minipage}\par\medskip
\noindent\begin{minipage}{\linewidth}\textbf{12. Restamos las dos ecuaciones preparadas}\\[2pt]Restamos la ecuación 1 preparada a la ecuación 3 preparada. Juntamos cada término solo con otro que tenga la misma letra. Los números sin letra se calculan por separado.
\par\smallskip {\small\renewcommand{\arraystretch}{1.3}\begin{tabular}{>{\raggedright\arraybackslash}p{2.80cm}>{\raggedright\arraybackslash}p{5.38cm}>{\raggedright\arraybackslash}p{3.03cm}}\hline\rowcolor{black!7} \textbf{Qué agrupamos} & \textbf{Cuenta} & \textbf{Qué significa}\\\hline Los términos con x & $9x - \left(9x\right)=0$ & Se cancela esta letra.\\ Los términos con y & $10y - \left(9y\right)=y$ & Quedan y.\\ Los términos con z & $6z - \left(9z\right)=-3z$ & Quedan -3z.\\ Los números a la derecha del igual & $221000 - \left(216000\right)=5000$ & Este será el número de la nueva igualdad.\\ \hline\end{tabular}}\par\smallskip
\par\smallskip Primera igualdad que usaremos:
\[9x + 10y + 6z = 221000\]
\par\smallskip Le restamos esta igualdad:
\[9x + 9y + 9z = 216000\]
\par\smallskip\textbf{Qué aprendimos.} Los términos con x dan cero. Los términos con letras distintas no se mezclan.
\end{minipage}\par\medskip
\noindent\begin{minipage}{\linewidth}\textbf{13. Mira lo que hemos conseguido}\\[2pt]La nueva ecuación es el resultado de la cuenta anterior. Todavía contiene y, z. Seguiremos reduciendo las condiciones.
\par\smallskip Escribimos la nueva condición:
\[y - 3z = 5000\]
\[\left\{\begin{aligned}x + y + z &= 24000\\3y + 2z &= 48000\\y - 3z &= 5000\end{aligned}\right.\]
\par\smallskip\textbf{Qué aprendimos.} Conservamos la ecuación 1 y reemplazamos la 3 por esta nueva condición. De esa manera conservamos toda la información necesaria para resolver el sistema.
\end{minipage}\par\medskip
\noindent\begin{minipage}{\linewidth}\textbf{14. Ponemos una ecuación más cómoda arriba}\\[2pt]Cambiamos de lugar las ecuaciones 2 y 3. Las condiciones siguen siendo las mismas; solamente cambia el orden en que las leemos.
\par\smallskip\textbf{Por qué.} La ecuación que subimos tiene un número más pequeño delante de y. Eso facilita las multiplicaciones que haremos a continuación.
\[\left\{\begin{aligned}x + y + z &= 24000\\y - 3z &= 5000\\3y + 2z &= 48000\end{aligned}\right.\]
\end{minipage}\par\medskip
\noindent\begin{minipage}{\linewidth}\textbf{15. ¿Cómo haremos desaparecer y?}\\[2pt]Trabajaremos con las ecuaciones 2 y 3. Queremos que los términos con y tengan el mismo tamaño. Así podremos restarlos.
\par\smallskip\textbf{Qué buscamos.} Conseguir términos con 3 veces y en las dos ecuaciones.
\par\smallskip\textbf{Por qué.} Tener 3 veces la misma letra en ambos lados de la resta nos permitirá cancelarla.
\par\smallskip {\small\renewcommand{\arraystretch}{1.3}\begin{tabular}{>{\raggedright\arraybackslash}p{2.58cm}>{\raggedright\arraybackslash}p{2.47cm}>{\raggedright\arraybackslash}p{6.17cm}}\hline\rowcolor{black!7} \textbf{Dónde miramos} & \textbf{Qué término hay} & \textbf{Qué necesitamos hacer}\\\hline Ecuación 2 & $y$ & Multiplicar TODA esta ecuación por 3\\ Ecuación 3 & $3y$ & La dejamos igual.\\ \hline\end{tabular}}\par\smallskip
\par\smallskip\textbf{Qué aprendimos.} Estos multiplicadores se eligen para igualar los términos con y. Los demás términos también se multiplican para conservar la igualdad.
\end{minipage}\par\medskip
\noindent\begin{minipage}{\linewidth}\textbf{16. Multiplicamos la ecuación 2 por 3}\\[2pt]Multiplicar una ecuación por 3 significa multiplicar CADA término de la izquierda y también el número de la derecha. La igualdad se conserva porque hacemos lo mismo en ambos lados.
\par\smallskip Esta es la ecuación antes de multiplicar:
\[y - 3z = 5000\]
\par\smallskip {\small\renewcommand{\arraystretch}{1.3}\begin{tabular}{>{\raggedright\arraybackslash}p{2.80cm}>{\raggedright\arraybackslash}p{1.91cm}>{\raggedright\arraybackslash}p{6.50cm}}\hline\rowcolor{black!7} \textbf{Parte de la ecuación} & \textbf{Antes} & \textbf{La cuenta y su resultado}\\\hline Término con y & $y$ & $3\cdot\left(y\right)=3y$\\ Término con z & $-3z$ & $3\cdot\left(-3z\right)=-9z$\\ Número a la derecha del igual & $5000$ & $3\cdot\left(5000\right)=15000$\\ \hline\end{tabular}}\par\smallskip
\par\smallskip Después de hacer todas las multiplicaciones queda:
\[3y - 9z = 15000\]
\par\smallskip\textbf{Qué aprendimos.} La misma condición ahora está escrita con los términos multiplicados. Ya tenemos el tamaño que buscábamos para y.
\par\smallskip\textbf{Comprueba que lo entiendes.} ¿Podríamos multiplicar solamente el término con y?\par\textit{Respuesta:} No. Eso cambiaría la condición. Debemos multiplicar todos los términos, también el número que está después del igual.
\end{minipage}\par\medskip
\noindent\begin{minipage}{\linewidth}\textbf{17. Restamos las dos ecuaciones preparadas}\\[2pt]Restamos la ecuación 2 preparada a la ecuación 3 preparada. Juntamos cada término solo con otro que tenga la misma letra. Los números sin letra se calculan por separado.
\par\smallskip {\small\renewcommand{\arraystretch}{1.3}\begin{tabular}{>{\raggedright\arraybackslash}p{2.80cm}>{\raggedright\arraybackslash}p{5.38cm}>{\raggedright\arraybackslash}p{3.03cm}}\hline\rowcolor{black!7} \textbf{Qué agrupamos} & \textbf{Cuenta} & \textbf{Qué significa}\\\hline Los términos con y & $3y - \left(3y\right)=0$ & Se cancela esta letra.\\ Los términos con z & $2z - \left(-9z\right)=11z$ & Quedan 11z.\\ Los números a la derecha del igual & $48000 - \left(15000\right)=33000$ & Este será el número de la nueva igualdad.\\ \hline\end{tabular}}\par\smallskip
\par\smallskip Primera igualdad que usaremos:
\[3y + 2z = 48000\]
\par\smallskip Le restamos esta igualdad:
\[3y - 9z = 15000\]
\par\smallskip\textbf{Qué aprendimos.} Los términos con y dan cero. Los términos con letras distintas no se mezclan.
\end{minipage}\par\medskip
\noindent\begin{minipage}{\linewidth}\textbf{18. Mira lo que hemos conseguido}\\[2pt]La nueva ecuación es el resultado de la cuenta anterior. Solo contiene z: después podremos despejarla.
\par\smallskip Escribimos la nueva condición:
\[11z = 33000\]
\[\left\{\begin{aligned}x + y + z &= 24000\\y - 3z &= 5000\\11z &= 33000\end{aligned}\right.\]
\par\smallskip\textbf{Qué aprendimos.} Conservamos la ecuación 2 y reemplazamos la 3 por esta nueva condición. De esa manera conservamos toda la información necesaria para resolver el sistema.
\end{minipage}\par\medskip
\noindent\begin{minipage}{\linewidth}\textbf{19. Ya tenemos una ecuación con una sola letra}\\[2pt]En esta condición solamente falta conocer z. El número 11 está multiplicando a esa letra. Ahora podemos deshacer esa multiplicación.
\par\smallskip Esta es la condición que vamos a resolver:
\[11z = 33000\]
\end{minipage}\par\medskip
\noindent\begin{minipage}{\linewidth}\textbf{20. Calculamos cuánto vale z}\\[2pt]Para deshacer una multiplicación por 11, dividimos AMBOS lados entre 11. En el lado izquierdo queda una sola z. En el derecho hacemos la división.
\par\smallskip Dividimos los dos lados entre el mismo número:
\[\frac{11z}{11}=\frac{33000}{11}\]
\par\smallskip Ya podemos leer el valor de la letra:
\[z=\frac{33000}{11}=3000\]
\par\smallskip\textbf{Qué aprendimos.} Este valor se usará en las ecuaciones donde todavía falten letras.
\end{minipage}\par\medskip
\noindent\begin{minipage}{\linewidth}\textbf{21. Usamos lo que ya sabemos para hallar y}\\[2pt]Ya sabemos z = 3000. Sustituir significa escribir ese valor cada vez que aparezca su letra. Usamos paréntesis para que se vea qué número estamos multiplicando.
\par\smallskip\textbf{Qué buscamos.} Dejar solamente y como letra desconocida en esta ecuación.
\par\smallskip Tomamos esta condición que obtuvimos antes:
\[y - 3z = 5000\]
\par\smallskip Reemplazamos las letras conocidas por sus valores:
\[y - 3\left(3000\right)=5000\]
\end{minipage}\par\medskip
\noindent\begin{minipage}{\linewidth}\textbf{22. Calculamos las multiplicaciones que ya se pueden hacer}\\[2pt]La letra y sigue sin conocerse. Las otras partes ya tienen un valor o una expresión conocida y podemos efectuar sus productos.
\par\smallskip {\small\renewcommand{\arraystretch}{1.3}\begin{tabular}{>{\raggedright\arraybackslash}p{3.36cm}>{\raggedright\arraybackslash}p{4.48cm}>{\raggedright\arraybackslash}p{3.36cm}}\hline\rowcolor{black!7} \textbf{Parte conocida} & \textbf{La multiplicación} & \textbf{Resultado}\\\hline El término -3z & $-3\cdot\left(3000\right)$ & $-9000$\\ \hline\end{tabular}}\par\smallskip
\par\smallskip Al reunir esas partes conocidas, la ecuación queda:
\[y - 9000=5000\]
\end{minipage}\par\medskip
\noindent\begin{minipage}{\linewidth}\textbf{23. Dejamos el término con y solo a un lado}\\[2pt]Para quitar la parte conocida, sumamos 9000 en AMBOS lados. Piensa en una balanza: si hacemos lo mismo a ambos lados, la igualdad se mantiene.
\par\smallskip Hacemos la misma operación a la izquierda y a la derecha:
\[y - 9000 + 9000=5000 + 9000\]
\par\smallskip Las partes que se compensan dan cero; queda:
\[y=14000\]
\par\smallskip\textbf{Qué aprendimos.} «Pasar al otro lado» es una forma abreviada de decir que hicimos la misma suma o resta en los dos lados.
\end{minipage}\par\medskip
\noindent\begin{minipage}{\linewidth}\textbf{24. Leemos el valor de y}\\[2pt]La letra ya está sola: su número multiplicador es 1, que no se escribe. La igualdad nos dice directamente cuánto vale.
\par\smallskip Anotamos el valor encontrado:
\[y=14000\]
\end{minipage}\par\medskip
\noindent\begin{minipage}{\linewidth}\textbf{25. Usamos lo que ya sabemos para hallar x}\\[2pt]Ya sabemos y = 14000; z = 3000. Sustituir significa escribir ese valor cada vez que aparezca su letra. Usamos paréntesis para que se vea qué número estamos multiplicando.
\par\smallskip\textbf{Qué buscamos.} Dejar solamente x como letra desconocida en esta ecuación.
\par\smallskip Tomamos esta condición que obtuvimos antes:
\[x + y + z = 24000\]
\par\smallskip Reemplazamos las letras conocidas por sus valores:
\[x + \left(14000\right) + \left(3000\right)=24000\]
\end{minipage}\par\medskip
\noindent\begin{minipage}{\linewidth}\textbf{26. Calculamos las multiplicaciones que ya se pueden hacer}\\[2pt]La letra x sigue sin conocerse. Las otras partes ya tienen un valor o una expresión conocida y podemos efectuar sus productos.
\par\smallskip {\small\renewcommand{\arraystretch}{1.3}\begin{tabular}{>{\raggedright\arraybackslash}p{3.36cm}>{\raggedright\arraybackslash}p{4.48cm}>{\raggedright\arraybackslash}p{3.36cm}}\hline\rowcolor{black!7} \textbf{Parte conocida} & \textbf{La multiplicación} & \textbf{Resultado}\\\hline El término y & $1\cdot\left(14000\right)$ & $14000$\\ El término z & $1\cdot\left(3000\right)$ & $3000$\\ \hline\end{tabular}}\par\smallskip
\par\smallskip Al reunir esas partes conocidas, la ecuación queda:
\[x + 17000=24000\]
\end{minipage}\par\medskip
\noindent\begin{minipage}{\linewidth}\textbf{27. Dejamos el término con x solo a un lado}\\[2pt]Para quitar la parte conocida, restamos 17000 en AMBOS lados. Piensa en una balanza: si hacemos lo mismo a ambos lados, la igualdad se mantiene.
\par\smallskip Hacemos la misma operación a la izquierda y a la derecha:
\[x + 17000 - 17000=24000 - 17000\]
\par\smallskip Las partes que se compensan dan cero; queda:
\[x=7000\]
\par\smallskip\textbf{Qué aprendimos.} «Pasar al otro lado» es una forma abreviada de decir que hicimos la misma suma o resta en los dos lados.
\end{minipage}\par\medskip
\noindent\begin{minipage}{\linewidth}\textbf{28. Leemos el valor de x}\\[2pt]La letra ya está sola: su número multiplicador es 1, que no se escribe. La igualdad nos dice directamente cuánto vale.
\par\smallskip Anotamos el valor encontrado:
\[x=7000\]
\end{minipage}\par\medskip
\noindent\begin{minipage}{\linewidth}\textbf{29. Reunimos la respuesta; todavía falta comprobarla}\\[2pt]Escribimos todos los valores en el orden x, y, z. Estos valores deben funcionar juntos, no cada uno en una ecuación diferente.
\[\left\{\begin{aligned}x&=7000\\y&=14000\\z&=3000\end{aligned}\right.\]
\par\smallskip\textbf{Qué aprendimos.} El siguiente paso es volver a las ecuaciones del enunciado y comprobar cada condición.
\end{minipage}\par\medskip
\noindent\begin{minipage}{\linewidth}\textbf{30. ¿La respuesta cumple la ecuación 1?}\\[2pt]Comprobar significa volver a la condición original y calcularla con nuestra respuesta. Si los dos lados dan el mismo número, esta condición sí se cumple. Deben cumplirse todas.
\par\smallskip La condición del enunciado era:
\[x + y + z = 24000\]
\par\smallskip Ponemos nuestros valores donde estaban las letras:
\[\left(7000\right) + \left(14000\right) + \left(3000\right)=24000\]
\par\smallskip Hacemos las multiplicaciones y luego la suma:
\[7000 + 14000 + 3000=24000\]
\par\smallskip Comparamos los dos lados:
\[24000=24000\]
\par\smallskip\textbf{Qué aprendimos.} La igualdad final es verdadera. Nuestra respuesta sí cumple la ecuación 1.
\end{minipage}\par\medskip
\noindent\begin{minipage}{\linewidth}\textbf{31. ¿La respuesta cumple la ecuación 2?}\\[2pt]Comprobar significa volver a la condición original y calcularla con nuestra respuesta. Si los dos lados dan el mismo número, esta condición sí se cumple. Deben cumplirse todas.
\par\smallskip La condición del enunciado era:
\[-2x + y = 0\]
\par\smallskip Ponemos nuestros valores donde estaban las letras:
\[-2\left(7000\right) + \left(14000\right)=0\]
\par\smallskip Hacemos las multiplicaciones y luego la suma:
\[-14000 + 14000=0\]
\par\smallskip Comparamos los dos lados:
\[0=0\]
\par\smallskip\textbf{Qué aprendimos.} La igualdad final es verdadera. Nuestra respuesta sí cumple la ecuación 2.
\end{minipage}\par\medskip
\noindent\begin{minipage}{\linewidth}\textbf{32. ¿La respuesta cumple la ecuación 3?}\\[2pt]Comprobar significa volver a la condición original y calcularla con nuestra respuesta. Si los dos lados dan el mismo número, esta condición sí se cumple. Deben cumplirse todas.
\par\smallskip La condición del enunciado era:
\[9x + 10y + 6z = 221000\]
\par\smallskip Ponemos nuestros valores donde estaban las letras:
\[9\left(7000\right) + 10\left(14000\right) + 6\left(3000\right)=221000\]
\par\smallskip Hacemos las multiplicaciones y luego la suma:
\[63000 + 140000 + 18000=221000\]
\par\smallskip Comparamos los dos lados:
\[221000=221000\]
\par\smallskip\textbf{Qué aprendimos.} La igualdad final es verdadera. Nuestra respuesta sí cumple la ecuación 3.
\end{minipage}\par\medskip
\noindent\begin{minipage}{\linewidth}\textbf{33. Interpretamos y verificamos las unidades}\\[2pt]Se invierten \$7,000, \$14,000 y \$3,000, respectivamente. Los intereses son \$630, \$1,400 y \$180; suman \$2,210. El capital total es \$24,000.
\end{minipage}\par\medskip
\respuesta{Primer fideicomiso: \$7,000. Segundo: \$14,000. Tercero: \$3,000. Interés anual total: \$2,210.}
\noindent\begin{minipage}{\linewidth}\centering\includegraphics[width=.96\linewidth,height=6.8cm,keepaspectratio]{../recursos/grafico_1_1-28.png}\\{\small Lectura visual complementaria; la demostración está en las operaciones.}\end{minipage}
\end{multicols*}
\newpage\fancyfoot[L]{\small Sección 1.1 | Ejercicio T.1 | Paso a paso}\hypertarget{1.1-T.1}{}
\pdfbookmark[1]{1.1 / T.1}{bm1.1-T.1}
\titulo{T.1 | Intercambiar dos ecuaciones conserva las soluciones}
\textbf{Enunciado}\par Demuestre que el sistema lineal que se obtiene al intercambiar dos ecuaciones en (2) tiene exactamente las mismas soluciones que (2).
\par\medskip\large
\noindent\begin{minipage}{\linewidth}\textbf{1. Hipótesis y notación del sistema general}\\[2pt]Sean m y n enteros positivos y sean reales todos los coeficientes y términos independientes. Para cada índice i, Li es la suma de los términos de la ecuación i. La letra en negrita representa una lista de n coordenadas. No fijamos valores numéricos: el argumento debe servir para cualquier sistema con estas características.
\par\smallskip\textbf{Qué buscamos.} Probar que el conjunto solución del sistema original S es exactamente el del sistema transformado S'.
\[L_i(\mathbf{x})=\sum_{j=1}^{n}a_{ij}x_j\]
\[S:\quad L_i(\mathbf{x})=b_i\quad (i=1,\ldots,m)\]
\end{minipage}\par\medskip
\noindent\begin{minipage}{\linewidth}\textbf{2. Qué significa demostrar que las soluciones son las mismas}\\[2pt]Llamamos V al conjunto de todas las soluciones de S y W al conjunto de todas las soluciones de S'. El símbolo de pertenencia significa «es elemento de». Para probar la igualdad de conjuntos hay que probar las dos inclusiones: toda solución original sirve en el nuevo sistema y toda solución nueva sirve en el original. Una inclusión sola no basta.
\[V=\{\mathbf{x}\in\mathbb{R}^{n}:L_i(\mathbf{x})=b_i\ (i=1,\ldots,m)\}\]
\[V\subseteq W\quad\mathrm{y}\quad W\subseteq V\quad\Longrightarrow\quad V=W\]
\par\smallskip\textbf{Qué aprendimos.} En lo que sigue elegimos una solución arbitraria. No suponemos que exista alguna: si el conjunto está vacío, las implicaciones siguen siendo válidas.
\end{minipage}\par\medskip
\noindent\begin{minipage}{\linewidth}\textbf{3. Operación: intercambiar dos ecuaciones}\\[2pt]Supongamos m mayor o igual que 2. Fijemos índices p y q distintos. S' se obtiene de S colocando la ecuación q en la posición p y la ecuación p en la posición q. Las demás ecuaciones conservan su lugar. Solo cambia el orden de las condiciones.
\[1\leq p,q\leq m,\quad p\ne q\]
\[\left\{\begin{aligned}L_q(\mathbf{x})&=b_q\\L_p(\mathbf{x})&=b_p\\L_i(\mathbf{x})&=b_i\quad(i\ne p,q)\end{aligned}\right.\]
\end{minipage}\par\medskip
\noindent\begin{minipage}{\linewidth}\textbf{4. Primera inclusión: de V hacia W}\\[2pt]Sea u un elemento arbitrario de V. Por definición cumple todas las ecuaciones originales. En particular cumple las ecuaciones p y q. Por tanto cumple la ecuación que ahora ocupa la posición p, la que ocupa la posición q y todas las que no cambiaron. Así cumple S' y pertenece a W.
\[\mathbf{u}\in V\quad\Longrightarrow\quad L_i(\mathbf{u})=b_i\quad(1\leq i\leq m)\]
\[L_q(\mathbf{u})=b_q,\quad L_p(\mathbf{u})=b_p\]
\[\mathbf{u}\in W\quad\Longrightarrow\quad V\subseteq W\]
\par\smallskip\textbf{Qué aprendimos.} La última inclusión se deduce porque u fue arbitrario: el razonamiento vale para cada elemento de V.
\end{minipage}\par\medskip
\noindent\begin{minipage}{\linewidth}\textbf{5. Segunda inclusión: de W hacia V}\\[2pt]Sea v un elemento arbitrario de W. Cumple las condiciones de S'. La condición en la posición q es precisamente la antigua ecuación p; la condición en la posición p es la antigua ecuación q. Las demás son idénticas a las de S. Por ello v cumple todas las ecuaciones de S y pertenece a V.
\[\mathbf{v}\in W\quad\Longrightarrow\quad L_p(\mathbf{v})=b_p,\quad L_q(\mathbf{v})=b_q\]
\[L_i(\mathbf{v})=b_i\quad(i\ne p,q)\]
\[\mathbf{v}\in V\quad\Longrightarrow\quad W\subseteq V\]
\end{minipage}\par\medskip
\noindent\begin{minipage}{\linewidth}\textbf{6. Conclusión formal}\\[2pt]Por las dos inclusiones y la propiedad de doble inclusión de conjuntos, V y W son iguales. En consecuencia, intercambiar dos ecuaciones conserva exactamente todas las soluciones. La prueba no depende de si hay cero, una o infinitas soluciones. Queda demostrado.
\[V\subseteq W,\quad W\subseteq V\quad\Longrightarrow\quad V=W\]
\par\smallskip\textbf{Qué aprendimos.} El orden en que se escriben las condiciones no modifica cuáles valores las cumplen simultáneamente.
\end{minipage}\par\medskip
\respuesta{Intercambiar dos ecuaciones conserva exactamente el conjunto solución.}
\noindent\begin{minipage}{\linewidth}\centering\includegraphics[width=.96\linewidth,height=6.8cm,keepaspectratio]{../recursos/grafico_1_1-T_1.png}\\{\small Lectura visual complementaria; la demostración está en las operaciones.}\end{minipage}
\normalsize
\newpage\fancyfoot[L]{\small Sección 1.1 | Ejercicio T.2 | Paso a paso}\hypertarget{1.1-T.2}{}
\pdfbookmark[1]{1.1 / T.2}{bm1.1-T.2}
\titulo{T.2 | Multiplicar una ecuación por una constante no nula}
\textbf{Enunciado}\par Demuestre que el sistema lineal obtenido al reemplazar una ecuación en (2) por un múltiplo constante de la ecuación diferente de cero, tiene exactamente las mismas soluciones que (2).
\par\medskip\large
\noindent\begin{minipage}{\linewidth}\textbf{1. Hipótesis y notación del sistema general}\\[2pt]Sean m y n enteros positivos y sean reales todos los coeficientes y términos independientes. Para cada índice i, Li es la suma de los términos de la ecuación i. La letra en negrita representa una lista de n coordenadas. No fijamos valores numéricos: el argumento debe servir para cualquier sistema con estas características.
\par\smallskip\textbf{Qué buscamos.} Probar que el conjunto solución del sistema original S es exactamente el del sistema transformado S'.
\[L_i(\mathbf{x})=\sum_{j=1}^{n}a_{ij}x_j\]
\[S:\quad L_i(\mathbf{x})=b_i\quad (i=1,\ldots,m)\]
\end{minipage}\par\medskip
\noindent\begin{minipage}{\linewidth}\textbf{2. Qué significa demostrar que las soluciones son las mismas}\\[2pt]Llamamos V al conjunto de todas las soluciones de S y W al conjunto de todas las soluciones de S'. El símbolo de pertenencia significa «es elemento de». Para probar la igualdad de conjuntos hay que probar las dos inclusiones: toda solución original sirve en el nuevo sistema y toda solución nueva sirve en el original. Una inclusión sola no basta.
\[V=\{\mathbf{x}\in\mathbb{R}^{n}:L_i(\mathbf{x})=b_i\ (i=1,\ldots,m)\}\]
\[V\subseteq W\quad\mathrm{y}\quad W\subseteq V\quad\Longrightarrow\quad V=W\]
\par\smallskip\textbf{Qué aprendimos.} En lo que sigue elegimos una solución arbitraria. No suponemos que exista alguna: si el conjunto está vacío, las implicaciones siguen siendo válidas.
\end{minipage}\par\medskip
\noindent\begin{minipage}{\linewidth}\textbf{3. Operación y condición necesaria}\\[2pt]Fijemos una ecuación p y un número real $\alpha$ distinto de cero. S' conserva las otras ecuaciones y reemplaza la ecuación p por su producto completo por $\alpha$. «Completo» incluye el lado derecho. La hipótesis $\alpha$ distinto de cero permitirá recuperar la condición original dividiendo entre $\alpha$.
\[\alpha\in\mathbb{R},\quad\alpha\ne0\]
\[\left\{\begin{aligned}\alpha L_p(\mathbf{x})&=\alpha b_p\\L_i(\mathbf{x})&=b_i\quad(i\ne p)\end{aligned}\right.\]
\end{minipage}\par\medskip
\noindent\begin{minipage}{\linewidth}\textbf{4. Primera inclusión: de V hacia W}\\[2pt]Sea u una solución arbitraria de S. Multiplicamos ambos miembros de su ecuación p por $\alpha$. La nueva igualdad es verdadera. Como las otras ecuaciones no cambiaron, u también cumple S'. Esto prueba la primera inclusión.
\[L_p(\mathbf{u})=b_p\]
\[\alpha L_p(\mathbf{u})=\alpha b_p\]
\[\mathbf{u}\in W,\quad V\subseteq W\]
\end{minipage}\par\medskip
\noindent\begin{minipage}{\linewidth}\textbf{5. Segunda inclusión: de W hacia V}\\[2pt]Sea v una solución arbitraria de S'. Su ecuación transformada es verdadera. Como $\alpha$ no es cero, dividimos ambos miembros entre $\alpha$ y recuperamos la ecuación p original. Las otras ecuaciones ya se cumplen. Entonces v es solución de S.
\[\alpha L_p(\mathbf{v})=\alpha b_p\]
\[\frac{\alpha L_p(\mathbf{v})}{\alpha}=\frac{\alpha b_p}{\alpha}\]
\[L_p(\mathbf{v})=b_p\]
\[\mathbf{v}\in V,\quad W\subseteq V\]
\end{minipage}\par\medskip
\noindent\begin{minipage}{\linewidth}\textbf{6. Conclusión formal y necesidad de la hipótesis}\\[2pt]Las dos inclusiones prueban V = W. Si se permitiera $\alpha$ = 0, no podríamos dividir y la afirmación sería falsa en general: la ecuación x = 1 se convertiría en 0 = 0 y admitiría, por ejemplo, x = 2. Así queda demostrada la equivalencia bajo la hipótesis indicada.
\[\alpha\ne0\quad\Longrightarrow\quad V=W\]
\[\alpha=0:\quad x=1\quad\longrightarrow\quad 0=0\]
\par\smallskip\textbf{Qué aprendimos.} Multiplicar por un número real distinto de cero es reversible; multiplicar por cero puede borrar información.
\end{minipage}\par\medskip
\respuesta{Los sistemas son equivalentes porque la operación tiene inversa: dividir por la misma constante no nula.}
\noindent\begin{minipage}{\linewidth}\centering\includegraphics[width=.96\linewidth,height=6.8cm,keepaspectratio]{../recursos/grafico_1_1-T_2.png}\\{\small Lectura visual complementaria; la demostración está en las operaciones.}\end{minipage}
\normalsize
\newpage\fancyfoot[L]{\small Sección 1.1 | Ejercicio T.3 | Paso a paso}\hypertarget{1.1-T.3}{}
\pdfbookmark[1]{1.1 / T.3}{bm1.1-T.3}
\titulo{T.3 | Sumar un múltiplo de otra ecuación}
\textbf{Enunciado}\par Demuestre que el sistema lineal que se obtiene al reemplazar una ecuación en (2) por la misma más un múltiplo de otra ecuación en (2) tiene exactamente las mismas soluciones que (2).
\par\medskip\large
\noindent\begin{minipage}{\linewidth}\textbf{1. Hipótesis y notación del sistema general}\\[2pt]Sean m y n enteros positivos y sean reales todos los coeficientes y términos independientes. Para cada índice i, Li es la suma de los términos de la ecuación i. La letra en negrita representa una lista de n coordenadas. No fijamos valores numéricos: el argumento debe servir para cualquier sistema con estas características.
\par\smallskip\textbf{Qué buscamos.} Probar que el conjunto solución del sistema original S es exactamente el del sistema transformado S'.
\[L_i(\mathbf{x})=\sum_{j=1}^{n}a_{ij}x_j\]
\[S:\quad L_i(\mathbf{x})=b_i\quad (i=1,\ldots,m)\]
\end{minipage}\par\medskip
\noindent\begin{minipage}{\linewidth}\textbf{2. Qué significa demostrar que las soluciones son las mismas}\\[2pt]Llamamos V al conjunto de todas las soluciones de S y W al conjunto de todas las soluciones de S'. El símbolo de pertenencia significa «es elemento de». Para probar la igualdad de conjuntos hay que probar las dos inclusiones: toda solución original sirve en el nuevo sistema y toda solución nueva sirve en el original. Una inclusión sola no basta.
\[V=\{\mathbf{x}\in\mathbb{R}^{n}:L_i(\mathbf{x})=b_i\ (i=1,\ldots,m)\}\]
\[V\subseteq W\quad\mathrm{y}\quad W\subseteq V\quad\Longrightarrow\quad V=W\]
\par\smallskip\textbf{Qué aprendimos.} En lo que sigue elegimos una solución arbitraria. No suponemos que exista alguna: si el conjunto está vacío, las implicaciones siguen siendo válidas.
\end{minipage}\par\medskip
\noindent\begin{minipage}{\linewidth}\textbf{3. Operación: reemplazar una ecuación y conservar la otra}\\[2pt]Supongamos m mayor o igual que 2. Fijemos p y q distintos y $\lambda$ real. Reemplazamos solo la ecuación p por ella más $\lambda$ veces la ecuación q. La ecuación q permanece en S', junto a todas las demás. Esta permanencia es necesaria para el argumento inverso.
\[p\ne q,\quad\lambda\in\mathbb{R}\]
\[\left\{\begin{aligned}L_p(\mathbf{x})+\lambda L_q(\mathbf{x})&=b_p+\lambda b_q\\L_i(\mathbf{x})&=b_i\quad(i\ne p)\end{aligned}\right.\]
\end{minipage}\par\medskip
\noindent\begin{minipage}{\linewidth}\textbf{4. Primera inclusión: de V hacia W}\\[2pt]Sea u una solución arbitraria de S. Satisface las ecuaciones p y q. Multiplicamos la segunda igualdad por $\lambda$ y la sumamos a la primera, miembro a miembro. Resulta la condición transformada. Las otras ecuaciones se conservan; luego u pertenece a W.
\[L_p(\mathbf{u})=b_p,\quad L_q(\mathbf{u})=b_q\]
\[L_p(\mathbf{u})+\lambda L_q(\mathbf{u})=b_p+\lambda b_q\]
\[\mathbf{u}\in W,\quad V\subseteq W\]
\end{minipage}\par\medskip
\noindent\begin{minipage}{\linewidth}\textbf{5. Segunda inclusión: recuperar la ecuación original}\\[2pt]Sea v una solución arbitraria de S'. Cumple la nueva ecuación p y también la ecuación q, que no cambió. Restamos $\lambda$ veces esta última igualdad a la nueva ecuación. Los términos añadidos se cancelan. Obtenemos la ecuación p original; las demás ya se cumplían. Por ello v pertenece a V.
\[L_p(\mathbf{v})+\lambda L_q(\mathbf{v})=b_p+\lambda b_q\]
\[L_q(\mathbf{v})=b_q\]
\[L_p(\mathbf{v})+\lambda L_q(\mathbf{v})-\lambda L_q(\mathbf{v})=b_p+\lambda b_q-\lambda b_q\]
\[L_p(\mathbf{v})=b_p\]
\[\mathbf{v}\in V,\quad W\subseteq V\]
\end{minipage}\par\medskip
\noindent\begin{minipage}{\linewidth}\textbf{6. Conclusión formal}\\[2pt]Por doble inclusión, los conjuntos solución coinciden. La operación inversa consiste en restar $\lambda$ veces la ecuación q a la nueva ecuación p. Es válida para cualquier $\lambda$ real, incluso cero; aquí no dividimos entre $\lambda$. Queda demostrado.
\[V\subseteq W,\quad W\subseteq V\quad\Longrightarrow\quad V=W\]
\par\smallskip\textbf{Qué aprendimos.} La nueva ecuación puede representar otra recta o plano, pero la intersección con todas las condiciones conservadas es exactamente la misma.
\end{minipage}\par\medskip
\respuesta{Sumar a una ecuación un múltiplo de otra distinta, conservando esta última, mantiene exactamente las soluciones.}
\noindent\begin{minipage}{\linewidth}\centering\includegraphics[width=.96\linewidth,height=6.8cm,keepaspectratio]{../recursos/grafico_1_1-T_3.png}\\{\small Lectura visual complementaria; la demostración está en las operaciones.}\end{minipage}
\normalsize
\newpage\fancyfoot[L]{\small Sección 1.1 | Ejercicio T.4 | Paso a paso}\hypertarget{1.1-T.4}{}
\pdfbookmark[1]{1.1 / T.4}{bm1.1-T.4}
\titulo{T.4 | Todo sistema homogéneo tiene solución}
\textbf{Enunciado}\par ¿El sistema lineal ax + by = 0, cx + dy = 0 siempre tiene solución para cualesquiera valores de a, b, c y d?
\[\left\{\begin{aligned}ax+by&=0\\cx+dy&=0\end{aligned}\right.\]
\par\medskip\large
\noindent\begin{minipage}{\linewidth}\textbf{1. Proposición e hipótesis}\\[2pt]Debemos decidir si, para cualquier elección de cuatro coeficientes reales, existe al menos una pareja real que cumpla las dos ecuaciones. «Para cualquier» se refiere a los coeficientes; «existe» se refiere a la pareja de incógnitas. La pareja puede ser la misma para todos los coeficientes.
\par\smallskip\textbf{Qué buscamos.} Dar una solución y verificarla sin imponer condiciones adicionales a los coeficientes.
\[\forall\,(a,b,c,d)\in\mathbb{R}^{4},\quad\exists\,(x,y)\in\mathbb{R}^{2}\]
\[\left\{\begin{aligned}ax+by&=0\\cx+dy&=0\end{aligned}\right.\]
\end{minipage}\par\medskip
\noindent\begin{minipage}{\linewidth}\textbf{2. Demostración: elegimos coeficientes arbitrarios}\\[2pt]Sean a, b, c y d números reales arbitrarios. Proponemos x = 0 e y = 0. Esta pareja pertenece a R\textsuperscript{2}. Por la propiedad del cero en la multiplicación, cualquier coeficiente real multiplicado por cero da cero.
\[(x,y)=(0,0)\in\mathbb{R}^{2}\]
\end{minipage}\par\medskip
\noindent\begin{minipage}{\linewidth}\textbf{3. Verificación de ambas ecuaciones}\\[2pt]Sustituimos la misma pareja en las dos condiciones. Cada igualdad se reduce a 0 = 0, así que ambas son verdaderas para los coeficientes arbitrarios que elegimos.
\[a\cdot0+b\cdot0=0+0=0\]
\[c\cdot0+d\cdot0=0+0=0\]
\par\smallskip\textbf{Qué aprendimos.} Por definición de solución de un sistema, (0, 0) es una solución.
\end{minipage}\par\medskip
\noindent\begin{minipage}{\linewidth}\textbf{4. Conclusión con los cuantificadores}\\[2pt]Como no restringimos los coeficientes, lo probado vale para toda elección real de a, b, c y d. Hemos exhibido una pareja que cumple el sistema, de modo que siempre existe solución. Queda demostrado.
\[\forall\,(a,b,c,d)\in\mathbb{R}^{4},\quad a\cdot0+b\cdot0=c\cdot0+d\cdot0=0\]
\par\smallskip\textbf{Qué aprendimos.} La existencia queda probada. No hemos afirmado unicidad: si las dos ecuaciones son x + y = 0 y 2x + 2y = 0, todas las parejas (t, -t), con t real, son soluciones.
\end{minipage}\par\medskip
\respuesta{Sí. (x, y) = (0, 0) siempre es solución, para cualquier a, b, c y d.}
\noindent\begin{minipage}{\linewidth}\centering\includegraphics[width=.96\linewidth,height=6.8cm,keepaspectratio]{../recursos/grafico_1_1-T_4.png}\\{\small Lectura visual complementaria; la demostración está en las operaciones.}\end{minipage}
\normalsize
\newpage\fancyfoot[L]{\small Sección 1.2 | Ejercicio 1 | Paso a paso}\hypertarget{1.2-1}{}
\pdfbookmark[1]{1.2 / 1}{bm1.2-1}
\titulo{1 | Leer una entrada de una matriz}
\textbf{Enunciado}\par Sean las matrices A, B y C indicadas. (a) Halle a12, a22 y a23. (b) Halle b11 y b31. (c) Halle c13, c31 y c33.
\par\smallskip \begin{minipage}[t]{.47\linewidth}\centering \textbf{A}\[\begin{bmatrix}2 & -3 & 5\\6 & -5 & 4\end{bmatrix}\]\end{minipage}\hfill\begin{minipage}[t]{.47\linewidth}\centering \textbf{B}\[\begin{bmatrix}4\\-3\\5\end{bmatrix}\]\end{minipage}\par\smallskip
\par\smallskip \begin{minipage}[t]{.47\linewidth}\centering \textbf{C}\[\begin{bmatrix}7 & 3 & 2\\-4 & 3 & 5\\6 & 1 & -1\end{bmatrix}\]\end{minipage}\par\smallskip
\begin{multicols*}{2}
\noindent\begin{minipage}{\linewidth}\textbf{1. Cómo se leen los dos índices}\\[2pt]El primer índice indica la fila, que se cuenta de arriba abajo; el segundo indica la columna, que se cuenta de izquierda a derecha. Se empieza a contar en 1. Así, a12 significa la entrada de A en la fila 1 y la columna 2.
\[a_{ij}=\mathrm{entrada\ de\ fila}\ i\ \mathrm{y\ columna}\ j\]
\end{minipage}\par\medskip
\noindent\begin{minipage}{\linewidth}\textbf{2. Inciso (a): nos ubicamos en A}\\[2pt]Buscamos primero la fila y después avanzamos hasta la columna.
\par\smallskip {\small\renewcommand{\arraystretch}{1.3}\begin{tabular}{>{\raggedright\arraybackslash}p{3.74cm}>{\raggedright\arraybackslash}p{3.74cm}>{\raggedright\arraybackslash}p{3.74cm}}\hline\rowcolor{black!7} \textbf{Entrada} & \textbf{Ubicación} & \textbf{Valor}\\\hline a12 & Fila 1, columna 2 & -3\\ a22 & Fila 2, columna 2 & -5\\ a23 & Fila 2, columna 3 & 4\\ \hline\end{tabular}}\par\smallskip
\end{minipage}\par\medskip
\noindent\begin{minipage}{\linewidth}\textbf{3. Inciso (b): B solo tiene una columna}\\[2pt]B es una matriz de 3 filas y 1 columna. El segundo índice solo puede ser 1.
\[b_{11}=4,\quad b_{31}=5\]
\end{minipage}\par\medskip
\noindent\begin{minipage}{\linewidth}\textbf{4. Inciso (c): buscamos en C}\\[2pt]La esquina superior derecha es c13; la inferior izquierda es c31; la inferior derecha es c33.
\[c_{13}=2,\quad c_{31}=6,\quad c_{33}=-1\]
\end{minipage}\par\medskip
\respuesta{a12 = -3, a22 = -5, a23 = 4; b11 = 4, b31 = 5; c13 = 2, c31 = 6, c33 = -1.}
\noindent\begin{minipage}{\linewidth}\centering\includegraphics[width=.96\linewidth,height=6.8cm,keepaspectratio]{../recursos/grafico_1_2-1.png}\\{\small Lectura visual complementaria; la demostración está en las operaciones.}\end{minipage}
\end{multicols*}
\newpage\fancyfoot[L]{\small Sección 1.2 | Ejercicio 2 | Paso a paso}\hypertarget{1.2-2}{}
\pdfbookmark[1]{1.2 / 2}{bm1.2-2}
\titulo{2 | Igualdad de matrices: hallar a, b, c y d}
\textbf{Enunciado}\par Si las dos matrices indicadas son iguales, determine a, b, c y d. La igualdad de matrices exige igualdad entrada por entrada.
\par\smallskip \begin{minipage}[t]{.47\linewidth}\centering \textbf{Lado izquierdo}\[\begin{bmatrix}a + b & c + d\\c - d & a - b\end{bmatrix}\]\end{minipage}\hfill\begin{minipage}[t]{.47\linewidth}\centering \textbf{Lado derecho}\[\begin{bmatrix}4 & 6\\10 & 2\end{bmatrix}\]\end{minipage}\par\smallskip
\begin{multicols*}{2}
\noindent\begin{minipage}{\linewidth}\textbf{1. Igualamos casillas que ocupan la misma posición}\\[2pt]Los dos tamaños son 2 por 2. Por definición de igualdad, la casilla superior izquierda se iguala con la superior izquierda, y lo mismo para las otras tres. Así obtenemos cuatro ecuaciones.
\[\left\{\begin{aligned}a+b&=4\\c+d&=6\\c-d&=10\\a-b&=2\end{aligned}\right.\]
\end{minipage}\par\medskip
\noindent\begin{minipage}{\linewidth}\textbf{2. Primero: ¿qué significa resolver este sistema?}\\[2pt]Las letras a, b, c, d representan números que todavía no conocemos. Una solución es una elección de esos números que hace verdaderas TODAS las ecuaciones al mismo tiempo. La llave de la izquierda significa que las condiciones se deben cumplir juntas.
\par\smallskip\textbf{Qué buscamos.} Encontrar un mismo grupo de valores para las 4 letras.
\par\smallskip El número 3 delante de a indica tres veces ese mismo valor; es decir, sumar tres copias:
\[3a=a+a+a\]
\par\smallskip Ejemplo numérico de la notación: si a valiera 2, tres veces ese valor sería 6. Este 2 es solo un ejemplo; todavía no es la respuesta del ejercicio. El punto centrado indica multiplicación:
\[3\cdot 2=2+2+2=6\]
\par\smallskip\textbf{Comprueba que lo entiendes.} ¿Basta con que los valores funcionen en una sola ecuación?\par\textit{Respuesta:} No. Deben funcionar en todas las ecuaciones del sistema, usando los mismos valores.
\end{minipage}\par\medskip
\noindent\begin{minipage}{\linewidth}\textbf{3. Nuestro plan: ir quitando letras}\\[2pt]Una ecuación con una sola letra es más fácil de resolver. Vamos a combinar ecuaciones para que una letra desaparezca. Cuando sepamos un valor, lo pondremos en las ecuaciones que faltan.
\par\smallskip\textbf{Qué buscamos.} Primero quitar a de las otras ecuaciones. Después seguir con las letras restantes.
\par\smallskip\textbf{Por qué.} Si una cantidad se resta de sí misma, queda cero. Usaremos esta idea con términos que tengan la misma letra.
\par\smallskip Ejemplo de la idea que utilizaremos:
\[3x-3x=0\]
\par\smallskip\textbf{Qué aprendimos.} Eliminar una letra de una ecuación NO significa que su valor sea cero; significa que sus términos se cancelaron en esa cuenta.
\end{minipage}\par\medskip
\noindent\begin{minipage}{\linewidth}\textbf{4. ¿Cómo haremos desaparecer a?}\\[2pt]Trabajaremos con las ecuaciones 1 y 2. Queremos que los términos con a tengan el mismo tamaño. Así podremos restarlos.
\par\smallskip\textbf{Qué buscamos.} Conseguir términos con 1 veces a en las dos ecuaciones.
\par\smallskip\textbf{Por qué.} Tener 1 veces la misma letra en ambos lados de la resta nos permitirá cancelarla.
\par\smallskip {\small\renewcommand{\arraystretch}{1.3}\begin{tabular}{>{\raggedright\arraybackslash}p{2.58cm}>{\raggedright\arraybackslash}p{2.47cm}>{\raggedright\arraybackslash}p{6.17cm}}\hline\rowcolor{black!7} \textbf{Dónde miramos} & \textbf{Qué término hay} & \textbf{Qué necesitamos hacer}\\\hline Ecuación 1 & $a$ & La dejamos igual.\\ Ecuación 2 & $a$ & La dejamos igual.\\ \hline\end{tabular}}\par\smallskip
\par\smallskip\textbf{Qué aprendimos.} Estos multiplicadores se eligen para igualar los términos con a. Los demás términos también se multiplican para conservar la igualdad.
\end{minipage}\par\medskip
\noindent\begin{minipage}{\linewidth}\textbf{5. Restamos las dos ecuaciones preparadas}\\[2pt]Restamos la ecuación 1 preparada a la ecuación 2 preparada. Juntamos cada término solo con otro que tenga la misma letra. Los números sin letra se calculan por separado.
\par\smallskip {\small\renewcommand{\arraystretch}{1.3}\begin{tabular}{>{\raggedright\arraybackslash}p{2.80cm}>{\raggedright\arraybackslash}p{5.38cm}>{\raggedright\arraybackslash}p{3.03cm}}\hline\rowcolor{black!7} \textbf{Qué agrupamos} & \textbf{Cuenta} & \textbf{Qué significa}\\\hline Los términos con a & $a - \left(a\right)=0$ & Se cancela esta letra.\\ Los términos con b & $-b - \left(b\right)=-2b$ & Quedan -2b.\\ Los números a la derecha del igual & $2 - \left(4\right)=-2$ & Este será el número de la nueva igualdad.\\ \hline\end{tabular}}\par\smallskip
\par\smallskip Primera igualdad que usaremos:
\[a - b = 2\]
\par\smallskip Le restamos esta igualdad:
\[a + b = 4\]
\par\smallskip\textbf{Qué aprendimos.} Los términos con a dan cero. Los términos con letras distintas no se mezclan.
\end{minipage}\par\medskip
\noindent\begin{minipage}{\linewidth}\textbf{6. Mira lo que hemos conseguido}\\[2pt]La nueva ecuación es el resultado de la cuenta anterior. Solo contiene b: después podremos despejarla.
\par\smallskip Escribimos la nueva condición:
\[-2b = -2\]
\[\left\{\begin{aligned}a + b &= 4\\-2b &= -2\\c + d &= 6\\c - d &= 10\end{aligned}\right.\]
\par\smallskip\textbf{Qué aprendimos.} Conservamos la ecuación 1 y reemplazamos la 2 por esta nueva condición. De esa manera conservamos toda la información necesaria para resolver el sistema.
\end{minipage}\par\medskip
\noindent\begin{minipage}{\linewidth}\textbf{7. ¿Cómo haremos desaparecer c?}\\[2pt]Trabajaremos con las ecuaciones 3 y 4. Queremos que los términos con c tengan el mismo tamaño. Así podremos restarlos.
\par\smallskip\textbf{Qué buscamos.} Conseguir términos con 1 veces c en las dos ecuaciones.
\par\smallskip\textbf{Por qué.} Tener 1 veces la misma letra en ambos lados de la resta nos permitirá cancelarla.
\par\smallskip {\small\renewcommand{\arraystretch}{1.3}\begin{tabular}{>{\raggedright\arraybackslash}p{2.58cm}>{\raggedright\arraybackslash}p{2.47cm}>{\raggedright\arraybackslash}p{6.17cm}}\hline\rowcolor{black!7} \textbf{Dónde miramos} & \textbf{Qué término hay} & \textbf{Qué necesitamos hacer}\\\hline Ecuación 3 & $c$ & La dejamos igual.\\ Ecuación 4 & $c$ & La dejamos igual.\\ \hline\end{tabular}}\par\smallskip
\par\smallskip\textbf{Qué aprendimos.} Estos multiplicadores se eligen para igualar los términos con c. Los demás términos también se multiplican para conservar la igualdad.
\end{minipage}\par\medskip
\noindent\begin{minipage}{\linewidth}\textbf{8. Restamos las dos ecuaciones preparadas}\\[2pt]Restamos la ecuación 3 preparada a la ecuación 4 preparada. Juntamos cada término solo con otro que tenga la misma letra. Los números sin letra se calculan por separado.
\par\smallskip {\small\renewcommand{\arraystretch}{1.3}\begin{tabular}{>{\raggedright\arraybackslash}p{2.80cm}>{\raggedright\arraybackslash}p{5.38cm}>{\raggedright\arraybackslash}p{3.03cm}}\hline\rowcolor{black!7} \textbf{Qué agrupamos} & \textbf{Cuenta} & \textbf{Qué significa}\\\hline Los términos con c & $c - \left(c\right)=0$ & Se cancela esta letra.\\ Los términos con d & $-d - \left(d\right)=-2d$ & Quedan -2d.\\ Los números a la derecha del igual & $10 - \left(6\right)=4$ & Este será el número de la nueva igualdad.\\ \hline\end{tabular}}\par\smallskip
\par\smallskip Primera igualdad que usaremos:
\[c - d = 10\]
\par\smallskip Le restamos esta igualdad:
\[c + d = 6\]
\par\smallskip\textbf{Qué aprendimos.} Los términos con c dan cero. Los términos con letras distintas no se mezclan.
\end{minipage}\par\medskip
\noindent\begin{minipage}{\linewidth}\textbf{9. Mira lo que hemos conseguido}\\[2pt]La nueva ecuación es el resultado de la cuenta anterior. Solo contiene d: después podremos despejarla.
\par\smallskip Escribimos la nueva condición:
\[-2d = 4\]
\[\left\{\begin{aligned}a + b &= 4\\-2b &= -2\\c + d &= 6\\-2d &= 4\end{aligned}\right.\]
\par\smallskip\textbf{Qué aprendimos.} Conservamos la ecuación 3 y reemplazamos la 4 por esta nueva condición. De esa manera conservamos toda la información necesaria para resolver el sistema.
\end{minipage}\par\medskip
\noindent\begin{minipage}{\linewidth}\textbf{10. Ya tenemos una ecuación con una sola letra}\\[2pt]En esta condición solamente falta conocer d. El número -2 está multiplicando a esa letra. Ahora podemos deshacer esa multiplicación.
\par\smallskip Esta es la condición que vamos a resolver:
\[-2d = 4\]
\end{minipage}\par\medskip
\noindent\begin{minipage}{\linewidth}\textbf{11. Calculamos cuánto vale d}\\[2pt]Para deshacer una multiplicación por -2, dividimos AMBOS lados entre -2. En el lado izquierdo queda una sola d. En el derecho hacemos la división.
\par\smallskip Dividimos los dos lados entre el mismo número:
\[\frac{-2d}{-2}=\frac{4}{-2}\]
\par\smallskip Ya podemos leer el valor de la letra:
\[d=\frac{4}{-2}=-2\]
\par\smallskip\textbf{Qué aprendimos.} Este valor se usará en las ecuaciones donde todavía falten letras.
\end{minipage}\par\medskip
\noindent\begin{minipage}{\linewidth}\textbf{12. Usamos lo que ya sabemos para hallar c}\\[2pt]Ya sabemos d = -2. Sustituir significa escribir ese valor cada vez que aparezca su letra. Usamos paréntesis para que se vea qué número estamos multiplicando.
\par\smallskip\textbf{Qué buscamos.} Dejar solamente c como letra desconocida en esta ecuación.
\par\smallskip Tomamos esta condición que obtuvimos antes:
\[c + d = 6\]
\par\smallskip Reemplazamos las letras conocidas por sus valores:
\[c + \left(-2\right)=6\]
\end{minipage}\par\medskip
\noindent\begin{minipage}{\linewidth}\textbf{13. Calculamos las multiplicaciones que ya se pueden hacer}\\[2pt]La letra c sigue sin conocerse. Las otras partes ya tienen un valor o una expresión conocida y podemos efectuar sus productos.
\par\smallskip {\small\renewcommand{\arraystretch}{1.3}\begin{tabular}{>{\raggedright\arraybackslash}p{3.36cm}>{\raggedright\arraybackslash}p{4.48cm}>{\raggedright\arraybackslash}p{3.36cm}}\hline\rowcolor{black!7} \textbf{Parte conocida} & \textbf{La multiplicación} & \textbf{Resultado}\\\hline El término d & $1\cdot\left(-2\right)$ & $-2$\\ \hline\end{tabular}}\par\smallskip
\par\smallskip Al reunir esas partes conocidas, la ecuación queda:
\[c - 2=6\]
\end{minipage}\par\medskip
\noindent\begin{minipage}{\linewidth}\textbf{14. Dejamos el término con c solo a un lado}\\[2pt]Para quitar la parte conocida, sumamos 2 en AMBOS lados. Piensa en una balanza: si hacemos lo mismo a ambos lados, la igualdad se mantiene.
\par\smallskip Hacemos la misma operación a la izquierda y a la derecha:
\[c - 2 + 2=6 + 2\]
\par\smallskip Las partes que se compensan dan cero; queda:
\[c=8\]
\par\smallskip\textbf{Qué aprendimos.} «Pasar al otro lado» es una forma abreviada de decir que hicimos la misma suma o resta en los dos lados.
\end{minipage}\par\medskip
\noindent\begin{minipage}{\linewidth}\textbf{15. Leemos el valor de c}\\[2pt]La letra ya está sola: su número multiplicador es 1, que no se escribe. La igualdad nos dice directamente cuánto vale.
\par\smallskip Anotamos el valor encontrado:
\[c=8\]
\end{minipage}\par\medskip
\noindent\begin{minipage}{\linewidth}\textbf{16. Ya tenemos una ecuación con una sola letra}\\[2pt]En esta condición solamente falta conocer b. El número -2 está multiplicando a esa letra. Ahora podemos deshacer esa multiplicación.
\par\smallskip Esta es la condición que vamos a resolver:
\[-2b = -2\]
\end{minipage}\par\medskip
\noindent\begin{minipage}{\linewidth}\textbf{17. Calculamos cuánto vale b}\\[2pt]Para deshacer una multiplicación por -2, dividimos AMBOS lados entre -2. En el lado izquierdo queda una sola b. En el derecho hacemos la división.
\par\smallskip Dividimos los dos lados entre el mismo número:
\[\frac{-2b}{-2}=\frac{-2}{-2}\]
\par\smallskip Ya podemos leer el valor de la letra:
\[b=\frac{-2}{-2}=1\]
\par\smallskip\textbf{Qué aprendimos.} Un número negativo dividido entre otro negativo da positivo.
\end{minipage}\par\medskip
\noindent\begin{minipage}{\linewidth}\textbf{18. Usamos lo que ya sabemos para hallar a}\\[2pt]Ya sabemos b = 1. Sustituir significa escribir ese valor cada vez que aparezca su letra. Usamos paréntesis para que se vea qué número estamos multiplicando.
\par\smallskip\textbf{Qué buscamos.} Dejar solamente a como letra desconocida en esta ecuación.
\par\smallskip Tomamos esta condición que obtuvimos antes:
\[a + b = 4\]
\par\smallskip Reemplazamos las letras conocidas por sus valores:
\[a + \left(1\right)=4\]
\end{minipage}\par\medskip
\noindent\begin{minipage}{\linewidth}\textbf{19. Calculamos las multiplicaciones que ya se pueden hacer}\\[2pt]La letra a sigue sin conocerse. Las otras partes ya tienen un valor o una expresión conocida y podemos efectuar sus productos.
\par\smallskip {\small\renewcommand{\arraystretch}{1.3}\begin{tabular}{>{\raggedright\arraybackslash}p{3.36cm}>{\raggedright\arraybackslash}p{4.48cm}>{\raggedright\arraybackslash}p{3.36cm}}\hline\rowcolor{black!7} \textbf{Parte conocida} & \textbf{La multiplicación} & \textbf{Resultado}\\\hline El término b & $1\cdot\left(1\right)$ & $1$\\ \hline\end{tabular}}\par\smallskip
\par\smallskip Al reunir esas partes conocidas, la ecuación queda:
\[a + 1=4\]
\end{minipage}\par\medskip
\noindent\begin{minipage}{\linewidth}\textbf{20. Dejamos el término con a solo a un lado}\\[2pt]Para quitar la parte conocida, restamos 1 en AMBOS lados. Piensa en una balanza: si hacemos lo mismo a ambos lados, la igualdad se mantiene.
\par\smallskip Hacemos la misma operación a la izquierda y a la derecha:
\[a + 1 - 1=4 - 1\]
\par\smallskip Las partes que se compensan dan cero; queda:
\[a=3\]
\par\smallskip\textbf{Qué aprendimos.} «Pasar al otro lado» es una forma abreviada de decir que hicimos la misma suma o resta en los dos lados.
\end{minipage}\par\medskip
\noindent\begin{minipage}{\linewidth}\textbf{21. Leemos el valor de a}\\[2pt]La letra ya está sola: su número multiplicador es 1, que no se escribe. La igualdad nos dice directamente cuánto vale.
\par\smallskip Anotamos el valor encontrado:
\[a=3\]
\end{minipage}\par\medskip
\noindent\begin{minipage}{\linewidth}\textbf{22. Reunimos la respuesta; todavía falta comprobarla}\\[2pt]Escribimos todos los valores en el orden a, b, c, d. Estos valores deben funcionar juntos, no cada uno en una ecuación diferente.
\[\left\{\begin{aligned}a&=3\\b&=1\\c&=8\\d&=-2\end{aligned}\right.\]
\par\smallskip\textbf{Qué aprendimos.} El siguiente paso es volver a las ecuaciones del enunciado y comprobar cada condición.
\end{minipage}\par\medskip
\noindent\begin{minipage}{\linewidth}\textbf{23. ¿La respuesta cumple la ecuación 1?}\\[2pt]Comprobar significa volver a la condición original y calcularla con nuestra respuesta. Si los dos lados dan el mismo número, esta condición sí se cumple. Deben cumplirse todas.
\par\smallskip La condición del enunciado era:
\[a + b = 4\]
\par\smallskip Ponemos nuestros valores donde estaban las letras:
\[\left(3\right) + \left(1\right)=4\]
\par\smallskip Hacemos las multiplicaciones y luego la suma:
\[3 + 1=4\]
\par\smallskip Comparamos los dos lados:
\[4=4\]
\par\smallskip\textbf{Qué aprendimos.} La igualdad final es verdadera. Nuestra respuesta sí cumple la ecuación 1.
\end{minipage}\par\medskip
\noindent\begin{minipage}{\linewidth}\textbf{24. ¿La respuesta cumple la ecuación 2?}\\[2pt]Comprobar significa volver a la condición original y calcularla con nuestra respuesta. Si los dos lados dan el mismo número, esta condición sí se cumple. Deben cumplirse todas.
\par\smallskip La condición del enunciado era:
\[a - b = 2\]
\par\smallskip Ponemos nuestros valores donde estaban las letras:
\[\left(3\right) - \left(1\right)=2\]
\par\smallskip Hacemos las multiplicaciones y luego la suma:
\[3 - 1=2\]
\par\smallskip Comparamos los dos lados:
\[2=2\]
\par\smallskip\textbf{Qué aprendimos.} La igualdad final es verdadera. Nuestra respuesta sí cumple la ecuación 2.
\end{minipage}\par\medskip
\noindent\begin{minipage}{\linewidth}\textbf{25. ¿La respuesta cumple la ecuación 3?}\\[2pt]Comprobar significa volver a la condición original y calcularla con nuestra respuesta. Si los dos lados dan el mismo número, esta condición sí se cumple. Deben cumplirse todas.
\par\smallskip La condición del enunciado era:
\[c + d = 6\]
\par\smallskip Ponemos nuestros valores donde estaban las letras:
\[\left(8\right) + \left(-2\right)=6\]
\par\smallskip Hacemos las multiplicaciones y luego la suma:
\[8 - 2=6\]
\par\smallskip Comparamos los dos lados:
\[6=6\]
\par\smallskip\textbf{Qué aprendimos.} La igualdad final es verdadera. Nuestra respuesta sí cumple la ecuación 3.
\end{minipage}\par\medskip
\noindent\begin{minipage}{\linewidth}\textbf{26. ¿La respuesta cumple la ecuación 4?}\\[2pt]Comprobar significa volver a la condición original y calcularla con nuestra respuesta. Si los dos lados dan el mismo número, esta condición sí se cumple. Deben cumplirse todas.
\par\smallskip La condición del enunciado era:
\[c - d = 10\]
\par\smallskip Ponemos nuestros valores donde estaban las letras:
\[\left(8\right) - \left(-2\right)=10\]
\par\smallskip Hacemos las multiplicaciones y luego la suma:
\[8 + 2=10\]
\par\smallskip Comparamos los dos lados:
\[10=10\]
\par\smallskip\textbf{Qué aprendimos.} La igualdad final es verdadera. Nuestra respuesta sí cumple la ecuación 4.
\end{minipage}\par\medskip
\noindent\begin{minipage}{\linewidth}\textbf{27. Volvemos a las matrices originales}\\[2pt]Sustituimos los cuatro valores calculados en cada entrada de la matriz izquierda. Las cuatro casillas deben coincidir con la matriz derecha.
\par\smallskip \begin{minipage}[t]{.47\linewidth}\centering \textbf{Izquierda después de sustituir}\[\begin{bmatrix}4 & 6\\10 & 2\end{bmatrix}\]\end{minipage}\hfill\begin{minipage}[t]{.47\linewidth}\centering \textbf{Derecha}\[\begin{bmatrix}4 & 6\\10 & 2\end{bmatrix}\]\end{minipage}\par\smallskip
\end{minipage}\par\medskip
\respuesta{a = 3, b = 1, c = 8, d = -2.}
\noindent\begin{minipage}{\linewidth}\centering\includegraphics[width=.96\linewidth,height=6.8cm,keepaspectratio]{../recursos/grafico_1_2-2.png}\\{\small Lectura visual complementaria; la demostración está en las operaciones.}\end{minipage}
\end{multicols*}
\newpage\fancyfoot[L]{\small Sección 1.2 | Ejercicio 3 | Paso a paso}\hypertarget{1.2-3}{}
\pdfbookmark[1]{1.2 / 3}{bm1.2-3}
\titulo{3 | Igualdad de matrices: hallar a, b, c y d}
\textbf{Enunciado}\par Si las dos matrices indicadas son iguales, determine a, b, c y d. La igualdad de matrices exige igualdad entrada por entrada.
\par\smallskip \begin{minipage}[t]{.47\linewidth}\centering \textbf{Lado izquierdo}\[\begin{bmatrix}a + 2 b & 2 a - b\\2 c + d & c - 2 d\end{bmatrix}\]\end{minipage}\hfill\begin{minipage}[t]{.47\linewidth}\centering \textbf{Lado derecho}\[\begin{bmatrix}4 & -2\\4 & -3\end{bmatrix}\]\end{minipage}\par\smallskip
\begin{multicols*}{2}
\noindent\begin{minipage}{\linewidth}\textbf{1. Igualamos casillas que ocupan la misma posición}\\[2pt]Los dos tamaños son 2 por 2. Por definición de igualdad, la casilla superior izquierda se iguala con la superior izquierda, y lo mismo para las otras tres. Así obtenemos cuatro ecuaciones.
\[\left\{\begin{aligned}a+2b&=4\\2a-b&=-2\\2c+d&=4\\c-2d&=-3\end{aligned}\right.\]
\end{minipage}\par\medskip
\noindent\begin{minipage}{\linewidth}\textbf{2. Primero: ¿qué significa resolver este sistema?}\\[2pt]Las letras a, b, c, d representan números que todavía no conocemos. Una solución es una elección de esos números que hace verdaderas TODAS las ecuaciones al mismo tiempo. La llave de la izquierda significa que las condiciones se deben cumplir juntas.
\par\smallskip\textbf{Qué buscamos.} Encontrar un mismo grupo de valores para las 4 letras.
\par\smallskip El número 3 delante de a indica tres veces ese mismo valor; es decir, sumar tres copias:
\[3a=a+a+a\]
\par\smallskip Ejemplo numérico de la notación: si a valiera 2, tres veces ese valor sería 6. Este 2 es solo un ejemplo; todavía no es la respuesta del ejercicio. El punto centrado indica multiplicación:
\[3\cdot 2=2+2+2=6\]
\par\smallskip\textbf{Comprueba que lo entiendes.} ¿Basta con que los valores funcionen en una sola ecuación?\par\textit{Respuesta:} No. Deben funcionar en todas las ecuaciones del sistema, usando los mismos valores.
\end{minipage}\par\medskip
\noindent\begin{minipage}{\linewidth}\textbf{3. Nuestro plan: ir quitando letras}\\[2pt]Una ecuación con una sola letra es más fácil de resolver. Vamos a combinar ecuaciones para que una letra desaparezca. Cuando sepamos un valor, lo pondremos en las ecuaciones que faltan.
\par\smallskip\textbf{Qué buscamos.} Primero quitar a de las otras ecuaciones. Después seguir con las letras restantes.
\par\smallskip\textbf{Por qué.} Si una cantidad se resta de sí misma, queda cero. Usaremos esta idea con términos que tengan la misma letra.
\par\smallskip Ejemplo de la idea que utilizaremos:
\[3x-3x=0\]
\par\smallskip\textbf{Qué aprendimos.} Eliminar una letra de una ecuación NO significa que su valor sea cero; significa que sus términos se cancelaron en esa cuenta.
\end{minipage}\par\medskip
\noindent\begin{minipage}{\linewidth}\textbf{4. ¿Cómo haremos desaparecer a?}\\[2pt]Trabajaremos con las ecuaciones 1 y 2. Queremos que los términos con a tengan el mismo tamaño. Así podremos restarlos.
\par\smallskip\textbf{Qué buscamos.} Conseguir términos con 2 veces a en las dos ecuaciones.
\par\smallskip\textbf{Por qué.} Tener 2 veces la misma letra en ambos lados de la resta nos permitirá cancelarla.
\par\smallskip {\small\renewcommand{\arraystretch}{1.3}\begin{tabular}{>{\raggedright\arraybackslash}p{2.58cm}>{\raggedright\arraybackslash}p{2.47cm}>{\raggedright\arraybackslash}p{6.17cm}}\hline\rowcolor{black!7} \textbf{Dónde miramos} & \textbf{Qué término hay} & \textbf{Qué necesitamos hacer}\\\hline Ecuación 1 & $a$ & Multiplicar TODA esta ecuación por 2\\ Ecuación 2 & $2a$ & La dejamos igual.\\ \hline\end{tabular}}\par\smallskip
\par\smallskip\textbf{Qué aprendimos.} Estos multiplicadores se eligen para igualar los términos con a. Los demás términos también se multiplican para conservar la igualdad.
\end{minipage}\par\medskip
\noindent\begin{minipage}{\linewidth}\textbf{5. Multiplicamos la ecuación 1 por 2}\\[2pt]Multiplicar una ecuación por 2 significa multiplicar CADA término de la izquierda y también el número de la derecha. La igualdad se conserva porque hacemos lo mismo en ambos lados.
\par\smallskip Esta es la ecuación antes de multiplicar:
\[a + 2b = 4\]
\par\smallskip {\small\renewcommand{\arraystretch}{1.3}\begin{tabular}{>{\raggedright\arraybackslash}p{2.80cm}>{\raggedright\arraybackslash}p{1.91cm}>{\raggedright\arraybackslash}p{6.50cm}}\hline\rowcolor{black!7} \textbf{Parte de la ecuación} & \textbf{Antes} & \textbf{La cuenta y su resultado}\\\hline Término con a & $a$ & $2\cdot\left(a\right)=2a$\\ Término con b & $2b$ & $2\cdot\left(2b\right)=4b$\\ Número a la derecha del igual & $4$ & $2\cdot\left(4\right)=8$\\ \hline\end{tabular}}\par\smallskip
\par\smallskip Después de hacer todas las multiplicaciones queda:
\[2a + 4b = 8\]
\par\smallskip\textbf{Qué aprendimos.} La misma condición ahora está escrita con los términos multiplicados. Ya tenemos el tamaño que buscábamos para a.
\par\smallskip\textbf{Comprueba que lo entiendes.} ¿Podríamos multiplicar solamente el término con a?\par\textit{Respuesta:} No. Eso cambiaría la condición. Debemos multiplicar todos los términos, también el número que está después del igual.
\end{minipage}\par\medskip
\noindent\begin{minipage}{\linewidth}\textbf{6. Restamos las dos ecuaciones preparadas}\\[2pt]Restamos la ecuación 1 preparada a la ecuación 2 preparada. Juntamos cada término solo con otro que tenga la misma letra. Los números sin letra se calculan por separado.
\par\smallskip {\small\renewcommand{\arraystretch}{1.3}\begin{tabular}{>{\raggedright\arraybackslash}p{2.80cm}>{\raggedright\arraybackslash}p{5.38cm}>{\raggedright\arraybackslash}p{3.03cm}}\hline\rowcolor{black!7} \textbf{Qué agrupamos} & \textbf{Cuenta} & \textbf{Qué significa}\\\hline Los términos con a & $2a - \left(2a\right)=0$ & Se cancela esta letra.\\ Los términos con b & $-b - \left(4b\right)=-5b$ & Quedan -5b.\\ Los números a la derecha del igual & $-2 - \left(8\right)=-10$ & Este será el número de la nueva igualdad.\\ \hline\end{tabular}}\par\smallskip
\par\smallskip Primera igualdad que usaremos:
\[2a - b = -2\]
\par\smallskip Le restamos esta igualdad:
\[2a + 4b = 8\]
\par\smallskip\textbf{Qué aprendimos.} Los términos con a dan cero. Los términos con letras distintas no se mezclan.
\end{minipage}\par\medskip
\noindent\begin{minipage}{\linewidth}\textbf{7. Mira lo que hemos conseguido}\\[2pt]La nueva ecuación es el resultado de la cuenta anterior. Solo contiene b: después podremos despejarla.
\par\smallskip Escribimos la nueva condición:
\[-5b = -10\]
\[\left\{\begin{aligned}a + 2b &= 4\\-5b &= -10\\2c + d &= 4\\c - 2d &= -3\end{aligned}\right.\]
\par\smallskip\textbf{Qué aprendimos.} Conservamos la ecuación 1 y reemplazamos la 2 por esta nueva condición. De esa manera conservamos toda la información necesaria para resolver el sistema.
\end{minipage}\par\medskip
\noindent\begin{minipage}{\linewidth}\textbf{8. Ponemos una ecuación más cómoda arriba}\\[2pt]Cambiamos de lugar las ecuaciones 3 y 4. Las condiciones siguen siendo las mismas; solamente cambia el orden en que las leemos.
\par\smallskip\textbf{Por qué.} La ecuación que subimos tiene un número más pequeño delante de c. Eso facilita las multiplicaciones que haremos a continuación.
\[\left\{\begin{aligned}a + 2b &= 4\\-5b &= -10\\c - 2d &= -3\\2c + d &= 4\end{aligned}\right.\]
\end{minipage}\par\medskip
\noindent\begin{minipage}{\linewidth}\textbf{9. ¿Cómo haremos desaparecer c?}\\[2pt]Trabajaremos con las ecuaciones 3 y 4. Queremos que los términos con c tengan el mismo tamaño. Así podremos restarlos.
\par\smallskip\textbf{Qué buscamos.} Conseguir términos con 2 veces c en las dos ecuaciones.
\par\smallskip\textbf{Por qué.} Tener 2 veces la misma letra en ambos lados de la resta nos permitirá cancelarla.
\par\smallskip {\small\renewcommand{\arraystretch}{1.3}\begin{tabular}{>{\raggedright\arraybackslash}p{2.58cm}>{\raggedright\arraybackslash}p{2.47cm}>{\raggedright\arraybackslash}p{6.17cm}}\hline\rowcolor{black!7} \textbf{Dónde miramos} & \textbf{Qué término hay} & \textbf{Qué necesitamos hacer}\\\hline Ecuación 3 & $c$ & Multiplicar TODA esta ecuación por 2\\ Ecuación 4 & $2c$ & La dejamos igual.\\ \hline\end{tabular}}\par\smallskip
\par\smallskip\textbf{Qué aprendimos.} Estos multiplicadores se eligen para igualar los términos con c. Los demás términos también se multiplican para conservar la igualdad.
\end{minipage}\par\medskip
\noindent\begin{minipage}{\linewidth}\textbf{10. Multiplicamos la ecuación 3 por 2}\\[2pt]Multiplicar una ecuación por 2 significa multiplicar CADA término de la izquierda y también el número de la derecha. La igualdad se conserva porque hacemos lo mismo en ambos lados.
\par\smallskip Esta es la ecuación antes de multiplicar:
\[c - 2d = -3\]
\par\smallskip {\small\renewcommand{\arraystretch}{1.3}\begin{tabular}{>{\raggedright\arraybackslash}p{2.80cm}>{\raggedright\arraybackslash}p{1.91cm}>{\raggedright\arraybackslash}p{6.50cm}}\hline\rowcolor{black!7} \textbf{Parte de la ecuación} & \textbf{Antes} & \textbf{La cuenta y su resultado}\\\hline Término con c & $c$ & $2\cdot\left(c\right)=2c$\\ Término con d & $-2d$ & $2\cdot\left(-2d\right)=-4d$\\ Número a la derecha del igual & $-3$ & $2\cdot\left(-3\right)=-6$\\ \hline\end{tabular}}\par\smallskip
\par\smallskip Después de hacer todas las multiplicaciones queda:
\[2c - 4d = -6\]
\par\smallskip\textbf{Qué aprendimos.} La misma condición ahora está escrita con los términos multiplicados. Ya tenemos el tamaño que buscábamos para c.
\par\smallskip\textbf{Comprueba que lo entiendes.} ¿Podríamos multiplicar solamente el término con c?\par\textit{Respuesta:} No. Eso cambiaría la condición. Debemos multiplicar todos los términos, también el número que está después del igual.
\end{minipage}\par\medskip
\noindent\begin{minipage}{\linewidth}\textbf{11. Restamos las dos ecuaciones preparadas}\\[2pt]Restamos la ecuación 3 preparada a la ecuación 4 preparada. Juntamos cada término solo con otro que tenga la misma letra. Los números sin letra se calculan por separado.
\par\smallskip {\small\renewcommand{\arraystretch}{1.3}\begin{tabular}{>{\raggedright\arraybackslash}p{2.80cm}>{\raggedright\arraybackslash}p{5.38cm}>{\raggedright\arraybackslash}p{3.03cm}}\hline\rowcolor{black!7} \textbf{Qué agrupamos} & \textbf{Cuenta} & \textbf{Qué significa}\\\hline Los términos con c & $2c - \left(2c\right)=0$ & Se cancela esta letra.\\ Los términos con d & $d - \left(-4d\right)=5d$ & Quedan 5d.\\ Los números a la derecha del igual & $4 - \left(-6\right)=4+6=10$ & Este será el número de la nueva igualdad.\\ \hline\end{tabular}}\par\smallskip
\par\smallskip Primera igualdad que usaremos:
\[2c + d = 4\]
\par\smallskip Le restamos esta igualdad:
\[2c - 4d = -6\]
\par\smallskip\textbf{Qué aprendimos.} Los términos con c dan cero. Recuerda: restar un número negativo equivale a sumar su opuesto.
\end{minipage}\par\medskip
\noindent\begin{minipage}{\linewidth}\textbf{12. Mira lo que hemos conseguido}\\[2pt]La nueva ecuación es el resultado de la cuenta anterior. Solo contiene d: después podremos despejarla.
\par\smallskip Escribimos la nueva condición:
\[5d = 10\]
\[\left\{\begin{aligned}a + 2b &= 4\\-5b &= -10\\c - 2d &= -3\\5d &= 10\end{aligned}\right.\]
\par\smallskip\textbf{Qué aprendimos.} Conservamos la ecuación 3 y reemplazamos la 4 por esta nueva condición. De esa manera conservamos toda la información necesaria para resolver el sistema.
\end{minipage}\par\medskip
\noindent\begin{minipage}{\linewidth}\textbf{13. Ya tenemos una ecuación con una sola letra}\\[2pt]En esta condición solamente falta conocer d. El número 5 está multiplicando a esa letra. Ahora podemos deshacer esa multiplicación.
\par\smallskip Esta es la condición que vamos a resolver:
\[5d = 10\]
\end{minipage}\par\medskip
\noindent\begin{minipage}{\linewidth}\textbf{14. Calculamos cuánto vale d}\\[2pt]Para deshacer una multiplicación por 5, dividimos AMBOS lados entre 5. En el lado izquierdo queda una sola d. En el derecho hacemos la división.
\par\smallskip Dividimos los dos lados entre el mismo número:
\[\frac{5d}{5}=\frac{10}{5}\]
\par\smallskip Ya podemos leer el valor de la letra:
\[d=\frac{10}{5}=2\]
\par\smallskip\textbf{Qué aprendimos.} Este valor se usará en las ecuaciones donde todavía falten letras.
\end{minipage}\par\medskip
\noindent\begin{minipage}{\linewidth}\textbf{15. Usamos lo que ya sabemos para hallar c}\\[2pt]Ya sabemos d = 2. Sustituir significa escribir ese valor cada vez que aparezca su letra. Usamos paréntesis para que se vea qué número estamos multiplicando.
\par\smallskip\textbf{Qué buscamos.} Dejar solamente c como letra desconocida en esta ecuación.
\par\smallskip Tomamos esta condición que obtuvimos antes:
\[c - 2d = -3\]
\par\smallskip Reemplazamos las letras conocidas por sus valores:
\[c - 2\left(2\right)=-3\]
\end{minipage}\par\medskip
\noindent\begin{minipage}{\linewidth}\textbf{16. Calculamos las multiplicaciones que ya se pueden hacer}\\[2pt]La letra c sigue sin conocerse. Las otras partes ya tienen un valor o una expresión conocida y podemos efectuar sus productos.
\par\smallskip {\small\renewcommand{\arraystretch}{1.3}\begin{tabular}{>{\raggedright\arraybackslash}p{3.36cm}>{\raggedright\arraybackslash}p{4.48cm}>{\raggedright\arraybackslash}p{3.36cm}}\hline\rowcolor{black!7} \textbf{Parte conocida} & \textbf{La multiplicación} & \textbf{Resultado}\\\hline El término -2d & $-2\cdot\left(2\right)$ & $-4$\\ \hline\end{tabular}}\par\smallskip
\par\smallskip Al reunir esas partes conocidas, la ecuación queda:
\[c - 4=-3\]
\end{minipage}\par\medskip
\noindent\begin{minipage}{\linewidth}\textbf{17. Dejamos el término con c solo a un lado}\\[2pt]Para quitar la parte conocida, sumamos 4 en AMBOS lados. Piensa en una balanza: si hacemos lo mismo a ambos lados, la igualdad se mantiene.
\par\smallskip Hacemos la misma operación a la izquierda y a la derecha:
\[c - 4 + 4=-3 + 4\]
\par\smallskip Las partes que se compensan dan cero; queda:
\[c=1\]
\par\smallskip\textbf{Qué aprendimos.} «Pasar al otro lado» es una forma abreviada de decir que hicimos la misma suma o resta en los dos lados.
\end{minipage}\par\medskip
\noindent\begin{minipage}{\linewidth}\textbf{18. Leemos el valor de c}\\[2pt]La letra ya está sola: su número multiplicador es 1, que no se escribe. La igualdad nos dice directamente cuánto vale.
\par\smallskip Anotamos el valor encontrado:
\[c=1\]
\end{minipage}\par\medskip
\noindent\begin{minipage}{\linewidth}\textbf{19. Ya tenemos una ecuación con una sola letra}\\[2pt]En esta condición solamente falta conocer b. El número -5 está multiplicando a esa letra. Ahora podemos deshacer esa multiplicación.
\par\smallskip Esta es la condición que vamos a resolver:
\[-5b = -10\]
\end{minipage}\par\medskip
\noindent\begin{minipage}{\linewidth}\textbf{20. Calculamos cuánto vale b}\\[2pt]Para deshacer una multiplicación por -5, dividimos AMBOS lados entre -5. En el lado izquierdo queda una sola b. En el derecho hacemos la división.
\par\smallskip Dividimos los dos lados entre el mismo número:
\[\frac{-5b}{-5}=\frac{-10}{-5}\]
\par\smallskip Ya podemos leer el valor de la letra:
\[b=\frac{-10}{-5}=2\]
\par\smallskip\textbf{Qué aprendimos.} Un número negativo dividido entre otro negativo da positivo.
\end{minipage}\par\medskip
\noindent\begin{minipage}{\linewidth}\textbf{21. Usamos lo que ya sabemos para hallar a}\\[2pt]Ya sabemos b = 2. Sustituir significa escribir ese valor cada vez que aparezca su letra. Usamos paréntesis para que se vea qué número estamos multiplicando.
\par\smallskip\textbf{Qué buscamos.} Dejar solamente a como letra desconocida en esta ecuación.
\par\smallskip Tomamos esta condición que obtuvimos antes:
\[a + 2b = 4\]
\par\smallskip Reemplazamos las letras conocidas por sus valores:
\[a + 2\left(2\right)=4\]
\end{minipage}\par\medskip
\noindent\begin{minipage}{\linewidth}\textbf{22. Calculamos las multiplicaciones que ya se pueden hacer}\\[2pt]La letra a sigue sin conocerse. Las otras partes ya tienen un valor o una expresión conocida y podemos efectuar sus productos.
\par\smallskip {\small\renewcommand{\arraystretch}{1.3}\begin{tabular}{>{\raggedright\arraybackslash}p{3.36cm}>{\raggedright\arraybackslash}p{4.48cm}>{\raggedright\arraybackslash}p{3.36cm}}\hline\rowcolor{black!7} \textbf{Parte conocida} & \textbf{La multiplicación} & \textbf{Resultado}\\\hline El término 2b & $2\cdot\left(2\right)$ & $4$\\ \hline\end{tabular}}\par\smallskip
\par\smallskip Al reunir esas partes conocidas, la ecuación queda:
\[a + 4=4\]
\end{minipage}\par\medskip
\noindent\begin{minipage}{\linewidth}\textbf{23. Dejamos el término con a solo a un lado}\\[2pt]Para quitar la parte conocida, restamos 4 en AMBOS lados. Piensa en una balanza: si hacemos lo mismo a ambos lados, la igualdad se mantiene.
\par\smallskip Hacemos la misma operación a la izquierda y a la derecha:
\[a + 4 - 4=4 - 4\]
\par\smallskip Las partes que se compensan dan cero; queda:
\[a=0\]
\par\smallskip\textbf{Qué aprendimos.} «Pasar al otro lado» es una forma abreviada de decir que hicimos la misma suma o resta en los dos lados.
\end{minipage}\par\medskip
\noindent\begin{minipage}{\linewidth}\textbf{24. Leemos el valor de a}\\[2pt]La letra ya está sola: su número multiplicador es 1, que no se escribe. La igualdad nos dice directamente cuánto vale.
\par\smallskip Anotamos el valor encontrado:
\[a=0\]
\end{minipage}\par\medskip
\noindent\begin{minipage}{\linewidth}\textbf{25. Reunimos la respuesta; todavía falta comprobarla}\\[2pt]Escribimos todos los valores en el orden a, b, c, d. Estos valores deben funcionar juntos, no cada uno en una ecuación diferente.
\[\left\{\begin{aligned}a&=0\\b&=2\\c&=1\\d&=2\end{aligned}\right.\]
\par\smallskip\textbf{Qué aprendimos.} El siguiente paso es volver a las ecuaciones del enunciado y comprobar cada condición.
\end{minipage}\par\medskip
\noindent\begin{minipage}{\linewidth}\textbf{26. ¿La respuesta cumple la ecuación 1?}\\[2pt]Comprobar significa volver a la condición original y calcularla con nuestra respuesta. Si los dos lados dan el mismo número, esta condición sí se cumple. Deben cumplirse todas.
\par\smallskip La condición del enunciado era:
\[a + 2b = 4\]
\par\smallskip Ponemos nuestros valores donde estaban las letras:
\[\left(0\right) + 2\left(2\right)=4\]
\par\smallskip Hacemos las multiplicaciones y luego la suma:
\[4=4\]
\par\smallskip Comparamos los dos lados:
\[4=4\]
\par\smallskip\textbf{Qué aprendimos.} La igualdad final es verdadera. Nuestra respuesta sí cumple la ecuación 1.
\end{minipage}\par\medskip
\noindent\begin{minipage}{\linewidth}\textbf{27. ¿La respuesta cumple la ecuación 2?}\\[2pt]Comprobar significa volver a la condición original y calcularla con nuestra respuesta. Si los dos lados dan el mismo número, esta condición sí se cumple. Deben cumplirse todas.
\par\smallskip La condición del enunciado era:
\[2a - b = -2\]
\par\smallskip Ponemos nuestros valores donde estaban las letras:
\[2\left(0\right) - \left(2\right)=-2\]
\par\smallskip Hacemos las multiplicaciones y luego la suma:
\[-2=-2\]
\par\smallskip Comparamos los dos lados:
\[-2=-2\]
\par\smallskip\textbf{Qué aprendimos.} La igualdad final es verdadera. Nuestra respuesta sí cumple la ecuación 2.
\end{minipage}\par\medskip
\noindent\begin{minipage}{\linewidth}\textbf{28. ¿La respuesta cumple la ecuación 3?}\\[2pt]Comprobar significa volver a la condición original y calcularla con nuestra respuesta. Si los dos lados dan el mismo número, esta condición sí se cumple. Deben cumplirse todas.
\par\smallskip La condición del enunciado era:
\[2c + d = 4\]
\par\smallskip Ponemos nuestros valores donde estaban las letras:
\[2\left(1\right) + \left(2\right)=4\]
\par\smallskip Hacemos las multiplicaciones y luego la suma:
\[2 + 2=4\]
\par\smallskip Comparamos los dos lados:
\[4=4\]
\par\smallskip\textbf{Qué aprendimos.} La igualdad final es verdadera. Nuestra respuesta sí cumple la ecuación 3.
\end{minipage}\par\medskip
\noindent\begin{minipage}{\linewidth}\textbf{29. ¿La respuesta cumple la ecuación 4?}\\[2pt]Comprobar significa volver a la condición original y calcularla con nuestra respuesta. Si los dos lados dan el mismo número, esta condición sí se cumple. Deben cumplirse todas.
\par\smallskip La condición del enunciado era:
\[c - 2d = -3\]
\par\smallskip Ponemos nuestros valores donde estaban las letras:
\[\left(1\right) - 2\left(2\right)=-3\]
\par\smallskip Hacemos las multiplicaciones y luego la suma:
\[1 - 4=-3\]
\par\smallskip Comparamos los dos lados:
\[-3=-3\]
\par\smallskip\textbf{Qué aprendimos.} La igualdad final es verdadera. Nuestra respuesta sí cumple la ecuación 4.
\end{minipage}\par\medskip
\noindent\begin{minipage}{\linewidth}\textbf{30. Volvemos a las matrices originales}\\[2pt]Sustituimos los cuatro valores calculados en cada entrada de la matriz izquierda. Las cuatro casillas deben coincidir con la matriz derecha.
\par\smallskip \begin{minipage}[t]{.47\linewidth}\centering \textbf{Izquierda después de sustituir}\[\begin{bmatrix}4 & -2\\4 & -3\end{bmatrix}\]\end{minipage}\hfill\begin{minipage}[t]{.47\linewidth}\centering \textbf{Derecha}\[\begin{bmatrix}4 & -2\\4 & -3\end{bmatrix}\]\end{minipage}\par\smallskip
\end{minipage}\par\medskip
\respuesta{a = 0, b = 2, c = 1, d = 2.}
\noindent\begin{minipage}{\linewidth}\centering\includegraphics[width=.96\linewidth,height=6.8cm,keepaspectratio]{../recursos/grafico_1_2-3.png}\\{\small Lectura visual complementaria; la demostración está en las operaciones.}\end{minipage}
\end{multicols*}
\newpage\fancyfoot[L]{\small Sección 1.2 | Ejercicio 4 | Paso a paso}\hypertarget{1.2-4}{}
\pdfbookmark[1]{1.2 / 4}{bm1.2-4}
\titulo{4 | Operaciones con matrices}
\textbf{Enunciado}\par De ser posible, calcule cada expresión indicada. Se trabaja con las operaciones usuales de matrices reales.
\par\smallskip \begin{minipage}[t]{.47\linewidth}\centering \textbf{A}\[\begin{bmatrix}1 & 2 & 3\\2 & 1 & 4\end{bmatrix}\]\end{minipage}\hfill\begin{minipage}[t]{.47\linewidth}\centering \textbf{B}\[\begin{bmatrix}1 & 0\\2 & 1\\3 & 2\end{bmatrix}\]\end{minipage}\par\smallskip
\par\smallskip \begin{minipage}[t]{.47\linewidth}\centering \textbf{C}\[\begin{bmatrix}3 & -1 & 3\\4 & 1 & 5\\2 & 1 & 3\end{bmatrix}\]\end{minipage}\hfill\begin{minipage}[t]{.47\linewidth}\centering \textbf{D}\[\begin{bmatrix}3 & -2\\2 & 4\end{bmatrix}\]\end{minipage}\par\smallskip
\par\smallskip \begin{minipage}[t]{.47\linewidth}\centering \textbf{E}\[\begin{bmatrix}2 & -4 & 5\\0 & 1 & 4\\3 & 2 & 1\end{bmatrix}\]\end{minipage}\hfill\begin{minipage}[t]{.47\linewidth}\centering \textbf{F}\[\begin{bmatrix}-4 & 5\\2 & 3\end{bmatrix}\]\end{minipage}\par\smallskip
\par\smallskip \begin{minipage}[t]{.47\linewidth}\centering \textbf{O}\[\begin{bmatrix}0 & 0 & 0\\0 & 0 & 0\\0 & 0 & 0\end{bmatrix}\]\end{minipage}\par\smallskip
\textbf{Expresiones del enunciado}\par
a1. \(C+E\)\par
a2. \(E+C\)\par
b. \(A+B\)\par
c. \(D-F\)\par
d. \(-3C+5O\)\par
e. \(2C-3E\)\par
f. \(2B+F\)\par
\begin{multicols*}{2}
\noindent\begin{minipage}{\linewidth}\textbf{1. Antes de calcular: filas, columnas y operaciones}\\[2pt]El tamaño de una matriz se escribe filas por columnas. Para sumar o restar, los dos tamaños deben coincidir exactamente. Multiplicar por un número afecta a cada entrada. Transponer intercambia las posiciones: la entrada de fila i y columna j pasa a la fila j y columna i.
\[(M^T)_{ij}=m_{ji}\]
\par\smallskip\textbf{Qué aprendimos.} Los paréntesis indican qué operación hay que realizar primero. Una suma de tamaños distintos no se vuelve válida por transponerla después.
\end{minipage}\par\medskip
\noindent\begin{minipage}{\linewidth}\textbf{2. Inciso a1: C+E: qué vamos a calcular}\\[2pt]Identificamos la operación y comprobamos sus dimensiones antes de efectuar las cuentas. Las matrices originales permanecen escritas en el enunciado.
\[C+E\]
\end{minipage}\par\medskip
\noindent\begin{minipage}{\linewidth}\textbf{3. Operamos las entradas que ocupan la misma posición}\\[2pt]La fila y la columna nos dicen qué valores combinar. Hacemos una cuenta por casilla y después reunimos todas las casillas en la matriz resultado.
\par\smallskip {\small\renewcommand{\arraystretch}{1.3}\begin{tabular}{>{\raggedright\arraybackslash}p{3.74cm}>{\raggedright\arraybackslash}p{3.74cm}>{\raggedright\arraybackslash}p{3.74cm}}\hline\rowcolor{black!7} \textbf{Casilla} & \textbf{Cuenta} & \textbf{Resultado}\\\hline Fila 1, columna 1 & $(3)+(2)$ & $5$\\ Fila 1, columna 2 & $(-1)+(-4)$ & $-5$\\ Fila 1, columna 3 & $(3)+(5)$ & $8$\\ Fila 2, columna 1 & $(4)+(0)$ & $4$\\ Fila 2, columna 2 & $(1)+(1)$ & $2$\\ Fila 2, columna 3 & $(5)+(4)$ & $9$\\ Fila 3, columna 1 & $(2)+(3)$ & $5$\\ Fila 3, columna 2 & $(1)+(2)$ & $3$\\ Fila 3, columna 3 & $(3)+(1)$ & $4$\\ \hline\end{tabular}}\par\smallskip
\par\smallskip \textbf{Resultado de esta operación}\[\begin{bmatrix}5 & -5 & 8\\4 & 2 & 9\\5 & 3 & 4\end{bmatrix}\]
\par\smallskip\textbf{Qué aprendimos.} El resultado tiene 3 filas y 3 columnas.
\end{minipage}\par\medskip
\noindent\begin{minipage}{\linewidth}\textbf{4. Resultado del inciso a1}\\[2pt]Reunimos el resultado. Comprueba que cada casilla procede de la operación indicada, sin cambiar su posición.
\par\smallskip \textbf{Resultado}\[\begin{bmatrix}5 & -5 & 8\\4 & 2 & 9\\5 & 3 & 4\end{bmatrix}\]
\end{minipage}\par\medskip
\noindent\begin{minipage}{\linewidth}\textbf{5. Inciso a2: E+C: qué vamos a calcular}\\[2pt]Identificamos la operación y comprobamos sus dimensiones antes de efectuar las cuentas. Las matrices originales permanecen escritas en el enunciado.
\[E+C\]
\end{minipage}\par\medskip
\noindent\begin{minipage}{\linewidth}\textbf{6. Operamos las entradas que ocupan la misma posición}\\[2pt]La fila y la columna nos dicen qué valores combinar. Hacemos una cuenta por casilla y después reunimos todas las casillas en la matriz resultado.
\par\smallskip {\small\renewcommand{\arraystretch}{1.3}\begin{tabular}{>{\raggedright\arraybackslash}p{3.74cm}>{\raggedright\arraybackslash}p{3.74cm}>{\raggedright\arraybackslash}p{3.74cm}}\hline\rowcolor{black!7} \textbf{Casilla} & \textbf{Cuenta} & \textbf{Resultado}\\\hline Fila 1, columna 1 & $(2)+(3)$ & $5$\\ Fila 1, columna 2 & $(-4)+(-1)$ & $-5$\\ Fila 1, columna 3 & $(5)+(3)$ & $8$\\ Fila 2, columna 1 & $(0)+(4)$ & $4$\\ Fila 2, columna 2 & $(1)+(1)$ & $2$\\ Fila 2, columna 3 & $(4)+(5)$ & $9$\\ Fila 3, columna 1 & $(3)+(2)$ & $5$\\ Fila 3, columna 2 & $(2)+(1)$ & $3$\\ Fila 3, columna 3 & $(1)+(3)$ & $4$\\ \hline\end{tabular}}\par\smallskip
\par\smallskip \textbf{Resultado de esta operación}\[\begin{bmatrix}5 & -5 & 8\\4 & 2 & 9\\5 & 3 & 4\end{bmatrix}\]
\par\smallskip\textbf{Qué aprendimos.} El resultado tiene 3 filas y 3 columnas.
\end{minipage}\par\medskip
\noindent\begin{minipage}{\linewidth}\textbf{7. Resultado del inciso a2}\\[2pt]Reunimos el resultado. Comprueba que cada casilla procede de la operación indicada, sin cambiar su posición.
\par\smallskip \textbf{Resultado}\[\begin{bmatrix}5 & -5 & 8\\4 & 2 & 9\\5 & 3 & 4\end{bmatrix}\]
\end{minipage}\par\medskip
\noindent\begin{minipage}{\linewidth}\textbf{8. Inciso b: A+B: qué vamos a calcular}\\[2pt]Identificamos la operación y comprobamos sus dimensiones antes de efectuar las cuentas. Las matrices originales permanecen escritas en el enunciado.
\[A+B\]
\end{minipage}\par\medskip
\noindent\begin{minipage}{\linewidth}\textbf{9. Este inciso no está definido}\\[2pt]El primer operando tiene tamaño 2 por 3 y el segundo 3 por 2. La suma o resta no está definida. No se agregan ceros ni se recortan filas para forzarla.
\par\smallskip\textbf{Qué aprendimos.} La comprobación de dimensiones es parte de la solución, no un error del programa.
\end{minipage}\par\medskip
\noindent\begin{minipage}{\linewidth}\textbf{10. Inciso c: D-F: qué vamos a calcular}\\[2pt]Identificamos la operación y comprobamos sus dimensiones antes de efectuar las cuentas. Las matrices originales permanecen escritas en el enunciado.
\[D-F\]
\end{minipage}\par\medskip
\noindent\begin{minipage}{\linewidth}\textbf{11. Operamos las entradas que ocupan la misma posición}\\[2pt]La fila y la columna nos dicen qué valores combinar. Hacemos una cuenta por casilla y después reunimos todas las casillas en la matriz resultado.
\par\smallskip {\small\renewcommand{\arraystretch}{1.3}\begin{tabular}{>{\raggedright\arraybackslash}p{3.74cm}>{\raggedright\arraybackslash}p{3.74cm}>{\raggedright\arraybackslash}p{3.74cm}}\hline\rowcolor{black!7} \textbf{Casilla} & \textbf{Cuenta} & \textbf{Resultado}\\\hline Fila 1, columna 1 & $(3)-(-4)$ & $7$\\ Fila 1, columna 2 & $(-2)-(5)$ & $-7$\\ Fila 2, columna 1 & $(2)-(2)$ & $0$\\ Fila 2, columna 2 & $(4)-(3)$ & $1$\\ \hline\end{tabular}}\par\smallskip
\par\smallskip \textbf{Resultado de esta operación}\[\begin{bmatrix}7 & -7\\0 & 1\end{bmatrix}\]
\par\smallskip\textbf{Qué aprendimos.} El resultado tiene 2 filas y 2 columnas.
\end{minipage}\par\medskip
\noindent\begin{minipage}{\linewidth}\textbf{12. Resultado del inciso c}\\[2pt]Reunimos el resultado. Comprueba que cada casilla procede de la operación indicada, sin cambiar su posición.
\par\smallskip \textbf{Resultado}\[\begin{bmatrix}7 & -7\\0 & 1\end{bmatrix}\]
\end{minipage}\par\medskip
\noindent\begin{minipage}{\linewidth}\textbf{13. Inciso d: -3C+5O: qué vamos a calcular}\\[2pt]Identificamos la operación y comprobamos sus dimensiones antes de efectuar las cuentas. Las matrices originales permanecen escritas en el enunciado.
\[-3C+5O\]
\end{minipage}\par\medskip
\noindent\begin{minipage}{\linewidth}\textbf{14. Multiplicamos cada entrada por el escalar}\\[2pt]La fila y la columna nos dicen qué valores combinar. Hacemos una cuenta por casilla y después reunimos todas las casillas en la matriz resultado.
\par\smallskip {\small\renewcommand{\arraystretch}{1.3}\begin{tabular}{>{\raggedright\arraybackslash}p{3.74cm}>{\raggedright\arraybackslash}p{3.74cm}>{\raggedright\arraybackslash}p{3.74cm}}\hline\rowcolor{black!7} \textbf{Casilla} & \textbf{Cuenta} & \textbf{Resultado}\\\hline Fila 1, columna 1 & $-3\cdot(3)$ & $-9$\\ Fila 1, columna 2 & $-3\cdot(-1)$ & $3$\\ Fila 1, columna 3 & $-3\cdot(3)$ & $-9$\\ Fila 2, columna 1 & $-3\cdot(4)$ & $-12$\\ Fila 2, columna 2 & $-3\cdot(1)$ & $-3$\\ Fila 2, columna 3 & $-3\cdot(5)$ & $-15$\\ Fila 3, columna 1 & $-3\cdot(2)$ & $-6$\\ Fila 3, columna 2 & $-3\cdot(1)$ & $-3$\\ Fila 3, columna 3 & $-3\cdot(3)$ & $-9$\\ \hline\end{tabular}}\par\smallskip
\par\smallskip \textbf{Resultado de esta operación}\[\begin{bmatrix}-9 & 3 & -9\\-12 & -3 & -15\\-6 & -3 & -9\end{bmatrix}\]
\par\smallskip\textbf{Qué aprendimos.} El resultado tiene 3 filas y 3 columnas.
\end{minipage}\par\medskip
\noindent\begin{minipage}{\linewidth}\textbf{15. Multiplicamos cada entrada por el escalar}\\[2pt]La fila y la columna nos dicen qué valores combinar. Hacemos una cuenta por casilla y después reunimos todas las casillas en la matriz resultado.
\par\smallskip {\small\renewcommand{\arraystretch}{1.3}\begin{tabular}{>{\raggedright\arraybackslash}p{3.74cm}>{\raggedright\arraybackslash}p{3.74cm}>{\raggedright\arraybackslash}p{3.74cm}}\hline\rowcolor{black!7} \textbf{Casilla} & \textbf{Cuenta} & \textbf{Resultado}\\\hline Fila 1, columna 1 & $5\cdot(0)$ & $0$\\ Fila 1, columna 2 & $5\cdot(0)$ & $0$\\ Fila 1, columna 3 & $5\cdot(0)$ & $0$\\ Fila 2, columna 1 & $5\cdot(0)$ & $0$\\ Fila 2, columna 2 & $5\cdot(0)$ & $0$\\ Fila 2, columna 3 & $5\cdot(0)$ & $0$\\ Fila 3, columna 1 & $5\cdot(0)$ & $0$\\ Fila 3, columna 2 & $5\cdot(0)$ & $0$\\ Fila 3, columna 3 & $5\cdot(0)$ & $0$\\ \hline\end{tabular}}\par\smallskip
\par\smallskip \textbf{Resultado de esta operación}\[\begin{bmatrix}0 & 0 & 0\\0 & 0 & 0\\0 & 0 & 0\end{bmatrix}\]
\par\smallskip\textbf{Qué aprendimos.} El resultado tiene 3 filas y 3 columnas.
\end{minipage}\par\medskip
\noindent\begin{minipage}{\linewidth}\textbf{16. Operamos las entradas que ocupan la misma posición}\\[2pt]La fila y la columna nos dicen qué valores combinar. Hacemos una cuenta por casilla y después reunimos todas las casillas en la matriz resultado.
\par\smallskip {\small\renewcommand{\arraystretch}{1.3}\begin{tabular}{>{\raggedright\arraybackslash}p{3.74cm}>{\raggedright\arraybackslash}p{3.74cm}>{\raggedright\arraybackslash}p{3.74cm}}\hline\rowcolor{black!7} \textbf{Casilla} & \textbf{Cuenta} & \textbf{Resultado}\\\hline Fila 1, columna 1 & $(-9)+(0)$ & $-9$\\ Fila 1, columna 2 & $(3)+(0)$ & $3$\\ Fila 1, columna 3 & $(-9)+(0)$ & $-9$\\ Fila 2, columna 1 & $(-12)+(0)$ & $-12$\\ Fila 2, columna 2 & $(-3)+(0)$ & $-3$\\ Fila 2, columna 3 & $(-15)+(0)$ & $-15$\\ Fila 3, columna 1 & $(-6)+(0)$ & $-6$\\ Fila 3, columna 2 & $(-3)+(0)$ & $-3$\\ Fila 3, columna 3 & $(-9)+(0)$ & $-9$\\ \hline\end{tabular}}\par\smallskip
\par\smallskip \textbf{Resultado de esta operación}\[\begin{bmatrix}-9 & 3 & -9\\-12 & -3 & -15\\-6 & -3 & -9\end{bmatrix}\]
\par\smallskip\textbf{Qué aprendimos.} El resultado tiene 3 filas y 3 columnas.
\end{minipage}\par\medskip
\noindent\begin{minipage}{\linewidth}\textbf{17. Resultado del inciso d}\\[2pt]Reunimos el resultado. Comprueba que cada casilla procede de la operación indicada, sin cambiar su posición.
\par\smallskip \textbf{Resultado}\[\begin{bmatrix}-9 & 3 & -9\\-12 & -3 & -15\\-6 & -3 & -9\end{bmatrix}\]
\end{minipage}\par\medskip
\noindent\begin{minipage}{\linewidth}\textbf{18. Inciso e: 2C-3E: qué vamos a calcular}\\[2pt]Identificamos la operación y comprobamos sus dimensiones antes de efectuar las cuentas. Las matrices originales permanecen escritas en el enunciado.
\[2C-3E\]
\end{minipage}\par\medskip
\noindent\begin{minipage}{\linewidth}\textbf{19. Multiplicamos cada entrada por el escalar}\\[2pt]La fila y la columna nos dicen qué valores combinar. Hacemos una cuenta por casilla y después reunimos todas las casillas en la matriz resultado.
\par\smallskip {\small\renewcommand{\arraystretch}{1.3}\begin{tabular}{>{\raggedright\arraybackslash}p{3.74cm}>{\raggedright\arraybackslash}p{3.74cm}>{\raggedright\arraybackslash}p{3.74cm}}\hline\rowcolor{black!7} \textbf{Casilla} & \textbf{Cuenta} & \textbf{Resultado}\\\hline Fila 1, columna 1 & $2\cdot(3)$ & $6$\\ Fila 1, columna 2 & $2\cdot(-1)$ & $-2$\\ Fila 1, columna 3 & $2\cdot(3)$ & $6$\\ Fila 2, columna 1 & $2\cdot(4)$ & $8$\\ Fila 2, columna 2 & $2\cdot(1)$ & $2$\\ Fila 2, columna 3 & $2\cdot(5)$ & $10$\\ Fila 3, columna 1 & $2\cdot(2)$ & $4$\\ Fila 3, columna 2 & $2\cdot(1)$ & $2$\\ Fila 3, columna 3 & $2\cdot(3)$ & $6$\\ \hline\end{tabular}}\par\smallskip
\par\smallskip \textbf{Resultado de esta operación}\[\begin{bmatrix}6 & -2 & 6\\8 & 2 & 10\\4 & 2 & 6\end{bmatrix}\]
\par\smallskip\textbf{Qué aprendimos.} El resultado tiene 3 filas y 3 columnas.
\end{minipage}\par\medskip
\noindent\begin{minipage}{\linewidth}\textbf{20. Multiplicamos cada entrada por el escalar}\\[2pt]La fila y la columna nos dicen qué valores combinar. Hacemos una cuenta por casilla y después reunimos todas las casillas en la matriz resultado.
\par\smallskip {\small\renewcommand{\arraystretch}{1.3}\begin{tabular}{>{\raggedright\arraybackslash}p{3.74cm}>{\raggedright\arraybackslash}p{3.74cm}>{\raggedright\arraybackslash}p{3.74cm}}\hline\rowcolor{black!7} \textbf{Casilla} & \textbf{Cuenta} & \textbf{Resultado}\\\hline Fila 1, columna 1 & $3\cdot(2)$ & $6$\\ Fila 1, columna 2 & $3\cdot(-4)$ & $-12$\\ Fila 1, columna 3 & $3\cdot(5)$ & $15$\\ Fila 2, columna 1 & $3\cdot(0)$ & $0$\\ Fila 2, columna 2 & $3\cdot(1)$ & $3$\\ Fila 2, columna 3 & $3\cdot(4)$ & $12$\\ Fila 3, columna 1 & $3\cdot(3)$ & $9$\\ Fila 3, columna 2 & $3\cdot(2)$ & $6$\\ Fila 3, columna 3 & $3\cdot(1)$ & $3$\\ \hline\end{tabular}}\par\smallskip
\par\smallskip \textbf{Resultado de esta operación}\[\begin{bmatrix}6 & -12 & 15\\0 & 3 & 12\\9 & 6 & 3\end{bmatrix}\]
\par\smallskip\textbf{Qué aprendimos.} El resultado tiene 3 filas y 3 columnas.
\end{minipage}\par\medskip
\noindent\begin{minipage}{\linewidth}\textbf{21. Operamos las entradas que ocupan la misma posición}\\[2pt]La fila y la columna nos dicen qué valores combinar. Hacemos una cuenta por casilla y después reunimos todas las casillas en la matriz resultado.
\par\smallskip {\small\renewcommand{\arraystretch}{1.3}\begin{tabular}{>{\raggedright\arraybackslash}p{3.74cm}>{\raggedright\arraybackslash}p{3.74cm}>{\raggedright\arraybackslash}p{3.74cm}}\hline\rowcolor{black!7} \textbf{Casilla} & \textbf{Cuenta} & \textbf{Resultado}\\\hline Fila 1, columna 1 & $(6)-(6)$ & $0$\\ Fila 1, columna 2 & $(-2)-(-12)$ & $10$\\ Fila 1, columna 3 & $(6)-(15)$ & $-9$\\ Fila 2, columna 1 & $(8)-(0)$ & $8$\\ Fila 2, columna 2 & $(2)-(3)$ & $-1$\\ Fila 2, columna 3 & $(10)-(12)$ & $-2$\\ Fila 3, columna 1 & $(4)-(9)$ & $-5$\\ Fila 3, columna 2 & $(2)-(6)$ & $-4$\\ Fila 3, columna 3 & $(6)-(3)$ & $3$\\ \hline\end{tabular}}\par\smallskip
\par\smallskip \textbf{Resultado de esta operación}\[\begin{bmatrix}0 & 10 & -9\\8 & -1 & -2\\-5 & -4 & 3\end{bmatrix}\]
\par\smallskip\textbf{Qué aprendimos.} El resultado tiene 3 filas y 3 columnas.
\end{minipage}\par\medskip
\noindent\begin{minipage}{\linewidth}\textbf{22. Resultado del inciso e}\\[2pt]Reunimos el resultado. Comprueba que cada casilla procede de la operación indicada, sin cambiar su posición.
\par\smallskip \textbf{Resultado}\[\begin{bmatrix}0 & 10 & -9\\8 & -1 & -2\\-5 & -4 & 3\end{bmatrix}\]
\end{minipage}\par\medskip
\noindent\begin{minipage}{\linewidth}\textbf{23. Inciso f: 2B+F: qué vamos a calcular}\\[2pt]Identificamos la operación y comprobamos sus dimensiones antes de efectuar las cuentas. Las matrices originales permanecen escritas en el enunciado.
\[2B+F\]
\end{minipage}\par\medskip
\noindent\begin{minipage}{\linewidth}\textbf{24. Multiplicamos cada entrada por el escalar}\\[2pt]La fila y la columna nos dicen qué valores combinar. Hacemos una cuenta por casilla y después reunimos todas las casillas en la matriz resultado.
\par\smallskip {\small\renewcommand{\arraystretch}{1.3}\begin{tabular}{>{\raggedright\arraybackslash}p{3.74cm}>{\raggedright\arraybackslash}p{3.74cm}>{\raggedright\arraybackslash}p{3.74cm}}\hline\rowcolor{black!7} \textbf{Casilla} & \textbf{Cuenta} & \textbf{Resultado}\\\hline Fila 1, columna 1 & $2\cdot(1)$ & $2$\\ Fila 1, columna 2 & $2\cdot(0)$ & $0$\\ Fila 2, columna 1 & $2\cdot(2)$ & $4$\\ Fila 2, columna 2 & $2\cdot(1)$ & $2$\\ Fila 3, columna 1 & $2\cdot(3)$ & $6$\\ Fila 3, columna 2 & $2\cdot(2)$ & $4$\\ \hline\end{tabular}}\par\smallskip
\par\smallskip \textbf{Resultado de esta operación}\[\begin{bmatrix}2 & 0\\4 & 2\\6 & 4\end{bmatrix}\]
\par\smallskip\textbf{Qué aprendimos.} El resultado tiene 3 filas y 2 columnas.
\end{minipage}\par\medskip
\noindent\begin{minipage}{\linewidth}\textbf{25. Este inciso no está definido}\\[2pt]El primer operando tiene tamaño 3 por 2 y el segundo 2 por 2. La suma o resta no está definida. No se agregan ceros ni se recortan filas para forzarla.
\par\smallskip\textbf{Qué aprendimos.} La comprobación de dimensiones es parte de la solución, no un error del programa.
\end{minipage}\par\medskip
\respuesta{a1: C+E: [[5, -5, 8], [4, 2, 9], [5, 3, 4]]; a2: E+C: [[5, -5, 8], [4, 2, 9], [5, 3, 4]]; b: A+B: no definida; c: D-F: [[7, -7], [0, 1]]; d: -3C+5O: [[-9, 3, -9], [-12, -3, -15], [-6, -3, -9]]; e: 2C-3E: [[0, 10, -9], [8, -1, -2], [-5, -4, 3]]; f: 2B+F: no definida.}
\noindent\begin{minipage}{\linewidth}\centering\includegraphics[width=.96\linewidth,height=6.8cm,keepaspectratio]{../recursos/grafico_1_2-4.png}\\{\small Lectura visual complementaria; la demostración está en las operaciones.}\end{minipage}
\end{multicols*}
\newpage\fancyfoot[L]{\small Sección 1.2 | Ejercicio 5 | Paso a paso}\hypertarget{1.2-5}{}
\pdfbookmark[1]{1.2 / 5}{bm1.2-5}
\titulo{5 | Operaciones con matrices}
\textbf{Enunciado}\par De ser posible, calcule cada expresión indicada. Se trabaja con las operaciones usuales de matrices reales.
\par\smallskip \begin{minipage}[t]{.47\linewidth}\centering \textbf{A}\[\begin{bmatrix}1 & 2 & 3\\2 & 1 & 4\end{bmatrix}\]\end{minipage}\hfill\begin{minipage}[t]{.47\linewidth}\centering \textbf{B}\[\begin{bmatrix}1 & 0\\2 & 1\\3 & 2\end{bmatrix}\]\end{minipage}\par\smallskip
\par\smallskip \begin{minipage}[t]{.47\linewidth}\centering \textbf{C}\[\begin{bmatrix}3 & -1 & 3\\4 & 1 & 5\\2 & 1 & 3\end{bmatrix}\]\end{minipage}\hfill\begin{minipage}[t]{.47\linewidth}\centering \textbf{D}\[\begin{bmatrix}3 & -2\\2 & 4\end{bmatrix}\]\end{minipage}\par\smallskip
\par\smallskip \begin{minipage}[t]{.47\linewidth}\centering \textbf{E}\[\begin{bmatrix}2 & -4 & 5\\0 & 1 & 4\\3 & 2 & 1\end{bmatrix}\]\end{minipage}\hfill\begin{minipage}[t]{.47\linewidth}\centering \textbf{F}\[\begin{bmatrix}-4 & 5\\2 & 3\end{bmatrix}\]\end{minipage}\par\smallskip
\par\smallskip \begin{minipage}[t]{.47\linewidth}\centering \textbf{O}\[\begin{bmatrix}0 & 0 & 0\\0 & 0 & 0\\0 & 0 & 0\end{bmatrix}\]\end{minipage}\par\smallskip
\textbf{Expresiones del enunciado}\par
a. \(3D+2F\)\par
b1. \(3(2A)\)\par
b2. \(6A\)\par
c1. \(3A+2A\)\par
c2. \(5A\)\par
d1. \(2(D+F)\)\par
d2. \(2D+2F\)\par
e1. \((2+3)D\)\par
e2. \(2D+3D\)\par
f. \(3(B+D)\)\par
\begin{multicols*}{2}
\noindent\begin{minipage}{\linewidth}\textbf{1. Antes de calcular: filas, columnas y operaciones}\\[2pt]El tamaño de una matriz se escribe filas por columnas. Para sumar o restar, los dos tamaños deben coincidir exactamente. Multiplicar por un número afecta a cada entrada. Transponer intercambia las posiciones: la entrada de fila i y columna j pasa a la fila j y columna i.
\[(M^T)_{ij}=m_{ji}\]
\par\smallskip\textbf{Qué aprendimos.} Los paréntesis indican qué operación hay que realizar primero. Una suma de tamaños distintos no se vuelve válida por transponerla después.
\end{minipage}\par\medskip
\noindent\begin{minipage}{\linewidth}\textbf{2. Inciso a: 3D+2F: qué vamos a calcular}\\[2pt]Identificamos la operación y comprobamos sus dimensiones antes de efectuar las cuentas. Las matrices originales permanecen escritas en el enunciado.
\[3D+2F\]
\end{minipage}\par\medskip
\noindent\begin{minipage}{\linewidth}\textbf{3. Multiplicamos cada entrada por el escalar}\\[2pt]La fila y la columna nos dicen qué valores combinar. Hacemos una cuenta por casilla y después reunimos todas las casillas en la matriz resultado.
\par\smallskip {\small\renewcommand{\arraystretch}{1.3}\begin{tabular}{>{\raggedright\arraybackslash}p{3.74cm}>{\raggedright\arraybackslash}p{3.74cm}>{\raggedright\arraybackslash}p{3.74cm}}\hline\rowcolor{black!7} \textbf{Casilla} & \textbf{Cuenta} & \textbf{Resultado}\\\hline Fila 1, columna 1 & $3\cdot(3)$ & $9$\\ Fila 1, columna 2 & $3\cdot(-2)$ & $-6$\\ Fila 2, columna 1 & $3\cdot(2)$ & $6$\\ Fila 2, columna 2 & $3\cdot(4)$ & $12$\\ \hline\end{tabular}}\par\smallskip
\par\smallskip \textbf{Resultado de esta operación}\[\begin{bmatrix}9 & -6\\6 & 12\end{bmatrix}\]
\par\smallskip\textbf{Qué aprendimos.} El resultado tiene 2 filas y 2 columnas.
\end{minipage}\par\medskip
\noindent\begin{minipage}{\linewidth}\textbf{4. Multiplicamos cada entrada por el escalar}\\[2pt]La fila y la columna nos dicen qué valores combinar. Hacemos una cuenta por casilla y después reunimos todas las casillas en la matriz resultado.
\par\smallskip {\small\renewcommand{\arraystretch}{1.3}\begin{tabular}{>{\raggedright\arraybackslash}p{3.74cm}>{\raggedright\arraybackslash}p{3.74cm}>{\raggedright\arraybackslash}p{3.74cm}}\hline\rowcolor{black!7} \textbf{Casilla} & \textbf{Cuenta} & \textbf{Resultado}\\\hline Fila 1, columna 1 & $2\cdot(-4)$ & $-8$\\ Fila 1, columna 2 & $2\cdot(5)$ & $10$\\ Fila 2, columna 1 & $2\cdot(2)$ & $4$\\ Fila 2, columna 2 & $2\cdot(3)$ & $6$\\ \hline\end{tabular}}\par\smallskip
\par\smallskip \textbf{Resultado de esta operación}\[\begin{bmatrix}-8 & 10\\4 & 6\end{bmatrix}\]
\par\smallskip\textbf{Qué aprendimos.} El resultado tiene 2 filas y 2 columnas.
\end{minipage}\par\medskip
\noindent\begin{minipage}{\linewidth}\textbf{5. Operamos las entradas que ocupan la misma posición}\\[2pt]La fila y la columna nos dicen qué valores combinar. Hacemos una cuenta por casilla y después reunimos todas las casillas en la matriz resultado.
\par\smallskip {\small\renewcommand{\arraystretch}{1.3}\begin{tabular}{>{\raggedright\arraybackslash}p{3.74cm}>{\raggedright\arraybackslash}p{3.74cm}>{\raggedright\arraybackslash}p{3.74cm}}\hline\rowcolor{black!7} \textbf{Casilla} & \textbf{Cuenta} & \textbf{Resultado}\\\hline Fila 1, columna 1 & $(9)+(-8)$ & $1$\\ Fila 1, columna 2 & $(-6)+(10)$ & $4$\\ Fila 2, columna 1 & $(6)+(4)$ & $10$\\ Fila 2, columna 2 & $(12)+(6)$ & $18$\\ \hline\end{tabular}}\par\smallskip
\par\smallskip \textbf{Resultado de esta operación}\[\begin{bmatrix}1 & 4\\10 & 18\end{bmatrix}\]
\par\smallskip\textbf{Qué aprendimos.} El resultado tiene 2 filas y 2 columnas.
\end{minipage}\par\medskip
\noindent\begin{minipage}{\linewidth}\textbf{6. Resultado del inciso a}\\[2pt]Reunimos el resultado. Comprueba que cada casilla procede de la operación indicada, sin cambiar su posición.
\par\smallskip \textbf{Resultado}\[\begin{bmatrix}1 & 4\\10 & 18\end{bmatrix}\]
\end{minipage}\par\medskip
\noindent\begin{minipage}{\linewidth}\textbf{7. Inciso b1: 3(2A): qué vamos a calcular}\\[2pt]Identificamos la operación y comprobamos sus dimensiones antes de efectuar las cuentas. Las matrices originales permanecen escritas en el enunciado.
\[3(2A)\]
\end{minipage}\par\medskip
\noindent\begin{minipage}{\linewidth}\textbf{8. Multiplicamos cada entrada por el escalar}\\[2pt]La fila y la columna nos dicen qué valores combinar. Hacemos una cuenta por casilla y después reunimos todas las casillas en la matriz resultado.
\par\smallskip {\small\renewcommand{\arraystretch}{1.3}\begin{tabular}{>{\raggedright\arraybackslash}p{3.74cm}>{\raggedright\arraybackslash}p{3.74cm}>{\raggedright\arraybackslash}p{3.74cm}}\hline\rowcolor{black!7} \textbf{Casilla} & \textbf{Cuenta} & \textbf{Resultado}\\\hline Fila 1, columna 1 & $2\cdot(1)$ & $2$\\ Fila 1, columna 2 & $2\cdot(2)$ & $4$\\ Fila 1, columna 3 & $2\cdot(3)$ & $6$\\ Fila 2, columna 1 & $2\cdot(2)$ & $4$\\ Fila 2, columna 2 & $2\cdot(1)$ & $2$\\ Fila 2, columna 3 & $2\cdot(4)$ & $8$\\ \hline\end{tabular}}\par\smallskip
\par\smallskip \textbf{Resultado de esta operación}\[\begin{bmatrix}2 & 4 & 6\\4 & 2 & 8\end{bmatrix}\]
\par\smallskip\textbf{Qué aprendimos.} El resultado tiene 2 filas y 3 columnas.
\end{minipage}\par\medskip
\noindent\begin{minipage}{\linewidth}\textbf{9. Multiplicamos cada entrada por el escalar}\\[2pt]La fila y la columna nos dicen qué valores combinar. Hacemos una cuenta por casilla y después reunimos todas las casillas en la matriz resultado.
\par\smallskip {\small\renewcommand{\arraystretch}{1.3}\begin{tabular}{>{\raggedright\arraybackslash}p{3.74cm}>{\raggedright\arraybackslash}p{3.74cm}>{\raggedright\arraybackslash}p{3.74cm}}\hline\rowcolor{black!7} \textbf{Casilla} & \textbf{Cuenta} & \textbf{Resultado}\\\hline Fila 1, columna 1 & $3\cdot(2)$ & $6$\\ Fila 1, columna 2 & $3\cdot(4)$ & $12$\\ Fila 1, columna 3 & $3\cdot(6)$ & $18$\\ Fila 2, columna 1 & $3\cdot(4)$ & $12$\\ Fila 2, columna 2 & $3\cdot(2)$ & $6$\\ Fila 2, columna 3 & $3\cdot(8)$ & $24$\\ \hline\end{tabular}}\par\smallskip
\par\smallskip \textbf{Resultado de esta operación}\[\begin{bmatrix}6 & 12 & 18\\12 & 6 & 24\end{bmatrix}\]
\par\smallskip\textbf{Qué aprendimos.} El resultado tiene 2 filas y 3 columnas.
\end{minipage}\par\medskip
\noindent\begin{minipage}{\linewidth}\textbf{10. Resultado del inciso b1}\\[2pt]Reunimos el resultado. Comprueba que cada casilla procede de la operación indicada, sin cambiar su posición.
\par\smallskip \textbf{Resultado}\[\begin{bmatrix}6 & 12 & 18\\12 & 6 & 24\end{bmatrix}\]
\end{minipage}\par\medskip
\noindent\begin{minipage}{\linewidth}\textbf{11. Inciso b2: 6A: qué vamos a calcular}\\[2pt]Identificamos la operación y comprobamos sus dimensiones antes de efectuar las cuentas. Las matrices originales permanecen escritas en el enunciado.
\[6A\]
\end{minipage}\par\medskip
\noindent\begin{minipage}{\linewidth}\textbf{12. Multiplicamos cada entrada por el escalar}\\[2pt]La fila y la columna nos dicen qué valores combinar. Hacemos una cuenta por casilla y después reunimos todas las casillas en la matriz resultado.
\par\smallskip {\small\renewcommand{\arraystretch}{1.3}\begin{tabular}{>{\raggedright\arraybackslash}p{3.74cm}>{\raggedright\arraybackslash}p{3.74cm}>{\raggedright\arraybackslash}p{3.74cm}}\hline\rowcolor{black!7} \textbf{Casilla} & \textbf{Cuenta} & \textbf{Resultado}\\\hline Fila 1, columna 1 & $6\cdot(1)$ & $6$\\ Fila 1, columna 2 & $6\cdot(2)$ & $12$\\ Fila 1, columna 3 & $6\cdot(3)$ & $18$\\ Fila 2, columna 1 & $6\cdot(2)$ & $12$\\ Fila 2, columna 2 & $6\cdot(1)$ & $6$\\ Fila 2, columna 3 & $6\cdot(4)$ & $24$\\ \hline\end{tabular}}\par\smallskip
\par\smallskip \textbf{Resultado de esta operación}\[\begin{bmatrix}6 & 12 & 18\\12 & 6 & 24\end{bmatrix}\]
\par\smallskip\textbf{Qué aprendimos.} El resultado tiene 2 filas y 3 columnas.
\end{minipage}\par\medskip
\noindent\begin{minipage}{\linewidth}\textbf{13. Resultado del inciso b2}\\[2pt]Reunimos el resultado. Comprueba que cada casilla procede de la operación indicada, sin cambiar su posición.
\par\smallskip \textbf{Resultado}\[\begin{bmatrix}6 & 12 & 18\\12 & 6 & 24\end{bmatrix}\]
\end{minipage}\par\medskip
\noindent\begin{minipage}{\linewidth}\textbf{14. Inciso c1: 3A+2A: qué vamos a calcular}\\[2pt]Identificamos la operación y comprobamos sus dimensiones antes de efectuar las cuentas. Las matrices originales permanecen escritas en el enunciado.
\[3A+2A\]
\end{minipage}\par\medskip
\noindent\begin{minipage}{\linewidth}\textbf{15. Multiplicamos cada entrada por el escalar}\\[2pt]La fila y la columna nos dicen qué valores combinar. Hacemos una cuenta por casilla y después reunimos todas las casillas en la matriz resultado.
\par\smallskip {\small\renewcommand{\arraystretch}{1.3}\begin{tabular}{>{\raggedright\arraybackslash}p{3.74cm}>{\raggedright\arraybackslash}p{3.74cm}>{\raggedright\arraybackslash}p{3.74cm}}\hline\rowcolor{black!7} \textbf{Casilla} & \textbf{Cuenta} & \textbf{Resultado}\\\hline Fila 1, columna 1 & $3\cdot(1)$ & $3$\\ Fila 1, columna 2 & $3\cdot(2)$ & $6$\\ Fila 1, columna 3 & $3\cdot(3)$ & $9$\\ Fila 2, columna 1 & $3\cdot(2)$ & $6$\\ Fila 2, columna 2 & $3\cdot(1)$ & $3$\\ Fila 2, columna 3 & $3\cdot(4)$ & $12$\\ \hline\end{tabular}}\par\smallskip
\par\smallskip \textbf{Resultado de esta operación}\[\begin{bmatrix}3 & 6 & 9\\6 & 3 & 12\end{bmatrix}\]
\par\smallskip\textbf{Qué aprendimos.} El resultado tiene 2 filas y 3 columnas.
\end{minipage}\par\medskip
\noindent\begin{minipage}{\linewidth}\textbf{16. Multiplicamos cada entrada por el escalar}\\[2pt]La fila y la columna nos dicen qué valores combinar. Hacemos una cuenta por casilla y después reunimos todas las casillas en la matriz resultado.
\par\smallskip {\small\renewcommand{\arraystretch}{1.3}\begin{tabular}{>{\raggedright\arraybackslash}p{3.74cm}>{\raggedright\arraybackslash}p{3.74cm}>{\raggedright\arraybackslash}p{3.74cm}}\hline\rowcolor{black!7} \textbf{Casilla} & \textbf{Cuenta} & \textbf{Resultado}\\\hline Fila 1, columna 1 & $2\cdot(1)$ & $2$\\ Fila 1, columna 2 & $2\cdot(2)$ & $4$\\ Fila 1, columna 3 & $2\cdot(3)$ & $6$\\ Fila 2, columna 1 & $2\cdot(2)$ & $4$\\ Fila 2, columna 2 & $2\cdot(1)$ & $2$\\ Fila 2, columna 3 & $2\cdot(4)$ & $8$\\ \hline\end{tabular}}\par\smallskip
\par\smallskip \textbf{Resultado de esta operación}\[\begin{bmatrix}2 & 4 & 6\\4 & 2 & 8\end{bmatrix}\]
\par\smallskip\textbf{Qué aprendimos.} El resultado tiene 2 filas y 3 columnas.
\end{minipage}\par\medskip
\noindent\begin{minipage}{\linewidth}\textbf{17. Operamos las entradas que ocupan la misma posición}\\[2pt]La fila y la columna nos dicen qué valores combinar. Hacemos una cuenta por casilla y después reunimos todas las casillas en la matriz resultado.
\par\smallskip {\small\renewcommand{\arraystretch}{1.3}\begin{tabular}{>{\raggedright\arraybackslash}p{3.74cm}>{\raggedright\arraybackslash}p{3.74cm}>{\raggedright\arraybackslash}p{3.74cm}}\hline\rowcolor{black!7} \textbf{Casilla} & \textbf{Cuenta} & \textbf{Resultado}\\\hline Fila 1, columna 1 & $(3)+(2)$ & $5$\\ Fila 1, columna 2 & $(6)+(4)$ & $10$\\ Fila 1, columna 3 & $(9)+(6)$ & $15$\\ Fila 2, columna 1 & $(6)+(4)$ & $10$\\ Fila 2, columna 2 & $(3)+(2)$ & $5$\\ Fila 2, columna 3 & $(12)+(8)$ & $20$\\ \hline\end{tabular}}\par\smallskip
\par\smallskip \textbf{Resultado de esta operación}\[\begin{bmatrix}5 & 10 & 15\\10 & 5 & 20\end{bmatrix}\]
\par\smallskip\textbf{Qué aprendimos.} El resultado tiene 2 filas y 3 columnas.
\end{minipage}\par\medskip
\noindent\begin{minipage}{\linewidth}\textbf{18. Resultado del inciso c1}\\[2pt]Reunimos el resultado. Comprueba que cada casilla procede de la operación indicada, sin cambiar su posición.
\par\smallskip \textbf{Resultado}\[\begin{bmatrix}5 & 10 & 15\\10 & 5 & 20\end{bmatrix}\]
\end{minipage}\par\medskip
\noindent\begin{minipage}{\linewidth}\textbf{19. Inciso c2: 5A: qué vamos a calcular}\\[2pt]Identificamos la operación y comprobamos sus dimensiones antes de efectuar las cuentas. Las matrices originales permanecen escritas en el enunciado.
\[5A\]
\end{minipage}\par\medskip
\noindent\begin{minipage}{\linewidth}\textbf{20. Multiplicamos cada entrada por el escalar}\\[2pt]La fila y la columna nos dicen qué valores combinar. Hacemos una cuenta por casilla y después reunimos todas las casillas en la matriz resultado.
\par\smallskip {\small\renewcommand{\arraystretch}{1.3}\begin{tabular}{>{\raggedright\arraybackslash}p{3.74cm}>{\raggedright\arraybackslash}p{3.74cm}>{\raggedright\arraybackslash}p{3.74cm}}\hline\rowcolor{black!7} \textbf{Casilla} & \textbf{Cuenta} & \textbf{Resultado}\\\hline Fila 1, columna 1 & $5\cdot(1)$ & $5$\\ Fila 1, columna 2 & $5\cdot(2)$ & $10$\\ Fila 1, columna 3 & $5\cdot(3)$ & $15$\\ Fila 2, columna 1 & $5\cdot(2)$ & $10$\\ Fila 2, columna 2 & $5\cdot(1)$ & $5$\\ Fila 2, columna 3 & $5\cdot(4)$ & $20$\\ \hline\end{tabular}}\par\smallskip
\par\smallskip \textbf{Resultado de esta operación}\[\begin{bmatrix}5 & 10 & 15\\10 & 5 & 20\end{bmatrix}\]
\par\smallskip\textbf{Qué aprendimos.} El resultado tiene 2 filas y 3 columnas.
\end{minipage}\par\medskip
\noindent\begin{minipage}{\linewidth}\textbf{21. Resultado del inciso c2}\\[2pt]Reunimos el resultado. Comprueba que cada casilla procede de la operación indicada, sin cambiar su posición.
\par\smallskip \textbf{Resultado}\[\begin{bmatrix}5 & 10 & 15\\10 & 5 & 20\end{bmatrix}\]
\end{minipage}\par\medskip
\noindent\begin{minipage}{\linewidth}\textbf{22. Inciso d1: 2(D+F): qué vamos a calcular}\\[2pt]Identificamos la operación y comprobamos sus dimensiones antes de efectuar las cuentas. Las matrices originales permanecen escritas en el enunciado.
\[2(D+F)\]
\end{minipage}\par\medskip
\noindent\begin{minipage}{\linewidth}\textbf{23. Operamos las entradas que ocupan la misma posición}\\[2pt]La fila y la columna nos dicen qué valores combinar. Hacemos una cuenta por casilla y después reunimos todas las casillas en la matriz resultado.
\par\smallskip {\small\renewcommand{\arraystretch}{1.3}\begin{tabular}{>{\raggedright\arraybackslash}p{3.74cm}>{\raggedright\arraybackslash}p{3.74cm}>{\raggedright\arraybackslash}p{3.74cm}}\hline\rowcolor{black!7} \textbf{Casilla} & \textbf{Cuenta} & \textbf{Resultado}\\\hline Fila 1, columna 1 & $(3)+(-4)$ & $-1$\\ Fila 1, columna 2 & $(-2)+(5)$ & $3$\\ Fila 2, columna 1 & $(2)+(2)$ & $4$\\ Fila 2, columna 2 & $(4)+(3)$ & $7$\\ \hline\end{tabular}}\par\smallskip
\par\smallskip \textbf{Resultado de esta operación}\[\begin{bmatrix}-1 & 3\\4 & 7\end{bmatrix}\]
\par\smallskip\textbf{Qué aprendimos.} El resultado tiene 2 filas y 2 columnas.
\end{minipage}\par\medskip
\noindent\begin{minipage}{\linewidth}\textbf{24. Multiplicamos cada entrada por el escalar}\\[2pt]La fila y la columna nos dicen qué valores combinar. Hacemos una cuenta por casilla y después reunimos todas las casillas en la matriz resultado.
\par\smallskip {\small\renewcommand{\arraystretch}{1.3}\begin{tabular}{>{\raggedright\arraybackslash}p{3.74cm}>{\raggedright\arraybackslash}p{3.74cm}>{\raggedright\arraybackslash}p{3.74cm}}\hline\rowcolor{black!7} \textbf{Casilla} & \textbf{Cuenta} & \textbf{Resultado}\\\hline Fila 1, columna 1 & $2\cdot(-1)$ & $-2$\\ Fila 1, columna 2 & $2\cdot(3)$ & $6$\\ Fila 2, columna 1 & $2\cdot(4)$ & $8$\\ Fila 2, columna 2 & $2\cdot(7)$ & $14$\\ \hline\end{tabular}}\par\smallskip
\par\smallskip \textbf{Resultado de esta operación}\[\begin{bmatrix}-2 & 6\\8 & 14\end{bmatrix}\]
\par\smallskip\textbf{Qué aprendimos.} El resultado tiene 2 filas y 2 columnas.
\end{minipage}\par\medskip
\noindent\begin{minipage}{\linewidth}\textbf{25. Resultado del inciso d1}\\[2pt]Reunimos el resultado. Comprueba que cada casilla procede de la operación indicada, sin cambiar su posición.
\par\smallskip \textbf{Resultado}\[\begin{bmatrix}-2 & 6\\8 & 14\end{bmatrix}\]
\end{minipage}\par\medskip
\noindent\begin{minipage}{\linewidth}\textbf{26. Inciso d2: 2D+2F: qué vamos a calcular}\\[2pt]Identificamos la operación y comprobamos sus dimensiones antes de efectuar las cuentas. Las matrices originales permanecen escritas en el enunciado.
\[2D+2F\]
\end{minipage}\par\medskip
\noindent\begin{minipage}{\linewidth}\textbf{27. Multiplicamos cada entrada por el escalar}\\[2pt]La fila y la columna nos dicen qué valores combinar. Hacemos una cuenta por casilla y después reunimos todas las casillas en la matriz resultado.
\par\smallskip {\small\renewcommand{\arraystretch}{1.3}\begin{tabular}{>{\raggedright\arraybackslash}p{3.74cm}>{\raggedright\arraybackslash}p{3.74cm}>{\raggedright\arraybackslash}p{3.74cm}}\hline\rowcolor{black!7} \textbf{Casilla} & \textbf{Cuenta} & \textbf{Resultado}\\\hline Fila 1, columna 1 & $2\cdot(3)$ & $6$\\ Fila 1, columna 2 & $2\cdot(-2)$ & $-4$\\ Fila 2, columna 1 & $2\cdot(2)$ & $4$\\ Fila 2, columna 2 & $2\cdot(4)$ & $8$\\ \hline\end{tabular}}\par\smallskip
\par\smallskip \textbf{Resultado de esta operación}\[\begin{bmatrix}6 & -4\\4 & 8\end{bmatrix}\]
\par\smallskip\textbf{Qué aprendimos.} El resultado tiene 2 filas y 2 columnas.
\end{minipage}\par\medskip
\noindent\begin{minipage}{\linewidth}\textbf{28. Multiplicamos cada entrada por el escalar}\\[2pt]La fila y la columna nos dicen qué valores combinar. Hacemos una cuenta por casilla y después reunimos todas las casillas en la matriz resultado.
\par\smallskip {\small\renewcommand{\arraystretch}{1.3}\begin{tabular}{>{\raggedright\arraybackslash}p{3.74cm}>{\raggedright\arraybackslash}p{3.74cm}>{\raggedright\arraybackslash}p{3.74cm}}\hline\rowcolor{black!7} \textbf{Casilla} & \textbf{Cuenta} & \textbf{Resultado}\\\hline Fila 1, columna 1 & $2\cdot(-4)$ & $-8$\\ Fila 1, columna 2 & $2\cdot(5)$ & $10$\\ Fila 2, columna 1 & $2\cdot(2)$ & $4$\\ Fila 2, columna 2 & $2\cdot(3)$ & $6$\\ \hline\end{tabular}}\par\smallskip
\par\smallskip \textbf{Resultado de esta operación}\[\begin{bmatrix}-8 & 10\\4 & 6\end{bmatrix}\]
\par\smallskip\textbf{Qué aprendimos.} El resultado tiene 2 filas y 2 columnas.
\end{minipage}\par\medskip
\noindent\begin{minipage}{\linewidth}\textbf{29. Operamos las entradas que ocupan la misma posición}\\[2pt]La fila y la columna nos dicen qué valores combinar. Hacemos una cuenta por casilla y después reunimos todas las casillas en la matriz resultado.
\par\smallskip {\small\renewcommand{\arraystretch}{1.3}\begin{tabular}{>{\raggedright\arraybackslash}p{3.74cm}>{\raggedright\arraybackslash}p{3.74cm}>{\raggedright\arraybackslash}p{3.74cm}}\hline\rowcolor{black!7} \textbf{Casilla} & \textbf{Cuenta} & \textbf{Resultado}\\\hline Fila 1, columna 1 & $(6)+(-8)$ & $-2$\\ Fila 1, columna 2 & $(-4)+(10)$ & $6$\\ Fila 2, columna 1 & $(4)+(4)$ & $8$\\ Fila 2, columna 2 & $(8)+(6)$ & $14$\\ \hline\end{tabular}}\par\smallskip
\par\smallskip \textbf{Resultado de esta operación}\[\begin{bmatrix}-2 & 6\\8 & 14\end{bmatrix}\]
\par\smallskip\textbf{Qué aprendimos.} El resultado tiene 2 filas y 2 columnas.
\end{minipage}\par\medskip
\noindent\begin{minipage}{\linewidth}\textbf{30. Resultado del inciso d2}\\[2pt]Reunimos el resultado. Comprueba que cada casilla procede de la operación indicada, sin cambiar su posición.
\par\smallskip \textbf{Resultado}\[\begin{bmatrix}-2 & 6\\8 & 14\end{bmatrix}\]
\end{minipage}\par\medskip
\noindent\begin{minipage}{\linewidth}\textbf{31. Inciso e1: (2+3)D: qué vamos a calcular}\\[2pt]Identificamos la operación y comprobamos sus dimensiones antes de efectuar las cuentas. Las matrices originales permanecen escritas en el enunciado.
\[(2+3)D\]
\end{minipage}\par\medskip
\noindent\begin{minipage}{\linewidth}\textbf{32. Multiplicamos cada entrada por el escalar}\\[2pt]La fila y la columna nos dicen qué valores combinar. Hacemos una cuenta por casilla y después reunimos todas las casillas en la matriz resultado.
\par\smallskip {\small\renewcommand{\arraystretch}{1.3}\begin{tabular}{>{\raggedright\arraybackslash}p{3.74cm}>{\raggedright\arraybackslash}p{3.74cm}>{\raggedright\arraybackslash}p{3.74cm}}\hline\rowcolor{black!7} \textbf{Casilla} & \textbf{Cuenta} & \textbf{Resultado}\\\hline Fila 1, columna 1 & $5\cdot(3)$ & $15$\\ Fila 1, columna 2 & $5\cdot(-2)$ & $-10$\\ Fila 2, columna 1 & $5\cdot(2)$ & $10$\\ Fila 2, columna 2 & $5\cdot(4)$ & $20$\\ \hline\end{tabular}}\par\smallskip
\par\smallskip \textbf{Resultado de esta operación}\[\begin{bmatrix}15 & -10\\10 & 20\end{bmatrix}\]
\par\smallskip\textbf{Qué aprendimos.} El resultado tiene 2 filas y 2 columnas.
\end{minipage}\par\medskip
\noindent\begin{minipage}{\linewidth}\textbf{33. Resultado del inciso e1}\\[2pt]Reunimos el resultado. Comprueba que cada casilla procede de la operación indicada, sin cambiar su posición.
\par\smallskip \textbf{Resultado}\[\begin{bmatrix}15 & -10\\10 & 20\end{bmatrix}\]
\end{minipage}\par\medskip
\noindent\begin{minipage}{\linewidth}\textbf{34. Inciso e2: 2D+3D: qué vamos a calcular}\\[2pt]Identificamos la operación y comprobamos sus dimensiones antes de efectuar las cuentas. Las matrices originales permanecen escritas en el enunciado.
\[2D+3D\]
\end{minipage}\par\medskip
\noindent\begin{minipage}{\linewidth}\textbf{35. Multiplicamos cada entrada por el escalar}\\[2pt]La fila y la columna nos dicen qué valores combinar. Hacemos una cuenta por casilla y después reunimos todas las casillas en la matriz resultado.
\par\smallskip {\small\renewcommand{\arraystretch}{1.3}\begin{tabular}{>{\raggedright\arraybackslash}p{3.74cm}>{\raggedright\arraybackslash}p{3.74cm}>{\raggedright\arraybackslash}p{3.74cm}}\hline\rowcolor{black!7} \textbf{Casilla} & \textbf{Cuenta} & \textbf{Resultado}\\\hline Fila 1, columna 1 & $2\cdot(3)$ & $6$\\ Fila 1, columna 2 & $2\cdot(-2)$ & $-4$\\ Fila 2, columna 1 & $2\cdot(2)$ & $4$\\ Fila 2, columna 2 & $2\cdot(4)$ & $8$\\ \hline\end{tabular}}\par\smallskip
\par\smallskip \textbf{Resultado de esta operación}\[\begin{bmatrix}6 & -4\\4 & 8\end{bmatrix}\]
\par\smallskip\textbf{Qué aprendimos.} El resultado tiene 2 filas y 2 columnas.
\end{minipage}\par\medskip
\noindent\begin{minipage}{\linewidth}\textbf{36. Multiplicamos cada entrada por el escalar}\\[2pt]La fila y la columna nos dicen qué valores combinar. Hacemos una cuenta por casilla y después reunimos todas las casillas en la matriz resultado.
\par\smallskip {\small\renewcommand{\arraystretch}{1.3}\begin{tabular}{>{\raggedright\arraybackslash}p{3.74cm}>{\raggedright\arraybackslash}p{3.74cm}>{\raggedright\arraybackslash}p{3.74cm}}\hline\rowcolor{black!7} \textbf{Casilla} & \textbf{Cuenta} & \textbf{Resultado}\\\hline Fila 1, columna 1 & $3\cdot(3)$ & $9$\\ Fila 1, columna 2 & $3\cdot(-2)$ & $-6$\\ Fila 2, columna 1 & $3\cdot(2)$ & $6$\\ Fila 2, columna 2 & $3\cdot(4)$ & $12$\\ \hline\end{tabular}}\par\smallskip
\par\smallskip \textbf{Resultado de esta operación}\[\begin{bmatrix}9 & -6\\6 & 12\end{bmatrix}\]
\par\smallskip\textbf{Qué aprendimos.} El resultado tiene 2 filas y 2 columnas.
\end{minipage}\par\medskip
\noindent\begin{minipage}{\linewidth}\textbf{37. Operamos las entradas que ocupan la misma posición}\\[2pt]La fila y la columna nos dicen qué valores combinar. Hacemos una cuenta por casilla y después reunimos todas las casillas en la matriz resultado.
\par\smallskip {\small\renewcommand{\arraystretch}{1.3}\begin{tabular}{>{\raggedright\arraybackslash}p{3.74cm}>{\raggedright\arraybackslash}p{3.74cm}>{\raggedright\arraybackslash}p{3.74cm}}\hline\rowcolor{black!7} \textbf{Casilla} & \textbf{Cuenta} & \textbf{Resultado}\\\hline Fila 1, columna 1 & $(6)+(9)$ & $15$\\ Fila 1, columna 2 & $(-4)+(-6)$ & $-10$\\ Fila 2, columna 1 & $(4)+(6)$ & $10$\\ Fila 2, columna 2 & $(8)+(12)$ & $20$\\ \hline\end{tabular}}\par\smallskip
\par\smallskip \textbf{Resultado de esta operación}\[\begin{bmatrix}15 & -10\\10 & 20\end{bmatrix}\]
\par\smallskip\textbf{Qué aprendimos.} El resultado tiene 2 filas y 2 columnas.
\end{minipage}\par\medskip
\noindent\begin{minipage}{\linewidth}\textbf{38. Resultado del inciso e2}\\[2pt]Reunimos el resultado. Comprueba que cada casilla procede de la operación indicada, sin cambiar su posición.
\par\smallskip \textbf{Resultado}\[\begin{bmatrix}15 & -10\\10 & 20\end{bmatrix}\]
\end{minipage}\par\medskip
\noindent\begin{minipage}{\linewidth}\textbf{39. Inciso f: 3(B+D): qué vamos a calcular}\\[2pt]Identificamos la operación y comprobamos sus dimensiones antes de efectuar las cuentas. Las matrices originales permanecen escritas en el enunciado.
\[3(B+D)\]
\end{minipage}\par\medskip
\noindent\begin{minipage}{\linewidth}\textbf{40. Este inciso no está definido}\\[2pt]El primer operando tiene tamaño 3 por 2 y el segundo 2 por 2. La suma o resta no está definida. No se agregan ceros ni se recortan filas para forzarla.
\par\smallskip\textbf{Qué aprendimos.} La comprobación de dimensiones es parte de la solución, no un error del programa.
\end{minipage}\par\medskip
\respuesta{a: 3D+2F: [[1, 4], [10, 18]]; b1: 3(2A): [[6, 12, 18], [12, 6, 24]]; b2: 6A: [[6, 12, 18], [12, 6, 24]]; c1: 3A+2A: [[5, 10, 15], [10, 5, 20]]; c2: 5A: [[5, 10, 15], [10, 5, 20]]; d1: 2(D+F): [[-2, 6], [8, 14]]; d2: 2D+2F: [[-2, 6], [8, 14]]; e1: (2+3)D: [[15, -10], [10, 20]]; e2: 2D+3D: [[15, -10], [10, 20]]; f: 3(B+D): no definida.}
\noindent\begin{minipage}{\linewidth}\centering\includegraphics[width=.96\linewidth,height=6.8cm,keepaspectratio]{../recursos/grafico_1_2-5.png}\\{\small Lectura visual complementaria; la demostración está en las operaciones.}\end{minipage}
\end{multicols*}
\newpage\fancyfoot[L]{\small Sección 1.2 | Ejercicio 6 | Paso a paso}\hypertarget{1.2-6}{}
\pdfbookmark[1]{1.2 / 6}{bm1.2-6}
\titulo{6 | Transposición y dimensiones}
\textbf{Enunciado}\par De ser posible, calcule cada expresión indicada. Se trabaja con las operaciones usuales de matrices reales.
\par\smallskip \begin{minipage}[t]{.47\linewidth}\centering \textbf{A}\[\begin{bmatrix}1 & 2 & 3\\2 & 1 & 4\end{bmatrix}\]\end{minipage}\hfill\begin{minipage}[t]{.47\linewidth}\centering \textbf{B}\[\begin{bmatrix}1 & 0\\2 & 1\\3 & 2\end{bmatrix}\]\end{minipage}\par\smallskip
\par\smallskip \begin{minipage}[t]{.47\linewidth}\centering \textbf{C}\[\begin{bmatrix}3 & -1 & 3\\4 & 1 & 5\\2 & 1 & 3\end{bmatrix}\]\end{minipage}\hfill\begin{minipage}[t]{.47\linewidth}\centering \textbf{D}\[\begin{bmatrix}3 & -2\\2 & 4\end{bmatrix}\]\end{minipage}\par\smallskip
\par\smallskip \begin{minipage}[t]{.47\linewidth}\centering \textbf{E}\[\begin{bmatrix}2 & -4 & 5\\0 & 1 & 4\\3 & 2 & 1\end{bmatrix}\]\end{minipage}\hfill\begin{minipage}[t]{.47\linewidth}\centering \textbf{F}\[\begin{bmatrix}-4 & 5\\2 & 3\end{bmatrix}\]\end{minipage}\par\smallskip
\par\smallskip \begin{minipage}[t]{.47\linewidth}\centering \textbf{O}\[\begin{bmatrix}0 & 0 & 0\\0 & 0 & 0\\0 & 0 & 0\end{bmatrix}\]\end{minipage}\par\smallskip
\textbf{Expresiones del enunciado}\par
a1. \(A^T\)\par
a2. \((A^T)^T\)\par
b1. \((C+E)^T\)\par
b2. \(C^T+E^T\)\par
c. \((2D+3F)^T\)\par
d. \(D-D^T\)\par
e. \(2A^T+B\)\par
f. \((3D-2F)^T\)\par
\begin{multicols*}{2}
\noindent\begin{minipage}{\linewidth}\textbf{1. Antes de calcular: filas, columnas y operaciones}\\[2pt]El tamaño de una matriz se escribe filas por columnas. Para sumar o restar, los dos tamaños deben coincidir exactamente. Multiplicar por un número afecta a cada entrada. Transponer intercambia las posiciones: la entrada de fila i y columna j pasa a la fila j y columna i.
\[(M^T)_{ij}=m_{ji}\]
\par\smallskip\textbf{Qué aprendimos.} Los paréntesis indican qué operación hay que realizar primero. Una suma de tamaños distintos no se vuelve válida por transponerla después.
\end{minipage}\par\medskip
\noindent\begin{minipage}{\linewidth}\textbf{2. Inciso a1: A\textasciicircum{}T: qué vamos a calcular}\\[2pt]Identificamos la operación y comprobamos sus dimensiones antes de efectuar las cuentas. Las matrices originales permanecen escritas en el enunciado.
\[A^T\]
\end{minipage}\par\medskip
\noindent\begin{minipage}{\linewidth}\textbf{3. Transponemos: las filas se convierten en columnas}\\[2pt]No cambiamos los valores. Solo intercambiamos los dos índices de cada entrada. Una matriz de m filas y n columnas produce otra de n filas y m columnas.
\par\smallskip {\small\renewcommand{\arraystretch}{1.3}\begin{tabular}{>{\raggedright\arraybackslash}p{3.74cm}>{\raggedright\arraybackslash}p{3.74cm}>{\raggedright\arraybackslash}p{3.74cm}}\hline\rowcolor{black!7} \textbf{Posición antes} & \textbf{Posición después} & \textbf{Valor conservado}\\\hline (1, 1) & (1, 1) & $1$\\ (1, 2) & (2, 1) & $2$\\ (1, 3) & (3, 1) & $3$\\ (2, 1) & (1, 2) & $2$\\ (2, 2) & (2, 2) & $1$\\ (2, 3) & (3, 2) & $4$\\ \hline\end{tabular}}\par\smallskip
\par\smallskip \textbf{Resultado de esta operación}\[\begin{bmatrix}1 & 2\\2 & 1\\3 & 4\end{bmatrix}\]
\par\smallskip\textbf{Qué aprendimos.} El resultado tiene 3 filas y 2 columnas.
\end{minipage}\par\medskip
\noindent\begin{minipage}{\linewidth}\textbf{4. Resultado del inciso a1}\\[2pt]Reunimos el resultado. Comprueba que cada casilla procede de la operación indicada, sin cambiar su posición.
\par\smallskip \textbf{Resultado}\[\begin{bmatrix}1 & 2\\2 & 1\\3 & 4\end{bmatrix}\]
\end{minipage}\par\medskip
\noindent\begin{minipage}{\linewidth}\textbf{5. Inciso a2: (A\textasciicircum{}T)\textasciicircum{}T: qué vamos a calcular}\\[2pt]Identificamos la operación y comprobamos sus dimensiones antes de efectuar las cuentas. Las matrices originales permanecen escritas en el enunciado.
\[(A^T)^T\]
\end{minipage}\par\medskip
\noindent\begin{minipage}{\linewidth}\textbf{6. Transponemos: las filas se convierten en columnas}\\[2pt]No cambiamos los valores. Solo intercambiamos los dos índices de cada entrada. Una matriz de m filas y n columnas produce otra de n filas y m columnas.
\par\smallskip {\small\renewcommand{\arraystretch}{1.3}\begin{tabular}{>{\raggedright\arraybackslash}p{3.74cm}>{\raggedright\arraybackslash}p{3.74cm}>{\raggedright\arraybackslash}p{3.74cm}}\hline\rowcolor{black!7} \textbf{Posición antes} & \textbf{Posición después} & \textbf{Valor conservado}\\\hline (1, 1) & (1, 1) & $1$\\ (1, 2) & (2, 1) & $2$\\ (1, 3) & (3, 1) & $3$\\ (2, 1) & (1, 2) & $2$\\ (2, 2) & (2, 2) & $1$\\ (2, 3) & (3, 2) & $4$\\ \hline\end{tabular}}\par\smallskip
\par\smallskip \textbf{Resultado de esta operación}\[\begin{bmatrix}1 & 2\\2 & 1\\3 & 4\end{bmatrix}\]
\par\smallskip\textbf{Qué aprendimos.} El resultado tiene 3 filas y 2 columnas.
\end{minipage}\par\medskip
\noindent\begin{minipage}{\linewidth}\textbf{7. Transponemos: las filas se convierten en columnas}\\[2pt]No cambiamos los valores. Solo intercambiamos los dos índices de cada entrada. Una matriz de m filas y n columnas produce otra de n filas y m columnas.
\par\smallskip {\small\renewcommand{\arraystretch}{1.3}\begin{tabular}{>{\raggedright\arraybackslash}p{3.74cm}>{\raggedright\arraybackslash}p{3.74cm}>{\raggedright\arraybackslash}p{3.74cm}}\hline\rowcolor{black!7} \textbf{Posición antes} & \textbf{Posición después} & \textbf{Valor conservado}\\\hline (1, 1) & (1, 1) & $1$\\ (1, 2) & (2, 1) & $2$\\ (2, 1) & (1, 2) & $2$\\ (2, 2) & (2, 2) & $1$\\ (3, 1) & (1, 3) & $3$\\ (3, 2) & (2, 3) & $4$\\ \hline\end{tabular}}\par\smallskip
\par\smallskip \textbf{Resultado de esta operación}\[\begin{bmatrix}1 & 2 & 3\\2 & 1 & 4\end{bmatrix}\]
\par\smallskip\textbf{Qué aprendimos.} El resultado tiene 2 filas y 3 columnas.
\end{minipage}\par\medskip
\noindent\begin{minipage}{\linewidth}\textbf{8. Resultado del inciso a2}\\[2pt]Reunimos el resultado. Comprueba que cada casilla procede de la operación indicada, sin cambiar su posición.
\par\smallskip \textbf{Resultado}\[\begin{bmatrix}1 & 2 & 3\\2 & 1 & 4\end{bmatrix}\]
\end{minipage}\par\medskip
\noindent\begin{minipage}{\linewidth}\textbf{9. Inciso b1: (C+E)\textasciicircum{}T: qué vamos a calcular}\\[2pt]Identificamos la operación y comprobamos sus dimensiones antes de efectuar las cuentas. Las matrices originales permanecen escritas en el enunciado.
\[(C+E)^T\]
\end{minipage}\par\medskip
\noindent\begin{minipage}{\linewidth}\textbf{10. Operamos las entradas que ocupan la misma posición}\\[2pt]La fila y la columna nos dicen qué valores combinar. Hacemos una cuenta por casilla y después reunimos todas las casillas en la matriz resultado.
\par\smallskip {\small\renewcommand{\arraystretch}{1.3}\begin{tabular}{>{\raggedright\arraybackslash}p{3.74cm}>{\raggedright\arraybackslash}p{3.74cm}>{\raggedright\arraybackslash}p{3.74cm}}\hline\rowcolor{black!7} \textbf{Casilla} & \textbf{Cuenta} & \textbf{Resultado}\\\hline Fila 1, columna 1 & $(3)+(2)$ & $5$\\ Fila 1, columna 2 & $(-1)+(-4)$ & $-5$\\ Fila 1, columna 3 & $(3)+(5)$ & $8$\\ Fila 2, columna 1 & $(4)+(0)$ & $4$\\ Fila 2, columna 2 & $(1)+(1)$ & $2$\\ Fila 2, columna 3 & $(5)+(4)$ & $9$\\ Fila 3, columna 1 & $(2)+(3)$ & $5$\\ Fila 3, columna 2 & $(1)+(2)$ & $3$\\ Fila 3, columna 3 & $(3)+(1)$ & $4$\\ \hline\end{tabular}}\par\smallskip
\par\smallskip \textbf{Resultado de esta operación}\[\begin{bmatrix}5 & -5 & 8\\4 & 2 & 9\\5 & 3 & 4\end{bmatrix}\]
\par\smallskip\textbf{Qué aprendimos.} El resultado tiene 3 filas y 3 columnas.
\end{minipage}\par\medskip
\noindent\begin{minipage}{\linewidth}\textbf{11. Transponemos: las filas se convierten en columnas}\\[2pt]No cambiamos los valores. Solo intercambiamos los dos índices de cada entrada. Una matriz de m filas y n columnas produce otra de n filas y m columnas.
\par\smallskip {\small\renewcommand{\arraystretch}{1.3}\begin{tabular}{>{\raggedright\arraybackslash}p{3.74cm}>{\raggedright\arraybackslash}p{3.74cm}>{\raggedright\arraybackslash}p{3.74cm}}\hline\rowcolor{black!7} \textbf{Posición antes} & \textbf{Posición después} & \textbf{Valor conservado}\\\hline (1, 1) & (1, 1) & $5$\\ (1, 2) & (2, 1) & $-5$\\ (1, 3) & (3, 1) & $8$\\ (2, 1) & (1, 2) & $4$\\ (2, 2) & (2, 2) & $2$\\ (2, 3) & (3, 2) & $9$\\ (3, 1) & (1, 3) & $5$\\ (3, 2) & (2, 3) & $3$\\ (3, 3) & (3, 3) & $4$\\ \hline\end{tabular}}\par\smallskip
\par\smallskip \textbf{Resultado de esta operación}\[\begin{bmatrix}5 & 4 & 5\\-5 & 2 & 3\\8 & 9 & 4\end{bmatrix}\]
\par\smallskip\textbf{Qué aprendimos.} El resultado tiene 3 filas y 3 columnas.
\end{minipage}\par\medskip
\noindent\begin{minipage}{\linewidth}\textbf{12. Resultado del inciso b1}\\[2pt]Reunimos el resultado. Comprueba que cada casilla procede de la operación indicada, sin cambiar su posición.
\par\smallskip \textbf{Resultado}\[\begin{bmatrix}5 & 4 & 5\\-5 & 2 & 3\\8 & 9 & 4\end{bmatrix}\]
\end{minipage}\par\medskip
\noindent\begin{minipage}{\linewidth}\textbf{13. Inciso b2: C\textasciicircum{}T+E\textasciicircum{}T: qué vamos a calcular}\\[2pt]Identificamos la operación y comprobamos sus dimensiones antes de efectuar las cuentas. Las matrices originales permanecen escritas en el enunciado.
\[C^T+E^T\]
\end{minipage}\par\medskip
\noindent\begin{minipage}{\linewidth}\textbf{14. Transponemos: las filas se convierten en columnas}\\[2pt]No cambiamos los valores. Solo intercambiamos los dos índices de cada entrada. Una matriz de m filas y n columnas produce otra de n filas y m columnas.
\par\smallskip {\small\renewcommand{\arraystretch}{1.3}\begin{tabular}{>{\raggedright\arraybackslash}p{3.74cm}>{\raggedright\arraybackslash}p{3.74cm}>{\raggedright\arraybackslash}p{3.74cm}}\hline\rowcolor{black!7} \textbf{Posición antes} & \textbf{Posición después} & \textbf{Valor conservado}\\\hline (1, 1) & (1, 1) & $3$\\ (1, 2) & (2, 1) & $-1$\\ (1, 3) & (3, 1) & $3$\\ (2, 1) & (1, 2) & $4$\\ (2, 2) & (2, 2) & $1$\\ (2, 3) & (3, 2) & $5$\\ (3, 1) & (1, 3) & $2$\\ (3, 2) & (2, 3) & $1$\\ (3, 3) & (3, 3) & $3$\\ \hline\end{tabular}}\par\smallskip
\par\smallskip \textbf{Resultado de esta operación}\[\begin{bmatrix}3 & 4 & 2\\-1 & 1 & 1\\3 & 5 & 3\end{bmatrix}\]
\par\smallskip\textbf{Qué aprendimos.} El resultado tiene 3 filas y 3 columnas.
\end{minipage}\par\medskip
\noindent\begin{minipage}{\linewidth}\textbf{15. Transponemos: las filas se convierten en columnas}\\[2pt]No cambiamos los valores. Solo intercambiamos los dos índices de cada entrada. Una matriz de m filas y n columnas produce otra de n filas y m columnas.
\par\smallskip {\small\renewcommand{\arraystretch}{1.3}\begin{tabular}{>{\raggedright\arraybackslash}p{3.74cm}>{\raggedright\arraybackslash}p{3.74cm}>{\raggedright\arraybackslash}p{3.74cm}}\hline\rowcolor{black!7} \textbf{Posición antes} & \textbf{Posición después} & \textbf{Valor conservado}\\\hline (1, 1) & (1, 1) & $2$\\ (1, 2) & (2, 1) & $-4$\\ (1, 3) & (3, 1) & $5$\\ (2, 1) & (1, 2) & $0$\\ (2, 2) & (2, 2) & $1$\\ (2, 3) & (3, 2) & $4$\\ (3, 1) & (1, 3) & $3$\\ (3, 2) & (2, 3) & $2$\\ (3, 3) & (3, 3) & $1$\\ \hline\end{tabular}}\par\smallskip
\par\smallskip \textbf{Resultado de esta operación}\[\begin{bmatrix}2 & 0 & 3\\-4 & 1 & 2\\5 & 4 & 1\end{bmatrix}\]
\par\smallskip\textbf{Qué aprendimos.} El resultado tiene 3 filas y 3 columnas.
\end{minipage}\par\medskip
\noindent\begin{minipage}{\linewidth}\textbf{16. Operamos las entradas que ocupan la misma posición}\\[2pt]La fila y la columna nos dicen qué valores combinar. Hacemos una cuenta por casilla y después reunimos todas las casillas en la matriz resultado.
\par\smallskip {\small\renewcommand{\arraystretch}{1.3}\begin{tabular}{>{\raggedright\arraybackslash}p{3.74cm}>{\raggedright\arraybackslash}p{3.74cm}>{\raggedright\arraybackslash}p{3.74cm}}\hline\rowcolor{black!7} \textbf{Casilla} & \textbf{Cuenta} & \textbf{Resultado}\\\hline Fila 1, columna 1 & $(3)+(2)$ & $5$\\ Fila 1, columna 2 & $(4)+(0)$ & $4$\\ Fila 1, columna 3 & $(2)+(3)$ & $5$\\ Fila 2, columna 1 & $(-1)+(-4)$ & $-5$\\ Fila 2, columna 2 & $(1)+(1)$ & $2$\\ Fila 2, columna 3 & $(1)+(2)$ & $3$\\ Fila 3, columna 1 & $(3)+(5)$ & $8$\\ Fila 3, columna 2 & $(5)+(4)$ & $9$\\ Fila 3, columna 3 & $(3)+(1)$ & $4$\\ \hline\end{tabular}}\par\smallskip
\par\smallskip \textbf{Resultado de esta operación}\[\begin{bmatrix}5 & 4 & 5\\-5 & 2 & 3\\8 & 9 & 4\end{bmatrix}\]
\par\smallskip\textbf{Qué aprendimos.} El resultado tiene 3 filas y 3 columnas.
\end{minipage}\par\medskip
\noindent\begin{minipage}{\linewidth}\textbf{17. Resultado del inciso b2}\\[2pt]Reunimos el resultado. Comprueba que cada casilla procede de la operación indicada, sin cambiar su posición.
\par\smallskip \textbf{Resultado}\[\begin{bmatrix}5 & 4 & 5\\-5 & 2 & 3\\8 & 9 & 4\end{bmatrix}\]
\end{minipage}\par\medskip
\noindent\begin{minipage}{\linewidth}\textbf{18. Inciso c: (2D+3F)\textasciicircum{}T: qué vamos a calcular}\\[2pt]Identificamos la operación y comprobamos sus dimensiones antes de efectuar las cuentas. Las matrices originales permanecen escritas en el enunciado.
\[(2D+3F)^T\]
\end{minipage}\par\medskip
\noindent\begin{minipage}{\linewidth}\textbf{19. Multiplicamos cada entrada por el escalar}\\[2pt]La fila y la columna nos dicen qué valores combinar. Hacemos una cuenta por casilla y después reunimos todas las casillas en la matriz resultado.
\par\smallskip {\small\renewcommand{\arraystretch}{1.3}\begin{tabular}{>{\raggedright\arraybackslash}p{3.74cm}>{\raggedright\arraybackslash}p{3.74cm}>{\raggedright\arraybackslash}p{3.74cm}}\hline\rowcolor{black!7} \textbf{Casilla} & \textbf{Cuenta} & \textbf{Resultado}\\\hline Fila 1, columna 1 & $2\cdot(3)$ & $6$\\ Fila 1, columna 2 & $2\cdot(-2)$ & $-4$\\ Fila 2, columna 1 & $2\cdot(2)$ & $4$\\ Fila 2, columna 2 & $2\cdot(4)$ & $8$\\ \hline\end{tabular}}\par\smallskip
\par\smallskip \textbf{Resultado de esta operación}\[\begin{bmatrix}6 & -4\\4 & 8\end{bmatrix}\]
\par\smallskip\textbf{Qué aprendimos.} El resultado tiene 2 filas y 2 columnas.
\end{minipage}\par\medskip
\noindent\begin{minipage}{\linewidth}\textbf{20. Multiplicamos cada entrada por el escalar}\\[2pt]La fila y la columna nos dicen qué valores combinar. Hacemos una cuenta por casilla y después reunimos todas las casillas en la matriz resultado.
\par\smallskip {\small\renewcommand{\arraystretch}{1.3}\begin{tabular}{>{\raggedright\arraybackslash}p{3.74cm}>{\raggedright\arraybackslash}p{3.74cm}>{\raggedright\arraybackslash}p{3.74cm}}\hline\rowcolor{black!7} \textbf{Casilla} & \textbf{Cuenta} & \textbf{Resultado}\\\hline Fila 1, columna 1 & $3\cdot(-4)$ & $-12$\\ Fila 1, columna 2 & $3\cdot(5)$ & $15$\\ Fila 2, columna 1 & $3\cdot(2)$ & $6$\\ Fila 2, columna 2 & $3\cdot(3)$ & $9$\\ \hline\end{tabular}}\par\smallskip
\par\smallskip \textbf{Resultado de esta operación}\[\begin{bmatrix}-12 & 15\\6 & 9\end{bmatrix}\]
\par\smallskip\textbf{Qué aprendimos.} El resultado tiene 2 filas y 2 columnas.
\end{minipage}\par\medskip
\noindent\begin{minipage}{\linewidth}\textbf{21. Operamos las entradas que ocupan la misma posición}\\[2pt]La fila y la columna nos dicen qué valores combinar. Hacemos una cuenta por casilla y después reunimos todas las casillas en la matriz resultado.
\par\smallskip {\small\renewcommand{\arraystretch}{1.3}\begin{tabular}{>{\raggedright\arraybackslash}p{3.74cm}>{\raggedright\arraybackslash}p{3.74cm}>{\raggedright\arraybackslash}p{3.74cm}}\hline\rowcolor{black!7} \textbf{Casilla} & \textbf{Cuenta} & \textbf{Resultado}\\\hline Fila 1, columna 1 & $(6)+(-12)$ & $-6$\\ Fila 1, columna 2 & $(-4)+(15)$ & $11$\\ Fila 2, columna 1 & $(4)+(6)$ & $10$\\ Fila 2, columna 2 & $(8)+(9)$ & $17$\\ \hline\end{tabular}}\par\smallskip
\par\smallskip \textbf{Resultado de esta operación}\[\begin{bmatrix}-6 & 11\\10 & 17\end{bmatrix}\]
\par\smallskip\textbf{Qué aprendimos.} El resultado tiene 2 filas y 2 columnas.
\end{minipage}\par\medskip
\noindent\begin{minipage}{\linewidth}\textbf{22. Transponemos: las filas se convierten en columnas}\\[2pt]No cambiamos los valores. Solo intercambiamos los dos índices de cada entrada. Una matriz de m filas y n columnas produce otra de n filas y m columnas.
\par\smallskip {\small\renewcommand{\arraystretch}{1.3}\begin{tabular}{>{\raggedright\arraybackslash}p{3.74cm}>{\raggedright\arraybackslash}p{3.74cm}>{\raggedright\arraybackslash}p{3.74cm}}\hline\rowcolor{black!7} \textbf{Posición antes} & \textbf{Posición después} & \textbf{Valor conservado}\\\hline (1, 1) & (1, 1) & $-6$\\ (1, 2) & (2, 1) & $11$\\ (2, 1) & (1, 2) & $10$\\ (2, 2) & (2, 2) & $17$\\ \hline\end{tabular}}\par\smallskip
\par\smallskip \textbf{Resultado de esta operación}\[\begin{bmatrix}-6 & 10\\11 & 17\end{bmatrix}\]
\par\smallskip\textbf{Qué aprendimos.} El resultado tiene 2 filas y 2 columnas.
\end{minipage}\par\medskip
\noindent\begin{minipage}{\linewidth}\textbf{23. Resultado del inciso c}\\[2pt]Reunimos el resultado. Comprueba que cada casilla procede de la operación indicada, sin cambiar su posición.
\par\smallskip \textbf{Resultado}\[\begin{bmatrix}-6 & 10\\11 & 17\end{bmatrix}\]
\end{minipage}\par\medskip
\noindent\begin{minipage}{\linewidth}\textbf{24. Inciso d: D-D\textasciicircum{}T: qué vamos a calcular}\\[2pt]Identificamos la operación y comprobamos sus dimensiones antes de efectuar las cuentas. Las matrices originales permanecen escritas en el enunciado.
\[D-D^T\]
\end{minipage}\par\medskip
\noindent\begin{minipage}{\linewidth}\textbf{25. Transponemos: las filas se convierten en columnas}\\[2pt]No cambiamos los valores. Solo intercambiamos los dos índices de cada entrada. Una matriz de m filas y n columnas produce otra de n filas y m columnas.
\par\smallskip {\small\renewcommand{\arraystretch}{1.3}\begin{tabular}{>{\raggedright\arraybackslash}p{3.74cm}>{\raggedright\arraybackslash}p{3.74cm}>{\raggedright\arraybackslash}p{3.74cm}}\hline\rowcolor{black!7} \textbf{Posición antes} & \textbf{Posición después} & \textbf{Valor conservado}\\\hline (1, 1) & (1, 1) & $3$\\ (1, 2) & (2, 1) & $-2$\\ (2, 1) & (1, 2) & $2$\\ (2, 2) & (2, 2) & $4$\\ \hline\end{tabular}}\par\smallskip
\par\smallskip \textbf{Resultado de esta operación}\[\begin{bmatrix}3 & 2\\-2 & 4\end{bmatrix}\]
\par\smallskip\textbf{Qué aprendimos.} El resultado tiene 2 filas y 2 columnas.
\end{minipage}\par\medskip
\noindent\begin{minipage}{\linewidth}\textbf{26. Operamos las entradas que ocupan la misma posición}\\[2pt]La fila y la columna nos dicen qué valores combinar. Hacemos una cuenta por casilla y después reunimos todas las casillas en la matriz resultado.
\par\smallskip {\small\renewcommand{\arraystretch}{1.3}\begin{tabular}{>{\raggedright\arraybackslash}p{3.74cm}>{\raggedright\arraybackslash}p{3.74cm}>{\raggedright\arraybackslash}p{3.74cm}}\hline\rowcolor{black!7} \textbf{Casilla} & \textbf{Cuenta} & \textbf{Resultado}\\\hline Fila 1, columna 1 & $(3)-(3)$ & $0$\\ Fila 1, columna 2 & $(-2)-(2)$ & $-4$\\ Fila 2, columna 1 & $(2)-(-2)$ & $4$\\ Fila 2, columna 2 & $(4)-(4)$ & $0$\\ \hline\end{tabular}}\par\smallskip
\par\smallskip \textbf{Resultado de esta operación}\[\begin{bmatrix}0 & -4\\4 & 0\end{bmatrix}\]
\par\smallskip\textbf{Qué aprendimos.} El resultado tiene 2 filas y 2 columnas.
\end{minipage}\par\medskip
\noindent\begin{minipage}{\linewidth}\textbf{27. Resultado del inciso d}\\[2pt]Reunimos el resultado. Comprueba que cada casilla procede de la operación indicada, sin cambiar su posición.
\par\smallskip \textbf{Resultado}\[\begin{bmatrix}0 & -4\\4 & 0\end{bmatrix}\]
\end{minipage}\par\medskip
\noindent\begin{minipage}{\linewidth}\textbf{28. Inciso e: 2A\textasciicircum{}T+B: qué vamos a calcular}\\[2pt]Identificamos la operación y comprobamos sus dimensiones antes de efectuar las cuentas. Las matrices originales permanecen escritas en el enunciado.
\[2A^T+B\]
\end{minipage}\par\medskip
\noindent\begin{minipage}{\linewidth}\textbf{29. Transponemos: las filas se convierten en columnas}\\[2pt]No cambiamos los valores. Solo intercambiamos los dos índices de cada entrada. Una matriz de m filas y n columnas produce otra de n filas y m columnas.
\par\smallskip {\small\renewcommand{\arraystretch}{1.3}\begin{tabular}{>{\raggedright\arraybackslash}p{3.74cm}>{\raggedright\arraybackslash}p{3.74cm}>{\raggedright\arraybackslash}p{3.74cm}}\hline\rowcolor{black!7} \textbf{Posición antes} & \textbf{Posición después} & \textbf{Valor conservado}\\\hline (1, 1) & (1, 1) & $1$\\ (1, 2) & (2, 1) & $2$\\ (1, 3) & (3, 1) & $3$\\ (2, 1) & (1, 2) & $2$\\ (2, 2) & (2, 2) & $1$\\ (2, 3) & (3, 2) & $4$\\ \hline\end{tabular}}\par\smallskip
\par\smallskip \textbf{Resultado de esta operación}\[\begin{bmatrix}1 & 2\\2 & 1\\3 & 4\end{bmatrix}\]
\par\smallskip\textbf{Qué aprendimos.} El resultado tiene 3 filas y 2 columnas.
\end{minipage}\par\medskip
\noindent\begin{minipage}{\linewidth}\textbf{30. Multiplicamos cada entrada por el escalar}\\[2pt]La fila y la columna nos dicen qué valores combinar. Hacemos una cuenta por casilla y después reunimos todas las casillas en la matriz resultado.
\par\smallskip {\small\renewcommand{\arraystretch}{1.3}\begin{tabular}{>{\raggedright\arraybackslash}p{3.74cm}>{\raggedright\arraybackslash}p{3.74cm}>{\raggedright\arraybackslash}p{3.74cm}}\hline\rowcolor{black!7} \textbf{Casilla} & \textbf{Cuenta} & \textbf{Resultado}\\\hline Fila 1, columna 1 & $2\cdot(1)$ & $2$\\ Fila 1, columna 2 & $2\cdot(2)$ & $4$\\ Fila 2, columna 1 & $2\cdot(2)$ & $4$\\ Fila 2, columna 2 & $2\cdot(1)$ & $2$\\ Fila 3, columna 1 & $2\cdot(3)$ & $6$\\ Fila 3, columna 2 & $2\cdot(4)$ & $8$\\ \hline\end{tabular}}\par\smallskip
\par\smallskip \textbf{Resultado de esta operación}\[\begin{bmatrix}2 & 4\\4 & 2\\6 & 8\end{bmatrix}\]
\par\smallskip\textbf{Qué aprendimos.} El resultado tiene 3 filas y 2 columnas.
\end{minipage}\par\medskip
\noindent\begin{minipage}{\linewidth}\textbf{31. Operamos las entradas que ocupan la misma posición}\\[2pt]La fila y la columna nos dicen qué valores combinar. Hacemos una cuenta por casilla y después reunimos todas las casillas en la matriz resultado.
\par\smallskip {\small\renewcommand{\arraystretch}{1.3}\begin{tabular}{>{\raggedright\arraybackslash}p{3.74cm}>{\raggedright\arraybackslash}p{3.74cm}>{\raggedright\arraybackslash}p{3.74cm}}\hline\rowcolor{black!7} \textbf{Casilla} & \textbf{Cuenta} & \textbf{Resultado}\\\hline Fila 1, columna 1 & $(2)+(1)$ & $3$\\ Fila 1, columna 2 & $(4)+(0)$ & $4$\\ Fila 2, columna 1 & $(4)+(2)$ & $6$\\ Fila 2, columna 2 & $(2)+(1)$ & $3$\\ Fila 3, columna 1 & $(6)+(3)$ & $9$\\ Fila 3, columna 2 & $(8)+(2)$ & $10$\\ \hline\end{tabular}}\par\smallskip
\par\smallskip \textbf{Resultado de esta operación}\[\begin{bmatrix}3 & 4\\6 & 3\\9 & 10\end{bmatrix}\]
\par\smallskip\textbf{Qué aprendimos.} El resultado tiene 3 filas y 2 columnas.
\end{minipage}\par\medskip
\noindent\begin{minipage}{\linewidth}\textbf{32. Resultado del inciso e}\\[2pt]Reunimos el resultado. Comprueba que cada casilla procede de la operación indicada, sin cambiar su posición.
\par\smallskip \textbf{Resultado}\[\begin{bmatrix}3 & 4\\6 & 3\\9 & 10\end{bmatrix}\]
\end{minipage}\par\medskip
\noindent\begin{minipage}{\linewidth}\textbf{33. Inciso f: (3D-2F)\textasciicircum{}T: qué vamos a calcular}\\[2pt]Identificamos la operación y comprobamos sus dimensiones antes de efectuar las cuentas. Las matrices originales permanecen escritas en el enunciado.
\[(3D-2F)^T\]
\end{minipage}\par\medskip
\noindent\begin{minipage}{\linewidth}\textbf{34. Multiplicamos cada entrada por el escalar}\\[2pt]La fila y la columna nos dicen qué valores combinar. Hacemos una cuenta por casilla y después reunimos todas las casillas en la matriz resultado.
\par\smallskip {\small\renewcommand{\arraystretch}{1.3}\begin{tabular}{>{\raggedright\arraybackslash}p{3.74cm}>{\raggedright\arraybackslash}p{3.74cm}>{\raggedright\arraybackslash}p{3.74cm}}\hline\rowcolor{black!7} \textbf{Casilla} & \textbf{Cuenta} & \textbf{Resultado}\\\hline Fila 1, columna 1 & $3\cdot(3)$ & $9$\\ Fila 1, columna 2 & $3\cdot(-2)$ & $-6$\\ Fila 2, columna 1 & $3\cdot(2)$ & $6$\\ Fila 2, columna 2 & $3\cdot(4)$ & $12$\\ \hline\end{tabular}}\par\smallskip
\par\smallskip \textbf{Resultado de esta operación}\[\begin{bmatrix}9 & -6\\6 & 12\end{bmatrix}\]
\par\smallskip\textbf{Qué aprendimos.} El resultado tiene 2 filas y 2 columnas.
\end{minipage}\par\medskip
\noindent\begin{minipage}{\linewidth}\textbf{35. Multiplicamos cada entrada por el escalar}\\[2pt]La fila y la columna nos dicen qué valores combinar. Hacemos una cuenta por casilla y después reunimos todas las casillas en la matriz resultado.
\par\smallskip {\small\renewcommand{\arraystretch}{1.3}\begin{tabular}{>{\raggedright\arraybackslash}p{3.74cm}>{\raggedright\arraybackslash}p{3.74cm}>{\raggedright\arraybackslash}p{3.74cm}}\hline\rowcolor{black!7} \textbf{Casilla} & \textbf{Cuenta} & \textbf{Resultado}\\\hline Fila 1, columna 1 & $2\cdot(-4)$ & $-8$\\ Fila 1, columna 2 & $2\cdot(5)$ & $10$\\ Fila 2, columna 1 & $2\cdot(2)$ & $4$\\ Fila 2, columna 2 & $2\cdot(3)$ & $6$\\ \hline\end{tabular}}\par\smallskip
\par\smallskip \textbf{Resultado de esta operación}\[\begin{bmatrix}-8 & 10\\4 & 6\end{bmatrix}\]
\par\smallskip\textbf{Qué aprendimos.} El resultado tiene 2 filas y 2 columnas.
\end{minipage}\par\medskip
\noindent\begin{minipage}{\linewidth}\textbf{36. Operamos las entradas que ocupan la misma posición}\\[2pt]La fila y la columna nos dicen qué valores combinar. Hacemos una cuenta por casilla y después reunimos todas las casillas en la matriz resultado.
\par\smallskip {\small\renewcommand{\arraystretch}{1.3}\begin{tabular}{>{\raggedright\arraybackslash}p{3.74cm}>{\raggedright\arraybackslash}p{3.74cm}>{\raggedright\arraybackslash}p{3.74cm}}\hline\rowcolor{black!7} \textbf{Casilla} & \textbf{Cuenta} & \textbf{Resultado}\\\hline Fila 1, columna 1 & $(9)-(-8)$ & $17$\\ Fila 1, columna 2 & $(-6)-(10)$ & $-16$\\ Fila 2, columna 1 & $(6)-(4)$ & $2$\\ Fila 2, columna 2 & $(12)-(6)$ & $6$\\ \hline\end{tabular}}\par\smallskip
\par\smallskip \textbf{Resultado de esta operación}\[\begin{bmatrix}17 & -16\\2 & 6\end{bmatrix}\]
\par\smallskip\textbf{Qué aprendimos.} El resultado tiene 2 filas y 2 columnas.
\end{minipage}\par\medskip
\noindent\begin{minipage}{\linewidth}\textbf{37. Transponemos: las filas se convierten en columnas}\\[2pt]No cambiamos los valores. Solo intercambiamos los dos índices de cada entrada. Una matriz de m filas y n columnas produce otra de n filas y m columnas.
\par\smallskip {\small\renewcommand{\arraystretch}{1.3}\begin{tabular}{>{\raggedright\arraybackslash}p{3.74cm}>{\raggedright\arraybackslash}p{3.74cm}>{\raggedright\arraybackslash}p{3.74cm}}\hline\rowcolor{black!7} \textbf{Posición antes} & \textbf{Posición después} & \textbf{Valor conservado}\\\hline (1, 1) & (1, 1) & $17$\\ (1, 2) & (2, 1) & $-16$\\ (2, 1) & (1, 2) & $2$\\ (2, 2) & (2, 2) & $6$\\ \hline\end{tabular}}\par\smallskip
\par\smallskip \textbf{Resultado de esta operación}\[\begin{bmatrix}17 & 2\\-16 & 6\end{bmatrix}\]
\par\smallskip\textbf{Qué aprendimos.} El resultado tiene 2 filas y 2 columnas.
\end{minipage}\par\medskip
\noindent\begin{minipage}{\linewidth}\textbf{38. Resultado del inciso f}\\[2pt]Reunimos el resultado. Comprueba que cada casilla procede de la operación indicada, sin cambiar su posición.
\par\smallskip \textbf{Resultado}\[\begin{bmatrix}17 & 2\\-16 & 6\end{bmatrix}\]
\end{minipage}\par\medskip
\respuesta{a1: A\textasciicircum{}T: [[1, 2], [2, 1], [3, 4]]; a2: (A\textasciicircum{}T)\textasciicircum{}T: [[1, 2, 3], [2, 1, 4]]; b1: (C+E)\textasciicircum{}T: [[5, 4, 5], [-5, 2, 3], [8, 9, 4]]; b2: C\textasciicircum{}T+E\textasciicircum{}T: [[5, 4, 5], [-5, 2, 3], [8, 9, 4]]; c: (2D+3F)\textasciicircum{}T: [[-6, 10], [11, 17]]; d: D-D\textasciicircum{}T: [[0, -4], [4, 0]]; e: 2A\textasciicircum{}T+B: [[3, 4], [6, 3], [9, 10]]; f: (3D-2F)\textasciicircum{}T: [[17, 2], [-16, 6]].}
\noindent\begin{minipage}{\linewidth}\centering\includegraphics[width=.96\linewidth,height=6.8cm,keepaspectratio]{../recursos/grafico_1_2-6.png}\\{\small Lectura visual complementaria; la demostración está en las operaciones.}\end{minipage}
\end{multicols*}
\newpage\fancyfoot[L]{\small Sección 1.2 | Ejercicio 7 | Paso a paso}\hypertarget{1.2-7}{}
\pdfbookmark[1]{1.2 / 7}{bm1.2-7}
\titulo{7 | Transposición y dimensiones}
\textbf{Enunciado}\par De ser posible, calcule cada expresión indicada. Se trabaja con las operaciones usuales de matrices reales.
\par\smallskip \begin{minipage}[t]{.47\linewidth}\centering \textbf{A}\[\begin{bmatrix}1 & 2 & 3\\2 & 1 & 4\end{bmatrix}\]\end{minipage}\hfill\begin{minipage}[t]{.47\linewidth}\centering \textbf{B}\[\begin{bmatrix}1 & 0\\2 & 1\\3 & 2\end{bmatrix}\]\end{minipage}\par\smallskip
\par\smallskip \begin{minipage}[t]{.47\linewidth}\centering \textbf{C}\[\begin{bmatrix}3 & -1 & 3\\4 & 1 & 5\\2 & 1 & 3\end{bmatrix}\]\end{minipage}\hfill\begin{minipage}[t]{.47\linewidth}\centering \textbf{D}\[\begin{bmatrix}3 & -2\\2 & 4\end{bmatrix}\]\end{minipage}\par\smallskip
\par\smallskip \begin{minipage}[t]{.47\linewidth}\centering \textbf{E}\[\begin{bmatrix}2 & -4 & 5\\0 & 1 & 4\\3 & 2 & 1\end{bmatrix}\]\end{minipage}\hfill\begin{minipage}[t]{.47\linewidth}\centering \textbf{F}\[\begin{bmatrix}-4 & 5\\2 & 3\end{bmatrix}\]\end{minipage}\par\smallskip
\par\smallskip \begin{minipage}[t]{.47\linewidth}\centering \textbf{O}\[\begin{bmatrix}0 & 0 & 0\\0 & 0 & 0\\0 & 0 & 0\end{bmatrix}\]\end{minipage}\par\smallskip
\textbf{Expresiones del enunciado}\par
a. \((2A)^T\)\par
b. \((A-B)^T\)\par
c. \((3B^T-2A)^T\)\par
d. \((3A^T-5B)^T\)\par
e1. \((-A)^T\)\par
e2. \(-(A^T)\)\par
f. \((C+E+F^T)^T\)\par
\begin{multicols*}{2}
\noindent\begin{minipage}{\linewidth}\textbf{1. Antes de calcular: filas, columnas y operaciones}\\[2pt]El tamaño de una matriz se escribe filas por columnas. Para sumar o restar, los dos tamaños deben coincidir exactamente. Multiplicar por un número afecta a cada entrada. Transponer intercambia las posiciones: la entrada de fila i y columna j pasa a la fila j y columna i.
\[(M^T)_{ij}=m_{ji}\]
\par\smallskip\textbf{Qué aprendimos.} Los paréntesis indican qué operación hay que realizar primero. Una suma de tamaños distintos no se vuelve válida por transponerla después.
\end{minipage}\par\medskip
\noindent\begin{minipage}{\linewidth}\textbf{2. Inciso a: (2A)\textasciicircum{}T: qué vamos a calcular}\\[2pt]Identificamos la operación y comprobamos sus dimensiones antes de efectuar las cuentas. Las matrices originales permanecen escritas en el enunciado.
\[(2A)^T\]
\end{minipage}\par\medskip
\noindent\begin{minipage}{\linewidth}\textbf{3. Multiplicamos cada entrada por el escalar}\\[2pt]La fila y la columna nos dicen qué valores combinar. Hacemos una cuenta por casilla y después reunimos todas las casillas en la matriz resultado.
\par\smallskip {\small\renewcommand{\arraystretch}{1.3}\begin{tabular}{>{\raggedright\arraybackslash}p{3.74cm}>{\raggedright\arraybackslash}p{3.74cm}>{\raggedright\arraybackslash}p{3.74cm}}\hline\rowcolor{black!7} \textbf{Casilla} & \textbf{Cuenta} & \textbf{Resultado}\\\hline Fila 1, columna 1 & $2\cdot(1)$ & $2$\\ Fila 1, columna 2 & $2\cdot(2)$ & $4$\\ Fila 1, columna 3 & $2\cdot(3)$ & $6$\\ Fila 2, columna 1 & $2\cdot(2)$ & $4$\\ Fila 2, columna 2 & $2\cdot(1)$ & $2$\\ Fila 2, columna 3 & $2\cdot(4)$ & $8$\\ \hline\end{tabular}}\par\smallskip
\par\smallskip \textbf{Resultado de esta operación}\[\begin{bmatrix}2 & 4 & 6\\4 & 2 & 8\end{bmatrix}\]
\par\smallskip\textbf{Qué aprendimos.} El resultado tiene 2 filas y 3 columnas.
\end{minipage}\par\medskip
\noindent\begin{minipage}{\linewidth}\textbf{4. Transponemos: las filas se convierten en columnas}\\[2pt]No cambiamos los valores. Solo intercambiamos los dos índices de cada entrada. Una matriz de m filas y n columnas produce otra de n filas y m columnas.
\par\smallskip {\small\renewcommand{\arraystretch}{1.3}\begin{tabular}{>{\raggedright\arraybackslash}p{3.74cm}>{\raggedright\arraybackslash}p{3.74cm}>{\raggedright\arraybackslash}p{3.74cm}}\hline\rowcolor{black!7} \textbf{Posición antes} & \textbf{Posición después} & \textbf{Valor conservado}\\\hline (1, 1) & (1, 1) & $2$\\ (1, 2) & (2, 1) & $4$\\ (1, 3) & (3, 1) & $6$\\ (2, 1) & (1, 2) & $4$\\ (2, 2) & (2, 2) & $2$\\ (2, 3) & (3, 2) & $8$\\ \hline\end{tabular}}\par\smallskip
\par\smallskip \textbf{Resultado de esta operación}\[\begin{bmatrix}2 & 4\\4 & 2\\6 & 8\end{bmatrix}\]
\par\smallskip\textbf{Qué aprendimos.} El resultado tiene 3 filas y 2 columnas.
\end{minipage}\par\medskip
\noindent\begin{minipage}{\linewidth}\textbf{5. Resultado del inciso a}\\[2pt]Reunimos el resultado. Comprueba que cada casilla procede de la operación indicada, sin cambiar su posición.
\par\smallskip \textbf{Resultado}\[\begin{bmatrix}2 & 4\\4 & 2\\6 & 8\end{bmatrix}\]
\end{minipage}\par\medskip
\noindent\begin{minipage}{\linewidth}\textbf{6. Inciso b: (A-B)\textasciicircum{}T: qué vamos a calcular}\\[2pt]Identificamos la operación y comprobamos sus dimensiones antes de efectuar las cuentas. Las matrices originales permanecen escritas en el enunciado.
\[(A-B)^T\]
\end{minipage}\par\medskip
\noindent\begin{minipage}{\linewidth}\textbf{7. Este inciso no está definido}\\[2pt]El primer operando tiene tamaño 2 por 3 y el segundo 3 por 2. La suma o resta no está definida. No se agregan ceros ni se recortan filas para forzarla.
\par\smallskip\textbf{Qué aprendimos.} La comprobación de dimensiones es parte de la solución, no un error del programa.
\end{minipage}\par\medskip
\noindent\begin{minipage}{\linewidth}\textbf{8. Inciso c: (3B\textasciicircum{}T-2A)\textasciicircum{}T: qué vamos a calcular}\\[2pt]Identificamos la operación y comprobamos sus dimensiones antes de efectuar las cuentas. Las matrices originales permanecen escritas en el enunciado.
\[(3B^T-2A)^T\]
\end{minipage}\par\medskip
\noindent\begin{minipage}{\linewidth}\textbf{9. Transponemos: las filas se convierten en columnas}\\[2pt]No cambiamos los valores. Solo intercambiamos los dos índices de cada entrada. Una matriz de m filas y n columnas produce otra de n filas y m columnas.
\par\smallskip {\small\renewcommand{\arraystretch}{1.3}\begin{tabular}{>{\raggedright\arraybackslash}p{3.74cm}>{\raggedright\arraybackslash}p{3.74cm}>{\raggedright\arraybackslash}p{3.74cm}}\hline\rowcolor{black!7} \textbf{Posición antes} & \textbf{Posición después} & \textbf{Valor conservado}\\\hline (1, 1) & (1, 1) & $1$\\ (1, 2) & (2, 1) & $0$\\ (2, 1) & (1, 2) & $2$\\ (2, 2) & (2, 2) & $1$\\ (3, 1) & (1, 3) & $3$\\ (3, 2) & (2, 3) & $2$\\ \hline\end{tabular}}\par\smallskip
\par\smallskip \textbf{Resultado de esta operación}\[\begin{bmatrix}1 & 2 & 3\\0 & 1 & 2\end{bmatrix}\]
\par\smallskip\textbf{Qué aprendimos.} El resultado tiene 2 filas y 3 columnas.
\end{minipage}\par\medskip
\noindent\begin{minipage}{\linewidth}\textbf{10. Multiplicamos cada entrada por el escalar}\\[2pt]La fila y la columna nos dicen qué valores combinar. Hacemos una cuenta por casilla y después reunimos todas las casillas en la matriz resultado.
\par\smallskip {\small\renewcommand{\arraystretch}{1.3}\begin{tabular}{>{\raggedright\arraybackslash}p{3.74cm}>{\raggedright\arraybackslash}p{3.74cm}>{\raggedright\arraybackslash}p{3.74cm}}\hline\rowcolor{black!7} \textbf{Casilla} & \textbf{Cuenta} & \textbf{Resultado}\\\hline Fila 1, columna 1 & $3\cdot(1)$ & $3$\\ Fila 1, columna 2 & $3\cdot(2)$ & $6$\\ Fila 1, columna 3 & $3\cdot(3)$ & $9$\\ Fila 2, columna 1 & $3\cdot(0)$ & $0$\\ Fila 2, columna 2 & $3\cdot(1)$ & $3$\\ Fila 2, columna 3 & $3\cdot(2)$ & $6$\\ \hline\end{tabular}}\par\smallskip
\par\smallskip \textbf{Resultado de esta operación}\[\begin{bmatrix}3 & 6 & 9\\0 & 3 & 6\end{bmatrix}\]
\par\smallskip\textbf{Qué aprendimos.} El resultado tiene 2 filas y 3 columnas.
\end{minipage}\par\medskip
\noindent\begin{minipage}{\linewidth}\textbf{11. Multiplicamos cada entrada por el escalar}\\[2pt]La fila y la columna nos dicen qué valores combinar. Hacemos una cuenta por casilla y después reunimos todas las casillas en la matriz resultado.
\par\smallskip {\small\renewcommand{\arraystretch}{1.3}\begin{tabular}{>{\raggedright\arraybackslash}p{3.74cm}>{\raggedright\arraybackslash}p{3.74cm}>{\raggedright\arraybackslash}p{3.74cm}}\hline\rowcolor{black!7} \textbf{Casilla} & \textbf{Cuenta} & \textbf{Resultado}\\\hline Fila 1, columna 1 & $2\cdot(1)$ & $2$\\ Fila 1, columna 2 & $2\cdot(2)$ & $4$\\ Fila 1, columna 3 & $2\cdot(3)$ & $6$\\ Fila 2, columna 1 & $2\cdot(2)$ & $4$\\ Fila 2, columna 2 & $2\cdot(1)$ & $2$\\ Fila 2, columna 3 & $2\cdot(4)$ & $8$\\ \hline\end{tabular}}\par\smallskip
\par\smallskip \textbf{Resultado de esta operación}\[\begin{bmatrix}2 & 4 & 6\\4 & 2 & 8\end{bmatrix}\]
\par\smallskip\textbf{Qué aprendimos.} El resultado tiene 2 filas y 3 columnas.
\end{minipage}\par\medskip
\noindent\begin{minipage}{\linewidth}\textbf{12. Operamos las entradas que ocupan la misma posición}\\[2pt]La fila y la columna nos dicen qué valores combinar. Hacemos una cuenta por casilla y después reunimos todas las casillas en la matriz resultado.
\par\smallskip {\small\renewcommand{\arraystretch}{1.3}\begin{tabular}{>{\raggedright\arraybackslash}p{3.74cm}>{\raggedright\arraybackslash}p{3.74cm}>{\raggedright\arraybackslash}p{3.74cm}}\hline\rowcolor{black!7} \textbf{Casilla} & \textbf{Cuenta} & \textbf{Resultado}\\\hline Fila 1, columna 1 & $(3)-(2)$ & $1$\\ Fila 1, columna 2 & $(6)-(4)$ & $2$\\ Fila 1, columna 3 & $(9)-(6)$ & $3$\\ Fila 2, columna 1 & $(0)-(4)$ & $-4$\\ Fila 2, columna 2 & $(3)-(2)$ & $1$\\ Fila 2, columna 3 & $(6)-(8)$ & $-2$\\ \hline\end{tabular}}\par\smallskip
\par\smallskip \textbf{Resultado de esta operación}\[\begin{bmatrix}1 & 2 & 3\\-4 & 1 & -2\end{bmatrix}\]
\par\smallskip\textbf{Qué aprendimos.} El resultado tiene 2 filas y 3 columnas.
\end{minipage}\par\medskip
\noindent\begin{minipage}{\linewidth}\textbf{13. Transponemos: las filas se convierten en columnas}\\[2pt]No cambiamos los valores. Solo intercambiamos los dos índices de cada entrada. Una matriz de m filas y n columnas produce otra de n filas y m columnas.
\par\smallskip {\small\renewcommand{\arraystretch}{1.3}\begin{tabular}{>{\raggedright\arraybackslash}p{3.74cm}>{\raggedright\arraybackslash}p{3.74cm}>{\raggedright\arraybackslash}p{3.74cm}}\hline\rowcolor{black!7} \textbf{Posición antes} & \textbf{Posición después} & \textbf{Valor conservado}\\\hline (1, 1) & (1, 1) & $1$\\ (1, 2) & (2, 1) & $2$\\ (1, 3) & (3, 1) & $3$\\ (2, 1) & (1, 2) & $-4$\\ (2, 2) & (2, 2) & $1$\\ (2, 3) & (3, 2) & $-2$\\ \hline\end{tabular}}\par\smallskip
\par\smallskip \textbf{Resultado de esta operación}\[\begin{bmatrix}1 & -4\\2 & 1\\3 & -2\end{bmatrix}\]
\par\smallskip\textbf{Qué aprendimos.} El resultado tiene 3 filas y 2 columnas.
\end{minipage}\par\medskip
\noindent\begin{minipage}{\linewidth}\textbf{14. Resultado del inciso c}\\[2pt]Reunimos el resultado. Comprueba que cada casilla procede de la operación indicada, sin cambiar su posición.
\par\smallskip \textbf{Resultado}\[\begin{bmatrix}1 & -4\\2 & 1\\3 & -2\end{bmatrix}\]
\end{minipage}\par\medskip
\noindent\begin{minipage}{\linewidth}\textbf{15. Inciso d: (3A\textasciicircum{}T-5B)\textasciicircum{}T: qué vamos a calcular}\\[2pt]Identificamos la operación y comprobamos sus dimensiones antes de efectuar las cuentas. Las matrices originales permanecen escritas en el enunciado.
\[(3A^T-5B)^T\]
\end{minipage}\par\medskip
\noindent\begin{minipage}{\linewidth}\textbf{16. Transponemos: las filas se convierten en columnas}\\[2pt]No cambiamos los valores. Solo intercambiamos los dos índices de cada entrada. Una matriz de m filas y n columnas produce otra de n filas y m columnas.
\par\smallskip {\small\renewcommand{\arraystretch}{1.3}\begin{tabular}{>{\raggedright\arraybackslash}p{3.74cm}>{\raggedright\arraybackslash}p{3.74cm}>{\raggedright\arraybackslash}p{3.74cm}}\hline\rowcolor{black!7} \textbf{Posición antes} & \textbf{Posición después} & \textbf{Valor conservado}\\\hline (1, 1) & (1, 1) & $1$\\ (1, 2) & (2, 1) & $2$\\ (1, 3) & (3, 1) & $3$\\ (2, 1) & (1, 2) & $2$\\ (2, 2) & (2, 2) & $1$\\ (2, 3) & (3, 2) & $4$\\ \hline\end{tabular}}\par\smallskip
\par\smallskip \textbf{Resultado de esta operación}\[\begin{bmatrix}1 & 2\\2 & 1\\3 & 4\end{bmatrix}\]
\par\smallskip\textbf{Qué aprendimos.} El resultado tiene 3 filas y 2 columnas.
\end{minipage}\par\medskip
\noindent\begin{minipage}{\linewidth}\textbf{17. Multiplicamos cada entrada por el escalar}\\[2pt]La fila y la columna nos dicen qué valores combinar. Hacemos una cuenta por casilla y después reunimos todas las casillas en la matriz resultado.
\par\smallskip {\small\renewcommand{\arraystretch}{1.3}\begin{tabular}{>{\raggedright\arraybackslash}p{3.74cm}>{\raggedright\arraybackslash}p{3.74cm}>{\raggedright\arraybackslash}p{3.74cm}}\hline\rowcolor{black!7} \textbf{Casilla} & \textbf{Cuenta} & \textbf{Resultado}\\\hline Fila 1, columna 1 & $3\cdot(1)$ & $3$\\ Fila 1, columna 2 & $3\cdot(2)$ & $6$\\ Fila 2, columna 1 & $3\cdot(2)$ & $6$\\ Fila 2, columna 2 & $3\cdot(1)$ & $3$\\ Fila 3, columna 1 & $3\cdot(3)$ & $9$\\ Fila 3, columna 2 & $3\cdot(4)$ & $12$\\ \hline\end{tabular}}\par\smallskip
\par\smallskip \textbf{Resultado de esta operación}\[\begin{bmatrix}3 & 6\\6 & 3\\9 & 12\end{bmatrix}\]
\par\smallskip\textbf{Qué aprendimos.} El resultado tiene 3 filas y 2 columnas.
\end{minipage}\par\medskip
\noindent\begin{minipage}{\linewidth}\textbf{18. Multiplicamos cada entrada por el escalar}\\[2pt]La fila y la columna nos dicen qué valores combinar. Hacemos una cuenta por casilla y después reunimos todas las casillas en la matriz resultado.
\par\smallskip {\small\renewcommand{\arraystretch}{1.3}\begin{tabular}{>{\raggedright\arraybackslash}p{3.74cm}>{\raggedright\arraybackslash}p{3.74cm}>{\raggedright\arraybackslash}p{3.74cm}}\hline\rowcolor{black!7} \textbf{Casilla} & \textbf{Cuenta} & \textbf{Resultado}\\\hline Fila 1, columna 1 & $5\cdot(1)$ & $5$\\ Fila 1, columna 2 & $5\cdot(0)$ & $0$\\ Fila 2, columna 1 & $5\cdot(2)$ & $10$\\ Fila 2, columna 2 & $5\cdot(1)$ & $5$\\ Fila 3, columna 1 & $5\cdot(3)$ & $15$\\ Fila 3, columna 2 & $5\cdot(2)$ & $10$\\ \hline\end{tabular}}\par\smallskip
\par\smallskip \textbf{Resultado de esta operación}\[\begin{bmatrix}5 & 0\\10 & 5\\15 & 10\end{bmatrix}\]
\par\smallskip\textbf{Qué aprendimos.} El resultado tiene 3 filas y 2 columnas.
\end{minipage}\par\medskip
\noindent\begin{minipage}{\linewidth}\textbf{19. Operamos las entradas que ocupan la misma posición}\\[2pt]La fila y la columna nos dicen qué valores combinar. Hacemos una cuenta por casilla y después reunimos todas las casillas en la matriz resultado.
\par\smallskip {\small\renewcommand{\arraystretch}{1.3}\begin{tabular}{>{\raggedright\arraybackslash}p{3.74cm}>{\raggedright\arraybackslash}p{3.74cm}>{\raggedright\arraybackslash}p{3.74cm}}\hline\rowcolor{black!7} \textbf{Casilla} & \textbf{Cuenta} & \textbf{Resultado}\\\hline Fila 1, columna 1 & $(3)-(5)$ & $-2$\\ Fila 1, columna 2 & $(6)-(0)$ & $6$\\ Fila 2, columna 1 & $(6)-(10)$ & $-4$\\ Fila 2, columna 2 & $(3)-(5)$ & $-2$\\ Fila 3, columna 1 & $(9)-(15)$ & $-6$\\ Fila 3, columna 2 & $(12)-(10)$ & $2$\\ \hline\end{tabular}}\par\smallskip
\par\smallskip \textbf{Resultado de esta operación}\[\begin{bmatrix}-2 & 6\\-4 & -2\\-6 & 2\end{bmatrix}\]
\par\smallskip\textbf{Qué aprendimos.} El resultado tiene 3 filas y 2 columnas.
\end{minipage}\par\medskip
\noindent\begin{minipage}{\linewidth}\textbf{20. Transponemos: las filas se convierten en columnas}\\[2pt]No cambiamos los valores. Solo intercambiamos los dos índices de cada entrada. Una matriz de m filas y n columnas produce otra de n filas y m columnas.
\par\smallskip {\small\renewcommand{\arraystretch}{1.3}\begin{tabular}{>{\raggedright\arraybackslash}p{3.74cm}>{\raggedright\arraybackslash}p{3.74cm}>{\raggedright\arraybackslash}p{3.74cm}}\hline\rowcolor{black!7} \textbf{Posición antes} & \textbf{Posición después} & \textbf{Valor conservado}\\\hline (1, 1) & (1, 1) & $-2$\\ (1, 2) & (2, 1) & $6$\\ (2, 1) & (1, 2) & $-4$\\ (2, 2) & (2, 2) & $-2$\\ (3, 1) & (1, 3) & $-6$\\ (3, 2) & (2, 3) & $2$\\ \hline\end{tabular}}\par\smallskip
\par\smallskip \textbf{Resultado de esta operación}\[\begin{bmatrix}-2 & -4 & -6\\6 & -2 & 2\end{bmatrix}\]
\par\smallskip\textbf{Qué aprendimos.} El resultado tiene 2 filas y 3 columnas.
\end{minipage}\par\medskip
\noindent\begin{minipage}{\linewidth}\textbf{21. Resultado del inciso d}\\[2pt]Reunimos el resultado. Comprueba que cada casilla procede de la operación indicada, sin cambiar su posición.
\par\smallskip \textbf{Resultado}\[\begin{bmatrix}-2 & -4 & -6\\6 & -2 & 2\end{bmatrix}\]
\end{minipage}\par\medskip
\noindent\begin{minipage}{\linewidth}\textbf{22. Inciso e1: (-A)\textasciicircum{}T: qué vamos a calcular}\\[2pt]Identificamos la operación y comprobamos sus dimensiones antes de efectuar las cuentas. Las matrices originales permanecen escritas en el enunciado.
\[(-A)^T\]
\end{minipage}\par\medskip
\noindent\begin{minipage}{\linewidth}\textbf{23. Multiplicamos cada entrada por el escalar}\\[2pt]La fila y la columna nos dicen qué valores combinar. Hacemos una cuenta por casilla y después reunimos todas las casillas en la matriz resultado.
\par\smallskip {\small\renewcommand{\arraystretch}{1.3}\begin{tabular}{>{\raggedright\arraybackslash}p{3.74cm}>{\raggedright\arraybackslash}p{3.74cm}>{\raggedright\arraybackslash}p{3.74cm}}\hline\rowcolor{black!7} \textbf{Casilla} & \textbf{Cuenta} & \textbf{Resultado}\\\hline Fila 1, columna 1 & $-1\cdot(1)$ & $-1$\\ Fila 1, columna 2 & $-1\cdot(2)$ & $-2$\\ Fila 1, columna 3 & $-1\cdot(3)$ & $-3$\\ Fila 2, columna 1 & $-1\cdot(2)$ & $-2$\\ Fila 2, columna 2 & $-1\cdot(1)$ & $-1$\\ Fila 2, columna 3 & $-1\cdot(4)$ & $-4$\\ \hline\end{tabular}}\par\smallskip
\par\smallskip \textbf{Resultado de esta operación}\[\begin{bmatrix}-1 & -2 & -3\\-2 & -1 & -4\end{bmatrix}\]
\par\smallskip\textbf{Qué aprendimos.} El resultado tiene 2 filas y 3 columnas.
\end{minipage}\par\medskip
\noindent\begin{minipage}{\linewidth}\textbf{24. Transponemos: las filas se convierten en columnas}\\[2pt]No cambiamos los valores. Solo intercambiamos los dos índices de cada entrada. Una matriz de m filas y n columnas produce otra de n filas y m columnas.
\par\smallskip {\small\renewcommand{\arraystretch}{1.3}\begin{tabular}{>{\raggedright\arraybackslash}p{3.74cm}>{\raggedright\arraybackslash}p{3.74cm}>{\raggedright\arraybackslash}p{3.74cm}}\hline\rowcolor{black!7} \textbf{Posición antes} & \textbf{Posición después} & \textbf{Valor conservado}\\\hline (1, 1) & (1, 1) & $-1$\\ (1, 2) & (2, 1) & $-2$\\ (1, 3) & (3, 1) & $-3$\\ (2, 1) & (1, 2) & $-2$\\ (2, 2) & (2, 2) & $-1$\\ (2, 3) & (3, 2) & $-4$\\ \hline\end{tabular}}\par\smallskip
\par\smallskip \textbf{Resultado de esta operación}\[\begin{bmatrix}-1 & -2\\-2 & -1\\-3 & -4\end{bmatrix}\]
\par\smallskip\textbf{Qué aprendimos.} El resultado tiene 3 filas y 2 columnas.
\end{minipage}\par\medskip
\noindent\begin{minipage}{\linewidth}\textbf{25. Resultado del inciso e1}\\[2pt]Reunimos el resultado. Comprueba que cada casilla procede de la operación indicada, sin cambiar su posición.
\par\smallskip \textbf{Resultado}\[\begin{bmatrix}-1 & -2\\-2 & -1\\-3 & -4\end{bmatrix}\]
\end{minipage}\par\medskip
\noindent\begin{minipage}{\linewidth}\textbf{26. Inciso e2: -(A\textasciicircum{}T): qué vamos a calcular}\\[2pt]Identificamos la operación y comprobamos sus dimensiones antes de efectuar las cuentas. Las matrices originales permanecen escritas en el enunciado.
\[-(A^T)\]
\end{minipage}\par\medskip
\noindent\begin{minipage}{\linewidth}\textbf{27. Transponemos: las filas se convierten en columnas}\\[2pt]No cambiamos los valores. Solo intercambiamos los dos índices de cada entrada. Una matriz de m filas y n columnas produce otra de n filas y m columnas.
\par\smallskip {\small\renewcommand{\arraystretch}{1.3}\begin{tabular}{>{\raggedright\arraybackslash}p{3.74cm}>{\raggedright\arraybackslash}p{3.74cm}>{\raggedright\arraybackslash}p{3.74cm}}\hline\rowcolor{black!7} \textbf{Posición antes} & \textbf{Posición después} & \textbf{Valor conservado}\\\hline (1, 1) & (1, 1) & $1$\\ (1, 2) & (2, 1) & $2$\\ (1, 3) & (3, 1) & $3$\\ (2, 1) & (1, 2) & $2$\\ (2, 2) & (2, 2) & $1$\\ (2, 3) & (3, 2) & $4$\\ \hline\end{tabular}}\par\smallskip
\par\smallskip \textbf{Resultado de esta operación}\[\begin{bmatrix}1 & 2\\2 & 1\\3 & 4\end{bmatrix}\]
\par\smallskip\textbf{Qué aprendimos.} El resultado tiene 3 filas y 2 columnas.
\end{minipage}\par\medskip
\noindent\begin{minipage}{\linewidth}\textbf{28. Multiplicamos cada entrada por el escalar}\\[2pt]La fila y la columna nos dicen qué valores combinar. Hacemos una cuenta por casilla y después reunimos todas las casillas en la matriz resultado.
\par\smallskip {\small\renewcommand{\arraystretch}{1.3}\begin{tabular}{>{\raggedright\arraybackslash}p{3.74cm}>{\raggedright\arraybackslash}p{3.74cm}>{\raggedright\arraybackslash}p{3.74cm}}\hline\rowcolor{black!7} \textbf{Casilla} & \textbf{Cuenta} & \textbf{Resultado}\\\hline Fila 1, columna 1 & $-1\cdot(1)$ & $-1$\\ Fila 1, columna 2 & $-1\cdot(2)$ & $-2$\\ Fila 2, columna 1 & $-1\cdot(2)$ & $-2$\\ Fila 2, columna 2 & $-1\cdot(1)$ & $-1$\\ Fila 3, columna 1 & $-1\cdot(3)$ & $-3$\\ Fila 3, columna 2 & $-1\cdot(4)$ & $-4$\\ \hline\end{tabular}}\par\smallskip
\par\smallskip \textbf{Resultado de esta operación}\[\begin{bmatrix}-1 & -2\\-2 & -1\\-3 & -4\end{bmatrix}\]
\par\smallskip\textbf{Qué aprendimos.} El resultado tiene 3 filas y 2 columnas.
\end{minipage}\par\medskip
\noindent\begin{minipage}{\linewidth}\textbf{29. Resultado del inciso e2}\\[2pt]Reunimos el resultado. Comprueba que cada casilla procede de la operación indicada, sin cambiar su posición.
\par\smallskip \textbf{Resultado}\[\begin{bmatrix}-1 & -2\\-2 & -1\\-3 & -4\end{bmatrix}\]
\end{minipage}\par\medskip
\noindent\begin{minipage}{\linewidth}\textbf{30. Inciso f: (C+E+F\textasciicircum{}T)\textasciicircum{}T: qué vamos a calcular}\\[2pt]Identificamos la operación y comprobamos sus dimensiones antes de efectuar las cuentas. Las matrices originales permanecen escritas en el enunciado.
\[(C+E+F^T)^T\]
\end{minipage}\par\medskip
\noindent\begin{minipage}{\linewidth}\textbf{31. Operamos las entradas que ocupan la misma posición}\\[2pt]La fila y la columna nos dicen qué valores combinar. Hacemos una cuenta por casilla y después reunimos todas las casillas en la matriz resultado.
\par\smallskip {\small\renewcommand{\arraystretch}{1.3}\begin{tabular}{>{\raggedright\arraybackslash}p{3.74cm}>{\raggedright\arraybackslash}p{3.74cm}>{\raggedright\arraybackslash}p{3.74cm}}\hline\rowcolor{black!7} \textbf{Casilla} & \textbf{Cuenta} & \textbf{Resultado}\\\hline Fila 1, columna 1 & $(3)+(2)$ & $5$\\ Fila 1, columna 2 & $(-1)+(-4)$ & $-5$\\ Fila 1, columna 3 & $(3)+(5)$ & $8$\\ Fila 2, columna 1 & $(4)+(0)$ & $4$\\ Fila 2, columna 2 & $(1)+(1)$ & $2$\\ Fila 2, columna 3 & $(5)+(4)$ & $9$\\ Fila 3, columna 1 & $(2)+(3)$ & $5$\\ Fila 3, columna 2 & $(1)+(2)$ & $3$\\ Fila 3, columna 3 & $(3)+(1)$ & $4$\\ \hline\end{tabular}}\par\smallskip
\par\smallskip \textbf{Resultado de esta operación}\[\begin{bmatrix}5 & -5 & 8\\4 & 2 & 9\\5 & 3 & 4\end{bmatrix}\]
\par\smallskip\textbf{Qué aprendimos.} El resultado tiene 3 filas y 3 columnas.
\end{minipage}\par\medskip
\noindent\begin{minipage}{\linewidth}\textbf{32. Transponemos: las filas se convierten en columnas}\\[2pt]No cambiamos los valores. Solo intercambiamos los dos índices de cada entrada. Una matriz de m filas y n columnas produce otra de n filas y m columnas.
\par\smallskip {\small\renewcommand{\arraystretch}{1.3}\begin{tabular}{>{\raggedright\arraybackslash}p{3.74cm}>{\raggedright\arraybackslash}p{3.74cm}>{\raggedright\arraybackslash}p{3.74cm}}\hline\rowcolor{black!7} \textbf{Posición antes} & \textbf{Posición después} & \textbf{Valor conservado}\\\hline (1, 1) & (1, 1) & $-4$\\ (1, 2) & (2, 1) & $5$\\ (2, 1) & (1, 2) & $2$\\ (2, 2) & (2, 2) & $3$\\ \hline\end{tabular}}\par\smallskip
\par\smallskip \textbf{Resultado de esta operación}\[\begin{bmatrix}-4 & 2\\5 & 3\end{bmatrix}\]
\par\smallskip\textbf{Qué aprendimos.} El resultado tiene 2 filas y 2 columnas.
\end{minipage}\par\medskip
\noindent\begin{minipage}{\linewidth}\textbf{33. Este inciso no está definido}\\[2pt]El primer operando tiene tamaño 3 por 3 y el segundo 2 por 2. La suma o resta no está definida. No se agregan ceros ni se recortan filas para forzarla.
\par\smallskip\textbf{Qué aprendimos.} La comprobación de dimensiones es parte de la solución, no un error del programa.
\end{minipage}\par\medskip
\respuesta{a: (2A)\textasciicircum{}T: [[2, 4], [4, 2], [6, 8]]; b: (A-B)\textasciicircum{}T: no definida; c: (3B\textasciicircum{}T-2A)\textasciicircum{}T: [[1, -4], [2, 1], [3, -2]]; d: (3A\textasciicircum{}T-5B)\textasciicircum{}T: [[-2, -4, -6], [6, -2, 2]]; e1: (-A)\textasciicircum{}T: [[-1, -2], [-2, -1], [-3, -4]]; e2: -(A\textasciicircum{}T): [[-1, -2], [-2, -1], [-3, -4]]; f: (C+E+F\textasciicircum{}T)\textasciicircum{}T: no definida.}
\noindent\begin{minipage}{\linewidth}\centering\includegraphics[width=.96\linewidth,height=6.8cm,keepaspectratio]{../recursos/grafico_1_2-7.png}\\{\small Lectura visual complementaria; la demostración está en las operaciones.}\end{minipage}
\end{multicols*}
\newpage\fancyfoot[L]{\small Sección 1.2 | Ejercicio 8 | Paso a paso}\hypertarget{1.2-8}{}
\pdfbookmark[1]{1.2 / 8}{bm1.2-8}
\titulo{8 | ¿Es una combinación lineal?}
\textbf{Enunciado}\par Determine si la matriz indicada es combinación lineal de I y J. Justifique su respuesta.
\par\smallskip \begin{minipage}[t]{.47\linewidth}\centering \textbf{Objetivo}\[\begin{bmatrix}3 & 0\\0 & 2\end{bmatrix}\]\end{minipage}\hfill\begin{minipage}[t]{.47\linewidth}\centering \textbf{I}\[\begin{bmatrix}1 & 0\\0 & 1\end{bmatrix}\]\end{minipage}\par\smallskip
\par\smallskip \begin{minipage}[t]{.47\linewidth}\centering \textbf{J}\[\begin{bmatrix}1 & 0\\0 & 0\end{bmatrix}\]\end{minipage}\par\smallskip
\begin{multicols*}{2}
\noindent\begin{minipage}{\linewidth}\textbf{1. Qué significa combinación lineal}\\[2pt]Debemos encontrar dos números reales $\alpha$ y $\beta$ que multipliquen a I y J y cuya suma produzca exactamente la matriz objetivo.
\par\smallskip \textbf{$\alpha$I + $\beta$J}\[\begin{bmatrix}\alpha + \beta & 0\\0 & \alpha\end{bmatrix}\]
\end{minipage}\par\medskip
\noindent\begin{minipage}{\linewidth}\textbf{2. Comparamos las casillas}\\[2pt]En todas las combinaciones de I y J las dos entradas fuera de la diagonal son cero. Esa condición es obligatoria, sin importar qué valores elijamos para $\alpha$ y $\beta$.
\par\smallskip\textbf{Qué aprendimos.} La casilla superior derecha de la matriz objetivo es 0; no impide la combinación.
\end{minipage}\par\medskip
\noindent\begin{minipage}{\linewidth}\textbf{3. Igualamos las dos entradas diagonales}\\[2pt]La casilla inferior derecha exige $\alpha$ = 2. La superior izquierda exige $\alpha$ + $\beta$ = 3. Sustituyendo $\alpha$, obtenemos $\beta$ = 1.
\[\alpha=2,\quad 2+\beta=3,\quad\beta=1\]
\end{minipage}\par\medskip
\noindent\begin{minipage}{\linewidth}\textbf{4. Verificación}\\[2pt]Multiplicamos y sumamos por casillas.
\par\smallskip \begin{minipage}[t]{.47\linewidth}\centering \textbf{2I + J}\[\begin{bmatrix}3 & 0\\0 & 2\end{bmatrix}\]\end{minipage}\hfill\begin{minipage}[t]{.47\linewidth}\centering \textbf{Objetivo}\[\begin{bmatrix}3 & 0\\0 & 2\end{bmatrix}\]\end{minipage}\par\smallskip
\par\smallskip\textbf{Qué aprendimos.} Las cuatro casillas coinciden: sí es combinación lineal.
\end{minipage}\par\medskip
\respuesta{Sí: 2I + J.}
\noindent\begin{minipage}{\linewidth}\centering\includegraphics[width=.96\linewidth,height=6.8cm,keepaspectratio]{../recursos/grafico_1_2-8.png}\\{\small Lectura visual complementaria; la demostración está en las operaciones.}\end{minipage}
\end{multicols*}
\newpage\fancyfoot[L]{\small Sección 1.2 | Ejercicio 9 | Paso a paso}\hypertarget{1.2-9}{}
\pdfbookmark[1]{1.2 / 9}{bm1.2-9}
\titulo{9 | ¿Es una combinación lineal?}
\textbf{Enunciado}\par Determine si la matriz indicada es combinación lineal de I y J. Justifique su respuesta.
\par\smallskip \begin{minipage}[t]{.47\linewidth}\centering \textbf{Objetivo}\[\begin{bmatrix}4 & 1\\0 & -3\end{bmatrix}\]\end{minipage}\hfill\begin{minipage}[t]{.47\linewidth}\centering \textbf{I}\[\begin{bmatrix}1 & 0\\0 & 1\end{bmatrix}\]\end{minipage}\par\smallskip
\par\smallskip \begin{minipage}[t]{.47\linewidth}\centering \textbf{J}\[\begin{bmatrix}1 & 0\\0 & 0\end{bmatrix}\]\end{minipage}\par\smallskip
\begin{multicols*}{2}
\noindent\begin{minipage}{\linewidth}\textbf{1. Qué significa combinación lineal}\\[2pt]Debemos encontrar dos números reales $\alpha$ y $\beta$ que multipliquen a I y J y cuya suma produzca exactamente la matriz objetivo.
\par\smallskip \textbf{$\alpha$I + $\beta$J}\[\begin{bmatrix}\alpha + \beta & 0\\0 & \alpha\end{bmatrix}\]
\end{minipage}\par\medskip
\noindent\begin{minipage}{\linewidth}\textbf{2. Comparamos las casillas}\\[2pt]En todas las combinaciones de I y J las dos entradas fuera de la diagonal son cero. Esa condición es obligatoria, sin importar qué valores elijamos para $\alpha$ y $\beta$.
\par\smallskip\textbf{Qué aprendimos.} La casilla superior derecha de la matriz objetivo es 1. Como la combinación siempre tiene 0 allí, es imposible.
\end{minipage}\par\medskip
\noindent\begin{minipage}{\linewidth}\textbf{3. Conclusión justificada}\\[2pt]No basta intentar varios coeficientes. La entrada (1, 2) demuestra que no puede existir ninguno: tendría que cumplirse 0 = 1.
\[0\ne1\]
\end{minipage}\par\medskip
\respuesta{No: toda combinación tiene entrada (1, 2) igual a 0 y la matriz objetivo tiene 1.}
\noindent\begin{minipage}{\linewidth}\centering\includegraphics[width=.96\linewidth,height=6.8cm,keepaspectratio]{../recursos/grafico_1_2-9.png}\\{\small Lectura visual complementaria; la demostración está en las operaciones.}\end{minipage}
\end{multicols*}
\newpage\fancyfoot[L]{\small Sección 1.2 | Ejercicio 10 | Paso a paso}\hypertarget{1.2-10}{}
\pdfbookmark[1]{1.2 / 10}{bm1.2-10}
\titulo{10 | Una matriz que depende de $\lambda$}
\textbf{Enunciado}\par Sean A e I3 las matrices indicadas. Si $\lambda$ es un número real, calcule $\lambda$I3 - A.
\par\smallskip \begin{minipage}[t]{.47\linewidth}\centering \textbf{A}\[\begin{bmatrix}1 & 2 & 3\\6 & -2 & 3\\5 & 2 & 4\end{bmatrix}\]\end{minipage}\hfill\begin{minipage}[t]{.47\linewidth}\centering \textbf{I3}\[\begin{bmatrix}1 & 0 & 0\\0 & 1 & 0\\0 & 0 & 1\end{bmatrix}\]\end{minipage}\par\smallskip
\begin{multicols*}{2}
\noindent\begin{minipage}{\linewidth}\textbf{1. Primero multiplicamos la identidad por $\lambda$}\\[2pt]La identidad tiene unos en la diagonal y ceros fuera. Multiplicar por $\lambda$ coloca $\lambda$ en la diagonal y mantiene los ceros restantes.
\par\smallskip \textbf{$\lambda$I3}\[\begin{bmatrix}\lambda & 0 & 0\\0 & \lambda & 0\\0 & 0 & \lambda\end{bmatrix}\]
\end{minipage}\par\medskip
\noindent\begin{minipage}{\linewidth}\textbf{2. Restamos entrada por entrada}\\[2pt]En la diagonal queda $\lambda$ menos la entrada de A. Fuera de la diagonal, calculamos 0 menos la entrada de A; por eso cambia su signo.
\par\smallskip \textbf{$\lambda$I3 - A}\[\begin{bmatrix}\lambda - 1 & -2 & -3\\-6 & \lambda + 2 & -3\\-5 & -2 & \lambda - 4\end{bmatrix}\]
\par\smallskip\textbf{Qué aprendimos.} La casilla central es $\lambda$ - (-2) = $\lambda$ + 2; restar un negativo equivale a sumar.
\end{minipage}\par\medskip
\respuesta{[[$\lambda$-1, -2, -3], [-6, $\lambda$+2, -3], [-5, -2, $\lambda$-4]].}
\noindent\begin{minipage}{\linewidth}\centering\includegraphics[width=.96\linewidth,height=6.8cm,keepaspectratio]{../recursos/grafico_1_2-10.png}\\{\small Lectura visual complementaria; la demostración está en las operaciones.}\end{minipage}
\end{multicols*}
\newpage\fancyfoot[L]{\small Sección 1.2 | Ejercicio 11 | Paso a paso}\hypertarget{1.2-11}{}
\pdfbookmark[1]{1.2 / 11}{bm1.2-11}
\titulo{11 | Operaciones con matrices binarias}
\textbf{Enunciado}\par De ser posible, calcule cada expresión indicada. Trabaje con matrices binarias: las sumas y diferencias se calculan módulo 2.
\par\smallskip \begin{minipage}[t]{.47\linewidth}\centering \textbf{A}\[\begin{bmatrix}1 & 0 & 1\\1 & 1 & 0\\0 & 1 & 1\end{bmatrix}\]\end{minipage}\hfill\begin{minipage}[t]{.47\linewidth}\centering \textbf{B}\[\begin{bmatrix}0 & 1 & 1\\1 & 0 & 1\\1 & 1 & 0\end{bmatrix}\]\end{minipage}\par\smallskip
\par\smallskip \begin{minipage}[t]{.47\linewidth}\centering \textbf{C}\[\begin{bmatrix}1 & 1 & 0\\0 & 1 & 1\\1 & 0 & 1\end{bmatrix}\]\end{minipage}\par\smallskip
\textbf{Expresiones del enunciado}\par
a. \(A+B\)\par
b. \(B+C\)\par
c. \(A+B+C\)\par
d. \(A+C^T\)\par
e. \(B-C\)\par
\begin{multicols*}{2}
\noindent\begin{minipage}{\linewidth}\textbf{1. Antes de calcular: filas, columnas y operaciones}\\[2pt]El tamaño de una matriz se escribe filas por columnas. Para sumar o restar, los dos tamaños deben coincidir exactamente. Multiplicar por un número afecta a cada entrada. Transponer intercambia las posiciones: la entrada de fila i y columna j pasa a la fila j y columna i.
\[(M^T)_{ij}=m_{ji}\]
\par\smallskip\textbf{Qué aprendimos.} Los paréntesis indican qué operación hay que realizar primero. Una suma de tamaños distintos no se vuelve válida por transponerla después.
\end{minipage}\par\medskip
\noindent\begin{minipage}{\linewidth}\textbf{2. Regla de la aritmética binaria}\\[2pt]Aquí los resultados solo pueden ser 0 o 1. Calculamos el resto al dividir la cuenta entre 2; por eso 1 + 1 se convierte en 0. En esta aritmética, restar coincide con sumar: el opuesto de 1 es 1.
\[0+0=0,\quad 0+1=1,\quad1+0=1,\quad1+1=0\quad(\mathrm{mod}\ 2)\]
\par\smallskip\textbf{Qué aprendimos.} Estas reglas se aplican en los ejercicios 11–15 y en los teóricos binarios, no en las operaciones reales de los ejercicios 4–10.
\end{minipage}\par\medskip
\noindent\begin{minipage}{\linewidth}\textbf{3. Inciso a: A+B: qué vamos a calcular}\\[2pt]Identificamos la operación y comprobamos sus dimensiones antes de efectuar las cuentas. Las matrices originales permanecen escritas en el enunciado.
\[A+B\]
\end{minipage}\par\medskip
\noindent\begin{minipage}{\linewidth}\textbf{4. Operamos las entradas que ocupan la misma posición}\\[2pt]La fila y la columna nos dicen qué valores combinar. Hacemos una cuenta por casilla y después reunimos todas las casillas en la matriz resultado. En cada casilla reducimos el resultado módulo 2: un número par da 0 y un número impar da 1.
\par\smallskip {\small\renewcommand{\arraystretch}{1.3}\begin{tabular}{>{\raggedright\arraybackslash}p{3.74cm}>{\raggedright\arraybackslash}p{3.74cm}>{\raggedright\arraybackslash}p{3.74cm}}\hline\rowcolor{black!7} \textbf{Casilla} & \textbf{Cuenta} & \textbf{Resultado}\\\hline Fila 1, columna 1 & $\mathrm{mod}((1)+(0),2)$ & $1$\\ Fila 1, columna 2 & $\mathrm{mod}((0)+(1),2)$ & $1$\\ Fila 1, columna 3 & $\mathrm{mod}((1)+(1),2)$ & $0$\\ Fila 2, columna 1 & $\mathrm{mod}((1)+(1),2)$ & $0$\\ Fila 2, columna 2 & $\mathrm{mod}((1)+(0),2)$ & $1$\\ Fila 2, columna 3 & $\mathrm{mod}((0)+(1),2)$ & $1$\\ Fila 3, columna 1 & $\mathrm{mod}((0)+(1),2)$ & $1$\\ Fila 3, columna 2 & $\mathrm{mod}((1)+(1),2)$ & $0$\\ Fila 3, columna 3 & $\mathrm{mod}((1)+(0),2)$ & $1$\\ \hline\end{tabular}}\par\smallskip
\par\smallskip \textbf{Resultado de esta operación}\[\begin{bmatrix}1 & 1 & 0\\0 & 1 & 1\\1 & 0 & 1\end{bmatrix}\]
\par\smallskip\textbf{Qué aprendimos.} El resultado tiene 3 filas y 3 columnas.
\end{minipage}\par\medskip
\noindent\begin{minipage}{\linewidth}\textbf{5. Resultado del inciso a}\\[2pt]Reunimos el resultado. Comprueba que cada casilla procede de la operación indicada, sin cambiar su posición.
\par\smallskip \textbf{Resultado}\[\begin{bmatrix}1 & 1 & 0\\0 & 1 & 1\\1 & 0 & 1\end{bmatrix}\]
\end{minipage}\par\medskip
\noindent\begin{minipage}{\linewidth}\textbf{6. Inciso b: B+C: qué vamos a calcular}\\[2pt]Identificamos la operación y comprobamos sus dimensiones antes de efectuar las cuentas. Las matrices originales permanecen escritas en el enunciado.
\[B+C\]
\end{minipage}\par\medskip
\noindent\begin{minipage}{\linewidth}\textbf{7. Operamos las entradas que ocupan la misma posición}\\[2pt]La fila y la columna nos dicen qué valores combinar. Hacemos una cuenta por casilla y después reunimos todas las casillas en la matriz resultado. En cada casilla reducimos el resultado módulo 2: un número par da 0 y un número impar da 1.
\par\smallskip {\small\renewcommand{\arraystretch}{1.3}\begin{tabular}{>{\raggedright\arraybackslash}p{3.74cm}>{\raggedright\arraybackslash}p{3.74cm}>{\raggedright\arraybackslash}p{3.74cm}}\hline\rowcolor{black!7} \textbf{Casilla} & \textbf{Cuenta} & \textbf{Resultado}\\\hline Fila 1, columna 1 & $\mathrm{mod}((0)+(1),2)$ & $1$\\ Fila 1, columna 2 & $\mathrm{mod}((1)+(1),2)$ & $0$\\ Fila 1, columna 3 & $\mathrm{mod}((1)+(0),2)$ & $1$\\ Fila 2, columna 1 & $\mathrm{mod}((1)+(0),2)$ & $1$\\ Fila 2, columna 2 & $\mathrm{mod}((0)+(1),2)$ & $1$\\ Fila 2, columna 3 & $\mathrm{mod}((1)+(1),2)$ & $0$\\ Fila 3, columna 1 & $\mathrm{mod}((1)+(1),2)$ & $0$\\ Fila 3, columna 2 & $\mathrm{mod}((1)+(0),2)$ & $1$\\ Fila 3, columna 3 & $\mathrm{mod}((0)+(1),2)$ & $1$\\ \hline\end{tabular}}\par\smallskip
\par\smallskip \textbf{Resultado de esta operación}\[\begin{bmatrix}1 & 0 & 1\\1 & 1 & 0\\0 & 1 & 1\end{bmatrix}\]
\par\smallskip\textbf{Qué aprendimos.} El resultado tiene 3 filas y 3 columnas.
\end{minipage}\par\medskip
\noindent\begin{minipage}{\linewidth}\textbf{8. Resultado del inciso b}\\[2pt]Reunimos el resultado. Comprueba que cada casilla procede de la operación indicada, sin cambiar su posición.
\par\smallskip \textbf{Resultado}\[\begin{bmatrix}1 & 0 & 1\\1 & 1 & 0\\0 & 1 & 1\end{bmatrix}\]
\end{minipage}\par\medskip
\noindent\begin{minipage}{\linewidth}\textbf{9. Inciso c: A+B+C: qué vamos a calcular}\\[2pt]Identificamos la operación y comprobamos sus dimensiones antes de efectuar las cuentas. Las matrices originales permanecen escritas en el enunciado.
\[A+B+C\]
\end{minipage}\par\medskip
\noindent\begin{minipage}{\linewidth}\textbf{10. Operamos las entradas que ocupan la misma posición}\\[2pt]La fila y la columna nos dicen qué valores combinar. Hacemos una cuenta por casilla y después reunimos todas las casillas en la matriz resultado. En cada casilla reducimos el resultado módulo 2: un número par da 0 y un número impar da 1.
\par\smallskip {\small\renewcommand{\arraystretch}{1.3}\begin{tabular}{>{\raggedright\arraybackslash}p{3.74cm}>{\raggedright\arraybackslash}p{3.74cm}>{\raggedright\arraybackslash}p{3.74cm}}\hline\rowcolor{black!7} \textbf{Casilla} & \textbf{Cuenta} & \textbf{Resultado}\\\hline Fila 1, columna 1 & $\mathrm{mod}((1)+(0),2)$ & $1$\\ Fila 1, columna 2 & $\mathrm{mod}((0)+(1),2)$ & $1$\\ Fila 1, columna 3 & $\mathrm{mod}((1)+(1),2)$ & $0$\\ Fila 2, columna 1 & $\mathrm{mod}((1)+(1),2)$ & $0$\\ Fila 2, columna 2 & $\mathrm{mod}((1)+(0),2)$ & $1$\\ Fila 2, columna 3 & $\mathrm{mod}((0)+(1),2)$ & $1$\\ Fila 3, columna 1 & $\mathrm{mod}((0)+(1),2)$ & $1$\\ Fila 3, columna 2 & $\mathrm{mod}((1)+(1),2)$ & $0$\\ Fila 3, columna 3 & $\mathrm{mod}((1)+(0),2)$ & $1$\\ \hline\end{tabular}}\par\smallskip
\par\smallskip \textbf{Resultado de esta operación}\[\begin{bmatrix}1 & 1 & 0\\0 & 1 & 1\\1 & 0 & 1\end{bmatrix}\]
\par\smallskip\textbf{Qué aprendimos.} El resultado tiene 3 filas y 3 columnas.
\end{minipage}\par\medskip
\noindent\begin{minipage}{\linewidth}\textbf{11. Operamos las entradas que ocupan la misma posición}\\[2pt]La fila y la columna nos dicen qué valores combinar. Hacemos una cuenta por casilla y después reunimos todas las casillas en la matriz resultado. En cada casilla reducimos el resultado módulo 2: un número par da 0 y un número impar da 1.
\par\smallskip {\small\renewcommand{\arraystretch}{1.3}\begin{tabular}{>{\raggedright\arraybackslash}p{3.74cm}>{\raggedright\arraybackslash}p{3.74cm}>{\raggedright\arraybackslash}p{3.74cm}}\hline\rowcolor{black!7} \textbf{Casilla} & \textbf{Cuenta} & \textbf{Resultado}\\\hline Fila 1, columna 1 & $\mathrm{mod}((1)+(1),2)$ & $0$\\ Fila 1, columna 2 & $\mathrm{mod}((1)+(1),2)$ & $0$\\ Fila 1, columna 3 & $\mathrm{mod}((0)+(0),2)$ & $0$\\ Fila 2, columna 1 & $\mathrm{mod}((0)+(0),2)$ & $0$\\ Fila 2, columna 2 & $\mathrm{mod}((1)+(1),2)$ & $0$\\ Fila 2, columna 3 & $\mathrm{mod}((1)+(1),2)$ & $0$\\ Fila 3, columna 1 & $\mathrm{mod}((1)+(1),2)$ & $0$\\ Fila 3, columna 2 & $\mathrm{mod}((0)+(0),2)$ & $0$\\ Fila 3, columna 3 & $\mathrm{mod}((1)+(1),2)$ & $0$\\ \hline\end{tabular}}\par\smallskip
\par\smallskip \textbf{Resultado de esta operación}\[\begin{bmatrix}0 & 0 & 0\\0 & 0 & 0\\0 & 0 & 0\end{bmatrix}\]
\par\smallskip\textbf{Qué aprendimos.} El resultado tiene 3 filas y 3 columnas.
\end{minipage}\par\medskip
\noindent\begin{minipage}{\linewidth}\textbf{12. Resultado del inciso c}\\[2pt]Reunimos el resultado. Comprueba que cada casilla procede de la operación indicada, sin cambiar su posición.
\par\smallskip \textbf{Resultado}\[\begin{bmatrix}0 & 0 & 0\\0 & 0 & 0\\0 & 0 & 0\end{bmatrix}\]
\end{minipage}\par\medskip
\noindent\begin{minipage}{\linewidth}\textbf{13. Inciso d: A+C\textasciicircum{}T: qué vamos a calcular}\\[2pt]Identificamos la operación y comprobamos sus dimensiones antes de efectuar las cuentas. Las matrices originales permanecen escritas en el enunciado.
\[A+C^T\]
\end{minipage}\par\medskip
\noindent\begin{minipage}{\linewidth}\textbf{14. Transponemos: las filas se convierten en columnas}\\[2pt]No cambiamos los valores. Solo intercambiamos los dos índices de cada entrada. Una matriz de m filas y n columnas produce otra de n filas y m columnas.
\par\smallskip {\small\renewcommand{\arraystretch}{1.3}\begin{tabular}{>{\raggedright\arraybackslash}p{3.74cm}>{\raggedright\arraybackslash}p{3.74cm}>{\raggedright\arraybackslash}p{3.74cm}}\hline\rowcolor{black!7} \textbf{Posición antes} & \textbf{Posición después} & \textbf{Valor conservado}\\\hline (1, 1) & (1, 1) & $1$\\ (1, 2) & (2, 1) & $1$\\ (1, 3) & (3, 1) & $0$\\ (2, 1) & (1, 2) & $0$\\ (2, 2) & (2, 2) & $1$\\ (2, 3) & (3, 2) & $1$\\ (3, 1) & (1, 3) & $1$\\ (3, 2) & (2, 3) & $0$\\ (3, 3) & (3, 3) & $1$\\ \hline\end{tabular}}\par\smallskip
\par\smallskip \textbf{Resultado de esta operación}\[\begin{bmatrix}1 & 0 & 1\\1 & 1 & 0\\0 & 1 & 1\end{bmatrix}\]
\par\smallskip\textbf{Qué aprendimos.} El resultado tiene 3 filas y 3 columnas.
\end{minipage}\par\medskip
\noindent\begin{minipage}{\linewidth}\textbf{15. Operamos las entradas que ocupan la misma posición}\\[2pt]La fila y la columna nos dicen qué valores combinar. Hacemos una cuenta por casilla y después reunimos todas las casillas en la matriz resultado. En cada casilla reducimos el resultado módulo 2: un número par da 0 y un número impar da 1.
\par\smallskip {\small\renewcommand{\arraystretch}{1.3}\begin{tabular}{>{\raggedright\arraybackslash}p{3.74cm}>{\raggedright\arraybackslash}p{3.74cm}>{\raggedright\arraybackslash}p{3.74cm}}\hline\rowcolor{black!7} \textbf{Casilla} & \textbf{Cuenta} & \textbf{Resultado}\\\hline Fila 1, columna 1 & $\mathrm{mod}((1)+(1),2)$ & $0$\\ Fila 1, columna 2 & $\mathrm{mod}((0)+(0),2)$ & $0$\\ Fila 1, columna 3 & $\mathrm{mod}((1)+(1),2)$ & $0$\\ Fila 2, columna 1 & $\mathrm{mod}((1)+(1),2)$ & $0$\\ Fila 2, columna 2 & $\mathrm{mod}((1)+(1),2)$ & $0$\\ Fila 2, columna 3 & $\mathrm{mod}((0)+(0),2)$ & $0$\\ Fila 3, columna 1 & $\mathrm{mod}((0)+(0),2)$ & $0$\\ Fila 3, columna 2 & $\mathrm{mod}((1)+(1),2)$ & $0$\\ Fila 3, columna 3 & $\mathrm{mod}((1)+(1),2)$ & $0$\\ \hline\end{tabular}}\par\smallskip
\par\smallskip \textbf{Resultado de esta operación}\[\begin{bmatrix}0 & 0 & 0\\0 & 0 & 0\\0 & 0 & 0\end{bmatrix}\]
\par\smallskip\textbf{Qué aprendimos.} El resultado tiene 3 filas y 3 columnas.
\end{minipage}\par\medskip
\noindent\begin{minipage}{\linewidth}\textbf{16. Resultado del inciso d}\\[2pt]Reunimos el resultado. Comprueba que cada casilla procede de la operación indicada, sin cambiar su posición.
\par\smallskip \textbf{Resultado}\[\begin{bmatrix}0 & 0 & 0\\0 & 0 & 0\\0 & 0 & 0\end{bmatrix}\]
\end{minipage}\par\medskip
\noindent\begin{minipage}{\linewidth}\textbf{17. Inciso e: B-C: qué vamos a calcular}\\[2pt]Identificamos la operación y comprobamos sus dimensiones antes de efectuar las cuentas. Las matrices originales permanecen escritas en el enunciado.
\[B-C\]
\end{minipage}\par\medskip
\noindent\begin{minipage}{\linewidth}\textbf{18. Operamos las entradas que ocupan la misma posición}\\[2pt]La fila y la columna nos dicen qué valores combinar. Hacemos una cuenta por casilla y después reunimos todas las casillas en la matriz resultado. En cada casilla reducimos el resultado módulo 2: un número par da 0 y un número impar da 1.
\par\smallskip {\small\renewcommand{\arraystretch}{1.3}\begin{tabular}{>{\raggedright\arraybackslash}p{3.74cm}>{\raggedright\arraybackslash}p{3.74cm}>{\raggedright\arraybackslash}p{3.74cm}}\hline\rowcolor{black!7} \textbf{Casilla} & \textbf{Cuenta} & \textbf{Resultado}\\\hline Fila 1, columna 1 & $\mathrm{mod}((0)-(1),2)$ & $1$\\ Fila 1, columna 2 & $\mathrm{mod}((1)-(1),2)$ & $0$\\ Fila 1, columna 3 & $\mathrm{mod}((1)-(0),2)$ & $1$\\ Fila 2, columna 1 & $\mathrm{mod}((1)-(0),2)$ & $1$\\ Fila 2, columna 2 & $\mathrm{mod}((0)-(1),2)$ & $1$\\ Fila 2, columna 3 & $\mathrm{mod}((1)-(1),2)$ & $0$\\ Fila 3, columna 1 & $\mathrm{mod}((1)-(1),2)$ & $0$\\ Fila 3, columna 2 & $\mathrm{mod}((1)-(0),2)$ & $1$\\ Fila 3, columna 3 & $\mathrm{mod}((0)-(1),2)$ & $1$\\ \hline\end{tabular}}\par\smallskip
\par\smallskip \textbf{Resultado de esta operación}\[\begin{bmatrix}1 & 0 & 1\\1 & 1 & 0\\0 & 1 & 1\end{bmatrix}\]
\par\smallskip\textbf{Qué aprendimos.} El resultado tiene 3 filas y 3 columnas.
\end{minipage}\par\medskip
\noindent\begin{minipage}{\linewidth}\textbf{19. Resultado del inciso e}\\[2pt]Reunimos el resultado. Comprueba que cada casilla procede de la operación indicada, sin cambiar su posición.
\par\smallskip \textbf{Resultado}\[\begin{bmatrix}1 & 0 & 1\\1 & 1 & 0\\0 & 1 & 1\end{bmatrix}\]
\end{minipage}\par\medskip
\respuesta{a: A+B: [[1, 1, 0], [0, 1, 1], [1, 0, 1]]; b: B+C: [[1, 0, 1], [1, 1, 0], [0, 1, 1]]; c: A+B+C: [[0, 0, 0], [0, 0, 0], [0, 0, 0]]; d: A+C\textasciicircum{}T: [[0, 0, 0], [0, 0, 0], [0, 0, 0]]; e: B-C: [[1, 0, 1], [1, 1, 0], [0, 1, 1]].}
\noindent\begin{minipage}{\linewidth}\centering\includegraphics[width=.96\linewidth,height=6.8cm,keepaspectratio]{../recursos/grafico_1_2-11.png}\\{\small Lectura visual complementaria; la demostración está en las operaciones.}\end{minipage}
\end{multicols*}
\newpage\fancyfoot[L]{\small Sección 1.2 | Ejercicio 12 | Paso a paso}\hypertarget{1.2-12}{}
\pdfbookmark[1]{1.2 / 12}{bm1.2-12}
\titulo{12 | Operaciones con matrices binarias}
\textbf{Enunciado}\par De ser posible, calcule cada expresión indicada. Trabaje con matrices binarias: las sumas y diferencias se calculan módulo 2.
\par\smallskip \begin{minipage}[t]{.47\linewidth}\centering \textbf{A}\[\begin{bmatrix}1 & 0\\1 & 0\end{bmatrix}\]\end{minipage}\hfill\begin{minipage}[t]{.47\linewidth}\centering \textbf{B}\[\begin{bmatrix}1 & 0\\0 & 1\end{bmatrix}\]\end{minipage}\par\smallskip
\par\smallskip \begin{minipage}[t]{.47\linewidth}\centering \textbf{C}\[\begin{bmatrix}1 & 1\\0 & 0\end{bmatrix}\]\end{minipage}\hfill\begin{minipage}[t]{.47\linewidth}\centering \textbf{D}\[\begin{bmatrix}0 & 0\\1 & 0\end{bmatrix}\]\end{minipage}\par\smallskip
\textbf{Expresiones del enunciado}\par
a. \(A+B\)\par
b. \(C+D\)\par
c. \(A+B+(C+D)^T\)\par
d. \(C-B\)\par
e. \(A-B+C-D\)\par
\begin{multicols*}{2}
\noindent\begin{minipage}{\linewidth}\textbf{1. Antes de calcular: filas, columnas y operaciones}\\[2pt]El tamaño de una matriz se escribe filas por columnas. Para sumar o restar, los dos tamaños deben coincidir exactamente. Multiplicar por un número afecta a cada entrada. Transponer intercambia las posiciones: la entrada de fila i y columna j pasa a la fila j y columna i.
\[(M^T)_{ij}=m_{ji}\]
\par\smallskip\textbf{Qué aprendimos.} Los paréntesis indican qué operación hay que realizar primero. Una suma de tamaños distintos no se vuelve válida por transponerla después.
\end{minipage}\par\medskip
\noindent\begin{minipage}{\linewidth}\textbf{2. Regla de la aritmética binaria}\\[2pt]Aquí los resultados solo pueden ser 0 o 1. Calculamos el resto al dividir la cuenta entre 2; por eso 1 + 1 se convierte en 0. En esta aritmética, restar coincide con sumar: el opuesto de 1 es 1.
\[0+0=0,\quad 0+1=1,\quad1+0=1,\quad1+1=0\quad(\mathrm{mod}\ 2)\]
\par\smallskip\textbf{Qué aprendimos.} Estas reglas se aplican en los ejercicios 11–15 y en los teóricos binarios, no en las operaciones reales de los ejercicios 4–10.
\end{minipage}\par\medskip
\noindent\begin{minipage}{\linewidth}\textbf{3. Inciso a: A+B: qué vamos a calcular}\\[2pt]Identificamos la operación y comprobamos sus dimensiones antes de efectuar las cuentas. Las matrices originales permanecen escritas en el enunciado.
\[A+B\]
\end{minipage}\par\medskip
\noindent\begin{minipage}{\linewidth}\textbf{4. Operamos las entradas que ocupan la misma posición}\\[2pt]La fila y la columna nos dicen qué valores combinar. Hacemos una cuenta por casilla y después reunimos todas las casillas en la matriz resultado. En cada casilla reducimos el resultado módulo 2: un número par da 0 y un número impar da 1.
\par\smallskip {\small\renewcommand{\arraystretch}{1.3}\begin{tabular}{>{\raggedright\arraybackslash}p{3.74cm}>{\raggedright\arraybackslash}p{3.74cm}>{\raggedright\arraybackslash}p{3.74cm}}\hline\rowcolor{black!7} \textbf{Casilla} & \textbf{Cuenta} & \textbf{Resultado}\\\hline Fila 1, columna 1 & $\mathrm{mod}((1)+(1),2)$ & $0$\\ Fila 1, columna 2 & $\mathrm{mod}((0)+(0),2)$ & $0$\\ Fila 2, columna 1 & $\mathrm{mod}((1)+(0),2)$ & $1$\\ Fila 2, columna 2 & $\mathrm{mod}((0)+(1),2)$ & $1$\\ \hline\end{tabular}}\par\smallskip
\par\smallskip \textbf{Resultado de esta operación}\[\begin{bmatrix}0 & 0\\1 & 1\end{bmatrix}\]
\par\smallskip\textbf{Qué aprendimos.} El resultado tiene 2 filas y 2 columnas.
\end{minipage}\par\medskip
\noindent\begin{minipage}{\linewidth}\textbf{5. Resultado del inciso a}\\[2pt]Reunimos el resultado. Comprueba que cada casilla procede de la operación indicada, sin cambiar su posición.
\par\smallskip \textbf{Resultado}\[\begin{bmatrix}0 & 0\\1 & 1\end{bmatrix}\]
\end{minipage}\par\medskip
\noindent\begin{minipage}{\linewidth}\textbf{6. Inciso b: C+D: qué vamos a calcular}\\[2pt]Identificamos la operación y comprobamos sus dimensiones antes de efectuar las cuentas. Las matrices originales permanecen escritas en el enunciado.
\[C+D\]
\end{minipage}\par\medskip
\noindent\begin{minipage}{\linewidth}\textbf{7. Operamos las entradas que ocupan la misma posición}\\[2pt]La fila y la columna nos dicen qué valores combinar. Hacemos una cuenta por casilla y después reunimos todas las casillas en la matriz resultado. En cada casilla reducimos el resultado módulo 2: un número par da 0 y un número impar da 1.
\par\smallskip {\small\renewcommand{\arraystretch}{1.3}\begin{tabular}{>{\raggedright\arraybackslash}p{3.74cm}>{\raggedright\arraybackslash}p{3.74cm}>{\raggedright\arraybackslash}p{3.74cm}}\hline\rowcolor{black!7} \textbf{Casilla} & \textbf{Cuenta} & \textbf{Resultado}\\\hline Fila 1, columna 1 & $\mathrm{mod}((1)+(0),2)$ & $1$\\ Fila 1, columna 2 & $\mathrm{mod}((1)+(0),2)$ & $1$\\ Fila 2, columna 1 & $\mathrm{mod}((0)+(1),2)$ & $1$\\ Fila 2, columna 2 & $\mathrm{mod}((0)+(0),2)$ & $0$\\ \hline\end{tabular}}\par\smallskip
\par\smallskip \textbf{Resultado de esta operación}\[\begin{bmatrix}1 & 1\\1 & 0\end{bmatrix}\]
\par\smallskip\textbf{Qué aprendimos.} El resultado tiene 2 filas y 2 columnas.
\end{minipage}\par\medskip
\noindent\begin{minipage}{\linewidth}\textbf{8. Resultado del inciso b}\\[2pt]Reunimos el resultado. Comprueba que cada casilla procede de la operación indicada, sin cambiar su posición.
\par\smallskip \textbf{Resultado}\[\begin{bmatrix}1 & 1\\1 & 0\end{bmatrix}\]
\end{minipage}\par\medskip
\noindent\begin{minipage}{\linewidth}\textbf{9. Inciso c: A+B+(C+D)\textasciicircum{}T: qué vamos a calcular}\\[2pt]Identificamos la operación y comprobamos sus dimensiones antes de efectuar las cuentas. Las matrices originales permanecen escritas en el enunciado.
\[A+B+(C+D)^T\]
\end{minipage}\par\medskip
\noindent\begin{minipage}{\linewidth}\textbf{10. Operamos las entradas que ocupan la misma posición}\\[2pt]La fila y la columna nos dicen qué valores combinar. Hacemos una cuenta por casilla y después reunimos todas las casillas en la matriz resultado. En cada casilla reducimos el resultado módulo 2: un número par da 0 y un número impar da 1.
\par\smallskip {\small\renewcommand{\arraystretch}{1.3}\begin{tabular}{>{\raggedright\arraybackslash}p{3.74cm}>{\raggedright\arraybackslash}p{3.74cm}>{\raggedright\arraybackslash}p{3.74cm}}\hline\rowcolor{black!7} \textbf{Casilla} & \textbf{Cuenta} & \textbf{Resultado}\\\hline Fila 1, columna 1 & $\mathrm{mod}((1)+(1),2)$ & $0$\\ Fila 1, columna 2 & $\mathrm{mod}((0)+(0),2)$ & $0$\\ Fila 2, columna 1 & $\mathrm{mod}((1)+(0),2)$ & $1$\\ Fila 2, columna 2 & $\mathrm{mod}((0)+(1),2)$ & $1$\\ \hline\end{tabular}}\par\smallskip
\par\smallskip \textbf{Resultado de esta operación}\[\begin{bmatrix}0 & 0\\1 & 1\end{bmatrix}\]
\par\smallskip\textbf{Qué aprendimos.} El resultado tiene 2 filas y 2 columnas.
\end{minipage}\par\medskip
\noindent\begin{minipage}{\linewidth}\textbf{11. Operamos las entradas que ocupan la misma posición}\\[2pt]La fila y la columna nos dicen qué valores combinar. Hacemos una cuenta por casilla y después reunimos todas las casillas en la matriz resultado. En cada casilla reducimos el resultado módulo 2: un número par da 0 y un número impar da 1.
\par\smallskip {\small\renewcommand{\arraystretch}{1.3}\begin{tabular}{>{\raggedright\arraybackslash}p{3.74cm}>{\raggedright\arraybackslash}p{3.74cm}>{\raggedright\arraybackslash}p{3.74cm}}\hline\rowcolor{black!7} \textbf{Casilla} & \textbf{Cuenta} & \textbf{Resultado}\\\hline Fila 1, columna 1 & $\mathrm{mod}((1)+(0),2)$ & $1$\\ Fila 1, columna 2 & $\mathrm{mod}((1)+(0),2)$ & $1$\\ Fila 2, columna 1 & $\mathrm{mod}((0)+(1),2)$ & $1$\\ Fila 2, columna 2 & $\mathrm{mod}((0)+(0),2)$ & $0$\\ \hline\end{tabular}}\par\smallskip
\par\smallskip \textbf{Resultado de esta operación}\[\begin{bmatrix}1 & 1\\1 & 0\end{bmatrix}\]
\par\smallskip\textbf{Qué aprendimos.} El resultado tiene 2 filas y 2 columnas.
\end{minipage}\par\medskip
\noindent\begin{minipage}{\linewidth}\textbf{12. Transponemos: las filas se convierten en columnas}\\[2pt]No cambiamos los valores. Solo intercambiamos los dos índices de cada entrada. Una matriz de m filas y n columnas produce otra de n filas y m columnas.
\par\smallskip {\small\renewcommand{\arraystretch}{1.3}\begin{tabular}{>{\raggedright\arraybackslash}p{3.74cm}>{\raggedright\arraybackslash}p{3.74cm}>{\raggedright\arraybackslash}p{3.74cm}}\hline\rowcolor{black!7} \textbf{Posición antes} & \textbf{Posición después} & \textbf{Valor conservado}\\\hline (1, 1) & (1, 1) & $1$\\ (1, 2) & (2, 1) & $1$\\ (2, 1) & (1, 2) & $1$\\ (2, 2) & (2, 2) & $0$\\ \hline\end{tabular}}\par\smallskip
\par\smallskip \textbf{Resultado de esta operación}\[\begin{bmatrix}1 & 1\\1 & 0\end{bmatrix}\]
\par\smallskip\textbf{Qué aprendimos.} El resultado tiene 2 filas y 2 columnas.
\end{minipage}\par\medskip
\noindent\begin{minipage}{\linewidth}\textbf{13. Operamos las entradas que ocupan la misma posición}\\[2pt]La fila y la columna nos dicen qué valores combinar. Hacemos una cuenta por casilla y después reunimos todas las casillas en la matriz resultado. En cada casilla reducimos el resultado módulo 2: un número par da 0 y un número impar da 1.
\par\smallskip {\small\renewcommand{\arraystretch}{1.3}\begin{tabular}{>{\raggedright\arraybackslash}p{3.74cm}>{\raggedright\arraybackslash}p{3.74cm}>{\raggedright\arraybackslash}p{3.74cm}}\hline\rowcolor{black!7} \textbf{Casilla} & \textbf{Cuenta} & \textbf{Resultado}\\\hline Fila 1, columna 1 & $\mathrm{mod}((0)+(1),2)$ & $1$\\ Fila 1, columna 2 & $\mathrm{mod}((0)+(1),2)$ & $1$\\ Fila 2, columna 1 & $\mathrm{mod}((1)+(1),2)$ & $0$\\ Fila 2, columna 2 & $\mathrm{mod}((1)+(0),2)$ & $1$\\ \hline\end{tabular}}\par\smallskip
\par\smallskip \textbf{Resultado de esta operación}\[\begin{bmatrix}1 & 1\\0 & 1\end{bmatrix}\]
\par\smallskip\textbf{Qué aprendimos.} El resultado tiene 2 filas y 2 columnas.
\end{minipage}\par\medskip
\noindent\begin{minipage}{\linewidth}\textbf{14. Resultado del inciso c}\\[2pt]Reunimos el resultado. Comprueba que cada casilla procede de la operación indicada, sin cambiar su posición.
\par\smallskip \textbf{Resultado}\[\begin{bmatrix}1 & 1\\0 & 1\end{bmatrix}\]
\end{minipage}\par\medskip
\noindent\begin{minipage}{\linewidth}\textbf{15. Inciso d: C-B: qué vamos a calcular}\\[2pt]Identificamos la operación y comprobamos sus dimensiones antes de efectuar las cuentas. Las matrices originales permanecen escritas en el enunciado.
\[C-B\]
\end{minipage}\par\medskip
\noindent\begin{minipage}{\linewidth}\textbf{16. Operamos las entradas que ocupan la misma posición}\\[2pt]La fila y la columna nos dicen qué valores combinar. Hacemos una cuenta por casilla y después reunimos todas las casillas en la matriz resultado. En cada casilla reducimos el resultado módulo 2: un número par da 0 y un número impar da 1.
\par\smallskip {\small\renewcommand{\arraystretch}{1.3}\begin{tabular}{>{\raggedright\arraybackslash}p{3.74cm}>{\raggedright\arraybackslash}p{3.74cm}>{\raggedright\arraybackslash}p{3.74cm}}\hline\rowcolor{black!7} \textbf{Casilla} & \textbf{Cuenta} & \textbf{Resultado}\\\hline Fila 1, columna 1 & $\mathrm{mod}((1)-(1),2)$ & $0$\\ Fila 1, columna 2 & $\mathrm{mod}((1)-(0),2)$ & $1$\\ Fila 2, columna 1 & $\mathrm{mod}((0)-(0),2)$ & $0$\\ Fila 2, columna 2 & $\mathrm{mod}((0)-(1),2)$ & $1$\\ \hline\end{tabular}}\par\smallskip
\par\smallskip \textbf{Resultado de esta operación}\[\begin{bmatrix}0 & 1\\0 & 1\end{bmatrix}\]
\par\smallskip\textbf{Qué aprendimos.} El resultado tiene 2 filas y 2 columnas.
\end{minipage}\par\medskip
\noindent\begin{minipage}{\linewidth}\textbf{17. Resultado del inciso d}\\[2pt]Reunimos el resultado. Comprueba que cada casilla procede de la operación indicada, sin cambiar su posición.
\par\smallskip \textbf{Resultado}\[\begin{bmatrix}0 & 1\\0 & 1\end{bmatrix}\]
\end{minipage}\par\medskip
\noindent\begin{minipage}{\linewidth}\textbf{18. Inciso e: A-B+C-D: qué vamos a calcular}\\[2pt]Identificamos la operación y comprobamos sus dimensiones antes de efectuar las cuentas. Las matrices originales permanecen escritas en el enunciado.
\[A-B+C-D\]
\end{minipage}\par\medskip
\noindent\begin{minipage}{\linewidth}\textbf{19. Operamos las entradas que ocupan la misma posición}\\[2pt]La fila y la columna nos dicen qué valores combinar. Hacemos una cuenta por casilla y después reunimos todas las casillas en la matriz resultado. En cada casilla reducimos el resultado módulo 2: un número par da 0 y un número impar da 1.
\par\smallskip {\small\renewcommand{\arraystretch}{1.3}\begin{tabular}{>{\raggedright\arraybackslash}p{3.74cm}>{\raggedright\arraybackslash}p{3.74cm}>{\raggedright\arraybackslash}p{3.74cm}}\hline\rowcolor{black!7} \textbf{Casilla} & \textbf{Cuenta} & \textbf{Resultado}\\\hline Fila 1, columna 1 & $\mathrm{mod}((1)-(1),2)$ & $0$\\ Fila 1, columna 2 & $\mathrm{mod}((0)-(0),2)$ & $0$\\ Fila 2, columna 1 & $\mathrm{mod}((1)-(0),2)$ & $1$\\ Fila 2, columna 2 & $\mathrm{mod}((0)-(1),2)$ & $1$\\ \hline\end{tabular}}\par\smallskip
\par\smallskip \textbf{Resultado de esta operación}\[\begin{bmatrix}0 & 0\\1 & 1\end{bmatrix}\]
\par\smallskip\textbf{Qué aprendimos.} El resultado tiene 2 filas y 2 columnas.
\end{minipage}\par\medskip
\noindent\begin{minipage}{\linewidth}\textbf{20. Operamos las entradas que ocupan la misma posición}\\[2pt]La fila y la columna nos dicen qué valores combinar. Hacemos una cuenta por casilla y después reunimos todas las casillas en la matriz resultado. En cada casilla reducimos el resultado módulo 2: un número par da 0 y un número impar da 1.
\par\smallskip {\small\renewcommand{\arraystretch}{1.3}\begin{tabular}{>{\raggedright\arraybackslash}p{3.74cm}>{\raggedright\arraybackslash}p{3.74cm}>{\raggedright\arraybackslash}p{3.74cm}}\hline\rowcolor{black!7} \textbf{Casilla} & \textbf{Cuenta} & \textbf{Resultado}\\\hline Fila 1, columna 1 & $\mathrm{mod}((0)+(1),2)$ & $1$\\ Fila 1, columna 2 & $\mathrm{mod}((0)+(1),2)$ & $1$\\ Fila 2, columna 1 & $\mathrm{mod}((1)+(0),2)$ & $1$\\ Fila 2, columna 2 & $\mathrm{mod}((1)+(0),2)$ & $1$\\ \hline\end{tabular}}\par\smallskip
\par\smallskip \textbf{Resultado de esta operación}\[\begin{bmatrix}1 & 1\\1 & 1\end{bmatrix}\]
\par\smallskip\textbf{Qué aprendimos.} El resultado tiene 2 filas y 2 columnas.
\end{minipage}\par\medskip
\noindent\begin{minipage}{\linewidth}\textbf{21. Operamos las entradas que ocupan la misma posición}\\[2pt]La fila y la columna nos dicen qué valores combinar. Hacemos una cuenta por casilla y después reunimos todas las casillas en la matriz resultado. En cada casilla reducimos el resultado módulo 2: un número par da 0 y un número impar da 1.
\par\smallskip {\small\renewcommand{\arraystretch}{1.3}\begin{tabular}{>{\raggedright\arraybackslash}p{3.74cm}>{\raggedright\arraybackslash}p{3.74cm}>{\raggedright\arraybackslash}p{3.74cm}}\hline\rowcolor{black!7} \textbf{Casilla} & \textbf{Cuenta} & \textbf{Resultado}\\\hline Fila 1, columna 1 & $\mathrm{mod}((1)-(0),2)$ & $1$\\ Fila 1, columna 2 & $\mathrm{mod}((1)-(0),2)$ & $1$\\ Fila 2, columna 1 & $\mathrm{mod}((1)-(1),2)$ & $0$\\ Fila 2, columna 2 & $\mathrm{mod}((1)-(0),2)$ & $1$\\ \hline\end{tabular}}\par\smallskip
\par\smallskip \textbf{Resultado de esta operación}\[\begin{bmatrix}1 & 1\\0 & 1\end{bmatrix}\]
\par\smallskip\textbf{Qué aprendimos.} El resultado tiene 2 filas y 2 columnas.
\end{minipage}\par\medskip
\noindent\begin{minipage}{\linewidth}\textbf{22. Resultado del inciso e}\\[2pt]Reunimos el resultado. Comprueba que cada casilla procede de la operación indicada, sin cambiar su posición.
\par\smallskip \textbf{Resultado}\[\begin{bmatrix}1 & 1\\0 & 1\end{bmatrix}\]
\end{minipage}\par\medskip
\respuesta{a: A+B: [[0, 0], [1, 1]]; b: C+D: [[1, 1], [1, 0]]; c: A+B+(C+D)\textasciicircum{}T: [[1, 1], [0, 1]]; d: C-B: [[0, 1], [0, 1]]; e: A-B+C-D: [[1, 1], [0, 1]].}
\noindent\begin{minipage}{\linewidth}\centering\includegraphics[width=.96\linewidth,height=6.8cm,keepaspectratio]{../recursos/grafico_1_2-12.png}\\{\small Lectura visual complementaria; la demostración está en las operaciones.}\end{minipage}
\end{multicols*}
\newpage\fancyfoot[L]{\small Sección 1.2 | Ejercicio 13 | Paso a paso}\hypertarget{1.2-13}{}
\pdfbookmark[1]{1.2 / 13}{bm1.2-13}
\titulo{13 | Encontrar matrices binarias B y C}
\textbf{Enunciado}\par Sea A la matriz indicada. (a) Determine B de modo que A+B sea la matriz cero. (b) Determine C de modo que A+C tenga todas sus entradas iguales a 1. Las operaciones son módulo 2.
\par\smallskip \textbf{A}\[\begin{bmatrix}1 & 0\\0 & 0\end{bmatrix}\]
\begin{multicols*}{2}
\noindent\begin{minipage}{\linewidth}\textbf{1. Regla binaria que usaremos}\\[2pt]Si sumamos dos bits iguales obtenemos 0; si son distintos obtenemos 1. Para producir ceros copiamos el bit original. Para producir unos elegimos el bit contrario.
\par\smallskip {\small\renewcommand{\arraystretch}{1.3}\begin{tabular}{>{\raggedright\arraybackslash}p{3.74cm}>{\raggedright\arraybackslash}p{3.74cm}>{\raggedright\arraybackslash}p{3.74cm}}\hline\rowcolor{black!7} \textbf{a} & \textbf{b para obtener 0} & \textbf{c para obtener 1}\\\hline 0 & 0 & 1\\ 1 & 1 & 0\\ \hline\end{tabular}}\par\smallskip
\end{minipage}\par\medskip
\noindent\begin{minipage}{\linewidth}\textbf{2. Inciso (a): B debe coincidir con A}\\[2pt]Cada casilla de A se suma con la misma casilla de B. Copiamos las entradas, porque a+a = 0 módulo 2.
\par\smallskip \begin{minipage}[t]{.47\linewidth}\centering \textbf{B}\[\begin{bmatrix}1 & 0\\0 & 0\end{bmatrix}\]\end{minipage}\hfill\begin{minipage}[t]{.47\linewidth}\centering \textbf{A + B, módulo 2}\[\begin{bmatrix}0 & 0\\0 & 0\end{bmatrix}\]\end{minipage}\par\smallskip
\end{minipage}\par\medskip
\noindent\begin{minipage}{\linewidth}\textbf{3. Inciso (b): C debe complementar A}\\[2pt]Cambiamos 1 por 0 y 0 por 1 en cada casilla.
\par\smallskip \begin{minipage}[t]{.47\linewidth}\centering \textbf{C}\[\begin{bmatrix}0 & 1\\1 & 1\end{bmatrix}\]\end{minipage}\hfill\begin{minipage}[t]{.47\linewidth}\centering \textbf{A + C, módulo 2}\[\begin{bmatrix}1 & 1\\1 & 1\end{bmatrix}\]\end{minipage}\par\smallskip
\end{minipage}\par\medskip
\respuesta{B = [[1,0],[0,0]]; C = [[0,1],[1,1]].}
\noindent\begin{minipage}{\linewidth}\centering\includegraphics[width=.96\linewidth,height=6.8cm,keepaspectratio]{../recursos/grafico_1_2-13.png}\\{\small Lectura visual complementaria; la demostración está en las operaciones.}\end{minipage}
\end{multicols*}
\newpage\fancyfoot[L]{\small Sección 1.2 | Ejercicio 14 | Paso a paso}\hypertarget{1.2-14}{}
\pdfbookmark[1]{1.2 / 14}{bm1.2-14}
\titulo{14 | Completar un vector binario}
\textbf{Enunciado}\par Sea u el vector indicado. Determine el 4-vector v tal que u+v sea el vector objetivo, con aritmética binaria.
\par\smallskip \begin{minipage}[t]{.47\linewidth}\centering \textbf{u}\[\begin{bmatrix}1 & 1 & 0 & 0\end{bmatrix}\]\end{minipage}\hfill\begin{minipage}[t]{.47\linewidth}\centering \textbf{Objetivo}\[\begin{bmatrix}1 & 1 & 0 & 0\end{bmatrix}\]\end{minipage}\par\smallskip
\begin{multicols*}{2}
\noindent\begin{minipage}{\linewidth}\textbf{1. Una ecuación por coordenada}\\[2pt]Cada casilla se resuelve de manera independiente. En módulo 2, v se obtiene sumando u al objetivo: u + u = 0.
\[v_i=w_i+u_i\quad(\mathrm{mod}\ 2)\]
\end{minipage}\par\medskip
\noindent\begin{minipage}{\linewidth}\textbf{2. Elegimos cada bit de v}\\[2pt]Si los bits de u y del objetivo son iguales, necesitamos 0. Si son distintos, necesitamos 1.
\par\smallskip {\small\renewcommand{\arraystretch}{1.3}\begin{tabular}{>{\raggedright\arraybackslash}p{3.74cm}>{\raggedright\arraybackslash}p{3.74cm}>{\raggedright\arraybackslash}p{3.74cm}}\hline\rowcolor{black!7} \textbf{Coordenada} & \textbf{u, objetivo} & \textbf{v}\\\hline 1 & 1, 1 & 0\\ 2 & 1, 1 & 0\\ 3 & 0, 0 & 0\\ 4 & 0, 0 & 0\\ \hline\end{tabular}}\par\smallskip
\par\smallskip \textbf{v}\[\begin{bmatrix}0 & 0 & 0 & 0\end{bmatrix}\]
\end{minipage}\par\medskip
\noindent\begin{minipage}{\linewidth}\textbf{3. Comprobamos la suma}\\[2pt]Sumamos los bits de u y v y reducimos cada casilla módulo 2. El resultado coincide con el objetivo.
\par\smallskip \textbf{u + v, módulo 2}\[\begin{bmatrix}1 & 1 & 0 & 0\end{bmatrix}\]
\end{minipage}\par\medskip
\respuesta{[0, 0, 0, 0]}
\noindent\begin{minipage}{\linewidth}\centering\includegraphics[width=.96\linewidth,height=6.8cm,keepaspectratio]{../recursos/grafico_1_2-14.png}\\{\small Lectura visual complementaria; la demostración está en las operaciones.}\end{minipage}
\end{multicols*}
\newpage\fancyfoot[L]{\small Sección 1.2 | Ejercicio 15 | Paso a paso}\hypertarget{1.2-15}{}
\pdfbookmark[1]{1.2 / 15}{bm1.2-15}
\titulo{15 | Completar un vector binario}
\textbf{Enunciado}\par Sea u el vector indicado. Determine el 4-vector v tal que u+v sea el vector objetivo, con aritmética binaria.
\par\smallskip \begin{minipage}[t]{.47\linewidth}\centering \textbf{u}\[\begin{bmatrix}0 & 1 & 0 & 1\end{bmatrix}\]\end{minipage}\hfill\begin{minipage}[t]{.47\linewidth}\centering \textbf{Objetivo}\[\begin{bmatrix}1 & 1 & 1 & 1\end{bmatrix}\]\end{minipage}\par\smallskip
\begin{multicols*}{2}
\noindent\begin{minipage}{\linewidth}\textbf{1. Una ecuación por coordenada}\\[2pt]Cada casilla se resuelve de manera independiente. En módulo 2, v se obtiene sumando u al objetivo: u + u = 0.
\[v_i=w_i+u_i\quad(\mathrm{mod}\ 2)\]
\end{minipage}\par\medskip
\noindent\begin{minipage}{\linewidth}\textbf{2. Elegimos cada bit de v}\\[2pt]Si los bits de u y del objetivo son iguales, necesitamos 0. Si son distintos, necesitamos 1.
\par\smallskip {\small\renewcommand{\arraystretch}{1.3}\begin{tabular}{>{\raggedright\arraybackslash}p{3.74cm}>{\raggedright\arraybackslash}p{3.74cm}>{\raggedright\arraybackslash}p{3.74cm}}\hline\rowcolor{black!7} \textbf{Coordenada} & \textbf{u, objetivo} & \textbf{v}\\\hline 1 & 0, 1 & 1\\ 2 & 1, 1 & 0\\ 3 & 0, 1 & 1\\ 4 & 1, 1 & 0\\ \hline\end{tabular}}\par\smallskip
\par\smallskip \textbf{v}\[\begin{bmatrix}1 & 0 & 1 & 0\end{bmatrix}\]
\end{minipage}\par\medskip
\noindent\begin{minipage}{\linewidth}\textbf{3. Comprobamos la suma}\\[2pt]Sumamos los bits de u y v y reducimos cada casilla módulo 2. El resultado coincide con el objetivo.
\par\smallskip \textbf{u + v, módulo 2}\[\begin{bmatrix}1 & 1 & 1 & 1\end{bmatrix}\]
\end{minipage}\par\medskip
\respuesta{[1, 0, 1, 0]}
\noindent\begin{minipage}{\linewidth}\centering\includegraphics[width=.96\linewidth,height=6.8cm,keepaspectratio]{../recursos/grafico_1_2-15.png}\\{\small Lectura visual complementaria; la demostración está en las operaciones.}\end{minipage}
\end{multicols*}
\newpage\fancyfoot[L]{\small Sección 1.2 | Ejercicio T.1 | Paso a paso}\hypertarget{1.2-T.1}{}
\pdfbookmark[1]{1.2 / T.1}{bm1.2-T.1}
\titulo{T.1 | Suma y diferencia de matrices diagonales}
\textbf{Enunciado}\par Demuestre que la suma y la diferencia de dos matrices diagonales son matrices diagonales.
\par\smallskip \begin{minipage}[t]{.47\linewidth}\centering \textbf{A}\[\begin{bmatrix}2 & 0 & 0\\0 & -1 & 0\\0 & 0 & 3\end{bmatrix}\]\end{minipage}\hfill\begin{minipage}[t]{.47\linewidth}\centering \textbf{B}\[\begin{bmatrix}4 & 0 & 0\\0 & 0 & 0\\0 & 0 & -2\end{bmatrix}\]\end{minipage}\par\smallskip
\par\medskip\large
\noindent\begin{minipage}{\linewidth}\textbf{1. Hipótesis y significado de diagonal}\\[2pt]Sean A y B matrices diagonales reales del mismo orden n. Por definición, sus entradas fuera de la diagonal principal valen cero. Debemos probar que lo mismo ocurre con su suma y con su diferencia.
\[i\ne j\quad\Longrightarrow\quad a_{ij}=b_{ij}=0\]
\end{minipage}\par\medskip
\noindent\begin{minipage}{\linewidth}\textbf{2. Demostración para la suma}\\[2pt]Fijemos índices arbitrarios i y j distintos. Por la definición de suma de matrices, la entrada correspondiente es la suma de las dos entradas originales. Ambas son cero.
\[(A+B)_{ij}=a_{ij}+b_{ij}=0+0=0\]
\end{minipage}\par\medskip
\noindent\begin{minipage}{\linewidth}\textbf{3. Demostración para la diferencia}\\[2pt]Para los mismos índices arbitrarios, aplicamos la definición de resta.
\[(A-B)_{ij}=a_{ij}-b_{ij}=0-0=0\]
\end{minipage}\par\medskip
\noindent\begin{minipage}{\linewidth}\textbf{4. Conclusión formal}\\[2pt]Todas las entradas fuera de la diagonal son cero en A+B y A-B. Por definición, las dos matrices son diagonales. Las entradas de la diagonal pueden tomar cualquier valor. Queda demostrado.
\end{minipage}\par\medskip
\respuesta{Ambas son diagonales, siempre que los dos órdenes coincidan.}
\noindent\begin{minipage}{\linewidth}\centering\includegraphics[width=.96\linewidth,height=6.8cm,keepaspectratio]{../recursos/grafico_1_2-T_1.png}\\{\small Lectura visual complementaria; la demostración está en las operaciones.}\end{minipage}
\normalsize
\newpage\fancyfoot[L]{\small Sección 1.2 | Ejercicio T.2 | Paso a paso}\hypertarget{1.2-T.2}{}
\pdfbookmark[1]{1.2 / T.2}{bm1.2-T.2}
\titulo{T.2 | Suma y diferencia de matrices escalares}
\textbf{Enunciado}\par Demuestre que la suma y la diferencia de dos matrices escalares son matrices escalares.
\par\smallskip \begin{minipage}[t]{.47\linewidth}\centering \textbf{A}\[\begin{bmatrix}2 & 0 & 0\\0 & 2 & 0\\0 & 0 & 2\end{bmatrix}\]\end{minipage}\hfill\begin{minipage}[t]{.47\linewidth}\centering \textbf{B}\[\begin{bmatrix}-3 & 0 & 0\\0 & -3 & 0\\0 & 0 & -3\end{bmatrix}\]\end{minipage}\par\smallskip
\par\medskip\large
\noindent\begin{minipage}{\linewidth}\textbf{1. Hipótesis: qué es una matriz escalar}\\[2pt]Sean A y B matrices escalares del mismo orden n. Existen números reales $\alpha$ y $\beta$ tales que A = $\alpha$I y B = $\beta$I, donde I es la identidad de orden n.
\[A=\alpha I_n,\quad B=\beta I_n\]
\end{minipage}\par\medskip
\noindent\begin{minipage}{\linewidth}\textbf{2. Demostramos la afirmación para la suma}\\[2pt]Usamos la distributividad de la multiplicación de un escalar sobre la suma. El coeficiente $\alpha$+$\beta$ es un número real.
\[A+B=\alpha I_n+\beta I_n=(\alpha+\beta)I_n\]
\end{minipage}\par\medskip
\noindent\begin{minipage}{\linewidth}\textbf{3. Demostramos la afirmación para la diferencia}\\[2pt]Aplicamos la misma propiedad con el coeficiente -$\beta$.
\[A-B=\alpha I_n-\beta I_n=(\alpha-\beta)I_n\]
\end{minipage}\par\medskip
\noindent\begin{minipage}{\linewidth}\textbf{4. Conclusión formal}\\[2pt]Ambos resultados son múltiplos reales de la identidad. Esa es precisamente la definición de matriz escalar; por tanto la suma y la diferencia son escalares. Queda demostrado.
\end{minipage}\par\medskip
\respuesta{Ambas son escalares, siempre que los dos órdenes coincidan.}
\noindent\begin{minipage}{\linewidth}\centering\includegraphics[width=.96\linewidth,height=6.8cm,keepaspectratio]{../recursos/grafico_1_2-T_2.png}\\{\small Lectura visual complementaria; la demostración está en las operaciones.}\end{minipage}
\normalsize
\newpage\fancyfoot[L]{\small Sección 1.2 | Ejercicio T.3 | Paso a paso}\hypertarget{1.2-T.3}{}
\pdfbookmark[1]{1.2 / T.3}{bm1.2-T.3}
\titulo{T.3 | Una matriz simbólica y su transpuesta}
\textbf{Enunciado}\par Sea A la matriz simbólica indicada. (a) Calcule A-A\textsuperscript{T}. (b) Calcule A+A\textsuperscript{T}. (c) Calcule (A+A\textsuperscript{T})\textsuperscript{T}.
\par\smallskip \textbf{A}\[\begin{bmatrix}a & b & c\\c & d & e\\e & e & f\end{bmatrix}\]
\par\medskip\large
\noindent\begin{minipage}{\linewidth}\textbf{1. Primero transponemos A}\\[2pt]La entrada (i, j) pasa a la posición (j, i). Observa que las entradas (2, 3) y (3, 2) son ambas e, pero las entradas (1, 2) y (2, 1) son b y c.
\par\smallskip \textbf{A\textsuperscript{T}}\[\begin{bmatrix}a & c & e\\b & d & e\\c & e & f\end{bmatrix}\]
\end{minipage}\par\medskip
\noindent\begin{minipage}{\linewidth}\textbf{2. Inciso (a): restamos casilla por casilla}\\[2pt]En la diagonal aparecen a-a, d-d y f-f: dan cero. Las entradas (2, 3) y (3, 2) también dan e-e = 0.
\par\smallskip \textbf{A - A\textsuperscript{T}}\[\begin{bmatrix}0 & b - c & c - e\\- b + c & 0 & 0\\- c + e & 0 & 0\end{bmatrix}\]
\end{minipage}\par\medskip
\noindent\begin{minipage}{\linewidth}\textbf{3. Inciso (b): sumamos las mismas casillas}\\[2pt]En la diagonal se duplican los valores; en cada pareja fuera de la diagonal aparece la misma suma, en orden inverso.
\par\smallskip \textbf{A + A\textsuperscript{T}}\[\begin{bmatrix}2 a & b + c & c + e\\b + c & 2 d & 2 e\\c + e & 2 e & 2 f\end{bmatrix}\]
\end{minipage}\par\medskip
\noindent\begin{minipage}{\linewidth}\textbf{4. Inciso (c): transponemos el resultado}\\[2pt]Sea M = A+A\textsuperscript{T}. Para índices arbitrarios i, j, mji = aji+aij = aij+aji = mij, por la conmutatividad de los reales. Por definición, M\textsuperscript{T} = M.
\[(A+A^T)^T_{ij}=a_{ji}+a_{ij}=a_{ij}+a_{ji}=(A+A^T)_{ij}\]
\par\smallskip \textbf{(A + A\textsuperscript{T})\textsuperscript{T}}\[\begin{bmatrix}2 a & b + c & c + e\\b + c & 2 d & 2 e\\c + e & 2 e & 2 f\end{bmatrix}\]
\par\smallskip\textbf{Qué aprendimos.} La matriz A+A\textsuperscript{T} es simétrica. No hemos supuesto que A lo sea.
\end{minipage}\par\medskip
\respuesta{(a) [[0,b-c,c-e],[c-b,0,0],[e-c,0,0]]. (b) y (c) [[2a,b+c,c+e],[b+c,2d,2e],[c+e,2e,2f]].}
\noindent\begin{minipage}{\linewidth}\centering\includegraphics[width=.96\linewidth,height=6.8cm,keepaspectratio]{../recursos/grafico_1_2-T_3.png}\\{\small Lectura visual complementaria; la demostración está en las operaciones.}\end{minipage}
\normalsize
\newpage\fancyfoot[L]{\small Sección 1.2 | Ejercicio T.4 | Paso a paso}\hypertarget{1.2-T.4}{}
\pdfbookmark[1]{1.2 / T.4}{bm1.2-T.4}
\titulo{T.4 | Si kA es cero, k es cero o A es cero}
\textbf{Enunciado}\par Sea O la matriz n por n cuyas entradas son cero. Demuestre que, si k es real y kA = O, entonces k = 0 o A = O.
\par\smallskip \textbf{A}\[\begin{bmatrix}1 & 2\\0 & 3\end{bmatrix}\]
\par\medskip\large
\noindent\begin{minipage}{\linewidth}\textbf{1. Hipótesis y objetivo}\\[2pt]Sea A una matriz real de orden n y sea k real. Suponemos kA = O. Debemos probar una disyunción: al menos una de las afirmaciones k = 0 y A = O es verdadera.
\end{minipage}\par\medskip
\noindent\begin{minipage}{\linewidth}\textbf{2. Primer caso: k es cero}\\[2pt]Si k = 0, la primera afirmación de la conclusión ya se cumple. No hace falta exigir que A sea cero.
\end{minipage}\par\medskip
\noindent\begin{minipage}{\linewidth}\textbf{3. Segundo caso: k no es cero}\\[2pt]Por igualdad de matrices, cada entrada de kA es cero. Fijemos una posición arbitraria (i, j). Como k no es cero, podemos dividir ambos miembros entre k.
\[ka_{ij}=0,\quad k\ne0\]
\[a_{ij}=\frac{0}{k}=0\]
\end{minipage}\par\medskip
\noindent\begin{minipage}{\linewidth}\textbf{4. Conclusión formal}\\[2pt]El segundo caso muestra que todas las entradas de A son cero; luego A = O. Los casos k = 0 y k distinto de cero agotan todos los reales. Por tanto, siempre se cumple k = 0 o A = O. Queda demostrado.
\par\smallskip\textbf{Qué aprendimos.} La palabra «o» permite que ambas afirmaciones sean verdaderas. No estamos afirmando que exactamente una lo sea.
\end{minipage}\par\medskip
\respuesta{kA = O implica k = 0 o A = O.}
\noindent\begin{minipage}{\linewidth}\centering\includegraphics[width=.96\linewidth,height=6.8cm,keepaspectratio]{../recursos/grafico_1_2-T_4.png}\\{\small Lectura visual complementaria; la demostración está en las operaciones.}\end{minipage}
\normalsize
\newpage\fancyfoot[L]{\small Sección 1.2 | Ejercicio T.5 | Paso a paso}\hypertarget{1.2-T.5}{}
\pdfbookmark[1]{1.2 / T.5}{bm1.2-T.5}
\titulo{T.5 | Matrices triangulares y diagonales}
\textbf{Enunciado}\par Una matriz es triangular superior si aij = 0 para i > j, e inferior si aij = 0 para i < j. (a) Pruebe el cierre de las superiores bajo suma y resta. (b) Pruebe lo mismo para las inferiores. (c) Pruebe que ser de ambos tipos implica ser diagonal.
\par\smallskip \textbf{Ejemplo superior}\[\begin{bmatrix}1 & 2 & 3\\0 & 4 & 5\\0 & 0 & 6\end{bmatrix}\]
\par\medskip\large
\noindent\begin{minipage}{\linewidth}\textbf{1. Hipótesis comunes}\\[2pt]En los incisos (a) y (b), A y B son matrices reales cuadradas del mismo orden n. La condición «superior» anula las casillas debajo de la diagonal; la condición «inferior» anula las que están encima.
\end{minipage}\par\medskip
\noindent\begin{minipage}{\linewidth}\textbf{2. Inciso (a): fijamos una posición debajo de la diagonal}\\[2pt]Supongamos A y B triangulares superiores. Si i > j, entonces aij = bij = 0. Aplicando la definición de suma y de diferencia:
\[(A+B)_{ij}=0+0=0,\quad (A-B)_{ij}=0-0=0\quad(i>j)\]
\par\smallskip\textbf{Qué aprendimos.} Como los índices fueron arbitrarios, todas las entradas debajo de la diagonal son cero. A+B y A-B son superiores.
\end{minipage}\par\medskip
\noindent\begin{minipage}{\linewidth}\textbf{3. Inciso (b): fijamos una posición encima de la diagonal}\\[2pt]Supongamos A y B triangulares inferiores. Si i < j, entonces aij = bij = 0. Por tanto:
\[(A+B)_{ij}=0+0=0,\quad (A-B)_{ij}=0-0=0\quad(i<j)\]
\par\smallskip\textbf{Qué aprendimos.} Toda entrada encima de la diagonal es cero. Los dos resultados son inferiores.
\end{minipage}\par\medskip
\noindent\begin{minipage}{\linewidth}\textbf{4. Inciso (c): combinamos las dos condiciones}\\[2pt]Sea A superior e inferior a la vez. En una posición fuera de la diagonal se cumple i > j o i < j. En el primer caso la condición superior da aij = 0; en el segundo, la inferior da aij = 0. Por tanto toda entrada fuera de la diagonal es cero: A es diagonal. Quedan demostrados los tres incisos.
\end{minipage}\par\medskip
\respuesta{La suma y resta conservan cada tipo triangular; ser superior e inferior a la vez equivale a ser diagonal.}
\noindent\begin{minipage}{\linewidth}\centering\includegraphics[width=.96\linewidth,height=6.8cm,keepaspectratio]{../recursos/grafico_1_2-T_5.png}\\{\small Lectura visual complementaria; la demostración está en las operaciones.}\end{minipage}
\normalsize
\newpage\fancyfoot[L]{\small Sección 1.2 | Ejercicio T.6 | Paso a paso}\hypertarget{1.2-T.6}{}
\pdfbookmark[1]{1.2 / T.6}{bm1.2-T.6}
\titulo{T.6 | La transpuesta intercambia los tipos triangulares}
\textbf{Enunciado}\par (a) Demuestre que, si A es triangular superior, A\textsuperscript{T} es triangular inferior. (b) Demuestre que, si A es triangular inferior, A\textsuperscript{T} es triangular superior.
\par\medskip\large
\noindent\begin{minipage}{\linewidth}\textbf{1. Definición que usaremos}\\[2pt]Sea A cuadrada de orden n. La definición de transpuesta dice que la entrada (i, j) de A\textsuperscript{T} es la entrada (j, i) de A.
\[(A^T)_{ij}=a_{ji}\]
\end{minipage}\par\medskip
\noindent\begin{minipage}{\linewidth}\textbf{2. Inciso (a): demostrar que A\textsuperscript{T} es inferior}\\[2pt]Supongamos A superior. Para probar que A\textsuperscript{T} es inferior, debemos anular sus posiciones con i < j. Fijemos una de ellas. Entonces j > i: la entrada aji está debajo de la diagonal de A y vale cero.
\[i<j\quad\Longrightarrow\quad j>i\quad\Longrightarrow\quad (A^T)_{ij}=a_{ji}=0\]
\par\smallskip\textbf{Qué aprendimos.} A\textsuperscript{T} no tiene entradas no nulas encima de la diagonal; es triangular inferior.
\end{minipage}\par\medskip
\noindent\begin{minipage}{\linewidth}\textbf{3. Inciso (b): demostrar que A\textsuperscript{T} es superior}\\[2pt]Supongamos A inferior. Ahora fijamos i > j. Entonces j < i: la entrada aji está encima de la diagonal de A y vale cero.
\[i>j\quad\Longrightarrow\quad j<i\quad\Longrightarrow\quad (A^T)_{ij}=a_{ji}=0\]
\par\smallskip\textbf{Qué aprendimos.} A\textsuperscript{T} no tiene entradas no nulas debajo de la diagonal; es triangular superior. Queda demostrado.
\end{minipage}\par\medskip
\respuesta{Superior $\to$ transpuesta inferior; inferior $\to$ transpuesta superior.}
\noindent\begin{minipage}{\linewidth}\centering\includegraphics[width=.96\linewidth,height=6.8cm,keepaspectratio]{../recursos/grafico_1_2-T_6.png}\\{\small Lectura visual complementaria; la demostración está en las operaciones.}\end{minipage}
\normalsize
\newpage\fancyfoot[L]{\small Sección 1.2 | Ejercicio T.7 | Paso a paso}\hypertarget{1.2-T.7}{}
\pdfbookmark[1]{1.2 / T.7}{bm1.2-T.7}
\titulo{T.7 | La diagonal de A-A\textsuperscript{T}}
\textbf{Enunciado}\par Si A es una matriz n por n, ¿cuáles son las entradas de la diagonal principal de A-A\textsuperscript{T}? Justifique.
\par\medskip\large
\noindent\begin{minipage}{\linewidth}\textbf{1. Elegimos una entrada arbitraria de la diagonal}\\[2pt]Sea i cualquier índice entre 1 y n. En la diagonal, fila y columna tienen el mismo índice. Transponer no mueve esa casilla.
\[(A^T)_{ii}=a_{ii}\]
\end{minipage}\par\medskip
\noindent\begin{minipage}{\linewidth}\textbf{2. Calculamos la diferencia en esa casilla}\\[2pt]Aplicamos la definición de resta de matrices.
\[(A-A^T)_{ii}=a_{ii}-(A^T)_{ii}=a_{ii}-a_{ii}=0\]
\end{minipage}\par\medskip
\noindent\begin{minipage}{\linewidth}\textbf{3. Conclusión formal}\\[2pt]Como i fue arbitrario, todas las entradas de la diagonal principal son cero. No hace falta que A sea simétrica: la conclusión vale para cualquier matriz cuadrada real.
\end{minipage}\par\medskip
\respuesta{Todas son cero.}
\noindent\begin{minipage}{\linewidth}\centering\includegraphics[width=.96\linewidth,height=6.8cm,keepaspectratio]{../recursos/grafico_1_2-T_7.png}\\{\small Lectura visual complementaria; la demostración está en las operaciones.}\end{minipage}
\normalsize
\newpage\fancyfoot[L]{\small Sección 1.2 | Ejercicio T.8 | Paso a paso}\hypertarget{1.2-T.8}{}
\pdfbookmark[1]{1.2 / T.8}{bm1.2-T.8}
\titulo{T.8 | El vector cero es neutro para la suma}
\textbf{Enunciado}\par Si x es un n-vector, demuestre que x + 0 = x.
\par\medskip\large
\noindent\begin{minipage}{\linewidth}\textbf{1. Escribimos las coordenadas}\\[2pt]Sea x un vector real arbitrario de n coordenadas. El vector cero tiene n coordenadas, todas iguales a cero.
\[\mathbf{x}=(x_1,\ldots,x_n),\quad\mathbf{0}=(0,\ldots,0)\]
\end{minipage}\par\medskip
\noindent\begin{minipage}{\linewidth}\textbf{2. Sumamos por definición}\\[2pt]La suma de vectores se realiza coordenada a coordenada. Para cada i, el cero es neutro en los números reales.
\[(\mathbf{x}+\mathbf{0})_i=x_i+0=x_i\quad(i=1,\ldots,n)\]
\end{minipage}\par\medskip
\noindent\begin{minipage}{\linewidth}\textbf{3. Concluimos la igualdad de vectores}\\[2pt]Todas las coordenadas de x+0 coinciden con las de x. Por definición de igualdad de vectores, x+0 = x. El vector x y la dimensión n fueron arbitrarios. Queda demostrado.
\end{minipage}\par\medskip
\respuesta{x + 0 = x para todo n-vector real x.}
\noindent\begin{minipage}{\linewidth}\centering\includegraphics[width=.96\linewidth,height=6.8cm,keepaspectratio]{../recursos/grafico_1_2-T_8.png}\\{\small Lectura visual complementaria; la demostración está en las operaciones.}\end{minipage}
\normalsize
\newpage\fancyfoot[L]{\small Sección 1.2 | Ejercicio T.9 | Paso a paso}\hypertarget{1.2-T.9}{}
\pdfbookmark[1]{1.2 / T.9}{bm1.2-T.9}
\titulo{T.9 | Todos los 2-vectores binarios}
\textbf{Enunciado}\par Haga una lista de todos los posibles 2-vectores binarios. ¿Cuántos hay?
\par\medskip\large
\noindent\begin{minipage}{\linewidth}\textbf{1. Contamos las elecciones por coordenada}\\[2pt]Cada coordenada puede elegirse de dos maneras: 0 o 1. Las elecciones son independientes. Multiplicamos 2 por sí mismo una vez por cada posición.
\[2^{2}=4\]
\end{minipage}\par\medskip
\noindent\begin{minipage}{\linewidth}\textbf{2. Lista completa, sin repeticiones}\\[2pt]Ordenamos las listas de bits desde todo ceros hasta todo unos. Cada vector corresponde a una elección distinta de sus coordenadas.
\par\smallskip {\small\renewcommand{\arraystretch}{1.3}\begin{tabular}{>{\raggedright\arraybackslash}p{5.85cm}>{\raggedright\arraybackslash}p{5.85cm}}\hline\rowcolor{black!7} \textbf{Vector} & \textbf{Vector}\\\hline (0, 0) & (0, 1)\\ (1, 0) & (1, 1)\\ \hline\end{tabular}}\par\smallskip
\end{minipage}\par\medskip
\noindent\begin{minipage}{\linewidth}\textbf{3. Por qué no falta ningún vector}\\[2pt]Para cualquier vector, su primera coordenada es 0 o 1; dentro de cada grupo hacemos lo mismo con la siguiente, hasta completar las n posiciones. Este procedimiento recorre todas las elecciones una sola vez.
\end{minipage}\par\medskip
\respuesta{Hay 4 vectores distintos.}
\noindent\begin{minipage}{\linewidth}\centering\includegraphics[width=.96\linewidth,height=6.8cm,keepaspectratio]{../recursos/grafico_1_2-T_9.png}\\{\small Lectura visual complementaria; la demostración está en las operaciones.}\end{minipage}
\normalsize
\newpage\fancyfoot[L]{\small Sección 1.2 | Ejercicio T.10 | Paso a paso}\hypertarget{1.2-T.10}{}
\pdfbookmark[1]{1.2 / T.10}{bm1.2-T.10}
\titulo{T.10 | Todos los 3-vectores binarios}
\textbf{Enunciado}\par Haga una lista de todos los posibles 3-vectores binarios. ¿Cuántos hay?
\par\medskip\large
\noindent\begin{minipage}{\linewidth}\textbf{1. Contamos las elecciones por coordenada}\\[2pt]Cada coordenada puede elegirse de dos maneras: 0 o 1. Las elecciones son independientes. Multiplicamos 2 por sí mismo una vez por cada posición.
\[2^{3}=8\]
\end{minipage}\par\medskip
\noindent\begin{minipage}{\linewidth}\textbf{2. Lista completa, sin repeticiones}\\[2pt]Ordenamos las listas de bits desde todo ceros hasta todo unos. Cada vector corresponde a una elección distinta de sus coordenadas.
\par\smallskip {\small\renewcommand{\arraystretch}{1.3}\begin{tabular}{>{\raggedright\arraybackslash}p{5.85cm}>{\raggedright\arraybackslash}p{5.85cm}}\hline\rowcolor{black!7} \textbf{Vector} & \textbf{Vector}\\\hline (0, 0, 0) & (0, 0, 1)\\ (0, 1, 0) & (0, 1, 1)\\ (1, 0, 0) & (1, 0, 1)\\ (1, 1, 0) & (1, 1, 1)\\ \hline\end{tabular}}\par\smallskip
\end{minipage}\par\medskip
\noindent\begin{minipage}{\linewidth}\textbf{3. Por qué no falta ningún vector}\\[2pt]Para cualquier vector, su primera coordenada es 0 o 1; dentro de cada grupo hacemos lo mismo con la siguiente, hasta completar las n posiciones. Este procedimiento recorre todas las elecciones una sola vez.
\end{minipage}\par\medskip
\respuesta{Hay 8 vectores distintos.}
\noindent\begin{minipage}{\linewidth}\centering\includegraphics[width=.96\linewidth,height=6.8cm,keepaspectratio]{../recursos/grafico_1_2-T_10.png}\\{\small Lectura visual complementaria; la demostración está en las operaciones.}\end{minipage}
\normalsize
\newpage\fancyfoot[L]{\small Sección 1.2 | Ejercicio T.11 | Paso a paso}\hypertarget{1.2-T.11}{}
\pdfbookmark[1]{1.2 / T.11}{bm1.2-T.11}
\titulo{T.11 | Todos los 4-vectores binarios}
\textbf{Enunciado}\par Haga una lista de todos los posibles 4-vectores binarios. ¿Cuántos hay?
\par\medskip\large
\noindent\begin{minipage}{\linewidth}\textbf{1. Contamos las elecciones por coordenada}\\[2pt]Cada coordenada puede elegirse de dos maneras: 0 o 1. Las elecciones son independientes. Multiplicamos 2 por sí mismo una vez por cada posición.
\[2^{4}=16\]
\end{minipage}\par\medskip
\noindent\begin{minipage}{\linewidth}\textbf{2. Lista completa, sin repeticiones}\\[2pt]Ordenamos las listas de bits desde todo ceros hasta todo unos. Cada vector corresponde a una elección distinta de sus coordenadas.
\par\smallskip {\small\renewcommand{\arraystretch}{1.3}\begin{tabular}{>{\raggedright\arraybackslash}p{5.85cm}>{\raggedright\arraybackslash}p{5.85cm}}\hline\rowcolor{black!7} \textbf{Vector} & \textbf{Vector}\\\hline (0, 0, 0, 0) & (0, 0, 0, 1)\\ (0, 0, 1, 0) & (0, 0, 1, 1)\\ (0, 1, 0, 0) & (0, 1, 0, 1)\\ (0, 1, 1, 0) & (0, 1, 1, 1)\\ (1, 0, 0, 0) & (1, 0, 0, 1)\\ (1, 0, 1, 0) & (1, 0, 1, 1)\\ (1, 1, 0, 0) & (1, 1, 0, 1)\\ (1, 1, 1, 0) & (1, 1, 1, 1)\\ \hline\end{tabular}}\par\smallskip
\end{minipage}\par\medskip
\noindent\begin{minipage}{\linewidth}\textbf{3. Por qué no falta ningún vector}\\[2pt]Para cualquier vector, su primera coordenada es 0 o 1; dentro de cada grupo hacemos lo mismo con la siguiente, hasta completar las n posiciones. Este procedimiento recorre todas las elecciones una sola vez.
\end{minipage}\par\medskip
\respuesta{Hay 16 vectores distintos.}
\noindent\begin{minipage}{\linewidth}\centering\includegraphics[width=.96\linewidth,height=6.8cm,keepaspectratio]{../recursos/grafico_1_2-T_11.png}\\{\small Lectura visual complementaria; la demostración está en las operaciones.}\end{minipage}
\normalsize
\newpage\fancyfoot[L]{\small Sección 1.2 | Ejercicio T.12 | Paso a paso}\hypertarget{1.2-T.12}{}
\pdfbookmark[1]{1.2 / T.12}{bm1.2-T.12}
\titulo{T.12 | Contar vectores binarios de longitud n}
\textbf{Enunciado}\par ¿Cuántos 5-vectores binarios hay? ¿Cuántos n-vectores binarios existen?
\par\medskip\large
\noindent\begin{minipage}{\linewidth}\textbf{1. Regla del producto}\\[2pt]Cada una de las n coordenadas tiene 2 elecciones, independientes de las demás. Por el principio multiplicativo:
\[2\cdot2\cdots2=2^n\]
\end{minipage}\par\medskip
\noindent\begin{minipage}{\linewidth}\textbf{2. Caso de cinco coordenadas}\\[2pt]Sustituimos n = 5 en la fórmula.
\[2^5=32\]
\end{minipage}\par\medskip
\noindent\begin{minipage}{\linewidth}\textbf{3. Justificación por inducción alternativa}\\[2pt]Para longitud 1 hay 2 vectores. Si hay 2\textsuperscript{n} vectores de longitud n, a cada uno se le puede añadir 0 o 1; eso produce 2·2\textsuperscript{n} = 2\textsuperscript{n}\textsuperscript{+}\textsuperscript{1} vectores distintos. Así la fórmula vale para toda longitud positiva. Para longitud cero, si se admite el vector vacío, hay uno, coherente con 2\textsuperscript{0} = 1.
\end{minipage}\par\medskip
\respuesta{32 vectores de longitud 5; en general, 2\textasciicircum{}n.}
\noindent\begin{minipage}{\linewidth}\centering\includegraphics[width=.96\linewidth,height=6.8cm,keepaspectratio]{../recursos/grafico_1_2-T_12.png}\\{\small Lectura visual complementaria; la demostración está en las operaciones.}\end{minipage}
\normalsize
\newpage\fancyfoot[L]{\small Sección 1.2 | Ejercicio T.13 | Paso a paso}\hypertarget{1.2-T.13}{}
\pdfbookmark[1]{1.2 / T.13}{bm1.2-T.13}
\titulo{T.13 | Todas las matrices binarias de orden 2}
\textbf{Enunciado}\par Haga una lista de todas las posibles matrices binarias de 2 por 2. ¿Cuántas hay?
\par\medskip\large
\noindent\begin{minipage}{\linewidth}\textbf{1. Pasamos de cuatro casillas a cuatro bits}\\[2pt]Leemos las casillas por filas: (a11, a12, a21, a22). Cada elección de cuatro bits produce una matriz, y cada matriz produce esa única lista. Esta correspondencia no repite ni omite matrices.
\[2^4=16\]
\end{minipage}\par\medskip
\noindent\begin{minipage}{\linewidth}\textbf{2. Matrices 1 a 4}\\[2pt]Cada bloque muestra cuatro elecciones consecutivas, desde 0000 hasta 1111. Las entradas se acomodan primero en la fila superior y después en la inferior.
\par\smallskip \begin{minipage}[t]{.47\linewidth}\centering \textbf{1}\[\begin{bmatrix}0 & 0\\0 & 0\end{bmatrix}\]\end{minipage}\hfill\begin{minipage}[t]{.47\linewidth}\centering \textbf{2}\[\begin{bmatrix}0 & 0\\0 & 1\end{bmatrix}\]\end{minipage}\par\smallskip
\par\smallskip \begin{minipage}[t]{.47\linewidth}\centering \textbf{3}\[\begin{bmatrix}0 & 0\\1 & 0\end{bmatrix}\]\end{minipage}\hfill\begin{minipage}[t]{.47\linewidth}\centering \textbf{4}\[\begin{bmatrix}0 & 0\\1 & 1\end{bmatrix}\]\end{minipage}\par\smallskip
\end{minipage}\par\medskip
\noindent\begin{minipage}{\linewidth}\textbf{3. Matrices 5 a 8}\\[2pt]Cada bloque muestra cuatro elecciones consecutivas, desde 0000 hasta 1111. Las entradas se acomodan primero en la fila superior y después en la inferior.
\par\smallskip \begin{minipage}[t]{.47\linewidth}\centering \textbf{5}\[\begin{bmatrix}0 & 1\\0 & 0\end{bmatrix}\]\end{minipage}\hfill\begin{minipage}[t]{.47\linewidth}\centering \textbf{6}\[\begin{bmatrix}0 & 1\\0 & 1\end{bmatrix}\]\end{minipage}\par\smallskip
\par\smallskip \begin{minipage}[t]{.47\linewidth}\centering \textbf{7}\[\begin{bmatrix}0 & 1\\1 & 0\end{bmatrix}\]\end{minipage}\hfill\begin{minipage}[t]{.47\linewidth}\centering \textbf{8}\[\begin{bmatrix}0 & 1\\1 & 1\end{bmatrix}\]\end{minipage}\par\smallskip
\end{minipage}\par\medskip
\noindent\begin{minipage}{\linewidth}\textbf{4. Matrices 9 a 12}\\[2pt]Cada bloque muestra cuatro elecciones consecutivas, desde 0000 hasta 1111. Las entradas se acomodan primero en la fila superior y después en la inferior.
\par\smallskip \begin{minipage}[t]{.47\linewidth}\centering \textbf{9}\[\begin{bmatrix}1 & 0\\0 & 0\end{bmatrix}\]\end{minipage}\hfill\begin{minipage}[t]{.47\linewidth}\centering \textbf{10}\[\begin{bmatrix}1 & 0\\0 & 1\end{bmatrix}\]\end{minipage}\par\smallskip
\par\smallskip \begin{minipage}[t]{.47\linewidth}\centering \textbf{11}\[\begin{bmatrix}1 & 0\\1 & 0\end{bmatrix}\]\end{minipage}\hfill\begin{minipage}[t]{.47\linewidth}\centering \textbf{12}\[\begin{bmatrix}1 & 0\\1 & 1\end{bmatrix}\]\end{minipage}\par\smallskip
\end{minipage}\par\medskip
\noindent\begin{minipage}{\linewidth}\textbf{5. Matrices 13 a 16}\\[2pt]Cada bloque muestra cuatro elecciones consecutivas, desde 0000 hasta 1111. Las entradas se acomodan primero en la fila superior y después en la inferior.
\par\smallskip \begin{minipage}[t]{.47\linewidth}\centering \textbf{13}\[\begin{bmatrix}1 & 1\\0 & 0\end{bmatrix}\]\end{minipage}\hfill\begin{minipage}[t]{.47\linewidth}\centering \textbf{14}\[\begin{bmatrix}1 & 1\\0 & 1\end{bmatrix}\]\end{minipage}\par\smallskip
\par\smallskip \begin{minipage}[t]{.47\linewidth}\centering \textbf{15}\[\begin{bmatrix}1 & 1\\1 & 0\end{bmatrix}\]\end{minipage}\hfill\begin{minipage}[t]{.47\linewidth}\centering \textbf{16}\[\begin{bmatrix}1 & 1\\1 & 1\end{bmatrix}\]\end{minipage}\par\smallskip
\end{minipage}\par\medskip
\noindent\begin{minipage}{\linewidth}\textbf{6. Conclusión}\\[2pt]Se han listado las 16 matrices posibles. Contamos casillas, no filas: aquí hay cuatro elecciones independientes.
\end{minipage}\par\medskip
\respuesta{Hay 16 matrices binarias de 2 por 2.}
\noindent\begin{minipage}{\linewidth}\centering\includegraphics[width=.96\linewidth,height=6.8cm,keepaspectratio]{../recursos/grafico_1_2-T_13.png}\\{\small Lectura visual complementaria; la demostración está en las operaciones.}\end{minipage}
\normalsize
\newpage\fancyfoot[L]{\small Sección 1.2 | Ejercicio T.14 | Paso a paso}\hypertarget{1.2-T.14}{}
\pdfbookmark[1]{1.2 / T.14}{bm1.2-T.14}
\titulo{T.14 | Contar matrices binarias de orden 3}
\textbf{Enunciado}\par ¿Cuántas matrices binarias de 3 por 3 existen?
\par\medskip\large
\noindent\begin{minipage}{\linewidth}\textbf{1. Contamos las posiciones, no solo las filas}\\[2pt]Una matriz de n filas y n columnas tiene n\textsuperscript{2} entradas. Cada entrada es un bit, elegido independientemente de las demás.
\[n\cdot n=n^2\]
\end{minipage}\par\medskip
\noindent\begin{minipage}{\linewidth}\textbf{2. Aplicamos el principio multiplicativo}\\[2pt]Hay dos posibilidades por entrada. Al multiplicarlas una vez por cada casilla, obtenemos:
\[2^{3\cdot3}=2^9=512\]
\end{minipage}\par\medskip
\noindent\begin{minipage}{\linewidth}\textbf{3. Por qué la cuenta es exacta}\\[2pt]Leer las casillas por filas identifica cada matriz con una lista binaria de longitud n\textsuperscript{2}. La lectura es reversible: de la lista recuperamos una única matriz. Por eso contar matrices equivale a contar esas listas.
\end{minipage}\par\medskip
\respuesta{512 matrices.}
\noindent\begin{minipage}{\linewidth}\centering\includegraphics[width=.96\linewidth,height=6.8cm,keepaspectratio]{../recursos/grafico_1_2-T_14.png}\\{\small Lectura visual complementaria; la demostración está en las operaciones.}\end{minipage}
\normalsize
\newpage\fancyfoot[L]{\small Sección 1.2 | Ejercicio T.15 | Paso a paso}\hypertarget{1.2-T.15}{}
\pdfbookmark[1]{1.2 / T.15}{bm1.2-T.15}
\titulo{T.15 | Contar matrices binarias de orden n}
\textbf{Enunciado}\par ¿Cuántas matrices binarias de n por n existen?
\par\medskip\large
\noindent\begin{minipage}{\linewidth}\textbf{1. Contamos las posiciones, no solo las filas}\\[2pt]Una matriz de n filas y n columnas tiene n\textsuperscript{2} entradas. Cada entrada es un bit, elegido independientemente de las demás.
\[n\cdot n=n^2\]
\end{minipage}\par\medskip
\noindent\begin{minipage}{\linewidth}\textbf{2. Aplicamos el principio multiplicativo}\\[2pt]Hay dos posibilidades por entrada. Al multiplicarlas una vez por cada casilla, obtenemos:
\[2^{n^2}\]
\end{minipage}\par\medskip
\noindent\begin{minipage}{\linewidth}\textbf{3. Por qué la cuenta es exacta}\\[2pt]Leer las casillas por filas identifica cada matriz con una lista binaria de longitud n\textsuperscript{2}. La lectura es reversible: de la lista recuperamos una única matriz. Por eso contar matrices equivale a contar esas listas.
\end{minipage}\par\medskip
\respuesta{2\textasciicircum{}(n\textsuperscript{2}) matrices.}
\noindent\begin{minipage}{\linewidth}\centering\includegraphics[width=.96\linewidth,height=6.8cm,keepaspectratio]{../recursos/grafico_1_2-T_15.png}\\{\small Lectura visual complementaria; la demostración está en las operaciones.}\end{minipage}
\normalsize
\newpage\fancyfoot[L]{\small Sección 1.2 | Ejercicio T.16 | Paso a paso}\hypertarget{1.2-T.16}{}
\pdfbookmark[1]{1.2 / T.16}{bm1.2-T.16}
\titulo{T.16 | Interruptores: obtener todos OFF}
\textbf{Enunciado}\par Represente OFF con 0 y ON con 1. Sea A = [[ON,ON,OFF],[OFF,ON,OFF],[OFF,ON,ON]]. Determine B en estados ON/OFF para que A+B tenga todas sus entradas OFF.
\par\smallskip \textbf{A}\[\begin{bmatrix}1 & 1 & 0\\0 & 1 & 0\\0 & 1 & 1\end{bmatrix}\]
\par\medskip\large
\noindent\begin{minipage}{\linewidth}\textbf{1. Convertimos palabras en bits}\\[2pt]Cada ON se escribe 1 y cada OFF se escribe 0. Las operaciones que siguen son binarias, módulo 2.
\par\smallskip \textbf{A}\[\begin{bmatrix}1 & 1 & 0\\0 & 1 & 0\\0 & 1 & 1\end{bmatrix}\]
\end{minipage}\par\medskip
\noindent\begin{minipage}{\linewidth}\textbf{2. Elegimos el bit que necesitamos en cada posición}\\[2pt]Para obtener 0, los dos bits deben ser iguales: B coincide con A.
\par\smallskip \textbf{B en bits}\[\begin{bmatrix}1 & 1 & 0\\0 & 1 & 0\\0 & 1 & 1\end{bmatrix}\]
\par\smallskip\textbf{Qué aprendimos.} B en ON/OFF: [['ON', 'ON', 'OFF'], ['OFF', 'ON', 'OFF'], ['OFF', 'ON', 'ON']]
\end{minipage}\par\medskip
\noindent\begin{minipage}{\linewidth}\textbf{3. Comprobación de todas las posiciones}\\[2pt]Sumamos A y B módulo 2. El resultado tiene el mismo estado objetivo en las nueve casillas.
\par\smallskip \textbf{A + B, módulo 2}\[\begin{bmatrix}0 & 0 & 0\\0 & 0 & 0\\0 & 0 & 0\end{bmatrix}\]
\end{minipage}\par\medskip
\respuesta{B = [['ON', 'ON', 'OFF'], ['OFF', 'ON', 'OFF'], ['OFF', 'ON', 'ON']]}
\noindent\begin{minipage}{\linewidth}\centering\includegraphics[width=.96\linewidth,height=6.8cm,keepaspectratio]{../recursos/grafico_1_2-T_16.png}\\{\small Lectura visual complementaria; la demostración está en las operaciones.}\end{minipage}
\normalsize
\newpage\fancyfoot[L]{\small Sección 1.2 | Ejercicio T.17 | Paso a paso}\hypertarget{1.2-T.17}{}
\pdfbookmark[1]{1.2 / T.17}{bm1.2-T.17}
\titulo{T.17 | Interruptores: obtener todos ON}
\textbf{Enunciado}\par Represente OFF con 0 y ON con 1. Sea A = [[ON,ON,OFF],[OFF,ON,OFF],[OFF,ON,ON]]. Determine B en estados ON/OFF para que A+B tenga todas sus entradas ON.
\par\smallskip \textbf{A}\[\begin{bmatrix}1 & 1 & 0\\0 & 1 & 0\\0 & 1 & 1\end{bmatrix}\]
\par\medskip\large
\noindent\begin{minipage}{\linewidth}\textbf{1. Convertimos palabras en bits}\\[2pt]Cada ON se escribe 1 y cada OFF se escribe 0. Las operaciones que siguen son binarias, módulo 2.
\par\smallskip \textbf{A}\[\begin{bmatrix}1 & 1 & 0\\0 & 1 & 0\\0 & 1 & 1\end{bmatrix}\]
\end{minipage}\par\medskip
\noindent\begin{minipage}{\linewidth}\textbf{2. Elegimos el bit que necesitamos en cada posición}\\[2pt]Para obtener 1, los dos bits deben ser distintos: B es el complemento de A.
\par\smallskip \textbf{B en bits}\[\begin{bmatrix}0 & 0 & 1\\1 & 0 & 1\\1 & 0 & 0\end{bmatrix}\]
\par\smallskip\textbf{Qué aprendimos.} B en ON/OFF: [['OFF', 'OFF', 'ON'], ['ON', 'OFF', 'ON'], ['ON', 'OFF', 'OFF']]
\end{minipage}\par\medskip
\noindent\begin{minipage}{\linewidth}\textbf{3. Comprobación de todas las posiciones}\\[2pt]Sumamos A y B módulo 2. El resultado tiene el mismo estado objetivo en las nueve casillas.
\par\smallskip \textbf{A + B, módulo 2}\[\begin{bmatrix}1 & 1 & 1\\1 & 1 & 1\\1 & 1 & 1\end{bmatrix}\]
\end{minipage}\par\medskip
\respuesta{B = [['OFF', 'OFF', 'ON'], ['ON', 'OFF', 'ON'], ['ON', 'OFF', 'OFF']]}
\noindent\begin{minipage}{\linewidth}\centering\includegraphics[width=.96\linewidth,height=6.8cm,keepaspectratio]{../recursos/grafico_1_2-T_17.png}\\{\small Lectura visual complementaria; la demostración está en las operaciones.}\end{minipage}
\normalsize
\newpage\fancyfoot[L]{\small Sección 1.2 | Ejercicio T.18 | Paso a paso}\hypertarget{1.2-T.18}{}
\pdfbookmark[1]{1.2 / T.18}{bm1.2-T.18}
\titulo{T.18 | Invertir todos los interruptores con la misma máscara}
\textbf{Enunciado}\par Sea A = [[1,0],[0,1],[1,1]], con 0 = apagado y 1 = encendido. (a) Halle B para que A+B invierta cada estado. (b) Para C = [[1,1],[0,0],[1,0]], ¿la misma B invierte C? Verifique. (c) Generalice a una matriz binaria m por n.
\par\smallskip \begin{minipage}[t]{.47\linewidth}\centering \textbf{A}\[\begin{bmatrix}1 & 0\\0 & 1\\1 & 1\end{bmatrix}\]\end{minipage}\hfill\begin{minipage}[t]{.47\linewidth}\centering \textbf{C}\[\begin{bmatrix}1 & 1\\0 & 0\\1 & 0\end{bmatrix}\]\end{minipage}\par\smallskip
\par\medskip\large
\noindent\begin{minipage}{\linewidth}\textbf{1. Inciso (a): qué operación invierte un bit}\\[2pt]En módulo 2, sumar 1 cambia ambos estados posibles: 0 pasa a 1 y 1 pasa a 0. Por tanto se necesita un 1 en cada posición de B.
\[0+1=1,\quad1+1=0\quad(\mathrm{mod}\ 2)\]
\par\smallskip \begin{minipage}[t]{.47\linewidth}\centering \textbf{B}\[\begin{bmatrix}1 & 1\\1 & 1\\1 & 1\end{bmatrix}\]\end{minipage}\hfill\begin{minipage}[t]{.47\linewidth}\centering \textbf{A + B, módulo 2}\[\begin{bmatrix}0 & 1\\1 & 0\\0 & 0\end{bmatrix}\]\end{minipage}\par\smallskip
\end{minipage}\par\medskip
\noindent\begin{minipage}{\linewidth}\textbf{2. Inciso (b): probamos exactamente la misma B}\\[2pt]No cambiamos B al cambiar el estado inicial. Sumamos la matriz de unos a C y reducimos las seis entradas módulo 2.
\par\smallskip \textbf{C + B, módulo 2}\[\begin{bmatrix}0 & 0\\1 & 1\\0 & 1\end{bmatrix}\]
\par\smallskip\textbf{Qué aprendimos.} Sí. Cada cero de C se convirtió en uno y cada uno en cero.
\end{minipage}\par\medskip
\noindent\begin{minipage}{\linewidth}\textbf{3. Inciso (c): demostración para cualquier tamaño}\\[2pt]Sean A binaria de m por n e i, j índices arbitrarios. Elegimos bij = 1 en todas las casillas. Si aij = 0, aij+bij = 1; si aij = 1, aij+bij = 0 módulo 2. Son los únicos dos casos posibles. Así todos los estados se invierten, para cualquier A y cualquier tamaño.
\[B=J_{m\times n},\quad b_{ij}=1\]
\par\smallskip\textbf{Qué aprendimos.} La máscara que invierte todos los estados es la matriz de unos; no el complemento del estado inicial. El complemento sería el resultado de aplicar la máscara.
\end{minipage}\par\medskip
\respuesta{B es la matriz de unos. Para C, el resultado es [[0,0],[1,1],[0,1]]. En general, bij = 1.}
\noindent\begin{minipage}{\linewidth}\centering\includegraphics[width=.96\linewidth,height=6.8cm,keepaspectratio]{../recursos/grafico_1_2-T_18.png}\\{\small Lectura visual complementaria; la demostración está en las operaciones.}\end{minipage}
\normalsize
\newpage\fancyfoot[L]{\small Sección 1.2 | Ejercicio ML.1 | Paso a paso}\hypertarget{1.2-ML.1}{}
\pdfbookmark[1]{1.2 / ML.1}{bm1.2-ML.1}
\titulo{ML.1 | MATLAB: entradas y formatos}
\textbf{Enunciado}\par Introduzca A = [[5,1,2],[-3,0,1],[2,4,1]] y B = [[4·2,2/3],[1/201,5-8.2],[0.00001,(9+4)/3]]. (a) Despliegue a23, b23 y b12. (b) Despliegue fila1(A), columna3(A) y fila2(B). (c) Despliegue B con format long, compare los valores del inciso (a) y restablezca format short.
\par\smallskip \begin{minipage}[t]{.47\linewidth}\centering \textbf{A}\[\begin{bmatrix}5 & 1 & 2\\-3 & 0 & 1\\2 & 4 & 1\end{bmatrix}\]\end{minipage}\hfill\begin{minipage}[t]{.47\linewidth}\centering \textbf{B}\[\begin{bmatrix}8 & \frac{2}{3}\\\frac{1}{201} & - \frac{16}{5}\\\frac{1}{100000} & \frac{13}{3}\end{bmatrix}\]\end{minipage}\par\smallskip
\par\medskip
\noindent\begin{minipage}{\linewidth}\textbf{1. Introducir matrices en MATLAB}\\[2pt]Los espacios separan columnas y los puntos y coma separan filas. El acceso A(i,j) usa el primer índice como fila y el segundo como columna. MATLAB empieza a contar en 1.
\par\smallskip \begin{minipage}[t]{.47\linewidth}\centering \textbf{A}\[\begin{bmatrix}5 & 1 & 2\\-3 & 0 & 1\\2 & 4 & 1\end{bmatrix}\]\end{minipage}\hfill\begin{minipage}[t]{.47\linewidth}\centering \textbf{B en valores exactos}\[\begin{bmatrix}8 & \frac{2}{3}\\\frac{1}{201} & - \frac{16}{5}\\\frac{1}{100000} & \frac{13}{3}\end{bmatrix}\]\end{minipage}\par\smallskip
\end{minipage}\par\medskip
\noindent\begin{minipage}{\linewidth}\textbf{2. Inciso (a): dos accesos válidos y uno inexistente}\\[2pt]A(2,3) es 1. B(1,2) es 2/3. El enunciado pide b23, pero B es de tamaño 3 por 2: no existe su columna 3. Escribimos esa observación y no sustituimos el índice por otro sin autorización.
\[a_{23}=1,\quad b_{12}=\frac{2}{3}\]
\par\smallskip\textbf{Qué aprendimos.} B(2,3) produciría un error de índice en MATLAB. El archivo entregado detecta este caso y continúa con los demás incisos.
\end{minipage}\par\medskip
\noindent\begin{minipage}{\linewidth}\textbf{3. Inciso (b): seleccionar una fila o una columna}\\[2pt]El símbolo : significa todas las posiciones de ese índice. A(1,:) toma la primera fila; A(:,3), la tercera columna; B(2,:), la segunda fila.
\par\smallskip \begin{minipage}[t]{.47\linewidth}\centering \textbf{A(1,:)}\[\begin{bmatrix}5 & 1 & 2\end{bmatrix}\]\end{minipage}\hfill\begin{minipage}[t]{.47\linewidth}\centering \textbf{A(:,3)}\[\begin{bmatrix}2\\1\\1\end{bmatrix}\]\end{minipage}\par\smallskip
\par\smallskip \begin{minipage}[t]{.47\linewidth}\centering \textbf{B(2,:)}\[\begin{bmatrix}\frac{1}{201} & - \frac{16}{5}\end{bmatrix}\]\end{minipage}\par\smallskip
\end{minipage}\par\medskip
\noindent\begin{minipage}{\linewidth}\textbf{4. Inciso (c): formato y precisión no son lo mismo}\\[2pt]format long muestra más cifras que format short. Ninguno cambia los valores almacenados ni recupera una fracción exacta perdida por representación en coma flotante.
\par\smallskip {\small\renewcommand{\arraystretch}{1.3}\begin{tabular}{>{\raggedright\arraybackslash}p{3.74cm}>{\raggedright\arraybackslash}p{3.74cm}>{\raggedright\arraybackslash}p{3.74cm}}\hline\rowcolor{black!7} \textbf{Entrada} & \textbf{Long: ejemplo de salida} & \textbf{Short: ejemplo de salida}\\\hline b12 & 0.666666666666667 & 0.6667\\ b21 & 0.004975124378109 & 0.0050\\ b31 & 0.000010000000000 & 0.0000\\ \hline\end{tabular}}\par\smallskip
\par\smallskip\textbf{Qué aprendimos.} b31 no se vuelve cero: solo se ve redondeado como 0.0000 en esta presentación. Al final ejecutamos format short.
\end{minipage}\par\medskip
\noindent\begin{minipage}{\linewidth}\textbf{5. Ejecutarlo y abrir la simulación en MATLAB}\\[2pt]Abre la carpeta matlab del proyecto como carpeta actual de MATLAB. Ejecuta ml\_1 para ver los comandos y resultados. Ejecuta simular\_capitulo\_1\_2 y elige ML.1 para abrir el laboratorio visual. Mantén bingen.m y binadd.m en esa misma carpeta. No necesitas pegar comandos de Python en MATLAB.
\end{minipage}\par\medskip
\par\medskip\textbf{Código MATLAB: ml\_1.m}\par{\small\begin{verbatim}
%% ML.1 - Entradas, filas, columnas y formatos
clear; clc;
A = [5 1 2; -3 0 1; 2 4 1];
B = [4*2 2/3; 1/201 5-8.2; 0.00001 (9+4)/3];
%% (a) Primer indice: fila; segundo indice: columna
fprintf('a23 = %.15g\n', A(2,3));
% El enunciado pide b23, pero B solo tiene DOS columnas.
if size(B,2) < 3
    disp('b23 no existe: B tiene 3 filas y 2 columnas.');
else
    disp(B(2,3));
end
fprintf('b12 = %.15g\n', B(1,2));
%% (b) Dos puntos significa: tomar todas las posiciones
fila1_A = A(1,:)
columna3_A = A(:,3)
fila2_B = B(2,:)
%% (c) Cambiar el formato no cambia los valores almacenados
B_guardada = B;
format long
B
format short
B
assert(isequal(B, B_guardada));
fprintf('b12 con 15 decimales: %.15f\n', B(1,2));
fprintf('b12 con 4 decimales: %.4f\n', B(1,2));
format short

\end{verbatim}}
\respuesta{a23 = 1; b23 no existe; b12 = 2/3. Filas y columna indicadas en la solución. format solo cambia la presentación.}
\noindent\begin{minipage}{\linewidth}\centering\includegraphics[width=.96\linewidth,height=6.8cm,keepaspectratio]{../recursos/grafico_1_2-ML_1.png}\\{\small Lectura visual complementaria; la demostración está en las operaciones.}\end{minipage}
\normalsize
\newpage\fancyfoot[L]{\small Sección 1.2 | Ejercicio ML.2 | Paso a paso}\hypertarget{1.2-ML.2}{}
\pdfbookmark[1]{1.2 / ML.2}{bm1.2-ML.2}
\titulo{ML.2 | MATLAB: matriz de Hilbert}
\textbf{Enunciado}\par Escriba H = hilb(5);. (a) Determine el tamaño de H. (b) Despliegue H. (c) Despliegue H como números racionales. (d) Extraiga las tres primeras columnas. (e) Extraiga las dos últimas filas.
\par\smallskip \textbf{H}\[\begin{bmatrix}1 & \frac{1}{2} & \frac{1}{3} & \frac{1}{4} & \frac{1}{5}\\\frac{1}{2} & \frac{1}{3} & \frac{1}{4} & \frac{1}{5} & \frac{1}{6}\\\frac{1}{3} & \frac{1}{4} & \frac{1}{5} & \frac{1}{6} & \frac{1}{7}\\\frac{1}{4} & \frac{1}{5} & \frac{1}{6} & \frac{1}{7} & \frac{1}{8}\\\frac{1}{5} & \frac{1}{6} & \frac{1}{7} & \frac{1}{8} & \frac{1}{9}\end{bmatrix}\]
\par\medskip
\noindent\begin{minipage}{\linewidth}\textbf{1. Construimos H y entendemos el punto y coma}\\[2pt]hilb(5) crea una matriz de 5 por 5. El punto y coma suprime la salida automática del comando; no impide crear la variable. Cada entrada se define mediante la suma de sus índices.
\[H_{ij}=\frac{1}{i+j-1}\quad(i,j=1,\ldots,5)\]
\[H_{23}=\frac{1}{2+3-1}=\frac14\]
\end{minipage}\par\medskip
\noindent\begin{minipage}{\linewidth}\textbf{2. Incisos (a), (b) y (c)}\\[2pt]size(H) devuelve [5 5]. Escribir H muestra su contenido. format rat seguido de H presenta razones de enteros; la matriz numérica de MATLAB sigue almacenada como double.
\par\smallskip \textbf{H en fracciones}\[\begin{bmatrix}1 & \frac{1}{2} & \frac{1}{3} & \frac{1}{4} & \frac{1}{5}\\\frac{1}{2} & \frac{1}{3} & \frac{1}{4} & \frac{1}{5} & \frac{1}{6}\\\frac{1}{3} & \frac{1}{4} & \frac{1}{5} & \frac{1}{6} & \frac{1}{7}\\\frac{1}{4} & \frac{1}{5} & \frac{1}{6} & \frac{1}{7} & \frac{1}{8}\\\frac{1}{5} & \frac{1}{6} & \frac{1}{7} & \frac{1}{8} & \frac{1}{9}\end{bmatrix}\]
\end{minipage}\par\medskip
\noindent\begin{minipage}{\linewidth}\textbf{3. Inciso (d): tres primeras columnas}\\[2pt]H(:,1:3) conserva todas las filas y las columnas 1, 2 y 3. El tamaño resultante es 5 por 3.
\par\smallskip \textbf{H(:,1:3)}\[\begin{bmatrix}1 & \frac{1}{2} & \frac{1}{3}\\\frac{1}{2} & \frac{1}{3} & \frac{1}{4}\\\frac{1}{3} & \frac{1}{4} & \frac{1}{5}\\\frac{1}{4} & \frac{1}{5} & \frac{1}{6}\\\frac{1}{5} & \frac{1}{6} & \frac{1}{7}\end{bmatrix}\]
\end{minipage}\par\medskip
\noindent\begin{minipage}{\linewidth}\textbf{4. Inciso (e): dos últimas filas}\\[2pt]H(end-1:end,:) toma las filas 4 y 5 y todas las columnas. end representa el último índice del eje donde aparece. El resultado es 2 por 5.
\par\smallskip \textbf{H(end-1:end,:)}\[\begin{bmatrix}\frac{1}{4} & \frac{1}{5} & \frac{1}{6} & \frac{1}{7} & \frac{1}{8}\\\frac{1}{5} & \frac{1}{6} & \frac{1}{7} & \frac{1}{8} & \frac{1}{9}\end{bmatrix}\]
\end{minipage}\par\medskip
\noindent\begin{minipage}{\linewidth}\textbf{5. Ejecutarlo y abrir la simulación en MATLAB}\\[2pt]Abre la carpeta matlab del proyecto como carpeta actual de MATLAB. Ejecuta ml\_2 para ver los comandos y resultados. Ejecuta simular\_capitulo\_1\_2 y elige ML.2 para abrir el laboratorio visual. Mantén bingen.m y binadd.m en esa misma carpeta. No necesitas pegar comandos de Python en MATLAB.
\end{minipage}\par\medskip
\par\medskip\textbf{Código MATLAB: ml\_2.m}\par{\small\begin{verbatim}
%% ML.2 - Matriz de Hilbert
clear; clc;
H = hilb(5); % El punto y coma evita la salida automatica.
%% (a) Tamano
size(H)
%% (b) Contenido decimal
format short
H
%% (c) Presentacion racional
% Es presentacion racional de numeros double; no exige sym.
format rat
H
%% (d) Todas las filas y las tres primeras columnas
primeras_tres_columnas = H(:,1:3)
%% (e) Dos ultimas filas y todas las columnas
ultimas_dos_filas = H(end-1:end,:)
format short
% Ejemplo de la regla exacta que define las entradas:
fprintf('H(2,3) = 1/(2+3-1) = %.4f\n', H(2,3));

\end{verbatim}}
\respuesta{size(H) = [5 5]; H(:,1:3) tiene tamaño 5 por 3; H(end-1:end,:) tiene tamaño 2 por 5.}
\noindent\begin{minipage}{\linewidth}\centering\includegraphics[width=.96\linewidth,height=6.8cm,keepaspectratio]{../recursos/grafico_1_2-ML_2.png}\\{\small Lectura visual complementaria; la demostración está en las operaciones.}\end{minipage}
\normalsize
\newpage\fancyfoot[L]{\small Sección 1.2 | Ejercicio ML.3 | Paso a paso}\hypertarget{1.2-ML.3}{}
\pdfbookmark[1]{1.2 / ML.3}{bm1.2-ML.3}
\titulo{ML.3 | MATLAB: bingen para vectores binarios}
\textbf{Enunciado}\par Utilice bingen para resolver los ejercicios T.10 y T.11: liste todos los vectores binarios de longitudes 3 y 4.
\par\medskip
\noindent\begin{minipage}{\linewidth}\textbf{1. La función auxiliar que se entrega}\\[2pt]bingen.m es una implementación incluida en el proyecto. Enumera los números de 0 a 2\textsuperscript{n}-1, escribe cada uno con n bits y devuelve un vector por fila. Guarda bingen.m junto a ml\_3.m.
\end{minipage}\par\medskip
\noindent\begin{minipage}{\linewidth}\textbf{2. Longitud 3}\\[2pt]V3 = bingen(3) produce 8 filas y 3 columnas; cada fila es un vector distinto. No contamos las columnas como si fueran vectores.
\par\smallskip \textbf{V3}\[\begin{bmatrix}0 & 0 & 0\\0 & 0 & 1\\0 & 1 & 0\\0 & 1 & 1\\1 & 0 & 0\\1 & 0 & 1\\1 & 1 & 0\\1 & 1 & 1\end{bmatrix}\]
\end{minipage}\par\medskip
\noindent\begin{minipage}{\linewidth}\textbf{3. Longitud 4}\\[2pt]V4 = bingen(4) produce 16 filas y 4 columnas. Sus dos bloques se muestran para facilitar la lectura.
\par\smallskip \begin{minipage}[t]{.47\linewidth}\centering \textbf{V4: filas 1 a 8}\[\begin{bmatrix}0 & 0 & 0 & 0\\0 & 0 & 0 & 1\\0 & 0 & 1 & 0\\0 & 0 & 1 & 1\\0 & 1 & 0 & 0\\0 & 1 & 0 & 1\\0 & 1 & 1 & 0\\0 & 1 & 1 & 1\end{bmatrix}\]\end{minipage}\hfill\begin{minipage}[t]{.47\linewidth}\centering \textbf{V4: filas 9 a 16}\[\begin{bmatrix}1 & 0 & 0 & 0\\1 & 0 & 0 & 1\\1 & 0 & 1 & 0\\1 & 0 & 1 & 1\\1 & 1 & 0 & 0\\1 & 1 & 0 & 1\\1 & 1 & 1 & 0\\1 & 1 & 1 & 1\end{bmatrix}\]\end{minipage}\par\smallskip
\end{minipage}\par\medskip
\noindent\begin{minipage}{\linewidth}\textbf{4. Comprobación}\\[2pt]size(V3,1) da 8 y size(V4,1) da 16. unique(V4,'rows') conserva las 16 filas: ninguna se repite.
\[2^3=8,\quad2^4=16\]
\end{minipage}\par\medskip
\noindent\begin{minipage}{\linewidth}\textbf{5. Ejecutarlo y abrir la simulación en MATLAB}\\[2pt]Abre la carpeta matlab del proyecto como carpeta actual de MATLAB. Ejecuta ml\_3 para ver los comandos y resultados. Ejecuta simular\_capitulo\_1\_2 y elige ML.3 para abrir el laboratorio visual. Mantén bingen.m y binadd.m en esa misma carpeta. No necesitas pegar comandos de Python en MATLAB.
\end{minipage}\par\medskip
\par\medskip\textbf{Código MATLAB: ml\_3.m}\par{\small\begin{verbatim}
%% ML.3 - Todos los vectores binarios de longitudes 3 y 4
clear; clc;
% bingen.m se entrega en esta misma carpeta.
V3 = bingen(3)
V4 = bingen(4)
fprintf('Vectores de longitud 3: %d\n', size(V3,1));
fprintf('Vectores de longitud 4: %d\n', size(V4,1));
assert(isequal(size(V3),[8 3]));
assert(isequal(size(V4),[16 4]));
assert(size(unique(V4,'rows'),1)==16);

\end{verbatim}}
\respuesta{8 vectores de longitud 3 y 16 de longitud 4.}
\noindent\begin{minipage}{\linewidth}\centering\includegraphics[width=.96\linewidth,height=6.8cm,keepaspectratio]{../recursos/grafico_1_2-ML_3.png}\\{\small Lectura visual complementaria; la demostración está en las operaciones.}\end{minipage}
\normalsize
\newpage\fancyfoot[L]{\small Sección 1.2 | Ejercicio ML.4 | Paso a paso}\hypertarget{1.2-ML.4}{}
\pdfbookmark[1]{1.2 / ML.4}{bm1.2-ML.4}
\titulo{ML.4 | MATLAB: enumerar matrices binarias}
\textbf{Enunciado}\par Utilice bingen para resolver el ejercicio T.13. Una matriz n por n tiene el mismo número de entradas que un n\textsuperscript{2}-vector.
\par\medskip
\noindent\begin{minipage}{\linewidth}\textbf{1. Cuatro casillas independientes}\\[2pt]Una matriz 2 por 2 tiene cuatro entradas. Por tanto usamos bingen(4), que genera 16 filas de cuatro bits; bingen(2) sería insuficiente.
\[2\cdot2=4,\quad2^4=16\]
\end{minipage}\par\medskip
\noindent\begin{minipage}{\linewidth}\textbf{2. Convertimos cada vector en una matriz}\\[2pt]Elegimos el orden por filas [a11 a12 a21 a22]. MATLAB reshape rellena por columnas. Por eso usamos reshape(V(k,:),2,2).' para conservar el orden por filas.
\par\smallskip \begin{minipage}[t]{.47\linewidth}\centering \textbf{Vector de ejemplo}\[\begin{bmatrix}0 & 1 & 0 & 0\end{bmatrix}\]\end{minipage}\hfill\begin{minipage}[t]{.47\linewidth}\centering \textbf{Matriz correspondiente}\[\begin{bmatrix}0 & 1\\0 & 0\end{bmatrix}\]\end{minipage}\par\smallskip
\end{minipage}\par\medskip
\noindent\begin{minipage}{\linewidth}\textbf{3. Recorremos los 16 casos}\\[2pt]El bucle for usa una fila diferente de V en cada iteración. El resultado se guarda en M(:,:,k). Como la conversión es reversible y las filas de V no se repiten, las 16 matrices son distintas.
\end{minipage}\par\medskip
\noindent\begin{minipage}{\linewidth}\textbf{4. Ejecutarlo y abrir la simulación en MATLAB}\\[2pt]Abre la carpeta matlab del proyecto como carpeta actual de MATLAB. Ejecuta ml\_4 para ver los comandos y resultados. Ejecuta simular\_capitulo\_1\_2 y elige ML.4 para abrir el laboratorio visual. Mantén bingen.m y binadd.m en esa misma carpeta. No necesitas pegar comandos de Python en MATLAB.
\end{minipage}\par\medskip
\par\medskip\textbf{Código MATLAB: ml\_4.m}\par{\small\begin{verbatim}
%% ML.4 - Las 16 matrices binarias de 2 por 2
clear; clc;
V = bingen(4); % Cuatro entradas independientes, no dos.
M = zeros(2,2,16);
for k = 1:size(V,1)
    % V(k,:) = [a11 a12 a21 a22], lectura por filas.
    % reshape rellena por columnas: transponemos para conservar filas.
    M(:,:,k) = reshape(V(k,:),2,2).';
    fprintf('Matriz %d de 16\n', k);
    disp(M(:,:,k));
end
assert(size(M,3)==16);

\end{verbatim}}
\respuesta{Se generan las 16 matrices binarias de 2 por 2, sin omisiones ni repeticiones.}
\noindent\begin{minipage}{\linewidth}\centering\includegraphics[width=.96\linewidth,height=6.8cm,keepaspectratio]{../recursos/grafico_1_2-ML_4.png}\\{\small Lectura visual complementaria; la demostración está en las operaciones.}\end{minipage}
\normalsize
\newpage\fancyfoot[L]{\small Sección 1.2 | Ejercicio ML.5 | Paso a paso}\hypertarget{1.2-ML.5}{}
\pdfbookmark[1]{1.2 / ML.5}{bm1.2-ML.5}
\titulo{ML.5 | MATLAB: binadd y aritmética módulo 2}
\textbf{Enunciado}\par Resuelva el ejercicio 11 utilizando binadd. Se entrega binadd.m para sumar matrices binarias módulo 2.
\par\smallskip \begin{minipage}[t]{.47\linewidth}\centering \textbf{A}\[\begin{bmatrix}1 & 0 & 1\\1 & 1 & 0\\0 & 1 & 1\end{bmatrix}\]\end{minipage}\hfill\begin{minipage}[t]{.47\linewidth}\centering \textbf{B}\[\begin{bmatrix}0 & 1 & 1\\1 & 0 & 1\\1 & 1 & 0\end{bmatrix}\]\end{minipage}\par\smallskip
\par\smallskip \begin{minipage}[t]{.47\linewidth}\centering \textbf{C}\[\begin{bmatrix}1 & 1 & 0\\0 & 1 & 1\\1 & 0 & 1\end{bmatrix}\]\end{minipage}\par\smallskip
\textbf{Expresiones del enunciado}\par
a. \(A+B\)\par
b. \(B+C\)\par
c. \(A+B+C\)\par
d. \(A+C^T\)\par
e. \(B-C\)\par
\par\medskip
\noindent\begin{minipage}{\linewidth}\textbf{1. Antes de calcular: filas, columnas y operaciones}\\[2pt]El tamaño de una matriz se escribe filas por columnas. Para sumar o restar, los dos tamaños deben coincidir exactamente. Multiplicar por un número afecta a cada entrada. Transponer intercambia las posiciones: la entrada de fila i y columna j pasa a la fila j y columna i.
\[(M^T)_{ij}=m_{ji}\]
\par\smallskip\textbf{Qué aprendimos.} Los paréntesis indican qué operación hay que realizar primero. Una suma de tamaños distintos no se vuelve válida por transponerla después.
\end{minipage}\par\medskip
\noindent\begin{minipage}{\linewidth}\textbf{2. Regla de la aritmética binaria}\\[2pt]Aquí los resultados solo pueden ser 0 o 1. Calculamos el resto al dividir la cuenta entre 2; por eso 1 + 1 se convierte en 0. En esta aritmética, restar coincide con sumar: el opuesto de 1 es 1.
\[0+0=0,\quad 0+1=1,\quad1+0=1,\quad1+1=0\quad(\mathrm{mod}\ 2)\]
\par\smallskip\textbf{Qué aprendimos.} Estas reglas se aplican en los ejercicios 11–15 y en los teóricos binarios, no en las operaciones reales de los ejercicios 4–10.
\end{minipage}\par\medskip
\noindent\begin{minipage}{\linewidth}\textbf{3. Inciso a: A+B: qué vamos a calcular}\\[2pt]Identificamos la operación y comprobamos sus dimensiones antes de efectuar las cuentas. Las matrices originales permanecen escritas en el enunciado.
\[A+B\]
\end{minipage}\par\medskip
\noindent\begin{minipage}{\linewidth}\textbf{4. Operamos las entradas que ocupan la misma posición}\\[2pt]La fila y la columna nos dicen qué valores combinar. Hacemos una cuenta por casilla y después reunimos todas las casillas en la matriz resultado. En cada casilla reducimos el resultado módulo 2: un número par da 0 y un número impar da 1.
\par\smallskip {\small\renewcommand{\arraystretch}{1.3}\begin{tabular}{>{\raggedright\arraybackslash}p{3.74cm}>{\raggedright\arraybackslash}p{3.74cm}>{\raggedright\arraybackslash}p{3.74cm}}\hline\rowcolor{black!7} \textbf{Casilla} & \textbf{Cuenta} & \textbf{Resultado}\\\hline Fila 1, columna 1 & $\mathrm{mod}((1)+(0),2)$ & $1$\\ Fila 1, columna 2 & $\mathrm{mod}((0)+(1),2)$ & $1$\\ Fila 1, columna 3 & $\mathrm{mod}((1)+(1),2)$ & $0$\\ Fila 2, columna 1 & $\mathrm{mod}((1)+(1),2)$ & $0$\\ Fila 2, columna 2 & $\mathrm{mod}((1)+(0),2)$ & $1$\\ Fila 2, columna 3 & $\mathrm{mod}((0)+(1),2)$ & $1$\\ Fila 3, columna 1 & $\mathrm{mod}((0)+(1),2)$ & $1$\\ Fila 3, columna 2 & $\mathrm{mod}((1)+(1),2)$ & $0$\\ Fila 3, columna 3 & $\mathrm{mod}((1)+(0),2)$ & $1$\\ \hline\end{tabular}}\par\smallskip
\par\smallskip \textbf{Resultado de esta operación}\[\begin{bmatrix}1 & 1 & 0\\0 & 1 & 1\\1 & 0 & 1\end{bmatrix}\]
\par\smallskip\textbf{Qué aprendimos.} El resultado tiene 3 filas y 3 columnas.
\end{minipage}\par\medskip
\noindent\begin{minipage}{\linewidth}\textbf{5. Resultado del inciso a}\\[2pt]Reunimos el resultado. Comprueba que cada casilla procede de la operación indicada, sin cambiar su posición.
\par\smallskip \textbf{Resultado}\[\begin{bmatrix}1 & 1 & 0\\0 & 1 & 1\\1 & 0 & 1\end{bmatrix}\]
\end{minipage}\par\medskip
\noindent\begin{minipage}{\linewidth}\textbf{6. Inciso b: B+C: qué vamos a calcular}\\[2pt]Identificamos la operación y comprobamos sus dimensiones antes de efectuar las cuentas. Las matrices originales permanecen escritas en el enunciado.
\[B+C\]
\end{minipage}\par\medskip
\noindent\begin{minipage}{\linewidth}\textbf{7. Operamos las entradas que ocupan la misma posición}\\[2pt]La fila y la columna nos dicen qué valores combinar. Hacemos una cuenta por casilla y después reunimos todas las casillas en la matriz resultado. En cada casilla reducimos el resultado módulo 2: un número par da 0 y un número impar da 1.
\par\smallskip {\small\renewcommand{\arraystretch}{1.3}\begin{tabular}{>{\raggedright\arraybackslash}p{3.74cm}>{\raggedright\arraybackslash}p{3.74cm}>{\raggedright\arraybackslash}p{3.74cm}}\hline\rowcolor{black!7} \textbf{Casilla} & \textbf{Cuenta} & \textbf{Resultado}\\\hline Fila 1, columna 1 & $\mathrm{mod}((0)+(1),2)$ & $1$\\ Fila 1, columna 2 & $\mathrm{mod}((1)+(1),2)$ & $0$\\ Fila 1, columna 3 & $\mathrm{mod}((1)+(0),2)$ & $1$\\ Fila 2, columna 1 & $\mathrm{mod}((1)+(0),2)$ & $1$\\ Fila 2, columna 2 & $\mathrm{mod}((0)+(1),2)$ & $1$\\ Fila 2, columna 3 & $\mathrm{mod}((1)+(1),2)$ & $0$\\ Fila 3, columna 1 & $\mathrm{mod}((1)+(1),2)$ & $0$\\ Fila 3, columna 2 & $\mathrm{mod}((1)+(0),2)$ & $1$\\ Fila 3, columna 3 & $\mathrm{mod}((0)+(1),2)$ & $1$\\ \hline\end{tabular}}\par\smallskip
\par\smallskip \textbf{Resultado de esta operación}\[\begin{bmatrix}1 & 0 & 1\\1 & 1 & 0\\0 & 1 & 1\end{bmatrix}\]
\par\smallskip\textbf{Qué aprendimos.} El resultado tiene 3 filas y 3 columnas.
\end{minipage}\par\medskip
\noindent\begin{minipage}{\linewidth}\textbf{8. Resultado del inciso b}\\[2pt]Reunimos el resultado. Comprueba que cada casilla procede de la operación indicada, sin cambiar su posición.
\par\smallskip \textbf{Resultado}\[\begin{bmatrix}1 & 0 & 1\\1 & 1 & 0\\0 & 1 & 1\end{bmatrix}\]
\end{minipage}\par\medskip
\noindent\begin{minipage}{\linewidth}\textbf{9. Inciso c: A+B+C: qué vamos a calcular}\\[2pt]Identificamos la operación y comprobamos sus dimensiones antes de efectuar las cuentas. Las matrices originales permanecen escritas en el enunciado.
\[A+B+C\]
\end{minipage}\par\medskip
\noindent\begin{minipage}{\linewidth}\textbf{10. Operamos las entradas que ocupan la misma posición}\\[2pt]La fila y la columna nos dicen qué valores combinar. Hacemos una cuenta por casilla y después reunimos todas las casillas en la matriz resultado. En cada casilla reducimos el resultado módulo 2: un número par da 0 y un número impar da 1.
\par\smallskip {\small\renewcommand{\arraystretch}{1.3}\begin{tabular}{>{\raggedright\arraybackslash}p{3.74cm}>{\raggedright\arraybackslash}p{3.74cm}>{\raggedright\arraybackslash}p{3.74cm}}\hline\rowcolor{black!7} \textbf{Casilla} & \textbf{Cuenta} & \textbf{Resultado}\\\hline Fila 1, columna 1 & $\mathrm{mod}((1)+(0),2)$ & $1$\\ Fila 1, columna 2 & $\mathrm{mod}((0)+(1),2)$ & $1$\\ Fila 1, columna 3 & $\mathrm{mod}((1)+(1),2)$ & $0$\\ Fila 2, columna 1 & $\mathrm{mod}((1)+(1),2)$ & $0$\\ Fila 2, columna 2 & $\mathrm{mod}((1)+(0),2)$ & $1$\\ Fila 2, columna 3 & $\mathrm{mod}((0)+(1),2)$ & $1$\\ Fila 3, columna 1 & $\mathrm{mod}((0)+(1),2)$ & $1$\\ Fila 3, columna 2 & $\mathrm{mod}((1)+(1),2)$ & $0$\\ Fila 3, columna 3 & $\mathrm{mod}((1)+(0),2)$ & $1$\\ \hline\end{tabular}}\par\smallskip
\par\smallskip \textbf{Resultado de esta operación}\[\begin{bmatrix}1 & 1 & 0\\0 & 1 & 1\\1 & 0 & 1\end{bmatrix}\]
\par\smallskip\textbf{Qué aprendimos.} El resultado tiene 3 filas y 3 columnas.
\end{minipage}\par\medskip
\noindent\begin{minipage}{\linewidth}\textbf{11. Operamos las entradas que ocupan la misma posición}\\[2pt]La fila y la columna nos dicen qué valores combinar. Hacemos una cuenta por casilla y después reunimos todas las casillas en la matriz resultado. En cada casilla reducimos el resultado módulo 2: un número par da 0 y un número impar da 1.
\par\smallskip {\small\renewcommand{\arraystretch}{1.3}\begin{tabular}{>{\raggedright\arraybackslash}p{3.74cm}>{\raggedright\arraybackslash}p{3.74cm}>{\raggedright\arraybackslash}p{3.74cm}}\hline\rowcolor{black!7} \textbf{Casilla} & \textbf{Cuenta} & \textbf{Resultado}\\\hline Fila 1, columna 1 & $\mathrm{mod}((1)+(1),2)$ & $0$\\ Fila 1, columna 2 & $\mathrm{mod}((1)+(1),2)$ & $0$\\ Fila 1, columna 3 & $\mathrm{mod}((0)+(0),2)$ & $0$\\ Fila 2, columna 1 & $\mathrm{mod}((0)+(0),2)$ & $0$\\ Fila 2, columna 2 & $\mathrm{mod}((1)+(1),2)$ & $0$\\ Fila 2, columna 3 & $\mathrm{mod}((1)+(1),2)$ & $0$\\ Fila 3, columna 1 & $\mathrm{mod}((1)+(1),2)$ & $0$\\ Fila 3, columna 2 & $\mathrm{mod}((0)+(0),2)$ & $0$\\ Fila 3, columna 3 & $\mathrm{mod}((1)+(1),2)$ & $0$\\ \hline\end{tabular}}\par\smallskip
\par\smallskip \textbf{Resultado de esta operación}\[\begin{bmatrix}0 & 0 & 0\\0 & 0 & 0\\0 & 0 & 0\end{bmatrix}\]
\par\smallskip\textbf{Qué aprendimos.} El resultado tiene 3 filas y 3 columnas.
\end{minipage}\par\medskip
\noindent\begin{minipage}{\linewidth}\textbf{12. Resultado del inciso c}\\[2pt]Reunimos el resultado. Comprueba que cada casilla procede de la operación indicada, sin cambiar su posición.
\par\smallskip \textbf{Resultado}\[\begin{bmatrix}0 & 0 & 0\\0 & 0 & 0\\0 & 0 & 0\end{bmatrix}\]
\end{minipage}\par\medskip
\noindent\begin{minipage}{\linewidth}\textbf{13. Inciso d: A+C\textasciicircum{}T: qué vamos a calcular}\\[2pt]Identificamos la operación y comprobamos sus dimensiones antes de efectuar las cuentas. Las matrices originales permanecen escritas en el enunciado.
\[A+C^T\]
\end{minipage}\par\medskip
\noindent\begin{minipage}{\linewidth}\textbf{14. Transponemos: las filas se convierten en columnas}\\[2pt]No cambiamos los valores. Solo intercambiamos los dos índices de cada entrada. Una matriz de m filas y n columnas produce otra de n filas y m columnas.
\par\smallskip {\small\renewcommand{\arraystretch}{1.3}\begin{tabular}{>{\raggedright\arraybackslash}p{3.74cm}>{\raggedright\arraybackslash}p{3.74cm}>{\raggedright\arraybackslash}p{3.74cm}}\hline\rowcolor{black!7} \textbf{Posición antes} & \textbf{Posición después} & \textbf{Valor conservado}\\\hline (1, 1) & (1, 1) & $1$\\ (1, 2) & (2, 1) & $1$\\ (1, 3) & (3, 1) & $0$\\ (2, 1) & (1, 2) & $0$\\ (2, 2) & (2, 2) & $1$\\ (2, 3) & (3, 2) & $1$\\ (3, 1) & (1, 3) & $1$\\ (3, 2) & (2, 3) & $0$\\ (3, 3) & (3, 3) & $1$\\ \hline\end{tabular}}\par\smallskip
\par\smallskip \textbf{Resultado de esta operación}\[\begin{bmatrix}1 & 0 & 1\\1 & 1 & 0\\0 & 1 & 1\end{bmatrix}\]
\par\smallskip\textbf{Qué aprendimos.} El resultado tiene 3 filas y 3 columnas.
\end{minipage}\par\medskip
\noindent\begin{minipage}{\linewidth}\textbf{15. Operamos las entradas que ocupan la misma posición}\\[2pt]La fila y la columna nos dicen qué valores combinar. Hacemos una cuenta por casilla y después reunimos todas las casillas en la matriz resultado. En cada casilla reducimos el resultado módulo 2: un número par da 0 y un número impar da 1.
\par\smallskip {\small\renewcommand{\arraystretch}{1.3}\begin{tabular}{>{\raggedright\arraybackslash}p{3.74cm}>{\raggedright\arraybackslash}p{3.74cm}>{\raggedright\arraybackslash}p{3.74cm}}\hline\rowcolor{black!7} \textbf{Casilla} & \textbf{Cuenta} & \textbf{Resultado}\\\hline Fila 1, columna 1 & $\mathrm{mod}((1)+(1),2)$ & $0$\\ Fila 1, columna 2 & $\mathrm{mod}((0)+(0),2)$ & $0$\\ Fila 1, columna 3 & $\mathrm{mod}((1)+(1),2)$ & $0$\\ Fila 2, columna 1 & $\mathrm{mod}((1)+(1),2)$ & $0$\\ Fila 2, columna 2 & $\mathrm{mod}((1)+(1),2)$ & $0$\\ Fila 2, columna 3 & $\mathrm{mod}((0)+(0),2)$ & $0$\\ Fila 3, columna 1 & $\mathrm{mod}((0)+(0),2)$ & $0$\\ Fila 3, columna 2 & $\mathrm{mod}((1)+(1),2)$ & $0$\\ Fila 3, columna 3 & $\mathrm{mod}((1)+(1),2)$ & $0$\\ \hline\end{tabular}}\par\smallskip
\par\smallskip \textbf{Resultado de esta operación}\[\begin{bmatrix}0 & 0 & 0\\0 & 0 & 0\\0 & 0 & 0\end{bmatrix}\]
\par\smallskip\textbf{Qué aprendimos.} El resultado tiene 3 filas y 3 columnas.
\end{minipage}\par\medskip
\noindent\begin{minipage}{\linewidth}\textbf{16. Resultado del inciso d}\\[2pt]Reunimos el resultado. Comprueba que cada casilla procede de la operación indicada, sin cambiar su posición.
\par\smallskip \textbf{Resultado}\[\begin{bmatrix}0 & 0 & 0\\0 & 0 & 0\\0 & 0 & 0\end{bmatrix}\]
\end{minipage}\par\medskip
\noindent\begin{minipage}{\linewidth}\textbf{17. Inciso e: B-C: qué vamos a calcular}\\[2pt]Identificamos la operación y comprobamos sus dimensiones antes de efectuar las cuentas. Las matrices originales permanecen escritas en el enunciado.
\[B-C\]
\end{minipage}\par\medskip
\noindent\begin{minipage}{\linewidth}\textbf{18. Operamos las entradas que ocupan la misma posición}\\[2pt]La fila y la columna nos dicen qué valores combinar. Hacemos una cuenta por casilla y después reunimos todas las casillas en la matriz resultado. En cada casilla reducimos el resultado módulo 2: un número par da 0 y un número impar da 1.
\par\smallskip {\small\renewcommand{\arraystretch}{1.3}\begin{tabular}{>{\raggedright\arraybackslash}p{3.74cm}>{\raggedright\arraybackslash}p{3.74cm}>{\raggedright\arraybackslash}p{3.74cm}}\hline\rowcolor{black!7} \textbf{Casilla} & \textbf{Cuenta} & \textbf{Resultado}\\\hline Fila 1, columna 1 & $\mathrm{mod}((0)-(1),2)$ & $1$\\ Fila 1, columna 2 & $\mathrm{mod}((1)-(1),2)$ & $0$\\ Fila 1, columna 3 & $\mathrm{mod}((1)-(0),2)$ & $1$\\ Fila 2, columna 1 & $\mathrm{mod}((1)-(0),2)$ & $1$\\ Fila 2, columna 2 & $\mathrm{mod}((0)-(1),2)$ & $1$\\ Fila 2, columna 3 & $\mathrm{mod}((1)-(1),2)$ & $0$\\ Fila 3, columna 1 & $\mathrm{mod}((1)-(1),2)$ & $0$\\ Fila 3, columna 2 & $\mathrm{mod}((1)-(0),2)$ & $1$\\ Fila 3, columna 3 & $\mathrm{mod}((0)-(1),2)$ & $1$\\ \hline\end{tabular}}\par\smallskip
\par\smallskip \textbf{Resultado de esta operación}\[\begin{bmatrix}1 & 0 & 1\\1 & 1 & 0\\0 & 1 & 1\end{bmatrix}\]
\par\smallskip\textbf{Qué aprendimos.} El resultado tiene 3 filas y 3 columnas.
\end{minipage}\par\medskip
\noindent\begin{minipage}{\linewidth}\textbf{19. Resultado del inciso e}\\[2pt]Reunimos el resultado. Comprueba que cada casilla procede de la operación indicada, sin cambiar su posición.
\par\smallskip \textbf{Resultado}\[\begin{bmatrix}1 & 0 & 1\\1 & 1 & 0\\0 & 1 & 1\end{bmatrix}\]
\end{minipage}\par\medskip
\noindent\begin{minipage}{\linewidth}\textbf{20. Ejecutarlo y abrir la simulación en MATLAB}\\[2pt]Abre la carpeta matlab del proyecto como carpeta actual de MATLAB. Ejecuta ml\_5 para ver los comandos y resultados. Ejecuta simular\_capitulo\_1\_2 y elige ML.5 para abrir el laboratorio visual. Mantén bingen.m y binadd.m en esa misma carpeta. No necesitas pegar comandos de Python en MATLAB.
\end{minipage}\par\medskip
\par\medskip\textbf{Código MATLAB: ml\_5.m}\par{\small\begin{verbatim}
%% ML.5 - Ejercicio 11 con binadd (modulo 2)
clear; clc;
A = [1 0 1; 1 1 0; 0 1 1];
B = [0 1 1; 1 0 1; 1 1 0];
C = [1 1 0; 0 1 1; 1 0 1];
%% (a) A+B
Ra = binadd(A,B)
%% (b) B+C
Rb = binadd(B,C)
%% (c) A+B+C
Rc = binadd(binadd(A,B),C)
%% (d) A+transpuesta(C), sin conjugacion
Rd = binadd(A,C.')
%% (e) En modulo 2, restar y sumar coinciden
Re = binadd(B,C)
assert(isequal(Re, mod(B-C,2)));
assert(isequal(Ra, [1 1 0;0 1 1;1 0 1]));
assert(isequal(Rb, [1 0 1;1 1 0;0 1 1]));
assert(isequal(Rc, zeros(3)));
format short

\end{verbatim}}
\respuesta{a: A+B: [[1, 1, 0], [0, 1, 1], [1, 0, 1]]; b: B+C: [[1, 0, 1], [1, 1, 0], [0, 1, 1]]; c: A+B+C: [[0, 0, 0], [0, 0, 0], [0, 0, 0]]; d: A+C\textasciicircum{}T: [[0, 0, 0], [0, 0, 0], [0, 0, 0]]; e: B-C: [[1, 0, 1], [1, 1, 0], [0, 1, 1]].}
\noindent\begin{minipage}{\linewidth}\centering\includegraphics[width=.96\linewidth,height=6.8cm,keepaspectratio]{../recursos/grafico_1_2-ML_5.png}\\{\small Lectura visual complementaria; la demostración está en las operaciones.}\end{minipage}
\normalsize
\newpage\pdfbookmark[0]{Respuestas para repasar}{resumen}\fancyfoot[L]{\small Colección | Respuestas por sección}\titulo{Respuestas para repasar}\begin{multicols}{2}
\textbf{1.1 / 1.} x = 4; y = 2.\par\smallskip
\textbf{1.1 / 2.} x = 1; y = 2; z = -2.\par\smallskip
\textbf{1.1 / 3.} x = -4; y = 2; z = 10.\par\smallskip
\textbf{1.1 / 4.} Sin solución.\par\smallskip
\textbf{1.1 / 5.} x = 2; y = -1; z = -2.\par\smallskip
\textbf{1.1 / 6.} x = 10*t + 13; y = -8*t - 8; z = t. Parámetro real.\par\smallskip
\textbf{1.1 / 7.} x = -20; y = t/4 + 8; z = t. Parámetro real.\par\smallskip
\textbf{1.1 / 8.} Sin solución.\par\smallskip
\textbf{1.1 / 9.} Sin solución.\par\smallskip
\textbf{1.1 / 10.} x = 2; y = -1.\par\smallskip
\textbf{1.1 / 11.} x = 5; y = 1.\par\smallskip
\textbf{1.1 / 12.} Sin solución.\par\smallskip
\textbf{1.1 / 13.} Sin solución.\par\smallskip
\textbf{1.1 / 14.} x = -1; y = 2; z = -2.\par\smallskip
\textbf{1.1 / 15.} (a) t = 10: infinitas soluciones. (b) Cualquier t distinto de 10; por ejemplo 0. (c) Infinitos valores.\par\smallskip
\textbf{1.1 / 16.} (a), (b), (c) y (d): sí son soluciones; la suma es (-1, 1, 1) y su triple es (-3, 3, 3).\par\smallskip
\textbf{1.1 / 17.} (x, y, z) = (1, 1, 4).\par\smallskip
\textbf{1.1 / 18.} (x, y, z) = (2, 1, 0).\par\smallskip
\textbf{1.1 / 19.} Sí: r = -3.\par\smallskip
\textbf{1.1 / 20.} No existe tal valor de r.\par\smallskip
\textbf{1.1 / 21.} (a) Un punto. (b) Ningún punto. (c) Infinitos puntos sobre una recta.\par\smallskip
\textbf{1.1 / 22.} (a) Cero. (b) Infinitos, sobre una recta. (c) Cero.\par\smallskip
\textbf{1.1 / 23.} x = 20; y = 20. Se producen 20 toneladas sin azufre y 20 con azufre. Mezcladora: 100 + 80 = 180 min; refinación: 80 + 40 = 120 min.\par\smallskip
\textbf{1.1 / 24.} x = 3/2; y = 5/2. Se producen 1.5 toneladas de regular y 2.5 toneladas de especial al día. En A: 3 + 5 = 8 h. En B: 7.5 + 7.5 = 15 h. Las toneladas pueden ser fraccionarias.\par\smallskip
\textbf{1.1 / 25.} Alimento A: 16/5 = 3.2 onzas. Alimento B: 21/5 = 4.2 onzas. Alimento C: 2 onzas.\par\smallskip
\textbf{1.1 / 26.} x = 20; y = 40 - t; z = t. Parámetro real. El revelador de 2 minutos debe producirse en 20 toneladas al día. Los otros dos deben sumar 40 toneladas: y = 40 - t, z = t, con 0 <= t <= 40. Hay infinitos planes no negativos.\par\smallskip
\textbf{1.1 / 27.} a = 5, b = -3 y c = -7. La parábola es p(x) = 5x\textsuperscript{2} - 3x - 7.\par\smallskip
\textbf{1.1 / 28.} Primer fideicomiso: \$7,000. Segundo: \$14,000. Tercero: \$3,000. Interés anual total: \$2,210.\par\smallskip
\textbf{1.1 / T.1.} Intercambiar dos ecuaciones conserva exactamente el conjunto solución.\par\smallskip
\textbf{1.1 / T.2.} Los sistemas son equivalentes porque la operación tiene inversa: dividir por la misma constante no nula.\par\smallskip
\textbf{1.1 / T.3.} Sumar a una ecuación un múltiplo de otra distinta, conservando esta última, mantiene exactamente las soluciones.\par\smallskip
\textbf{1.1 / T.4.} Sí. (x, y) = (0, 0) siempre es solución, para cualquier a, b, c y d.\par\smallskip
\textbf{1.2 / 1.} a12 = -3, a22 = -5, a23 = 4; b11 = 4, b31 = 5; c13 = 2, c31 = 6, c33 = -1.\par\smallskip
\textbf{1.2 / 2.} a = 3, b = 1, c = 8, d = -2.\par\smallskip
\textbf{1.2 / 3.} a = 0, b = 2, c = 1, d = 2.\par\smallskip
\textbf{1.2 / 4.} a1: C+E: [[5, -5, 8], [4, 2, 9], [5, 3, 4]]; a2: E+C: [[5, -5, 8], [4, 2, 9], [5, 3, 4]]; b: A+B: no definida; c: D-F: [[7, -7], [0, 1]]; d: -3C+5O: [[-9, 3, -9], [-12, -3, -15], [-6, -3, -9]]; e: 2C-3E: [[0, 10, -9], [8, -1, -2], [-5, -4, 3]]; f: 2B+F: no definida.\par\smallskip
\textbf{1.2 / 5.} a: 3D+2F: [[1, 4], [10, 18]]; b1: 3(2A): [[6, 12, 18], [12, 6, 24]]; b2: 6A: [[6, 12, 18], [12, 6, 24]]; c1: 3A+2A: [[5, 10, 15], [10, 5, 20]]; c2: 5A: [[5, 10, 15], [10, 5, 20]]; d1: 2(D+F): [[-2, 6], [8, 14]]; d2: 2D+2F: [[-2, 6], [8, 14]]; e1: (2+3)D: [[15, -10], [10, 20]]; e2: 2D+3D: [[15, -10], [10, 20]]; f: 3(B+D): no definida.\par\smallskip
\textbf{1.2 / 6.} a1: A\textasciicircum{}T: [[1, 2], [2, 1], [3, 4]]; a2: (A\textasciicircum{}T)\textasciicircum{}T: [[1, 2, 3], [2, 1, 4]]; b1: (C+E)\textasciicircum{}T: [[5, 4, 5], [-5, 2, 3], [8, 9, 4]]; b2: C\textasciicircum{}T+E\textasciicircum{}T: [[5, 4, 5], [-5, 2, 3], [8, 9, 4]]; c: (2D+3F)\textasciicircum{}T: [[-6, 10], [11, 17]]; d: D-D\textasciicircum{}T: [[0, -4], [4, 0]]; e: 2A\textasciicircum{}T+B: [[3, 4], [6, 3], [9, 10]]; f: (3D-2F)\textasciicircum{}T: [[17, 2], [-16, 6]].\par\smallskip
\textbf{1.2 / 7.} a: (2A)\textasciicircum{}T: [[2, 4], [4, 2], [6, 8]]; b: (A-B)\textasciicircum{}T: no definida; c: (3B\textasciicircum{}T-2A)\textasciicircum{}T: [[1, -4], [2, 1], [3, -2]]; d: (3A\textasciicircum{}T-5B)\textasciicircum{}T: [[-2, -4, -6], [6, -2, 2]]; e1: (-A)\textasciicircum{}T: [[-1, -2], [-2, -1], [-3, -4]]; e2: -(A\textasciicircum{}T): [[-1, -2], [-2, -1], [-3, -4]]; f: (C+E+F\textasciicircum{}T)\textasciicircum{}T: no definida.\par\smallskip
\textbf{1.2 / 8.} Sí: 2I + J.\par\smallskip
\textbf{1.2 / 9.} No: toda combinación tiene entrada (1, 2) igual a 0 y la matriz objetivo tiene 1.\par\smallskip
\textbf{1.2 / 10.} [[$\lambda$-1, -2, -3], [-6, $\lambda$+2, -3], [-5, -2, $\lambda$-4]].\par\smallskip
\textbf{1.2 / 11.} a: A+B: [[1, 1, 0], [0, 1, 1], [1, 0, 1]]; b: B+C: [[1, 0, 1], [1, 1, 0], [0, 1, 1]]; c: A+B+C: [[0, 0, 0], [0, 0, 0], [0, 0, 0]]; d: A+C\textasciicircum{}T: [[0, 0, 0], [0, 0, 0], [0, 0, 0]]; e: B-C: [[1, 0, 1], [1, 1, 0], [0, 1, 1]].\par\smallskip
\textbf{1.2 / 12.} a: A+B: [[0, 0], [1, 1]]; b: C+D: [[1, 1], [1, 0]]; c: A+B+(C+D)\textasciicircum{}T: [[1, 1], [0, 1]]; d: C-B: [[0, 1], [0, 1]]; e: A-B+C-D: [[1, 1], [0, 1]].\par\smallskip
\textbf{1.2 / 13.} B = [[1,0],[0,0]]; C = [[0,1],[1,1]].\par\smallskip
\textbf{1.2 / 14.} [0, 0, 0, 0]\par\smallskip
\textbf{1.2 / 15.} [1, 0, 1, 0]\par\smallskip
\textbf{1.2 / T.1.} Ambas son diagonales, siempre que los dos órdenes coincidan.\par\smallskip
\textbf{1.2 / T.2.} Ambas son escalares, siempre que los dos órdenes coincidan.\par\smallskip
\textbf{1.2 / T.3.} (a) [[0,b-c,c-e],[c-b,0,0],[e-c,0,0]]. (b) y (c) [[2a,b+c,c+e],[b+c,2d,2e],[c+e,2e,2f]].\par\smallskip
\textbf{1.2 / T.4.} kA = O implica k = 0 o A = O.\par\smallskip
\textbf{1.2 / T.5.} La suma y resta conservan cada tipo triangular; ser superior e inferior a la vez equivale a ser diagonal.\par\smallskip
\textbf{1.2 / T.6.} Superior $\to$ transpuesta inferior; inferior $\to$ transpuesta superior.\par\smallskip
\textbf{1.2 / T.7.} Todas son cero.\par\smallskip
\textbf{1.2 / T.8.} x + 0 = x para todo n-vector real x.\par\smallskip
\textbf{1.2 / T.9.} Hay 4 vectores distintos.\par\smallskip
\textbf{1.2 / T.10.} Hay 8 vectores distintos.\par\smallskip
\textbf{1.2 / T.11.} Hay 16 vectores distintos.\par\smallskip
\textbf{1.2 / T.12.} 32 vectores de longitud 5; en general, 2\textasciicircum{}n.\par\smallskip
\textbf{1.2 / T.13.} Hay 16 matrices binarias de 2 por 2.\par\smallskip
\textbf{1.2 / T.14.} 512 matrices.\par\smallskip
\textbf{1.2 / T.15.} 2\textasciicircum{}(n\textsuperscript{2}) matrices.\par\smallskip
\textbf{1.2 / T.16.} B = [['ON', 'ON', 'OFF'], ['OFF', 'ON', 'OFF'], ['OFF', 'ON', 'ON']]\par\smallskip
\textbf{1.2 / T.17.} B = [['OFF', 'OFF', 'ON'], ['ON', 'OFF', 'ON'], ['ON', 'OFF', 'OFF']]\par\smallskip
\textbf{1.2 / T.18.} B es la matriz de unos. Para C, el resultado es [[0,0],[1,1],[0,1]]. En general, bij = 1.\par\smallskip
\textbf{1.2 / ML.1.} a23 = 1; b23 no existe; b12 = 2/3. Filas y columna indicadas en la solución. format solo cambia la presentación.\par\smallskip
\textbf{1.2 / ML.2.} size(H) = [5 5]; H(:,1:3) tiene tamaño 5 por 3; H(end-1:end,:) tiene tamaño 2 por 5.\par\smallskip
\textbf{1.2 / ML.3.} 8 vectores de longitud 3 y 16 de longitud 4.\par\smallskip
\textbf{1.2 / ML.4.} Se generan las 16 matrices binarias de 2 por 2, sin omisiones ni repeticiones.\par\smallskip
\textbf{1.2 / ML.5.} a: A+B: [[1, 1, 0], [0, 1, 1], [1, 0, 1]]; b: B+C: [[1, 0, 1], [1, 1, 0], [0, 1, 1]]; c: A+B+C: [[0, 0, 0], [0, 0, 0], [0, 0, 0]]; d: A+C\textasciicircum{}T: [[0, 0, 0], [0, 0, 0], [0, 0, 0]]; e: B-C: [[1, 0, 1], [1, 1, 0], [0, 1, 1]].\par\smallskip
\end{multicols}\end{document}